\documentclass[11pt]{article}
%% ======================================================================
%% A proof of the Erdos--Gallai cycle decomposition conjecture.
%% Main file: preamble, front matter and the inputs s1.tex ... s7.tex,
%% formalization.tex (Section 8); bibliography refs.bib (amsplain).
%% Packages: amsmath, amssymb, amsthm, geometry, hyperref only. Every macro
%% is defined here with \newcommand; section files define no macros,
%% environments or counters.
%% ======================================================================
\usepackage{amsmath,amssymb,amsthm}
\usepackage[margin=1in]{geometry}
\usepackage[colorlinks=true,linkcolor=blue,citecolor=blue,urlcolor=blue]{hyperref}
\hypersetup{pdftitle={A proof of the Erd\H{o}s--Gallai cycle decomposition conjecture},pdfauthor={Ryan Coffey}}

\numberwithin{equation}{section}

%% ----------------------------------------------------------------------
%% Theorem environments (one shared counter, numbered within sections)
%% ----------------------------------------------------------------------
\theoremstyle{plain}
\newtheorem{theorem}{Theorem}[section]
\newtheorem{lemma}[theorem]{Lemma}
\newtheorem{proposition}[theorem]{Proposition}
\newtheorem{corollary}[theorem]{Corollary}
\newtheorem{fact}[theorem]{Fact}
\newtheorem{citedresult}[theorem]{Cited result}
\newtheorem{claim}{Claim}[theorem]
\theoremstyle{definition}
\newtheorem{definition}[theorem]{Definition}
\newtheorem{construction}[theorem]{Construction}
\newtheorem{condition}[theorem]{Condition}
\newtheorem{convention}[theorem]{Convention}
\theoremstyle{remark}
\newtheorem{remark}[theorem]{Remark}

%% ----------------------------------------------------------------------
%% Labels with custom text, e.g. \namedlabel{s1:condG1}{\Gc{1}}
%% ----------------------------------------------------------------------
\makeatletter
\newcommand{\namedlabel}[2]{\begingroup\protected@edef\@currentlabel{#2}\phantomsection\label{#1}\endgroup}
\makeatother

%% ----------------------------------------------------------------------
%% Dependency lines and names
%% ----------------------------------------------------------------------
\newcommand{\deps}[1]{\par\smallskip\noindent{\small\textbf{Dependencies:} #1}\par\smallskip}
\newcommand{\EG}{Erd\H{o}s--Gallai}
\newcommand{\BMnames}{Buci\'c--Montgomery}

%% ----------------------------------------------------------------------
%% General mathematics
%% ----------------------------------------------------------------------
\newcommand{\f}{f}                      % f(G), f(F): least number of objects
\newcommand{\defeq}{\mathrel{:=}}
\newcommand{\Prob}{\mathbb{P}}
\newcommand{\Ex}{\mathbb{E}}
\newcommand{\ind}[1]{\mathbf{1}[#1]}
\newcommand{\logs}{\log^{\star}}         % iterated logarithm (base 2)
\newcommand{\Bin}{\mathrm{Bin}}
\newcommand{\Nbr}{N}                     % N_G(U): neighbours outside U
\newcommand{\Ball}[3]{B^{#1}_{#2}(#3)}   % \Ball{\ell}{G-F}{U,V}

%% ----------------------------------------------------------------------
%% Constants and their order (s1:defConstants, s1:condG1--G4)
%% ----------------------------------------------------------------------
\newcommand{\eps}{\varepsilon}          % eps := 2^{-5}
\newcommand{\Dstar}{D_{*}}
\newcommand{\Nzero}{N_{0}}
\newcommand{\Cp}{C'}                     % C' := 103
\newcommand{\Aexp}{A}                    % A := 105 (degree-recursion exponent)
\newcommand{\Czero}{C_{0}}
\newcommand{\cEG}{c_{\mathrm{EG}}}
\newcommand{\epsOne}{\varepsilon_{1}}
\newcommand{\epsTwo}{\varepsilon_{2}}
\newcommand{\thetaQ}{\theta_{Q}}
\newcommand{\epsA}{\varepsilon_{A}}
\newcommand{\epsU}{\varepsilon_{U}}
\newcommand{\epsK}{\varepsilon_{K}}
\newcommand{\epsX}{\varepsilon_{X}}
\newcommand{\epsCONC}{\varepsilon_{\mathrm{CONC}}}
\newcommand{\epsChain}{\varepsilon_{\mathrm{ch}}}
\newcommand{\epsD}{\varepsilon(\Dstar)}  % generic: depends on D_* only, ->0
\newcommand{\Gc}[1]{\ensuremath{(\Gamma\mathrm{#1})}}   % galactic condition
\newcommand{\Gstar}{\ensuremath{(\mathrm{G}^{*})}}
\newcommand{\cOV}{c_{\mathrm{OV}}}       % log_2(4/3)
%% numerical constants of the cost line
\newcommand{\cVX}{80}
\newcommand{\cTPV}{169}
\newcommand{\cPV}{369}
\newcommand{\cKRED}{745}
\newcommand{\cKREDexact}{739}
\newcommand{\cFresh}{338}
\newcommand{\cTot}{1085}
\newcommand{\cTotR}{1091}

%% ----------------------------------------------------------------------
%% The hierarchy (s2)
%% ----------------------------------------------------------------------
\newcommand{\HBs}{\mathrm{HB}^{*}}
\newcommand{\HBt}{\mathrm{HB}^{*\tau}}
\newcommand{\HBtp}{\mathrm{HB}^{*\tau{+}}}
\newcommand{\tHB}{t^{\mathrm{HB}}}
\newcommand{\Cyc}{\mathrm{Cyc}}
\newcommand{\Std}{\mathrm{Std}}
\newcommand{\home}{\mathrm{home}}
\newcommand{\anc}{\mathrm{anc}}
\newcommand{\Dup}{\mathrm{Dup}}
\newcommand{\Dups}{\mathrm{Dup}^{*}}
\newcommand{\dup}{\mathrm{dup}}
\newcommand{\mult}{\mathrm{mult}}
\newcommand{\piece}{\mathcal{Q}}          % pieces of the s=0 Lemma-14 run
\newcommand{\Leaf}{\mathfrak{L}}          % a leaf of a Lemma-14^tau recursion
\newcommand{\Hout}{H_{\mathrm{out}}}
\newcommand{\Hin}{H_{\mathrm{in}}}
\newcommand{\thGC}{\theta^{\mathrm{GC}}}
\newcommand{\Rt}{\mathrm{Rt}}             % rooted vertices of a light child
\newcommand{\Pl}{\mathrm{Pl}}             % parentless vertices of a light child
\newcommand{\parU}{\mathrm{par}}          % the fixed parent of a rooted vertex
\newcommand{\Kstd}{K^{\mathrm{std}}}      % union of standalone edge sets
\newcommand{\Bad}{\mathrm{Bad}}            % demoted or parent-bad light parts

%% ----------------------------------------------------------------------
%% Colouring, lending, zones (s3, s5)
%% ----------------------------------------------------------------------
\newcommand{\Own}{\mathrm{Own}}
\newcommand{\Lend}{\mathrm{Lend}}
\newcommand{\Lent}{\mathrm{Lent}}
\newcommand{\LentU}{\mathrm{Lent}^{\mathrm{U}}}
\newcommand{\LentJS}{\mathrm{Lent}^{\mathrm{JS}}}
\newcommand{\LentJV}{\mathrm{Lent}^{\mathrm{JV}}}
\newcommand{\Lentext}{\mathrm{Lent}^{\mathrm{ext}}}
\newcommand{\klend}{k_{\mathrm{lend}}}
\newcommand{\kown}{k_{\mathrm{own}}}
\newcommand{\IU}{\mathcal{I}_{\mathrm{U}}}
\newcommand{\IJS}{\mathcal{I}_{\mathrm{JS}}}
\newcommand{\IJV}{\mathcal{I}_{\mathrm{JV}}}
\newcommand{\LU}{\mathcal{L}^{\mathrm{U}}}
\newcommand{\LJS}{\mathcal{L}^{\mathrm{JS}}}
\newcommand{\LJV}{\mathcal{L}^{\mathrm{JV}}}
\newcommand{\KJS}{K^{\mathrm{JS}}}
\newcommand{\tJS}{t^{\mathrm{JS}}}
\newcommand{\KHUB}{K^{\mathrm{HUB}}}
\newcommand{\lab}{\mathrm{lab}}
\newcommand{\Zone}{\mathrm{Zone}}
\newcommand{\uslot}{\varsigma}            % Theorem-U slot index
\newcommand{\Tslot}{T^{\mathrm{sl}}}      % number of slots T^sl_Y
\newcommand{\zp}{\varpi}                  % zone-choice probability varpi_Y
\newcommand{\Bdl}{\mathrm{Bdl}}           % Theorem-U chaining bundle (l,c)
\newcommand{\Aw}{\mathcal{A}}             % hub-avoiding sets A(w) of Lemma HB

%% ----------------------------------------------------------------------
%% Vortex-local notation (sans serif = local to one vortex run; s4)
%% ----------------------------------------------------------------------
\newcommand{\vJ}{\mathsf{J}}
\newcommand{\vR}{\mathsf{R}}
\newcommand{\vM}{\mathsf{M}}
\newcommand{\vU}{\mathsf{U}}
\newcommand{\vW}{\mathsf{W}}
\newcommand{\vZ}{\mathsf{Z}}
\newcommand{\vV}{\mathsf{V}}
\newcommand{\vH}{\mathsf{H}}
\newcommand{\vF}{\mathsf{F}}
\newcommand{\vP}{\mathsf{P}}
\newcommand{\vQ}{\mathsf{Q}}
\newcommand{\vC}{\mathsf{C}}
\newcommand{\vlev}{\mathsf{lv}}
\newcommand{\vkap}{\mathsf{k}}
\newcommand{\vm}{\mathsf{m}}
\newcommand{\vb}{\mathsf{b}}
\newcommand{\vt}{\mathsf{t}}
\newcommand{\vPaths}{\mathcal{P}^{\mathsf{v}}}
\newcommand{\vPathsQ}{\widetilde{\mathcal{P}}^{\mathsf{v}}}

%% ----------------------------------------------------------------------
%% Standalone parts, TPV interface, JS-LC, chaining (s6)
%% ----------------------------------------------------------------------
\newcommand{\Lost}{\mathrm{Lost}}
\newcommand{\Ret}{\mathrm{Ret}}           % Ret_Z := A_Z u F_Z u Lost_Z
\newcommand{\Qs}{Q^{*}}                   % Q*_Z := Q_Z \ Lost_Z
\newcommand{\Erem}{E^{\mathrm{rem}}}
\newcommand{\EV}{E^{\mathrm{V}}}
\newcommand{\EQ}{E^{\mathrm{Q}}}
\newcommand{\Bead}{\mathfrak{B}}          % deterministic bead graph B_{Y,l}
\newcommand{\cagg}{c^{\mathrm{agg}}}
\newcommand{\cpp}{c_{\mathrm{pp}}}
\newcommand{\sco}{\mathrm{sc}}            % split cost sc_{Y,l}
\newcommand{\clu}{\mathfrak{K}}           % a cluster
\newcommand{\exc}{\mathrm{exc}}
\newcommand{\dem}{\mathrm{dm}}            % HCC-P demands dm^-(u), dm^+(u)
\newcommand{\pad}{\mathrm{pad}}           % HCC-P padding
\newcommand{\MED}{\mathrm{MED}}
\newcommand{\Jlost}{J^{\mathrm{lost}}}
\newcommand{\Jhub}{J^{\mathrm{hub}}}
\newcommand{\Jfr}{J^{\mathrm{fr}}}
\newcommand{\Jpar}{J^{\mathrm{par}}}
\newcommand{\Jp}[1]{\textup{(J#1)}}       % properties (J1)--(J3)
\newcommand{\JC}[1]{\textup{(JC#1)}}      % J-consumer interface (JC1)--(JC3)

%% ----------------------------------------------------------------------
%% Junction-vertex quotient, ultra hubs, induction (s7)
%% ----------------------------------------------------------------------
\newcommand{\plab}{\mathrm{pl}}           % round-split pool label pl(v)
\newcommand{\Pool}{\mathrm{Pool}}
\newcommand{\Cand}{\mathrm{Cand}}
\newcommand{\Hcd}{H^{\mathrm{cd}}}        % candidate threshold H^cd_l
\newcommand{\thult}{\theta^{\mathrm{ult}}}
\newcommand{\clive}{c^{\mathrm{live}}}
\newcommand{\tCC}{t^{\mathrm{CC}}}
\newcommand{\Used}{\mathrm{Used}}
\newcommand{\Live}{\mathrm{Live}}
\newcommand{\BP}{B^{\mathrm{P}}}
\newcommand{\BH}{B^{\mathrm{H}}}
\newcommand{\Dec}{\mathcal{D}}
\newcommand{\Past}{\mathfrak{F}^{-}}
\newcommand{\XU}{X_{\mathrm{U}}}
\newcommand{\Xpool}{X_{\mathrm{pool}}}
\newcommand{\XV}{X_{\mathrm{V}}}
\newcommand{\Xp}{X'}
\newcommand{\dett}{\mathrm{det}}
\newcommand{\copies}{\mathrm{copies}}
\newcommand{\payrd}{\mathrm{pay}^{\mathrm{rd}}}
\newcommand{\Obj}{\mathcal{O}}            % objects output by the round-l J-consumer

%% ----------------------------------------------------------------------
%% Local helpers used in several sections
%% ----------------------------------------------------------------------
\newcommand{\UQ}{\mathcal{M}}             % Theorem-U chaining quotient multigraph
\newcommand{\LayP}{\mathcal{X}}           % HCC-P: port set of layer j
\newcommand{\pot}{\mathrm{pot}}           % MED / HCC-P potential
\newcommand{\extph}{\varphi}              % P-vortex external phase map

\title{A proof of the Erd\H{o}s--Gallai cycle decomposition conjecture}
\author{Ryan Coffey\\
\small Yale University}
\date{1 October 2026}

\begin{document}
%% ======================================================================
%% Section 1, preceded by the front matter (\maketitle and the abstract).
%% Uses only the macros and environments of the preamble; defines nothing.
%% ======================================================================
\maketitle

\begin{abstract}
The Erd\H{o}s--Gallai conjecture (Erd\H{o}s Problem \#184) asserts that the edge
set of every $n$-vertex simple graph can be partitioned into $O(n)$ cycles and
single edges. We prove this conjecture: every $n$-vertex graph can be
decomposed into $O(n)$ cycles and edges. The best previous bound, $O(n\logs
n)$, is due to \BMnames{} \cite{BM}.

Our argument builds on the method of \cite{BM} and has three new ingredients.
First, \emph{$\tau$-rules} placed inside the recursion of \BMnames's Lemma~14,
together with a cap on the degrees of guest vertices, produce a hierarchy $\HBtp$
in which hub concentration is negligible for every run. Second, \emph{standalone
parts are chained} into cycles through junction paths that run inside lent
random colour classes of ancestors at least two rounds back (a transparent
vortex, TPV, together with a junction synchronisation that splits every large
class into cherries, JS-LC); only parity leftovers, single edges, remain. Third,
\emph{each end of a leftover single edge receives a private junction edge}, which
joins it to a junction vertex (possibly shared with other ends) in a sparse
random pool inside its ancestor's expander; the resulting simple quotient graphs
have at most a vanishing fraction of $n$ vertices in total (the fraction tends to
$0$ as the degree threshold $\Dstar$ grows), every decomposition of a quotient
lifts back, and the quotients are handled by induction on the number of
vertices.

The constants are galactic and are not computed: the degree threshold $\Dstar$ is
only required to be sufficiently large (each condition on it holds for all
sufficiently large $\Dstar$), and it is of tower type ($\log_2\log_2\Dstar\ge2^8$
is one of these conditions). The bound $f(n)=O(n)$ has also been formally verified in the
Lean~4 proof assistant.
\end{abstract}
\section{Introduction, conventions, cited inputs, constants, dependency structure}\label{s1:sec}

For a graph $G$ let $\f(G)$ be the least number of cycles and single edges that
partition $E(G)$, and let $\f(n)$ be the maximum of $\f(G)$ over graphs $G$ on
$n$ vertices (Definition~\ref{s1:defObject}). Erd\H{o}s and Gallai
conjectured that $\f(n)=O(n)$ (see \cite{EG66}). Repeatedly removing a longest cycle gives
$\f(n)=O(n\log n)$ \cite{EG66}; Conlon, Fox and Sudakov \cite{CFS14} proved
$\f(n)=O(n\log\log n)$; and \BMnames{} \cite{BM} proved $\f(n)=O(n\logs n)$,
which was the best previously known bound. Lov\'asz \cite{Lov68} proved that every
$n$-vertex graph decomposes into at most $n/2$ paths and cycles.

In this paper we prove that $\f(n)=O(n)$, confirming the conjecture of
Erd\H{o}s and Gallai. The precise statement is Theorem~\ref{s1:thmMain}, which
is proved in Section~\ref{s7:sec}. The proof builds on the method of \BMnames{}
\cite{BM}: it keeps the round structure of their argument, and it uses, and in
several places modifies, their lemmas on expanders and path connectivity.
Remark~\ref{s1:remStatus} says where the new ingredients lie, and the overview
that follows it describes the architecture of the argument.
Section~\ref{s1:sec} contains the overview, the conventions, every cited result
(each published result of \cite{BM} is quoted in our notation, with its number in
\cite{BM}), one elementary decomposition bound and one standard concentration
inequality that are proved in full instead of being cited (Fact~\ref{s1:factEG0}
and Lemma~\ref{s1:lemBBD}), the constants and the single list of galactic
conditions, the symbol table, and the dependency table of all numbered
statements. The bound $f(n)=O(n)$ has also been formally verified in Lean~4 with Mathlib;
Section~\ref{s8:sec} describes the formalization.

\begin{theorem}[Main Theorem]\label{s1:thmMain}
Let $\eps,\sigma,\Cp,\Aexp,\Nzero,\Dstar$ be as in
Definition~\ref{s1:defConstants}, with $\Dstar$ satisfying
\ref{s1:condG1}--\ref{s1:condG4} of Condition~\ref{s1:condGamma}. Put
\[
\Czero\defeq \frac{\Dstar}{2}+\cTot+\epsOne(\Dstar)
\qquad\text{and}\qquad
\cEG\defeq\max\Bigl\{\frac{\Czero}{1-2\thetaQ},\,\frac{\Nzero}{2}\Bigr\},
\]
where $\epsOne(\Dstar)$ and $\thetaQ=\epsTwo(\Dstar)$ are the explicit functions
of $\Dstar$ defined in Proposition~\ref{s7:propCost} and
Lemma~\ref{s7:lemUHsplit}. Then every graph $G$ on $n$ vertices satisfies
\[
\f(G)\le\cEG\, n .
\]
In particular $\f(n)=O(n)$.
\deps{\ref{s7:thmMainProof} (forward pointer; in Table~\ref{s1:tabDAG} the proof
is the node \texttt{s7:thmMainProof}).}
\end{theorem}

The proof of Theorem~\ref{s1:thmMain} is given in Section~\ref{s7:sec}, as
Corollary~\ref{s7:thmMainProof}. Theorem~\ref{s1:thmMain} is the only statement
of this paper that is stated before the statements it depends on.

\begin{remark}[organization of the proof; the new ingredients]\label{s1:remStatus}
\begin{enumerate}
\item[(i)] Compared with \cite{BM}, the new ingredients of the argument are
concentrated in the following places:
\begin{enumerate}
\item[(a)] the routing engine inside JS-LC (Lemmas \ref{s6:lemHCCP},
\ref{s6:lemPAR}, \ref{s6:lemMED} and \ref{s6:lemEQLPT} as used for cherry
systems at junction multiplicity $2M_l-2$ in Lemma~\ref{s6:lemJSLC});
\item[(b)] Lemma 14$^\tau$ (Lemma~\ref{s2:lem14tau}) and part~(iv) of
Theorem~\ref{s6:thmCONCL}: together they modify the core Lemma~14 of
\cite{BM} instead of using it as a black box;
\item[(c)] the composition JV$^{+*}$ of Section~\ref{s7:sec}; within it, the
most delicate points are property \Jp{1} of Lemma~\ref{s6:lemJplus}, aggregated
over all MIX-parts, on which Lemma~\ref{s7:lemMULT} rests; the conditioning in
Lemma~\ref{s7:lemCC} and in the ultra-hub bound; the lift, including the
cross-round use of junction edges through the $\Lent$ sets; and the existence of
one joint outcome;
\item[(d)] every place where a per-vertex multiplicity feeds a path-connectivity
parameter~$t$: property (P4) of the vortex runs of Section~\ref{s4:sec}, Step~6
of Lemma~\ref{s5:lemParent}, and the joint routing of JS-LC at junction
multiplicity $2M_l-2$. In particular, Lemma~\ref{s4:lemPV}(c) bounds the number
of arc ends at a vertex by its degree, so that a vertex lies in at most $M_l-1$
pairs, and the U-lent classes are path connected for the multiplicity parameter
$t_Y=\lceil\lambda_r^{1.6}\rceil\ge M_l$ (Definition~\ref{s3:defCOL}; rows~4
and~7 of Table~\ref{s3:tabCOLJV}).
\end{enumerate}

\item[(ii)] The argument relies in particular on the following properties, each
proved where indicated:
\begin{enumerate}
\item[(1)] the hub concentration is negligible: for every valid $\HBtp$ run
(Definition~\ref{s2:defHBtp}) and every designation,
$\sum_{(Y,l)}m_{Y,l}\le\epsCONC(\Dstar)\,n$ (Theorem~\ref{s6:thmCONCL}(iv));
\item[(2)] in every round, JS-LC emits at most one $\Jhub$-edge of any given class at
any given hub (property \Jp{1} of Lemma~\ref{s6:lemJplus});
\item[(3)] the lift of every cycle of a decomposition of a quotient is a cycle
of $G$, and the lifted objects are pairwise edge-disjoint
(Lemma~\ref{s7:lemLift}(ii)); every edge of $G$ is covered exactly once, and in
particular no junction edge is used both by a lifted cycle and by the device of
its owner (Lemma~\ref{s7:lemLift}(iii) and Theorem~\ref{s6:thmMIXC}(b));
\item[(4)] every per-round cost term is of one of the three kinds listed in the
paragraph ``Per-round costs here'' of \ref{s1:ssecOverview}; in particular, no
per-round cost term of order $\Theta(n)$ occurs other than a once-per-vertex
charge.
\end{enumerate}
\end{enumerate}
\deps{none.}
\end{remark}

\subsection*{Overview: the architecture, and where the $\logs$ disappears}
\namedlabel{s1:ssecOverview}{the overview}

This overview is prose; it contains no statement. Every claim it makes about the
present argument is proved at the place it points
to; the paragraph on \cite{BM} describes their published proof.

\paragraph{The chain.} The argument is a chain of five links:
\[
\text{UNI-MIX}\ \longrightarrow\ \text{MIX-C on }\HBtp\text{ with no VX-parts}
\ \longrightarrow\ \text{JV}^{+*}\ \longrightarrow\ \text{UH}^{*}
\ \longrightarrow\ \text{layered HI}''.
\]
UNI-MIX and MIX-C decompose all edges except a set of parity leftovers; JV$^{+*}$
joins the ends of the leftovers by private junction edges to pooled junction
vertices, and so routes them into simple quotient graphs; UH$^{*}$ controls the size of these quotients (round-split pools and
sub-layered ultra hubs); the layered induction HI$''$ closes the argument.
Section~\ref{s2:sec} builds the hierarchy $\HBtp$ (Lemma~14$^\tau$, Lemma SEP,
the overlap bounds, rule GC, ORIGIN$^\tau$, the tower facts);
Section~\ref{s3:sec} proves the random splitting and linking tools
(Lemma~15$^+$, Theorem~16$^*$, Lemma HB) and defines the lending colouring
COL-JV; Section~\ref{s4:sec} proves the vortex tools (TPV, the four-phase
P-vortex, Theorem VX$^+$); Section~\ref{s5:sec} treats light parts (zones, the
child and parent sides of the Theorem-U machinery, Lemma K-RED);
Section~\ref{s6:sec} treats standalone parts (GATE, MED, EQ-LPT, PAR, HCC-P,
Theorems CONC and CONC-L, Lemma JS-LC, Lemma J$^+$ and the MIX-C assembly for an
abstract J-consumer); Section~\ref{s7:sec} constructs the junction-vertex
quotient, bounds its size and the cost, and proves Theorem JV$^{+*}$,
Theorem HI$''$ and the Main Theorem.

\paragraph{Where \BMnames{} lose the factor $\logs n$.} The proof of Theorem~2
of \cite{BM} iterates their Lemma~26 about $\logs n$ times: each iteration
removes long cycles, splits the remaining edges into expanders by Lemma~14 of
\cite{BM} (with $s=0$), and decomposes each (necessarily small) expander into
cycles and few edges; the average degree of the leftover drops from $d$ to a
polylogarithm of $d$. Each iteration costs $\Theta(n)$ objects, charged to all
vertices that are active in it, and there are $\logs n$ iterations.

\paragraph{Per-round costs here.} The hierarchy $\HBtp$ has the same round
structure (rounds $l=1,\dots,R$, with $R=O(\logs n)$). In the present
argument every per-round cost term is of one of three kinds.
\begin{enumerate}
\item[(a)] \emph{Decaying in the round index.} It is bounded by a constant
multiple of one of $n/M_l$, $n/\log P_l$, $n\lambda_r^{-1/2}$, $n/P_r$,
$n M_l^3\log\lambda_{l-2}/\lambda_{l-2}^{95}$ (the parameters are those of
Definition~\ref{s2:defHBtp}) and $n\,\eta(\lambda_{l-2})$, where
$\eta(x)=\log_2\bigl(2\Aexp\log_2(\Aexp\log_2x)\bigr)/(\log_2x)^2$; the last
bounds the number $\mathrm{cap}_l$ of visit-capping pieces of the Theorem-U
chaining (Lemma~\ref{s5:lemParent}(ii), Step~4 of its proof). The hub charge is
of this kind: a vertex can be a hub in many rounds, and the round-$l$ term is
$\sum_{Z\in\Std_l}|A_Z|\le15.2\,\eps n/\log P_l$ (Proposition~\ref{s2:propOV},
(K2)), which sums over all rounds to at most $\epsA n$
(Lemma~\ref{s2:lemTower}(e)). Because the round
parameters grow like a tower towards earlier rounds, each of these sums over all
rounds to at most a quantity that depends only on $\Dstar$, tends to $0$ as
$\Dstar\to\infty$, and is uniform in $R$ (tower-lacunary sums,
Lemma~\ref{s2:lemLacunary}; the terms
$nM_l^3\log\lambda_{l-2}/\lambda_{l-2}^{95}$ sum to at most
$2F(\log_2\Dstar)\,n$, with $F(x)=(3184+30400\log_2x)\,x^{-90.2}$, by
Lemma~\ref{s7:lemUHsplit}(iv), and the terms $\mathrm{cap}_l$ are part of
$\epsChain(\Dstar)\,n$ in Lemma~\ref{s5:lemParent}(ii)).
\item[(b)] \emph{A once-per-vertex global charge.} Fresh ports (a vertex is fresh
in at most two rounds, and TPV pays $\cTPV$ per fresh port: $\cTPV\cdot2n=\cFresh
n$); parentless light incidences (at most two per vertex, each paying $\cPV$:
$\cPV\cdot 2n$); unpaired fresh legs (at most $2n$); long cycles (at most $n$); and the final leftover
$E_0$ (fewer than $\Dstar n/2$ edges).
\item[(c)] \emph{Per-ancestor terms with a factor $R-r$.} Some terms are charged
to an ancestor $Y$ of round $r$ once in every later round, so that in total they
carry a factor $R-r$: the $\tau$-, $d^*$- and guest-cap terms of the
concentration bound (Theorems~\ref{s6:thmCONC}(iii) and~\ref{s6:thmCONCL}(iv)),
and the lost-parent mass $\mathrm{lp}=\sum_{Y\in\Bad}|Y|\,(R-r(Y))$
(Definition~\ref{s5:defStages}). The factor satisfies
$R-r\le2\logs d_r+2\le\log\lambda_r$ (Lemma~\ref{s2:lemTower}(a),(d)), so it
depends on $d_r$ alone, and it is outweighed by the decay of the rest of the term
in $\lambda_r$ (for $\mathrm{lp}$, by $\Prob(Y\in\Bad)\le|Y|^{-2}$,
Lemma~\ref{s5:lemE1}(c)). Hence these terms, too, are at most (for
$\mathrm{lp}$: in expectation) $n$ times a function of $\Dstar$ that tends to
$0$ as $\Dstar\to\infty$ and is uniform in $R$
(Theorem~\ref{s6:thmCONCL}(iv), Lemma~\ref{s5:lemExpect}). This is why
$\logs\Dstar$, and not $\logs n$, appears in $\epsCONC$.
\end{enumerate}

\paragraph{The three mechanisms.}
\begin{enumerate}
\item[(a)] \emph{No hub concentration.} The $\tau$-rules inside the recursion of
Lemma~14 of \cite{BM}, together with rule GC, give
$\sum_{(Y,l)}m_{Y,l}\le\epsCONC(\Dstar)\,n$ for every valid run and every
designation (Theorem~\ref{s6:thmCONCL}(iv)). The chaining cost of JS-LC is at
most $1.5\sum_{(Y,l)}m_{Y,l}$ plus a decaying term, hence a vanishing fraction
of $n$.
\item[(b)] \emph{TPV and JS-LC.} Classed ports are never retired by the vortex
of their own part; their edges become beads, which are chained into cycles
through junction paths inside lent classes of an ancestor of round at most
$l-2$. Only parity leftovers remain; they form the J-set $J_l$ of
Lemma~\ref{s6:lemJplus}.
\item[(c)] \emph{Private junction edges.} Each end of a J-item receives a
private junction edge, which joins it to a junction vertex in a pool of density
$M_l^{-2}$ inside the expander of its ancestor, which has minimum degree larger
than $s_r/2\gg M_l$. A junction vertex may serve several ends, but every
junction edge serves at most one (Lemma~\ref{s7:lemLift}(iii)). The quotient
$Q_l$ is simple, every decomposition of $Q_l$ lifts to a decomposition of the
corresponding edges of $G$ (Lemma~\ref{s7:lemLift}), and
$\sum_l|V(Q_l)|\le\thetaQ n\le n/4$ (Lemma~\ref{s7:lemUHsplit}). The
Erd\H{o}s--Gallai bound for the smaller graphs $Q_l$ is supplied by minimality in
the layered induction HI$''$ (Theorem~\ref{s7:thmHI}).
\end{enumerate}

\paragraph{The cost line.} On the outcome chosen in Section~\ref{s7:sec}
(Proposition~\ref{s7:propCost}),
\[
\f(G)\ \le\ \Bigl(\frac{\Dstar}{2}+\underbrace{\cKRED}_{\text{K-RED}}
+\underbrace{\cFresh}_{\text{TPV, fresh}}+\underbrace{2}_{\text{fresh legs}}
+\epsOne(\Dstar)\Bigr)n\;+\;2\sum_{l}\f(Q_l),
\]
with $\cKRED+\cFresh+2=\cTot$. Together with $\sum_l|V(Q_l)|\le\thetaQ n$ and
$\thetaQ\le1/4$, the induction of Theorem~\ref{s7:thmHI} gives
$\f(G)\le\Czero n+2\cEG\thetaQ n\le\cEG n$.

\subsection{Conventions}\label{s1:ssecConv}

\begin{convention}[graphs, logarithms, neighbourhoods, balls]\label{s1:convGraphs}
\begin{enumerate}
\item[(a)] Graphs are finite and simple, unless they are explicitly called
multigraphs (auxiliary multigraphs occur only in Sections~\ref{s5:sec}
and~\ref{s7:sec}). $G$ always denotes the input graph and $n\defeq|V(G)|$. For
any graph $H$ we write $|H|\defeq|V(H)|$, as in \cite{BM}. For $k\in\mathbb N$,
$[k]\defeq\{1,\dots,k\}$.
\item[(b)] $\log=\log_2$ throughout, as in \cite{BM}. We put
$\log^{[0]}x\defeq x$ and $\log^{[k]}x\defeq\log\bigl(\log^{[k-1]}x\bigr)$, and
$\logs x$ is the least integer $k\ge0$ with $\log^{[k]}x\le1$, as in
\cite{BM}.
\item[(c)] For a vertex $v$ of a graph $H$, $N_H(v)$ is its set of neighbours
and $d_H(v)=\deg_H(v)$ its degree; $\delta(H)$ and $\Delta(H)$ are the minimum
and maximum degree. For $U\subseteq V(H)$, $\Nbr_H(U)$ is the set of vertices of
$V(H)\setminus U$ that have a neighbour in $U$. $H[U]$ is the induced subgraph;
$H-U$ and $H\setminus U$ both denote $H[V(H)\setminus U]$; for
$F\subseteq E(H)$, $H-F$ is obtained by deleting the edges of $F$ (keeping all
vertices). For disjoint $A,B\subseteq V(H)$, $E_H(A,B)$ is the set of edges with
one end in $A$ and the other in $B$, and $e_H(A,B)\defeq|E_H(A,B)|$; for an
edge set $F$ and a vertex $v$, $\deg_F(v)$ is the number of edges of $F$ at $v$.
\item[(d)] A \emph{path through} $V$ is a path all of whose interior vertices lie
in $V$; its ends are unrestricted. For $U,V\subseteq V(H)$ and an integer
$i\ge0$, $\Ball{i}{H}{U,V}$ is the set of vertices of $V$ that can be reached by
a path through $V$ of length at most $i$ starting at a vertex of $U$ (the
starting vertex need not lie in $V$); thus $\Ball{i}{H}{U,V}\subseteq V$
(Cited result~\ref{s1:citNotation}).
\item[(e)] Families of pairs of vertices are \emph{multisets}: a pair may occur
several times, and ``every vertex lies in at most $t$ pairs'' counts occurrences
(Cited result~\ref{s1:citDef7}).
\end{enumerate}
\deps{none.}
\end{convention}

\begin{definition}[objects and $\f$]\label{s1:defObject}
An \emph{object} is a cycle (of length at least $3$) or a single edge. A
\emph{decomposition} of an edge set $F$ is a partition of $F$ into the edge sets
of objects; a decomposition of a graph $H$ is a decomposition of $E(H)$.
$\f(F)$ is the least number of objects in a decomposition of $F$ (so
$\f(\emptyset)=0$); $\f(H)\defeq\f(E(H))$; and
$\f(n)\defeq\max\{\f(H):|V(H)|=n\}$.
\deps{none.}
\end{definition}

\begin{fact}[elementary properties of $\f$]\label{s1:factAdd}
\begin{enumerate}
\item[(a)] $\f(F)\le|F|$ for every edge set $F$.
\item[(b)] If $F_1,F_2$ are disjoint edge sets, then $\f(F_1\cup F_2)\le
\f(F_1)+\f(F_2)$.
\item[(c)] If $H$ is the vertex-disjoint union of graphs $H_1$ and $H_2$, then
$\f(H)=\f(H_1)+\f(H_2)$.
\item[(d)] Adding or deleting isolated vertices does not change $\f$. In
particular $\f(n)$ is non-decreasing in $n$.
\end{enumerate}
\deps{\ref{s1:defObject}.}
\end{fact}
\begin{proof}
(a) The single edges of $F$ form a decomposition of $F$ into $|F|$ objects.
(b) The union of a decomposition of $F_1$ and a decomposition of $F_2$ is a
decomposition of $F_1\cup F_2$, because $F_1\cap F_2=\emptyset$.
(c) By (b), $\f(H)\le\f(H_1)+\f(H_2)$. Conversely, the edge set of every object
induces a connected graph, so every object of a decomposition of $E(H)$ has all
its edges in $E(H_1)$ or all in $E(H_2)$; hence every decomposition of $E(H)$
splits into a decomposition of $E(H_1)$ and one of $E(H_2)$, and
$\f(H)\ge\f(H_1)+\f(H_2)$.
(d) $\f(H)$ depends only on $E(H)$, which does not change. Adding an isolated
vertex to an $n$-vertex graph gives an $(n+1)$-vertex graph with the same $\f$,
so $\f(n+1)\ge\f(n)$.
\end{proof}

\begin{fact}[the long-cycle bound]\label{s1:factEG0}
\begin{enumerate}
\item[(a)] Every graph $H$ with $h\ge1$ vertices has a decomposition into at
most $h(\log h+1)$ objects, at most $h-1$ of which are single edges. In
particular $\f(H)\le h(\log h+1)$.
\item[(b)] Let $N>1$ and $L\defeq\log N$, and let $\alpha$ and $\beta$ be real
numbers with $0\le\alpha\le N$, $\beta>0$ and $4\beta\le L$. If $F$ is the
edge set of a (simple) graph, so that $F$ has no loops and no parallel edges,
and every edge of $F$ has both ends in a set $W$ with $|W|\le\beta\alpha/L$,
then $F$ has a decomposition into at most $\beta\alpha$ objects, at most $|W|$
of which are single edges.
\end{enumerate}
The argument, iterated removal of a long cycle, is the elementary
$O(n\log n)$ bound (see \cite{EG66} and \cite{BM}); it is proved here in full. It uses that graphs are
simple (Convention~\ref{s1:convGraphs}(a)) in two places: in the claim, where a
vertex of degree at least $d$ has at least $d$ neighbours, and through
$|E(H)|\le\binom h2$.
Part~(b) is the form in which the vortex runs of Section~\ref{s4:sec} are
finished.
\deps{\ref{s1:convGraphs}, \ref{s1:defObject}.}
\end{fact}
\begin{proof}
(a) Put $m\defeq|E(H)|$, so that $m\le\binom h2\le h^2/2$.

\emph{Claim. If a graph $H'$ with vertex set $V(H)$ has $m'\ge h$ edges, then
$H'$ contains a cycle of length at least $m'/h$.} Put
$d\defeq\max\{2,\lceil m'/h\rceil\}$. Then $d-1\le m'/h$: if
$\lceil m'/h\rceil\ge2$ then $d-1=\lceil m'/h\rceil-1<m'/h$, and otherwise
$d-1=1\le m'/h$ because $m'\ge h$. Delete from $H'$, one at a time, a vertex
whose degree in the current graph is at most $d-1$, as long as such a vertex
exists. Every edge of $H'$ disappears exactly once, together with the first of
its two ends to be deleted. Suppose that all $h$ vertices are deleted. The last
of them has degree $0$ when it is deleted, and each of the others has degree at
most $d-1$ when it is deleted; hence $m'\le(h-1)(d-1)\le(h-1)m'/h<m'$, since
$m'\ge h\ge1$, a contradiction. So the deletions stop at a subgraph $H^*$ of $H'$
with at least one vertex and minimum degree at least $d\ge2$. Let
$v_0v_1\cdots v_r$ be a longest path in $H^*$. Every neighbour of $v_0$ in $H^*$
is one of $v_1,\dots,v_r$, since otherwise the path could be extended. As $v_0$
has at least $d$ neighbours in $H^*$, the largest index $i$ with
$v_0v_i\in E(H^*)$ satisfies $i\ge d\ge2$, so $v_0v_1\cdots v_iv_0$ is a cycle of
$H'$ of length $i+1>d\ge\lceil m'/h\rceil\ge m'/h$. This proves the claim.

\emph{Removal of long cycles.} Put $H_0\defeq H$ and $m_i\defeq|E(H_i)|$. For
$i=0,1,\dots$, as long as $m_i\ge h$, choose by the claim (applied to
$H'\defeq H_i$) a cycle $C_i$ of $H_i$ of length at least $m_i/h$, and put
$H_{i+1}\defeq H_i-E(C_i)$. As $m_{i+1}<m_i$, this stops, at some $t\ge0$ with
$m_t\le h-1$. The cycles $C_0,\dots,C_{t-1}$ together with the $m_t$ edges of
$H_t$, taken as single edges, form a decomposition of $E(H)$ into $t+m_t$
objects, of which $m_t\le h-1$ are single edges. For $i<t$ we have
$m_{i+1}=m_i-|E(C_i)|\le(1-1/h)m_i$; hence $m_i\le(1-1/h)^im$ for $0\le i\le t$.

We use $(1-1/h)^h\le1/2$. For $h=1$ the left-hand side is $0$. For $h\ge2$ put
$x\defeq1/h$; then $(1-x)^h(1+x)^h=(1-x^2)^h\le1$ and, by the binomial theorem,
$(1+x)^h\ge1+hx=2$, so $(1-x)^h\le1/2$.

If $t=0$, the decomposition has $m_0\le h-1\le h(\log h+1)$ objects, as
$\log h\ge0$. Let $t\ge1$, and write $t-1=qh+r'$ with integers $q\ge0$ and
$0\le r'\le h-1$. Since $C_{t-1}$ was chosen, $m_{t-1}\ge h$; as $t-1\ge qh$ and
$0\le1-1/h\le1$,
\[
h\ \le\ m_{t-1}\ \le\ (1-1/h)^{t-1}m\ \le\ \bigl((1-1/h)^h\bigr)^qm\ \le\ 2^{-q}m\
\le\ 2^{-q}h^2/2 .
\]
Hence $2^{q+1}\le h$, i.e.\ $q+1\le\log h$, and $t=qh+r'+1\le(q+1)h\le h\log h$.
So the decomposition has $t+m_t\le h\log h+h-1\le h(\log h+1)$ objects.

(b) Note that $L>0$ as $N>1$. If $W=\emptyset$, then $F=\emptyset$, and the
empty decomposition has $0\le\beta\alpha$ objects. Otherwise $h\defeq|W|\ge1$,
and $h\le\beta\alpha/L\le\beta N/L\le N/4$ because $\alpha\le N$, $\beta>0$ and
$4\beta\le L$; hence $\log h\le\log N-2=L-2$. By (a), applied to the graph with
vertex set $W$ and edge set $F$, the set $F$ has a decomposition into at most
$h(\log h+1)\le h(L-1)\le hL\le\beta\alpha$ objects, at most $h-1\le|W|$ of
which are single edges.
\end{proof}

\subsection{Cited results of \BMnames{}}\label{s1:ssecCited}

Every published result of \cite{BM} used in this paper is stated below in our
notation, with its number in \cite{BM}; numbers refer to the arXiv version
arXiv:2211.07689v2. Where this paper modifies a proof of \cite{BM}
(Lemma~14, Propositions~12 and~13, Lemma~17, Proposition~18, Lemma~19, Lemma~9)
or extracts an explicit bound from a proof (Lemma~25), the relevant steps of the
proof are summarised as well; the modified proofs themselves are written out in
full where they are used. All items of this subsection are published results.

\begin{citedresult}[{\cite[Section~2.1, Section~2.6]{BM}}]\label{s1:citNotation}
\cite{BM} uses the following notation, which is also the notation of this
paper (Convention~\ref{s1:convGraphs}). For a graph $G$ and $v\in V(G)$,
$d_G(v)$ is the degree and $N_G(v)$ the neighbourhood of $v$; $\Delta(G)$ is the
maximum degree; for $U\subseteq V(G)$, $\Nbr_G(U)$ is the set of vertices of
$V(G)\setminus U$ with a neighbour in $U$; $G[V]$ is the induced subgraph,
$G\setminus V\defeq G[V(G)\setminus V]$, and $G-F$ is obtained by deleting the
edge set $F$; $|G|$ is the number of vertices. All logarithms have base~$2$;
$\log^{[0]}n=n$, $\log^{[k]}n$ is the $k$-fold iterated logarithm, and $\logs n$
is the least $k\ge0$ with $\log^{[k]}n\le1$. A path through $V$ is a path with
all internal vertices in $V$, and $\Ball{i}{G}{U,V}$ is the set of vertices of
$V$ that can be reached by a path through $V$ of length at most $i$ starting at a
vertex of $U$ (the starting vertex is not required to lie in $V$), so that
$\Ball{i}{G}{U,V}\subseteq V$.
\deps{none (published).}
\end{citedresult}

\begin{citedresult}[{\cite[Theorem~4]{BM} (Chernoff)}]\label{s1:citChernoff}
Let $n$ be an integer, $0\le\delta,p\le1$, $X\sim\Bin(n,p)$ and $\mu\defeq\Ex
X=np$. Then
\[
\Prob\bigl(X>(1+\delta)\mu\bigr)\le e^{-\delta^2\mu/3}
\qquad\text{and}\qquad
\Prob\bigl(X<(1-\delta)\mu\bigr)\le e^{-\delta^2\mu/2}.
\]

\deps{none (published).}
\end{citedresult}

\begin{citedresult}[{\cite[Definition~7]{BM}}]\label{s1:citDef7}
A graph $G$ is \emph{$(\ell,t)$-path connected through} a vertex set $V\subseteq
V(G)$ if, for every collection $\mathcal P\subseteq\binom{V(G)}{2}$ in which every
vertex lies in at most $t$ pairs, there are edge-disjoint paths $P_{x,y}$,
$\{x,y\}\in\mathcal P$, such that each $P_{x,y}$ is an $xy$-path through $V$ of
length at most $\ell$.

\emph{Multiset clause} (quoted in substance from \cite{BM}): here
$\binom{V(G)}{2}$ denotes the \emph{multiset} of pairs of distinct vertices of
$G$; in particular the same pair may occur several times in $\mathcal P$. Thus a
pair occurring $k$ times receives $k$ pairwise edge-disjoint paths, and the
bound $t$ counts occurrences. The ends $x,y$ are unrestricted (they need not lie
in $V$); only interior vertices must lie in $V$.
\deps{none (published).}
\end{citedresult}

\begin{citedresult}[{\cite[Proposition~8]{BM}}]\label{s1:citProp8}
Let $1\le\ell,t\le n$, let $G$ be an $n$-vertex graph and let $V\subseteq V(G)$
satisfy $|V|\ge4t-2$ and $|\Ball{\ell}{G}{U,V}|>|V|/2$ for every $U\subseteq
V(G)$ with $|U|=t$. If $x_1,\dots,x_{2t-1},y_1,\dots,y_{2t-1}$ are distinct
vertices of $G$, then for some $j\in[2t-1]$ there is an $x_jy_j$-path in $G$
through $V$ of length at most $4\ell\log n$.
\deps{none (published).}
\end{citedresult}

\begin{citedresult}[{\cite[Lemma~9]{BM}}]\label{s1:citLem9}
Let $n\ge2$ and $1\le\ell,k\le n$. Let $G$ be an $n$-vertex graph and
$V\subseteq V(G)$ with $|V|\ge n/8+1$, and suppose that for every non-empty
$U\subseteq V(G)$ and every $F\subseteq E(G)$ with
$|F|\le2^9k|U|(\ell\log n)^2$ we have $|\Ball{\ell}{G-F}{U,V}|>|V|/2$. Then $G$
is $(4\ell\log n,k)$-path connected through $V$.

\emph{Structure of the proof in \cite{BM}} (modified in Lemma~\ref{s3:lemL9rho},
which uses Haxell's theorem, Cited result~\ref{s1:citHaxell}, in place of
Theorem~6).
Let $\mathcal P$ be a collection of pairs in which every vertex lies in at most
$k$ pairs, let $r\defeq|\mathcal P|$, and index the pairs by $i\in[r]$. Let
$H_i$ be the hypergraph on $E(G)$ whose edges are the edge sets of the paths
through $V$ of length at most $4\ell\log n$ joining the $i$-th pair. For
$I\subseteq[r]$ let $M_I$ be a maximal matching of $\bigcup_{i\in I}H_i$; the
claim is $|M_I|\ge4\ell\log n\cdot|I|$. If not, the set $F$ of edges covered by
$M_I$ has $|F|<(4\ell\log n)^2|I|$. Let $I'\subseteq I$ be maximal such that no
vertex lies in two pairs of $I'$; then $2k|I'|\ge|I|$ and $|I'|\le n/2$. With
$t\defeq\lceil|I'|/16\rceil$ one has $2t-1\le|I'|$, $4t-2\le\lceil n/8\rceil+1\le
|V|$ and $|F|<8kt(8\ell\log n)^2$, so every $t$-set $U$ satisfies
$|\Ball{\ell}{G-F}{U,V}|>|V|/2$; Proposition~8 (Cited
result~\ref{s1:citProp8}), applied to $G-F$ and $2t-1$ pairs of $I'$, gives a
path in $G-F$ joining a pair of $I'$, i.e.\ an edge of some $H_j$ disjoint from
$M_I$, contradicting maximality. Theorem~6 of \cite{BM} (Aharoni--Haxell
\cite{AH00}) then gives the paths. Since the pairs are indexed by $i\in[r]$, and the inequality
$2k|I'|\ge|I|$ counts the pairs of $I$ with multiplicity (every pair of $I$
meets a vertex of a pair of $I'$, and each of the $2|I'|$ vertices of pairs of
$I'$ lies in at most $k$ pairs of $I$), the argument applies verbatim to
multisets $\mathcal P$.
\deps{none (published).}
\end{citedresult}

\begin{citedresult}[{\cite[Definition~11]{BM}}]\label{s1:citDef11}
An $n$-vertex graph $G$ is an \emph{$(\varepsilon,s)$-expander} if for every
$U\subseteq V(G)$ and every $F\subseteq E(G)$ with $1\le|U|\le2n/3$ and
$|F|\le s|U|$ we have
\[
|\Nbr_{G-F}(U)|\ge\frac{\varepsilon|U|}{\log^2n}.
\]
\emph{Remark (from \cite{BM}; hypotheses made explicit):} if $n\ge2$ and $\varepsilon>0$ then, for every
$v\in V(G)$, the pair $U=\{v\}$, $F=\{\text{edges at }v\}$ violates this
inequality, so necessarily $|F|>s$; hence $\delta(G)>s$ for every
$(\varepsilon,s)$-expander $G$ with $\varepsilon>0$ on at least two vertices.
(For $\varepsilon=0$ every graph satisfies the inequality; wherever this remark
is used below, $\varepsilon\ge2^{-7}$.)
\deps{none (published).}
\end{citedresult}

\begin{citedresult}[{\cite[Proposition~12]{BM}}]\label{s1:citProp12}
For $U\subseteq V(G)$ and $d>0$ let $N_{G,d}(U)$ be the set of vertices of
$V(G)\setminus U$ with at least $d$ neighbours in $U$. Let $G$ be an $n$-vertex
$(\varepsilon,s)$-expander, $U\subseteq V(G)$ with $1\le|U|\le2n/3$, and $F$ a set
of at most $s|U|/2$ edges. Then for every $0<d\le s$, either (a)
$|\Nbr_{G-F}(U)|\ge s|U|/(2d)$, or (b) $|N_{G-F,d}(U)|\ge\varepsilon|U|/\log^2n$.

\emph{Proof in \cite{BM}} (its counting is modified in
Proposition~\ref{s3:propP13s}). If (a) fails, the set
$X\defeq\Nbr_{G-F}(U)\setminus N_{G-F,d}(U)$ has fewer than $s|U|/(2d)$
vertices, and the set $F'$ of edges of $G-F$ between $U$ and $X$ has
$|F'|\le d|X|\le s|U|/2$. Then $|F\cup F'|\le s|U|$ and
$\Nbr_{G-F-F'}(U)=N_{G-F,d}(U)$, so the expansion of $G$ gives
(b).
\deps{none (published).}
\end{citedresult}

\begin{citedresult}[{\cite[Proposition~13]{BM}}]\label{s1:citProp13}
There is $n_0$ such that the following holds for $n\ge n_0$,
$1\ge\varepsilon\ge2^{-9}$ and $s\ge8\log^{13}n$. Let $G$ be an $n$-vertex
$(\varepsilon,s)$-expander, $U\subseteq V(G)$ with $|U|\le2n/3$, and $F$ a set of
at most $s|U|/4$ edges. Then $G-F$ contains either
\begin{enumerate}
\item[(a)] $|U|/\log^7n$ vertex-disjoint stars, each with $\log^9n$ leaves, with
centre in $U$ and all leaves in $V(G)\setminus U$, or
\item[(b)] a bipartite subgraph $H$ with classes $U$ and $X\subseteq
V(G)\setminus U$ such that $|X|\ge\varepsilon|U|/(2\log^2n)$, every vertex of
$X$ has degree at least $\log^4n$ in $H$, and every vertex of $U$ has degree at
most $2\log^9n$ in $H$.
\end{enumerate}
\emph{Structure of the proof in \cite{BM}} (modified in
Proposition~\ref{s3:propP13s}). Take a maximal family of vertex-disjoint stars as
in (a), with centre set $C\subseteq U$ and leaf set $L$. If (a) fails, then
$|C|\le|U|/\log^7n$, $|L|\le|U|\log^2n$, and by maximality no vertex of
$U\setminus C$ has $\log^9n$ neighbours in $G-F$ outside $U\cup L$, whence
$|\Nbr_{G-F}(U\setminus C)|<2|U|\log^9n$. With $d\defeq\log^4n$ and
$\Delta\defeq2\log^9n$, go through the vertices $v$ of $V(G)\setminus U$ in turn
and add a star with centre $v$ and $d$ leaves in $U$ whenever the degrees in $U$
stay at most $\Delta$; $X$ is the set of centres added and $H$ the union of the
stars. The saturated set $U'$ (vertices of $U\setminus C$ of $H$-degree exactly
$\Delta$) satisfies $|U'|\le2d|L|/\Delta\le\varepsilon|U|/(8\log^2n)$, because a
saturated vertex has at least $\Delta/2$ of its $H$-neighbours in $L$. With
$B\defeq C\cup U'$, Proposition~12 applied to $U\setminus B$ excludes option (a)
there, and every vertex of $N_{G-F,d}(U\setminus B)\setminus B$ lies in $X$,
giving $|X|\ge\varepsilon|U|/(2\log^2n)$.
\deps{none (published).}
\end{citedresult}

\begin{citedresult}[{\cite[Lemma~14]{BM}}]\label{s1:citLem14}
Given an $n$-vertex graph $G$, a non-negative integer $s$ and
$\varepsilon\le2^{-5}$, one can delete at most $4sn\log n$ edges of $G$ so that
the remaining edges can be partitioned into graphs $G_1,\dots,G_r$ with
$\sum_{i=1}^r|G_i|\le2n$, each $G_i$ an $(\varepsilon,s)$-expander. (For $s=0$
nothing is deleted.)

\emph{Structure of the proof in \cite{BM}} (modified in
Lemma~\ref{s2:lem14tau}). The proof is by induction on $n$ under the stronger
bound $\sum_i|G_i|\le2n-2n/(2+\log n)$, which is at least $n$; the case $n=1$ is
trivial, and if $G$ is an $(\varepsilon,s)$-expander one takes $G_1=G$.
Otherwise there is a witness $(U,F)$: $1\le|U|\le2n/3$, $F\subseteq E(G)$,
$|F|\le s|U|$ and $|\Nbr_{G-F}(U)|<\varepsilon|U|/\log^2n$. Put
\[
G_1\defeq G[U\cup\Nbr_{G-F}(U)]-F,\qquad
G_2\defeq G\setminus U-E(G_1)-F,
\]
which edge-partition $G-F$, and $n_1\defeq|G_1|=|U\cup\Nbr_{G-F}(U)|$,
$n_2\defeq|G_2|=n-|U|$. Then
\begin{align*}
&\text{(6)}\quad n_1+n_2=n+|\Nbr_{G-F}(U)|<n+\varepsilon n_1/\log^2n;\\
&\text{(7)}\quad n_2<n\ \text{ and }\ n_1\le|U|+\varepsilon|U|/\log^2n\le
2n/3+\varepsilon n\le3n/4<n.
\end{align*}
The induction hypothesis gives sets $E_i\subseteq E(G_i)$ with $|E_i|\le4sn_i\log
n_i$ and partitions of $G_i-E_i$ into $(\varepsilon,s)$-expanders $G_{i,j}$ with
$\sum_j|G_{i,j}|\le2n_i-2n_i/(2+\log n_i)$. Using $\log n_1\le\log(3n/4)<\log
n-2/5$:
\begin{align*}
&\text{(8)}\quad |F|+|E_1|+|E_2|\le s\bigl(|U|+4n_1\log n_1+4n_2\log n_2\bigr)
\le s\bigl(4(n_1+n_2)\log n-\tfrac35n_1\bigr)\le4sn\log n,\\
&\text{(9)}\quad \frac{2n_1}{2+\log n_1}>\frac{2n_1}{2+\log
n}+\frac{n_1}{10\log^2n},
\end{align*}
where the last step of (8) uses (6). For $n\ge3$, (9) follows from $\log
n_1<\log n-2/5$ and $n_1<n$: then $\frac2{2+\log n_1}-\frac2{2+\log
n}>\frac{0.8}{(2+\log n)^2}\ge\frac1{10\log^2n}$, where the last step holds
because $\log n\ge\log3>2/(2\sqrt2-1)$. For $n=2$ we have $n_1=1$ by (7), and
(9) reads $1>\frac23+\frac1{10}$. Combining (9) with (6) gives
$\sum_{i,j}|G_{i,j}|\le(n_1+n_2)(2-\frac{2}{2+\log n})-\frac{n_1}{10\log^2n}\le
2n-\frac{2n}{2+\log n}$. The deleted edges are $F\cup E_1\cup E_2$.

\deps{none (published).}
\end{citedresult}

\begin{citedresult}[{\cite[Lemma~15]{BM}}]\label{s1:citLem15}
Let $n,s\in\mathbb N$, $k\ge1$ an integer, and $0<\varepsilon\le1$, and let $G$ be an $n$-vertex
$(\varepsilon,s)$-expander with $s\ge2^{12}\varepsilon^{-1}k^2\log^4n$. Then
there are edge-disjoint graphs $G_1,\dots,G_k$ with $E(G)=\bigcup_iE(G_i)$, each
an $\bigl(\varepsilon/4,\sqrt{s\varepsilon}/(8k\log n)\bigr)$-expander with vertex
set $V(G)$.

This result is quoted for comparison only and is \emph{not used}: it is replaced
by Lemma~\ref{s3:lemL15p} (Lemma~15$^+$), which is proved directly and loses
neither the factor $4$ in $\varepsilon$ nor the square root in
$s$.
\deps{none (published).}
\end{citedresult}

\begin{citedresult}[{\cite[Theorem~16]{BM}}]\label{s1:citThm16}
Let $G$ be an $n$-vertex $(\varepsilon,s)$-expander with
$1\ge\varepsilon\ge2^{-7}$ and $s\ge\log^{135}n$, and let $V\subseteq V(G)$
contain each vertex independently with probability $1/3$. Then with high
probability $G$ is $(4\log^5n,2^8\log^5n)$-path connected through $V$.

This theorem is not used as a black box. Its proof (split $E(G)$ by Lemma~15
into $k=2^{17}\log^{42}n$ expanders; apply Lemma~19 in each; for a given $U$ and
$F$ with $|F|\le2^{17}|U|\log^{15}n$ some class contains at most
$|U|/\log^{27}n$ edges of $F$; conclude by Lemma~9 with $k=2^8\log^5n$ and
$\ell=\log^4n$) is modified in Theorem~\ref{s3:thmT16s}.
\deps{none (published).}
\end{citedresult}

\begin{citedresult}[{\cite[Lemma~17]{BM}}]\label{s1:citLem17}
Let $n\ge2$ and let $G$ be an $n$-vertex $(\varepsilon,s)$-expander with
$2^{-9}\le\varepsilon\le1$ and $s\ge8\log^{13}n$. Let $U\subseteq V(G)$ satisfy
$|\Nbr_G(U)|\ge|U|\log^{24}n$, let $F\subseteq E(G)$ satisfy $|F|\le|U|$, and let
$V\subseteq V(G)$ contain each vertex independently with probability $1/3$. Then
with probability $1-e^{-\Omega(|U|\log^2n)}$ we have
$|\Ball{\log^4n}{G-F}{U,V}|>|V|/2$.

\emph{Proof parameters and structure in \cite{BM}} (modified in
Lemma~\ref{s3:lemL17s}). $\ell\defeq\log^4n$, $q\defeq3/11$, and $p$ is defined
by $(1-p)^{\ell-1}=11/12$, so that $p\ge1/(15\log^4n)$; $V$ is realised as
$V_1\cup\dots\cup V_\ell$ with independent $V_i$, each vertex in $V_i$ with
probability $p$ for $i<\ell$ and $q$ for $i=\ell$. $B_i$ is the set of vertices
reachable from $U$ by a path of $G-F$ of length at most $i$ whose interior lies
in $V_1\cup\dots\cup V_{i-1}$; it is determined by $U,V_1,\dots,V_{i-1}$, every
vertex of $\Nbr_{G-F}(B_i)$ with a $G-F$-neighbour in $B_i\cap V_i$ lies in
$B_{i+1}$, and $B_\ell\cap V_\ell\subseteq\Ball{\ell}{G-F}{U,V}$. The growth of
$B_i$ is shown via Proposition~13: in case (a) by Chernoff's bound on the star
centres, in case (b) by a bounded-differences inequality for the number of
vertices of $X$ with an $H$-neighbour in $V_i$; the final stage uses Chernoff's
bound with density $q$.
\deps{none (published).}
\end{citedresult}

\begin{citedresult}[{\cite[Proposition~18]{BM}}]\label{s1:citProp18}
Let $n\ge2$, $0<\varepsilon\le1$ and $s\ge\log^{24}n$. Let $G$ be an $n$-vertex
$(\varepsilon,s)$-expander and $U\subseteq V(G)$ with $|U|\le2n/3$. Then there is
$U'\subseteq U$ with $|\Nbr_G(U')|\ge|U'|\log^{24}n$ and
$|U'|\ge\varepsilon|U|/(3\log^{26}n)$.

\emph{Proof in \cite{BM}} (modified in Lemma~\ref{s3:lemP18s}). Put
$\theta\defeq\log^{24}n$ and let $U'\subseteq U$ be maximal with
$|\Nbr_G(U')|\ge\theta|U'|$. Let $F$ be the set of edges from $U\setminus U'$ to
$V(G)\setminus(U'\cup\Nbr_G(U'))$. Expansion applied to $U$ and $F$ gives
$\varepsilon|U|/\log^2n\le|\Nbr_{G-F}(U)|\le|\Nbr_G(U')|<(|U'|+1)\theta+1\le
3|U'|\theta$.

\emph{Remark.} The maximality step gives, for every $v\in U\setminus
U'$, \emph{fewer than $\theta+1$} neighbours outside $U'\cup\Nbr_G(U')$ (adding
$v$ to $U'$ removes at most the one vertex $v$ from the neighbourhood); \cite{BM}
write ``at most $\theta$'', which is exact only when $\theta$ is an integer.
Hence $|F|\le(\theta+1)|U\setminus U'|\le(\theta+1)|U|$, and the proof needs
$s\ge\theta+1$ rather than $s\ge\theta$. Lemma~\ref{s3:lemP18s} is formulated
with $s_1\ge\theta_*+1$ accordingly.
\deps{none (published).}
\end{citedresult}

\begin{citedresult}[{\cite[Lemma~19]{BM}}]\label{s1:citLem19}
Let $G$ be an $n$-vertex $(\varepsilon,s)$-expander with
$2^{-9}\le\varepsilon\le1$ and $s\ge2\log^{24}n$, and let $V\subseteq V(G)$
contain each vertex independently with probability $1/3$. Then with probability
$1-o(1/n)$, for every non-empty $U\subseteq V(G)$ and every $F\subseteq E(G)$ with
$|F|\le|U|/\log^{27}n$, we have $|\Ball{\log^4n}{G-F}{U,V}|>|V|/2$. (The
statement in \cite{BM} says ``every $U$''; the empty set is excluded implicitly.)

\emph{Structure of the proof in \cite{BM}} (modified in
Lemma~\ref{s3:lemP18s}). Call $U'$ well-expanding if
$|\Nbr_G(U')|\ge|U'|\log^{24}n$. A union bound of the failure probability of
Lemma~17 over all pairs $(U',F)$ with $U'$ well-expanding and $|F|\le|U'|$ gives
failure probability $o(1/n)$. For $|U|\le2n/3$, Proposition~18 gives a
well-expanding $U'\subseteq U$ with $|U'|\ge\varepsilon|U|/(3\log^{26}n)\ge|F|$,
and the ball is monotone in $U$; for $|U|>2n/3$ one passes to a subset $\bar U$
with $n/2\le|\bar U|\le2n/3$.
\deps{none (published).}
\end{citedresult}

\begin{citedresult}[{\cite[Theorem~21]{BM} (Lov\'asz \cite{Lov68})}]\label{s1:citThm21}
Every $n$-vertex graph can be decomposed into at most $n/2$ paths and
cycles.
\deps{none (published).}
\end{citedresult}

\begin{citedresult}[{\cite[Corollary~22]{BM}}]\label{s1:citCor22}
Every graph can be decomposed into paths such that each vertex is an end of at
most two of the paths. (The proof in \cite{BM} adds a new vertex joined to every
vertex of even degree, applies Theorem~21, and deletes the new
vertex.)

\emph{Consequence used in this paper.} If every edge of a graph $H$ has both
ends in a set $W$, then $E(H)$ decomposes into at most $|W|$ paths. Indeed, apply
the corollary to the graph $(W,E(H))$: every path has two distinct ends in $W$
and every vertex of $W$ is an end of at most two paths, so twice the number of
paths is at most $2|W|$.
\deps{none (published).}
\end{citedresult}

\begin{citedresult}[{\cite[Lemma~25]{BM}, explicit form}]\label{s1:citLem25}
Let $\varepsilon\ge2^{-5}$ and $m\ge\max(2,2^{30}/\varepsilon^2)$ (so
$m\ge2^{40}$ suffices at $\varepsilon=2^{-5}$). Every $m$-vertex
$(\varepsilon,0)$-expander contains a cycle of length at least
$\varepsilon^2m/(18\log^4m)$. (The bound $m\ge2$, which excludes the
one-vertex graph, follows from $m\ge2^{30}/\varepsilon^2$ unless
$\varepsilon^2>2^{29}$; without it the statement would be false for such
$\varepsilon$ at $m=1$.)

\emph{Explicit form from the proof in \cite{BM}.} The statement of \cite{BM}
says ``a cycle of length $\Omega(m/\log^4m)$''. The proof runs a depth-first
search, finds a moment at which the unexplored set $U$ and the processed set $R$
have equal size, so that the current path $P$ contains
$\Nbr_G(U)$ and has at least $\varepsilon m/(3\log^2m)$ vertices; it splits $P$
into consecutive segments $X,Y,Z$ with $|X|,|Z|\ge|P|/3$ and
$\varepsilon^2m/(18\log^4m)\le|Y|<\varepsilon^2m/(9\log^4m)$. If $X$ and $Z$
are joined by a path avoiding $Y$, a shortest such path together with a segment
of $P$ is a cycle containing $Y$, of length at least $|Y|\ge
\varepsilon^2m/(18\log^4m)$; otherwise expansion of the side of size at most
$m/2$ gives a contradiction. This explicit bound is the form used in
Lemma~\ref{s2:lemCap}.
\deps{none (published).}
\end{citedresult}

\begin{remark}[where the results of \cite{BM} are used]\label{s1:remBMused}
\begin{center}
\small
\begin{tabular}{p{6.2cm}p{8.6cm}}
\textbf{Result of \cite{BM}} & \textbf{Used in}\\
\hline
Definition~7 (\ref{s1:citDef7}) & Sections~\ref{s3:sec}, \ref{s4:sec}, \ref{s5:sec}, \ref{s6:sec}\\
Lemma~9, Proposition~8 (\ref{s1:citLem9}, \ref{s1:citProp8}) & Lemma~\ref{s3:lemL9rho}\\
Theorem~6 (Aharoni--Haxell) & not used (Lemma~\ref{s3:lemL9rho} uses Haxell's theorem \cite{Hax95}, Cited result~\ref{s1:citHaxell}, instead)\\
Definition~11 (\ref{s1:citDef11}) & everywhere\\
Propositions~12--13, Lemma~17, Proposition~18, Lemma~19 & Section~\ref{s3:sec} (proof of Theorem~16$^*$)\\
Lemma~14 and its proof (\ref{s1:citLem14}) & Section~\ref{s2:sec}\\
Corollary~22 (\ref{s1:citCor22}) & Section~\ref{s4:sec}\\
Lemma~25 (\ref{s1:citLem25}) & Lemma~\ref{s2:lemCap}\\
Theorem~4 (\ref{s1:citChernoff}) & Sections~\ref{s3:sec}, \ref{s4:sec}\\
Lemma~15, Theorem~16, Theorem~21 & not used directly (Lemma~15 is replaced by Lemma~\ref{s3:lemL15p}; Theorem~16 by Theorem~\ref{s3:thmT16s}; Theorem~21 enters only through Corollary~22)\\
\hline
\end{tabular}
\end{center}
Theorem~2 of \cite{BM} (the bound $O(n\logs n)$) is not used, and neither are
Theorem~24 and Lemma~26 of \cite{BM}, which are used only in its proof: the vortex runs
of Section~\ref{s4:sec} are finished with the elementary long-cycle bound of
Fact~\ref{s1:factEG0}.
The three results of the row ``not used directly'' are not dependencies of this
remark.
\deps{\ref{s1:citNotation}, \ref{s1:citChernoff},
\ref{s1:citDef7}, \ref{s1:citProp8}, \ref{s1:citLem9}, \ref{s1:citDef11},
\ref{s1:citProp12}, \ref{s1:citProp13}, \ref{s1:citLem14},
\ref{s1:citLem17}, \ref{s1:citProp18}, \ref{s1:citLem19},
\ref{s1:citCor22}, \ref{s1:citLem25}, \ref{s1:factEG0}.}
\end{remark}

\subsection{Other cited and standard results}\label{s1:ssecOther}

The following standard results are used. Where a short derivation of the exact
form used here is needed, it is written out.
One of them, Lemma BBD below (a Bernstein inequality for functions of
independent indicators with bounded differences), is not cited but proved in
full. Its probabilistic part uses only finite sums. Its analytic part uses
only elementary properties of the exponential function (such as $e^v\ge1+v$
for real $v$) and the fact that a differentiable function whose derivative is
non-negative on an interval is non-decreasing there; the one elementary
one-variable inequality that it needs is proved there.

\begin{citedresult}[{Chernoff bounds for sums of independent indicators; \cite[Theorems~2.1 and~2.8]{JLR00}}]\label{s1:citChernoffGen}
Let $X=\sum_{i=1}^mI_i$ be a sum of independent indicator variables with
$\Prob(I_i=1)=p_i$ (a Poisson-binomial variable), and let $0\le\delta\le1$.
\begin{enumerate}
\item[(a)] If $\mu\ge\Ex X$ then $\Prob\bigl(X\ge(1+\delta)\mu\bigr)\le
e^{-\delta^2\mu/3}$; if $0\le\mu\le\Ex X$ then $\Prob\bigl(X\le(1-\delta)\mu\bigr)
\le e^{-\delta^2\mu/2}$. In particular the two bounds of Cited
result~\ref{s1:citChernoff} hold verbatim with $\mu=\Ex X$.
\item[(b)] If $\Ex X\le\mu$ then for every integer $j\ge1$,
\[
\Prob(X\ge j)\le\frac{\mu^j}{j!}\le\Bigl(\frac{e\mu}{j}\Bigr)^j.
\]
\end{enumerate}
\deps{none (standard).}
\end{citedresult}
\begin{proof}[Derivation of the forms used]
(a) By \cite[Theorems~2.1 and~2.8]{JLR00}, with $\lambda\defeq\Ex X$ and $t\ge0$,
$\Prob(X\ge\lambda+t)\le\exp\bigl(-t^2/(2(\lambda+t/3))\bigr)$ and
$\Prob(X\le\lambda-t)\le\exp\bigl(-t^2/(2\lambda)\bigr)$. For the upper tail put
$t\defeq(1+\delta)\mu-\lambda\ge\delta\mu$; the exponent
$t^2/(2(\lambda+t/3))$ is increasing in $t$ and, for fixed $\mu$, decreasing in
$\lambda$ (as $\lambda$ decreases, $t$ increases and
$\lambda+t/3=\frac23\lambda+\frac13(1+\delta)\mu$ decreases), so it is at least
its value at $\lambda=\mu$, namely $\delta^2\mu/(2(1+\delta/3))\ge\delta^2\mu/3$
since $\delta\le1$. For the lower tail write $\lambda=\mu+a$ with $a\ge0$ and
put $t\defeq\lambda-(1-\delta)\mu=\delta\mu+a\ge0$. Then
$t^2=\delta^2\mu^2+2\delta\mu a+a^2\ge\delta^2\mu^2+\delta^2\mu a=\delta^2\mu\lambda$
(as $\delta\le2$), so $t^2/(2\lambda)\ge\delta^2\mu/2$ (the case $\lambda=0$ is
trivial).
(b) If $X\ge j$ then some $j$-set $S\subseteq[m]$ has $I_i=1$ for all $i\in S$,
so $\Prob(X\ge j)\le\sum_{|S|=j}\prod_{i\in S}p_i\le(\sum_ip_i)^j/j!\le\mu^j/j!$,
because every product $\prod_{i\in S}p_i$ occurs $j!$ times in the expansion of
$(\sum_ip_i)^j$ and all terms are non-negative. Finally $j!\ge(j/e)^j$.
\end{proof}

\begin{citedresult}[{Markov's inequality and its two uses}]\label{s1:citMarkov}
If $X\ge0$ and $a>0$ then $\Prob(X\ge a)\le\Ex X/a$. Consequently:
\begin{enumerate}
\item[(a)] for $X\ge0$, $\Prob(X\le3\Ex X)\ge2/3$;
\item[(b)] for $X_1,X_2\ge0$, $\Prob(X_1\le4\Ex X_1\text{ and }X_2\le4\Ex X_2)\ge
1/2$;
\item[(c)] (a) and (b) hold with $\Prob$ and $\Ex$ replaced by the conditional
probability and expectation given any event of positive probability, or given
any fixed outcome of variables that are independent of the variables being
drawn.
\end{enumerate}
\deps{none (standard).}
\end{citedresult}
\begin{proof}[Derivation]
(a) If $\Ex X=0$ then $X=0$ almost surely. Otherwise
$\Prob(X>3\Ex X)\le\Prob(X\ge3\Ex X)\le1/3$.
(b) Similarly $\Prob(X_i>4\Ex X_i)\le1/4$ for $i=1,2$, and a union bound gives
failure probability at most $1/2$.
(c) The same argument in the conditional probability space; if the conditioning
fixes variables independent of those being drawn, the conditional law of the
latter is their unconditional law.
\end{proof}

\begin{lemma}[Lemma BBD: Bernstein inequality for bounded differences]\label{s1:lemBBD}
Let $m\ge0$ be an integer, let $p_1,\dots,p_m\in[0,1]$, and let
$I_1,\dots,I_m$ be independent random variables with $\Prob(I_k=1)=p_k$ and
$\Prob(I_k=0)=1-p_k$; write $I\defeq(I_1,\dots,I_m)$. Let
$\Psi\colon\{0,1\}^m\to\mathbb{R}$, and let $b\ge0$ and
$c_1,\dots,c_m\in[0,b]$ be such that $|\Psi(x)-\Psi(x')|\le c_k$ whenever
$x,x'\in\{0,1\}^m$ differ only in the $k$-th coordinate. Let $\beta>0$ with
$\beta\ge\sum_{k=1}^mp_k(1-p_k)c_k^2$. Then for every $a\ge0$,
\[
\Prob\bigl(\Psi(I)\ge\Ex\Psi(I)+a\bigr)\le\exp\Bigl(-\frac{a^2}{2(\beta+ba/3)}\Bigr)
\quad\text{and}\quad
\Prob\bigl(\Psi(I)\le\Ex\Psi(I)-a\bigr)\le\exp\Bigl(-\frac{a^2}{2(\beta+ba/3)}\Bigr).
\]
\deps{none (standard).}
\end{lemma}
\begin{proof}
This is a standard Bernstein-type inequality; we prove it by the
exponential-moment method. All sums in this proof are finite. For
$0\le k\le m$, $x=(x_1,\dots,x_k)\in\{0,1\}^k$ and
$z=(z_{k+1},\dots,z_m)\in\{0,1\}^{m-k}$ put
\[
\mathrm{wt}_k(x)\defeq\prod_{j=1}^{k}\bigl(p_jx_j+(1-p_j)(1-x_j)\bigr),\qquad
\mathrm{wt}'_k(z)\defeq\prod_{j=k+1}^{m}\bigl(p_jz_j+(1-p_j)(1-z_j)\bigr)
\]
(empty products are $1$), and write $(x,z)\in\{0,1\}^m$ for the
concatenation. These weights are non-negative,
$\sum_{x\in\{0,1\}^k}\mathrm{wt}_k(x)=\prod_{j\le k}\bigl(p_j+(1-p_j)\bigr)=1$, and
likewise $\sum_{z\in\{0,1\}^{m-k}}\mathrm{wt}'_k(z)=1$. By independence,
$\Prob(I=x)=\mathrm{wt}_m(x)$ for $x\in\{0,1\}^m$. Hence
$\Prob(I\in S)=\sum_{x\in S}\mathrm{wt}_m(x)$ for $S\subseteq\{0,1\}^m$, and
$\Ex g(I)=\sum_{x}\mathrm{wt}_m(x)g(x)$ for $g\colon\{0,1\}^m\to\mathbb{R}$.

\emph{Step 1: partial averages.} For $0\le k\le m$ and $x\in\{0,1\}^k$ put
\[
\Psi_k(x)\defeq\sum_{z\in\{0,1\}^{m-k}}\mathrm{wt}'_k(z)\,\Psi(x,z).
\]
Then
$\Psi_m=\Psi$, and, as $\mathrm{wt}'_0=\mathrm{wt}_m$, $\Psi_0$ is the constant
$\sum_x\mathrm{wt}_m(x)\Psi(x)=\Ex\Psi(I)$. Let $1\le k\le m$ and
$y\in\{0,1\}^{k-1}$. Split the sum defining $\Psi_{k-1}(y)$ according to
the first coordinate $z_k$ of $z=(z_k,z')$, and use
$\mathrm{wt}'_{k-1}(z_k,z')=\bigl(p_kz_k+(1-p_k)(1-z_k)\bigr)\mathrm{wt}'_k(z')$. This gives
\begin{equation}\label{s1:eqBBDavg}
\Psi_{k-1}(y)=p_k\Psi_k(y,1)+(1-p_k)\Psi_k(y,0).
\end{equation}
Put $\Xi_k(y)\defeq\Psi_k(y,1)-\Psi_k(y,0)
=\sum_{z}\mathrm{wt}'_k(z)\bigl(\Psi(y,1,z)-\Psi(y,0,z)\bigr)$. The points
$(y,1,z)$ and $(y,0,z)$ differ only in the $k$-th coordinate, so
$|\Xi_k(y)|\le\sum_z\mathrm{wt}'_k(z)c_k=c_k$. By~\eqref{s1:eqBBDavg},
\begin{equation}\label{s1:eqBBDdiff}
\Psi_k(y,x_k)-\Psi_{k-1}(y)=(x_k-p_k)\,\Xi_k(y)\qquad(x_k\in\{0,1\}).
\end{equation}

\emph{Step 2: an elementary inequality.} For every real $u$,
\begin{equation}\label{s1:eqBBDphi}
\Bigl(1-\frac u3\Bigr)e^u\le1+\frac{2u}3+\frac{u^2}6 .
\end{equation}
Indeed, let $\chi(u)$ be the right side minus the left side. Then
$\chi(0)=0$ and
$\chi'(u)=\frac23+\frac u3-\bigl(\frac23-\frac u3\bigr)e^u$, so
$\chi'(0)=0$; and $\chi''(u)=\frac13\bigl(1-(1-u)e^u\bigr)\ge0$,
because $1-u\le e^{-u}$. Hence $\chi'$ is non-decreasing, so
$\chi'\le0$ on $(-\infty,0]$ and $\chi'\ge0$ on $[0,\infty)$. Thus
$\chi$ is non-increasing on $(-\infty,0]$ and non-decreasing on
$[0,\infty)$, and $\chi\ge\chi(0)=0$. Now let $0\le\zeta<3$ and $u\le\zeta$.
Then $1-u/3\ge1-\zeta/3>0$ and
$1+\frac{2u}3+\frac{u^2}6=\bigl(1-\frac u3\bigr)(1+u)+\frac{u^2}2$. Dividing
\eqref{s1:eqBBDphi} by $1-u/3$ gives
\begin{equation}\label{s1:eqBBDexp}
e^u\le1+u+\frac{u^2}{2(1-u/3)}\le1+u+\frac{u^2}{2(1-\zeta/3)} .
\end{equation}

\emph{Step 3: one coordinate.} Let $\eta\ge0$ with $\eta b<3$ (it is chosen
in Step~5), and put $\zeta\defeq\eta b$. Let $1\le k\le m$ and
$y\in\{0,1\}^{k-1}$, and write $p\defeq p_k$ and $\Xi\defeq\Xi_k(y)$.
The numbers $\eta(1-p)\Xi$ and $-\eta p\Xi$ are at most
$\eta|\Xi|\le\eta c_k\le\eta b=\zeta$. Apply~\eqref{s1:eqBBDexp} to both,
multiply by $p\ge0$ and $1-p\ge0$ respectively, and add. Since
$p(1-p)-(1-p)p=0$ and $p(1-p)^2+(1-p)p^2=p(1-p)$, this gives
\begin{equation}\label{s1:eqBBDstep}
p\,e^{\eta(1-p)\Xi}+(1-p)\,e^{-\eta p\Xi}
\le1+\frac{\eta^2p(1-p)\Xi^2}{2(1-\zeta/3)}
\le\exp\Bigl(\frac{\eta^2p_k(1-p_k)c_k^2}{2(1-\zeta/3)}\Bigr),
\end{equation}
where the last step uses $\Xi^2\le c_k^2$ and $1+v\le e^v$.

\emph{Step 4: induction over the coordinates.} For $0\le k\le m$ put
\[
\mathcal{E}_k\defeq\sum_{x\in\{0,1\}^k}\mathrm{wt}_k(x)\exp\bigl(\eta(\Psi_k(x)-\Psi_0)\bigr).
\]
Then $\mathcal{E}_0=1$, and $\mathcal{E}_m=\Ex\exp\bigl(\eta(\Psi(I)-\Ex\Psi(I))\bigr)$.
Let $1\le k\le m$, and write $x\in\{0,1\}^k$ as $(y,x_k)$ with
$y\in\{0,1\}^{k-1}$; then $\mathrm{wt}_k(y,1)=p_k\mathrm{wt}_{k-1}(y)$ and
$\mathrm{wt}_k(y,0)=(1-p_k)\mathrm{wt}_{k-1}(y)$. By~\eqref{s1:eqBBDdiff}, and then
by~\eqref{s1:eqBBDstep} and the non-negativity of all weights,
\begin{align*}
\mathcal{E}_k&=\sum_{y\in\{0,1\}^{k-1}}\mathrm{wt}_{k-1}(y)\,e^{\eta(\Psi_{k-1}(y)-\Psi_0)}
\Bigl(p_ke^{\eta(1-p_k)\Xi_k(y)}+(1-p_k)e^{-\eta p_k\Xi_k(y)}\Bigr)\\
&\le\mathcal{E}_{k-1}\exp\Bigl(\frac{\eta^2p_k(1-p_k)c_k^2}{2(1-\zeta/3)}\Bigr).
\end{align*}
By induction on $k$, and since $\sum_kp_k(1-p_k)c_k^2\le\beta$,
$\mathcal{E}_m\le\exp\bigl(\eta^2\beta/(2(1-\zeta/3))\bigr)$.

\emph{Step 5: the upper tail.} Let
$S\defeq\{x\in\{0,1\}^m:\Psi(x)\ge\Psi_0+a\}$. For $x\in S$ we have
$\exp\bigl(\eta(\Psi(x)-\Psi_0-a)\bigr)\ge1$, as $\eta\ge0$, and all terms
below are non-negative. Hence
\begin{align*}
\Prob\bigl(\Psi(I)\ge\Ex\Psi(I)+a\bigr)&=\sum_{x\in S}\mathrm{wt}_m(x)
\le e^{-\eta a}\sum_{x\in\{0,1\}^m}\mathrm{wt}_m(x)e^{\eta(\Psi(x)-\Psi_0)}\\
&=e^{-\eta a}\mathcal{E}_m
\le\exp\Bigl(-\eta a+\frac{\eta^2\beta}{2(1-\eta b/3)}\Bigr).
\end{align*}
Now choose $\eta\defeq a/(\beta+ba/3)\ge0$. Since $\beta>0$, we have
$\eta b/3=(ba/3)/(\beta+ba/3)<1$, so $\eta b<3$ as required in Step~3, and
$1-\eta b/3=\beta/(\beta+ba/3)$. Therefore
\[
-\eta a+\frac{\eta^2\beta}{2(1-\eta b/3)}
=-\frac{a^2}{\beta+ba/3}+\frac{a^2}{(\beta+ba/3)^2}\cdot\frac{\beta\,(\beta+ba/3)}{2\beta}
=-\frac{a^2}{2(\beta+ba/3)},
\]
which is the first bound.

\emph{Step 6: the lower tail.} The function $-\Psi$ satisfies the hypotheses
with the same $b$, $c_1,\dots,c_m$ and $\beta$, and
$\Ex(-\Psi(I))=-\Ex\Psi(I)$. So the first bound for $-\Psi$ is the second
bound for $\Psi$.
\end{proof}

\begin{citedresult}[{Hall's theorem \cite{Hal35}}]\label{s1:citHall}
Let $\mathcal J$ be a finite set and, for every $a\in\mathcal J$, let
$\mathcal T(a)$ be a finite set. If
$\bigl|\bigcup_{a\in\mathcal J'}\mathcal T(a)\bigr|\ge|\mathcal J'|$ for every
$\mathcal J'\subseteq\mathcal J$, then there is an injective map $\chi$ from
$\mathcal J$ to $\bigcup_{a\in\mathcal J}\mathcal T(a)$ with
$\chi(a)\in\mathcal T(a)$ for every $a\in\mathcal J$ (a system of distinct
representatives).
\deps{none (standard).}
\end{citedresult}

\begin{citedresult}[{Haxell's condition for matchability \cite{Hax95}}]\label{s1:citHaxell}
Let $q\ge1$ be an integer, and let $\mathcal H$ be a hypergraph (a set of
subsets, called edges, of a finite vertex set) whose vertex set is the
disjoint union of two sets $\mathcal X$ and $\mathcal Y$, such that every edge
$e$ satisfies $|e\cap\mathcal X|=1$ and $|e\cap\mathcal Y|\le q$. Suppose that
for every $\mathcal X'\subseteq\mathcal X$ and every $Z\subseteq\mathcal Y$ with
$|Z|\le(2q-1)(|\mathcal X'|-1)$ there is an edge $e$ of $\mathcal H$ with
$e\cap\mathcal X\subseteq\mathcal X'$ and $e\cap Z=\emptyset$. Then
$\mathcal H$ has an \emph{$\mathcal X$-saturating matching}, i.e.\ a set of
pairwise disjoint edges such that every vertex of $\mathcal X$ lies in one of
them.

For $\mathcal X'=\emptyset$ there is no such $Z$, because $(2q-1)(0-1)<0$. For
$q=1$ (edges with at most two vertices) the result is Hall's theorem; a vertex
$a$ with $\{a\}\in\mathcal H$ is matched by that edge. In the later literature
the result of \cite{Hax95} is quoted for hypergraphs whose edges have at most
(in some sources: exactly) $r$ vertices, one of them in $\mathcal X$, with the
bound $(2r-3)(|\mathcal X'|-1)$. With $r=q+1$, the maximum edge size, this bound
is $(2(q+1)-3)(|\mathcal X'|-1)=(2q-1)(|\mathcal X'|-1)$: the stated form agrees
with Haxell's theorem as quoted there, and the case $q=1$ (Hall's theorem)
confirms this reading of the parameters. The only application, in the proof of
Lemma~\ref{s3:lemL9rho}, verifies the hypothesis for every $Z$ with
$|Z|<q^2|\mathcal X'|$; so that application stays valid if $2q-1$ is replaced
by any positive factor at most $q^2$. Unlike the other cited items of this
subsection, which are standard textbook results, this is a published research
result. The derivation below shows
that the stated form follows from the special case in which every edge meets
$\mathcal Y$ in exactly $q$ vertices; so the stated form holds if \cite{Hax95}
proves either of the two quoted forms. After the derivation, a complete proof of
the stated form is written out (Haxell's alternating-tree argument), so that
the application does not depend on the precise form in which the result is
stated in \cite{Hax95}.
\deps{none (published).}
\end{citedresult}
\begin{proof}[Derivation from the case $|e\cap\mathcal Y|=q$]
Let $\mathcal H$ be as stated, and let $\mathcal X_0$ be the set of
$a\in\mathcal X$ such that $\{a\}$ is an edge. Every edge $e$ with
$e\cap\mathcal X\not\subseteq\mathcal X_0$ has $e\cap\mathcal Y\ne\emptyset$
(otherwise $e=\{a\}$ with $a\in\mathcal X_0$). For each such $e$ take a set
$N_e$ of $q-|e\cap\mathcal Y|$ new vertices, the sets $N_e$ pairwise disjoint
and disjoint from $\mathcal X\cup\mathcal Y$, and put $e^+:=e\cup N_e$. Let
$\mathcal H^+$ be the hypergraph on
$(\mathcal X\setminus\mathcal X_0)\sqcup\mathcal Y^+$, where
$\mathcal Y^+:=\mathcal Y\cup\bigcup_eN_e$, whose edges are these sets $e^+$.
Every edge of $\mathcal H^+$ meets $\mathcal X\setminus\mathcal X_0$ in exactly
one vertex and $\mathcal Y^+$ in exactly $q$ vertices.

We check the hypothesis for $\mathcal H^+$. Let
$\mathcal X'\subseteq\mathcal X\setminus\mathcal X_0$ and
$Z^+\subseteq\mathcal Y^+$ with $|Z^+|\le(2q-1)(|\mathcal X'|-1)$. For every
$z\in Z^+\setminus\mathcal Y$ let $f_z$ be the unique edge of $\mathcal H$ with
$z\in N_{f_z}$, and choose a vertex $b_z\in f_z\cap\mathcal Y$. Put
$Z:=(Z^+\cap\mathcal Y)\cup\{b_z:z\in Z^+\setminus\mathcal Y\}\subseteq\mathcal Y$.
Then $|Z|\le|Z^+|$, so the hypothesis for $\mathcal H$ gives an edge $e$ with
$e\cap\mathcal X\subseteq\mathcal X'$ and $e\cap Z=\emptyset$. As
$|e\cap\mathcal X|=1$, the set $e\cap\mathcal X$ is a single vertex of
$\mathcal X'\subseteq\mathcal X\setminus\mathcal X_0$; so
$e\cap\mathcal X\not\subseteq\mathcal X_0$, the set $e^+$ is an edge of
$\mathcal H^+$, and
$e^+\cap(\mathcal X\setminus\mathcal X_0)=e\cap\mathcal X\subseteq\mathcal X'$.
Moreover $e^+\cap Z^+=\emptyset$: a vertex of $e^+\cap Z^+$ lies either in
$e\cap\mathcal Y$, and then in $Z\cap e$; or it is some $z\in N_e$, and then
$f_z=e$ and $b_z\in e\cap Z$. Both contradict $e\cap Z=\emptyset$.

By the case $|e\cap\mathcal Y|=q$, $\mathcal H^+$ has an
$(\mathcal X\setminus\mathcal X_0)$-saturating matching
$\{e_1^+,\dots,e_k^+\}$. The edges $e_1,\dots,e_k$ of $\mathcal H$ are pairwise
disjoint, as $e_i\subseteq e_i^+$, and they cover
$\mathcal X\setminus\mathcal X_0$, as
$e_i\cap\mathcal X=e_i^+\cap(\mathcal X\setminus\mathcal X_0)$. The edges
$\{a\}$, $a\in\mathcal X_0$, are disjoint from each other and from every $e_i$,
since $e_i\cap\mathcal X\subseteq\mathcal X\setminus\mathcal X_0$. So
$\{e_1,\dots,e_k\}\cup\{\{a\}:a\in\mathcal X_0\}$ is an $\mathcal X$-saturating
matching of $\mathcal H$.
\end{proof}
\begin{proof}[Proof of the stated form]
This is Haxell's alternating-tree argument. For an edge $e$ let $a(e)$ denote
the vertex of $e\cap\mathcal X$. We use induction on $|\mathcal X|$. If
$\mathcal X=\emptyset$, the empty set of edges is an $\mathcal X$-saturating
matching. Let $|\mathcal X|\ge1$ and fix $a_0\in\mathcal X$. The edges of
$\mathcal H$ that do not contain $a_0$ form a hypergraph on
$(\mathcal X\setminus\{a_0\})\sqcup\mathcal Y$ that satisfies the hypothesis
with $\mathcal X\setminus\{a_0\}$ in place of $\mathcal X$: for
$\mathcal X'\subseteq\mathcal X\setminus\{a_0\}$, the edge $e$ given by the
hypothesis for $\mathcal H$ has $e\cap\mathcal X\subseteq\mathcal X'$, so
$a_0\notin e$. By the induction hypothesis this hypergraph has an
$(\mathcal X\setminus\{a_0\})$-saturating matching.

Call a set $\mathcal M$ of pairwise disjoint edges of $\mathcal H$ \emph{good} if
no edge of $\mathcal M$ contains $a_0$ and every vertex of
$\mathcal X\setminus\{a_0\}$ lies in an edge of $\mathcal M$. Since every edge
meets $\mathcal X$ in exactly one vertex, the map $m\mapsto a(m)$ is then a
bijection from $\mathcal M$ onto $\mathcal X\setminus\{a_0\}$. We have just seen
that a good set exists. It suffices to find a good $\mathcal M$ and an edge $x$
with $a(x)=a_0$ and $x\cap m\cap\mathcal Y=\emptyset$ for every
$m\in\mathcal M$: then $x$ is disjoint from every $m\in\mathcal M$ (as
$a(m)\ne a_0$), and $\mathcal M\cup\{x\}$ is an $\mathcal X$-saturating
matching.

\emph{Configurations.} Let $\mathcal M$ be good. A \emph{configuration} for
$\mathcal M$ is a sequence $(x_1,\mathcal S_1,\dots,x_k,\mathcal S_k)$ with
$k\ge0$, of edges $x_i$ of $\mathcal H$ and sets
$\mathcal S_i\subseteq\mathcal M$, with the following properties. Put
$\mathcal X^{(0)}:=\{a_0\}$, $Z^{(0)}:=\emptyset$ and, for $1\le i\le k$,
\[
\mathcal X^{(i)}:=\{a_0\}\cup\{a(m):m\in\mathcal S_1\cup\dots\cup\mathcal S_i\},
\qquad
Z^{(i)}:=\mathcal Y\cap\Bigl(x_1\cup\dots\cup x_i\cup
\bigcup_{m\in\mathcal S_1\cup\dots\cup\mathcal S_i}m\Bigr).
\]
The properties are, for every $1\le i\le k$: (C1) $a(x_i)\in\mathcal X^{(i-1)}$ and
$x_i\cap Z^{(i-1)}=\emptyset$; (C2)
$\mathcal S_i=\{m\in\mathcal M:m\cap x_i\cap\mathcal Y\ne\emptyset\}$ and
$\mathcal S_i\ne\emptyset$. The \emph{signature} of the configuration is the
sequence $(|\mathcal S_1|,\dots,|\mathcal S_k|)$ of positive integers.

The sets $\mathcal S_i$ are pairwise disjoint: if $j<i$ and
$m\in\mathcal S_j$, then $m\cap\mathcal Y\subseteq Z^{(i-1)}$, which $x_i$
avoids by (C1), so $m\notin\mathcal S_i$. As $m\mapsto a(m)$ is injective on
$\mathcal M$ and $a(m)\ne a_0$, we get
$|\mathcal X^{(k)}|-1=\sum_{i\le k}|\mathcal S_i|\le|\mathcal M|=|\mathcal X|-1$.
Each $m\in\mathcal S_i$ shares a vertex of $\mathcal Y$ with $x_i$, so
$|(m\cap\mathcal Y)\setminus x_i|\le q-1$. Hence, using $|\mathcal S_i|\ge1$,
\begin{align*}
|Z^{(k)}|&\le\sum_{i\le k}\Bigl(|x_i\cap\mathcal Y|
+\sum_{m\in\mathcal S_i}|(m\cap\mathcal Y)\setminus x_i|\Bigr)
\le\sum_{i\le k}\bigl(q+(q-1)|\mathcal S_i|\bigr)\\
&\le(2q-1)\sum_{i\le k}|\mathcal S_i|=(2q-1)\bigl(|\mathcal X^{(k)}|-1\bigr).
\end{align*}
So the hypothesis, with $\mathcal X':=\mathcal X^{(k)}$ and $Z:=Z^{(k)}$, gives
an edge $x$ with $a(x)\in\mathcal X^{(k)}$ and $x\cap Z^{(k)}=\emptyset$. $(*)$

\emph{Swap step.} Suppose that $\mathcal M$ is good, that
$(x_1,\mathcal S_1,\dots,x_j,\mathcal S_j)$ is a configuration for
$\mathcal M$, and that $y$ is an edge with $a(y)\in\mathcal X^{(j)}$,
$y\cap Z^{(j)}=\emptyset$ and $y\cap m\cap\mathcal Y=\emptyset$ for every
$m\in\mathcal M$. If $a(y)=a_0$, then $\mathcal M$ and $x:=y$ are as required,
and the swap step ends. Otherwise $a(y)=a(m)$ for exactly one $m\in\mathcal M$, and
$m\in\mathcal S_i$ for exactly one $i\le j$. Put
$\mathcal M':=(\mathcal M\setminus\{m\})\cup\{y\}$. It is good: $y$ shares no
vertex of $\mathcal Y$ with an edge of $\mathcal M$, the vertex $a(y)=a(m)$
lies in no edge of $\mathcal M\setminus\{m\}$, $a_0\notin y$, and $a(m)$ is
covered by $y$. For $i'\le j$ we have $x_{i'}\cap\mathcal Y\subseteq Z^{(j)}$,
so $y\cap x_{i'}\cap\mathcal Y=\emptyset$; and $m\in\mathcal S_{i'}$ only for
$i'=i$. Hence
$\{m'\in\mathcal M':m'\cap x_{i'}\cap\mathcal Y\ne\emptyset\}$ equals
$\mathcal S_{i'}$ for $i'<i$ and $\mathcal S_i\setminus\{m\}$ for $i'=i$, and
the sets $\mathcal X^{(i')}$ and $Z^{(i')}$ with $i'<i$ are the same for
$\mathcal M'$ as for $\mathcal M$. So, if $\mathcal S_i\ne\{m\}$, then
$(x_1,\mathcal S_1,\dots,x_{i-1},\mathcal S_{i-1},x_i,\mathcal S_i\setminus\{m\})$
is a configuration for $\mathcal M'$, and the swap step ends. If $\mathcal S_i=\{m\}$,
then $(x_1,\mathcal S_1,\dots,x_{i-1},\mathcal S_{i-1})$ is a configuration for
$\mathcal M'$, the edge $x_i$ satisfies $a(x_i)\in\mathcal X^{(i-1)}$ and
$x_i\cap Z^{(i-1)}=\emptyset$ by (C1), and $x_i$ shares no vertex of
$\mathcal Y$ with an edge of $\mathcal M'$; we repeat the swap step with
$\mathcal M'$, this configuration and $y:=x_i$. As $i-1<j$, and as
$\mathcal X^{(0)}=\{a_0\}$, the swap step ends after at most $j+1$ rounds.

\emph{One move.} Let $\mathcal M$ be good, with a configuration of signature
$(s_1,\dots,s_k)$, and let $x$ be as in $(*)$. If $x$ shares a vertex of
$\mathcal Y$ with some edge of $\mathcal M$, append $x$ and
$\mathcal S_{k+1}:=\{m\in\mathcal M:m\cap x\cap\mathcal Y\ne\emptyset\}$; by
$(*)$ this is a configuration for $\mathcal M$, of signature
$(s_1,\dots,s_k,s_{k+1})$. Otherwise apply the swap step with $j:=k$ and
$y:=x$. Either it finds $\mathcal M$ and $x$ as required, or it ends with a
good $\mathcal M'$ and a configuration for $\mathcal M'$ of signature
$(s_1,\dots,s_{i-1},s_i-1)$ for some $i\le k$ with $s_i\ge2$.

\emph{Termination.} For finite sequences $\sigma,\tau$ of positive integers
write $\sigma\prec\tau$ if the word $\sigma\infty$ precedes the word
$\tau\infty$ in the lexicographic order on words over
$\{1,2,\dots\}\cup\{\infty\}$. As $\infty$ occurs only as the last letter,
neither of two distinct such words is a prefix of the other, so $\prec$ is a
strict total order. A move that does not find $\mathcal M$ and $x$ as required
replaces the signature $\sigma$
by some $\sigma'\prec\sigma$: an extension $\sigma'$ of $\sigma$ has a finite
entry where $\sigma\infty$ has $\infty$, and $(s_1,\dots,s_{i-1},s_i-1)$ is
smaller at position $i$. Every signature has positive entries with sum at most
$|\mathcal X|-1$, so there are finitely many signatures, and a sequence that
is strictly decreasing for $\prec$ has no repetition; hence it is finite.
Start with a good $\mathcal M$ and the configuration with $k=0$, and make
moves as long as possible. After finitely many moves a swap step finds a good
$\mathcal M$ and an edge $x$ as required.
\end{proof}

\begin{citedresult}[{$T$-joins and Euler circuits}]\label{s1:citEuler}
\begin{enumerate}
\item[(a)] Let $H$ be a connected graph and $T\subseteq V(H)$ with $|T|$ even.
Every spanning tree $S$ of $H$ contains a $T$-join, i.e.\ a set $J\subseteq E(S)$
such that the vertices of odd $J$-degree are exactly the vertices of $T$.
\item[(b)] A connected graph (or multigraph, possibly with loops) with at least one
edge in which every vertex has even degree has a closed trail using every edge
exactly once.
\item[(c)] A graph in which every vertex has even degree is an edge-disjoint union
of cycles; a directed graph without loops in which every vertex has in-degree
equal to out-degree is an arc-disjoint union of directed cycles.
\end{enumerate}
The Lov\'asz path decomposition is used only through Cited
result~\ref{s1:citCor22}.
\deps{none (standard).}
\end{citedresult}
\begin{proof}[Derivation of (a) and (c)]
(a) Root $S$ at a vertex $\rho$, and put the edge from a vertex $v\ne\rho$ to its
parent into $J$ if and only if the subtree of $S$ rooted at $v$ contains an odd
number of vertices of $T$. For $v\ne\rho$ the $J$-degree of $v$ is congruent
modulo $2$ to the number of $T$-vertices in its subtree minus those in the
subtrees of its children, i.e.\ to $\ind{v\in T}$; for $\rho$ it is congruent to
$|T|-\ind{\rho\in T}\equiv\ind{\rho\in T}$, since $|T|$ is even. (b) is Euler's
theorem. (c) If some edge (arc) remains, start at one of its ends and walk along
unused edges (leaving along unused out-arcs); since every vertex has even degree
(balanced degrees), the walk can always continue until it revisits a vertex,
which closes a cycle (directed cycle); delete it and repeat.
\end{proof}

\subsection{Constants, their order, and the galactic conditions}\label{s1:ssecConstants}

\begin{definition}[the constants]\label{s1:defConstants}
\begin{enumerate}
\item[(i)] $\eps\defeq2^{-5}$, $\sigma\defeq100$, $\Cp\defeq103$ and
$\Aexp\defeq105$.
\item[(ii)] $\Nzero$ is an absolute constant, i.e.\ it depends only on $\sigma$,
$\Cp$ and $\Aexp$, with $\Nzero\ge2^{40}$ (redundant: size condition~(i),
$L\ge2^{10}$, already forces every $N\ge\Nzero$ to satisfy $N\ge2^{1024}$,
i.e.\ $\lceil\Nzero\rceil\ge2^{1024}$) and $\Nzero$ at least
each of the absolute size thresholds required in the proofs of
Lemma~\ref{s4:lemTPV}, Lemma~\ref{s4:lemPV} and Theorem~\ref{s4:thmVXp}.
Every requirement on $\Nzero$ in this item is an \emph{eventuality}: it holds
for every sufficiently large value of $\Nzero$. Indeed $\Nzero\ge2^{40}$ is a
lower bound, and every other requirement asks that an explicit inequality in
$N$, which holds for all sufficiently large $N$, hold for every $N\ge\Nzero$. So
such an $\Nzero$ exists (finitely many eventualities have a common threshold),
and only its existence is used; no numerical value of $\Nzero$ enters any proof.
(Each of these three statements carries the hypothesis that its vertex set has
at least $\Nzero$ vertices, and its proof records the size conditions it
needs. The lemmas of Section~\ref{s5:sec} carry the hypothesis $n\ge\Nzero$
through the setting paragraph of that section and use $\Nzero$ only to apply
these three statements. Lemma~\ref{s3:lemL15p} and Theorem~\ref{s3:thmT16s}
need no lower bound on the number of vertices.)
\item[(iii)] $\Dstar$ is a constant satisfying \ref{s1:condG1}--\ref{s1:condG4}
of Condition~\ref{s1:condGamma}. Each of these conditions is an eventuality in
$\Dstar$: once $\Nzero$ is fixed, every sufficiently large $\Dstar$ satisfies all
of them (Lemma~\ref{s7:lemGammaSat}). So such a constant exists; only its
existence is used, and no numerical value of $\Dstar$ enters any proof.
\item[(iv)] $\epsOne(\Dstar)$ and $\epsTwo(\Dstar)$ are the explicit functions
of $\Dstar$ defined in Proposition~\ref{s7:propCost} and
Lemma~\ref{s7:lemUHsplit} respectively, and $\thetaQ\defeq\epsTwo(\Dstar)$. Each is a finite sum of explicit terms
depending only on $\Dstar$, $\eps$, $\sigma$, $\Cp$ and $\Aexp$.
\item[(v)] $\Czero\defeq\Dstar/2+\cTot+\epsOne(\Dstar)$; the function
$\epsOne$ is counted exactly once. By \ref{s1:condG4},
$\Czero\le\Dstar/2+\cTotR$.
\item[(vi)] $\cEG\defeq\max\bigl\{\Czero/(1-2\thetaQ),\ \Nzero/2\bigr\}$.
\item[(vii)] (Generic notation, proofs only; permitted, but not used in the
present text.) Inside a proof, $\epsD$ may be
used as shorthand for an expression written out explicitly in the same proof
that depends only on $\Dstar$, $\eps$, $\sigma$, $\Cp$, $\Aexp$ and tends to $0$ as
$\Dstar\to\infty$. It never appears in the final bound of a statement; every
named small quantity ($\epsA$, $\epsU$, $\epsK=\epsChain$, $\epsCONC$,
$\varepsilon_M$, $\epsX$, $\epsOne$, $\epsTwo$) is given by an explicit formula
in its defining statement.
\end{enumerate}
\deps{none.}
{\small\noindent(The references of this definition to later statements are not
dependencies; see the paragraph on forward references in
Section~\ref{s1:ssecDAG}. For instance, $\epsOne$ and $\epsTwo$ are defined in
Proposition~\ref{s7:propCost} and Lemma~\ref{s7:lemUHsplit}.)\par}
\end{definition}

\begin{remark}[order of the constants]\label{s1:remOrder}
The constants are fixed in the following order:
\begin{enumerate}
\item $\sigma$, $\Cp$, $\Aexp$, $\eps$;
\item $\Nzero$;
\item $\Dstar$;
\item $\epsOne(\Dstar)$, $\epsTwo(\Dstar)=\thetaQ$ and $\Czero$;
\item $\cEG$.
\end{enumerate}
Nothing earlier in this list depends on anything later. In particular $\Czero$
and $\thetaQ$ depend on $\Dstar$ only (given the fixed $\sigma,\Cp,\Aexp,\eps$):
they do not depend on $G$, on $n$, on the run of the hierarchy, on the
designation, or on $\cEG$. The constant $\cEG$ enters only as the constant $c$
of the induction, Theorem~\ref{s7:thmHI}, which is applied last; this is why
that induction is not circular. The constants of steps~2 and~3 are only required to be sufficiently
large: every condition on $\Nzero$ (Definition~\ref{s1:defConstants}(ii)) and on
$\Dstar$ (Condition~\ref{s1:condGamma}) is an eventuality in that constant, given
the constants of the earlier steps. Table~\ref{s1:tabOrder}, after
Condition~\ref{s1:condGamma}, lists every such condition with the variable in
which it is eventual, the constants it may depend on, and the values at which it
is used.
\deps{\ref{s1:defConstants}.}
\end{remark}

\begin{condition}[the galactic conditions on $\Dstar$]\label{s1:condGamma}
The constant $\Dstar$ satisfies the following conditions. This is the single list
of galactic conditions of the paper; every use cites the specific item.
Every condition is an \emph{eventuality} in $\Dstar$: once $\Nzero$ is fixed, it
holds for all sufficiently large $\Dstar$ (Lemma~\ref{s7:lemGammaSat}). Only the
existence of a suitable $\Dstar$ is used: no proof evaluates any of the
inequalities below at a particular numerical value.
\begin{enumerate}
\item[\Gc{1}]\namedlabel{s1:condG1}{\Gc{1}} (\emph{Tower condition.})
$\Dstar>2$, so that $\log_2\log_2\Dstar$ is defined and positive; and for every real
$\mu\ge\log_2\log_2\Dstar$, with $\lambda\defeq2^\mu$, the following hold:
\begin{enumerate}
\item[(a)] $\mu\ge2^8$;
\item[(b)] $2^\mu\ge2^{14}\Aexp\mu^3$;
\item[(c)] $2\Aexp\log_2(\Aexp\mu)\le1.6\,\mu$;
\item[(d)] $2\log_2\mu+8\le\mu$ (this follows from (a), and is listed
separately for direct citation);
\item[(e)]\namedlabel{s1:condGstar}{\Gstar} (the condition \Gstar)
$\lambda^{36}\ge2^{240}(\Aexp\mu)^{46\Aexp}$;
\item[(f)] for every row of the COL-JV table (Table~\ref{s3:tabCOLJV}), every
inequality in column~3 of that row holds at $\lambda$, where
$\bar M\defeq(\Aexp\mu)^{2\Aexp}$ and $\bar k\defeq192\lambda^3+\frac43\bar M^2$ as in
the caption of that table. (Row~13 is \Gstar, and the inequality of row~12 is
item~(c). This item points forward to explicit formulas, as \Gc{4} does for
$\epsOne$ and $\epsTwo$: column~3 consists of $18$ explicit inequalities, two in
row~1, five in row~8 and one in each other row, and they involve only $\mu$,
$\lambda=2^\mu$, $\bar M$ and $\bar k$.)
\end{enumerate}
Each item consists of explicit inequalities in the single real variable $\mu$,
and each of them holds for all sufficiently large $\mu$
(Lemma~\ref{s7:lemGammaSat}); \Gc{1} asks that
all of them hold on the whole ray $\mu\ge\log_2\log_2\Dstar$. So \Gc{1} is upward
closed in $\Dstar$. (The requirement $\Dstar>2$ only makes the ray well
defined: whenever $\log_2\log_2\Dstar$ is defined, (a) at
$\mu=\log_2\log_2\Dstar$ already gives $\Dstar\ge2^{2^{256}}$.)

Every round $r\le R$ of the hierarchy has average degree $d_r\ge\Dstar$
(Definition~\ref{s2:defHBtp}), so $\lambda_r\defeq\log_2d_r$ satisfies
$\log_2\lambda_r\ge\log_2\log_2\Dstar$, and (a)--(f) hold at
$\mu=\log_2\lambda_r$. They also hold at $\mu=\log_2\log_2d$ for every
$d\ge\Dstar$, and at $\mu=\log_2\Dstar$ (as $t\ge\log_2t$ for $t>0$). By (a) at
$\mu=\log_2\log_2\Dstar$, $\log_2\Dstar\ge2^{256}$. The consequences used later,
each proved where it is used, are: every requirement (column~2) and every
column-3 inequality of the COL-JV table (Lemma~\ref{s3:lemCOLJV}(ii)); the tower
facts of Proposition~\ref{s2:propDegRec}, Lemma~\ref{s2:lemLacunary} and
Lemma~\ref{s2:lemTower} (among them $\thGC_r>\tau_r$ and
$\lambda_r\ge2^{2\logs d_r+2}$); and Proposition~\ref{s2:propStructure} and
elementary estimates in Sections~\ref{s5:sec}--\ref{s7:sec}, each of which cites
the item it uses.

\item[\Gc{2}]\namedlabel{s1:condG2}{\Gc{2}}
\begin{enumerate}
\item[(a)] $\Dstar\ge2^{117}$; hence $\log_2(\tHB_ld_l)\ge117$ in every round
$l$, the fixed-point requirement of the Lemma-25 size cap
(Lemma~\ref{s2:lemCap});
\item[(b)] $P_{l-2}/2\ge M_l\log^4M_l$ for every $3\le l\le R$, including rounds
with $M_l=2^{40}$ (light parts two rounds back are admissible parents,
Proposition~\ref{s2:propParentless});
\item[(c)] every per-ancestor failure probability of Lemma~\ref{s3:lemCOL} and
of Lemma~\ref{s5:lemE1} is at most $|V(Y)|^{-2}$ (the failure-union lines of
Table~\ref{s3:tabCOLJV}).
\end{enumerate}
Each of (a)--(c) is implied by \ref{s1:condG1}: (a) is
immediate from \ref{s1:condG1}; (b) is proved in Lemma~\ref{s2:lemTower}(b); and
the bounds of (c) are proved in Lemma~\ref{s3:lemCOL} (with rows~6 and~7 of
Table~\ref{s3:tabCOLJV}) and in Lemma~\ref{s5:lemE1}(b),(c). None of these
proofs, and none of the results they use, uses (b) or (c). They are listed separately for traceability.

\item[\Gc{3}]\namedlabel{s1:condG3}{\Gc{3}} $(\log_2\Dstar)^{\Cp}\ge2\Nzero$.

Hence every pre-part has at least $P_R\ge(\log\Dstar)^{\Cp}\ge2\Nzero$ vertices,
every light part at least $P_R/2\ge\Nzero$, every ancestor, part and TPV,
P-vortex or VX$^+$ run meets the absolute size thresholds of its tools, and every
$o(1)$ term in a per-part success probability is at most $1/100$ (each proved
where used).

\item[\Gc{4}]\namedlabel{s1:condG4}{\Gc{4}} $\epsOne(\Dstar)\le6$ and
$\thetaQ=\epsTwo(\Dstar)\le1/4$, where $\epsOne$ and $\epsTwo$ are the explicit
expressions of Proposition~\ref{s7:propCost} and Lemma~\ref{s7:lemUHsplit}.

These expressions collect the following terms, with the coefficients stated;
here $\epsOne$ contains $3\,\epsX$ and $\epsTwo$ contains $12\,\epsX$, where
$\epsX=\epsU+10.3/\Dstar+2.1(6\log_2\Dstar+12)/\Dstar$
(Lemma~\ref{s7:lemEXprime}):
\begin{itemize}
\item the concentration term $1.5\,\epsCONC(\Dstar)$ of $\epsOne$ (it does
not occur in $\epsTwo$), where
\[
\epsCONC(\Dstar)=2^{\sigma+15}\frac{\logs\Dstar}{\log\Dstar}
+200\,\eps\,\frac{\logs\Dstar}{\Cp\log\log\Dstar}
+\frac{4(2\logs\Dstar+2)}{(\log\Dstar)^{1/2}}
\]
(Theorem~\ref{s6:thmCONCL}(iv); its first two terms bound the $\tau$- and
$d^*$-terms of all ancestors as in Theorem~\ref{s6:thmCONC}(iii), with a spare
factor $2$, and its last term bounds the guest-cap terms $\thGC_r(Y^0)$ of the
light parts);
\item $\epsA=31\eps/(\Cp\log\log\Dstar)$ (defined in
Lemma~\ref{s2:lemTower}(e)), with coefficient $169$ in $\epsOne$ and $12$ in
$\epsTwo$;
\item $\epsU$, the demotion and lost-parent expectations
(Lemma~\ref{s5:lemExpect}), with coefficient $3$ in $\epsOne$ and $12$ in
$\epsTwo$ (through $\epsX$);
\item $\epsK$, the U-chaining term (Lemma~\ref{s5:lemKRED}), with
coefficient $1$ in $\epsOne$ (it does not occur in $\epsTwo$);
\item the pool term
$\sum_l(6\log_2M_l+12)/M_l\le2.1(6\log_2\Dstar+12)/\Dstar$, whose bound
enters with coefficient $3$ in $\epsOne$ and $12$ in $\epsTwo$ (through
$\epsX$);
\item in $\epsTwo$ only, the quotient term $2F(\log_2\Dstar)$, with
coefficient $4$ (so $\epsTwo$ contains $8F(\log_2\Dstar)$), where
$F(x)=(3184+30400\log_2x)\,x^{-90.2}$ (Lemma~\ref{s7:lemUHsplit}(iv));
\item terms $O(1/\Dstar)$: $(252+9+3\cdot10.3)/\Dstar=291.9/\Dstar$ in
$\epsOne$ and $(12\cdot10.3+4\cdot60)/\Dstar=363.6/\Dstar$ in $\epsTwo$.
\end{itemize}
\end{enumerate}

\emph{Remarks.} (1) \emph{Eventualities.} Each of
\ref{s1:condG1}--\ref{s1:condG4} holds for all sufficiently large $\Dstar$, once
$\Nzero$ is fixed: \ref{s1:condG1} because each of its finitely many items holds
for all sufficiently large $\mu$ (and its range $\mu\ge\log_2\log_2\Dstar$
shrinks as $\Dstar$ grows); \ref{s1:condG2}(a) and \ref{s1:condG3} because they
are lower bounds on $\Dstar$, and \ref{s1:condG2}(b),(c) because they follow
from \ref{s1:condG1}; and \ref{s1:condG4} because
$\epsOne(\Dstar)\to0$ and $\epsTwo(\Dstar)\to0$ as $\Dstar\to\infty$. This is
proved in Lemma~\ref{s7:lemGammaSat}. Hence every sufficiently large $\Dstar$
satisfies all four conditions, and only the existence of such a $\Dstar$ is
used.

(2) \emph{Order of choice; no circularity.} Table~\ref{s1:tabOrder} lists the
constants in the order in which they are chosen (Remark~\ref{s1:remOrder}), the
conditions on them, the variable in which each condition is an eventuality, and
the constants it may depend on. A condition on the constant of a step depends
only on constants of earlier steps: the size conditions, \ref{s1:condG1},
\ref{s1:condG2}(a) and \ref{s1:condG4} depend on absolute constants only, and
\ref{s1:condG3} depends on $\Nzero$, which is chosen before $\Dstar$. Every
condition is used only at values of its variable that are at least the chosen
threshold: the size conditions at $N=|Z|$ for the vertex set $Z$ of a vortex
run, which is a pre-part or a light part and so has $N\ge P_R/2\ge\Nzero$ by
\ref{s1:condG3}; the items of \ref{s1:condG1} at values
$\mu\ge\log_2\log_2\Dstar$, as described there; and \ref{s1:condG2}(a),
\ref{s1:condG3} and \ref{s1:condG4} at $\Dstar$ itself. The quantities of a
run ($d_r$, $\lambda_r$, $|V(Y)|$, $L_Y$) are not chosen: they are universally
quantified, and they exceed the thresholds because $d_r\ge\Dstar$ in every round
$r\le R$ and by \ref{s1:condG3}. The hypotheses $n\ge\Nzero$ and
$d_1\ge\Dstar$ on the graph are not consequences of these conditions: a graph
with $n<\Nzero$ or with $d_1<\Dstar$ is handled directly in steps~(2) and~(3) of
the proof of Theorem~\ref{s7:thmHI}, because $\cEG\ge\Nzero/2$ and
$\cEG\ge\Czero\ge\Dstar/2$. Finally
$\epsOne$, $\epsTwo=\thetaQ$ and $\Czero$ are functions of $\Dstar$, and $\cEG$
is defined last; it is used only as the constant $c$ of
Theorem~\ref{s7:thmHI}, and no condition on $\Nzero$ or $\Dstar$ involves it.
So no choice is circular.
\deps{\ref{s1:defConstants}, \ref{s1:remOrder}.}
\end{condition}

\begin{table}[htbp]
\begin{center}
\small
\begin{tabular}{|p{0.75cm}|p{4.6cm}|p{2.3cm}|p{2.2cm}|p{3.4cm}|}
\hline
Step & Constant; conditions on it & Eventual in & Depends on & Used at\\\hline
1 & $\eps=2^{-5}$, $\sigma=100$, $\Cp=103$, $\Aexp=105$ & explicit & -- & everywhere\\\hline
2 & $\Nzero$: $\Nzero\ge2^{40}$ (redundant: size condition~(i), $L\ge2^{10}$, already forces every $N\ge\Nzero$ to satisfy $N\ge2^{1024}$, i.e.\ $\lceil\Nzero\rceil\ge2^{1024}$), and the size conditions of Lemma~\ref{s4:lemTPV}, Lemma~\ref{s4:lemPV} and Theorem~\ref{s4:thmVXp} hold for every $N\ge\Nzero$ & $N$ (so every requirement is upward closed in $\Nzero$) & absolute constants only & $N=|Z|$ for the vertex set $Z$ of a vortex run; $N\ge\Nzero$ by \Gc{3}\\\hline
3 & $\Dstar$: \Gc{1}, i.e.\ $\Dstar>2$ and (a)--(f), including \Gstar{} and the column-3 inequalities of Table~\ref{s3:tabCOLJV} & $\mu$, on the ray $\mu\ge\log_2\log_2\Dstar$; hence $\Dstar$ & absolute constants only & $\mu=\log_2\lambda_r$ ($r\le R$), $\log_2\log_2d$ ($d\ge\Dstar$; e.g.\ $\log_2x$ for $x\ge\log_2\Dstar$), $\log_2\Dstar$\\\cline{2-5}
 & \Gc{2}(a): $\Dstar\ge2^{117}$ & $\Dstar$ & absolute constants only & $d_l\ge\Dstar$ (Lemma~\ref{s2:lemCap})\\\cline{2-5}
 & \Gc{2}(b),(c) & consequences of \Gc{1} & absolute constants only; not a condition & every round, every ancestor\\\cline{2-5}
 & \Gc{3}: $(\log_2\Dstar)^{\Cp}\ge2\Nzero$ & $\Dstar$ & $\Nzero$ (step~2) & sizes of parts, $P_R/2\ge\Nzero$\\\cline{2-5}
 & \Gc{4}: $\epsOne(\Dstar)\le6$ and $\epsTwo(\Dstar)\le1/4$ & $\Dstar$ & absolute constants only & $\Dstar$ (cost line, Section~\ref{s7:sec}; Definition~\ref{s1:defConstants}(v); Remark~\ref{s5:remConstants}(b))\\\hline
4 & $\epsOne(\Dstar)$, $\thetaQ=\epsTwo(\Dstar)$, $\Czero$ & defined & $\Dstar$ & Theorem~\ref{s7:thmHI}\\\hline
5 & $\cEG=\max\{\Czero/(1-2\thetaQ),\Nzero/2\}$ & defined & $\Czero$, $\thetaQ$, $\Nzero$ & Theorem~\ref{s7:thmHI}, as its constant $c$\\\hline
\end{tabular}
\end{center}
\caption{The order in which the constants are chosen, and the conditions on
them (Condition~\ref{s1:condGamma}, Remark~(2)). ``Eventual in $x$'' means:
holds for all sufficiently large $x$ once the constants of the earlier steps are
fixed. A condition on the constant of a step depends only on earlier steps. For
steps~2 and~3, the last column gives the values at which a condition is used;
each of them is at least the chosen threshold.}
\label{s1:tabOrder}
\end{table}

\subsection{The symbol table}\label{s1:ssecSymbols}

Every symbol that is used outside a single proof has exactly one meaning, given in
Table~\ref{s1:tabSymbols}. A symbol listed there is not reused locally. Symbols
local to one proof or one run are introduced there and have no meaning elsewhere:
the sans-serif symbols of a vortex run (Convention~\ref{s4:convVortex}); the
starred parameters of the proof of Theorem~\ref{s3:thmT16s}; the node notation of
Lemma~14$^\tau$ and Lemma SEP (Section~\ref{s2:sec}); the cluster, layer,
demand, padding and potential notation of Lemmas~\ref{s6:lemMED}
and~\ref{s6:lemHCCP}; and the letters $S_w$, $\mu$, $\mu^*$, $T$, $m_w$ of
Lemmas~\ref{s7:lemCC} and~\ref{s7:lemUltra}. The letter $\eta$ is deliberately
local as well and is not in the table: it denotes the failure terms
$\eta_{\mathrm{TPV}}$, $\eta_{\mathrm{PV}}$ and $\eta_{\mathrm{VX}}$ of the three
vortex statements of Section~\ref{s4:sec}; the bijections $\eta_h$ of
Definition~\ref{s7:defSchedule} and Construction~\ref{s7:consRound}; the
functions $\eta(x)$ of Lemma~\ref{s2:lemLacunary}(iv) and of the proof of
Lemma~\ref{s5:lemParent}(ii) (the latter also in the overview); and numbers
inside single proofs (for instance those of Lemmas~\ref{s1:lemBBD}
and~\ref{s6:lemJSLC}). The column ``macro'' gives the
\LaTeX{} macro of the preamble (``--'' if the symbol has none), and the column
``defined in'' points to the defining statement (forward pointers are allowed in
this table).

\begin{center}
\refstepcounter{table}\label{s1:tabSymbols}
\textbf{Table~\thetable.} The symbol table.\\[4pt]
{\small
\setlength{\tabcolsep}{4pt}
\begin{tabular}{p{3.7cm}p{3.0cm}p{5.6cm}p{2.5cm}}
\textbf{Symbol} & \textbf{Macro} & \textbf{Meaning} & \textbf{Defined in}\\
\hline
\multicolumn{4}{l}{\emph{General}}\\
$\f(F)$, $\f(G)$, $\f(n)$ & \texttt{\string\f} & least number of objects (cycles, single edges) partitioning $F$, $E(G)$; maximum over $n$-vertex $G$ & \ref{s1:defObject}\\
$n$ & -- & number of vertices of the input graph $G$ & \ref{s1:convGraphs}\\
$\log$, $\logs$ & \texttt{\string\log}, \texttt{\string\logs} & base $2$; iterated logarithm & \ref{s1:convGraphs}\\
$\Nbr_G(U)$, $\Ball{i}{G}{U,V}$ & \texttt{\string\Nbr}, \texttt{\string\Ball} & outside neighbourhood; ball through $V$ & \ref{s1:citNotation}\\
$\Prob$, $\Ex$, $\ind{\cdot}$, $\Bin$ & \texttt{\string\Prob}, \texttt{\string\Ex}, \texttt{\string\ind}, \texttt{\string\Bin} & probability, expectation, indicator, binomial & --\\
\hline
\multicolumn{4}{l}{\emph{Constants}}\\
$\eps=2^{-5}$; $\sigma=100$; $\Cp=103$; $\Aexp=105$ & \texttt{\string\eps}, --, \texttt{\string\Cp}, \texttt{\string\Aexp} & fixed parameters & \ref{s1:defConstants}\\
$\Nzero$, $\Dstar$, $\Czero$, $\cEG$ & \texttt{\string\Nzero}, \texttt{\string\Dstar}, \texttt{\string\Czero}, \texttt{\string\cEG} & size threshold; galactic degree threshold; main constants & \ref{s1:defConstants}\\
$\epsOne$, $\epsTwo=\thetaQ$ & \texttt{\string\epsOne}, \texttt{\string\epsTwo}, \texttt{\string\thetaQ} & cost and quotient-size functions of $\Dstar$ & \ref{s7:propCost}, \ref{s7:lemUHsplit}\\
$\epsA$, $\epsU$, $\epsK=\epsChain$, $\epsCONC$, $\varepsilon_M$, $\epsX$ & \texttt{\string\epsA}, \texttt{\string\epsU}, \texttt{\string\epsK}, \texttt{\string\epsChain}, \texttt{\string\epsCONC}, --, \texttt{\string\epsX} & named small functions of $\Dstar$ & \ref{s2:lemTower}, \ref{s5:lemExpect}, \ref{s5:lemParent}/\ref{s5:lemKRED}, \ref{s6:thmCONCL}, \ref{s6:thmMIXC}, \ref{s7:lemEXprime}\\
$\epsD$ & \texttt{\string\epsD} & generic, inside proofs only & \ref{s1:defConstants}(vii)\\
\Gc{1}--\Gc{4}, \Gstar & \texttt{\string\Gc}, \texttt{\string\Gstar} & galactic conditions & \ref{s1:condGamma}\\
$\cOV=\log_2(4/3)$ & \texttt{\string\cOV} & overlap constant & \ref{s2:lemOVgeneric}\\
$\psi(x)=(6\log_2x+12)/x$ & -- & tower sum weight & \ref{s2:lemTower}\\
\hline
\end{tabular}}
\end{center}

\begin{center}
{\small
\setlength{\tabcolsep}{4pt}
\begin{tabular}{p{3.7cm}p{3.0cm}p{5.6cm}p{2.5cm}}
\multicolumn{4}{l}{(Table~\ref{s1:tabSymbols}, continued.) \emph{Hierarchy}}\\
\hline
$\HBtp$ ($\HBs$, $\HBt$ only in comparisons) & \texttt{\string\HBtp} (\texttt{\string\HBs}, \texttt{\string\HBt}) & the hierarchy & \ref{s2:defHBtp}\\
$G_l$, $G'_l$, $d_l$, $\lambda_l=\log d_l$, $R$, $E_0$, $\Cyc_l$ & --, \texttt{\string\Cyc} & round-$l$ graph, the same after removing long cycles; average degree; its logarithm; last round; leftover; long cycles & \ref{s2:defHBtp}\\
$\tHB_l=\log^2d_l$; $M_l\in\mathbb N$ (a ceiling, (R2)); $\Lambda_l=\log M_l$; $s_l=\lceil\Lambda_l^\sigma\rceil$; $P_l=\lceil\lambda_l^{\Cp}\rceil$; $\tau_l=\lceil128s_l\log^2M_l\rceil$ & \texttt{\string\tHB} & round parameters & \ref{s2:defHBtp}\\
$\piece$ & \texttt{\string\piece} & piece of the $s=0$ run of Lemma~14 & \ref{s2:defHBtp}\\
$Z^0$, $X^0_Z$, $S_Z$, $\home(v)$, $D_l$ & \texttt{\string\home} & pre-part, its leaf expander, guests, home, duplicated vertices & \ref{s2:defHBtp}\\
light part $Z=Z^0\setminus S_Z$, $X_Z=X^0_Z-S_Z$; $\Std_l$; GC-part; $\thGC_l(Z^0)=\lceil|Z^0|\lambda_l^{-1/2}\rceil$ & \texttt{\string\Std}, \texttt{\string\thGC} & light/standalone split; guest cap & \ref{s2:defHBtp}\\
$E_l(Z)$ & -- & edges assigned at round $l$ to $Z$ & \ref{s2:defHBtp}\\
witness $(U,F)$, $N$, $F_0$, $\Hout$, $U'$, $N'$, $F_1$, $\Hin$, $N''$, $F''$, $G_1$, $G_2$ & \texttt{\string\Hout}, \texttt{\string\Hin} & $\tau$-rules (local to Section~\ref{s2:sec}) & \ref{s2:defWitness}, \ref{s2:defTauRules}\\
leaf $\Leaf$, $\Dup$, $\Dups_l$ & \texttt{\string\Leaf}, \texttt{\string\Dup}, \texttt{\string\Dups} & leaf of a recursion; vertices in at least two leaves (of one recursion; of round $l$'s two-level recursion) & \ref{s2:lemSEP}, \ref{s2:defHBtp}\\
ancestor $Y$; $V(Y)$; $H_Y$; $\eps_Y$; $s_Y$; $L_Y=\log|V(Y)|$ & -- & ancestor of round $r$ and its expander & \ref{s2:defAncestors}\\
$\anc_l(x)$; $A_Z$; $U_Z$; $F_Z$; $Q_Z$ & \texttt{\string\anc} & ancestors; hubs; ports; fresh ports; classed ports & \ref{s2:defAncestors}\\
$\mu_r(w)$, $\dup_r$, $\mult_r(w)$ & \texttt{\string\dup}, \texttt{\string\mult} & number of round-$r$ pre-parts containing $w$; duplication; number of round-$r$ ancestors containing $w$ & \ref{s2:defAncestors}\\
$j_0(x)$, $j_1(x)$ & -- & first round in which $x$ lies in a pre-part; in a light part & \ref{s2:defAncestors}\\
$\nu_l$ & -- & number of ancestors of rounds $\le l-2$ (at most $2.74n/P_{l-2}$) & \ref{s2:propOV} (bound: \ref{s2:lemTower}(b))\\
$(\alpha)$, $(\beta)$ & -- & ORIGIN edge types & \ref{s2:propOrigin}\\
\hline
\end{tabular}}
\end{center}

\begin{center}
{\small
\setlength{\tabcolsep}{4pt}
\begin{tabular}{p{3.7cm}p{3.0cm}p{5.6cm}p{2.5cm}}
\multicolumn{4}{l}{(Table~\ref{s1:tabSymbols}, continued.) \emph{Colouring, lending, zones, light parts}}\\
\hline
$\Own_Y$, $\Lend_Y$, $\klend(Y)$, $\kown$ & \texttt{\string\Own}, \texttt{\string\Lend}, \texttt{\string\klend}, \texttt{\string\kown} & two-stage colouring & \ref{s3:defCOL}\\
$\IU$, $\IJS$, $\IJV$ & \texttt{\string\IU}, \texttt{\string\IJS}, \texttt{\string\IJV} & lent index families & \ref{s3:defCOL}\\
$\LU_{Y,l,c,\uslot}$, $\LJS_{Y,l,j}$, $\LJV_{Y,l}$ & \texttt{\string\LU}, \texttt{\string\LJS}, \texttt{\string\LJV} & U-, JS- and JV-lent classes & \ref{s3:defCOL}\\
$\KJS_l=M_l^2$; $\tJS_l=2M_l+2$; $\rho_l=M_l^{-4}$; $\lab_{Y,l}(y)$; $T_j(Y,l)$ & \texttt{\string\KJS}, \texttt{\string\tJS}, \texttt{\string\lab} & JS layers, multiplicity, label probability, labels, junction sets & \ref{s3:defCOL}\\
$p_Y=1/(2\klend(Y))$ & -- & probability that an edge of $H_Y$ lies in $\LJV_{Y,l}$ & \ref{s3:defCOL}\\
$\Tslot_Y=\lceil L_Y^2\rceil$ & \texttt{\string\Tslot} & number of U-slots & \ref{s3:defCOL}\\
$t_Y=\lceil\lambda_r^{1.6}\rceil$ & -- & U-multiplicity parameter of a light ancestor $Y$ of round $r$ & \ref{s3:defCOL}\\
$\Aw(w)$ & \texttt{\string\Aw} & candidate sets of Lemma HB & \ref{s3:lemHB}\\
$\mathrm{choice}(v)$, $\zp_Y=L_Y^{-2}$, sublabel $(l,c,\uslot)$, $\Zone_{Y,l,c,\uslot}$, $\rho_Y$ & \texttt{\string\zp}, \texttt{\string\uslot}, \texttt{\string\Zone} & Theorem-U zones & \ref{s5:defZones}, \ref{s5:lemZones}\\
(E1), demoted, parent-bad, good parent, $\Bad$, $\Rt(Z)$, $\Pl(Z)$, $\parU(v)$, $\mathrm{dem}$, $\mathrm{lp}$ & \texttt{\string\Bad}, \texttt{\string\Rt}, \texttt{\string\Pl}, \texttt{\string\parU} & stage-2 statuses of light parts & \ref{s5:defStages}\\
$\extph$ & \texttt{\string\extph} & external phase map: phase of a vertex's round-$l$ zone, or $*$ & \ref{s4:lemPV}, \ref{s5:lemChild}\\
U-bundle $\Bdl_{l,c}$; U-quotient multigraph $\UQ_{l,c,i}$ & \texttt{\string\Bdl}, \texttt{\string\UQ} & Theorem-U chaining & \ref{s5:lemParent}\\
$\LentU(Y)$, $\Lentext(Y)$, $\Kstd$ & \texttt{\string\LentU}, \texttt{\string\Lentext}, \texttt{\string\Kstd} & used U-lent edges; externally used lent edges; standalone edges & \ref{s5:lemKRED}\\
\hline
\end{tabular}}
\end{center}

\begin{center}
{\small
\setlength{\tabcolsep}{4pt}
\begin{tabular}{p{3.7cm}p{3.0cm}p{5.6cm}p{2.5cm}}
\multicolumn{4}{l}{(Table~\ref{s1:tabSymbols}, continued.) \emph{Standalone parts and chaining}}\\
\hline
designation $\delta$, class $Y(u)$ & -- & class of a classed port & \ref{s6:defDesign}\\
$d_{Y,l}(h)$, $m_{Y,l}$, $d^*_{Y,l}$, $\alpha_{Y,l}$ & -- & class degree, its maximum, duplicated class ports, guest core degree & \ref{s6:defDesign}\\
$\Bead_{Y,l}$, $\gamma_l=\lfloor P_{l-2}/M_l\rfloor$, giant & \texttt{\string\Bead} & deterministic bead graph; giant threshold ($(Y,l)$ is giant iff a component has more than $2\gamma_l$ edges) & \ref{s6:defDesign}\\
$\cagg_{h,l}$, $c_x(Z)$, $\cpp(u)$ & \texttt{\string\cagg}, \texttt{\string\cpp} & class counts at hubs, fresh centres, ports & \ref{s6:defDesign}\\
cluster $\clu=(A_\clu,U_\clu,B_\clu)$, $b(\clu)$, $\exc(u)$, $\Phi(\clu)$, $\MED(<)$, $\pot$ & \texttt{\string\clu}, \texttt{\string\exc}, \texttt{\string\MED}, \texttt{\string\pot} & clusters & \ref{s6:defCluster}, \ref{s6:lemMED}\\
$\LayP_j$, $\pad(u)$, $\dem^{\pm}(u)$, $\Phi_j$, $\mathcal P_j$, $F_j$ & \texttt{\string\LayP}, \texttt{\string\pad}, \texttt{\string\dem} & HCC-P data & \ref{s6:lemHCCP}\\
lend-bad; $\Lost_Z$, $\Lost_l$; $\Ret_Z=A_Z\cup F_Z\cup\Lost_Z$; $\Qs_Z=Q_Z\setminus\Lost_Z$; $O_Z$; $\XU$ & \texttt{\string\Lost}, \texttt{\string\Ret}, \texttt{\string\Qs}, \texttt{\string\XU} & stage-2 lending data and the functional $X_U$ & \ref{s6:defLending}\\
$\sco_{Y,l}$, $g_{Y,l}$, $S_0(Y,l)$ & \texttt{\string\sco} & cherry split cost, number of cherry classes, non-giant system (local to the proof) & \ref{s6:lemJSLC} (proof)\\
$J_l=\Jlost_l+\Jhub_l+\Jfr_l+\Jpar_l$; \Jp{1}--\Jp{3} & \texttt{\string\Jlost}, \texttt{\string\Jhub}, \texttt{\string\Jfr}, \texttt{\string\Jpar}, \texttt{\string\Jp} & the J-set & \ref{s6:lemJplus}\\
J-consumer, $\Obj_l$, $\LentJV_l$, \JC{1}--\JC{3} & \texttt{\string\Obj}, \texttt{\string\LentJV}, \texttt{\string\JC} & consumer interface & \ref{s6:defJconsumer}\\
$\Lent(Y)$, $\LentU(Y)$, $\LentJS(Y)$, $\LentJV(Y)$, $\Erem(Y)$, $H_0(Y)$ & \texttt{\string\Lent}, \texttt{\string\LentU}, \texttt{\string\LentJS}, \texttt{\string\LentJV}, \texttt{\string\Erem} & used lent edges and remainders (defined recursively along the processing order) & \ref{s6:consOrder}, \ref{s6:lemLent}\\
$\EV(Z)$, $\EQ(Z)$, $B_Z=\EQ(Z)$ & \texttt{\string\EV}, \texttt{\string\EQ} & TPV output & \ref{s4:lemTPV}, \ref{s6:consOrder}\\
\hline
\end{tabular}}
\end{center}

\begin{center}
{\small
\setlength{\tabcolsep}{4pt}
\begin{tabular}{p{3.7cm}p{3.0cm}p{5.6cm}p{2.5cm}}
\multicolumn{4}{l}{(Table~\ref{s1:tabSymbols}, continued.) \emph{Junction-vertex quotient}}\\
\hline
$\plab(v)$, $q_l=M_l^{-2}$, $\pi_{l,r}=2^{-(l-1-r)}$, $\Pool_{l,r}$, $\Pool_l$, $r(w)$ & \texttt{\string\plab}, \texttt{\string\Pool} & round-split pool & \ref{s7:defPool}\\
$\Cand_l(u)$, $\Hcd_l=\lambda_{l-2}^{95}/(8M_l^2)$, JV-bad & \texttt{\string\Cand}, \texttt{\string\Hcd} & candidates & \ref{s7:defCand}\\
$\xi_l$, $\Past_l$ & \texttt{\string\Past} & round-$l$ randomness; the fixed past & \ref{s7:defSchedule}\\
items, live items, $\Live$, $\clive_h$, $\thult_l=\lfloor M_l\Hcd_l/7\rfloor$, ultra, $k_h$, $\KHUB_l=4M_l$ & \texttt{\string\Live}, \texttt{\string\clive}, \texttt{\string\thult}, \texttt{\string\KHUB} & round-step data & \ref{s7:consRound}\\
PAR object, cherry, unpaired fresh leg & -- & PAR data & \ref{s7:consRound}\\
$\Used(u)$, $\Used_\kappa(h)$ & \texttt{\string\Used} & junctions already used & \ref{s7:consRound}\\
$\BP_\kappa$, $\BH_\kappa$, $B^{(\iota)}$, $Q_l$, $\Dec_l$ & \texttt{\string\BP}, \texttt{\string\BH}, \texttt{\string\Dec} & layers, sub-layers, quotient, its decomposition & \ref{s7:consRound}\\
$\copies_l$, $\payrd_l$ & \texttt{\string\copies}, \texttt{\string\payrd} & round-$l$ quotient size and random payments & \ref{s7:consRound}\\
$m_\kappa(h,w)$, $m_\kappa(w,w')$ & -- & multiplicities in HUB and PAR layers & \ref{s7:lemMULT}, \ref{s7:lemCC}\\
$\tCC_l=\lceil2\log_2M_l\rceil$ & \texttt{\string\tCC} & parameter of Lemma CC & \ref{s7:lemCC}\\
$X_{V,l}$, $\XV$ & \texttt{\string\XV} & junction-copy weight & \ref{s7:lemVstar}\\
$\omega_l(v)$ & -- & pooled-vertex weight ($\cagg_{v,l}$ for $v\in D_l$, $M_l$ otherwise) & \ref{s7:lemPay}\\
$\Xpool$, $\Xp=\XU+\Xpool+\XV$ & \texttt{\string\Xpool}, \texttt{\string\Xp} & stage-1 functional & \ref{s7:defXprime}\\
$\dett_l$ & \texttt{\string\dett} & deterministic part of the quotient size & \ref{s7:lemUHsplit}\\
\hline
\end{tabular}}
\end{center}

\subsection{The dependency table}\label{s1:ssecDAG}

Table~\ref{s1:tabDAG} has one row for every numbered statement of
Sections~\ref{s1:sec}--\ref{s7:sec}, in the order in which the statements appear
(sectioning units such as the overview are not statements and have no row). The
columns are: the label; a short name and the number of the statement; its
section; and the statements it depends on (this column copies the
``Dependencies'' line printed at the end of the statement).

\emph{What an entry means.} The entries of the row of a statement $A$ are
its logical inputs: the statements that the statement or the proof of $A$ uses
through a reference, by name (for instance ``by Hall's theorem'' or ``by
Markov's inequality''), or through a named item such as \Gc{3}, (K1), \Jp{2},
COL(g) or (R5). A fact that is proved in the running text between two
statements (for instance \eqref{s5:eqLY}, \eqref{s5:eqZp},
\eqref{s6:eqTowerHalf}, the observations of Section~\ref{s4:ssecConv}, or the
bounds on the numbers of PAR objects and hub items in the opening paragraph of
the subsection on quotient size and payments in Section~\ref{s7:sec}) has no
row; a statement that
uses such a fact lists the statements from which the fact is proved.
Convention~\ref{s1:convGraphs} (graphs are finite and simple, unless they are
called multigraphs) holds throughout; it is listed only in the rows of the
statements that cite it, directly or through an observation of
Section~\ref{s4:ssecConv}. Definition~\ref{s1:defConstants} fixes the constants
used throughout. It is listed in the row of every later statement that cites it
or invokes a requirement that it places on a constant (for instance ``by the
choice of $\Nzero$''); every other statement that uses one of these constants
reaches it through the entries of its row. The row of a remark lists the earlier
statements that the remark is about. A few rows also list the result of
\cite{BM} whose proof the statement modifies and writes out in full.
Theorem~\ref{s3:thmT16s} does not list \cite[Theorem~16]{BM}: its proof follows
that of \cite[Theorem~16]{BM}, but it is written out in full in
Section~\ref{s3:sec} and uses neither that theorem nor a step of its proof. A few
rows list a statement that is used only where the statement of the row is
applied, or only through another entry of the same row: for instance \Gc{3} in
the rows of Lemma~\ref{s4:lemTPV}, Lemma~\ref{s4:lemPV} and
Theorem~\ref{s4:thmVXp}, and \Gc{2} in the row of Lemma~\ref{s2:lemTower}. These
entries are harmless over-declarations. The following are not entries: pointers
that only say where a result is proved, used or modified, or that it is not
used (except in the row of Theorem~\ref{s1:thmMain}, whose one entry is its
proof, Corollary~\ref{s7:thmMainProof}; see below); comparisons with \cite{BM}; and the mere use of notation introduced
elsewhere. The items whose name is marked ``not used'' (or ``not used
directly'') are entries of no row.

\emph{Forward references.} Apart from pointers (see above) and the reference of
Theorem~\ref{s1:thmMain} to its proof, Corollary~\ref{s7:thmMainProof}, four
kinds of forward reference occur in the text. They are not entries, and none of
them is circular.
(1) Constants. Theorem~\ref{s1:thmMain}, Definition~\ref{s1:defConstants},
Remark~\ref{s1:remOrder} and Condition~\ref{s1:condGamma} use explicit
functions of $\Dstar$ ($\epsOne$, $\epsTwo$ and the terms they collect),
explicit size thresholds and, in \Gc{1}(f), the explicit inequalities of
column~3 of Table~\ref{s3:tabCOLJV}; these are written out in later
statements, proofs and tables.
These are explicit formulas; no lemma of the paper is needed to state them.
What is used about them is that finitely many explicit inequalities hold
eventually: in $N$ for $\Nzero$ (Definition~\ref{s1:defConstants}(ii) and the
size conditions of Section~\ref{s4:sec}), and in $\Dstar$ for $\Dstar$
(Lemma~\ref{s7:lemGammaSat}). The existence of $\Nzero$ and of such a $\Dstar$
is used only in Corollary~\ref{s7:thmMainProof}; every other statement uses only
the properties that Definition~\ref{s1:defConstants} requires of the fixed
constants. Remark~\ref{s1:remOrder} also says where
constants enter later (for instance, that $\cEG$ enters only through the
induction of Theorem~\ref{s7:thmHI}); this is a description, verified in those
statements.
(2) Announced consequences. Condition~\ref{s1:condGamma} records consequences
of its items that are proved later, where they are used. For instance,
\Gc{2}(b) is proved in Lemma~\ref{s2:lemTower}(b), and the bounds of \Gc{2}(c)
in Lemma~\ref{s3:lemCOL} (with rows~6 and~7 of Table~\ref{s3:tabCOLJV}) and in
Lemma~\ref{s5:lemE1}(b),(c). None of these lemmas uses \Gc{2}(b) or (c), and
\Gc{2}(a) enters them only through Lemma~\ref{s2:lemCap}; the only statement
that uses \Gc{2}(b),(c) as such, Lemma~\ref{s7:lemGammaSat}, lists all three
lemmas. Definition~\ref{s7:defCand} notes that every JV-good port $u$ has
$|\Cand_l(u)|\ge\Hcd_l$. This is Lemma~\ref{s7:lemCand}(iv), which is proved
from the definition of JV-good and from Lemma~\ref{s7:lemCand}(ii), not from the
note. Definition~\ref{s3:defCOL} notes, in its paragraph ``Consumers'', that
for $r\ge R-1$ the set $\Lent(Y)$ is empty, so that $\Erem(Y)=E_r(Y)$ or
$H_0(Y)=E_r(Y)$ (Construction~\ref{s6:consOrder}); and
Definition~\ref{s5:defStages} notes that for $r\ge R-1$
Lemma~\ref{s5:lemE1}(c) reduces to~(b). None of these conditions and
definitions uses these consequences.
(3) Constructions. That Construction~\ref{s6:consOrder} applies every device
within its hypotheses is shown in Lemma~\ref{s6:lemLent} and
Theorem~\ref{s6:thmMIXC}(a); that Construction~\ref{s7:consRound} is well defined
is Lemma~\ref{s7:lemWellDef}(i), (ii), (iv) and~(v). These results are proved
about the constructions. Besides earlier statements, each part uses only the
definitions of the steps it concerns and of the steps before them (for
Construction~\ref{s6:consOrder}, the steps at the rounds already processed), and
not the well-definedness of the step that it proves.
(4) Definition~\ref{s7:defSchedule} describes the choice of the stage-1 outcome
through the functional $\Xp$ of Definition~\ref{s7:defXprime}, and the choice of
each $\xi_l$ through step~(f) of Construction~\ref{s7:consRound}. These choices
are made in Lemma~\ref{s7:lemOneOutcome}, whose row lists
Construction~\ref{s7:consRound} and Lemma~\ref{s7:lemEXprime} (the latter rests
on Definition~\ref{s7:defXprime}). Moreover, $\Xp$ is a function of the run,
$\delta$ and stage~1 only (Definition~\ref{s7:defXprime}), so the selection~(1e)
of Definition~\ref{s7:defSchedule} depends neither on the later stages of the
schedule nor on the round step.

\begin{center}
\refstepcounter{table}\label{s1:tabDAG}
\textbf{Table~\thetable.} Dependency table of all numbered statements.\\[4pt]
{\footnotesize
\setlength{\tabcolsep}{3pt}
\begin{tabular}{p{3.2cm}p{3.6cm}p{0.5cm}p{7.6cm}}
\textbf{Label} & \textbf{Name} & \textbf{S.} & \textbf{Depends on}\\
\hline
\texttt{s1:thmMain} & Main Theorem (\ref{s1:thmMain}) & 1 & \ref{s7:thmMainProof}\\
\texttt{s1:remStatus} & organization of the proof (\ref{s1:remStatus}) & 1 & --\\
\texttt{s1:convGraphs} & graph conventions (\ref{s1:convGraphs}) & 1 & --\\
\texttt{s1:defObject} & objects; $\f$ (\ref{s1:defObject}) & 1 & --\\
\texttt{s1:factAdd} & properties of $\f$ (\ref{s1:factAdd}) & 1 & \ref{s1:defObject}\\
\texttt{s1:factEG0} & long-cycle bound (\ref{s1:factEG0}) & 1 & \ref{s1:convGraphs}, \ref{s1:defObject}\\
\texttt{s1:citNotation} & B--M notation (\ref{s1:citNotation}) & 1 & --\\
\texttt{s1:citChernoff} & B--M Thm 4 (Chernoff) (\ref{s1:citChernoff}) & 1 & --\\
\texttt{s1:citDef7} & B--M Def 7 (\ref{s1:citDef7}) & 1 & --\\
\texttt{s1:citProp8} & B--M Prop 8 (\ref{s1:citProp8}) & 1 & --\\
\texttt{s1:citLem9} & B--M Lemma 9 (\ref{s1:citLem9}) & 1 & --\\
\texttt{s1:citDef11} & B--M Def 11 (\ref{s1:citDef11}) & 1 & --\\
\texttt{s1:citProp12} & B--M Prop 12 (\ref{s1:citProp12}) & 1 & --\\
\texttt{s1:citProp13} & B--M Prop 13 (\ref{s1:citProp13}) & 1 & --\\
\hline
\end{tabular}}
\end{center}

\begin{center}
{\footnotesize
\setlength{\tabcolsep}{3pt}
\begin{tabular}{p{3.2cm}p{3.6cm}p{0.5cm}p{7.6cm}}
\multicolumn{4}{l}{(Table~\ref{s1:tabDAG}, continued.)}\\
\textbf{Label} & \textbf{Name} & \textbf{S.} & \textbf{Depends on}\\
\hline
\texttt{s1:citLem14} & B--M Lemma 14 and its proof (\ref{s1:citLem14}) & 1 & --\\
\texttt{s1:citLem15} & B--M Lemma 15 (not used) (\ref{s1:citLem15}) & 1 & --\\
\texttt{s1:citThm16} & B--M Thm 16 (not used; see Thm 16$^*$) (\ref{s1:citThm16}) & 1 & --\\
\hline
\end{tabular}}
\end{center}

\begin{center}
{\footnotesize
\setlength{\tabcolsep}{3pt}
\begin{tabular}{p{3.2cm}p{3.6cm}p{0.5cm}p{7.6cm}}
\multicolumn{4}{l}{(Table~\ref{s1:tabDAG}, continued.)}\\
\textbf{Label} & \textbf{Name} & \textbf{S.} & \textbf{Depends on}\\
\hline
\texttt{s1:citLem17} & B--M Lemma 17 (\ref{s1:citLem17}) & 1 & --\\
\texttt{s1:citProp18} & B--M Prop 18 (\ref{s1:citProp18}) & 1 & --\\
\texttt{s1:citLem19} & B--M Lemma 19 (\ref{s1:citLem19}) & 1 & --\\
\texttt{s1:citThm21} & B--M Thm 21 (not used directly) (\ref{s1:citThm21}) & 1 & --\\
\texttt{s1:citCor22} & B--M Cor 22 (\ref{s1:citCor22}) & 1 & --\\
\texttt{s1:citLem25} & B--M Lemma 25 (\ref{s1:citLem25}) & 1 & --\\
\texttt{s1:remBMused} & where B--M is used (\ref{s1:remBMused}) & 1 & \ref{s1:citNotation}, \ref{s1:citChernoff}, \ref{s1:citDef7}, \ref{s1:citProp8}, \ref{s1:citLem9}, \ref{s1:citDef11}, \ref{s1:citProp12}, \ref{s1:citProp13}, \ref{s1:citLem14}, \ref{s1:citLem17}, \ref{s1:citProp18}, \ref{s1:citLem19}, \ref{s1:citCor22}, \ref{s1:citLem25}, \ref{s1:factEG0}\\
\hline
\end{tabular}}
\end{center}

\begin{center}
{\footnotesize
\setlength{\tabcolsep}{3pt}
\begin{tabular}{p{3.2cm}p{3.6cm}p{0.5cm}p{7.6cm}}
\multicolumn{4}{l}{(Table~\ref{s1:tabDAG}, continued.)}\\
\textbf{Label} & \textbf{Name} & \textbf{S.} & \textbf{Depends on}\\
\hline
\texttt{s1:citChernoffGen} & Chernoff (Poisson-binomial) (\ref{s1:citChernoffGen}) & 1 & --\\
\texttt{s1:citMarkov} & Markov (\ref{s1:citMarkov}) & 1 & --\\
\texttt{s1:lemBBD} & Lemma BBD (Bernstein, bounded differences) (\ref{s1:lemBBD}) & 1 & --\\
\texttt{s1:citHall} & Hall's theorem (\ref{s1:citHall}) & 1 & --\\
\texttt{s1:citHaxell} & Haxell 1995 (\ref{s1:citHaxell}) & 1 & --\\
\hline
\end{tabular}}
\end{center}

\begin{center}
{\footnotesize
\setlength{\tabcolsep}{3pt}
\begin{tabular}{p{3.2cm}p{3.6cm}p{0.5cm}p{7.6cm}}
\multicolumn{4}{l}{(Table~\ref{s1:tabDAG}, continued.)}\\
\textbf{Label} & \textbf{Name} & \textbf{S.} & \textbf{Depends on}\\
\hline
\texttt{s1:citEuler} & $T$-joins, Euler (\ref{s1:citEuler}) & 1 & --\\
\texttt{s1:defConstants} & the constants (\ref{s1:defConstants}) & 1 & --\\
\texttt{s1:remOrder} & order of constants (\ref{s1:remOrder}) & 1 & \ref{s1:defConstants}\\
\texttt{s1:condGamma} & galactic conditions (\ref{s1:condGamma}) & 1 & \ref{s1:defConstants}, \ref{s1:remOrder}\\
\hline
\end{tabular}}
\end{center}

\begin{center}
{\footnotesize
\setlength{\tabcolsep}{3pt}
\begin{tabular}{p{3.2cm}p{3.6cm}p{0.5cm}p{7.6cm}}
\multicolumn{4}{l}{(Table~\ref{s1:tabDAG}, continued.)}\\
\textbf{Label} & \textbf{Name} & \textbf{S.} & \textbf{Depends on}\\
\hline
\texttt{s2:defWitness} & witness (\ref{s2:defWitness}) & 2 & \ref{s1:citDef11}, \ref{s1:defConstants}\\
\texttt{s2:defTauRules} & $\tau$-rules and split (\ref{s2:defTauRules}) & 2 & \ref{s2:defWitness}\\
\texttt{s2:lemSEP} & Lemma SEP (\ref{s2:lemSEP}) & 2 & \ref{s2:defTauRules}\\
\texttt{s2:lemThinCut} & thin cut (\ref{s2:lemThinCut}) & 2 & \ref{s2:lemSEP}, \ref{s2:defTauRules}\\
\texttt{s2:lemOVgeneric} & Lemma OV (generic overlap) (\ref{s2:lemOVgeneric}) & 2 & \ref{s2:lemSEP}, \ref{s2:defWitness}, \ref{s1:citLem14}\\
\texttt{s2:lem14tau} & Lemma 14$^\tau$ (\ref{s2:lem14tau}) & 2 & \ref{s1:citLem14}, \ref{s1:citDef11}, \ref{s2:defTauRules}, \ref{s2:lemThinCut}, \ref{s2:lemOVgeneric}, \ref{s2:lemSEP}, \ref{s2:defWitness}\\
\texttt{s2:defHBtp} & the hierarchy $\HBtp$ (\ref{s2:defHBtp}) & 2 & \ref{s1:citLem14}, \ref{s2:lem14tau}, \ref{s1:defConstants}, \ref{s2:lemOVgeneric}, \ref{s2:lemSEP}\\
\hline
\end{tabular}}
\end{center}

\begin{center}
{\footnotesize
\setlength{\tabcolsep}{3pt}
\begin{tabular}{p{3.2cm}p{3.6cm}p{0.5cm}p{7.6cm}}
\multicolumn{4}{l}{(Table~\ref{s1:tabDAG}, continued.)}\\
\textbf{Label} & \textbf{Name} & \textbf{S.} & \textbf{Depends on}\\
\hline
\texttt{s2:defAncestors} & ancestors, typing (\ref{s2:defAncestors}) & 2 & \ref{s2:defHBtp}\\
\texttt{s2:lemCap} & Lemma-25 size cap (\ref{s2:lemCap}) & 2 & \ref{s1:citLem25}, \ref{s1:condG2}, \ref{s2:defHBtp}, \ref{s2:lem14tau}\\
\texttt{s2:lemHS} & Lemma HS (\ref{s2:lemHS}) & 2 & \ref{s1:citDef11}\\
\texttt{s2:propStructure} & structure of a run (\ref{s2:propStructure}) & 2 & \ref{s2:defHBtp}, \ref{s2:defAncestors}, \ref{s2:lem14tau}, \ref{s2:lemCap}, \ref{s2:lemHS}, \ref{s1:citDef11}, \ref{s1:citLem14}, \ref{s2:lemSEP}, \ref{s1:condG1}\\
\texttt{s2:lemEL} & edge laminarity EL (\ref{s2:lemEL}) & 2 & \ref{s2:defHBtp}, \ref{s2:defAncestors}\\
\texttt{s2:propOV} & overlap constants (\ref{s2:propOV}) & 2 & \ref{s2:lemOVgeneric}, \ref{s2:lem14tau}, \ref{s2:propStructure}, \ref{s2:defAncestors}, \ref{s2:defHBtp}, \ref{s2:lemSEP}\\
\texttt{s2:propDegRec} & degree recursion (\ref{s2:propDegRec}) & 2 & \ref{s2:lem14tau}, \ref{s2:lemCap}, \ref{s2:propOV}, \ref{s2:defHBtp}, \ref{s2:lemSEP}, \ref{s1:condG1}, \ref{s2:lemOVgeneric}\\
\hline
\end{tabular}}
\end{center}

\begin{center}
{\footnotesize
\setlength{\tabcolsep}{3pt}
\begin{tabular}{p{3.2cm}p{3.6cm}p{0.5cm}p{7.6cm}}
\multicolumn{4}{l}{(Table~\ref{s1:tabDAG}, continued.)}\\
\textbf{Label} & \textbf{Name} & \textbf{S.} & \textbf{Depends on}\\
\hline
\texttt{s2:propExists} & existence, termination (\ref{s2:propExists}) & 2 & \ref{s2:lem14tau}, \ref{s2:propDegRec}, \ref{s2:defHBtp}, \ref{s2:lemCap}, \ref{s2:lemOVgeneric}, \ref{s2:defWitness}, \ref{s2:propOV}\\
\texttt{s2:lemLacunary} & tower-lacunary sums (\ref{s2:lemLacunary}) & 2 & \ref{s2:propDegRec}, \ref{s1:condG1}\\
\texttt{s2:lemTower} & tower facts (\ref{s2:lemTower}) & 2 & \ref{s2:propDegRec}, \ref{s2:lemLacunary}, \ref{s2:defHBtp}, \ref{s1:condG1}, \ref{s1:condG2}, \ref{s2:propStructure}, \ref{s2:lemCap}, \ref{s2:propOV}, \ref{s2:defAncestors}\\
\texttt{s2:lemGC} & rule GC (\ref{s2:lemGC}) & 2 & \ref{s2:defHBtp}, \ref{s2:lemTower}, \ref{s2:lemEL}, \ref{s2:defAncestors}\\
\texttt{s2:propOrigin} & ORIGIN$^\tau$ (\ref{s2:propOrigin}) & 2 & \ref{s2:lemThinCut}, \ref{s2:lemSEP}, \ref{s2:lemCap}, \ref{s2:propOV}, \ref{s2:lemEL}, \ref{s2:defHBtp}, \ref{s2:lemGC}, \ref{s2:lemOVgeneric}, \ref{s2:propStructure}\\
\texttt{s2:propParentless} & fresh, parentless mass (\ref{s2:propParentless}) & 2 & \ref{s2:propStructure}, \ref{s2:lemTower}, \ref{s2:lemCap}, \ref{s2:defAncestors}, \ref{s2:defHBtp}\\
\texttt{s2:remTauConstants} & $\tau{+}$ constants (\ref{s2:remTauConstants}) & 2 & \ref{s2:propOV}, \ref{s2:lem14tau}, \ref{s2:lemTower}, \ref{s2:defHBtp}, \ref{s2:lemOVgeneric}\\
\hline
\end{tabular}}
\end{center}

\begin{center}
{\footnotesize
\setlength{\tabcolsep}{3pt}
\begin{tabular}{p{3.2cm}p{3.6cm}p{0.5cm}p{7.6cm}}
\multicolumn{4}{l}{(Table~\ref{s1:tabDAG}, continued.)}\\
\textbf{Label} & \textbf{Name} & \textbf{S.} & \textbf{Depends on}\\
\hline
\texttt{s3:lemMonotone} & monotonicity (\ref{s3:lemMonotone}) & 3 & \ref{s1:citDef7}, \ref{s1:citDef11}\\
\texttt{s3:lemL15p} & Lemma 15$^+$ (\ref{s3:lemL15p}) & 3 & \ref{s1:citDef11}, \ref{s1:citChernoff}\\
\texttt{s3:remMultiset} & multiset remark (\ref{s3:remMultiset}) & 3 & \ref{s1:citDef7}, \ref{s1:citLem9}\\
\texttt{s3:propP13s} & Prop 13$^*$ (\ref{s3:propP13s}) & 3 & \ref{s1:citProp12}, \ref{s1:citProp13}, \ref{s1:citDef11}\\
\texttt{s3:lemL17s} & Lemma 17$^*$ (\ref{s3:lemL17s}) & 3 & \ref{s3:propP13s}, \ref{s1:citLem17}, \ref{s1:lemBBD}, \ref{s1:citChernoff}, \ref{s1:citChernoffGen}\\
\texttt{s3:lemP18s} & Prop 18$^*$, Lemma 19$^*$ (\ref{s3:lemP18s}) & 3 & \ref{s1:citProp18}, \ref{s1:citLem19}, \ref{s3:lemL17s}, \ref{s1:citDef11}\\
\hline
\end{tabular}}
\end{center}

\begin{center}
{\footnotesize
\setlength{\tabcolsep}{3pt}
\begin{tabular}{p{3.2cm}p{3.6cm}p{0.5cm}p{7.6cm}}
\multicolumn{4}{l}{(Table~\ref{s1:tabDAG}, continued.)}\\
\textbf{Label} & \textbf{Name} & \textbf{S.} & \textbf{Depends on}\\
\hline
\texttt{s3:lemL9rho} & Lemma 9$_\rho$ (\ref{s3:lemL9rho}) & 3 & \ref{s1:citLem9}, \ref{s1:citProp8}, \ref{s1:citHaxell}, \ref{s1:citDef7}, \ref{s3:remMultiset}\\
\texttt{s3:thmT16s} & Theorem 16$^*$ (\ref{s3:thmT16s}) & 3 & \ref{s3:lemL15p}, \ref{s3:propP13s}, \ref{s3:lemL17s}, \ref{s3:lemP18s}, \ref{s3:lemL9rho}, \ref{s3:lemMonotone}, \ref{s1:citDef11}, \ref{s1:citChernoff}\\
\texttt{s3:lemHB} & Lemma HB (\ref{s3:lemHB}) & 3 & \ref{s1:citHall}, \ref{s1:citDef11}\\
\texttt{s3:defCOL} & lending data COL-JV (\ref{s3:defCOL}) & 3 & \ref{s2:defAncestors}, \ref{s2:propStructure}, \ref{s2:lemTower}\\
\texttt{s3:lemCOLJV} & COL-JV table (\ref{s3:lemCOLJV}) & 3 & \ref{s1:condG1}, \ref{s1:condGstar}, \ref{s2:lemTower}, \ref{s3:defCOL}, \ref{s2:defAncestors}, \ref{s2:propStructure}, \ref{s3:lemL15p}, \ref{s3:thmT16s}\\
\hline
\end{tabular}}
\end{center}

\begin{center}
{\footnotesize
\setlength{\tabcolsep}{3pt}
\begin{tabular}{p{3.2cm}p{3.6cm}p{0.5cm}p{7.6cm}}
\multicolumn{4}{l}{(Table~\ref{s1:tabDAG}, continued.)}\\
\textbf{Label} & \textbf{Name} & \textbf{S.} & \textbf{Depends on}\\
\hline
\texttt{s3:lemCOLJVev} & eventual form of COL-JV (\ref{s3:lemCOLJVev}) & 3 & --\\
\texttt{s3:lemCOL} & Lemma COL (\ref{s3:lemCOL}) & 3 & \ref{s3:defCOL}, \ref{s3:lemL15p}, \ref{s3:thmT16s}, \ref{s3:lemMonotone}, \ref{s3:lemCOLJV}, \ref{s2:lemTower}, \ref{s2:propStructure}, \ref{s2:defAncestors}, \ref{s1:condG1}\\
\hline
\end{tabular}}
\end{center}

\begin{center}
{\footnotesize
\setlength{\tabcolsep}{3pt}
\begin{tabular}{p{3.2cm}p{3.6cm}p{0.5cm}p{7.6cm}}
\multicolumn{4}{l}{(Table~\ref{s1:tabDAG}, continued.)}\\
\textbf{Label} & \textbf{Name} & \textbf{S.} & \textbf{Depends on}\\
\hline
\texttt{s4:convVortex} & vortex notation (\ref{s4:convVortex}) & 4 & --\\
\texttt{s4:lemTPV} & Lemma TPV (\ref{s4:lemTPV}) & 4 & \ref{s3:lemL15p}, \ref{s3:thmT16s}, \ref{s3:lemHB}, \ref{s3:lemMonotone}, \ref{s1:citCor22}, \ref{s1:factEG0}, \ref{s1:citChernoff}, \ref{s1:citMarkov}, \ref{s1:condG3}, \ref{s4:convVortex}, \ref{s1:convGraphs}, \ref{s1:defConstants}\\
\hline
\end{tabular}}
\end{center}

\begin{center}
{\footnotesize
\setlength{\tabcolsep}{3pt}
\begin{tabular}{p{3.2cm}p{3.6cm}p{0.5cm}p{7.6cm}}
\multicolumn{4}{l}{(Table~\ref{s1:tabDAG}, continued.)}\\
\textbf{Label} & \textbf{Name} & \textbf{S.} & \textbf{Depends on}\\
\hline
\texttt{s4:lemPV} & Lemma PV (P-vortex) (\ref{s4:lemPV}) & 4 & \ref{s3:thmT16s}, \ref{s3:lemHB}, \ref{s3:lemMonotone}, \ref{s1:citCor22}, \ref{s1:factEG0}, \ref{s1:citChernoff}, \ref{s1:citMarkov}, \ref{s1:condG3}, \ref{s4:convVortex}, \ref{s4:lemTPV}, \ref{s1:convGraphs}, \ref{s1:defConstants}\\
\texttt{s4:thmVXp} & Theorem VX$^+$ (\ref{s4:thmVXp}) & 4 & \ref{s3:lemL15p}, \ref{s3:thmT16s}, \ref{s3:lemHB}, \ref{s3:lemMonotone}, \ref{s1:citCor22}, \ref{s1:factEG0}, \ref{s1:citChernoff}, \ref{s1:condG3}, \ref{s4:convVortex}, \ref{s1:convGraphs}, \ref{s1:defConstants}\\
\texttt{s4:remTPVVX} & TPV versus VX$^+$ (not used) (\ref{s4:remTPVVX}) & 4 & \ref{s4:lemTPV}, \ref{s4:thmVXp}\\
\hline
\end{tabular}}
\end{center}

\begin{center}
{\footnotesize
\setlength{\tabcolsep}{3pt}
\begin{tabular}{p{3.2cm}p{3.6cm}p{0.5cm}p{7.6cm}}
\multicolumn{4}{l}{(Table~\ref{s1:tabDAG}, continued.)}\\
\textbf{Label} & \textbf{Name} & \textbf{S.} & \textbf{Depends on}\\
\hline
\texttt{s5:defZones} & zones (\ref{s5:defZones}) & 5 & \ref{s2:defAncestors}, \ref{s3:defCOL}, \ref{s2:propStructure}, \ref{s2:lemTower}, \ref{s1:condG1}\\
\texttt{s5:lemZones} & zone lemma (\ref{s5:lemZones}) & 5 & \ref{s5:defZones}, \ref{s2:propStructure}, \ref{s2:lemTower}, \ref{s1:condG1}, \ref{s3:defCOL}\\
\texttt{s5:defStages} & stage-2 statuses (\ref{s5:defStages}) & 5 & \ref{s5:defZones}, \ref{s3:defCOL}, \ref{s3:lemHB}, \ref{s2:propParentless}, \ref{s2:defAncestors}, \ref{s2:propStructure}, \ref{s2:lemTower}, \ref{s1:citDef7}\\
\texttt{s5:lemE1} & bad probabilities (\ref{s5:lemE1}) & 5 & \ref{s3:lemCOL}, \ref{s3:lemHB}, \ref{s5:lemZones}, \ref{s5:defStages}, \ref{s1:citChernoffGen}, \ref{s3:defCOL}, \ref{s3:lemL15p}, \ref{s3:lemCOLJV}, \ref{s2:propStructure}, \ref{s2:lemTower}, \ref{s1:condG1}\\
\texttt{s5:lemExpect} & expected bad mass (\ref{s5:lemExpect}) & 5 & \ref{s5:lemE1}, \ref{s5:defStages}, \ref{s2:propOV}, \ref{s2:lemTower}, \ref{s2:lemLacunary}, \ref{s1:condG1}, \ref{s2:propStructure}\\
\texttt{s5:lemChild} & child side (\ref{s5:lemChild}) & 5 & \ref{s4:lemPV}, \ref{s5:defStages}, \ref{s5:lemZones}, \ref{s3:lemCOL}, \ref{s2:propStructure}, \ref{s3:defCOL}, \ref{s1:condG3}\\
\texttt{s5:lemParent} & parent side (\ref{s5:lemParent}) & 5 & \ref{s5:lemChild}, \ref{s5:defStages}, \ref{s5:lemZones}, \ref{s3:defCOL}, \ref{s3:lemCOLJV}, \ref{s3:lemMonotone}, \ref{s1:citEuler}, \ref{s2:propOV}, \ref{s2:lemTower}, \ref{s2:propStructure}, \ref{s1:condG1}, \ref{s1:citDef7}\\
\texttt{s5:lemDemoted} & demoted parts (\ref{s5:lemDemoted}) & 5 & \ref{s4:thmVXp}, \ref{s2:propStructure}, \ref{s3:lemCOLJV}, \ref{s2:lemTower}, \ref{s1:condG1}, \ref{s1:condG3}, \ref{s5:defStages}, \ref{s5:lemParent}\\
\texttt{s5:lemKRED} & Lemma K-RED (\ref{s5:lemKRED}) & 5 & \ref{s5:lemChild}, \ref{s5:lemParent}, \ref{s5:lemDemoted}, \ref{s2:propStructure}, \ref{s2:propParentless}, \ref{s2:lemEL}, \ref{s3:defCOL}, \ref{s5:defStages}\\
\texttt{s5:remConstants} & 369 and 745 (\ref{s5:remConstants}) & 5 & \ref{s4:lemPV}, \ref{s5:lemKRED}, \ref{s2:propStructure}, \ref{s2:propParentless}, \ref{s2:propOV}, \ref{s2:lemTower}, \ref{s4:thmVXp}, \ref{s1:condG4}, \ref{s1:factEG0}\\
\hline
\end{tabular}}
\end{center}

\begin{center}
{\footnotesize
\setlength{\tabcolsep}{3pt}
\begin{tabular}{p{3.2cm}p{3.6cm}p{0.5cm}p{7.6cm}}
\multicolumn{4}{l}{(Table~\ref{s1:tabDAG}, continued.)}\\
\textbf{Label} & \textbf{Name} & \textbf{S.} & \textbf{Depends on}\\
\hline
\texttt{s6:lemGATE} & Lemma GATE (\ref{s6:lemGATE}) & 6 & \ref{s1:defObject}\\
\texttt{s6:defCluster} & clusters (\ref{s6:defCluster}) & 6 & --\\
\texttt{s6:lemMED} & Lemma MED (\ref{s6:lemMED}) & 6 & \ref{s6:defCluster}\\
\texttt{s6:lemEQLPT} & Lemma EQ-LPT (\ref{s6:lemEQLPT}) & 6 & --\\
\texttt{s6:lemPAR} & Lemma PAR (\ref{s6:lemPAR}) & 6 & \ref{s1:citEuler}\\
\texttt{s6:lemHCCP} & Lemma HCC-P (\ref{s6:lemHCCP}) & 6 & \ref{s6:lemGATE}, \ref{s6:lemMED}, \ref{s6:defCluster}\\
\texttt{s6:lemHCCglob} & union of systems (\ref{s6:lemHCCglob}) & 6 & \ref{s1:factAdd}, \ref{s6:lemHCCP}\\
\texttt{s6:defDesign} & designation, classes (\ref{s6:defDesign}) & 6 & \ref{s2:defAncestors}, \ref{s2:defHBtp}\\
\texttt{s6:thmCONC} & Theorem CONC (\ref{s6:thmCONC}) & 6 & \ref{s2:propOrigin}, \ref{s2:lemTower}, \ref{s2:propOV}, \ref{s2:lemLacunary}, \ref{s6:defDesign}, \ref{s2:defHBtp}, \ref{s1:condG1}\\
\texttt{s6:thmCONCL} & Theorem CONC-L (\ref{s6:thmCONCL}) & 6 & \ref{s2:propOrigin}, \ref{s2:lemSEP}, \ref{s2:lemThinCut}, \ref{s2:lemGC}, \ref{s2:lemEL}, \ref{s2:lemTower}, \ref{s2:lemLacunary}, \ref{s6:thmCONC}, \ref{s2:propOV}, \ref{s2:defHBtp}, \ref{s6:defDesign}, \ref{s1:condG1}\\
\hline
\end{tabular}}
\end{center}

\begin{center}
{\footnotesize
\setlength{\tabcolsep}{3pt}
\begin{tabular}{p{3.2cm}p{3.6cm}p{0.5cm}p{7.6cm}}
\multicolumn{4}{l}{(Table~\ref{s1:tabDAG}, continued.)}\\
\textbf{Label} & \textbf{Name} & \textbf{S.} & \textbf{Depends on}\\
\hline
\texttt{s6:defLending} & lending statuses (\ref{s6:defLending}) & 6 & \ref{s3:defCOL}, \ref{s3:lemCOL}, \ref{s5:defStages}, \ref{s6:defDesign}, \ref{s5:defZones}\\
\texttt{s6:lemLost} & the sets Lost (\ref{s6:lemLost}) & 6 & \ref{s6:defLending}, \ref{s3:lemCOL}, \ref{s5:lemE1}, \ref{s5:lemExpect}, \ref{s2:propStructure}, \ref{s2:lemTower}, \ref{s3:defCOL}\\
\texttt{s6:lemJSLC} & Lemma JS-LC (\ref{s6:lemJSLC}) & 6 & \ref{s6:lemPAR}, \ref{s6:lemMED}, \ref{s6:lemEQLPT}, \ref{s6:lemHCCP}, \ref{s6:lemHCCglob}, \ref{s6:defDesign}, \ref{s6:defLending}, \ref{s3:lemCOL}, \ref{s3:lemMonotone}, \ref{s2:lemEL}, \ref{s2:lemCap}, \ref{s2:propOV}, \ref{s2:lemTower}, \ref{s2:defHBtp}, \ref{s2:propStructure}, \ref{s3:defCOL}, \ref{s1:citDef7}\\
\texttt{s6:remStar} & star split (not used) (\ref{s6:remStar}) & 6 & \ref{s6:lemJSLC}\\
\texttt{s6:lemJplus} & Lemma J$^+$ (\ref{s6:lemJplus}) & 6 & \ref{s6:lemJSLC}, \ref{s6:lemPAR}, \ref{s6:defDesign}, \ref{s2:propStructure}, \ref{s2:lemEL}, \ref{s2:lemCap}, \ref{s3:defCOL}, \ref{s2:defHBtp}, \ref{s6:lemHCCP}, \ref{s6:defLending}\\
\texttt{s6:defJconsumer} & J-consumer (\ref{s6:defJconsumer}) & 6 & \ref{s3:defCOL}, \ref{s6:lemJplus}\\
\texttt{s6:consOrder} & processing order (\ref{s6:consOrder}) & 6 & \ref{s4:lemTPV}, \ref{s5:lemChild}, \ref{s5:lemParent}, \ref{s5:lemDemoted}, \ref{s6:lemJSLC}, \ref{s6:defJconsumer}, \ref{s6:defLending}, \ref{s5:defStages}, \ref{s4:thmVXp}, \ref{s3:defCOL}\\
\texttt{s6:lemLent} & Lent sets (\ref{s6:lemLent}) & 6 & \ref{s6:consOrder}, \ref{s6:defLending}, \ref{s6:defJconsumer}, \ref{s6:lemJSLC}, \ref{s5:lemParent}, \ref{s5:lemKRED}, \ref{s3:defCOL}, \ref{s6:lemJplus}, \ref{s5:defStages}, \ref{s4:lemTPV}, \ref{s5:lemChild}, \ref{s5:lemDemoted}, \ref{s4:thmVXp}\\
\texttt{s6:thmMIXC} & Theorem MIX-C (\ref{s6:thmMIXC}) & 6 & \ref{s6:consOrder}, \ref{s4:lemTPV}, \ref{s5:lemKRED}, \ref{s5:lemDemoted}, \ref{s6:lemJSLC}, \ref{s6:lemJplus}, \ref{s6:defJconsumer}, \ref{s6:lemLent}, \ref{s6:thmCONCL}, \ref{s2:propStructure}, \ref{s2:propParentless}, \ref{s2:propOV}, \ref{s2:lemTower}, \ref{s3:defCOL}, \ref{s3:lemCOL}, \ref{s4:lemPV}, \ref{s5:lemChild}, \ref{s4:thmVXp}, \ref{s2:defHBtp}, \ref{s1:condG3}, \ref{s6:defDesign}, \ref{s6:defLending}\\
\hline
\end{tabular}}
\end{center}

\begin{center}
{\footnotesize
\setlength{\tabcolsep}{3pt}
\begin{tabular}{p{3.2cm}p{3.6cm}p{0.5cm}p{7.6cm}}
\multicolumn{4}{l}{(Table~\ref{s1:tabDAG}, continued.)}\\
\textbf{Label} & \textbf{Name} & \textbf{S.} & \textbf{Depends on}\\
\hline
\texttt{s7:defPool} & pool labels (\ref{s7:defPool}) & 7 & \ref{s2:lemTower}, \ref{s3:defCOL}\\
\texttt{s7:defCand} & candidates (\ref{s7:defCand}) & 7 & \ref{s7:defPool}, \ref{s3:defCOL}, \ref{s6:defDesign}\\
\texttt{s7:lemCand} & candidate counts (\ref{s7:lemCand}) & 7 & \ref{s7:defCand}, \ref{s7:defPool}, \ref{s3:defCOL}, \ref{s3:lemCOLJV}, \ref{s2:propStructure}, \ref{s2:lemTower}, \ref{s1:citChernoffGen}\\
\texttt{s7:defSchedule} & randomness schedule (\ref{s7:defSchedule}) & 7 & \ref{s3:defCOL}, \ref{s5:defZones}, \ref{s7:defPool}, \ref{s6:defLending}, \ref{s5:defStages}, \ref{s7:defCand}, \ref{s4:lemTPV}, \ref{s4:lemPV}, \ref{s4:thmVXp}, \ref{s5:lemChild}, \ref{s5:lemDemoted}\\
\texttt{s7:consRound} & JV$^{+*}$ round step (\ref{s7:consRound}) & 7 & \ref{s6:lemJplus}, \ref{s6:defJconsumer}, \ref{s7:defCand}, \ref{s7:defPool}, \ref{s7:defSchedule}, \ref{s1:citMarkov}\\
\texttt{s7:lemWellDef} & round step well defined (\ref{s7:lemWellDef}) & 7 & \ref{s7:consRound}, \ref{s6:lemJplus}, \ref{s7:lemCand}, \ref{s3:lemCOLJV}, \ref{s1:citHall}, \ref{s2:defHBtp}\\
\texttt{s7:lemMULT} & Lemma MULT (\ref{s7:lemMULT}) & 7 & \ref{s7:consRound}, \ref{s6:lemJplus}, \ref{s7:defCand}, \ref{s7:defPool}, \ref{s2:propOV}\\
\texttt{s7:lemSimple} & $Q_l$ is simple (\ref{s7:lemSimple}) & 7 & \ref{s7:consRound}\\
\hline
\end{tabular}}
\end{center}

\begin{center}
{\footnotesize
\setlength{\tabcolsep}{3pt}
\begin{tabular}{p{3.2cm}p{3.6cm}p{0.5cm}p{7.6cm}}
\multicolumn{4}{l}{(Table~\ref{s1:tabDAG}, continued.)}\\
\textbf{Label} & \textbf{Name} & \textbf{S.} & \textbf{Depends on}\\
\hline
\texttt{s7:lemLift} & Lemma JV-L (lift) (\ref{s7:lemLift}) & 7 & \ref{s7:consRound}, \ref{s7:lemSimple}, \ref{s7:lemWellDef}, \ref{s6:lemJplus}, \ref{s6:defJconsumer}, \ref{s3:defCOL}, \ref{s2:propStructure}, \ref{s6:defLending}\\
\texttt{s7:lemCC} & Lemma CC (\ref{s7:lemCC}) & 7 & \ref{s7:consRound}, \ref{s7:lemWellDef}, \ref{s1:citChernoffGen}, \ref{s6:lemJplus}\\
\texttt{s7:lemUltra} & ultra-hub copies (\ref{s7:lemUltra}) & 7 & \ref{s7:consRound}, \ref{s7:lemWellDef}, \ref{s6:lemJplus}, \ref{s1:citChernoffGen}, \ref{s7:lemMULT}, \ref{s7:lemCand}\\
\texttt{s7:lemVstar} & the weight $\XV$ (\ref{s7:lemVstar}) & 7 & \ref{s7:defPool}, \ref{s2:propOV}, \ref{s2:lemTower}\\
\texttt{s7:lemPay} & payments (\ref{s7:lemPay}) & 7 & \ref{s7:consRound}, \ref{s7:lemWellDef}, \ref{s6:lemJplus}, \ref{s2:propParentless}, \ref{s2:lemTower}\\
\texttt{s7:defXprime} & the functional $\Xp$ (\ref{s7:defXprime}) & 7 & \ref{s6:defLending}, \ref{s7:lemVstar}, \ref{s7:lemPay}\\
\texttt{s7:lemEXprime} & expectation of $\Xp$ (\ref{s7:lemEXprime}) & 7 & \ref{s7:defXprime}, \ref{s6:lemLost}, \ref{s7:lemCand}, \ref{s7:lemVstar}, \ref{s2:propOV}, \ref{s2:propStructure}, \ref{s2:lemTower}, \ref{s5:lemExpect}, \ref{s2:lemCap}\\
\texttt{s7:lemUHsplit} & Lemma UH$^*$-split (\ref{s7:lemUHsplit}) & 7 & \ref{s7:lemCC}, \ref{s7:lemMULT}, \ref{s7:lemUltra}, \ref{s7:lemVstar}, \ref{s7:lemEXprime}, \ref{s7:defSchedule}, \ref{s2:propOV}, \ref{s2:lemTower}, \ref{s2:lemLacunary}, \ref{s1:condG4}, \ref{s1:condG1}, \ref{s7:consRound}, \ref{s7:lemWellDef}, \ref{s2:defHBtp}, \ref{s6:lemJplus}\\
\hline
\end{tabular}}
\end{center}

\begin{center}
{\footnotesize
\setlength{\tabcolsep}{3pt}
\begin{tabular}{p{3.2cm}p{3.6cm}p{0.5cm}p{7.6cm}}
\multicolumn{4}{l}{(Table~\ref{s1:tabDAG}, continued.)}\\
\textbf{Label} & \textbf{Name} & \textbf{S.} & \textbf{Depends on}\\
\hline
\texttt{s7:lemOneOutcome} & one joint outcome (\ref{s7:lemOneOutcome}) & 7 & \ref{s7:defSchedule}, \ref{s7:consRound}, \ref{s6:consOrder}, \ref{s4:lemTPV}, \ref{s4:lemPV}, \ref{s4:thmVXp}, \ref{s5:lemChild}, \ref{s5:lemDemoted}, \ref{s7:lemEXprime}, \ref{s7:lemPay}, \ref{s7:lemUHsplit}, \ref{s1:citMarkov}, \ref{s1:condG3}\\
\texttt{s7:propCost} & cost (\ref{s7:propCost}) & 7 & \ref{s6:thmMIXC}, \ref{s6:thmCONCL}, \ref{s7:lemLift}, \ref{s7:lemPay}, \ref{s7:lemEXprime}, \ref{s7:lemOneOutcome}, \ref{s5:lemKRED}, \ref{s2:propParentless}, \ref{s2:propOV}, \ref{s7:consRound}, \ref{s2:lemTower}, \ref{s1:condG4}, \ref{s6:defLending}, \ref{s7:defXprime}\\
\texttt{s7:thmJVps} & Theorem JV$^{+*}$ (\ref{s7:thmJVps}) & 7 & \ref{s7:propCost}, \ref{s7:lemUHsplit}, \ref{s7:lemSimple}, \ref{s7:lemLift}, \ref{s7:lemOneOutcome}, \ref{s6:thmMIXC}, \ref{s2:propExists}, \ref{s7:lemPay}, \ref{s7:lemUltra}, \ref{s1:condGamma}\\
\texttt{s7:thmHI} & Theorem HI$''$ (\ref{s7:thmHI}) & 7 & \ref{s1:factAdd}, \ref{s1:defObject}, \ref{s1:defConstants}\\
\texttt{s7:lemGammaSat} & satisfiability of the conditions (\ref{s7:lemGammaSat}) & 7 & \ref{s1:condGamma}, \ref{s7:propCost}, \ref{s7:lemUHsplit}, \ref{s5:lemExpect}, \ref{s5:lemKRED}, \ref{s2:lemTower}, \ref{s7:lemEXprime}, \ref{s6:thmCONCL}, \ref{s1:defConstants}, \ref{s3:lemCOL}, \ref{s5:lemE1}, \ref{s3:lemCOLJVev}\\
\texttt{s7:thmMainProof} & proof of the Main Theorem (\ref{s7:thmMainProof}) & 7 & \ref{s7:thmHI}, \ref{s7:thmJVps}, \ref{s2:propExists}, \ref{s7:lemGammaSat}, \ref{s6:defDesign}, \ref{s1:condGamma}, \ref{s1:defConstants}\\
\texttt{s7:remNonCirc} & non-circularity (\ref{s7:remNonCirc}) & 7 & \ref{s7:thmHI}, \ref{s7:thmJVps}, \ref{s7:propCost}, \ref{s7:lemUHsplit}, \ref{s7:lemCC}, \ref{s7:lemUltra}, \ref{s7:lemPay}, \ref{s7:lemSimple}, \ref{s7:consRound}, \ref{s4:lemTPV}, \ref{s4:lemPV}, \ref{s4:thmVXp}, \ref{s1:condG4}, \ref{s1:factEG0}\\
\hline
\end{tabular}}
\end{center}

Every edge of this table points from a statement to one appearing earlier in the
paper; the only exception is the Main Theorem, stated in
Section~\ref{s1:sec} and proved in Section~\ref{s7:sec}.

\section{The hierarchy \texorpdfstring{$\HBtp$}{HB*tau+}}\label{s2:sec}

This section defines the hierarchy $\HBtp$ on which the whole construction
runs, and proves the structural facts about it that later sections use. The
hierarchy follows the round structure of the proof of \cite[Theorem~2]{BM}:
in each round, long cycles are removed, the remaining graph is cut by the
recursion of \cite[Lemma~14]{BM} with $s=0$ into pieces, and each large piece
is cut again into robust expanders. There are two changes inside a round.
First, the robust recursion is replaced by Lemma~14$^\tau$ below: at every
split, vertices with many deleted edges across the cut are moved into the
separator (the \emph{$\tau$-rules}), which caps the number of deleted edges
from any vertex into the non-duplicated vertices of a leaf (the thin-cut
lemma). Second, a \emph{guest cap}
(rule~(GC)) is added to the definition of light pre-parts. The leaves of this
recursion are still certified expanders, and the overlap between leaves stays
small (Lemma OV).

All constants in this section are the constants of $\HBtp$. The
hierarchy without the $\tau$-rules and without (GC) (``plain $\HBs$''), and the
hierarchy with the $\tau$-rules but without (GC) (``$\HBt$''), appear only in
the final remark of this section, for comparison. Round parameters are written
$\tHB_l$, $\tau_r$ and $\thGC_r$.

\emph{Standing assumption.} Throughout this section, as everywhere in this
text, $\Dstar$ denotes the constant of Definition~\ref{s1:defConstants}(iii):
it satisfies \Gc{1}--\Gc{4} (Condition~\ref{s1:condGamma}). The proofs use
these conditions either directly or through the results they apply; for
instance the termination part of Proposition~\ref{s2:propExists} uses
Proposition~\ref{s2:propDegRec}, and hence \Gc{1}.

\emph{Local notation.} In the first two subsections of this section the
letter $H$ denotes a node graph with $m=|H|$ vertices, and $(U,F)$, $N$,
$F_0$, $F_1$, $F''$, $\Hout$, $\Hin$, $U'$, $N'$, $N''$, $G_1$, $G_2$, $n_1$,
$n_2$, $\nu^*$ are local to these subsections, as are $S$ (a total leaf
size), $M$ (a size threshold) and $c$ (a constant) in
Lemma OV. Throughout, $\eps=2^{-5}$, $\log=\log_2$, and
$|H|\defeq|V(H)|$ (Convention~\ref{s1:convGraphs}).

\subsection{Witnesses, the \texorpdfstring{$\tau$}{tau}-rules and separation}

\begin{definition}[witness and its minimal part]\label{s2:defWitness}
Let $s\ge0$ and let $H$ be a graph on $m=|H|$ vertices that is not an
$(\eps,s)$-expander (Cited result~\ref{s1:citDef11}). Then $m\ge2$, since a
graph on one vertex has no set $U$ with $1\le|U|\le2m/3$. A \emph{witness} at
$H$ is a pair $(U,F)$ with $U\subseteq V(H)$ and $F\subseteq E(H)$ such that
\[
1\le|U|\le\tfrac23m,\qquad |F|\le s|U|,\qquad
|\Nbr_{H-F}(U)|<\frac{\eps|U|}{\log^2m}.
\]
Since $H$ is not an $(\eps,s)$-expander, a witness exists. For a witness
$(U,F)$ put
\[
N\defeq\Nbr_{H-F}(U),\qquad F_0\defeq E_H\bigl(U,\,V(H)\setminus(U\cup
N)\bigr).
\]
\emph{Fact.} $F_0\subseteq F$ and $\Nbr_{H-F_0}(U)=N$. In particular
$(U,F_0)$ is again a witness at $H$, with the same set $N$.
\deps{\ref{s1:citDef11}, \ref{s1:defConstants}}
\end{definition}

\begin{proof}[Proof of the Fact]
Let $uv\in F_0$ with $u\in U$ and $v\notin U\cup N$. If $uv\notin F$, then
$uv$ is an edge of $H-F$, so $v\in\Nbr_{H-F}(U)=N$, a contradiction. Hence
$F_0\subseteq F$. Consequently $H-F\subseteq H-F_0$, so
$N\subseteq\Nbr_{H-F_0}(U)$. Conversely, let $v\notin U$ have a neighbour
$u\in U$ with $uv\notin F_0$. If $v\notin N$ then $uv\in
E_H(U,V(H)\setminus(U\cup N))=F_0$, a contradiction; so $v\in N$. Thus
$\Nbr_{H-F_0}(U)=N$. Finally $|F_0|\le|F|\le s|U|$, so $(U,F_0)$ is a witness.
\end{proof}

\begin{definition}[the $\tau$-rules and the split]\label{s2:defTauRules}
Let $H$, $m$ and $s$ be as in Definition~\ref{s2:defWitness}, let $(U,F)$ be a
witness at $H$ with $N$ and $F_0$ as there, write $V\defeq V(H)$, and let
$\tau>s$ be real. The \emph{$\tau$-rules} are:
\begin{enumerate}
\item[(1)] $\Hout\defeq\{v\in U:\deg_{F_0}(v)\ge\tau\}$,\quad
$U'\defeq U\setminus\Hout$,\quad $N'\defeq N\cup\Hout$;
\item[(2)] $F_1\defeq E_H\bigl(U',V\setminus(U'\cup N')\bigr)$,\quad
$\Hin\defeq\{x\in V\setminus(U'\cup N'):\deg_{F_1}(x)\ge\tau\}$,\quad
$N''\defeq N'\cup\Hin$,\quad $F''\defeq E_H\bigl(U',V\setminus(U'\cup
N'')\bigr)$.
\end{enumerate}
The \emph{split of $H$ by the $\tau$-rules} has the two children
$G_1\defeq H[U'\cup N'']-F''$ and $G_2\defeq H-U'-E(G_1)-F''$, and its
\emph{deleted set} is $F''$. Every edge of $F''$ has one end in $U'$ and the
other outside $U'\cup N''$, so no edge of $F''$ lies inside $U'\cup N''$, and
every edge of $F''$ meets $U'$. Hence
\begin{equation}\label{s2:eqSplit}
G_1=H[U'\cup N''],\qquad G_2=H[V\setminus U']-E\bigl(H[N'']\bigr).
\end{equation}
\emph{Facts.}
\begin{enumerate}
\item[(a)] $U'\cap N''=\emptyset$, $U'\cup N'=U\cup N$, and
$F''\subseteq F_1\subseteq F_0\subseteq F$.
\item[(b)] $E(H)=E(G_1)\sqcup E(G_2)\sqcup F''$ (disjoint union),
$V(G_1)=U'\cup N''$ and $V(G_2)=V\setminus U'$. A vertex of $U'$ lies only in
$G_1$, a vertex of $N''$ lies in both children, and a vertex of
$V\setminus(U'\cup N'')$ lies only in $G_2$.
\item[(c)] Given the pair $(U,N)$, the objects $\Hout$, $U'$, $N'$,
$F_1$, $\Hin$, $N''$, $F''$, $G_1$, $G_2$ depend only on $(U,N)$, because
they are defined from $F_0=E_H(U,V\setminus(U\cup N))$. The choice of $F$
matters only through $N=\Nbr_{H-F}(U)$: every witness $(U,F)$ yields the
same split as its minimal part $(U,F_0)$. (It is not true that the choice of
$F$ is irrelevant, since $F$ determines $N$.)
\end{enumerate}
\deps{\ref{s2:defWitness}}
\end{definition}

\begin{proof}[Proof of the Facts]
(a) $U'\subseteq U$ and $N\cap U=\emptyset$, so $U'\cap N=\emptyset$;
$\Hout\cap U'=\emptyset$ by definition; and $\Hin\cap(U'\cup N')=\emptyset$ by
definition. Hence $U'\cap N''=\emptyset$. As $\Hout\subseteq U$, we get
$U'\cup N'=(U\setminus\Hout)\cup N\cup\Hout=U\cup N$. Therefore
$F_1=E_H(U',V\setminus(U\cup N))\subseteq E_H(U,V\setminus(U\cup N))=F_0$, and
since $N'\subseteq N''$, $F''=E_H(U',V\setminus(U'\cup N''))\subseteq
E_H(U',V\setminus(U'\cup N'))=F_1$. Finally $F_0\subseteq F$ by
Definition~\ref{s2:defWitness}.

(b) The vertex sets follow from \eqref{s2:eqSplit}, and the statement about
where vertices go follows from $U'\cap N''=\emptyset$. Let $e=xy\in E(H)$. If
$x,y\in U'\cup N''$ then $e\in E(G_1)$, and $e\notin F''$. Otherwise one end,
say $y$, lies outside $U'\cup N''$. If $x\in U'$ then $e\in F''$. If
$x\notin U'$ then both ends lie in $V\setminus U'$ and not both lie in $N''$,
so $e\in E(G_2)$; and $e\notin F''$ because $e$ does not meet $U'$. An edge in
$E(G_1)\cap E(G_2)$ would have both ends in $(U'\cup N'')\cap(V\setminus
U')=N''$, which \eqref{s2:eqSplit} excludes from $G_2$. So the three sets are
disjoint and cover $E(H)$.

(c) $F_0$ is determined by $(U,N)$, and every later object in (1)--(2) and in
\eqref{s2:eqSplit} is defined from $U$, $N$ and $F_0$ alone. By the Fact of
Definition~\ref{s2:defWitness}, $(U,F_0)$ is a witness with the same $N$.
\end{proof}

\begin{lemma}[Lemma SEP]\label{s2:lemSEP}
A \emph{split recursion} on a graph $H_0$ is a finite rooted binary tree whose
nodes $\nu$ carry graphs $H_\nu$, with $H_\nu=H_0$ at the root, such that at
every non-leaf node $\nu$ there are disjoint sets $U'_\nu,N''_\nu\subseteq
V(H_\nu)$ for which the two children $\nu_1,\nu_2$ of $\nu$ carry
\[
H_{\nu_1}=H_\nu[U'_\nu\cup N''_\nu],\qquad
H_{\nu_2}=H_\nu[V(H_\nu)\setminus U'_\nu]-E\bigl(H_\nu[N''_\nu]\bigr),
\]
and the \emph{deleted set} at $\nu$ is $F''_\nu\defeq
E_{H_\nu}(U'_\nu,V(H_\nu)\setminus(U'_\nu\cup N''_\nu))$. (By
\eqref{s2:eqSplit}, every split by the $\tau$-rules has this form.) We write
$|\nu|\defeq|H_\nu|$, say that a vertex $v$ \emph{lies in} $\nu$ if $v\in
V(H_\nu)$, and call the leaf nodes \emph{leaves} (distinct leaf nodes are
distinct leaves even if their vertex sets coincide). A \emph{deleted edge} is an
edge of some $F''_\nu$. Let $\Dup$ be the set of vertices lying in at least two
leaves. Then:
\begin{enumerate}
\item[(0)] at every non-leaf $\nu$, $E(H_\nu)=E(H_{\nu_1})\sqcup
E(H_{\nu_2})\sqcup F''_\nu$; a vertex of $U'_\nu$ lies only in $\nu_1$, a
vertex of $N''_\nu$ lies in both children, every other vertex of $\nu$ lies
only in $\nu_2$, and $|\nu_1|+|\nu_2|=|\nu|+|N''_\nu|$. Every vertex of a node
lies in at least one leaf below it, and the total size of the leaves is
$|H_0|+\sum_{\nu}|N''_\nu|$, the sum over the non-leaf nodes;
\item[(i)] every edge $xy$ of $H_0$ either lies in exactly one leaf, which
contains $x$ and $y$, or is deleted at exactly one node $\nu$. In the second
case $x$ and $y$ lie in $\nu$ and go to different children of $\nu$, each to
only one child, and $\nu$ is the first node on the path of nodes having $xy$
as an edge at which this happens;
\item[(ii)] if $u\notin\Dup$ and $u\in V(H_0)$, then the nodes in which $u$
lies are exactly the ancestors of the unique leaf $\Leaf_u$ containing $u$
($\Leaf_u$ included), and $u\notin N''_\nu$ at every node $\nu$ on that path;
\item[(iii)] if $u\notin\Dup$, $u,h\in V(H_0)$, $h\notin V(\Leaf_u)$, and $\nu^*$ is the
deepest ancestor of $\Leaf_u$ in which $h$ lies, then every deleted edge
$hu'$ with $u'\in V(\Leaf_u)\setminus\Dup$ is deleted at $\nu^*$.
\end{enumerate}
\deps{\ref{s2:defTauRules}}
\end{lemma}

\begin{proof}
(0) Fix a non-leaf $\nu$ and write $H=H_\nu$, $V=V(H)$, $U'=U'_\nu$,
$N''=N''_\nu$. The proof of Fact~(b) of Definition~\ref{s2:defTauRules} used
only that $U'\cap N''=\emptyset$ and the formulas \eqref{s2:eqSplit}, so it
gives $E(H)=E(H_{\nu_1})\sqcup E(H_{\nu_2})\sqcup F''_\nu$ and the statement
about where vertices go. As $U'\cap N''=\emptyset$,
$|\nu_1|+|\nu_2|=|U'|+|N''|+|V|-|U'|=|\nu|+|N''|$. Every vertex of $\nu$ lies
in a child of $\nu$, so by induction on the height of the subtree every vertex
of a node lies in some leaf below it. Summing $|\nu_1|+|\nu_2|-|\nu|=|N''_\nu|$
over all non-leaf nodes gives the total leaf size.

(i) Let $\Omega_{xy}$ be the set of nodes having $xy$ as an edge; the root is in $\Omega_{xy}$.
Every edge of a child is an edge of its parent, so $\Omega_{xy}$ is closed under taking
parents. By (0), at a non-leaf $\nu\in\Omega_{xy}$ the edge $xy$ lies in exactly one of
$E(H_{\nu_1})$, $E(H_{\nu_2})$ and $F''_\nu$; so at most one child of $\nu$ is
in $\Omega_{xy}$, and none is exactly when $xy\in F''_\nu$. Hence $\Omega_{xy}$ is a path from the
root ending at a node $\nu_\circ$. If $\nu_\circ$ is a leaf, then $xy$ is an
edge of exactly one leaf, namely $\nu_\circ$ (other leaves are not in $\Omega_{xy}$), and
$xy$ is deleted nowhere (at nodes of $\Omega_{xy}\setminus\{\nu_\circ\}$ it passes to a
child; at nodes outside $\Omega_{xy}$ it is not an edge). If $\nu_\circ$ is not a leaf,
then $xy\in F''_{\nu_\circ}$, $xy$ lies in no leaf, and it is deleted at no
other node. At $\nu_\circ$, one end lies in $U'_{\nu_\circ}$ and goes only to
the first child, and the other lies outside $U'_{\nu_\circ}\cup
N''_{\nu_\circ}$ and goes only to the second. At every earlier node of $\Omega_{xy}$ the
edge $xy$ passes to a child, which contains both $x$ and $y$.

(ii) By (0), $u$ lies in at least one leaf, hence in exactly one, $\Leaf_u$.
The nodes in which $u$ lies form a set closed under taking parents (vertex sets
shrink downwards), and at each non-leaf node of this set at least one child
contains $u$. If some node $\nu$ had both children containing $u$, then $u$
would lie in a leaf below each child, i.e.\ in two distinct leaves. So exactly
one child of every such non-leaf node contains $u$: the nodes containing $u$
form a path from the root to a leaf, which is $\Leaf_u$. A vertex of
$N''_\nu$ lies in both children of $\nu$, so $u\notin N''_\nu$ on this path.

(iii) As $h\in V(H_0)$, the root is an ancestor of $\Leaf_u$ in which $h$
lies, so $\nu^*$ exists. Let $u'\in V(\Leaf_u)\setminus\Dup$. By (ii) applied to $u'$, whose
unique leaf is $\Leaf_u$, the nodes containing $u'$ are exactly the ancestors
of $\Leaf_u$. Suppose $hu'$ is deleted at $\nu$. Then $h$ and $u'$ lie in
$\nu$, so $\nu$ is an ancestor of $\Leaf_u$ in which $h$ lies, and $\nu$ is
not below $\nu^*$. By (i), at $\nu$ the vertices $h$ and $u'$ go to different
children, each to only one. The vertex $u'$ goes to the child of $\nu$ on the
path to $\Leaf_u$, so $h$ does not lie in that child. Hence $\nu$ is the
deepest ancestor of $\Leaf_u$ containing $h$, that is, $\nu=\nu^*$.
\end{proof}

\begin{lemma}[thin cut]\label{s2:lemThinCut}
Let $\tau>0$ and consider a split recursion (Lemma~\ref{s2:lemSEP}) in which
every split is produced by the $\tau$-rules at threshold $\tau$: at each
non-leaf node $\nu$ there are $s_\nu<\tau$ and a witness at $H_\nu$ (with
parameter $s_\nu$) whose $\tau$-rules (Definition~\ref{s2:defTauRules}) give
$U'_\nu$ and $N''_\nu$. Then for every leaf $\Leaf$ and every vertex $h$,
\[
\#\bigl\{\text{deleted edges } hu:\ u\in V(\Leaf)\setminus\Dup\bigr\}\le
\lceil\tau\rceil-1,
\]
which is $\tau-1$ when $\tau$ is an integer (as in all applications), and this
number is $0$ if $h\in V(\Leaf)$. Moreover, if $u\notin\Dup$ then
every edge at $u$ is either deleted or an edge of the unique leaf containing
$u$.
\deps{\ref{s2:lemSEP}, \ref{s2:defTauRules}}
\end{lemma}

\begin{proof}
The condition $s_\nu<\tau$ only ensures that the $\tau$-rules are defined
(Definition~\ref{s2:defTauRules} asks for $\tau>s$); the argument below uses only
the structure of the $\tau$-rules.
Fix $\Leaf$ and $h$; we may assume $h$ lies in the root (otherwise no edge
of $H_0$ meets $h$) and that some deleted
edge $hu$ with $u\in V(\Leaf)\setminus\Dup$ exists. For such $u$ we have
$\Leaf_u=\Leaf$, so by Lemma~\ref{s2:lemSEP}(ii) the nodes containing $u$ are
the ancestors of $\Leaf$.

\emph{Case $h\in V(\Leaf)$.} Let $hu$ be deleted at $\nu$. Then $\nu$ is an
ancestor of $\Leaf$, and by Lemma~\ref{s2:lemSEP}(i) the vertices $h$ and $u$
go to different children of $\nu$, each to only one. But the child of $\nu$ on
the path to $\Leaf$ contains $V(\Leaf)\ni h,u$. This is a contradiction, so
the number is $0$.

\emph{Case $h\notin V(\Leaf)$.} Let $\nu^*$ be the deepest ancestor of
$\Leaf$ in which $h$ lies. By Lemma~\ref{s2:lemSEP}(iii), every deleted edge
$hu$ with $u\in V(\Leaf)\setminus\Dup$ is deleted at $\nu^*$, i.e.\ lies in
$F''=F''_{\nu^*}$. Use the notation of Definition~\ref{s2:defTauRules} at
$\nu^*$. The vertex $h$ does not lie in the child of $\nu^*$ on the path to
$\Leaf$ (the \emph{path child}), so $h\notin N''$, since vertices of $N''$ lie
in both children.

\emph{Case A: the path child is $G_1$.} Then $h\notin V(G_1)=U'\cup N''$; in
particular $h\notin U'\cup N'$ and $h\notin\Hin$, so by the definition of
$\Hin$ we have $\deg_{F_1}(h)<\tau$, i.e.\ $\deg_{F_1}(h)\le\lceil\tau\rceil-1$.
As $F''\subseteq F_1$, at most $\lceil\tau\rceil-1$ edges of $F''$ meet $h$.

\emph{Case B: the path child is $G_2$.} Then $h\notin V(G_2)=V\setminus U'$,
so $h\in U'=U\setminus\Hout$, and by the definition of $\Hout$,
$\deg_{F_0}(h)<\tau$. As $F''\subseteq F_0$, at most $\lceil\tau\rceil-1$ edges
of $F''$ meet $h$.

In both cases the count is at most $\lceil\tau\rceil-1$. The last assertion is
Lemma~\ref{s2:lemSEP}(i): an edge $uy$ that is not deleted lies in exactly one
leaf, which contains $u$ and is therefore the unique leaf $\Leaf_u$.
\end{proof}

\subsection{Overlap, and Lemma 14\texorpdfstring{$^\tau$}{-tau}}

\begin{lemma}[generic overlap bound, Lemma OV]\label{s2:lemOVgeneric}
Let $c>0$ with $3.42\,c\,\eps<1$, and put $\cOV\defeq\log_2(4/3)$. Consider a
split recursion (Lemma~\ref{s2:lemSEP}) on a root graph with $n_0\ge1$
vertices such that at every non-leaf node $\nu$, with $m\defeq|\nu|$,
\[
m\ge2,\qquad |\nu_1|\le\tfrac34m,\qquad
|N''_\nu|\le\frac{c\,\eps\,|U'_\nu|}{\log^2m}.
\]
(Recall that vertices of $U'_\nu$ go only to the first child $\nu_1$.) Let $S$
be the total size of the leaves. For $M\ge2$ let
$\Delta_{\ge M}\defeq\sum_\nu|N''_\nu|$, the sum over the non-leaf nodes $\nu$
with $|\nu|\ge M$ (the number of \emph{duplication events} at nodes of size at
least $M$; this is a number, not to be confused with the vertex set $\Dup$ of
Lemma~\ref{s2:lemSEP}), and for a vertex $v$ let $\dup_{\ge M}(v)$ be the number of
non-leaf nodes $\nu$ with $|\nu|\ge M$ and $v\in N''_\nu$ (so
$\sum_v\dup_{\ge M}(v)=\Delta_{\ge M}$). Then:
\begin{enumerate}
\item[(a)] $\displaystyle\Delta_{\ge M}\le
c\,\eps\Bigl(\frac1{\log^2M}+\frac1{\cOV\log M}\Bigr)S\le
\frac{3.42\,c\,\eps\,S}{\log M}$;
\item[(b)] for every vertex $v$, the number of leaves $\Leaf$ with
$|\Leaf|\ge M$ containing $v$ is at most $1+\dup_{\ge M}(v)$; consequently
$\sum_{\Leaf:|\Leaf|\ge M}|\Leaf|-\bigl|\bigcup_{\Leaf:|\Leaf|\ge
M}V(\Leaf)\bigr|\le\Delta_{\ge M}$;
\item[(c)] $S\le n_0/(1-3.42\,c\,\eps)$.
\end{enumerate}
\emph{Instance $c=1$ (the $s=0$ recursion).} Call a split recursion an
\emph{$s=0$ recursion} if at every non-leaf node $\nu$ there is a witness
$(U_\nu,F_\nu)$ at $H_\nu$ with parameter $s=0$ (Definition~\ref{s2:defWitness})
and $U'_\nu=U_\nu$, $N''_\nu=\Nbr_{H_\nu}(U_\nu)$. This is exactly the
recursion in the proof of \cite[Lemma~14]{BM} with $s=0$
(Cited result~\ref{s1:citLem14}): there $F_\nu=\emptyset$, the children are
$H_\nu[U_\nu\cup N_\nu]$ and $H_\nu\setminus U_\nu-E(H_\nu[U_\nu\cup
N_\nu])=H_\nu[V(H_\nu)\setminus U_\nu]-E(H_\nu[N_\nu])$, and the deleted set
$E_{H_\nu}(U_\nu,V(H_\nu)\setminus(U_\nu\cup N_\nu))$ is empty. (It is also
the split by the $\tau$-rules at $s=0$ for any $\tau>0$: then $F_0=\emptyset$
and $\Hout=\Hin=\emptyset$.) Every $s=0$ recursion satisfies the hypotheses
with $c=1$; hence its total leaf size is at most $n_0/(1-3.42\eps)\le1.12\,n_0$.
The lemma is applied with $c=1.6$ to the recursion of Lemma~14$^\tau$ below,
and with $c=1.6$ to the two-level recursion of a round of $\HBtp$.
\deps{\ref{s2:lemSEP}, \ref{s2:defWitness}, \ref{s1:citLem14}}
\end{lemma}

\begin{proof}
(a) If no non-leaf node has size at least $M$ there is nothing to prove. For
every non-leaf node $\nu$ with $|\nu|\ge M$ and every $u\in U'_\nu$, give the
pair $(\nu,u)$ the charge $c\eps/\log^2|\nu|$. The total charge at $\nu$ is
$c\eps|U'_\nu|/\log^2|\nu|\ge|N''_\nu|$, so $\Delta_{\ge M}$ is at most the
total charge. A \emph{leaf copy} is a pair $(\Leaf,v)$ with $\Leaf$ a leaf and
$v\in V(\Leaf)$; there are exactly $S$ leaf copies. For each charged pair
$(\nu,u)$, the vertex $u$ lies in $\nu_1$, hence (Lemma~\ref{s2:lemSEP}(0)) in
some leaf $\Leaf$ below $\nu_1$; fix one and map $(\nu,u)$ to the leaf copy
$(\Leaf,u)$. The pair $(\nu,u)$ is recovered from $u$ and $\nu$, and $\nu$ is
an ancestor of $\Leaf$; so each leaf copy $(\Leaf,u)$ receives the charges of
pairs $(\nu,u)$ for a set of ancestors $\nu$ of $\Leaf$ with $|\nu|\ge M$,
$u\in U'_\nu$ and $\Leaf$ below $\nu_1$. List these ancestors as
$\nu^{(1)},\nu^{(2)},\dots,\nu^{(k)}$ from $\Leaf$ upwards. For $i<k$, the
node $\nu^{(i)}$ lies strictly below $\nu^{(i+1)}$ on the path from
$\nu^{(i+1)}$ to $\Leaf$, which passes through the first child of
$\nu^{(i+1)}$; vertex sets shrink downwards, so
$|\nu^{(i)}|\le|\nu^{(i+1)}_1|\le\frac34|\nu^{(i+1)}|$. Hence
$|\nu^{(i)}|\ge(4/3)^{i-1}|\nu^{(1)}|\ge(4/3)^{i-1}M$, i.e.\
$\log|\nu^{(i)}|\ge\log M+(i-1)\cOV$. The charge received by $(\Leaf,u)$ is
therefore at most
\[
c\eps\sum_{j\ge0}\frac1{(\log M+j\cOV)^2}\le
c\eps\Bigl(\frac1{\log^2M}+\int_0^\infty\frac{dx}{(\log M+x\cOV)^2}\Bigr)
=c\eps\Bigl(\frac1{\log^2M}+\frac1{\cOV\log M}\Bigr),
\]
because the summand is decreasing in $j$. Summing over the $S$ leaf copies
gives the first bound. For $M\ge2$ we have $\log M\ge1$, so
$1/\log^2M\le1/\log M$, and $1+1/\cOV=3.4094\ldots\le3.42$.

(b) Let $\mathcal{C}_v$ be the set of nodes of size at least $M$ in which $v$ lies. It
is closed under taking parents (vertex sets, hence sizes, shrink downwards), so
if nonempty it is a rooted subtree containing the root. Every leaf of the
recursion of size at least $M$ that contains $v$ belongs to $\mathcal{C}_v$ and has no
children, so it is a leaf of $\mathcal{C}_v$. In a finite rooted tree the number of
leaves is $1+\sum_\nu(\kappa_\nu-1)$, the sum over the nodes with
$\kappa_\nu\ge1$ children; here $\kappa_\nu\le2$, and $\kappa_\nu=2$ means that
$v$ lies in both children of $\nu$, i.e.\ $v\in N''_\nu$
(Lemma~\ref{s2:lemSEP}(0)), with $|\nu|\ge M$. So the number of such leaves is
at most $1+\dup_{\ge M}(v)$. Summing $(\#\{\Leaf\ni v:|\Leaf|\ge M\}-1)^+\le
\dup_{\ge M}(v)$ over $v$ gives the consequence.

(c) All non-leaf nodes have size at least $2$, so by
Lemma~\ref{s2:lemSEP}(0), $S=n_0+\sum_\nu|N''_\nu|=n_0+\Delta_{\ge2}\le
n_0+3.42c\eps S$ by (a) with $M=2$. Rearranging gives (c).

\emph{Instance $c=1$.} By Definition~\ref{s2:defWitness} with $s=0$ we have
$F_\nu=\emptyset$, so $N_\nu=\Nbr_{H_\nu}(U_\nu)$, and every neighbour of
$U_\nu$ outside $U_\nu$ lies in $N_\nu$; this gives the three formulas stated.
The node has $m\ge2$ vertices (Definition~\ref{s2:defWitness}), and
$|N''_\nu|=|N_\nu|<\eps|U_\nu|/\log^2m$, which is the hypothesis with $c=1$;
and $|\nu_1|=|U_\nu|+|N_\nu|<(1+\eps)\frac23m<\frac34m$, using $\log m\ge1$.
Finally $1/(1-3.42/32)=1.1196\ldots\le1.12$.
\end{proof}

\begin{lemma}[Lemma 14$^\tau$]\label{s2:lem14tau}
Let $H_0$ be a graph on $n_0$ vertices, $s\ge1$ an integer, $\eps=2^{-5}$, and
$\tau\ge128\,s\log^2n_0$. The \emph{$\tau$-run} of $H_0$ (with parameters
$s,\tau$) is the recursion that, at a node $H$, stops if $H$ is an
$(\eps,s)$-expander (the node is then a leaf), and otherwise picks any witness
$(U,F)$ at $H$ with parameter $s$ (Definition~\ref{s2:defWitness}) and splits
$H$ by the $\tau$-rules at threshold $\tau$ (Definition~\ref{s2:defTauRules};
note $\tau\ge128s>s$ whenever $n_0\ge2$). Let $\Dup$ be the set of vertices
lying in at least two leaves. Then, for every choice of witnesses:
\begin{enumerate}
\item[(a)] at every split, at a node of size $m$,
\[
|\Hout|+|\Hin|\le\frac{|F_0|}{\tau}\le\frac{s|U|}{\tau}\le\frac{|U|}{128\log^2m},
\qquad |N''|<\frac{1.25\,\eps|U|}{\log^2m}<\frac{1.5\,\eps|U|}{\log^2m},
\]
$|U'|\ge\frac{127}{128}|U|>0$, $U\subseteq U'\cup N''$, $n_1\defeq|U'\cup
N''|\le0.698\,m<\frac34m$ and $n_2\defeq m-|U'|<m$. The recursion terminates;
it is a split recursion; at most $4sn_0\log n_0$ edges are deleted; the
remaining edges are partitioned into the leaves; every leaf is an
$(\eps,s)$-expander; and the total size of the leaves is at most
$2n_0-2n_0/(2+\log n_0)$;
\item[(b)] (OV at $1.6\eps$) at every split, $|N''|\le1.6\,\eps|U'|/\log^2m$,
the vertices of $U'$ go only to $G_1$, and $|G_1|\le\frac34m$. Hence, by
Lemma~\ref{s2:lemOVgeneric} with $c=1.6$, the total leaf size $S$ satisfies
$S\le n_0/(1-5.5\eps)\le1.21\,n_0$, and, by the first bound of
Lemma~\ref{s2:lemOVgeneric}(a) with $c=1.6$, for every $M\ge2$ the duplication at
nodes of size at least $M$ is $\Delta_{\ge M}\le5.46\,\eps S/\log M$;
\item[(c)] for every leaf $\Leaf$ and every vertex $h$, if $\tau>0$ then
$\#\{\text{deleted edges }hu:u\in V(\Leaf)\setminus\Dup\}\le\lceil\tau\rceil-1$
($=\tau-1$ for integer $\tau$); the number is $0$ if $h\in V(\Leaf)$;
\item[(d)] if $u\notin\Dup$, every edge at $u$ is deleted or lies in the unique
leaf containing $u$.
\end{enumerate}
\deps{\ref{s1:citLem14}, \ref{s1:citDef11}, \ref{s2:defTauRules}, \ref{s2:lemThinCut}, \ref{s2:lemOVgeneric}, \ref{s2:lemSEP}, \ref{s2:defWitness}}
\end{lemma}

\begin{proof}
If $n_0=1$ the root is an $(\eps,s)$-expander, no edge is deleted, and (a), (b),
(d) and the last clause of (c) are trivial (the bound in (c) assumes $\tau>0$);
let $n_0\ge2$. Every node of the recursion is a subgraph of $H_0$ on $m\le n_0$
vertices, so $\tau\ge128\,s\log^2m$ at every node. The argument follows the
proof of \cite[Lemma~14]{BM} (Cited result~\ref{s1:citLem14}) with the split
of Definition~\ref{s2:defTauRules} in place of the split of \cite{BM}.

\emph{Counting at one split.} Fix a non-leaf node $H$ on $m\ge2$ vertices,
with witness $(U,F)$, and use the notation of
Definition~\ref{s2:defTauRules}; put $V\defeq V(H)$. Every edge of $F_0$ has
exactly one end in $U$, so $\sum_{v\in U}\deg_{F_0}(v)=|F_0|$. Let
$a\defeq\sum_{v\in\Hout}\deg_{F_0}(v)$; then $\tau|\Hout|\le a$. By Fact~(a) of
Definition~\ref{s2:defTauRules},
$F_1=E_H(U',V\setminus(U\cup N))$ is $F_0$ minus the edges at $\Hout$, so
$|F_1|=|F_0|-a$. Every edge of $F_1$ has exactly one end outside $U'\cup
N'=U\cup N$, so $\tau|\Hin|\le\sum_{x\notin U\cup N}\deg_{F_1}(x)=|F_1|$.
Adding, $\tau(|\Hout|+|\Hin|)\le|F_0|\le|F|\le s|U|$. With
$\tau\ge128\,s\log^2m$ this gives $|\Hout|+|\Hin|\le|U|/(128\log^2m)$, and
$1/128=\eps/4$. Therefore
\[
|N''|=|N|+|\Hout|+|\Hin|<\frac{\eps|U|}{\log^2m}+\frac{\eps|U|}{4\log^2m}
=\frac{1.25\,\eps|U|}{\log^2m}.
\]
Since $|\Hout|\le|U|/128$, $|U'|\ge\frac{127}{128}|U|>0$. As
$\Hout\subseteq N'\subseteq N''$, $U=U'\cup\Hout\subseteq U'\cup N''$, so
$|U|\le n_1$. Also $n_1=|U'|+|N''|\le|U|+|N''|<(1+1.5\eps)|U|\le
\frac23(1+\frac{3}{64})m\le0.698\,m$, using $\log m\ge1$; and $n_2=|V\setminus
U'|=m-|U'|<m$. By Fact~(b) of Definition~\ref{s2:defTauRules},
$n_1+n_2=m+|N''|$. Finally $\log n_1\le\log m+\log0.698<\log m-\frac25$, since
$\log_20.698=-0.518\ldots$.

\emph{Termination.} Both children of every split have fewer vertices than
their parent, so every root-to-node path has at most $n_0$ nodes, and the
binary tree is finite. By \eqref{s2:eqSplit} the recursion is a split
recursion. Its leaves are $(\eps,s)$-expanders by the stopping rule, and by
Lemma~\ref{s2:lemSEP}(i) every edge of $H_0$ is either deleted (exactly once)
or lies in exactly one leaf.

\emph{The induction.} We show by induction on the height of the subtree below
a node $\nu$, with $m=|\nu|$, that the subtree deletes at most $4sm\log m$ edges
and its leaves have total size at most $2m-2m/(2+\log m)$. For a leaf this
holds, since nothing is deleted and $m\le2m-2m/(2+\log m)$ (as
$2+\log m\ge2$). Let $\nu$ be a non-leaf node, with $n_1,n_2$ as above
($1\le n_1,n_2<m$), and let $E^{\mathrm{del}}_i$ be the set of edges deleted in the subtree
below the $i$-th child; by induction $|E^{\mathrm{del}}_i|\le4sn_i\log n_i$. The edges
deleted in the subtree below $\nu$ are $F''\cup E^{\mathrm{del}}_1\cup E^{\mathrm{del}}_2$, and $|F''|\le
|F_0|\le s|U|\le sn_1$. As $|U|\le n_1$, we have $|N''|<1.5\eps n_1/\log^2m$.
Hence, as in \cite[inequality~(8)]{BM},
\begin{align*}
\frac{|F''|+|E^{\mathrm{del}}_1|+|E^{\mathrm{del}}_2|}{s}
&\le n_1+4n_1\log n_1+4n_2\log n_2
\le n_1+4n_1\bigl(\log m-\tfrac25\bigr)+4n_2\log m\\
&=4(n_1+n_2)\log m-\tfrac35n_1=4(m+|N''|)\log m-\tfrac35n_1\\
&<4m\log m+n_1\Bigl(\frac{6\eps}{\log m}-\frac35\Bigr)\le4m\log m,
\end{align*}
since $6\eps/\log m\le6/32<3/5$. For the leaf sizes put $L\defeq\log m\ge1$.
As $n_2<m$, $2n_2/(2+\log n_2)\ge2n_2/(2+L)$. As $\log n_1<L-\frac25$,
\[
\frac{2n_1}{2+\log n_1}>\frac{2n_1}{\frac85+L}=\frac{2n_1}{2+L}
+\frac{\frac45n_1}{(\frac85+L)(2+L)}\ge\frac{2n_1}{2+L}+\frac{n_1}{10L^2},
\]
where the last step is $8L^2\ge(\frac85+L)(2+L)$, i.e.\
$7L^2-3.6L-3.2\ge0$, which holds for $L\ge1$ (the left side is $0.2$ at
$L=1$ and increasing for $L\ge1$). By induction the total leaf size below
$\nu$ is therefore
\begin{align*}
\sum_{i=1}^2\Bigl(2n_i-\frac{2n_i}{2+\log n_i}\Bigr)
&<(n_1+n_2)\Bigl(2-\frac2{2+L}\Bigr)-\frac{n_1}{10L^2}
=(m+|N''|)\Bigl(2-\frac2{2+L}\Bigr)-\frac{n_1}{10L^2}\\
&\le2m-\frac{2m}{2+L}+2|N''|-\frac{n_1}{10L^2}\\
&<2m-\frac{2m}{2+L}+\Bigl(3\eps-\frac1{10}\Bigr)\frac{n_1}{L^2}
\le2m-\frac{2m}{2+L},
\end{align*}
as in \cite[inequality~(9)]{BM}, because $3\eps=3/32<1/10$. At the root this gives (a).

(b) By (a), $|N''|<1.5\eps|U|/\log^2m\le1.5\cdot\frac{128}{127}\eps|U'|/\log^2m
<1.52\,\eps|U'|/\log^2m$. (The first bound of (a) even gives
$|N''|<1.25\cdot\frac{128}{127}\eps|U'|/\log^2m<1.26\,\eps|U'|/\log^2m$; the
looser constant $1.6$ suffices.) Vertices of $U'$ go only to $G_1$ and $|G_1|=n_1\le
\frac34m$ (Fact~(b) of Definition~\ref{s2:defTauRules} and (a)); every
non-leaf node has $m\ge2$. As $3.42\cdot1.6\,\eps=0.171<1$,
Lemma~\ref{s2:lemOVgeneric} applies with $c=1.6$: by its part (c),
$S\le n_0/(1-5.472\eps)\le n_0/(1-5.5\eps)=1.2075\ldots\,n_0\le1.21\,n_0$; by its
part (a), $\Delta_{\ge M}\le1.6\eps(1+1/\cOV)S/\log M\le5.46\,\eps S/\log M$,
since $1.6\,(1+1/\cOV)=5.455\ldots$.

(c), (d) Every split is produced by the $\tau$-rules at threshold $\tau$ from a
witness with parameter $s<\tau$, so Lemma~\ref{s2:lemThinCut} applies.
\end{proof}

\subsection{The hierarchy}

\begin{definition}[the hierarchy $\HBtp$]\label{s2:defHBtp}
Let $G$ be a graph on $n$ vertices, and let $\eps=2^{-5}$, $\sigma=100$,
$\Cp=103$, $\Aexp=105$ and $\Dstar$ be as in Definition~\ref{s1:defConstants}.
Put $G_1\defeq G$. All graphs $G_l$, $G'_l$ below have vertex set $V(G)$. For
$l=1,2,\dots$:
\begin{enumerate}
\item[(R0)]\namedlabel{s2:R0}{(R0)} $d_l\defeq2|E(G_l)|/n$. If $d_l<\Dstar$,
stop, and put $E_0\defeq E(G_l)$ and $R\defeq l-1$. Otherwise put
$\lambda_l\defeq\log d_l$.
\item[(R1)]\namedlabel{s2:R1}{(R1)} $\tHB_l\defeq\log^2d_l$. As long as the
current graph contains a cycle of length at least $\tHB_ld_l$, delete one (the
lexicographically first, for a fixed order of $V(G)$). The deleted cycles form
$\Cyc_l$; the remaining graph is $G'_l$. Thus $G'_l$ has no cycle of length at
least $\tHB_ld_l$.
\item[(R2)]\namedlabel{s2:R2}{(R2)}
$M_l\defeq\bigl\lceil\max\bigl(2^{40},\,2^{16}\,\tHB_ld_l\log^4(\tHB_ld_l)\bigr)\bigr\rceil\in\mathbb N$,\quad
$\Lambda_l\defeq\log M_l$,\quad $s_l\defeq\lceil\Lambda_l^{\sigma}\rceil$,\quad
$P_l\defeq\lceil\lambda_l^{\Cp}\rceil$,\quad
$\tau_l\defeq\lceil128\,s_l\log^2M_l\rceil$.
The ceiling makes $M_l$ a natural number, as its later uses require (it counts
classes and colours, for instance $\KJS_l=M_l^2$ in Definition~\ref{s3:defCOL}
and the palettes $[3M_l]$, $[4M_l]$ of Section~\ref{s7:sec}); $\Lambda_l$ is the
logarithm of this integer. Besides integrality, only two facts about the ceiling are used:
$M_l\ge\max\bigl(2^{40},\,2^{16}\,\tHB_ld_l\log^4(\tHB_ld_l)\bigr)$, and
$M_l\le\max\bigl(2^{40},\,2^{16}\,\tHB_ld_l\log^4(\tHB_ld_l)\bigr)+1$.
\item[(R3)]\namedlabel{s2:R3}{(R3)} Run an $s=0$ recursion
(Lemma~\ref{s2:lemOVgeneric}) on $G'_l$ that stops exactly at
$(\eps,0)$-expanders: a node is a leaf iff its graph is an
$(\eps,0)$-expander, and otherwise any witness with parameter $0$ is used. This
is the recursion of the proof of \cite[Lemma~14]{BM} with $s=0$ (Cited
result~\ref{s1:citLem14}). Its leaves are the \emph{$s=0$ pieces} $\piece$ of
round $l$. Inside every piece with $|\piece|\ge P_l$, run the $\tau$-run of
$\piece$ with $s=s_l$ and $\tau=\tau_l$ (Lemma~\ref{s2:lem14tau}), with any
witnesses. The \emph{two-level recursion} of round $l$ is the split recursion
obtained from the $s=0$ recursion by attaching, at every piece $\piece$ with
$|\piece|\ge P_l$, the tree of its $\tau$-run. Its leaves are the pieces of
size less than $P_l$ and the leaves of the $\tau$-runs. The \emph{round-$l$
pre-parts} are the leaves of the $\tau$-runs with at least $P_l$ vertices;
equivalently, the leaves of the two-level recursion with at least $P_l$
vertices. For a pre-part named $Z$, $Z^0$ denotes its vertex set and $X^0_Z$
its leaf graph, a graph on $Z^0$. Finally
$\Dups_l\defeq$ the set of vertices lying in at least two leaves (of any size,
singletons included) of the two-level recursion of round $l$.
\item[(R4)]\namedlabel{s2:R4}{(R4)} $D_l\defeq$ the set of vertices lying in at
least two round-$l$ pre-parts. Fix an order of the round-$l$ pre-parts (the
\emph{home order}). For a vertex $v$ lying in some pre-part,
$\home(v)=\home_l(v)\defeq$ the first pre-part containing $v$. The
\emph{guests} of a pre-part $Z$ are $S_Z\defeq\{v\in Z^0\cap D_l:\home(v)\ne
Z\}$. The pre-part $Z$ is \emph{light} iff
\begin{enumerate}
\item[(L1)] $|S_Z|\le|Z^0|/2$;
\item[(L2)] every $u\in Z^0$ has at most $s_l/2$ neighbours in $S_Z$ in the
graph $G'_l[Z^0]$;
\item[(GC)]\namedlabel{s2:ruleGC}{(GC)} no $x\in S_Z$ has at least
\[
\thGC_l(Z^0)\defeq\bigl\lceil|Z^0|\,\lambda_l^{-1/2}\bigr\rceil
\]
neighbours in $Z^0\setminus S_Z$ in the graph $X^0_Z$ (the \emph{guest cap}).
\end{enumerate}
Otherwise $Z$ is \emph{standalone}. A pre-part satisfying (L1) and (L2) but
failing (GC) is a \emph{GC-part}; it is standalone. A light pre-part $Z$ gives
the \emph{light part} with vertex set $Z^0\setminus S_Z$ and graph
$X_Z\defeq X^0_Z-S_Z$. $\Std_l$ is the set of round-$l$ standalone pre-parts;
it contains the GC-parts and the pre-parts failing (L1) or (L2), and it is a
deterministic function of the run.
\item[(R5)]\namedlabel{s2:R5}{(R5)} The edges of $G'_l$ are assigned in this
order.
(1) For every light pre-part $Z$ the edges of $X_Z$ are assigned to the light
part $Z$; for every standalone pre-part $Z$ the edges of $X^0_Z$ are assigned
to $Z$. (The graphs $X^0_Z$ of distinct pre-parts are edge-disjoint, being
distinct leaves of the two-level recursion, Lemma~\ref{s2:lemSEP}(i), and
$X_Z\subseteq X^0_Z$; so no edge is assigned twice in this step.)
(2) Every edge not yet assigned whose ends both lie in some light part is
assigned to the first such light part (in the home order).
(3) Every edge not yet assigned whose ends both lie in $Z^0$ for some
standalone pre-part $Z$ is assigned to the first such pre-part (in the home
order).
(4) Every other edge of $G'_l$ passes down: $E(G_{l+1})$ is the set of these
edges.
$E_l(Z)$ denotes the set of edges assigned at round $l$ to $Z$ (a light part or
a standalone pre-part).
\end{enumerate}
\emph{Naming.} Pre-parts are leaves, so distinct pre-parts are counted
separately even if their vertex sets coincide. If a pre-part $Z$ is light, the
letter $Z$ also denotes the vertex set $Z^0\setminus S_Z$ of its light part.
A \emph{round-$l$ part} is a light part or a standalone pre-part of round $l$.
We write $\home_l$, $S_Z$ and so on with the round when it is not clear from
the context.

A \emph{valid $\HBtp$ run} is an execution of this procedure with any
witnesses in both levels of (R3), any vertex order in (R1), and any home
orders in (R4). At $s=0$ the $\tau$-rules are idle
(Lemma~\ref{s2:lemOVgeneric}), so the first level of (R3) is the recursion
of \cite{BM} unchanged; the $\tau$-rules act only inside the pieces.
\deps{\ref{s1:citLem14}, \ref{s2:lem14tau}, \ref{s1:defConstants}, \ref{s2:lemOVgeneric}, \ref{s2:lemSEP}}
\end{definition}

\begin{definition}[ancestors, typing, multiplicities]\label{s2:defAncestors}
Fix a valid $\HBtp$ run. An \emph{ancestor of round $r$} is either
\begin{itemize}
\item a light part $Y$ of round $r$; then $V(Y)\defeq Y$ (the vertex set
$Y^0\setminus S_Y$), $H_Y\defeq X_Y$, $\eps_Y\defeq2^{-6}$ and
$s_Y\defeq s_r/2$; or
\item a standalone pre-part $Y\in\Std_r$, GC-parts included; then
$V(Y)\defeq Y^0$, $H_Y\defeq X^0_Y$, $\eps_Y\defeq2^{-5}$ and $s_Y\defeq s_r$.
\end{itemize}
We write $r(Y)$ for the round of $Y$ and $L_Y\defeq\log|V(Y)|$. Every ancestor
corresponds to a distinct pre-part $Y^0\supseteq V(Y)$ of its round. For
$l\ge3$ and a vertex $x$,
\[
\anc_l(x)\defeq\{Y:\ Y\text{ an ancestor of a round }r\le l-2\text{ with }x\in
V(Y)\},
\]
and $\anc_l(x)\defeq\emptyset$ for $l\le2$. For $Z\in\Std_l$:
\begin{itemize}
\item the \emph{hubs} of $Z$ are $A_Z\defeq Z^0\cap D_l$;
\item the \emph{ports} of $Z$ are $U_Z\defeq Z^0\setminus D_l$;
\item the \emph{fresh ports} are $F_Z\defeq\{x\in U_Z:\anc_l(x)=\emptyset\}$
(so all ports are fresh for $l\le2$);
\item the \emph{classed ports} are $Q_Z\defeq U_Z\setminus F_Z$.
\end{itemize}
For a round $r$ and a vertex $w$: $\mu_r(w)$ is the number of round-$r$
pre-parts containing $w$ (so $D_r=\{w:\mu_r(w)\ge2\}$);
$\dup_r\defeq\sum_w(\mu_r(w)-1)^+$; and $\mult_r(w)$ is the number of ancestors
$Y$ of round $r$ with $w\in V(Y)$. Finally $j_0(x)$ is the first round in which
$x$ lies in a pre-part, and $j_1(x)$ the first round in which $x$ lies in a
light part ($\infty$ if there is none).
\deps{\ref{s2:defHBtp}}
\end{definition}

\subsection{Basic properties of valid runs}

\begin{lemma}[Lemma-25 size cap]\label{s2:lemCap}
\begin{enumerate}
\item[(i)] If $m\ge2^{40}$, $T\ge2^{117}$ and $m<18432\,T\log^4m$, then
$m\le2^{16}T\log^4T$.
\item[(ii)] In every round $l\le R$ of a valid $\HBtp$ run, every $s=0$ piece
$\piece$ satisfies $|\piece|\le M_l$. Hence every round-$l$ pre-part has
$|Z^0|\le M_l$; for every round-$l$ part $Z$ (light or standalone) every vertex
is incident with at most $M_l-1$ edges of $E_l(Z)$; and
$\tau_l\ge128\,s_l\log^2|\piece|$, so Lemma~\ref{s2:lem14tau} applies to the
$\tau$-run of every piece with $|\piece|\ge P_l$.
\end{enumerate}
\emph{Remark (not used).} Part (i) is false for small $T$: for $T=2^{20}$ and
$m=2^{55}$ we have $m\ge2^{40}$ and $18432\,T\log^4m=2^{57.29\ldots}>m$, but
$2^{16}T\log^4T=2^{53.28\ldots}<m$. The uniform form $m\le\max\bigl(2^{40},
2^{16}T(\log T+40)^4\bigr)$ holds for all $T\ge1$.
\deps{\ref{s1:citLem25}, \ref{s1:condG2}, \ref{s2:defHBtp}, \ref{s2:lem14tau}}
\end{lemma}

\begin{proof}
(i) Put $\chi(y)\defeq y/\log^4y$; since $\frac{d}{dy}\ln \chi(y)=\frac1y\bigl(1-
\frac4{\ln y}\bigr)$, $\chi$ is increasing for $y>e^4$. The hypothesis says
$\chi(m)<18432\,T$. Put $x\defeq\log T\ge117$ and $B\defeq2^{16}T\log^4T=2^{16}Tx^4$,
so $\log B=16+x+4\log x$. We have $(2^{16}/18432)^{1/4}=(32/9)^{1/4}=1.37317\ldots$,
and $0.37317\,x\ge16+4\log x$ for $x\ge117$: at $x=117$ the two sides are
$43.66\ldots$ and $43.48\ldots$, and for $x\ge117$ the derivative of the left
side ($0.373\ldots$) exceeds that of the right side ($4/(x\ln2)<0.05$). Hence
$(32/9)^{1/4}x\ge16+x+4\log x=\log B$, i.e.\ $2^{16}x^4\ge18432\,(\log B)^4$,
i.e.\ $\chi(B)\ge18432\,T$. If $m>B$, then, as $B>e^4$, $\chi(m)>\chi(B)\ge18432\,T$,
contradicting the hypothesis. So $m\le B$.

\emph{Remark.} The numbers in the counterexample are direct computations. For
the uniform form put $B'\defeq2^{16}T(x+40)^4$ with $x=\log T\ge0$; then $\log
B'=16+x+4\log(x+40)$ and $\chi(B')\ge18432\,T$ is equivalent to
$(32/9)^{1/4}(x+40)\ge16+x+4\log(x+40)$, which follows from
$0.37317\,x+38.9\ge4\log(x+40)$; this holds at $x=0$ ($38.9\ge21.3$) and the
left side grows faster ($0.373$ against at most $4/(40\ln2)<0.15$). As
$B'>e^4$, $m>\max(2^{40},B')$ would give $\chi(m)>\chi(B')\ge18432\,T$.

(ii) Let $l\le R$ and $T\defeq\tHB_ld_l=d_l\log^2d_l$. By \Gc{2}(a)
(Condition~\ref{s1:condG2}), $d_l\ge\Dstar\ge2^{117}$, so $\log^2d_l\ge1$ and
$T\ge2^{117}$; and $M_l\ge\max(2^{40},2^{16}T\log^4T)$ by (R2). Let $\piece$ be an
$s=0$ piece and $m\defeq|\piece|$. If $m<2^{40}$ then $m\le M_l$. Otherwise the
graph of $\piece$ is an $(\eps,0)$-expander (the stopping rule in (R3)) with
$\eps=2^{-5}$ on $m\ge2^{40}=2^{30}/\eps^2$ vertices, so by Cited
result~\ref{s1:citLem25} it contains a cycle of length at least
$\eps^2m/(18\log^4m)=m/(18432\log^4m)$. This cycle lies in $G'_l$, which has no
cycle of length at least $T$ by (R1). Hence $m<18432\,T\log^4m$, and (i) gives
$m\le2^{16}T\log^4T\le M_l$. Every pre-part is a leaf of the $\tau$-run of some
piece, so $Z^0\subseteq V(\piece)$ and $|Z^0|\le M_l$. By (R5) every edge of
$E_l(Z)$ has both ends in $Z^0$ (in the light part $Z\subseteq Z^0$ if $Z$ is
light), so every vertex meets at most $|Z^0|-1\le M_l-1$ of them. Finally
$\tau_l\ge128\,s_l\log^2M_l\ge128\,s_l\log^2|\piece|$.
\end{proof}

\begin{lemma}[Lemma HS]\label{s2:lemHS}
Let $\varepsilon'>0$ and $s\ge0$, let $X$ be an $(\varepsilon',s)$-expander on
a vertex set of size $z\ge11$, let $W\subseteq V(X)$ with $|W|\le z/2$, and
suppose every $u\in V(X)\setminus W$ has at most $s_W\le s$ neighbours in $W$ in
$X$. Then $X-W$ is an $\bigl(\varepsilon'(\log z'/\log z)^2,\,s-s_W\bigr)$-expander
on $z'\defeq z-|W|$ vertices, and $\varepsilon'(\log z'/\log z)^2\ge
\varepsilon'/2$. In particular, if $X$ is a $(2^{-5},s)$-expander and
$s_W=s/2$, then $X-W$ is a spanning $(2^{-6},s/2)$-expander of $V(X)\setminus
W$.
\deps{\ref{s1:citDef11}}
\end{lemma}

\begin{proof}
Note $z'\ge z/2\ge5.5$. Let $U\subseteq V(X)\setminus W$ with $1\le|U|\le2z'/3$
and $F'\subseteq E(X-W)$ with $|F'|\le(s-s_W)|U|$. Put $F\defeq F'\cup
E_X(U,W)$. Then $|E_X(U,W)|\le s_W|U|$, so $|F|\le s|U|$, and $|U|\le2z'/3\le
2z/3$. By Cited result~\ref{s1:citDef11} applied to $X$,
$|\Nbr_{X-F}(U)|\ge\varepsilon'|U|/\log^2z$. No vertex of $W$ lies in
$\Nbr_{X-F}(U)$, because all edges between $U$ and $W$ are in $F$. For
$v\notin U\cup W$ and $u\in U$, the edge $uv$ (if present) lies in $X-W$ and not
in $E_X(U,W)$, so $uv\notin F$ iff $uv\notin F'$. Hence
$\Nbr_{X-F}(U)=\Nbr_{(X-W)-F'}(U)$, and
\[
|\Nbr_{(X-W)-F'}(U)|\ge\frac{\varepsilon'|U|}{\log^2z}
=\varepsilon'\Bigl(\frac{\log z'}{\log z}\Bigr)^2\frac{|U|}{\log^2z'}.
\]
This is the expansion claimed. Next, $(\log z'/\log z)^2\ge\frac12$ iff $z'\ge
z^{1/\sqrt2}$. Since $z'\ge z/2$, it suffices that $z^{1-1/\sqrt2}\ge2$, i.e.\
$\log z\ge1/(1-1/\sqrt2)=3.414\ldots$, i.e.\ $z\ge10.67$; this holds as $z\ge11$.
The last sentence follows because an $(\varepsilon_a,s)$-expander is an
$(\varepsilon_b,s)$-expander whenever $\varepsilon_b\le\varepsilon_a$.
\end{proof}

\begin{proposition}[structure of a valid run]\label{s2:propStructure}
For every valid $\HBtp$ run on a graph $G$ with $n\ge\Nzero$ and $d_1\ge\Dstar$:
\begin{enumerate}
\item[(i)] each $X^0_Z$ is a $(2^{-5},s_l)$-expander on $Z^0$; each light $X_Z$
is a $(2^{-6},s_l/2)$-expander on the light part $Z$; every $H_Y$ has minimum
degree greater than $s_Y$; and $s_r\ge\lambda_r^{100}$ for every $r\le R$;
\item[(ii)] the conclusions of Lemma~\ref{s2:lemCap}(ii) hold;
\item[(iii)] $E(G_{l+1})\subseteq E(G'_l)\subseteq E(G_l)$ for every $l\le R$,
so $d_{l+1}\le d_l$; and
\[
E(G)=\bigsqcup_{l\le R}E(\Cyc_l)\ \sqcup\ E_0\ \sqcup
\bigsqcup_{Y\text{ light}}E_{r(Y)}(Y)\ \sqcup\bigsqcup_{l\le R}\
\bigsqcup_{Z\in\Std_l}E_l(Z)
\]
(disjoint unions); $E(H_Y)\subseteq E_{r(Y)}(Y)$ for every ancestor $Y$; every
edge of $E_{r(Y)}(Y)$ has both ends in $V(Y)$; the graphs $H_Y$ of all
ancestors $Y$ of all rounds are pairwise edge-disjoint, and the graphs $X^0_Z$
of the pre-parts of one round are pairwise edge-disjoint; $\sum_{l\le
R}|\Cyc_l|\le n$ (the number of long cycles); and $|E_0|<\Dstar n/2$;
\item[(iv)] light parts of one round are pairwise vertex-disjoint, and
$|Y|\ge|Y^0|/2\ge P_r/2$ for every light part $Y$ of round $r$; every vertex
outside $D_l$ lies in at most one round-$l$ pre-part; the port sets $U_Z$ of
distinct $Z\in\Std_l$ are pairwise disjoint; and each vertex lies in at most one
light part per round, namely in the light part of $\home_l(v)$ if that
pre-part is light.
\end{enumerate}
\deps{\ref{s2:defHBtp}, \ref{s2:defAncestors}, \ref{s2:lem14tau}, \ref{s2:lemCap}, \ref{s2:lemHS}, \ref{s1:citDef11}, \ref{s1:citLem14}, \ref{s2:lemSEP}, \ref{s1:condG1}}
\end{proposition}

\begin{proof}
By \Gc{1}(a) (Condition~\ref{s1:condG1}) at $\mu=\log\log\Dstar$, we have
$\log\log\Dstar\ge2^8$, so $\lambda_l\ge\log\Dstar\ge2^{256}$ for every $l\le R$;
we use only much weaker consequences, such as $P_l\ge11$ and $\Dstar\ge8$.

(iv) Let $v$ lie in the light part of a light pre-part $Z$ of round $l$, i.e.\
$v\in Z^0\setminus S_Z$. If $v\in D_l$ then, as $v\notin S_Z$,
$\home_l(v)=Z$. If $v\notin D_l$ then $Z$ is the only round-$l$ pre-part
containing $v$, so again $\home_l(v)=Z$. Hence $v$ lies only in the light part
of $\home_l(v)$ (if that pre-part is light), which proves the first and last
assertions. $|Y|=|Y^0|-|S_Y|\ge|Y^0|/2$ by (L1), and $|Y^0|\ge P_r$. By
definition of $D_l$, a vertex outside $D_l$ lies in at most one round-$l$
pre-part, so the sets $U_Z=Z^0\setminus D_l$ are pairwise disjoint.

(i) Each $X^0_Z$ is a leaf of a $\tau$-run, hence a $(2^{-5},s_l)$-expander by
the stopping rule (Lemma~\ref{s2:lem14tau}(a)). For a light pre-part $Z$ apply
Lemma~\ref{s2:lemHS} to $X=X^0_Z$, $z=|Z^0|\ge P_l\ge11$, $W=S_Z$ and
$s_W=s_l/2$: $|S_Z|\le|Z^0|/2$ by (L1), and $X^0_Z$ is a subgraph of
$G'_l[Z^0]$ (node graphs of a split recursion are subgraphs of the root), so by
(L2) every vertex has at most $s_l/2$ neighbours in $S_Z$ in $X^0_Z$. So
$X_Z=X^0_Z-S_Z$ is a $(2^{-6},s_l/2)$-expander on $Z^0\setminus S_Z$. For an
ancestor $Y$, $H_Y$ is an $(\eps_Y,s_Y)$-expander on $|V(Y)|\ge P_r/2\ge2$
vertices (by (iv)), so its minimum degree exceeds $s_Y$ by the remark in Cited
result~\ref{s1:citDef11}. Finally $M_r\ge2^{16}T\log^4T\ge T=d_r\log^2d_r\ge d_r$
(with $T=\tHB_rd_r$), so $\Lambda_r\ge\lambda_r$ and
$s_r\ge\Lambda_r^{100}\ge\lambda_r^{100}$.

(ii) is Lemma~\ref{s2:lemCap}(ii).

(iii) By (R1), $E(G_l)=E(\Cyc_l)\sqcup E(G'_l)$. By Lemma~\ref{s2:lemSEP}(i)
applied to the two-level recursion of round $l$, every edge of $G'_l$ is
deleted exactly once or lies in exactly one leaf. By (R5), every edge of $G'_l$
is assigned to exactly one light part or standalone pre-part, or passes to
$G_{l+1}$: step (1) assigns each edge at most once (the $X^0_Z$ are
edge-disjoint and $X_Z\subseteq X^0_Z$), and steps (2)--(4) treat each
remaining edge once. So $E(G'_l)=\bigsqcup_{Z\text{ light}}E_l(Z)\sqcup
\bigsqcup_{Z\in\Std_l}E_l(Z)\sqcup E(G_{l+1})$; in particular
$E(G_{l+1})\subseteq E(G'_l)\subseteq E(G_l)$. By (R0), $E(G_{R+1})=E_0$, and
iterating from $l=1$ to $l=R$ gives the displayed partition. By (R5)(1),
$E(X_Y)$ is assigned to a light $Y$ and $E(X^0_Y)$ to a standalone $Y$, so
$E(H_Y)\subseteq E_{r(Y)}(Y)$; and each step of (R5) assigns an edge to $Y$
only if both its ends lie in $V(Y)$. Two pre-parts of the same round are
distinct leaves of the two-level recursion, so their leaf graphs are
edge-disjoint by Lemma~\ref{s2:lemSEP}(i). For distinct ancestors $Y\ne Y'$ (of
any rounds), $E(H_Y)\subseteq E_{r(Y)}(Y)$ and $E(H_{Y'})\subseteq
E_{r(Y')}(Y')$, and these assigned sets are disjoint by the partition above.
(The graphs $X^0_Y$ of light pre-parts of different rounds need not be
edge-disjoint: an edge of $X^0_Y$ between a guest $x\in S_Y$ and the light part
$Y$ may be assigned by no step of (R5) at round $r(Y)$, pass down, and lie in a
leaf graph of a later round. Only the statement above is claimed.)

\emph{Long cycles.} For $l\le R$ put $y_l\defeq d_l\ge\Dstar$; by the first part
of (iii), $y_{l+1}\le y_l$. The cycles of $\Cyc_l$ are edge-disjoint, each has
length at least $\tHB_ld_l=y_l\log^2y_l$, and together they have at most
$|E(G_l)|-|E(G'_l)|\le|E(G_l)|-|E(G_{l+1})|=\frac n2(y_l-y_{l+1})$ edges. Hence
\[
|\Cyc_l|\le\frac n2\cdot\frac{y_l-y_{l+1}}{y_l\log^2y_l}.
\]
Group the rounds $l\le R$ by $k\defeq\lfloor\log y_l\rfloor\ge k_0\defeq
\lfloor\log\Dstar\rfloor$. The rounds with a given $k$ are consecutive, and for
them $y_l\ge2^k$ and $\log^2y_l\ge k^2$, while the sum of $y_l-y_{l+1}$ over
them is at most the first $y_l$ of the group, which is less than $2^{k+1}$. So
these rounds contribute at most $\frac n2\cdot\frac{2^{k+1}}{2^kk^2}=n/k^2$,
and $\sum_{l\le R}|\Cyc_l|\le n\sum_{k\ge k_0}k^{-2}\le n/(k_0-1)\le n$, as
$k_0\ge3$.

Finally $E_0=E(G_{R+1})$ and $d_{R+1}<\Dstar$ by (R0), so
$|E_0|=nd_{R+1}/2<\Dstar n/2$.
\end{proof}

\begin{lemma}[edge laminarity, EL]\label{s2:lemEL}
For every ancestor $Y$ of round $r$ and every round $l>r$, no edge of $G_l$ has
both ends in $V(Y)$.
\deps{\ref{s2:defHBtp}, \ref{s2:defAncestors}}
\end{lemma}

\begin{proof}
By (R1) and (R5)(4), $E(G_{j+1})\subseteq E(G'_j)\subseteq E(G_j)$ for every
round $j$, so $E(G_l)\subseteq E(G_{r+1})\subseteq E(G'_r)$. Let $e\in
E(G_{r+1})$; then $e$ is an edge of $G'_r$ that was not assigned at round $r$.
If $Y$ is light and both ends of $e$ lie in $V(Y)=Y$, then $e$, if not assigned
in step (1) of (R5), is assigned in step (2) to the first light part containing
both ends; a contradiction. If $Y$ is standalone (a GC-part or not) and both
ends lie in $V(Y)=Y^0$, then $e$, if not assigned in steps (1) or (2), is
assigned in step (3) to the first standalone pre-part whose vertex set contains
both ends; again a contradiction.
\end{proof}

\begin{proposition}[overlap constants on $\HBtp$]\label{s2:propOV}
Fix a valid $\HBtp$ run and a round $r\le R$. Let $S_r$ be the total size of the
leaves of the two-level recursion of round $r$, and $S^{\piece}_r$ the total
size of the $s=0$ pieces. Then $S^{\piece}_r\le1.12\,n$ and $S_r\le1.21\,n$,
and:
\begin{enumerate}
\item[(K1)]\namedlabel{s2:K1}{(K1)} $\sum|Z^0|\le1.37\,n$, the sum over the
round-$r$ pre-parts; $|\Std_r|\le1.37\,n/P_r$; the number of ancestors of round
$r$ is at most $1.37\,n/P_r$ (each ancestor corresponds to a distinct pre-part
$Y^0\supseteq V(Y)$); hence the number $\nu_l$ of ancestors of rounds at most
$l-2$ satisfies $\nu_l\le1.37\,n\sum_{r\le l-2}P_r^{-1}$;
\item[(K2)]\namedlabel{s2:K2}{(K2)} $|D_r|\le\dup_r\le7.6\,\eps n/\log P_r$,
and $\sum_{Z\in\Std_r}|A_Z|\le\sum|Z^0\cap D_r|\le2\dup_r\le15.2\,\eps n/\log
P_r$, the middle sum over all round-$r$ pre-parts;
\item[(K3)]\namedlabel{s2:K3}{(K3)} the total size $\sum|Z^0|$ of the round-$r$
pre-parts failing (L1) is at most $30.4\,\eps n/\log P_r$, and
$\sum_Y|Y^0\cap\Dups_r|\le16\,\eps n/\log P_r$, the sum over all round-$r$
pre-parts $Y$;
\item[(F11)] $\sum_w\mult_r(w)=\sum_Y|V(Y)|\le1.37\,n$, the second sum over the
ancestors $Y$ of round $r$, and $\mult_r(w)\le\mu_r(w)$ for every vertex $w$.
\end{enumerate}
\deps{\ref{s2:lemOVgeneric}, \ref{s2:lem14tau}, \ref{s2:propStructure}, \ref{s2:defAncestors}, \ref{s2:defHBtp}, \ref{s2:lemSEP}}
\end{proposition}

\begin{proof}
\emph{The composite tree.} The two-level recursion of round $r$ is a split
recursion on $G'_r$, whose vertex set $V(G)$ has $n$ elements. Its splits in the
first level are splits of an $s=0$ recursion, which satisfy the hypotheses of
Lemma~\ref{s2:lemOVgeneric} with $c=1$ (the instance $c=1$ there), hence with
$c=1.6$. Its splits in the second level are splits of the $\tau$-run of a piece
$\piece$ with $|\piece|\ge P_r$; since $\tau_r\ge128\,s_r\log^2|\piece|$
(Proposition~\ref{s2:propStructure}(ii)), Lemma~\ref{s2:lem14tau}(b) shows that
they satisfy the hypotheses with $c=1.6$. Every non-leaf node has at least two
vertices. So Lemma~\ref{s2:lemOVgeneric} applies to the whole two-level
recursion with $c=1.6$ and $n_0=n$: by its part (c),
$S_r\le n/(1-5.472\eps)\le1.21\,n$; by the first bound of its part (a) and
$\log M\ge1$, for $M\ge2$, $\Delta_{\ge M}\le1.6\eps(1+1/\cOV)S_r/\log M$, and
$1.6(1+1/\cOV)=5.455\ldots\le5.46$ (the rounded constant $3.42$ would only give
$5.472$):
\begin{equation}\label{s2:eqDupComposite}
\Delta_{\ge M}\le\frac{5.46\,\eps S_r}{\log M}\le\frac{6.61\,\eps n}{\log M}.
\end{equation}
Applied to the first level alone (an $s=0$ recursion on $n$ vertices, $c=1$),
the lemma gives $S^{\piece}_r\le1.12\,n$.

(K1) Pre-parts are leaves, so $\sum|Z^0|\le S_r\le1.21\,n\le1.37\,n$. Every
pre-part has at least $P_r$ vertices, so there are at most $1.37\,n/P_r$
pre-parts. $\Std_r$ is a set of pre-parts, and the map from ancestors of round
$r$ to pre-parts (a light part to its pre-part, a standalone pre-part to
itself) is injective. Summing over $r\le l-2$ gives the bound on $\nu_l$.

(K2) The round-$r$ pre-parts are exactly the leaves of the two-level recursion
with at least $P_r$ vertices (by (R3) its leaves are the pieces with fewer than
$P_r$ vertices and the leaves of the $\tau$-runs; a piece with at least $P_r$
vertices whose $\tau$-run is a single node is such a leaf, hence a pre-part). By
Lemma~\ref{s2:lemOVgeneric}(b) with $M=P_r\ge2$, $\mu_r(w)\le1+\dup_{\ge
P_r}(w)$ for every $w$. Hence $\dup_r=\sum_w(\mu_r(w)-1)^+\le\Delta_{\ge P_r}\le
6.61\,\eps n/\log P_r\le7.6\,\eps n/\log P_r$ by \eqref{s2:eqDupComposite}.
Every $w\in D_r$ contributes at least $1$ to $\dup_r$, so $|D_r|\le\dup_r$.
Next, $\sum_{Z\in\Std_r}|A_Z|=\sum_{Z\in\Std_r}|Z^0\cap D_r|\le\sum_Z|Z^0\cap
D_r|=\sum_{w\in D_r}\mu_r(w)\le\sum_{w\in D_r}2(\mu_r(w)-1)=2\dup_r$, using
$\mu\le2(\mu-1)$ for $\mu\ge2$.

(K3) If $Z$ fails (L1) then $|Z^0|<2|S_Z|\le2|Z^0\cap D_r|$; summing and using
(K2) gives at most $4\dup_r\le30.4\,\eps n/\log P_r$. For the second bound, let
$u\in\Dups_r$ lie in a pre-part $Y$. Then $u$ lies in another leaf $\Leaf\ne Y$
of the two-level recursion. Let $\nu$ be the lowest common ancestor of the
leaves $Y$ and $\Leaf$; they lie below different children of $\nu$, and $u$
lies in both (vertex sets shrink downwards), so $u\in N''_\nu$ by
Lemma~\ref{s2:lemSEP}(0); and $|\nu|\ge|Y^0|\ge P_r$. Hence
$\dup_{\ge P_r}(u)\ge1$ and $\mu_r(u)\le1+\dup_{\ge P_r}(u)\le2\dup_{\ge
P_r}(u)$. Summing over $u\in\Dups_r$,
\[
\sum_Y|Y^0\cap\Dups_r|=\sum_{u\in\Dups_r}\mu_r(u)\le2\Delta_{\ge P_r}
\le\frac{13.3\,\eps n}{\log P_r}\le\frac{16\,\eps n}{\log P_r}.
\]

(F11) Counting pairs $(w,Y)$ with $w\in V(Y)$ gives
$\sum_w\mult_r(w)=\sum_Y|V(Y)|$. The injective map $Y\mapsto Y^0$ of (K1)
satisfies $V(Y)\subseteq Y^0$, so $\sum_Y|V(Y)|\le\sum_{Z}|Z^0|\le1.37\,n$ and,
restricted to the ancestors containing $w$, it shows $\mult_r(w)\le\mu_r(w)$.
\end{proof}

\subsection{Degree recursion, existence, and tower facts}

\begin{proposition}[degree recursion]\label{s2:propDegRec}
Let $l$ be a round of a valid $\HBtp$ run with $d_l\ge\Dstar$ (that is,
$l\le R$), and put $x\defeq\lambda_l$. Then $M_l\le d_l^2$, $\Lambda_l\le2x$,
$s_l\le P_l$, and
\[
d_{l+1}\le1.21\,P_l+9\,s_l\log M_l\le6P_l+9\,s_l\log M_l+1.2\,s_l\le
7\,x^{103}\le x^{\Aexp}<d_l
\]
($\Aexp=104$ would suffice; $\Aexp=105$ is used). More precisely, every edge of
$G_{l+1}$ is of one of the following kinds:
\begin{enumerate}
\item[(a)] an edge deleted by the $\tau$-run of a piece $\piece$ with
$|\piece|\ge P_l$; there are at most $4s_l|\piece|\log|\piece|$ of them for
each piece and at most $4.5\,n\,s_l\log M_l$ in total;
\item[(b)] an edge of a leaf of the two-level recursion with fewer than $P_l$
vertices (a small leaf of a $\tau$-run, or a piece with $|\piece|<P_l$); such a
leaf, with $q<P_l$ vertices, has fewer than $qP_l/2$ edges;
\item[(c)] a \emph{guest edge}: an edge of $X^0_Z$ with an end in $S_Z$, for a
light pre-part $Z$; there are at most $|Z^0|s_l/2$ of them for each light $Z$.
\end{enumerate}
Standalone pre-parts, GC-parts included, pass none of the edges of their leaf
graphs down; in particular, declaring a pre-part a GC-part (instead of light)
can only decrease the set of passed-down edges.
\deps{\ref{s2:lem14tau}, \ref{s2:lemCap}, \ref{s2:propOV}, \ref{s2:defHBtp}, \ref{s2:lemSEP}, \ref{s1:condG1}, \ref{s2:lemOVgeneric}}
\end{proposition}

\begin{proof}
Put $\mu\defeq\log x$, where $x=\log d_l\ge\log\Dstar$; so $\mu\ge\log\log\Dstar$,
and \Gc{1}(a),(b) (Condition~\ref{s1:condG1}) hold at $\mu$: $\mu\ge2^8$ and
$x=2^\mu\ge2^{14}\Aexp\mu^3$. Hence $x\ge2^{256}$, so $x\ge7$, $x^2\ge2^{107}$ and
$x^3\ge2^{101}$; and $x\ge2^{14}\cdot105\log^3x$, so $x\ge21+6\log x$ and
$x>105\log x$.

\emph{Parameters.} Let $T\defeq\tHB_ld_l=2^xx^2$. Then $\log T=x+2\log x\le2x$,
so $2^{16}T\log^4T\le2^{16}\cdot2^xx^2\cdot16x^4=2^{20}x^62^x\le2^{2x-1}$ because
$x\ge21+6\log x$; also $2^{40}\le2^{2x-1}$. By (R2), $M_l$ is at most this
maximum plus $1$, so $M_l\le2^{2x-1}+1\le2^{2x}=d_l^2$ and
$\Lambda_l=\log M_l\le2x$. Then $s_l\le\Lambda_l^{100}+1\le2^{100}x^{100}+1\le
2^{101}x^{100}\le x^{103}\le P_l$.

\emph{Classification.} Let $e\in E(G_{l+1})$; by (R5), $e$ is an edge of $G'_l$
not assigned at round $l$. By Lemma~\ref{s2:lemSEP}(i) applied to the two-level
recursion, $e$ is deleted at some split, or lies in exactly one leaf $\Leaf$.
Splits of the first level delete nothing (Lemma~\ref{s2:lemOVgeneric}, instance
$c=1$), so a deleted $e$ is of kind (a). If $|\Leaf|<P_l$, $e$ is of kind (b).
Otherwise $\Leaf$ is a pre-part $Z$ and $e\in E(X^0_Z)$. If $Z$ is standalone,
$e$ is assigned to $Z$ by (R5)(1), which is impossible. If $Z$ is light and $e$
has no end in $S_Z$, then $e\in E(X_Z)$ is assigned by (R5)(1), again
impossible. So $e$ is of kind (c). The same argument shows the assertion about
standalone pre-parts and GC-parts.

\emph{Counting.} (a) By Lemma~\ref{s2:lemCap}(ii), Lemma~\ref{s2:lem14tau}(a)
applies to the $\tau$-run of $\piece$ and gives at most
$4s_l|\piece|\log|\piece|\le4s_l|\piece|\log M_l$ deleted edges. By
Proposition~\ref{s2:propOV}, the pieces have total size
$S^{\piece}_l\le1.12\,n$, so kind (a) has at most $4.48\,n\,s_l\log M_l$ edges.
(b) Every leaf of the two-level recursion is non-empty (both children of every
split, of the $s=0$ recursion and of a $\tau$-run, are non-empty). So a leaf
with $1\le q<P_l$ vertices has at most $q(q-1)/2<qP_l/2$ edges.
(c) Every guest edge of $X^0_Z$ joins some $u\in Z^0$ to a vertex of $S_Z$, so
it is counted in $\sum_{u\in Z^0}|N_{X^0_Z}(u)\cap S_Z|$; each term is at most
$s_l/2$ by (L2), since $X^0_Z\subseteq G'_l[Z^0]$. The leaves of kind (b) and
the light pre-parts are distinct leaves of the two-level recursion, so the sum
of $q$ over the leaves of kind (b) plus the sum of $|Z^0|$ over the light
pre-parts is at most $S_l\le1.21\,n$ (Proposition~\ref{s2:propOV}). As
$s_l\le P_l$, kinds (b) and (c) together have at most
$\frac12\cdot1.21\,n\cdot P_l$ edges. Therefore
$|E(G_{l+1})|\le4.5\,n\,s_l\log M_l+0.605\,n\,P_l$, that is,
$d_{l+1}\le9\,s_l\log M_l+1.21\,P_l$, which is at most
$6P_l+9s_l\log M_l+1.2s_l$.

\emph{Final bounds.} Using $P_l\le x^{103}+1$, $s_l\le2^{101}x^{100}$ and
$\log M_l\le2x$:
\[
6P_l+9s_l\log M_l+1.2s_l\le6x^{103}+6+9\cdot2^{102}x^{101}+1.2\cdot2^{101}x^{100}
\le6x^{103}+2^{107}x^{101}\le7x^{103},
\]
since $x^2\ge2^{107}$. Next $7x^{103}\le x^{104}\le x^{105}$ as $x\ge7$, and
$x^{105}<2^x=d_l$ as $105\log x<x$.
\end{proof}

\begin{proposition}[existence and termination]\label{s2:propExists}
For every graph $G$ a valid $\HBtp$ run exists, and every sequence of admissible
choices terminates. More precisely: within a round, (R1) and the first level of
(R3) terminate; the $\tau$-runs terminate with any witnesses ($U'\ne\emptyset$
and $n_1,n_2<m$ at every split); (GC) is a deterministic relabelling that feeds
back into neither $D_l$ nor $\home$ nor the guest sets; and $d_{l+1}<d_l$ while
$d_l\ge\Dstar$, so $R$ is finite.
\deps{\ref{s2:lem14tau}, \ref{s2:propDegRec}, \ref{s2:defHBtp}, \ref{s2:lemCap}, \ref{s2:lemOVgeneric}, \ref{s2:defWitness}, \ref{s2:propOV}}
\end{proposition}

\begin{proof}
The proofs of Lemma~\ref{s2:lemCap}(ii), Proposition~\ref{s2:propOV} and
Proposition~\ref{s2:propDegRec} use only the data of the round $l$ in question
and $d_l\ge\Dstar$; so they apply to every round $l$ with $d_l\ge\Dstar$ of any
(possibly unfinished) execution of the procedure.

Fix a round $l$ with $d_l\ge\Dstar$ (if $d_1<\Dstar$ the run stops at once
with $R=0$). (R1) deletes at least one edge per step, so it terminates. In the
first level of (R3), a node that is not an $(\eps,0)$-expander has a witness
(Definition~\ref{s2:defWitness}), and its children have $|U|+|N|<\frac34m$ and
$m-|U|<m$ vertices (Lemma~\ref{s2:lemOVgeneric}, instance $c=1$); so every
root-to-node path has at most $n$ nodes and this level terminates. By
Lemma~\ref{s2:lemCap}(ii), $\tau_l\ge128\,s_l\log^2|\piece|$ for every piece,
so Lemma~\ref{s2:lem14tau}(a) applies to every $\tau$-run with every choice of
witnesses: $U'\ne\emptyset$, $n_1,n_2<m$, and the $\tau$-run terminates.
In (R4), $D_l$, $\home$ and the sets $S_Z$ are determined by the pre-parts and
the home order alone; conditions (L1), (L2) and (GC) are evaluated afterwards,
and their outcome is not used in defining $D_l$, $\home$ or $S_Z$. (R5) is
deterministic. All required choices (witnesses, vertex orders, home orders)
exist. So each round is a finite procedure. By
Proposition~\ref{s2:propDegRec}, $d_{l+1}<d_l$ whenever $d_l\ge\Dstar$; as
$|E(G_l)|$ is a non-negative integer, (R0) stops after finitely many rounds.
\end{proof}

\begin{lemma}[tower-lacunary sums]\label{s2:lemLacunary}
Let $x_0\defeq\log\Dstar$.
\begin{enumerate}
\item[(i)] If $X_1,\dots,X_R\ge0$ and $X_r\le X_{r+1}/2$ for all $r<R$, then
$\sum_{r\le R}X_r\le2X_R$ and $\sum_{r\le R}(R-r+1)X_r\le4X_R$.
\item[(ii)] For a valid run with $R\ge1$, the degree recursion gives
$\lambda_r\ge2^{\lambda_{r+1}/\Aexp}$ for all $r<R$. Hence, if
$F:[x_0,\infty)\to[0,\infty)$ is non-increasing and $F(2^{x/\Aexp})\le F(x)/2$
for all $x\ge x_0$, then $\sum_{r\le R}F(\lambda_r)\le2F(\lambda_R)\le
2F(x_0)$ and $\sum_{r\le R}(R-r+1)F(\lambda_r)\le4F(x_0)$; for the shifted
sums (when $R\ge3$), $\sum_{l=3}^RF(\lambda_{l-2})\le2F(\lambda_{R-2})\le2F(x_0)$
and $\sum_{l=3}^R(R-l+1)F(\lambda_{l-2})\le4F(\lambda_{R-2})\le4F(x_0)$.
\item[(iii)] More generally, if $F:[x_0,\infty)\to[0,\infty)$ satisfies
\[
\text{(H)}\qquad F(y)\le F(x)/2\quad\text{whenever } x\ge x_0\text{ and }
y\ge2^{x/\Aexp},
\]
then $\sum_{r\le R}F(\lambda_r)\le2F(\lambda_R)$ and $\sum_{r\le
R}(R-r+1)F(\lambda_r)\le4F(\lambda_R)$ (and likewise for $F(\lambda_{l-2})$).
\item[(iv)] Under \Gc{1}, condition (H) holds for $F(x)=x^{-a}$ (every $a\ge
1/100$), for $F(x)=1/\log x$ and for $F(x)=(95\log x+8)\,x^{-b}$ (every $b\ge
1$; not used in this paper), all three non-increasing on $[x_0,\infty)$;
and for the three
functions $F(x)=(2\logs(2^x)+2)/\eta(x)$ with $\eta(x)\in\{x,\ x^{1/2},\
\log x\}$ (note $2\logs(2^x)+2=2\logs x+4$ for $x>1$; not used in this
paper). Each of these three
satisfies $F(x)\le2F(x_0)$ for all $x\ge x_0$ (not used in this paper).
\end{enumerate}
Round-$l$ quantities such as $1/M_l$ or
$M_l^3(95\log\lambda_{l-2}+8)/\lambda_{l-2}^{95}$ are summed either by (i),
once the halving of consecutive terms is known, or by bounding them by
$F(\lambda_{l-2})$ for an $F$ as in (ii); for instance the bound
$M_l\le\lambda_{l-2}^{1.6}$, proved in part~(b) of the next lemma (this
example is not used in the proof of the present lemma), gives
$M_l^3(95\log\lambda_{l-2}+8)/\lambda_{l-2}^{95}\le(95\log\lambda_{l-2}+8)\,
\lambda_{l-2}^{-90.2}$. (This example is not used in this paper:
Lemma~\ref{s7:lemUHsplit}(iv) bounds its round-$l$ terms by a function of its
own and sums them by~(i).)
\deps{\ref{s2:propDegRec}, \ref{s1:condG1}}
\end{lemma}

\begin{proof}
(i) By induction $X_r\le2^{-(R-r)}X_R$, so $\sum_rX_r\le X_R\sum_{k\ge0}2^{-k}
=2X_R$ and $\sum_r(R-r+1)X_r\le X_R\sum_{k\ge0}(k+1)2^{-k}=4X_R$.

(iii) For $r<R$, Proposition~\ref{s2:propDegRec} at round $r$ gives
$d_{r+1}\le\lambda_r^{\Aexp}$, i.e.\ $\lambda_{r+1}\le\Aexp\log\lambda_r$,
i.e.\ $\lambda_r\ge2^{\lambda_{r+1}/\Aexp}$; and $\lambda_{r+1}\ge x_0$. By (H)
with $x=\lambda_{r+1}$ and $y=\lambda_r$, $F(\lambda_r)\le F(\lambda_{r+1})/2$,
and (i) applies to $X_r\defeq F(\lambda_r)$. For sums of $F(\lambda_{l-2})$
over $3\le l\le R$ apply the same to $r=l-2\le R-2$.

(ii) Let $x\ge x_0$. Then $\mu\defeq\log x\ge\log x_0=\log\log\Dstar$, so \Gc{1}(a),(b)
at $\mu$ give $\mu\ge2^8$ and $x=2^\mu\ge2^{14}\Aexp\mu^3\ge\Aexp\mu=\Aexp\log x$, that is,
$2^{x/\Aexp}\ge x$. Hence $2^{x/\Aexp}\ge x\ge x_0$ for $x\ge x_0$.
If $F$ is non-increasing and $F(2^{x/\Aexp})\le F(x)/2$, then for $y\ge
2^{x/\Aexp}$ we get $F(y)\le F(2^{x/\Aexp})\le F(x)/2$, which is (H). So (iii)
applies, and $F(\lambda_R)\le F(x_0)$ as $\lambda_R\ge x_0$.

(iv) Let $x\ge x_0$ and $y\ge2^{x/\Aexp}$, and put $\mu\defeq\log x\ge\log\log\Dstar$.
By \Gc{1}(a),(b) at $\mu$, $\mu\ge2^8$ and $x=2^\mu\ge2^{14}\Aexp\mu^3$; hence
$x\ge2^{256}$, $x/\Aexp\ge101\log x$, $x\ge2\Aexp\log x$, $x\ge\Aexp\log^2x/4$ and
$x/\Aexp\ge2^{12}$, and all of this holds in particular for $x=x_0$. Then
$\log y\ge x/\Aexp$, and $2^{x/\Aexp}\ge x^{101}$ because $x/\Aexp\ge101\log x$; so
$y\ge x^{101}\ge2^{32}$.
We use two facts about $\logs$: $\logs z\le z$ for $z\ge1$ (if $z\le2$ then
$\logs z\le1$; if $z>2$ then $\logs z=1+\logs(\log z)\le1+\log z\le z$ by
induction), and hence, for $y\ge4$, $\logs y=2+\logs(\log\log y)\le2+\log\log y
\le1+\log y$.

\emph{$F(x)=x^{-a}$, $a\ge1/100$:} $F(y)/F(x)=(x/y)^a\le x^{-100a}\le x^{-1}\le
1/2$.

\emph{$F(x)=1/\log x$:} $F(y)=1/\log y\le\Aexp/x\le1/(2\log x)=F(x)/2$, as
$x\ge2\Aexp\log x$.

\emph{$F(x)=(95\log x+8)x^{-b}$, $b\ge1$:} $x^{b+1}F'(x)=95/\ln2-b(95\log x+8)
<0$, so $F$ is decreasing; and $F(y)\le y^{-b/2}$, because $95\log y+8\le
y^{1/2}$ for $y\ge2^{30}$ (true at $2^{30}$, and the right side grows faster).
Hence $F(y)\le x^{-50b}\le x^{-b}/2\le F(x)/2$.

\emph{$F(x)=(2\logs x+4)/\eta(x)$.} Here $F(x)\ge4/\eta(x)$, so it suffices to
show $F(y)\le2/\eta(x)$. For $\eta(x)=x$: $F(y)\le(2\log y+6)/y\le y^{-1/2}$
(the inequality $2\log y+6\le y^{1/2}$ holds at $y=2^{12}$, and the right side
grows faster beyond), and $y^{1/2}\ge x^{50}\ge x/2$. For $\eta(x)=x^{1/2}$:
$F(y)\le(2\log y+6)/y^{1/2}\le y^{-1/4}$ (valid for $y\ge2^{32}$ in the same
way), and $y^{1/4}\ge x^{25}\ge x^{1/2}/2$. For $\eta(x)=\log x$: with
$u\defeq\log y\ge x/\Aexp$, $F(y)\le(2\log u+8)/u\le u^{-1/2}$ (valid for $u\ge
2^{12}$), and $u^{-1/2}\le(\Aexp/x)^{1/2}\le2/\log x$ as $x\ge\Aexp\log^2x/4$.

\emph{The bound $F(x)\le2F(x_0)$ for the three $\logs$-functions.} Let
$\mathrm{tw}(0)\defeq1$ and $\mathrm{tw}(j+1)\defeq2^{\mathrm{tw}(j)}$. Since
$\log^{[i]}$ is increasing and $\log^{[j]}\mathrm{tw}(j)=1$, a real $z>1$ has
$\logs z=j$ iff $\mathrm{tw}(j-1)<z\le\mathrm{tw}(j)$. Let $k_0\defeq\logs x_0\ge
1$, let $x\ge x_0$, and write $\logs x=k_0+k$ with $k\ge0$. Then $x_0\le
\mathrm{tw}(k_0)$ and $x>\mathrm{tw}(k_0+k-1)$. If $k=0$, then $F(x)\le F(x_0)$
since $\eta$ is increasing. If $k=1$, then $2\logs x+4=2k_0+6\le2(2k_0+4)$ and
$\eta(x)\ge\eta(x_0)$, so $F(x)\le2F(x_0)$. Let $k\ge2$. We claim
$\eta(x)\ge2^k\eta(x_0)$. Note $\mathrm{tw}(k_0+1)=2^{\mathrm{tw}(k_0)}\ge2^{x_0}$,
and $\mathrm{tw}(j+1)=2^{\mathrm{tw}(j)}\ge4\,\mathrm{tw}(j)$ whenever
$\mathrm{tw}(j)\ge4$. Hence $x>\mathrm{tw}(k_0+k-1)\ge4^{k-2}2^{x_0}$, which gives
the claim for $\eta(x)=x$ (as $4^{k-2}2^{x_0}\ge2^kx_0$ for $k\ge2$) and for
$\eta(x)=x^{1/2}$ (as $2^{k-2}2^{x_0/2}\ge2^kx_0^{1/2}$). For
$\eta=\log$: $\log x>\mathrm{tw}(k_0+k-2)$, which is at least $\mathrm{tw}(k_0)\ge
x_0\ge4\log x_0$ if $k=2$, and at least $4^{k-3}2^{x_0}\ge2^k\log x_0$ if
$k\ge3$. Using the claim and $2k_0+2k+4\le2^k(2k_0+4)$ (as $2^k\ge1+k$),
$F(x)\le(2k_0+2k+4)/(2^k\eta(x_0))\le(2k_0+4)/\eta(x_0)=F(x_0)$.
\end{proof}

\begin{lemma}[tower facts]\label{s2:lemTower}
For every valid $\HBtp$ run with $d_1\ge\Dstar$ (so $R\ge1$):
\begin{enumerate}
\item[(a)] $\lambda_r\ge d_{r+1}^{1/\Aexp}$ for $r\le R$; $d_{r+2}\le\lambda_r$
for $r\le R-1$; $R-r\le2\logs d_r-1$, in particular
\[
R-r\le2\logs d_r+2\qquad(r\le R);
\]
and $\lambda_r\ge\lambda_{l-2}$ whenever $r\le l-2\le R$.
\item[(b)] For $r\le R$: $M_r\le d_r^2$; $\lambda_r\le\Lambda_r\le2\lambda_r$;
$\lambda_r^{100}\le s_r\le2\Lambda_r^{\sigma}$; $s_r\le P_r$; and $L_Y\le\log
M_r\le2\lambda_r$ for every ancestor $Y$ of round $r$. For $3\le l\le R$:
$M_l\le(\Aexp\log\lambda_{l-2})^{2\Aexp}\le\lambda_{l-2}^{1.6}$;
$P_{l-2}\ge M_l^{13}$; and $P_{l-2}/2\ge M_l\log^4M_l$. For $2\le l\le R$:
$M_{l-1}\ge2M_l$. For $r<R$: $P_r\ge2P_{r+1}$. Moreover
$\sum_{l\le R}1/M_l\le2/\Dstar$ and $\sum_{l\le R}\psi(M_l)\le2.1\,\psi(\Dstar)$
with $\psi(x)\defeq(6\log x+12)/x$. Finally, completing (K1) of
Proposition~\ref{s2:propOV}, the number $\nu_l$ of ancestors of rounds at most
$l-2$ satisfies $\nu_l\le2.74\,n/P_{l-2}$ for $3\le l\le R$.
\item[(c)] For $r\le R$: $\tau_r\le257\,\Lambda_r^{\sigma+2}\le
2^{\sigma+11}\lambda_r^{\sigma+2}$ and $\tau_r/P_r\le2^{\sigma+11}/\lambda_r$;
and $\thGC_r(Z^0)\ge P_r\lambda_r^{-1/2}>\tau_r$ for every round-$r$ pre-part
$Z$.
\item[(d)] For $r\le l\le R$: $\lambda_r\ge2^{2\logs d_r+2}\ge2^{R-r}\ge
2^{l-r}$.
\item[(e)] Put
\[
\epsA\defeq\frac{31\,\eps}{\Cp\log\log\Dstar},
\]
a function of $\Dstar$ alone (it depends neither on the run nor on $n$). Then
\[
\sum_{l\le R}\frac{15.2\,\eps}{\log P_l}\le\epsA.
\]
Consequently $\sum_{l\le R}|D_l|\le\epsA n$ and $\sum_{l\le
R}\sum_{Z\in\Std_l}|A_Z|\le\epsA n$.
\end{enumerate}
\deps{\ref{s2:propDegRec}, \ref{s2:lemLacunary}, \ref{s2:defHBtp}, \ref{s1:condG1}, \ref{s1:condG2}, \ref{s2:propStructure}, \ref{s2:lemCap}, \ref{s2:propOV}, \ref{s2:defAncestors}}
\end{lemma}

\begin{proof}
Every round $r\le R$ has $d_r\ge\Dstar$, so $\log\lambda_r\ge\log\log\Dstar$, and
the items of \Gc{1} (Condition~\ref{s1:condG1}) hold at $\mu=\log\lambda_r$ for
every $r\le R$; they also hold at $\mu=\log\log d$ for every $d\ge\Dstar$. Proposition~\ref{s2:propDegRec} applies at every round $r\le R$.

(a) Proposition~\ref{s2:propDegRec} at round $r$ gives $d_{r+1}\le
\lambda_r^{\Aexp}$. For $r\le R-1$ it also applies at round $r+1$, so
$d_{r+2}\le\lambda_{r+1}^{\Aexp}$ with $\lambda_{r+1}=\log d_{r+1}\le\Aexp\log
\lambda_r$; hence $d_{r+2}\le(\Aexp\log\lambda_r)^{\Aexp}\le\lambda_r$, because
with $u\defeq\log\lambda_r$ we have $\Aexp\log(\Aexp u)\le0.8u\le u$ by \Gc{1}(c)
at $u$; equivalently, $(\Aexp\log\lambda_r)^{\Aexp}\le\lambda_r$.

Let $k\defeq\logs d_r$; $k\ge1$ as $d_r>1$. Suppose $R\ge r+2k$. We show
$d_{r+2i}\le\log^{[i]}d_r$ for $0\le i\le k$ by induction: for $i<k$ we have
$r+2i\le R-2$, so $d_{r+2i+2}\le\log d_{r+2i}\le\log(\log^{[i]}d_r)$, where
$\log^{[i]}d_r>1$ as $i<\logs d_r$. Hence $d_{r+2k}\le\log^{[k]}d_r\le1<\Dstar$,
contradicting $r+2k\le R$. Therefore $R-r\le2k-1\le2\logs d_r+2$. (The weaker
form with ``$+2$'' is the one used in later sections.) Finally
$d_{j+1}\le d_j$ for all $j\le R$ (Proposition~\ref{s2:propStructure}(iii)), so
$\lambda$ is non-increasing in the round.

(b) $M_r\le d_r^2$, $\Lambda_r\le2\lambda_r$ and $s_r\le P_r$ are in
Proposition~\ref{s2:propDegRec}; $\Lambda_r\ge\lambda_r$ and
$s_r\ge\lambda_r^{100}$ are in (the proof of)
Proposition~\ref{s2:propStructure}(i); and $s_r\le\Lambda_r^\sigma+1\le
2\Lambda_r^\sigma$. For an ancestor $Y$ of round $r$, $|V(Y)|\le|Y^0|\le M_r$
by Lemma~\ref{s2:lemCap}(ii), so $L_Y\le\log M_r=\Lambda_r\le2\lambda_r$.

Let $3\le l\le R$. By Proposition~\ref{s2:propDegRec} at rounds $l-1$ and
$l-2$, $d_l\le\lambda_{l-1}^{\Aexp}$ and $\lambda_{l-1}\le\Aexp\log
\lambda_{l-2}$; with $M_l\le d_l^2$ this gives
$M_l\le(\Aexp\log\lambda_{l-2})^{2\Aexp}$. With $u\defeq\log\lambda_{l-2}$, \Gc{1}(c) at $u$ gives
$2\Aexp\log(\Aexp u)=210\log(105u)\le1.6u$, so
$M_l\le\lambda_{l-2}^{1.6}$. Hence $M_l^{13}\le\lambda_{l-2}^{20.8}\le
\lambda_{l-2}^{103}\le P_{l-2}$. As $\Lambda_l\ge40$, $\log M_l=\Lambda_l\le
2^{\Lambda_l/4}=M_l^{1/4}$, so $M_l\log^4M_l\le M_l^2\le M_l^{13}/2\le
P_{l-2}/2$.

For $2\le l\le R$: $M_{l-1}\ge d_{l-1}$ (as in
Proposition~\ref{s2:propStructure}(i)), and $M_l\le d_l^2\le
\lambda_{l-1}^{2\Aexp}=(\log d_{l-1})^{2\Aexp}$; with $y\defeq\log d_{l-1}$ and $v\defeq\log y\ge\log\log\Dstar$, \Gc{1}(a),(b) at $v$
give $v\ge2^8$ and $y=2^v\ge2^{14}\Aexp v^3\ge1+2\Aexp v$, that is, $2^y\ge2y^{2\Aexp}$; so
$M_{l-1}\ge2M_l$. For $r<R$:
$P_{r+1}\le\lambda_{r+1}^{103}+1\le(\Aexp\log\lambda_r)^{103}+1\le\lambda_r+1\le
\lambda_r^{103}/2\le P_r/2$, because
$(\Aexp\log\lambda_r)^{103}\le(\Aexp\log\lambda_r)^{\Aexp}\le\lambda_r$ (the
first inequality as $\Aexp\log\lambda_r\ge\Aexp\ge1$, since $\lambda_r\ge2$; the
second from the proof of (a)) and $\lambda_r\ge2$.

By $M_{l-1}\ge2M_l$ and Lemma~\ref{s2:lemLacunary}(i) applied to $X_l\defeq
1/M_l$, $\sum_{l\le R}1/M_l\le2/M_R\le2/d_R\le2/\Dstar$. The function $\psi$ is
decreasing on $[1,\infty)$, because $\psi'(x)=(6/\ln2-6\log x-12)/x^2$ and
$6/\ln2<12$; and for $2\le l\le R$,
\[
\psi(M_{l-1})\le\psi(2M_l)=\frac{6\log M_l+18}{2M_l}
=\psi(M_l)\Bigl(\frac12+\frac3{6\log M_l+12}\Bigr)\le0.512\,\psi(M_l),
\]
as $\log M_l\ge40$. Hence $\sum_{l\le R}\psi(M_l)\le\psi(M_R)/(1-0.512)\le
2.05\,\psi(M_R)\le2.1\,\psi(\Dstar)$, since $M_R\ge d_R\ge\Dstar$.

Finally, by (K1) of Proposition~\ref{s2:propOV},
$\nu_l\le1.37\,n\sum_{r\le l-2}1/P_r$, and $P_r\ge2P_{r+1}$ for $r<l-2\le R$;
Lemma~\ref{s2:lemLacunary}(i) applied to $X_r\defeq1/P_r$ ($r\le l-2$) gives
$\sum_{r\le l-2}1/P_r\le2/P_{l-2}$, so $\nu_l\le2.74\,n/P_{l-2}$.

(c) $\tau_r=\lceil128\,s_r\Lambda_r^2\rceil\le128\cdot2\Lambda_r^\sigma
\Lambda_r^2+1\le257\,\Lambda_r^{\sigma+2}\le257\cdot2^{\sigma+2}
\lambda_r^{\sigma+2}<2^{\sigma+11}\lambda_r^{\sigma+2}$, as $257\cdot4<2^{11}$.
Since $\Cp=\sigma+3$, $P_r\ge\lambda_r^{\sigma+3}$ and
$\tau_r/P_r\le2^{\sigma+11}/\lambda_r$. For a round-$r$ pre-part $Z$,
$\thGC_r(Z^0)\ge|Z^0|\lambda_r^{-1/2}\ge P_r\lambda_r^{-1/2}\ge
\lambda_r^{102.5}$, while $\tau_r<2^{111}\lambda_r^{102}$; and
$\lambda_r^{1/2}>2^{111}$ because $\log\lambda_r\ge2^8>222$ by \Gc{1}(a).

(d) Put $\mu\defeq\log\lambda_r$; so $\mu\ge2^8$ by \Gc{1}(a). As $d_r=2^{\lambda_r}$ and
$\lambda_r=2^\mu$ with $\lambda_r,\mu>1$, $\logs d_r=2+\logs\mu$. By the fact
$\logs z\le1+\log z$ ($z\ge4$) shown in the proof of
Lemma~\ref{s2:lemLacunary}(iv), $2\logs d_r+2=6+2\logs\mu\le8+2\log\mu\le\mu$,
the last step by \Gc{1}(d) at $\mu$.
So $\lambda_r=2^\mu\ge2^{2\logs d_r+2}$, which is at least $2^{R-r}$ by (a) and
at least $2^{l-r}$ for $l\le R$.

(e) We have $\log P_l\ge\Cp\log\lambda_l$. The function
$F(x)\defeq1/(\Cp\log x)$ is non-increasing and satisfies
$F(2^{x/\Aexp})\le F(x)/2$ for $x\ge\log\Dstar$
(Lemma~\ref{s2:lemLacunary}(iv)). By Lemma~\ref{s2:lemLacunary}(ii),
\[
\sum_{l\le R}\frac1{\log P_l}\le\sum_{l\le R}F(\lambda_l)\le2F(\log\Dstar)=
\frac2{\Cp\log\log\Dstar},
\]
so $\sum_{l\le R}15.2\,\eps/\log P_l\le30.4\,\eps/(\Cp\log\log\Dstar)\le
\epsA$. By (K2) of Proposition~\ref{s2:propOV},
$\sum_l|D_l|\le\sum_l7.6\,\eps n/\log P_l\le\epsA n$ and
$\sum_l\sum_{Z\in\Std_l}|A_Z|\le\sum_l15.2\,\eps n/\log P_l\le\epsA n$.
\end{proof}

\subsection{Rule (GC), ORIGIN\texorpdfstring{$^\tau$}{-tau}, and parentless mass}

\begin{lemma}[rule GC]\label{s2:lemGC}
For every valid $\HBtp$ run and every round $r\le R$:
\begin{enumerate}
\item[(i)] the pre-parts, their leaf graphs, $\Dups_r$, $D_r$, $\home_r$ and
the guest sets $S_Z$ are defined in (R1)--(R4) before and independently of
condition (GC); (GC) only decides whether a pre-part satisfying (L1) and (L2)
is light or standalone;
\item[(ii)] for every light part $Y$ of round $r$ and every guest $x\in S_Y$,
$e_{X^0_Y}(x,Y)<\thGC_r(Y^0)$, where $Y$ denotes the vertex set
$Y^0\setminus S_Y$;
\item[(iii)] GC-parts are standalone; for a GC-part $Y$, every edge of $G'_r$
with both ends in $Y^0$ is assigned at round $r$ (the edges of $X^0_Y$, those
between a guest and $Y^0\setminus S_Y$ included, by (R5)(1)), so no edge of
$G_{r+1}$ has both ends in $V(Y)=Y^0$; and every guest $x\in S_Y$ lies in
$D_r\subseteq\Dups_r$;
\item[(iv)] $\thGC_r(Z^0)>\tau_r$ for every round-$r$ pre-part $Z$.
\end{enumerate}
\deps{\ref{s2:defHBtp}, \ref{s2:lemTower}, \ref{s2:lemEL}, \ref{s2:defAncestors}}
\end{lemma}

\begin{proof}
(i) is the order of definitions in (R3)--(R4): (GC) is evaluated only after
the pre-parts, $D_r$, $\home_r$ and the sets $S_Z$ are fixed, and none of
these uses the outcome of (GC).

(ii) A light pre-part satisfies (GC): no $x\in S_Y$ has $\thGC_r(Y^0)$ or more
neighbours in $Y^0\setminus S_Y$ in $X^0_Y$.

(iii) A GC-part fails (GC) and is therefore standalone by (R4). By (R5)(1) all
edges of $X^0_Y$ are assigned to $Y$. By Lemma~\ref{s2:lemEL}, applied to the
ancestor $Y$ (with $V(Y)=Y^0$), no edge of $G_{r+1}$ has both ends in $Y^0$;
equivalently, every edge of $G'_r$ inside $Y^0$ is assigned at round $r$. A
guest $x\in S_Y$ lies in $D_r$ by definition, hence in at least two pre-parts,
which are two leaves of the two-level recursion; so $x\in\Dups_r$.

(iv) is Lemma~\ref{s2:lemTower}(c).
\end{proof}

\begin{proposition}[ORIGIN$^\tau$]\label{s2:propOrigin}
Fix a valid $\HBtp$ run, a round $r\le R$, a round-$r$ pre-part $Y$ and a
vertex $u\in Y^0\setminus\Dups_r$.
\begin{enumerate}
\item[(a)] $u$ lies in a unique $s=0$ piece $\piece$ and in a unique leaf of the
two-level recursion, namely $Y$ (so $Y^0\subseteq V(\piece)$). Every edge $ux$
of $G_{r+1}$ is of exactly one of the following two types:
\begin{itemize}
\item[($\beta$)] $x\notin Y^0$, and $ux$ was deleted by the $\tau$-run of
$\piece$ (so $x\in V(\piece)$);
\item[($\alpha$)] $Y$ is light, $x\in S_Y$, and $ux\in E(X^0_Y)$ (a
\emph{guest--core edge}: it is not deleted and not assigned at round $r$, and
passes down).
\end{itemize}
A standalone $Y$ (GC-parts included) has no edges of type ($\alpha$). Since
$E(G_l)\subseteq E(G_{r+1})$ for $l>r$, the same classification applies to the
edges $ux$ of $G_l$.
\item[(b)] For every vertex $h$ and every round $l>r$,
\[
\#\bigl\{u\in Y^0\setminus\Dups_r:\ hu\in E(G_l)\text{ is of type }(\beta)
\bigr\}\le\tau_r-1.
\]
\item[(c)] $\sum_Y|Y^0\cap\Dups_r|\le16\,\eps n/\log P_r$, the sum over all
round-$r$ pre-parts $Y$.
\end{enumerate}
An edge of type ($\alpha$) has its core end $u$ outside $\Dups_r$ (it is not
true that ($\alpha$)-edges occur only at vertices of $\Dups_r$); a guest $x$
of a light $Y$ is the guest end of fewer than $\thGC_r(Y^0)$ edges of type
($\alpha$) (Lemma~\ref{s2:lemGC}(ii)); and there is no third type.
\deps{\ref{s2:lemThinCut}, \ref{s2:lemSEP}, \ref{s2:lemCap}, \ref{s2:propOV}, \ref{s2:lemEL}, \ref{s2:defHBtp}, \ref{s2:lemGC}, \ref{s2:lemOVgeneric}, \ref{s2:propStructure}}
\end{proposition}

\begin{proof}
(a) As $u\notin\Dups_r$, $u$ lies in exactly one leaf of the two-level
recursion, and $u\in Y^0$, so that leaf is $Y$. If $u$ lay in two $s=0$ pieces,
it would lie in a leaf below each of them (Lemma~\ref{s2:lemSEP}(0)), i.e.\ in
two leaves; so $u$ lies in a unique piece $\piece$, and $Y$ is a leaf of the
$\tau$-run of $\piece$ (by Lemma~\ref{s2:lemCap}(ii) this is a
Lemma~14$^\tau$ run, although only its split structure is used below). By
Lemma~\ref{s2:lemSEP}(i),(ii) applied to the two-level recursion, every edge of
$G'_r$ at $u$ is either an edge of the leaf $Y$, i.e.\ of $X^0_Y$, or deleted
at some split on the path from the root to $Y$. The splits of the first level
delete nothing (Lemma~\ref{s2:lemOVgeneric}, instance $c=1$), so deleted edges
at $u$ are deleted by the $\tau$-run of $\piece$ and are edges of $\piece$. A
deleted edge $ux$ has $x\notin Y^0$: at the split where it is deleted, $u$ and
$x$ go to different children, each to only one, so they do not lie in a common
leaf. Hence every edge $ux$ of $G'_r$ with $x\in Y^0$ is an edge of $X^0_Y$.

Now let $ux\in E(G_{r+1})$, an edge of $G'_r$ not assigned at round $r$. Since
$u\notin D_r$ (as $D_r\subseteq\Dups_r$), $Y$ is the only pre-part containing
$u$, and $u\notin S_Y$.

\emph{Case $ux$ deleted.} Then $x\notin Y^0$ and $x\in V(\piece)$; this is type
($\beta$). (Such an edge is indeed never assigned: a light part or standalone
pre-part containing both ends would contain $u$, hence would be the light part
of $Y$ or the pre-part $Y$ itself, neither of which contains $x\notin Y^0$.)

\emph{Case $ux\in E(X^0_Y)$.} If $Y$ is standalone, $ux$ is assigned to $Y$ by
(R5)(1), a contradiction. So $Y$ is light. If $x\notin S_Y$, then $ux\in
E(X_Y)$ is assigned by (R5)(1), a contradiction. So $x\in S_Y$: type
($\alpha$). (Such an edge is indeed passed down: the only light part that could
contain $u$ is $Y$, which does not contain $x$, and the only pre-part
containing $u$ is the light pre-part $Y$, so steps (2) and (3) of (R5) do not
apply.)

The types are exclusive ($x\notin Y^0$ in ($\beta$), $x\in S_Y\subseteq Y^0$ in
($\alpha$)). For standalone $Y$, Lemma~\ref{s2:lemEL} also shows directly that
no edge of $G_{r+1}$ has both ends in $V(Y)=Y^0$. By
Proposition~\ref{s2:propStructure}(iii) (or directly from (R1) and (R5)),
$E(G_l)\subseteq E(G_{r+1})$ for $l>r$.

(b) Let $h$ be a vertex and $l>r$. By (a), each $u$ counted has $hu$ deleted by
the $\tau$-run of the unique piece $\piece$ containing $u$, and this piece is
the piece whose $\tau$-run has $Y$ as a leaf; so all counted edges are deleted
edges of one $\tau$-run, and each is of the form $hu$ with $u\in
Y^0\setminus\Dups_r$. A vertex lying in two leaves of the $\tau$-run of
$\piece$ lies in two leaves of the two-level recursion, so
$Y^0\setminus\Dups_r\subseteq Y^0\setminus\Dup$, where $\Dup$ is taken in the
$\tau$-run of $\piece$. Lemma~\ref{s2:lemThinCut}, applied to this $\tau$-run
(all of whose splits are $\tau$-rule splits at threshold $\tau_r$) and its
leaf $Y$, bounds the number of such edges by $\lceil\tau_r\rceil-1=\tau_r-1$
($\tau_r$ is an integer by (R2)), and by $0$ if $h\in Y^0$. Distinct $u$ give
distinct edges $hu$.

(c) is (K3) of Proposition~\ref{s2:propOV}.

The final paragraph of the statement: an ($\alpha$)-edge is an edge $ux$ with
$u\in Y^0\setminus\Dups_r$ by the setting of (a); the ($\alpha$)-edges at a
guest $x$ with core ends in $Y\setminus\Dups_r$ are edges of $X^0_Y$ between
$x$ and $Y=Y^0\setminus S_Y$, and there are fewer than $\thGC_r(Y^0)$ of them
by Lemma~\ref{s2:lemGC}(ii).
\end{proof}

\begin{proposition}[fresh and parentless mass]\label{s2:propParentless}
For every valid $\HBtp$ run with $d_1\ge\Dstar$:
\begin{enumerate}
\item[(i)] If $x$ lies in a round-$r$ pre-part, then $H\defeq\home_r(x)$ is an
ancestor of round $r$ containing $x$: $x\in H^0\setminus S_H=V(H)$ if $H$ is
light, and $x\in H^0=V(H)$ otherwise. Hence, for $Z\in\Std_l$ and $x\in U_Z$,
$x\in F_Z$ iff $j_0(x)\ge l-1$. A vertex is a port of at most one part per
round, so it is a fresh port only at rounds $j_0(x)$ and $j_0(x)+1$,
and\namedlabel{s2:K4}{(K4)}
\begin{equation*}
\sum_{l\le R}\ \sum_{Z\in\Std_l}|F_Z|\le2n.\tag*{(K4)}
\end{equation*}
\item[(ii)] (admissible parents) If $Y$ is a light part of round $r\le l-2$ and
$Z$ is a light part of round $l\le R$, then
\[
|Y|\ge P_r/2\ge P_{l-2}/2\ge M_l\log^4M_l\ge|Z|\,L_Z^4.
\]
\item[(iii)] Let $\Bad$ be \emph{any} set of light parts. Call a pair $(v,Z)$,
with $Z$ a light part of some round $l$ and $v\in Z$, \emph{parentless with
respect to $\Bad$} if no light part outside $\Bad$ of a round at most $l-2$
contains $v$. Then the number of such pairs is at most
\[
2n+\sum_{Y\in\Bad}|Y|\,\bigl(R-r(Y)\bigr).
\]
More precisely, a vertex $v$ is in parentless pairs only at its light parts of
rounds $j_1(v)$ and $j_1(v)+1$, and at the later rounds at which the light part
of round $j_1(v)$ containing $v$ lies in $\Bad$.
\end{enumerate}
\deps{\ref{s2:propStructure}, \ref{s2:lemTower}, \ref{s2:lemCap}, \ref{s2:defAncestors}, \ref{s2:defHBtp}}
\end{proposition}

\begin{proof}
(i) By definition $H=\home_r(x)$ is a pre-part containing $x$. If $H$ is light,
then $x\notin S_H$, because $S_H$ consists of vertices whose home is not $H$;
so $x\in H^0\setminus S_H=V(H)$, and the light part $H$ is an ancestor of round
$r$. If $H$ is standalone, it is an ancestor with $V(H)=H^0\ni x$.

Let $Z\in\Std_l$ and $x\in U_Z\subseteq Z^0$, so $j_0(x)\le l$. If $l\le2$,
then $x\in F_Z$ by definition and $j_0(x)\ge1\ge l-1$. Let $l\ge3$. If
$j_0(x)\le l-2$, then $\home_{j_0(x)}(x)$ is an ancestor of round $j_0(x)\le
l-2$ containing $x$, so $\anc_l(x)\ne\emptyset$ and $x\notin F_Z$. Conversely,
if $\anc_l(x)\ne\emptyset$, then $x\in V(Y)\subseteq Y^0$ for a pre-part $Y$ of
a round at most $l-2$, so $j_0(x)\le l-2$. Hence $x\in F_Z$ iff $j_0(x)\ge
l-1$.

The port sets $U_Z$ of distinct $Z\in\Std_l$ are disjoint
(Proposition~\ref{s2:propStructure}(iv)), so $x$ is a port of at most one part
per round. If $x$ is a fresh port at round $l$ then $j_0(x)\le l$ (it lies in
$Z^0$) and $j_0(x)\ge l-1$, so $l\in\{j_0(x),j_0(x)+1\}$. Summing over $x$
gives $\sum_l\sum_Z|F_Z|\le2n$.

(ii) $|Y|\ge P_r/2$ by Proposition~\ref{s2:propStructure}(iv). As $\lambda$ is
non-increasing in the round (Lemma~\ref{s2:lemTower}(a)) and $r\le l-2$,
$P_r\ge P_{l-2}$. As $l\ge r+2\ge3$, $P_{l-2}/2\ge M_l\log^4M_l$ by
Lemma~\ref{s2:lemTower}(b). Finally $|Z|\le|Z^0|\le M_l$ by
Lemma~\ref{s2:lemCap}(ii), and $L_Z=\log|Z|\le\log M_l$.

(iii) Fix a vertex $v$ that lies in some light part, and let $j_1\defeq
j_1(v)$ and $Y_1(v)$ be the light part of round $j_1$ containing $v$. By
Proposition~\ref{s2:propStructure}(iv), $v$ lies in at most one light part per
round, so the pairs $(v,Z)$ correspond to distinct rounds $l\ge j_1$. If $l\ge
j_1+2$ and $Y_1(v)\notin\Bad$, then $Y_1(v)$ is a light part outside $\Bad$ of
round $j_1\le l-2$ containing $v$, so $(v,Z)$ is not parentless. Hence the
parentless pairs of $v$ are at the rounds $j_1$ and $j_1+1$ (at most two
pairs), and, if $Y_1(v)\in\Bad$, at rounds $l$ with $j_1+2\le l\le R$ (at most
$R-j_1$ pairs). Summing over $v$, the number of parentless pairs is at most
\[
2n+\sum_{v:\,Y_1(v)\in\Bad}\bigl(R-r(Y_1(v))\bigr)\le2n+\sum_{Y\in\Bad}
\sum_{v\in Y}\bigl(R-r(Y)\bigr)=2n+\sum_{Y\in\Bad}|Y|\bigl(R-r(Y)\bigr),
\]
since $v\in Y_1(v)$.

Part (iii) holds for every set $\Bad$. Its argument uses only that the light
part of round $j_1(v)$ containing $v$ is a light part outside $\Bad$; so, when
it is applied, every reason for which a light part is not used as a parent must
place that part in $\Bad$. In the later application this means that $\Bad$
consists of the light parts that are demoted \emph{or} parent-bad, not only the
demoted ones.
\end{proof}

\begin{table}[ht]
\begin{center}
\begin{tabular}{p{4.6cm}p{2.6cm}p{3.4cm}p{3.2cm}}
\hline
quantity & plain $\HBs$ & $\HBtp$ (used here) & where proved\\
\hline
$\sum|Z^0|$ over round-$r$ pre-parts & $1.28\,n$ & $1.37\,n$ ($1.21\,n$ on
the composite tree) & Prop.~\ref{s2:propOV}, (K1)\\
duplication at nodes of size $\ge P$ & $3.5\,\eps S/\log P$ &
$5.5\,\eps S/\log P$ & Lemma~\ref{s2:lem14tau}(b)\\
ancestors of rounds $\le l-2$ ($\nu_l$) & $2.6\,n/P_{l-2}$ &
$2.74\,n/P_{l-2}$ & Lemma~\ref{s2:lemTower}(b)\\
$\dup_r$ (and $|D_r|$) & $4.5\,\eps n/\log P_r$ & $7.6\,\eps n/\log P_r$ &
Prop.~\ref{s2:propOV}, (K2)\\
$\sum_{Z\in\Std_r}|A_Z|$ & $9\,\eps n/\log P_r$ & $15.2\,\eps n/\log P_r$ &
Prop.~\ref{s2:propOV}, (K2)\\
mass of pre-parts failing (L1) & $18\,\eps n/\log P_r$ &
$30.4\,\eps n/\log P_r$ & Prop.~\ref{s2:propOV}, (K3)\\
\hline
\end{tabular}
\end{center}
\caption{Overlap constants of plain $\HBs$ and of $\HBtp$; only the
right-hand column is used in this paper.}\label{s2:tabTauConstants}
\end{table}

\begin{remark}[constants of $\HBtp$ versus plain $\HBs$]\label{s2:remTauConstants}
Plain $\HBs$ is the procedure of Definition~\ref{s2:defHBtp} in which the
second level of (R3) is the recursion of \cite[Lemma~14]{BM} at $s=s_l$ (no
$\tau$-rules) and (GC) is omitted; $\HBt$ has the $\tau$-rules but no (GC). The
overlap constants change because the separator of a $\tau$-split may exceed
the separator $N$ of \cite{BM} by $|\Hout|+|\Hin|$, which raises the charge in
Lemma~\ref{s2:lemOVgeneric} from $c=1$ to $c=1.6$ at the second level.
Table~\ref{s2:tabTauConstants} lists the values; only the right-hand column is
used in this paper. The left-hand column is recorded for comparison only;
it is neither used nor proved here.
\deps{\ref{s2:propOV}, \ref{s2:lem14tau}, \ref{s2:lemTower}, \ref{s2:defHBtp}, \ref{s2:lemOVgeneric}}
\end{remark}


\section{Random splitting, linking, and the lending colouring}\label{s3:sec}

This section provides the probabilistic tools that every later routing step
uses, together with the stage-1 lending data.
\begin{itemize}
\item The first subsection proves monotonicity facts and
  Lemma~15$^+$. A uniformly random $k$-colouring of the edges of an
  $(\varepsilon',s)$-expander leaves every colour class an
  $(\varepsilon',s/(2k))$-expander: the first parameter is untouched and
  the second loses only the factor $2k$. Lemma~15$^+$ replaces
  \cite[Lemma~15]{BM} (quoted as Cited result~\ref{s1:citLem15}), which is not
  used.
\item The second subsection proves Theorem~16$^*$. It says that a
  robust expander is path connected through a random vertex set, in the
  multiset, for-all form of \cite[Definition~7]{BM}. The requirement on
  the robustness is linear in the multiplicity and polynomial in the
  inverse density. Its proof re-runs the proof of \cite[Theorem~16]{BM}
  (Cited result~\ref{s1:citThm16}) with explicit parameters.
\item The third subsection proves Lemma~HB by one application of Hall's
  theorem.
\item The fourth subsection defines the two-stage lending colouring
  and the JS labels, verifies the numerical conditions they need (the
  COL-JV table), and proves Lemma~COL.
\end{itemize}
Every statement of this section is proved in full here, from the results
stated in Section~\ref{s1:sec} and, in the fourth subsection, from the hierarchy of
Section~\ref{s2:sec}.

\medskip\noindent\textbf{Conventions for this section.} Logarithms are to
base $2$ (Convention~\ref{s1:convGraphs}), and $\ln$ is the natural
logarithm. For $\rho\in[0,1]$, a \emph{$\rho$-random subset} of a finite
set $S$ is a random subset of $S$ that contains each element of $S$
independently with probability $\rho$. The letter $\varepsilon'$ always
denotes a generic expansion parameter in $(0,1]$. The fixed constant
$\eps=2^{-5}$ of Definition~\ref{s1:defConstants} enters this section only
through the values $\varepsilon_Y\in\{2^{-5},2^{-6}\}$ of
Definition~\ref{s2:defAncestors}. For a graph $G$, $U\subseteq V(G)$ and
$F\subseteq E(G)$, the set $\Nbr_{G-F}(U)$ is the outside neighbourhood of
$U$ in the graph $G-F$ (Convention~\ref{s1:convGraphs}). The ball
$\Ball{i}{G-F}{U,V}$ is the set of vertices of $V$ that can be reached
from $U$ by a path of length at most $i$ in $G-F$ whose interior vertices
all lie in $V$ (Cited result~\ref{s1:citNotation}). Path connectivity
(Cited result~\ref{s1:citDef7}) always refers to \emph{multisets} of pairs of
distinct vertices; see the remark on multisets at the end of the first subsection.

%% ----------------------------------------------------------------------
\subsection{Monotonicity and random splitting}

\begin{lemma}[monotonicity and the for-all property]\label{s3:lemMonotone}
\begin{itemize}
\item[(i)] Let $H$ be an $(\varepsilon',s)$-expander and let $H'\supseteq
  H$ be a graph with $V(H')=V(H)$. Then $H'$ is an
  $(\varepsilon',s)$-expander. Moreover, every $(\varepsilon',s)$-expander
  is an $(\varepsilon'',s')$-expander whenever $\varepsilon''\le
  \varepsilon'$ and $s'\le s$.
\item[(ii)] Path connectivity is monotone. Let $G\subseteq G'$ be graphs
  with $V(G)=V(G')$, let $V\subseteq V'\subseteq V(G)$, and suppose
  $\ell\le\ell'$ and $t'\le t$. If $G$ is $(\ell,t)$-path connected
  through $V$, then $G'$ is $(\ell',t')$-path connected through $V'$.
\item[(iii)] Whether a graph $X$ is $(\ell,t)$-path connected through a
  set $V$ depends only on the pair $(X,V)$. On this event, required paths
  exist for \emph{every} multiset $\mathcal P$ of pairs of distinct
  vertices of $X$ in which each vertex lies in at most $t$ pairs. In
  particular, $\mathcal P$ may be chosen after $X$ and $V$ are revealed,
  as an arbitrary function of them and of any further randomness
  (adaptively). If several multisets $\mathcal P_1,\dots,\mathcal P_q$ are
  given and every vertex lies in at most $t$ pairs of the union multiset
  $\mathcal P_1+\dots+\mathcal P_q$, then one application of the property
  to this union joins all their pairs by pairwise edge-disjoint paths
  through $V$ of length at most $\ell$ (joint routing).
\end{itemize}
\deps{\ref{s1:citDef7}, \ref{s1:citDef11}}
\end{lemma}

\begin{proof}
(i) Put $n:=|V(H)|$. Let $U\subseteq V(H')$ with $1\le|U|\le 2n/3$, and
let $F\subseteq E(H')$ with $|F|\le s|U|$. Then $F\cap E(H)$ is a set of
at most $s|U|$ edges of $H$, and $H-(F\cap E(H))\subseteq H'-F$. Hence
\[
|\Nbr_{H'-F}(U)|\ \ge\ |\Nbr_{H-(F\cap E(H))}(U)|\ \ge\ \varepsilon'|U|/\log^2 n
\]
by Definition~11 of \cite{BM} (Cited result~\ref{s1:citDef11}) applied to $H$.
Both graphs have $n$ vertices, so the logarithm in the bound is the same.
The second claim holds because the condition of Cited result~\ref{s1:citDef11}
becomes weaker when $\varepsilon'$ or $s$ decreases.

(ii) Let $\mathcal P$ be a multiset of pairs of distinct vertices of
$G'$ in which every vertex lies in at most $t'\le t$ pairs. The
hypothesis on $G$ applies to $\mathcal P$, since $V(G)=V(G')$. It gives
pairwise edge-disjoint paths of $G$, one for each pair, each of length at
most $\ell\le\ell'$ and with all interior vertices in $V\subseteq V'$.
These are paths of $G'$ with the required properties.

(iii) By Cited result~\ref{s1:citDef7}, the property says that for every
admissible multiset the paths exist. It mentions no other data, so it is
a property of $(X,V)$ alone, and on this event it applies to every
admissible multiset, however that multiset was chosen. The union
$\mathcal P_1+\dots+\mathcal P_q$ is itself an admissible multiset, which
gives the last sentence.
\end{proof}

\begin{lemma}[Lemma 15$^+$: random splitting without loss]\label{s3:lemL15p}
Let $X$ be an $N$-vertex $(\varepsilon',s)$-expander with
$0<\varepsilon'\le1$, put $L:=\log N$, and let $k\ge1$ be an integer with
$s\ge 40kL$. Give every edge of $X$ a colour from $[k]$, independently
and uniformly at random, and for $i\in[k]$ let $X_i$ be the graph with
vertex set $V(X)$ whose edges are the edges of colour $i$. Then for each
$i\in[k]$, $X_i$ is an $(\varepsilon',s/(2k))$-expander with probability
at least $1-2N^{-5}$. Hence all $k$ classes are
$(\varepsilon',s/(2k))$-expanders with probability at least $1-2kN^{-5}$.
\deps{\ref{s1:citDef11}, \ref{s1:citChernoff}}
\end{lemma}

\begin{proof}
If $N=1$, then no set $U$ satisfies $1\le|U|\le 2N/3$, so every graph on
one vertex is an expander for all parameters, and there is nothing to
prove. Assume $N\ge2$, so $L\ge1$. Fix $i\in[k]$ and write
$H:=X_i$.

Fix $U\subseteq V(X)$ with $1\le u:=|U|\le 2N/3$. Put
$m:=\lceil\varepsilon' u/L^2\rceil$ and
$\mathcal N:=\Nbr_X(U)$. Since $\varepsilon'\le1\le L^2$, we have
$1\le m\le u$.

\smallskip\noindent\emph{Step 1 (every small deletion leaves many edges).}
We claim that
\begin{equation}\label{s3:eqL15step1}
e_X(U,\mathcal N\setminus\mathcal T)>su\qquad\text{for every }\mathcal T\subseteq\mathcal N\text{ with }|\mathcal T|\le m-1 .
\end{equation}
Suppose instead that $e_X(U,\mathcal N\setminus\mathcal T)\le su$ for some
such $\mathcal T$, and let $F$ be the set of edges of $X$ between $U$ and
$\mathcal N\setminus\mathcal T$. Then $|F|\le su$, and every neighbour of
$U$ in $X-F$ outside $U$ lies in $\mathcal T$. Hence
$|\Nbr_{X-F}(U)|\le m-1<\varepsilon' u/L^2$, which contradicts the
expansion of $X$ (Cited result~\ref{s1:citDef11}).

\smallskip\noindent\emph{Step 2 (what a failure looks like).} Suppose
that $U$ violates the $(\varepsilon',s/(2k))$-condition in $H$. Then
there is $F\subseteq E(H)$ with $|F|\le su/(2k)$ and
$|\Nbr_{H-F}(U)|<\varepsilon'u/L^2$. Since $|\Nbr_{H-F}(U)|$ is an
integer, it is at most $m-1$. Put $\mathcal T:=\Nbr_{H-F}(U)$. Since
$H\subseteq X$, we have $\mathcal T\subseteq\mathcal N$ and
$|\mathcal T|\le m-1$, and every edge of $H$ between $U$ and
$\mathcal N\setminus\mathcal T$ lies in $F$. Hence
$e_H(U,\mathcal N\setminus\mathcal T)\le su/(2k)$.

\smallskip\noindent\emph{Step 3 (one pair $(U,\mathcal T)$).} Fix
$\mathcal T\subseteq\mathcal N$ with $|\mathcal T|\le m-1$. Since $U$ and
$\mathcal N\setminus\mathcal T$ are disjoint, the edges of $X$ between
them are $e:=e_X(U,\mathcal N\setminus\mathcal T)$ distinct edges, and
$e>su$ by \eqref{s3:eqL15step1}. Each of them lies in $H$ independently
with probability $1/k$. So
$Z:=e_H(U,\mathcal N\setminus\mathcal T)\sim\Bin(e,1/k)$, with mean
$\mu_Z:=e/k>su/k$. The event $\{Z\le su/(2k)\}$ is contained in
$\{Z<\mu_Z/2\}$. By the lower Chernoff bound (Cited result~\ref{s1:citChernoff})
with $\delta=1/2$,
\[
\Prob\big(Z\le su/(2k)\big)\le e^{-\mu_Z/8}\le e^{-su/(8k)}\le e^{-5uL},
\]
where the last step uses $s\ge 40kL$.

\smallskip\noindent\emph{Step 4 (union bound).} There are at most $N^u$
sets $U$ of size $u$. The number of sets $\mathcal T$ with
$|\mathcal T|\le m-1$ is at most $\sum_{j<m}N^j\le N^m\le N^u$. Since
$e^{-5uL}=N^{-5u\log e}\le N^{-7.2u}$, the probability that $H$ is not an
$(\varepsilon',s/(2k))$-expander is at most
\[
\sum_{u\ge1}N^{2u}\cdot N^{-7.2u}\ \le\ \frac{N^{-5.2}}{1-N^{-5.2}}\ \le\ 2N^{-5}.
\]
The final statement follows by a union bound over the $k$ colours.
Neither $\varepsilon'$ nor the logarithm $L$ changes, because $X_i$ has
the same vertex set as $X$.
\end{proof}

\begin{remark}[multisets]\label{s3:remMultiset}
In \cite[Definition~7]{BM} (Cited result~\ref{s1:citDef7}), the collection
$\mathcal P$ is a \emph{multiset} of pairs of distinct vertices.
\begin{itemize}
\item The same pair may occur several times, and each occurrence needs a
  path of its own.
\item The endpoints are arbitrary vertices of the graph; they need not
  lie in $V$.
\item The proof of \cite[Lemma~9]{BM} (Cited result~\ref{s1:citLem9}) indexes
  the pairs by $i\in[r]$, so it applies verbatim to multisets.
\item Its counting step also holds when pairs are counted with
  multiplicity. The step says that a maximal subfamily $I'\subseteq I$ of
  pairwise vertex-disjoint pairs satisfies $2t|I'|\ge|I|$. It holds
  because every pair of $I$ meets one of the $2|I'|$ vertices covered by
  $I'$, and each vertex lies in at most $t$ pairs, counted with
  multiplicity.
\end{itemize}
All path-connectivity statements of this paper are meant in this
multiset sense, and the proofs of this subsection and the next treat the
pairs as an indexed family. Repeated pairs genuinely occur later: in the
routing inside the vortex devices of Section~4, and in the joint routing
of the JS-LC systems in Section~6.
\deps{\ref{s1:citDef7}, \ref{s1:citLem9}}
\end{remark}

%% ----------------------------------------------------------------------
\subsection{Linking through random sets: Theorem 16\texorpdfstring{$^*$}{*}}

This subsection proves Theorem~16$^*$, stated at its end. The route is
that of \cite[Section~4]{BM}:
\begin{enumerate}
\item well-expanding sets grow into a random set by sprinkling
  (Lemma~17$^*$ below, modelled on \cite[Lemma~17]{BM});
\item every set contains a well-expanding subset (Proposition~18$^*$
  and Lemma~19$^*$ below, modelled on \cite[Proposition~18 and
  Lemma~19]{BM});
\item ball expansion implies multiset linking (Lemma~9$_\rho$ below,
  modelled on \cite[Lemma~9]{BM});
\item a random edge-splitting (Lemma~\ref{s3:lemL15p}) assembles these.
\end{enumerate}
Every parameter is explicit, and the density $\rho$ and the multiplicity
$t$ are free.

\medskip\noindent\textbf{Local parameters.} Let $n\ge2$ be the number of
vertices of the graph at hand, $L:=\log n$, $\varepsilon'\in[2^{-7},1]$
its expansion parameter, $\rho\in(0,1]$ a density and $t\ge1$ a
multiplicity. The following quantities are local to this subsection. All
carry a star subscript; the letters $\ell$, $g$, $p$, $q$, $d$,
$\lambda$, $\Delta$, $\sigma$, $\theta$, $\mu$, $\bar s$, $K$, $M$, $t_0$
without a star keep the meaning they have elsewhere. Throughout,
$\ell_*$, $d_*$, $\lambda_*$, $\Delta_*$, $M_*$ and $K_*$ are integers.
\begin{equation}\label{s3:eqStar}
\begin{aligned}
&\ell_*:=\lfloor 2^{10}L^3\rfloor,\qquad g_*:=\frac{\varepsilon'}{8L^2},\qquad q_*:=0.9\rho,\qquad
 p_*\in(0,1]\ \text{ given by }\ (1-p_*)^{\ell_*-1}=\frac{1-\rho}{1-0.9\rho},\\
&d_*:=\lceil 1/p_*\rceil,\qquad \lambda_*:=\lceil 6/p_*\rceil,\qquad
 \Delta_*:=\lambda_*+\lceil d_*\lambda_*/3\rceil,\qquad \sigma_*:=\frac{24L^2}{\varepsilon'},\\
&\theta_*:=\frac{2^{19}\ell_*^2L^3}{\varepsilon'\rho^2},\qquad
 \mu_*:=\frac{\varepsilon'}{6\theta_*L^2},\qquad
 \bar s_*:=\frac{2^{28}tL^8}{\rho},\qquad
 M_*:=\lceil 2.1/\rho\rceil,\qquad
 K_*:=\Big\lceil\frac{6\,\bar s_*\,\theta_*L^2}{\varepsilon'}\Big\rceil .
\end{aligned}
\end{equation}
The integer $t_{0*}$ is defined inside the proof of Lemma~9$_\rho$
below. We record six elementary facts, each with its
proof. Write $p:=p_*$.
\begin{itemize}
\item[(S1)]\namedlabel{s3:eqS1}{(S1)} $2^{10}L^3-1<\ell_*\le 2^{10}L^3$ and $\ell_*\ge 2^{10}$.
  \emph{Proof:} $L\ge1$ because $n\ge2$.
\item[(S2)]\namedlabel{s3:eqS2}{(S2)} The number $p_*$ exists, is unique, and satisfies $p_*\ge 0.1\rho/\ell_*$, so
  $p_*^{-2}\le 100\ell_*^2\rho^{-2}$. If $V_1,\dots,V_{\ell_*}$ are independent random subsets of a
  set $S$, where $V_i$ is $p_*$-random for $i<\ell_*$ and $V_{\ell_*}$ is $q_*$-random, then
  $V_1\cup\dots\cup V_{\ell_*}$ is a $\rho$-random subset of $S$.
  \emph{Proof:} since $\ell_*\ge2$, the map $x\mapsto(1-x)^{\ell_*-1}$ is a continuous decreasing
  bijection of $[0,1]$ onto $[0,1]$, and the value $(1-\rho)/(1-0.9\rho)$ lies in $[0,1)$. Hence
  $p_*\in(0,1]$ exists and is unique. By Bernoulli's inequality,
  $1-(\ell_*-1)p\le(1-p)^{\ell_*-1}=1-0.1\rho/(1-0.9\rho)$. Hence $(\ell_*-1)p\ge 0.1\rho$, which
  gives $p\ge 0.1\rho/(\ell_*-1)\ge 0.1\rho/\ell_*$. For the last claim, an element avoids the union
  with probability $(1-p)^{\ell_*-1}(1-0.9\rho)=1-\rho$, independently over elements.
\item[(S3)]\namedlabel{s3:eqS3}{(S3)} $\Delta_*\le 2p^{-2}+8.34p^{-1}+2.34$. Hence
  $\Delta_*\le 3p^{-2}$ if $p\le 0.11$, and $\Delta_*\le 13p^{-2}$ for every $p\in(0,1]$.
  \emph{Proof:} $d_*\le p^{-1}+1$ and $\lambda_*\le 6p^{-1}+1$, so
  $d_*\lambda_*\le 6p^{-2}+7p^{-1}+1$. Then
  $\Delta_*\le \lambda_*+d_*\lambda_*/3+1\le(6p^{-1}+1)+(2p^{-2}+\tfrac73p^{-1}+\tfrac13)+1
  =2p^{-2}+\tfrac{25}{3}p^{-1}+\tfrac73$.
  If $p\le0.11$, then $8.34p^{-1}\le 0.92p^{-2}$ and $2.34\le 0.03p^{-2}$. If $p\le1$, then
  $8.34p^{-1}+2.34\le 10.68p^{-2}$.
\item[(S4)]\namedlabel{s3:eqS4}{(S4)} $8d_*\lambda_*\le 112p^{-2}\le\theta_*$ and
  $\theta_*\le 2^{39}L^9/(\varepsilon'\rho^2)\le 2^{46}L^9\rho^{-2}$.
  \emph{Proof:} $8d_*\lambda_*\le 8(6+7+1)p^{-2}$ since $p\le1$. By~\ref{s3:eqS2},
  $112p^{-2}\le 11200\,\ell_*^2\rho^{-2}\le 2^{19}\ell_*^2L^3/(\varepsilon'\rho^2)$. The upper bound
  uses $\ell_*^2\le 2^{20}L^6$ and $\varepsilon'\ge2^{-7}$.
\item[(S5)]\namedlabel{s3:eqS5}{(S5)} $K_*\ge\bar s_*/\mu_*$;
  $K_*\le(6\cdot2^{81}+1)\,tL^{19}\rho^{-3}\le 2^{83.6}\,tL^{19}\rho^{-3}$;
  $2K_*(\theta_*+1)\le 2^{130.6}\,tL^{28}\rho^{-5}$; and $40K_*L\le 2K_*(\theta_*+1)$.
  \emph{Proof:} $\bar s_*/\mu_*=6\bar s_*\theta_*L^2/\varepsilon'\le K_*$. By~\ref{s3:eqS4},
  $6\bar s_*\theta_*L^2/\varepsilon'\le 6\cdot2^{28}tL^8\rho^{-1}\cdot 2^{39}L^9\rho^{-2}\cdot L^2\cdot
  2^{14}=6\cdot2^{81}tL^{19}\rho^{-3}$, since $\varepsilon'^{-2}\le2^{14}$. The rounding adds at most
  $1\le tL^{19}\rho^{-3}$. Likewise $\theta_*+1\le(2^{46}+1)L^9\rho^{-2}$, and
  $\log\big(2(6\cdot2^{81}+1)(2^{46}+1)\big)<130.6$. Finally
  $\theta_*\ge 2^{19}L^3\ge 20L$.
\item[(S6)]\namedlabel{s3:eqS6}{(S6)} $\sigma_*\ge24$, $\lambda_*p\ge6$ and $1\le d_*p\le 1+p\le2$.
  \emph{Proof:} immediate from~\eqref{s3:eqStar}, since $\varepsilon'\le1\le L$.
\end{itemize}
The values $\sigma_*,\lambda_*,\Delta_*,d_*$ of~\eqref{s3:eqStar} satisfy the parameter
hypotheses of Proposition~13$^*$ below: $\sigma_*=24L^2/\varepsilon'$ holds with equality, and
$\Delta_*-\lambda_*=\lceil d_*\lambda_*/3\rceil\ge d_*\lambda_*/3=8d_*\lambda_*L^2/(\varepsilon'\sigma_*)$.
The proposition itself is stated for arbitrary numbers with these properties.

\begin{proposition}[Proposition 13$^*$: stars or bipartite]\label{s3:propP13s}
Let $G$ be an $n$-vertex $(\varepsilon',s_1)$-expander with $n\ge2$ and
$2^{-7}\le\varepsilon'\le1$, and put $L:=\log n$. Let
$W\subseteq V(G)$ with $1\le|W|\le 2n/3$, and let $F\subseteq E(G)$ with
$|F|\le s_1|W|/4$. Let $\sigma_*>0$ be real and let
$\lambda_*,d_*\ge1$ and $\Delta_*\ge\lambda_*$ be integers such that
\[
\sigma_*\ge\frac{24L^2}{\varepsilon'},\qquad
\Delta_*-\lambda_*\ge\frac{8d_*\lambda_*L^2}{\varepsilon'\sigma_*},\qquad
s_1\ge 8d_*\lambda_* .
\]
Then $G-F$ contains one of the following:
\begin{itemize}
\item[(a)] at least $|W|/\sigma_*$ vertex-disjoint stars, each with
  exactly $\lambda_*$ leaves, its centre in $W$ and all its leaves in
  $V(G)\setminus W$; or
\item[(b)] a bipartite subgraph $H\subseteq G-F$ with sides $W$ and
  $X\subseteq V(G)\setminus W$, such that $|X|\ge\varepsilon'|W|/(2L^2)$,
  every vertex of $X$ has degree exactly $d_*$ in $H$, and every vertex
  of $W$ has degree at most $\Delta_*$ in $H$.
\end{itemize}
\deps{\ref{s1:citProp12}, \ref{s1:citProp13}, \ref{s1:citDef11}}
\end{proposition}

\begin{proof}
We follow the proof of \cite[Proposition~13]{BM}
(Cited result~\ref{s1:citProp13}) with general parameters. For $U\subseteq
V(G)$ and an integer $d\ge1$, write $\Nbr_{G-F,d}(U)$ for the set of
vertices outside $U$ that have at least $d$ neighbours in $U$ in the
graph $G-F$. We use \cite[Proposition~12]{BM} (Cited result~\ref{s1:citProp12}):
if $|U|\le 2n/3$, $|F|\le s_1|U|/2$ and $0<d\le s_1$, then
$|\Nbr_{G-F}(U)|\ge s_1|U|/(2d)$ or
$|\Nbr_{G-F,d}(U)|\ge\varepsilon'|U|/L^2$.

Assume that (a) fails. Take a maximal collection of vertex-disjoint stars
in $G-F$, each with exactly $\lambda_*$ leaves, its centre in $W$ and its
leaves in $V(G)\setminus W$. Let $\mathrm{Cen}\subseteq W$ be the set of
centres and $\mathrm{Lvs}\subseteq V(G)\setminus W$ the set of leaves.
Since (a) fails, $|\mathrm{Cen}|<|W|/\sigma_*$, and
$|\mathrm{Lvs}|=\lambda_*|\mathrm{Cen}|<\lambda_*|W|/\sigma_*$.
By maximality, every $w\in W\setminus\mathrm{Cen}$ has at most
$\lambda_*-1$ neighbours in $G-F$ in $V(G)\setminus(W\cup\mathrm{Lvs})$.
(Otherwise a new star with centre $w$ and $\lambda_*$ such leaves could be
added. It is disjoint from the others, because $w\notin\mathrm{Cen}$,
$w\notin\mathrm{Lvs}$ and its leaves avoid $W\cup\mathrm{Lvs}$.) Every
vertex of $\Nbr_{G-F}(W\setminus\mathrm{Cen})$ lies in $\mathrm{Cen}$, in
$\mathrm{Lvs}$, or in $V(G)\setminus(W\cup\mathrm{Lvs})$. Hence
\begin{equation}\label{s3:eqP13a}
|\Nbr_{G-F}(W\setminus\mathrm{Cen})|
<\frac{(1+\lambda_*)|W|}{\sigma_*}+(\lambda_*-1)|W|
\le\Big(\frac{1+\lambda_*}{24}+\lambda_*-1\Big)|W| ,
\end{equation}
where we used $\sigma_*\ge 24L^2/\varepsilon'\ge24$.

\smallskip\noindent\emph{The greedy bipartite graph.} Label the vertices
of $V(G)\setminus W$ as $v_1,\dots,v_{r'}$ and let $H_0$ be the graph on
$W$ with no edges. For $i=1,\dots,r'$, check whether $v_i$ has $d_*$
neighbours in $G-F$ inside $W$, each of degree less than $\Delta_*$ in
$H_{i-1}$. If so, let $H_i$ be $H_{i-1}$ together with $v_i$ and these
$d_*$ edges, and put $v_i$ into $X$. Otherwise put $H_i:=H_{i-1}$.
Finally let $H:=H_{r'}$.
\begin{itemize}
\item $H\subseteq G-F$ is bipartite with sides $W$ and $X$.
\item Every vertex of $X$ has degree exactly $d_*$ in $H$.
\item Every vertex of $W$ has degree at most $\Delta_*$ in $H$: a degree
  increases, by one, only when it is below the integer $\Delta_*$.
\end{itemize}
It remains to show $|X|\ge\varepsilon'|W|/(2L^2)$.

\smallskip\noindent\emph{Saturated vertices.} Let
$\mathrm{Sat}:=\{w\in W\setminus\mathrm{Cen}:\deg_H(w)=\Delta_*\}$. A
vertex $w\in\mathrm{Sat}$ has at most $\lambda_*-1$ neighbours of $G-F$
in $V(G)\setminus(W\cup\mathrm{Lvs})$. Hence it has at most
$\lambda_*-1$ $H$-neighbours in $X\setminus\mathrm{Lvs}$, and so more
than $\Delta_*-\lambda_*$ $H$-neighbours in $X\cap\mathrm{Lvs}$. Every
vertex of $X$ has exactly $d_*$ $H$-neighbours, so
\[
(\Delta_*-\lambda_*)|\mathrm{Sat}|\le e_H(\mathrm{Sat},X\cap\mathrm{Lvs})\le d_*|\mathrm{Lvs}|<\frac{d_*\lambda_*|W|}{\sigma_*}.
\]
The hypothesis on $\Delta_*-\lambda_*$ forces $\Delta_*>\lambda_*$, and it
gives $|\mathrm{Sat}|<\varepsilon'|W|/(8L^2)$. Put
$\mathrm{Bl}:=\mathrm{Cen}\cup\mathrm{Sat}$. Then
\begin{equation}\label{s3:eqP13b}
|\mathrm{Bl}|<\frac{|W|}{\sigma_*}+\frac{\varepsilon'|W|}{8L^2}\le\frac{\varepsilon'|W|}{24L^2}+\frac{\varepsilon'|W|}{8L^2}=\frac{\varepsilon'|W|}{6L^2}\le\frac{|W|}{6},
\end{equation}
so $|W\setminus\mathrm{Bl}|\ge 5|W|/6$.

\smallskip\noindent\emph{Applying Proposition~12.} The set
$U:=W\setminus\mathrm{Bl}$ has $1\le|U|\le 2n/3$, and
$|F|\le s_1|W|/4\le s_1|U|/2$. Also $1\le d_*\le s_1$, since
$s_1\ge 8d_*\lambda_*$. By Cited result~\ref{s1:citProp12} one of two cases
holds.

\emph{Case 1: $|\Nbr_{G-F}(U)|\ge s_1|U|/(2d_*)$.} Then
$|\Nbr_{G-F}(U)|\ge s_1|W|/(3d_*)\ge 8\lambda_*|W|/3$. Let
$x\in\Nbr_{G-F}(U)$. Either $x\notin W\setminus\mathrm{Cen}$, and then
$x\in\Nbr_{G-F}(W\setminus\mathrm{Cen})$; or $x\in W\setminus\mathrm{Cen}$
with $x\notin U$, and then $x\in\mathrm{Sat}\subseteq W$. Hence
$|\Nbr_{G-F}(W\setminus\mathrm{Cen})|\ge(8\lambda_*/3-1)|W|$. Comparing
with~\eqref{s3:eqP13a} gives $8\lambda_*/3<(1+\lambda_*)/24+\lambda_*$,
that is, $39\lambda_*<1$. This contradicts $\lambda_*\ge1$.

\emph{Case 2: $|\Nbr_{G-F,d_*}(U)|\ge\varepsilon'|U|/L^2$.} Let
$v\in\Nbr_{G-F,d_*}(U)\setminus\mathrm{Bl}$. Then $v\notin W$, since
$v\notin U$ and $v\notin\mathrm{Bl}$, and $v$ has at least $d_*$
neighbours of $G-F$ in $U$. Vertices of $U$ lie in
$W\setminus\mathrm{Cen}$ and not in $\mathrm{Sat}$, so they have
$H$-degree below $\Delta_*$, and their degree in $H_{i-1}$ is at most
that. So when $v=v_i$ was processed, a star was available, and $v\in X$.
Therefore, using~\eqref{s3:eqP13b},
\[
|X|\ge|\Nbr_{G-F,d_*}(U)|-|\mathrm{Bl}|\ge\frac{\varepsilon'}{L^2}\cdot\frac{5|W|}{6}-\frac{\varepsilon'|W|}{6L^2}=\frac{2\varepsilon'|W|}{3L^2}\ge\frac{\varepsilon'|W|}{2L^2}.
\]
So (b) holds.
\end{proof}

\begin{lemma}[Lemma 17$^*$: sprinkling at density $\rho$]\label{s3:lemL17s}
Let $G$ be an $n$-vertex $(\varepsilon',s_1)$-expander with $n\ge2$ and
$2^{-7}\le\varepsilon'\le1$, let $\rho\in(0,1]$, take the parameters
of~\eqref{s3:eqStar} with $L:=\log n$, and assume $s_1\ge 8d_*\lambda_*$.
Let $V_1,\dots,V_{\ell_*}$ be independent random subsets of $V(G)$,
where $V_i$ is $p_*$-random for $i<\ell_*$ and $V_{\ell_*}$ is
$q_*$-random, and put $V:=V_1\cup\dots\cup V_{\ell_*}$. By~\ref{s3:eqS2},
$V$ is a $\rho$-random subset of $V(G)$. Call $U\subseteq V(G)$
\emph{well-expanding} if $|\Nbr_G(U)|\ge\theta_*|U|$. Then for every
well-expanding $U$ and every $F\subseteq E(G)$ with $|F|\le|U|$,
\[
\Prob\Big(|\Ball{\ell_*}{G-F}{U,V}|\le|V|/2\Big)\le\exp(-7|U|L).
\]
Since this event depends only on $V$, the same bound holds for every
$\rho$-random subset $V$ of $V(G)$.
\deps{\ref{s3:propP13s}, \ref{s1:citLem17}, \ref{s1:lemBBD}, \ref{s1:citChernoff}, \ref{s1:citChernoffGen}}
\end{lemma}

\begin{proof}
The structure is that of the proof of \cite[Lemma~17]{BM}
(Cited result~\ref{s1:citLem17}). The batch probabilities are adapted to the
density $\rho$, and the concentration step in case~(b) uses the Bernstein
inequality for bounded differences (Lemma~\ref{s1:lemBBD}).

If $U=\emptyset$ the bound is $1$, so let $u:=|U|\ge1$.

\smallskip\noindent\emph{Step 0: $n$ and $\rho n$ are large.} Since $U$
is well-expanding, $\theta_*u\le|\Nbr_G(U)|<n$. Hence
$n>\theta_*\ge 2^{19}\ell_*^2L^3/\rho^2$ (as $\varepsilon'\le1$), and so
$\rho n\ge\rho^2n>2^{19}\ell_*^2L^3\ge 2^{39}$ by~\ref{s3:eqS1}. Also
$7uL<7nL/\theta_*\le 7\rho^2n/(2^{19}\ell_*^2L^2)\le 2^{-36}\rho n$ (using
$\varepsilon'\le1$, $\rho\le1$ and $\ell_*\ge2^{10}$).

\smallskip\noindent\emph{Step 1: the reached sets.} For
$0\le i\le\ell_*$, let $\mathcal B_i$ be the set of vertices $v\in V(G)$
that can be reached from $U$ by a path in $G-F$ of length at most $i$
whose interior vertices (if any) lie in $V_1\cup\dots\cup V_{i-1}$. The
vertices of $\mathcal B_i$ themselves need not lie in any $V_j$. Then
$\mathcal B_0=U$, $\mathcal B_1=U\cup\Nbr_{G-F}(U)$ and
$\mathcal B_0\subseteq\mathcal B_1\subseteq\dots\subseteq\mathcal
B_{\ell_*}$. The set $\mathcal B_i$ is determined by $U$, $F$ and
$V_1,\dots,V_{i-1}$, so it is independent of $V_i,\dots,V_{\ell_*}$.
\begin{itemize}
\item[(O1)] Let $v\in\Nbr_{G-F}(\mathcal B_i)$ have a neighbour $w$ in
  $G-F$ with $w\in\mathcal B_i\cap V_i$. Then
  $v\in\mathcal B_{i+1}$. Indeed, take a path from $U$ to $w$ as in the
  definition of $\mathcal B_i$. The vertex $v$ does not lie on it, since
  otherwise an initial segment would put $v$ into $\mathcal B_i$.
  Appending the edge $wv$ gives a path of length at most $i+1$ whose
  interior lies in $V_1\cup\dots\cup V_i$.
\item[(O2)] $\mathcal B_{\ell_*}\cap V_{\ell_*}\subseteq\Ball{\ell_*}{G-F}{U,V}$,
  since the interiors lie in $V$ and the end lies in $V_{\ell_*}\subseteq V$.
\item[(O3)] $|\mathcal B_1|\ge\theta_*u$. Deleting one edge removes at
  most one vertex from an outside neighbourhood, so
  $|\Nbr_{G-F}(U)|\ge|\Nbr_G(U)|-|F|\ge(\theta_*-1)u$.
\end{itemize}

\smallskip\noindent\emph{Step 2: growth into a fixed set.} Let
$1\le i\le\ell_*-1$, and let $W\subseteq V(G)$ be a fixed set with
$\mathcal B_1\subseteq W$ and $|W|\le 2n/3$. Put
\[
A_i(W):=\{v\in\Nbr_{G-F}(W):\ v\text{ has a neighbour in }G-F\text{ that lies in }W\cap V_i\}.
\]
We claim that
\begin{equation}\label{s3:eqL17claim}
\Prob\big(|A_i(W)|<g_*|W|\big)\le e^{-18uL}.
\end{equation}
By (O3), $|W|\ge\theta_*u$. Also $|F|\le u\le|W|\le s_1|W|/4$, because
$s_1\ge8d_*\lambda_*\ge48$. So Proposition~\ref{s3:propP13s} applies to
$W$ and $F$ with the values~\eqref{s3:eqStar}, and gives one of two
cases.

\emph{Case (a): stars.} Let $\mathrm{Cen}$ be the set of centres of at
least $|W|/\sigma_*$ vertex-disjoint stars as in (a). Every leaf of a
star whose centre lies in $V_i$ belongs to $A_i(W)$, and distinct stars
have distinct leaves. Now $|\mathrm{Cen}\cap V_i|\sim\Bin(|\mathrm{Cen}|,p_*)$,
so the Chernoff bound (Cited result~\ref{s1:citChernoff}, $\delta=1/2$) gives
$\Prob(|\mathrm{Cen}\cap V_i|<p_*|\mathrm{Cen}|/2)\le\exp(-p_*|\mathrm{Cen}|/8)$.
Outside this event, by~\ref{s3:eqS6},
\[
|A_i(W)|\ge\lambda_*p_*|\mathrm{Cen}|/2\ge3|\mathrm{Cen}|\ge3|W|/\sigma_*=g_*|W|.
\]
For the exponent, \ref{s3:eqS2} and $|W|\ge\theta_*u$ give
\begin{multline*}
\frac{p_*|\mathrm{Cen}|}{8}\ge\frac{p_*\varepsilon'|W|}{192L^2}\ge\frac{0.1\rho}{\ell_*}\cdot\frac{2^{19}\ell_*^2L^3u}{192\,\rho^2L^2}=\frac{0.1\cdot2^{19}\ell_*Lu}{192\,\rho}\\
\ge 273\,\ell_*uL\ \ (\text{using }\rho\le1)\ \ \ge 18uL.
\end{multline*}

\emph{Case (b): a bipartite graph.} Let $H\subseteq G-F$ be bipartite
with sides $W$ and $X$ as in (b). Every $x\in X$ with an $H$-neighbour in
$V_i$ lies in $A_i(W)$, because $x\notin W$ and $H\subseteq G-F$. Let
$Z:=|\{x\in X:\Nbr_H(x)\cap V_i\ne\emptyset\}|$. For $x\in X$,
by~\ref{s3:eqS6},
$\Prob(\Nbr_H(x)\cap V_i=\emptyset)=(1-p_*)^{d_*}\le e^{-p_*d_*}\le e^{-1}$,
so $\Ex Z\ge(1-e^{-1})|X|\ge0.63|X|$.

The count $Z$ is a function of the independent indicators
$\xi_w:=\ind{w\in V_i}$, $w\in W$, each Bernoulli with parameter $p_*$.
Order $W$ as $w_1,\dots,w_{|W|}$. For $y=(y_1,\dots,y_{|W|})\in\{0,1\}^{|W|}$
let $\Psi(y)$ be the number of $x\in X$ such that $y_k=1$ for some $k$ with
$w_k\in\Nbr_H(x)$. Since $H$ is bipartite with sides $W$ and $X$, we have
$\Nbr_H(x)\subseteq W$ for $x\in X$, so
$Z=\Psi(\xi_{w_1},\dots,\xi_{w_{|W|}})$. If $y,y'\in\{0,1\}^{|W|}$ differ
only in the $k$-th coordinate, then only the vertices $x\in\Nbr_H(w_k)$ can
change their contribution to $\Psi$, each by at most $1$, so
$|\Psi(y)-\Psi(y')|\le\deg_H(w_k)\le\Delta_*$ by
Proposition~\ref{s3:propP13s}(b). Using $1-p_*\le1$,
$\deg_H(w_k)\le\Delta_*$, $\sum_{w\in W}\deg_H(w)=e(H)=d_*|X|$ (every
$x\in X$ has exactly $d_*$ neighbours in $H$, all in $W$) and~\ref{s3:eqS6},
\[
\sum_k p_*(1-p_*)\deg_H(w_k)^2\le p_*\Delta_*\sum_{w\in W}\deg_H(w)=p_*d_*\Delta_*|X|\le 2\Delta_*|X|=:\beta ,
\]
and $\beta>0$, since $|X|\ge\varepsilon'|W|/(2L^2)>0$ (as $|W|\ge\theta_*u>0$)
and $\Delta_*\ge\lambda_*\ge1$.
Apply the second bound of Lemma~\ref{s1:lemBBD} with $m:=|W|$,
$I_k:=\xi_{w_k}$, $p_k:=p_*$, $c_k:=\deg_H(w_k)$, with $\Delta_*$ in the role
of the bound $b$ of that lemma, this $\beta$ and $a:=0.3|X|$. Since
$2(\beta+\Delta_*a/3)=4\Delta_*|X|+0.2\Delta_*|X|$ and
$0.09/4.2>1/47$, it gives
\[
\Prob\big(Z\le\Ex Z-0.3|X|\big)\le\exp\Big(-\frac{0.09|X|^2}{4\Delta_*|X|+0.2\Delta_*|X|}\Big)\le\exp\Big(-\frac{|X|}{47\Delta_*}\Big).
\]
Since $\Ex Z-0.3|X|\ge0.33|X|$, we get
$\Prob(Z<0.33|X|)\le\exp(-|X|/(47\Delta_*))$. Outside this event,
$|A_i(W)|\ge Z\ge0.33|X|\ge0.165\,\varepsilon'|W|/L^2\ge g_*|W|$.

For the exponent, $|X|\ge\varepsilon'|W|/(2L^2)$ and $|W|\ge\theta_*u$
give
\[
\frac{|X|}{47\Delta_*}\ge\frac{\varepsilon'\theta_*u}{94\Delta_*L^2}=\frac{2^{19}\ell_*^2Lu}{94\,\Delta_*\rho^2}.
\]
\begin{itemize}
\item If $p_*\le0.11$, then \ref{s3:eqS3} and~\ref{s3:eqS2} give
  $\Delta_*\le3p_*^{-2}\le300\,\ell_*^2\rho^{-2}$, and the exponent is at
  least $2^{19}uL/28200\ge18uL$ (indeed $2^{19}/28200>18.59$; this is the
  tightest inequality of the lemma).
\item If $p_*>0.11$, then \ref{s3:eqS3} gives
  $\Delta_*\le13p_*^{-2}<13/0.0121<1075$, and the exponent is at least
  $2^{19}\ell_*^2uL/(94\cdot1075)$ (using $\rho\le1$) $\ge5\ell_*^2uL\ge18uL$.
\end{itemize}
The large-$p_*$ regime occurs only for $\rho$ extremely close to $1$. It
is handled directly by the bound $\Delta_*\le13p_*^{-2}$, so no coupling
with a smaller density is needed. This proves~\eqref{s3:eqL17claim}.

\smallskip\noindent\emph{Step 3: all rounds.} For $1\le i\le\ell_*-1$ let
$E_i$ be the event that $|\mathcal B_i|\ge2n/3$ or
$|\mathcal B_{i+1}\setminus\mathcal B_i|\ge g_*|\mathcal B_i|$. Condition
on $V_1,\dots,V_{i-1}$. This fixes $W:=\mathcal B_i\supseteq\mathcal B_1$,
while $V_i$ is still $p_*$-random (by the independence of $V_i$ from
$(V_1,\dots,V_{i-1})$). If $|W|\ge2n/3$ then $E_i$ holds.
Otherwise $A_i(W)\subseteq\mathcal B_{i+1}\setminus\mathcal B_i$ by (O1),
since $A_i(W)\cap W=\emptyset$, and~\eqref{s3:eqL17claim} gives
$\Prob(E_i^c\mid V_1,\dots,V_{i-1})\le e^{-18uL}$. Hence, using $u\ge1$,
\ref{s3:eqS1}, and the fact that $L^3e^{-11L}\le e^{-11}$ for $L\ge1$,
\[
\Prob\Big(\bigcup_{i<\ell_*}E_i^c\Big)\le(\ell_*-1)e^{-18uL}\le2^{10}L^3e^{-11L}\,e^{-7uL}\le0.02\,e^{-7uL}.
\]

\smallskip\noindent\emph{Step 4: growth to two thirds.} Suppose all
$E_i$ hold, and suppose $|\mathcal B_{\ell_*}|<2n/3$. Then
$|\mathcal B_i|<2n/3$ for all $i\le\ell_*$, so
$|\mathcal B_{i+1}|\ge(1+g_*)|\mathcal B_i|$ for $1\le i\le\ell_*-1$, and
$|\mathcal B_{\ell_*}|\ge(1+g_*)^{\ell_*-1}$. Since $g_*\le1/8$, we have
$\ln(1+g_*)\ge g_*-g_*^2/2\ge15g_*/16$. By~\ref{s3:eqS1} and
$\varepsilon'\ge2^{-7}$,
$(\ell_*-1)g_*>(2^{10}L^3-2)\,2^{-10}L^{-2}\ge L-2^{-9}$. Hence
$\ln\big((1+g_*)^{\ell_*-1}\big)\ge\tfrac{15}{16}(L-2^{-9})\ge L\ln2=\ln n$,
so $|\mathcal B_{\ell_*}|\ge n$, a contradiction. Therefore
$|\mathcal B_{\ell_*}|\ge2n/3$.

\smallskip\noindent\emph{Step 5: the last batch.} Condition on
$V_1,\dots,V_{\ell_*-1}$ with $|\mathcal B_{\ell_*}|\ge2n/3$. The set
$V_{\ell_*}$ is still $0.9\rho$-random, so
$|\mathcal B_{\ell_*}\cap V_{\ell_*}|$ is binomial with mean at least
$0.6\rho n$. The lower-tail Chernoff bound for a mean bound $\mu:=0.6\rho n\le\Ex|\mathcal B_{\ell_*}\cap V_{\ell_*}|$
(Cited result~\ref{s1:citChernoffGen}(a)) with $\delta=0.09$ gives
$\Prob(|\mathcal B_{\ell_*}\cap V_{\ell_*}|<0.546\rho n)\le
\exp(-0.0081\cdot0.6\rho n/2)\le\exp(-0.00243\rho n)$. Since $V$ is
$\rho$-random, $|V|\sim\Bin(n,\rho)$, and
$\Prob(|V|>1.09\rho n)\le\exp(-0.0081\rho n/3)=\exp(-0.0027\rho n)$.
Outside both events,
$|\mathcal B_{\ell_*}\cap V_{\ell_*}|\ge0.546\rho n>0.545\rho n\ge|V|/2$,
and then $|\Ball{\ell_*}{G-F}{U,V}|>|V|/2$ by (O2). The probability of
the two events is at most
$2e^{-\rho n/412}\le 2e^{-\rho n/413}e^{-7uL}\le0.01\,e^{-7uL}$. Here we
used Step~0: $7uL\le2^{-36}\rho n\le\rho n(1/412-1/413)$ and
$\rho n>2^{39}$.

\smallskip\noindent\emph{Conclusion.} If
$|\Ball{\ell_*}{G-F}{U,V}|\le|V|/2$, then either some $E_i$ fails, or
(by Step~4) $|\mathcal B_{\ell_*}|\ge2n/3$ and one of the two events of
Step~5 occurs. The total probability is at most
$0.03\,e^{-7uL}\le e^{-7uL}$. The final sentence of the lemma holds
because the event is a function of the random set $V$ alone, and $V$ is
$\rho$-random.
\end{proof}

\begin{lemma}[Proposition 18$^*$ and Lemma 19$^*$]\label{s3:lemP18s}
Let $G$ be an $n$-vertex $(\varepsilon',s_1)$-expander with $n\ge2$ and
$2^{-7}\le\varepsilon'\le1$, let $\rho\in(0,1]$, take the parameters
of~\eqref{s3:eqStar} with $L:=\log n$, and assume
$s_1\ge\theta_*+1$.
\begin{itemize}
\item[(i)] (Proposition 18$^*$.) Every $U\subseteq V(G)$ with
  $1\le|U|\le2n/3$ contains a set $U'$ with
  $|\Nbr_G(U')|\ge\theta_*|U'|$ and
  $|U'|\ge\varepsilon'|U|/(3\theta_*L^2)$.
\item[(ii)] (Lemma 19$^*$.) Let $V$ be a $\rho$-random subset of
  $V(G)$. With probability at least $1-n^{-6}$ the following holds: for
  every nonempty $U\subseteq V(G)$ and every $F\subseteq E(G)$ with
  $|F|\le\mu_*|U|$,
  \[
  |\Ball{\ell_*}{G-F}{U,V}|>|V|/2 .
  \]
\end{itemize}
The hypothesis is $s_1\ge\theta_*+1$ and not $s_1\ge\theta_*$. The
maximality step of \cite[Proposition~18]{BM} (Cited result~\ref{s1:citProp18})
produces vertices with \emph{fewer than $\theta_*+1$} neighbours outside,
so the edge set $F$ built there satisfies $|F|<(\theta_*+1)|U|$.
\deps{\ref{s1:citProp18}, \ref{s1:citLem19}, \ref{s3:lemL17s}, \ref{s1:citDef11}}
\end{lemma}

\begin{proof}
By~\ref{s3:eqS4}, $\theta_*\ge 8d_*\lambda_*$, so $s_1\ge 8d_*\lambda_*$
and Lemma~\ref{s3:lemL17s} applies to $G$. Also $\theta_*\ge1$.

(i) We follow the proof of \cite[Proposition~18]{BM}. Let $U'\subseteq U$
be maximal subject to $|\Nbr_G(U')|\ge\theta_*|U'|$; the empty set
qualifies. If $U'=U$ we are done, since
$\varepsilon'/(3\theta_*L^2)\le1$. Otherwise let $v\in U\setminus U'$, and
let $a_v$ be the number of neighbours of $v$ outside
$U'\cup\Nbr_G(U')$. Then
\[
|\Nbr_G(U'\cup\{v\})|=|\Nbr_G(U')|-\ind{v\in\Nbr_G(U')}+a_v .
\]
By maximality this is less than $\theta_*(|U'|+1)$. Two consequences
follow:
\begin{itemize}
\item $a_v<\theta_*(|U'|+1)-|\Nbr_G(U')|+1\le\theta_*+1$;
\item $|\Nbr_G(U')|<\theta_*(|U'|+1)+1$.
\end{itemize}
Let $F$ be the set of edges $vx$ with $v\in U\setminus U'$ and
$x\notin U'\cup\Nbr_G(U')$. Then
$|F|\le\sum_{v\in U\setminus U'}a_v<(\theta_*+1)|U|\le s_1|U|$.
Every vertex $x\notin U$ with a neighbour in $U$ in $G-F$ lies in
$\Nbr_G(U')$: if the neighbour is in $U'$ this is clear, and if it is in
$U\setminus U'$ then, since the edge is not in $F$,
$x\in U'\cup\Nbr_G(U')$ with $x\notin U'$. By the expansion of $G$
(Cited result~\ref{s1:citDef11}),
\[
\varepsilon'|U|/L^2\le|\Nbr_{G-F}(U)|\le|\Nbr_G(U')|<\theta_*(|U'|+1)+1 .
\]
If $U'=\emptyset$, then the middle term is $0$, which is impossible
because $|U|\ge1$. So $|U'|\ge1$, and then
$\theta_*(|U'|+1)+1\le3\theta_*|U'|$ since $\theta_*\ge1$. Hence
$|U'|\ge\varepsilon'|U|/(3\theta_*L^2)$.

(ii) We follow the proof of \cite[Lemma~19]{BM}
(Cited result~\ref{s1:citLem19}).

\emph{The union bound.} Consider all pairs $(U',F')$ with $U'$
well-expanding and nonempty, $F'\subseteq E(G)$ and $|F'|\le|U'|$. Put
$e_G:=|E(G)|\le n^2/2$; since $G$ has minimum degree $>s_1\ge1$,
$e_G\ge2$. The number of pairs with $|U'|=u$ is at most
$n^u\sum_{f=0}^{u}e_G^f\le n^u\cdot2e_G^u\le n^{3u}$. By
Lemma~\ref{s3:lemL17s}, each pair fails
$|\Ball{\ell_*}{G-F'}{U',V}|>|V|/2$ with probability at most
$e^{-7uL}=n^{-7u\log e}\le n^{-10.09u}$. Hence, with probability at least
\[
1-\sum_{u\ge1}n^{3u}n^{-10.09u}\ \ge\ 1-\frac{n^{-7.09}}{1-n^{-7.09}}\ \ge\ 1-n^{-6},
\]
every such pair satisfies the ball bound. Fix an outcome in this event.

\emph{Deterministic consequence.} Let $U$ be nonempty, and let
$F\subseteq E(G)$.
\begin{itemize}
\item \emph{Case $|U|\le2n/3$ and $|F|\le2\mu_*|U|$.} By (i) there is a
  well-expanding $U'\subseteq U$ with
  $|U'|\ge\varepsilon'|U|/(3\theta_*L^2)=2\mu_*|U|\ge|F|$. The ball
  $\Ball{\ell_*}{G-F}{\cdot,V}$ is monotone in its first argument, so
  $|\Ball{\ell_*}{G-F}{U,V}|\ge|\Ball{\ell_*}{G-F}{U',V}|>|V|/2$.
\item \emph{Case $|U|>2n/3$ and $|F|\le\mu_*|U|$.} Choose
  $\bar U\subseteq U$ with $n/2\le|\bar U|\le2n/3$. This is possible
  for every $n\ge2$: the interval $[n/2,2n/3]$ contains an integer, as
  one checks directly for $n\le5$, and for $n\ge6$ it has length at
  least $1$. Then $|F|\le\mu_*n\le2\mu_*|\bar U|$, and the previous case
  gives $|\Ball{\ell_*}{G-F}{U,V}|\ge|\Ball{\ell_*}{G-F}{\bar U,V}|>|V|/2$.
\end{itemize}
In both cases the conclusion of (ii) holds.
\end{proof}

\begin{lemma}[Lemma 9$_\rho$: multiset linking from ball expansion]\label{s3:lemL9rho}
Let $G$ be an $n$-vertex graph with $n\ge2^{30}$, put $L:=\log n$, and
let $t\ge1$ and $\rho\in(0,1]$. Take $\ell_*$, $\bar s_*$ and
$M_*=\lceil2.1/\rho\rceil$ from~\eqref{s3:eqStar}. Let $V\subseteq V(G)$
with $|V|\ge\rho n/2$, and suppose $\rho n\ge84$. Then
$|V|\ge n/M_*+2$. Suppose that
\[
|\Ball{\ell_*}{G-F}{U,V}|>|V|/2
\]
for every nonempty $U\subseteq V(G)$ and every $F\subseteq E(G)$ with
$|F|\le\bar s_*|U|$. Then $G$ is $(2^{12}L^4,t)$-path connected through
$V$ in the multiset sense of \cite[Definition~7]{BM}
(Cited result~\ref{s1:citDef7}). That is, every multiset of pairs of distinct
vertices of $G$ in which each vertex lies in at most $t$ pairs can be
joined by pairwise edge-disjoint paths of length at most $2^{12}L^4$
whose interior vertices lie in $V$.
\deps{\ref{s1:citLem9}, \ref{s1:citProp8}, \ref{s1:citHaxell}, \ref{s1:citDef7}, \ref{s3:remMultiset}}
\end{lemma}

\begin{proof}
First, $n/M_*\le\rho n/2.1$ and
$\rho n/2-\rho n/2.1=\rho n/42\ge2$, so $|V|\ge n/M_*+2$. Next,
$1\le\ell_*\le2^{10}L^3\le n$ for $n\ge2^{30}$. We re-run the proof of
\cite[Lemma~9]{BM} (Cited result~\ref{s1:citLem9}) with its constants changed,
and with Haxell's theorem (Cited result~\ref{s1:citHaxell}) in place of the
Aharoni--Haxell theorem \cite[Theorem~6]{BM}.
It uses \cite[Proposition~8]{BM} (Cited result~\ref{s1:citProp8}): if
$1\le\ell,t_0\le n$, $|V|\ge4t_0-2$, and every set of $t_0$ vertices has
a ball of radius $\ell$ in $V$ of size greater than $|V|/2$, then among
any $2t_0-1$ pairs of vertices, all $4t_0-2$ of them distinct, some pair
is joined by a path through $V$ of length at most $4\ell\log n$.

Let $\mathcal P=(\{x_i,y_i\})_{i\in[r]}$ be a multiset of pairs of
distinct vertices, indexed by $[r]$, in which every vertex lies in at
most $t$ pairs (counted with multiplicity; see
Remark~\ref{s3:remMultiset}). Put $h:=\lfloor4\ell_*L\rfloor$; then
$h\ge1$, since $\ell_*\ge1$ and $L\ge1$. For $i\in[r]$, let $\mathcal Y_i$
be the set of $x_iy_i$-paths in $G$ that have length at most $h$ and all
interior vertices in $V$. Every $P\in\mathcal Y_i$ has at most $h$ edges.

\emph{Claim: for every nonempty $I\subseteq[r]$ and every $F\subseteq E(G)$
with $|F|<h^2|I|$, there are $j\in I$ and $P\in\mathcal Y_j$ with
$E(P)\cap F=\emptyset$.}
Let $I'\subseteq I$ be maximal such that the pairs $\{x_i,y_i\}$,
$i\in I'$, are pairwise vertex-disjoint; $I'\ne\emptyset$ because
$I\ne\emptyset$. By maximality every pair
indexed by $I$ contains one of the $2|I'|$ vertices covered by $I'$, and
each vertex lies in at most $t$ pairs. Hence $|I|\le2t|I'|$. Also
$|I'|\le n/2$. Put $t_{0*}:=\lceil|I'|/(2M_*)\rceil$. Then:
\begin{itemize}
\item $1\le t_{0*}$ and $2t_{0*}-1\le|I'|$ (if $|I'|\le2M_*$ then
  $t_{0*}=1$; otherwise $2t_{0*}-1\le|I'|/M_*+1\le|I'|$, as $M_*\ge3$);
\item $4t_{0*}-2\le n/M_*+2\le|V|$;
\item $t_{0*}\le n$;
\item $|F|<h^2|I|\le16(\ell_*L)^2\cdot2t|I'|\le64M_*t\,t_{0*}(\ell_*L)^2$.
\end{itemize}
Since $M_*\le2.1/\rho+1\le3.1/\rho$ and $(\ell_*L)^2\le2^{20}L^8$, we get
$64M_*t(\ell_*L)^2\le2^{28}tL^8/\rho=\bar s_*$. Hence $|F|<\bar s_*t_{0*}$.
So every set $U$ of $t_{0*}$ vertices satisfies $|F|\le\bar s_*|U|$, and
therefore $|\Ball{\ell_*}{G-F}{U,V}|>|V|/2$ by hypothesis.

Choose $2t_{0*}-1$ indices of $I'$; their $4t_{0*}-2$ endpoints are
distinct. Apply Cited result~\ref{s1:citProp8} to $G-F$, with $\ell:=\ell_*$
and $t_0:=t_{0*}$. It gives, for one of these indices $j$, an
$x_jy_j$-path $P$ in $G-F$ through $V$ of length at most $4\ell_*L$, hence
at most $h$, as lengths are integers. So $j\in I'\subseteq I$,
$P\in\mathcal Y_j$ and $E(P)\cap F=\emptyset$. This proves the claim.

Now let $\mathcal H$ be the hypergraph with vertex set $[r]\sqcup E(G)$ (a
disjoint union) whose edges are the sets $\{i\}\cup E(P)$ with $i\in[r]$
and $P\in\mathcal Y_i$. Every edge of $\mathcal H$ meets $[r]$ in exactly
one vertex and $E(G)$ in at most $h$ vertices. We apply Cited
result~\ref{s1:citHaxell} to $\mathcal H$, with $[r]$ in the role of its
class $\mathcal X$, $E(G)$ in the role of its other class, and with $h\ge1$ in
place of its parameter $q$. We check its hypothesis. Let $I\subseteq[r]$ and $Z\subseteq E(G)$
with $|Z|\le(2h-1)(|I|-1)$. Then $I\ne\emptyset$, since otherwise
$|Z|\le1-2h<0$. As $h^2-(2h-1)=(h-1)^2\ge0$ and $|I|-1\ge0$,
\[
|Z|\le(2h-1)(|I|-1)\le h^2(|I|-1)<h^2|I| .
\]
So the claim, with $F:=Z$, gives $j\in I$ and $P\in\mathcal Y_j$ with
$E(P)\cap Z=\emptyset$. The edge $e:=\{j\}\cup E(P)$ of $\mathcal H$
satisfies $e\cap[r]=\{j\}\subseteq I$ and, as $j\notin E(G)\supseteq Z$,
$e\cap Z=E(P)\cap Z=\emptyset$. (The claim covers every $Z$ with
$|Z|<h^2|I|$, which is more than the hypothesis of Cited
result~\ref{s1:citHaxell} requires.) By Cited result~\ref{s1:citHaxell},
$\mathcal H$ has a set of pairwise disjoint edges such that every
$i\in[r]$ lies in one of them. For $i\in[r]$ let $e_i$ be an edge of this
set containing $i$. By construction $e_i=\{k\}\cup E(Q)$ for some
$k\in[r]$ and $Q\in\mathcal Y_k$, so $e_i\cap[r]=\{k\}$; as
$i\in e_i\cap[r]$, we get $k=i$, and we put $P_i:=Q$, so that
$E(P_i)\subseteq e_i$. For $i\ne i'$ we have
$e_i\cap[r]=\{i\}\ne\{i'\}=e_{i'}\cap[r]$, so $e_i\ne e_{i'}$, and
therefore $e_i\cap e_{i'}=\emptyset$, as both lie in the set of pairwise
disjoint edges. Hence the paths $P_i$, $i\in[r]$, are pairwise
edge-disjoint $x_iy_i$-paths through $V$ of length at most
$h\le4\ell_*L\le2^{12}L^4$. Repeated pairs receive distinct paths,
because the paths are indexed by $i$.
\end{proof}

\begin{theorem}[Theorem 16$^*$: linking through random sets; for-all multiset form]\label{s3:thmT16s}
Let $X$ be an $N$-vertex $(\varepsilon',s)$-expander with
$\varepsilon'\in[2^{-7},1]$, and put $L:=\log N$. Let $\rho\in(0,1]$ and
$t\ge1$, and let $V$ be a $\rho$-random subset of $V(X)$. Here $X$ is
fixed; if $X$ was itself produced at random, $V$ is independent of that
randomness. Suppose that
\[
\rho N\ge L^2\qquad\text{and}\qquad s\ge 2^{135}\,t\,L^{28}\rho^{-5}.
\]
Then, with probability at least $1-2^{86}\,t\,L^{19}\rho^{-3}N^{-3}$, the
graph $X$ is $(2^{12}L^4,t)$-path connected through $V$. That is, for
\emph{every} multiset $\mathcal P$ of pairs of distinct vertices of $X$
(endpoints anywhere, in $V$ or not) in which each vertex lies in at most
$t$ pairs, there are pairwise edge-disjoint paths of $X$, one for each
pair of $\mathcal P$, each joining its pair, of length at most
$2^{12}L^4$, and with all interior vertices in $V$. By
Lemma~\ref{s3:lemMonotone}(iii), $\mathcal P$ may be chosen after $V$ is
revealed. Moreover:
\begin{itemize}
\item[(a)] The constant $2^{86}$ in the failure probability is absolute.
  No lower bound on $N$ is needed.
\item[(b)] If $N\ge2$, the hypotheses force
  $\rho N>2^{27}L^{5.6}N^{4/5}$, because $s<N$. In particular the
  hypothesis $\rho N\ge L^2$ is implied by the others; it is kept only
  for comparison with \cite[Theorem~16]{BM}. The last batch of the
  sprinkling argument needs $\rho N$ large compared with $L$ times the
  size of the well-expanding sets involved. This is automatic (Step~0 of
  the proof of Lemma~\ref{s3:lemL17s}), so the final-stage union bound
  needs no further hypothesis.
\item[(c)] The requirement on $s$ is exactly linear in $t$ and uniform in
  $\rho\in(0,1]$. Apart from Step~0 of the proof below, which uses the full hypothesis on $s$ to
  obtain (b), $N\ge2^{30}$ and $\rho N\ge84$ (from (b), Step~5 then draws
  $e^{-\rho N/8}\le N^{-3}$),
  the proof uses only $s\ge2K_*(\theta_*+1)$, and
  $2K_*(\theta_*+1)\le2^{130.6}\,tL^{28}\rho^{-5}$ by~\ref{s3:eqS5}. No
  coupling with a smaller density is used, also when $\rho$ is close to
  $1$.
\end{itemize}
\deps{\ref{s3:lemL15p}, \ref{s3:propP13s}, \ref{s3:lemL17s}, \ref{s3:lemP18s}, \ref{s3:lemL9rho}, \ref{s3:lemMonotone}, \ref{s1:citDef11}, \ref{s1:citChernoff}}
\end{theorem}

\begin{proof}
The proof is an assembly, following the proof of \cite[Theorem~16]{BM}
(Cited result~\ref{s1:citThm16}). The loss-free splitting
Lemma~\ref{s3:lemL15p} takes the place of \cite[Lemma~15]{BM}.

\smallskip\noindent\emph{Step 0: $N$ and $\rho N$ are large.}
If $N=1$ there are no pairs of distinct vertices, and there is nothing
to prove. Let $N\ge2$, so $L\ge1$. Taking $U=\{v\}$ in
Cited result~\ref{s1:citDef11} shows that $X$ has minimum degree greater than
$s$. Hence $N>s\ge2^{135}tL^{28}\rho^{-5}\ge2^{135}$. Then
$\rho^5N>\rho^5s\ge2^{135}L^{28}$, so $\rho>2^{27}L^{5.6}N^{-1/5}$, which
is (b). In particular $\rho N\ge84$ and $N\ge2^{30}$. Take the
parameters~\eqref{s3:eqStar} with $n:=N$ and the given $\varepsilon'$,
$\rho$ and $t$.

\smallskip\noindent\emph{Step 1: splitting.} Colour $E(X)$ with $K_*$
colours, uniformly and independently of $V$, and let
$X_1,\dots,X_{K_*}$ be the colour classes, each with vertex set $V(X)$.
By~\ref{s3:eqS5}, $s\ge2K_*(\theta_*+1)\ge40K_*L$. So by
Lemma~\ref{s3:lemL15p}, with probability at least $1-2K_*N^{-5}$, every
$X_i$ is an $(\varepsilon',s/(2K_*))$-expander, and
$s/(2K_*)\ge\theta_*+1$.

\smallskip\noindent\emph{Step 2: balls in each class.} Condition on a
colouring in which all classes are such expanders; $V$ is still
$\rho$-random. Each $X_i$ has $N$ vertices and parameter $\varepsilon'$,
so the parameters~\eqref{s3:eqStar} for $X_i$ are those for $X$. By
Lemma~\ref{s3:lemP18s}(ii) with $s_1:=s/(2K_*)$ and a union bound over
$i$, with probability at least $1-K_*N^{-6}$ the following holds for
every $i$: for every nonempty $U$ and every $F_i\subseteq E(X_i)$ with
$|F_i|\le\mu_*|U|$, we have $|\Ball{\ell_*}{X_i-F_i}{U,V}|>|V|/2$.

\smallskip\noindent\emph{Step 3: the size of $V$.} Since
$|V|\sim\Bin(N,\rho)$, the Chernoff bound (Cited result~\ref{s1:citChernoff})
gives $\Prob(|V|<\rho N/2)\le e^{-\rho N/8}$.

\smallskip\noindent\emph{Step 4: assembly.} Assume the events of
Steps~1--3. Let $U$ be nonempty and let $F\subseteq E(X)$ with
$|F|\le\bar s_*|U|$. The classes partition $E(X)$, so some $i$ has
$|F\cap E(X_i)|\le|F|/K_*\le\bar s_*|U|/K_*\le\mu_*|U|$
by~\ref{s3:eqS5}. Since $X_i-(F\cap E(X_i))\subseteq X-F$, balls can only
grow when passing to $X-F$, and
\[
|\Ball{\ell_*}{X-F}{U,V}|\ge|\Ball{\ell_*}{X_i-(F\cap E(X_i))}{U,V}|>|V|/2 .
\]
So the hypotheses of Lemma~\ref{s3:lemL9rho} hold for $G:=X$:
$N\ge2^{30}$, $|V|\ge\rho N/2$ and $\rho N\ge84$. Hence $X$ is
$(2^{12}L^4,t)$-path connected through $V$, in the multiset sense.

\smallskip\noindent\emph{Step 5: probability.} The failure probability
is at most $2K_*N^{-5}+K_*N^{-6}+e^{-\rho N/8}$. By (b),
$\rho N/8>2^{24}N^{4/5}\ge3\ln N$, so $e^{-\rho N/8}\le N^{-3}$. By
Lemma~\ref{s3:lemMonotone}(iii), the random colouring of Step~1, which is
auxiliary, does not affect the event, which depends only on $(X,V)$. So
the event fails with probability at most
\[
3K_*N^{-5}+N^{-3}\le\big(3\cdot2^{83.6}+1\big)tL^{19}\rho^{-3}N^{-3}\le2^{86}\,tL^{19}\rho^{-3}N^{-3},
\]
using~\ref{s3:eqS5} and $tL^{19}\rho^{-3}\ge1$. This proves the main
statement and~(a).

For (c): apart from Step~0, whose consequences $N\ge2^{30}$ and $\rho N\ge84$
(used in Step~4) and $e^{-\rho N/8}\le N^{-3}$ (used in Step~5, through (b)) rest
on the full hypothesis $s\ge2^{135}tL^{28}\rho^{-5}$, the requirements on $s$
used above are $s\ge40K_*L$ (Step~1) and $s/(2K_*)\ge\theta_*+1$ (Step~2). Both follow from
$s\ge2K_*(\theta_*+1)$, and $2K_*(\theta_*+1)\le2^{130.6}tL^{28}\rho^{-5}$
by~\ref{s3:eqS5}. Only $\bar s_*$, and hence $K_*$, depends on $t$, and
linearly up to rounding. The parameters~\eqref{s3:eqStar} are defined for
every $\rho\in(0,1]$. The regime $p_*>0.11$, reached only for $\rho$ very
close to $1$, is treated directly in Lemma~\ref{s3:lemL17s} through
$\Delta_*\le13p_*^{-2}$ (fact~\ref{s3:eqS3}).
\end{proof}

%% ----------------------------------------------------------------------
\subsection{Candidate sets of bounded multiplicity}

\begin{lemma}[Lemma HB: candidate sets of bounded multiplicity; Hall argument]\label{s3:lemHB}
Let $X$ be an $N$-vertex $(\varepsilon',s)$-expander with $N\ge2$ and
$2^{-7}\le\varepsilon'\le1$, and put $L:=\log N$. Let $m$ be an integer
with $1\le m$ and $2m\le s$, and put $b:=\lceil2mL^2/\varepsilon'\rceil$.
Thus $b=\lceil64L^2m\rceil$, $\lceil128L^2m\rceil$ or
$\lceil256L^2m\rceil$ at $\varepsilon'=2^{-5}$, $2^{-6}$ or $2^{-7}$.
Then every vertex $w$ has a set $\Aw(w)\subseteq\Nbr_X(w)$ with
$|\Aw(w)|=m$, such that every vertex of $X$ lies in at most $b$ of the
sets $\Aw(w)$, $w\in V(X)$.
\deps{\ref{s1:citHall}, \ref{s1:citDef11}}
\end{lemma}

\begin{proof}
For $w\in V(X)$, its degree is $d_X(w)=|\Nbr_X(w)|$. As $N\ge2$, we have
$L\ge1$ and $d_X(w)\ge\delta(X)>s\ge2m$ for every $w$ (Cited
result~\ref{s1:citDef11}); in particular $d_X(w)-m>m\ge1$ (not needed
below). Let
$\vec E(X):=\{(w,u)\in V(X)\times V(X):\ wu\in E(X)\}$ be the set of
ordered pairs along edges (every edge gives two), and for $w\in V(X)$ let
$\vec E_X(w):=\{(w,u):\ u\in\Nbr_X(w)\}$. Then $|\vec E_X(w)|=d_X(w)$, and
distinct elements of $\vec E_X(w)$ have distinct second coordinates.

\emph{The auxiliary bipartite graph.} Its left side is $\vec E(X)$. Its
right side is the set of \emph{slots}
\[
\mathcal B:=\{(u,j,1):\ u\in V(X),\ j\in[b]\}\ \cup\
\{(w,i,2):\ w\in V(X),\ i\in[d_X(w)-m]\} ;
\]
$(u,j,1)$ is a \emph{load slot} of $u$, and $(w,i,2)$ is a \emph{spare
slot} of $w$. A pair $a=(w,u)\in\vec E(X)$ is joined to the $b$ load slots
of its second coordinate $u$ and to the $d_X(w)-m$ spare slots of its first
coordinate $w$:
\[
\mathcal T(a):=\{(u,j,1):\ j\in[b]\}\ \cup\ \{(w,i,2):\ i\in[d_X(w)-m]\}.
\]

\emph{Reduction to Hall's condition.} Suppose that
$\chi:\vec E(X)\to\mathcal B$ is injective and $\chi(a)\in\mathcal T(a)$ for
every $a\in\vec E(X)$. Fix $w\in V(X)$. If $a\in\vec E_X(w)$ and $\chi(a)$ is
a spare slot, then $\chi(a)$ is one of the $d_X(w)-m$ spare slots of $w$; as
$\chi$ is injective, at most $d_X(w)-m$ elements of $\vec E_X(w)$ are of
this kind. So at least $d_X(w)-(d_X(w)-m)=m$ elements
$(w,u)\in\vec E_X(w)$ have $\chi(w,u)$ a load slot, and then it is a load
slot of $u$. Let $\Aw(w)$ be the set of second coordinates of $m$ such
elements (the first $m$ in a fixed ordering of $V(X)$). Then
$\Aw(w)\subseteq\Nbr_X(w)$, and $|\Aw(w)|=m$ because distinct elements of
$\vec E_X(w)$ have distinct second coordinates. Now fix $u\in V(X)$. If
$u\in\Aw(w)$, then $\chi(w,u)$ is one of the $b$ load slots of $u$. The
pairs $(w,u)$ with distinct $w$ are distinct and $\chi$ is injective, so
$u\in\Aw(w)$ holds for at most $b$ vertices $w$. This is the statement. By
Hall's theorem (Cited result~\ref{s1:citHall}, applied to the finite index
set $\vec E(X)$ and the finite sets $\mathcal T(a)$), such a $\chi$ exists
provided
\begin{equation}\label{s3:eqHBhall}
|\mathcal E|\ \le\ \Bigl|\bigcup_{a\in\mathcal E}\mathcal T(a)\Bigr|\qquad\text{for every }\mathcal E\subseteq\vec E(X).
\end{equation}

\emph{Proof of~\eqref{s3:eqHBhall}.} Fix $\mathcal E\subseteq\vec E(X)$.
Let $\mathcal S$ be the set of first coordinates and $\mathcal R$ the set of
second coordinates of the elements of $\mathcal E$. For $w\in\mathcal S$ put
$\mathcal E_w:=\mathcal E\cap\vec E_X(w)$, so that
$1\le|\mathcal E_w|\le d_X(w)$ and
$|\mathcal E|=\sum_{w\in\mathcal S}|\mathcal E_w|$. The union
in~\eqref{s3:eqHBhall} consists of the $b|\mathcal R|$ load slots of the
vertices of $\mathcal R$ and the $\sum_{w\in\mathcal S}(d_X(w)-m)$ spare
slots of the vertices of $\mathcal S$, and all these slots are distinct.
So \eqref{s3:eqHBhall} is equivalent to
\begin{equation}\label{s3:eqHBdef}
\sum_{w\in\mathcal S}\bigl(|\mathcal E_w|-d_X(w)+m\bigr)\ \le\ b|\mathcal R| .
\end{equation}
Let $\mathcal S':=\{w\in\mathcal S:\ d_X(w)-|\mathcal E_w|\le m-1\}$. A
summand of~\eqref{s3:eqHBdef} with $w\in\mathcal S\setminus\mathcal S'$ is
at most $0$, and one with $w\in\mathcal S'$ is at most $m$, as
$|\mathcal E_w|\le d_X(w)$. So the left side of~\eqref{s3:eqHBdef} is at
most $m|\mathcal S'|$, and \eqref{s3:eqHBdef} holds if
$\mathcal S'=\emptyset$ (this includes $\mathcal E=\emptyset$). Assume
$\mathcal S'\ne\emptyset$.

Let $w\in\mathcal S'$ and $u\in\Nbr_X(w)\setminus\mathcal R$. Then
$(w,u)\in\vec E_X(w)\setminus\mathcal E_w$, since otherwise
$u\in\mathcal R$; distinct such $u$ give distinct pairs $(w,u)$. As
$|\vec E_X(w)\setminus\mathcal E_w|=d_X(w)-|\mathcal E_w|\le m-1$, we get
\begin{equation}\label{s3:eqHBout}
|\Nbr_X(w)\setminus\mathcal R|\ \le\ m-1\qquad\text{for every }w\in\mathcal S'.
\end{equation}
If $|\mathcal S'|\le2N/3$, put $\mathcal U:=\mathcal S'$. Otherwise let
$\mathcal U\subseteq\mathcal S'$ be any set with
$|\mathcal U|=\lceil|\mathcal S'|/2\rceil$; then
$1\le|\mathcal U|\le\lceil N/2\rceil\le2N/3$, because $N\ge2$ (for $N=2$,
$\lceil N/2\rceil=1\le4/3$; for $N\ge3$, $\lceil N/2\rceil\le(N+1)/2\le2N/3$).
In both cases
\[
1\le|\mathcal U|\le2N/3\qquad\text{and}\qquad|\mathcal U|\ge|\mathcal S'|/2 .
\]
Let $F$ be the set of edges $wu\in E(X)$ with $w\in\mathcal U$ and
$u\notin\mathcal R$. Every edge of $F$ is of the form $wu$ with
$w\in\mathcal U$ and $u\in\Nbr_X(w)\setminus\mathcal R$, so,
by~\eqref{s3:eqHBout},
$|F|\le\sum_{w\in\mathcal U}|\Nbr_X(w)\setminus\mathcal R|\le(m-1)|\mathcal U|\le s|\mathcal U|$,
since $m-1<2m\le s$. Let $v\in\Nbr_{X-F}(\mathcal U)$: then
$v\notin\mathcal U$ and $vw\in E(X)\setminus F$ for some $w\in\mathcal U$.
As $w\in\mathcal U$ and $vw\notin F$, we have $v\in\mathcal R$. Hence
$\Nbr_{X-F}(\mathcal U)\subseteq\mathcal R$, and the expansion of $X$
(Cited result~\ref{s1:citDef11} with $n=N$, applied to $\mathcal U$ and
$F$) gives
\[
|\mathcal R|\ \ge\ |\Nbr_{X-F}(\mathcal U)|\ \ge\ \frac{\varepsilon'|\mathcal U|}{L^2}\ \ge\ \frac{\varepsilon'|\mathcal S'|}{2L^2}.
\]
Since $b\ge2mL^2/\varepsilon'$, it follows that
\[
b|\mathcal R|\ \ge\ \frac{2mL^2}{\varepsilon'}\cdot\frac{\varepsilon'|\mathcal S'|}{2L^2}\ =\ m|\mathcal S'| ,
\]
which is at least the left side of~\eqref{s3:eqHBdef}. This
proves~\eqref{s3:eqHBhall}, and hence the lemma.
\end{proof}

%% ----------------------------------------------------------------------
\subsection{The lending colouring}

Fix a valid $\HBtp$ run on $G$ (Definition~\ref{s2:defHBtp}). Every
ancestor $Y$ (Definition~\ref{s2:defAncestors}) splits the edges of its
expander $H_Y$ into two halves:
\begin{itemize}
\item an \emph{own} half, which only $Y$'s own vortex device uses;
\item a \emph{lend} half, which is coloured into exclusive classes, each
  lent to exactly one later consumer.
\end{itemize}
The random data of this subsection form part of stage~1 of the
randomness. The other stage-1 data, introduced in later sections, are
drawn independently of them.

\begin{definition}[stage-1 lending data: the two-stage colouring COL-JV and the JS labels]\label{s3:defCOL}
For every ancestor $Y$ of round $r$, with $V(Y)$, $H_Y$,
$\varepsilon_Y$, $s_Y$ and $L_Y=\log|V(Y)|$ as in
Definition~\ref{s2:defAncestors}, the random data (i)--(iv) below are
drawn. The data of distinct ancestors are independent.
\begin{itemize}
\item[(i)] \emph{Own/Lend split.} Every edge of $H_Y$ independently lies
  in $\Own_Y$ or in $\Lend_Y$, with probability $1/2$ each.
\item[(ii)] \emph{Lent classes.} Consider the three index families
  \[
  \begin{aligned}
  \IU(Y)&:=\{(l,c,\uslot):\ r+2\le l\le R,\ c\in[4],\ \uslot\in[\Tslot_Y]\}\ \text{ if $Y$ is light, where } \Tslot_Y:=\lceil L_Y^2\rceil,\\
  \IU(Y)&:=\emptyset\ \text{ if $Y$ is standalone,}\\
  \IJS(Y)&:=\{(l,j):\ r+2\le l\le R,\ 0\le j<\KJS_l\},\ \text{ where } \KJS_l:=M_l^2,\\
  \IJV(Y)&:=\{l:\ r+2\le l\le R\}.
  \end{aligned}
  \]
  These families are regarded as pairwise disjoint (their indices carry
  the tags U, JS, JV). Put
  $\klend(Y):=|\IU(Y)|+|\IJS(Y)|+|\IJV(Y)|$.
  \begin{itemize}
  \item If $r\le R-2$, every edge of $\Lend_Y$ independently receives an
    index chosen uniformly from $\IU(Y)\sqcup\IJS(Y)\sqcup\IJV(Y)$. The
    edges with index $(l,c,\uslot)\in\IU(Y)$ form the \emph{U-lent class}
    $\LU_{Y,l,c,\uslot}$. The edges with index $(l,j)\in\IJS(Y)$ form the
    \emph{JS-lent class} $\LJS_{Y,l,j}$. The edges with index
    $l\in\IJV(Y)$ form the \emph{JV-lent class} $\LJV_{Y,l}$. Every
    class is regarded as a graph with vertex set $V(Y)$.
  \item If $r\ge R-1$, all three families are empty, $\klend(Y)=0$, and
    $\Lend_Y$ has no classes.
  \end{itemize}
\item[(iii)] \emph{Own classes (light $Y$ only).} Put
  $J_Y:=\lfloor\log(L_Y/8)\rfloor$ and $\kown:=4J_Y+1$. Every edge of
  $\Own_Y$ independently receives one of $\kown$ labels, chosen
  uniformly. This yields the own classes $\vR_{j,c}$ ($0\le j<J_Y$,
  $c\in[4]$) and $\vM$ of $Y$, each regarded as a graph on $V(Y)$. They
  are the own classes of the four-phase P-vortex run on $Y$ in later
  sections. For standalone $Y$, $\Own_Y$ is not subdivided at stage~1.
\item[(iv)] \emph{JS labels.} For every $r+2\le l\le R$ and every
  $y\in V(Y)$, independently, draw a label
  $\lab_{Y,l}(y)\in\{\ast\}\cup\{0,1,\dots,\KJS_l-1\}$ with
  $\Prob(\lab_{Y,l}(y)=j)=\rho_l:=M_l^{-4}$ for each $j$, so that
  $\Prob(\lab_{Y,l}(y)=\ast)=1-M_l^{-2}$. Put
  $T_j(Y,l):=\{y\in V(Y):\lab_{Y,l}(y)=j\}$.
\end{itemize}
Formally, every edge $e$ of $H_Y$ carries three independent uniform
random variables (a fair bit, a lent index and an own label), and the
labels of (iv) are independent of all edge variables. Consequently:
\begin{itemize}
\item colours of distinct edges are independent;
\item the colourings of distinct ancestors are independent;
\item the labels are independent of the colours;
\item given $\Lend_Y$, stage (ii) is a uniform $\klend(Y)$-colouring of
  $\Lend_Y$, and given $\Own_Y$, stage (iii) is a uniform
  $\kown$-colouring of $\Own_Y$;
\item for fixed $(l,j)$, the set $T_j(Y,l)$ is a $\rho_l$-random subset
  of $V(Y)$, independent of the colouring, and the sets $T_j(Y,l)$,
  $0\le j<\KJS_l$, are pairwise disjoint.
\end{itemize}

\emph{Derived quantities.} For $\klend(Y)\ge1$ put
$p_Y:=1/(2\klend(Y))$. For each fixed $l$ with $r+2\le l\le R$, this is
the probability that a fixed edge of $H_Y$ lies in $\LJV_{Y,l}$. Also put
$\tJS_l:=2M_l+2$. For a light ancestor $Y$ of round $r$ put
$t_Y:=\lceil\lambda_r^{1.6}\rceil$, the \emph{U-multiplicity parameter}
of $Y$: it is the multiplicity bound $t$ at which the U-lent classes of
$Y$ are path connected (Lemma~\ref{s3:lemCOL}(c)). It depends only on
the round of $Y$, and $t_Y\ge M_l$ for every $r+2\le l\le R$, because
$M_l\le\lambda_{l-2}^{1.6}\le\lambda_r^{1.6}$
(Lemma~\ref{s2:lemTower}(a),(b)).

\emph{Consumers} (property COL(g)). If $r\le R-2$, each part of the data
has exactly one consumer:
\begin{itemize}
\item $\Own_Y$, together with its own classes if $Y$ is light, is used
  only by the vortex device of $Y$ itself. This is the transparent
  P-vortex if $Y$ is standalone. If $Y$ is light, it is the four-phase
  P-vortex, or VX$^+$ if $Y$ is demoted.
\item $\LU_{Y,l,c,\uslot}$ is used only by the Theorem-U chaining of the
  U-bundle $\Bdl_{l,c}$, in slot $\uslot$.
\item $\LJS_{Y,l,j}$ is used only by the JS-LC systems of the pair
  $(Y,l)$, at junction $j$.
\item $\LJV_{Y,l}$ is used only by the round-$l$ J-consumer.
\end{itemize}
Lent edges that their consumer does not use are junk for $Y$'s own
device: they are added to the edge set that this device decomposes.
If $r\ge R-1$, then $\Own_Y$ has the same consumer, and $\Lend_Y$ has no
classes and no other consumer: all of $\Lend_Y$ is junk for $Y$'s own
device. Indeed $\Lent(Y)=\emptyset$, since it collects edges used at
rounds $l\ge r+2>R$, so $\Erem(Y)=E_r(Y)$ for standalone $Y$ and
$H_0(Y)=E_r(Y)$ for light non-demoted $Y$ (Construction~\ref{s6:consOrder});
a demoted $Y$ runs on $E_r(Y)$ itself.
\deps{\ref{s2:defAncestors}, \ref{s2:propStructure}, \ref{s2:lemTower}}
\end{definition}

The graphs $H_Y$ of all ancestors $Y$ of all rounds are pairwise
edge-disjoint (Proposition~\ref{s2:propStructure}(iii)). So the
colourings of distinct ancestors act on disjoint edge sets. The label
distribution in (iv) is well defined because
$\KJS_l\rho_l=M_l^{-2}\le1$.

\begin{table}[htbp]
\begin{center}
\small
\begin{tabular}{|p{0.55cm}|p{4.6cm}|p{5.7cm}|p{2.4cm}|}
\hline
No. & Requirement (use) & Sufficient inequality in $\lambda=2^\mu$ & $(a_i,b_i,c_i)$\\\hline
1 & $\thGC_r(Z^0)>\tau_r$ (rule GC) & $\lambda^{102.5}>257(\lambda+6\mu+20)^{102}$; crude form $\lambda^{1/2}>2^{111}$ & $(0.5,112,0)$\\\hline
2 & $P_{l-2}\ge M_l^{13}$, with $\lambda:=\lambda_{l-2}$ & $\lambda^{103}\ge\bar M^{13}$ & $(103,0,26A)$\\\hline
3 & Lemma 15$^+$, lent level: $s_Y/4\ge40\,\klend L_Y$ & $\lambda^{99}\ge640\,\bar k$ & $(96,19,4A)$\\\hline
4 & T16$^*$ for U-lent classes: $s_Y/(8\klend)\ge2^{135}\,t_Y\,L_Y^{28}(12L_Y^5)^5$, $t_Y=\lceil\lambda^{1.6}\rceil$ & $\lambda^{42.1}\ge2^{193}\cdot12^5$ (with row~9)$^{\dagger}$ & $(42.1,211,0)$\\\hline
5 & T16$^*$ for JS-lent classes: $s_Y/(8\klend)\ge2^{135}(3M)L_Y^{28}M^{20}$ & $\lambda^{72}\ge3\cdot2^{167}\cdot\bar k\,\bar M^{21}$ & $(69,178,46A)$\\\hline
6 & T16$^*$ failures over $\IJS(Y)$ sum to at most $|V(Y)|^{-2}/4$ & $\lambda^{84}\ge2^{110}\bar M^{15}$ & $(84,110,30A)$\\\hline
7 & T16$^*$ failures over $\IU(Y)$ sum to at most $|V(Y)|^{-2}/4$ & $\lambda^{64.4}\ge2^{127}\cdot12^4$ & $(64.4,142,0)$\\\hline
8 & own devices: (a) $s_r/4\ge2^{150}L_Y^{42}$; (b) $s_r/4\ge2^{146}L_Y^{38}\log L_Y$ (the threshold of VX$^+$ for standalone parts; not used, as VX$^+$ is applied only to demoted light parts, Lemma~\ref{s5:lemDemoted}; cf.\ Remark~\ref{s4:remTPVVX}); (c) $s_r/2\ge2^{151}L_Y^{38}\log L_Y$; (d) $s':=s_r/(16\kown)\ge2^{145}L_Y^{41}$ and $s'\ge2\lceil L_Y^6\rceil$ & (a) $\lambda^{58}\ge2^{194}$; (b) $\lambda^{62}\ge2^{186}(\mu+1)$; (c) $\lambda^{62}\ge2^{190}(\mu+1)$; (d) $\lambda^{58}\ge2^{191}$ and $\lambda^{99}\ge64((2\lambda)^6+1)$ & $(58,194,1)$\\\hline
9 & $\klend(Y)\le\lambda^{3.3}$ & $\bar k\le\lambda^{3.3}$ & $(0.3,9,4A)$\\\hline
10 & $p_Y\ge\lambda^{-4}$ & $2\bar k\le\lambda^{4}$ & $(1,10,4A)$\\\hline
11 & $\lambda_{l-2}^{95}/(8M_l^2)\ge2^{10}M_l^{10}$, with $\lambda:=\lambda_{l-2}$ & $\lambda^{95}\ge2^{13}\bar M^{12}$ & $(95,13,24A)$\\\hline
12 & $M_l\le\lambda_{l-2}^{1.6}$, with $\lambda:=\lambda_{l-2}$ & $\bar M\le\lambda^{1.6}$ & $(1.6,0,2A)$\\\hline
13 & \Gstar & $\lambda^{36}\ge2^{240}(A\mu)^{46A}$ & $(36,240,46A)$\\\hline
\end{tabular}
\end{center}
\caption{The COL-JV table (Lemma~COL-JV below). Here $\lambda:=\lambda_r$
and $\mu:=\log\lambda$, except in rows 2, 11 and 12, where
$\lambda:=\lambda_{l-2}$ for a round $3\le l\le R$. We write
$\bar M:=(A\mu)^{2A}$, with $A=105$, and $\bar k:=192\lambda^3+\frac43\bar M^2$.
Column~3 is a sufficient form of the requirement, obtained from the
standing bounds (B1)--(B7) of the proof. Column~4 gives, for row~$i$, reals
$a_i>0$ and $b_i,c_i\ge0$ such that, for $\mu\ge6$, the inequality
$a_i\mu-b_i-c_i\log(A\mu)\ge0$ implies every inequality in column~3 of row~$i$
(Lemma~\ref{s3:lemCOLJVev}, which has no hypothesis); hence every
column-3 inequality holds for all sufficiently large $\mu$, and \Gc{1}(f)
requires it for every $\mu\ge\log\log\Dstar$. The inequalities
$a_i\mu-b_i-c_i\log(A\mu)\ge0$ defined by column~4 are cruder than column~3; they
are not part of \Gc{1}, are not asserted in Lemma~COL-JV, and may fail at some
$\mu\ge\log\log\Dstar$ although \Gc{1} holds. No galactic threshold is
evaluated: the thresholds of Lemma~\ref{s3:lemCOLJVev} are explicit, but their
numerical values are never computed. Rows for off-route devices are omitted.
$^{\dagger}$Row~4: column~3 implies the column-2 requirement together with
row~9.}
\label{s3:tabCOLJV}
\end{table}

\begin{lemma}[the COL-JV table]\label{s3:lemCOLJV}
Assume condition~\ref{s1:condG1}. Let $Y$ be an ancestor of round $r\le R$, and write
$\lambda:=\lambda_r$ and $\mu:=\log\lambda$; if $r\le R-2$, write also $M:=M_{r+2}$. Then:
\begin{itemize}
\item[(i)] If $r\ge R-1$, then $\klend(Y)=0$. If $r\le R-2$, then
  $|\IU(Y)|\le12L_Y^3$, $|\IJS(Y)|\le\frac43M^2$ and
  $|\IJV(Y)|=R-r-1\le2\logs d_r+1$. Hence
  \[
  \klend(Y)\le12L_Y^3+\tfrac43M^2+2\logs d_r+2\le24L_Y^3+\tfrac43M^2\le\bar k\le\lambda^{3.3},
  \]
  and consequently $p_Y\ge\lambda^{-4}$. The JV family has
  $R-r-1\le2\logs d_r+1$ members. Its size is not bounded by an
  absolute constant, since $d_1$ may be of order $n$; the count
  ``$+16$'' must not be used.
\item[(ii)] Every requirement in column~2 and every inequality in column~3
  of Table~\ref{s3:tabCOLJV} holds. (Column~4 is not part of this assertion:
  its triples serve only Lemma~\ref{s3:lemCOLJVev}, and the inequalities
  $a_i\mu-b_i-c_i\log(A\mu)\ge0$ that they define need not hold at $\mu$.) More
  precisely, for each row every inequality in column~3 holds at every
  $\mu\ge\log\log\Dstar$ by \Gc{1}(f); in particular it holds at
  $\mu=\log\lambda_r$ and, for rows~2, 11 and~12, also at
  $\mu=\log\lambda_{l-2}$ for every round $3\le l\le R$, both of which are
  at least $\log\log\Dstar$. At the value of $\lambda$ fixed in the caption of
  Table~\ref{s3:tabCOLJV}, these inequalities imply the requirement in
  column~2 (for row~4, together with row~9). So every requirement holds, in
  every round.
  The scope of the rows is as follows.
  \begin{itemize}
  \item Rows~1, 3 and~8 (rule GC, the three levels of Lemma~15$^+$, and
    the own-device thresholds) hold for every round $r\le R$, that is,
    also for $r\in\{R-1,R\}$: row~1 for every round-$r$ pre-part, and
    rows~3 and~8 for every ancestor $Y$ of round $r$. They involve
    neither $M_{r+2}$ nor any lent class; for $r\ge R-1$ the lent level
    of row~3 is void, because $Y$ has no lent classes.
  \item Rows~4--7 and~10 concern lent classes or $p_Y$. They are
    asserted for $r\le R-2$; for $r\ge R-1$ there are no lent classes
    and $p_Y$ is not defined. Row~9 holds for every $r\le R$, trivially
    when $\klend(Y)=0$.
  \item Rows~2, 11 and~12 concern $\lambda_{l-2}$ for rounds
    $3\le l\le R$, and row~13 is \Gstar; they do not depend on $Y$.
  \end{itemize}
\end{itemize}
\deps{\ref{s1:condG1}, \ref{s1:condGstar}, \ref{s2:lemTower}, \ref{s3:defCOL}, \ref{s2:defAncestors}, \ref{s2:propStructure}, \ref{s3:lemL15p}, \ref{s3:thmT16s}}
\end{lemma}

\begin{proof}
Every round $r\le R$ has $d_r\ge\Dstar$, so $\mu=\log\lambda_r\ge\log\log\Dstar$, and
the items of \Gc{1} hold at $\mu$. The same holds for $\lambda_{l-2}$ with
$3\le l\le R$. We use
the following standing bounds.
\begin{itemize}
\item[(B1)] $L_Y\le\log M_r\le2\lambda$ (Lemma~\ref{s2:lemTower}(b)).
\item[(B2)] If $r\le R-2$: $M_l\le M\le\bar M$ for $r+2\le l\le R$, and
  $M_{l+1}\le M_l/2$ for $r+2\le l<R$ (Lemma~\ref{s2:lemTower}(b)). Likewise
  $M_l\le(A\log\lambda_{l-2})^{2A}$ for every $3\le l\le R$.
\item[(B3)] $s_r\ge\lambda^{100}$ and $s_r\le2\Lambda_r^{100}$ (Lemma~\ref{s2:lemTower}(b)),
  and $s_Y\ge s_r/2$ (Definition~\ref{s2:defAncestors}).
\item[(B4)] $|V(Y)|\ge P_r/2\ge\lambda^{103}/2$. A standalone $Y$ has
  $|V(Y)|=|Y^0|\ge P_r$, and a light $Y$ has $|Y|\ge|Y^0|/2\ge P_r/2$
  (Proposition~\ref{s2:propStructure}(iv)); and $P_r=\lceil\lambda^{C'}\rceil$
  with $C'=103$.
\item[(B5)] $R-r\le2\logs d_r+2$ and $2^{2\logs d_r+2}\le\lambda$
  (Lemma~\ref{s2:lemTower}(a),(d)), so $2\logs d_r+2\le\mu$.
\item[(B6)] $\Lambda_r\le\lambda+6\mu+20$. Indeed, with
  $T:=\tHB_rd_r=\lambda^2\,2^{\lambda}$ we have $\log T=\lambda+2\mu$. Put
  $B:=2^{16}T\log^4T$. As $\lambda=2^\mu\ge2^{256}$ ($\mu\ge2^8$ by \Gc{1}(a)) and
  $\log T=\lambda+2\mu\ge1$, we have $B\ge2^{16}T\ge2^{16+\lambda}>2^{161}\ge2^{40}$,
  so $M_r\le B+1$ by (R2) (Definition~\ref{s2:defHBtp}), and
  $\Lambda_r=\log M_r\le\log B+\log(1+1/B)\le\log B+1/(B\ln2)\le\log B+2^{-160}$,
  since $1/(B\ln2)<2^{-161}/0.69<2^{-160}$. Next
  $\log B=16+\lambda+2\mu+4\log(\lambda+2\mu)$, and
  $\log(\lambda+2\mu)=\mu+\log(1+2\mu/\lambda)\le\mu+2\mu/(\lambda\ln2)\le\mu+0.0025$,
  because $\mu\ge2^8$ and $2\mu/\lambda=2\mu2^{-\mu}\le2^{-240}$. Hence
  $\Lambda_r\le\lambda+6\mu+16.01+2^{-160}\le\lambda+6\mu+20$.
\item[(B7)] If $r\le R-2$ and $r+2\le l\le R$: $\tJS_l=2M_l+2\le3M_l\le3M$
  (as $M_l\ge2^{40}$), and $\rho_l^{-1}=M_l^4\le M^4$.
\end{itemize}
Only (B2) and (B7) involve $M$, and they are used only for $r\le R-2$.
(B1) and (B3)--(B6) hold for every round $r\le R$.

\emph{(i) Counting.} If $r\ge R-1$, the three index families are empty
(Definition~\ref{s3:defCOL}(ii)), so $\klend(Y)=0$. Let now $r\le R-2$.
\begin{itemize}
\item $|\IU(Y)|=4\Tslot_Y(R-r-1)\le4(L_Y^2+1)(2\logs d_r+1)$. By (B5)
  and (B4), $2\logs d_r+1<\mu\le L_Y$, since $L_Y\ge103\mu-1$. Hence
  $|\IU(Y)|\le4(L_Y^2+1)L_Y\le12L_Y^3$.
\item $|\IJS(Y)|=\sum_{l=r+2}^RM_l^2\le M^2\sum_{i\ge0}4^{-i}=\frac43M^2$
  by (B2).
\item $|\IJV(Y)|=R-r-1\le2\logs d_r+1$ by (B5).
\end{itemize}
The second inequality of the display holds because $2\logs d_r+2\le\mu\le12L_Y^3$.
The third follows from (B1) and (B2). The fourth is the sufficient
inequality of row~9, verified in (ii) below.
Then $p_Y=1/(2\klend(Y))\ge1/(2\lambda^{3.3})\ge\lambda^{-4}$, and row~10
gives the same conclusion directly.

\emph{(ii) The rows.} In each row, the passage from column~2 to column~3
uses only (B1)--(B7) and the fact that the requirement is monotone in
$L_Y$, $M_l$, $\klend$, $s_r$ and $|V(Y)|$ in the appropriate direction.
Rows~1, 3 and~8 use only (B1), (B3)--(B6), Lemma~\ref{s2:lemTower}(c)
and the bound $\klend(Y)\le\bar k$ of (i), which is trivial when
$\klend(Y)=0$; so they hold for every round $r\le R$. Rows~4--7, 9 and~10
are derived for $r\le R-2$, where (B2) and (B7) are available; for
$r\ge R-1$, row~9 reads $0\le\lambda^{3.3}$.
We now give the derivations.
\begin{itemize}
\item[1.] For a round-$r$ pre-part $Z^0$ we have $|Z^0|\ge P_r\ge\lambda^{103}$, so
  $\thGC_r(Z^0)=\lceil|Z^0|\lambda^{-1/2}\rceil\ge\lambda^{102.5}$,
  and $\tau_r\le128s_r\Lambda_r^2+1\le257\Lambda_r^{102}$ by (B3). Now
  use (B6). The crude form uses $\tau_r\le2^{111}\lambda^{102}$
  (Lemma~\ref{s2:lemTower}(c)).
\item[2.] $P_{l-2}\ge\lambda_{l-2}^{103}$, and $M_l\le\bar M$ by (B2).
\item[3.] $s_Y/4\ge\lambda^{100}/8$ and $40\klend L_Y\le80\lambda\bar k$.
  The other two levels of Lemma~\ref{s3:lemL15p} are $s_Y\ge80L_Y$ (the
  Own/Lend split, $k=2$) and $s_r/8\ge40\kown L_Y$ (own classes). They
  follow from $\lambda^{100}/2\ge160\lambda$ and
  $\lambda^{100}/8\ge160\lambda^2$, using $\kown\le L_Y\le2\lambda$, and
  both hold for every $\mu\ge1$.
\item[4.] U-lent classes exist only for light $Y$. Theorem~\ref{s3:thmT16s},
  with $t=t_Y=\lceil\lambda^{1.6}\rceil\le2\lambda^{1.6}$ and
  $\rho\ge1/(12L_Y^5)$, requires $s\ge2^{135}tL_Y^{28}\rho^{-5}$, and
  $2^{135}tL_Y^{28}\rho^{-5}\le2^{136}12^5\lambda^{1.6}L_Y^{53}$.
  With $s=s_Y/(8\klend)\ge\lambda^{100}/(16\bar k)$ and $L_Y\le2\lambda$,
  it suffices that
  $\lambda^{100}\ge2^{4+136+53}12^5\bar k\lambda^{54.6}$, i.e.\
  $\lambda^{45.4}\ge2^{193}\cdot12^5\cdot\bar k$. By row~9,
  $\bar k\le\lambda^{3.3}$, so the sufficient inequality of column~3,
  $\lambda^{42.1}\ge2^{193}\cdot12^5$, implies it. (Row~9 holds at
  $\mu$ under \Gc{1}; so column~3 implies column~2 under \Gc{1}.)
\item[5.] Theorem~\ref{s3:thmT16s} with $t=\tJS_l\le3M$ and
  $\rho=\rho_l\ge M^{-4}$ requires $s\ge2^{135}\cdot3M\cdot L_Y^{28}M^{20}$.
  With $s\ge\lambda^{100}/(16\bar k)$, $L_Y\le2\lambda$ and $M\le\bar M$,
  this becomes $\lambda^{100}\ge3\cdot2^{4+135+28}\bar k\bar M^{21}\lambda^{28}$.
\item[6.] The failure bound of Theorem~\ref{s3:thmT16s} is
  $2^{86}tL^{19}\rho^{-3}N^{-3}$ per class, with $N:=|V(Y)|$. Summed over
  the at most $\frac43M^2$ JS classes, with $t\le3M$, $\rho^{-3}\le M^{12}$
  and $L_Y\le2\lambda$, it is at most $2^{107}\lambda^{19}M^{15}N^{-3}$.
  This is at most $N^{-2}/4$ if $N\ge2^{109}\lambda^{19}M^{15}$, which by
  (B4) follows from $\lambda^{84}\ge2^{110}\bar M^{15}$. This line
  needs \Gc{1}: the class count and the multiplicity grow with $M$,
  which is bounded only by $\bar M=(A\mu)^{2A}$, not by a fixed power of
  $\lambda$ for moderate $\mu$. It
  does not follow from $C'>34$ alone.
\item[7.] There are at most $12L_Y^3$ U-classes, each with
  $t=t_Y\le2\lambda^{1.6}$ and $\rho^{-3}\le1728L_Y^{15}$. The sum is at most
  $12L_Y^3\cdot2^{86}\cdot2\lambda^{1.6}L_Y^{19}\cdot12^3L_Y^{15}N^{-3}
  =12^4\cdot2^{87}\lambda^{1.6}L_Y^{37}N^{-3}\le12^4\cdot2^{124}\lambda^{38.6}N^{-3}$.
  This is at most $N^{-2}/4$ if $N\ge12^4\cdot2^{126}\lambda^{38.6}$,
  which by (B4) follows from
  $\lambda^{103}/2\ge12^4\cdot2^{126}\lambda^{38.6}$, i.e.\
  $\lambda^{64.4}\ge2^{127}\cdot12^4$.
\item[8.] By (B3), $s_r/4\ge\lambda^{100}/4$ and $s_r/2\ge\lambda^{100}/2$.
  Also $\kown\le L_Y\le2\lambda$, so $s'\ge\lambda^{99}/32$, and
  $\log L_Y\le\mu+1$. The right sides are increasing in $L_Y$.
  Substituting $L_Y\le2\lambda$ gives the four inequalities.
\item[9.--10.] By the first three inequalities of (i),
  $\klend(Y)\le\bar k$. So $\bar k\le\lambda^{3.3}$ suffices for row~9,
  and, since $p_Y=1/(2\klend)\ge1/(2\bar k)$, $2\bar k\le\lambda^4$
  suffices for row~10.
\item[11.] $\lambda^{95}/(8M_l^2)\ge2^{10}M_l^{10}$ is equivalent to
  $\lambda^{95}\ge2^{13}M_l^{12}$; use (B2).
\item[12.--13.] Row~12 is (B2) together with the inequality shown. Row~13
  is the condition \Gstar of Condition~\ref{s1:condGamma} itself.
\end{itemize}
\end{proof}

The next lemma shows that every column-3 inequality of
Table~\ref{s3:tabCOLJV} holds for all sufficiently large $\mu$. It has no
hypothesis. It is used only in the proof of Lemma~\ref{s7:lemGammaSat}, to show
that item~(f) of \Gc{1} can be satisfied.

\begin{lemma}[eventual form of the COL-JV table]\label{s3:lemCOLJVev}
Call an inequality in a real variable $\mu$ \emph{of type (E)} if it reads
$a\mu-b-c\log(A\mu)\ge0$ with reals $a>0$ and $b,c\ge0$ (here $A=105$ and
$\log=\log_2$).
\begin{itemize}
\item[(i)] An inequality of type (E) holds for every
  $\mu\ge\max\{1,v_0^2\}$, where
  $v_0\defeq\bigl(c+(c^2+a(b+c\log A))^{1/2}\bigr)/a$. In particular, it holds
  for all sufficiently large $\mu$.
\item[(ii)] Let $\mu\ge6$ be real, and put $\lambda\defeq2^\mu$,
  $\bar M\defeq(A\mu)^{2A}$ and $\bar k\defeq192\lambda^3+\frac43\bar M^2$. Let
  $1\le i\le13$, and let $(a_i,b_i,c_i)$ be the triple in column~4 of row~$i$
  of Table~\ref{s3:tabCOLJV}. If $a_i\mu-b_i-c_i\log(A\mu)\ge0$, then every
  inequality in column~3 of row~$i$ holds at $\lambda$.
\item[(iii)] Consequently, for each row of Table~\ref{s3:tabCOLJV}, every
  inequality in column~3 holds at $\lambda=2^\mu$ for all sufficiently large
  real $\mu$.
\end{itemize}
The lemma concerns explicit functions of the real variable $\mu$ only. It
involves no graph, no run of the hierarchy and no condition on $\Dstar$.
\deps{none (the inequalities are those printed in Table~\ref{s3:tabCOLJV},
which is not a statement).}
\end{lemma}

\begin{proof}
(i) Let $\mu\ge\max\{1,v_0^2\}$ and $v\defeq\mu^{1/2}$, so that $v\ge1$ and
$v\ge v_0$. For $v\ge1$ we have $2^v\ge1+v>v$: equality $2^v=1+v$ holds at
$v=1$, and for $v\ge1$ the derivative $2^v\ln2\ge2\ln2>1$ of the left side
exceeds the derivative $1$ of the right side. Hence $\log v<v$, and
$\log\mu=2\log v\le2v$. As $c\ge0$,
\[
a\mu-b-c\log(A\mu)\ =\ av^2-(b+c\log A)-c\log\mu\ \ge\ av^2-2cv-(b+c\log A).
\]
The right side is a quadratic polynomial in $v$ with leading coefficient
$a>0$. Its discriminant $4c^2+4a(b+c\log A)$ is non-negative, as $b,c\ge0$ and
$\log A>0$, and its larger root is $v_0$. So it is non-negative for every
$v\ge v_0$, and the inequality of type (E) holds at $\mu$.

(ii) We use three elementary bounds, valid for $\mu\ge6$. First,
$\log\mu\le\log(A\mu)$ and $\log(\mu+1)\le1+\log\mu\le1+\log(A\mu)$ (as
$\mu+1\le2\mu$), and $\log(A\mu)\ge0$. Second, $6\mu+20\le2^\mu$ (at $\mu=6$
this reads $56\le64$, and for $\mu\ge6$ the derivative $2^\mu\ln2\ge44$ of the
right side exceeds the derivative $6$ of the left side), so
$\log(2^\mu+6\mu+20)\le\log(2^{\mu+1})=\mu+1$; and
$(2\lambda)^6+1\le2\cdot(2\lambda)^6=2^7\lambda^6$. Third, both summands of
$\bar k$ are at least $1$ (as $\lambda\ge1$ and $A\mu\ge1$), and
$x+y\le2\max\{x,y\}\le2xy$ for $x,y\ge1$; so
$\bar k\le2\cdot192\lambda^3\cdot\frac43\bar M^2=2^9\lambda^3\bar M^2$ and
\[
\log\bar k\ \le\ 3\mu+4A\log(A\mu)+9 .
\]
Also $\log\bar M=2A\log(A\mu)$. For an inequality of column~3 put
$\varphi\defeq\log(\text{left side})-\log(\text{right side})$; the inequality
holds if and only if $\varphi\ge0$ (in row~1, where it is strict, if and only
if $\varphi>0$). The bounds above give, row by row:
\begin{itemize}
\item[1.] $\varphi=102.5\mu-\log257-102\log(2^\mu+6\mu+20)\ge102.5\mu-\log257-102(\mu+1)\ge0.5\mu-111$,
  as $\log257<9$; so $\varphi\ge(0.5\mu-112)+1>0$ if $0.5\mu-112\ge0$. The crude
  form $\lambda^{1/2}>2^{111}$ reads $\mu>222$, which also holds if
  $0.5\mu-112\ge0$, since then $\mu\ge224$;
\item[2.] $\varphi=103\mu-13\log\bar M=103\mu-26A\log(A\mu)$;
\item[3.] $\varphi=99\mu-\log640-\log\bar k\ge96\mu-19-4A\log(A\mu)$, as
  $\log640<10$;
\item[4.] $\varphi=42.1\mu-193-5\log12\ge42.1\mu-211$, as $5\log12<18$;
\item[5.] $\varphi=72\mu-\log3-167-\log\bar k-21\log\bar M\ge
  69\mu-178-46A\log(A\mu)$, as $\log3<2$;
\item[6.] $\varphi=84\mu-110-15\log\bar M=84\mu-110-30A\log(A\mu)$;
\item[7.] $\varphi=64.4\mu-127-4\log12\ge64.4\mu-142$, as $4\log12<15$;
\item[8.] put $E\defeq58\mu-194-\log(A\mu)$. The five inequalities of row~8
  hold if the following five differences are non-negative:
  $58\mu-194\ge E$;
  $62\mu-186-\log(\mu+1)\ge62\mu-187-\log(A\mu)=E+4\mu+7$;
  $62\mu-190-\log(\mu+1)\ge62\mu-191-\log(A\mu)=E+4\mu+3$;
  $58\mu-191\ge58\mu-194\ge E$; and
  $99\mu-6-\log((2\lambda)^6+1)\ge99\mu-6-(7+6\mu)=93\mu-13>0$. So all five
  hold if $E\ge0$;
\item[9.] $\varphi=3.3\mu-\log\bar k\ge0.3\mu-9-4A\log(A\mu)$;
\item[10.] $\varphi=4\mu-1-\log\bar k\ge\mu-10-4A\log(A\mu)$;
\item[11.] $\varphi=95\mu-13-12\log\bar M=95\mu-13-24A\log(A\mu)$;
\item[12.] $\varphi=1.6\mu-\log\bar M=1.6\mu-2A\log(A\mu)$;
\item[13.] $\varphi=36\mu-240-46A\log(A\mu)$.
\end{itemize}
In each row the lower bound obtained is $a_i\mu-b_i-c_i\log(A\mu)$ for the
triple $(a_i,b_i,c_i)$ of column~4 (in row~1 it is this quantity plus $1$, and
in row~8 it is at least $E=a_8\mu-b_8-c_8\log(A\mu)$ for the four inequalities
that depend on it). So if $a_i\mu-b_i-c_i\log(A\mu)\ge0$, every inequality in
column~3 of row~$i$ holds at $\lambda$. This proves (ii).

(iii) Fix a row~$i$, and let $v_0$ be as in (i) for
$(a,b,c)=(a_i,b_i,c_i)$. By (i), the inequality
$a_i\mu-b_i-c_i\log(A\mu)\ge0$ holds for every $\mu\ge\max\{1,v_0^2\}$; by (ii),
every inequality in column~3 of row~$i$ holds at every
$\mu\ge\max\{6,v_0^2\}$. This threshold is explicit, but its numerical value is
never computed: only its existence is used.
\end{proof}

\begin{lemma}[Lemma COL]\label{s3:lemCOL}
Assume condition~\ref{s1:condG1}. Let $Y$ be an ancestor of round $r$, with the stage-1
lending data of Definition~\ref{s3:defCOL}. If $r\ge R-1$, then
$\klend(Y)=0$ and every statement below about lent classes is vacuous.
\begin{itemize}
\item[(a)] $\Own_Y$ and $\Lend_Y$ are $(\varepsilon_Y,s_Y/4)$-expanders
  on $V(Y)$. Every lent class, as a graph on $V(Y)$, is an
  $(\varepsilon_Y,\,s_Y/(8\klend(Y)))$-expander. If $Y$ is light, every
  own class $\vR_{j,c}$, $\vM$, as a graph on $V(Y)$, is a
  $(2^{-6},\,s_r/(16\kown))$-expander.
\item[(b)] For all $r+2\le l\le R$ and $0\le j<\KJS_l$, the class
  $\LJS_{Y,l,j}$ is $(2^{12}L_Y^4,\tJS_l)$-path connected through
  $T_j(Y,l)$. By Lemma~\ref{s3:lemMonotone}(iii), one application routes
  any multiset of pairs of distinct vertices in which every vertex lies
  in at most $\tJS_l$ pairs. For instance, the union of the junction-$j$
  pair multisets of several systems, with multiplicities counted over
  the union, is joined by pairwise edge-disjoint paths of
  $\LJS_{Y,l,j}$ of length at most $2^{12}L_Y^4$ with interiors in
  $T_j(Y,l)$.
\item[(c)] Let $Y$ be light. Let $(V_{l,c,\uslot})_{(l,c,\uslot)\in\IU(Y)}$
  be random subsets of $V(Y)$, independent of the stage-1 lending data
  of $Y$, such that each $V_{l,c,\uslot}$ is a
  $\rho_{l,c,\uslot}$-random subset of $V(Y)$ with
  $\rho_{l,c,\uslot}\ge1/(12L_Y^5)$. The family may be dependent across
  indices. Then, with probability at least $1-|V(Y)|^{-2}/2$ over the
  lending data and the sets, every U-lent class $\LU_{Y,l,c,\uslot}$ is
  $(2^{12}L_Y^4,t_Y)$-path connected through $V_{l,c,\uslot}$, where
  $t_Y=\lceil\lambda_r^{1.6}\rceil$ (Definition~\ref{s3:defCOL}).
\item[(e)] The own-device thresholds hold. If $Y$ is standalone, then
  $s_Y/4\ge2^{150}L_Y^{42}$ and $s_Y/4\ge2^{146}L_Y^{38}\log L_Y$. On (a),
  $\Own_Y$ is therefore a spanning $(2^{-5},s_r/4)$-expander of $V(Y)$
  whose robustness meets the thresholds of the transparent P-vortex
  and of VX$^+$ (the second of these is not used: VX$^+$ is applied only to
  demoted light parts, Lemma~\ref{s5:lemDemoted}; cf.\
  Remark~\ref{s4:remTPVVX}). If $Y$ is light, then
  $s_r/2\ge2^{151}L_Y^{38}\log L_Y$; this is the threshold for VX$^+$ at
  expansion $2^{-6}$ on $H_Y=X_Y$, which is a
  $(2^{-6},s_r/2)$-expander. Moreover, on (a) every own class is a
  $(2^{-6},s')$-expander with $s':=s_r/(16\kown)\ge2^{145}L_Y^{41}$ and
  $2\lceil L_Y^6\rceil\le s'$.
\item[(g)] (COL(g).) The consumers listed in Definition~\ref{s3:defCOL}
  are exclusive. The own classes partition $\Own_Y$, and if $r\le R-2$,
  the lent classes are pairwise edge-disjoint and partition $\Lend_Y$. So
  no edge of $H_Y$ is available to two consumers, and every lent edge
  that its consumer does not use is junk for $Y$'s own device. If
  $r\ge R-1$, then $\Lend_Y$ has no classes, and all of $\Lend_Y$ is junk
  for $Y$'s own device (Definition~\ref{s3:defCOL}, Consumers).
\end{itemize}
Items (a), (b) and (e) hold simultaneously with probability at least
$1-|V(Y)|^{-2}/2$ over the stage-1 lending data of $Y$. For every family
as in (c), items (a), (b), (c) and (e) hold simultaneously with
probability at least $1-|V(Y)|^{-2}$. Item (g) always holds. The failure
unions are bounded through rows~6 and~7 of Table~\ref{s3:tabCOLJV}, which
need \Gc{1}. There are no items (d) and (f); the letters
follow the references used in later sections.
\deps{\ref{s3:defCOL}, \ref{s3:lemL15p}, \ref{s3:thmT16s}, \ref{s3:lemMonotone}, \ref{s3:lemCOLJV}, \ref{s2:lemTower}, \ref{s2:propStructure}, \ref{s2:defAncestors}, \ref{s1:condG1}}
\end{lemma}

\begin{proof}
Write $N:=|V(Y)|$, $L:=L_Y$, $k:=\klend(Y)$, $\lambda:=\lambda_r$ and
$\mu:=\log\lambda$. By Proposition~\ref{s2:propStructure}(i),
$H_Y$ is an $(\varepsilon_Y,s_Y)$-expander on $V(Y)$. Here
$(\varepsilon_Y,s_Y)=(2^{-5},s_r)$ if $Y$ is standalone and
$(2^{-6},s_r/2)$ if $Y$ is light
(Definition~\ref{s2:defAncestors}), and $s_r\ge\lambda^{100}$
(Lemma~\ref{s2:lemTower}(b)). The standing bounds (B1) and (B3)--(B6) of
the proof of Lemma~\ref{s3:lemCOLJV} are available in every round, and
(B2), (B7) when $r\le R-2$. Rows~1, 3, 8 and~9 of
Table~\ref{s3:tabCOLJV} hold for every $r\le R$, and the other rows
involving $Y$ hold for $r\le R-2$ (Lemma~\ref{s3:lemCOLJV}(ii)); for
$r\ge R-1$ only rows~3, 8 and~9 are used below. In particular
$N\ge\lambda^{103}/2$ and $L\le2\lambda$.

\emph{(a)} Stage (i) is a uniform $2$-colouring of $E(H_Y)$. By
Lemma~\ref{s3:lemL15p} with $k=2$, which needs $s_Y\ge80L$ (row~3 of
Table~\ref{s3:tabCOLJV}), $\Own_Y$ and $\Lend_Y$ are
$(\varepsilon_Y,s_Y/4)$-expanders on $V(Y)$, except with probability at
most $4N^{-5}$.

Condition on stage (i) with $\Lend_Y$ such an expander. Stage (ii) is
then a uniform $k$-colouring of $\Lend_Y$. By Lemma~\ref{s3:lemL15p},
which needs $s_Y/4\ge40kL$ (row~3), every lent class is an
$(\varepsilon_Y,s_Y/(8k))$-expander on $V(Y)$, except with probability
at most $2kN^{-5}$.

If $Y$ is light, condition on $\Own_Y$ being a
$(2^{-6},s_r/8)$-expander. Stage (iii) is a uniform $\kown$-colouring.
Lemma~\ref{s3:lemL15p}, which needs $s_r/8\ge40\kown L$ (row~3), makes
every own class an expander with parameters $2^{-6}$ and
$s_r/(16\kown)$, except with probability at most $2\kown N^{-5}$.

Using row~9 ($k\le\lambda^{3.3}$), $\kown\le L\le2\lambda$ and
$N^3\ge\lambda^{309}/8$, the total failure probability of (a) is at most
\[
2(2+k+\kown)N^{-5}\le2(2+\lambda^{3.3}+2\lambda)N^{-5}\le N^{-2}/4 .
\]

\emph{(b)} Condition on a stage-(i)/(ii) outcome in which (a) holds. The
labels are independent of the colouring, so for fixed $(l,j)$ the set
$T_j(Y,l)$ is still a $\rho_l$-random subset of $V(Y)$. The class
$\LJS_{Y,l,j}$ is now a fixed $N$-vertex $(\varepsilon_Y,s)$-expander
with $s:=s_Y/(8k)$ and $\varepsilon_Y\in[2^{-7},1]$. Apply
Theorem~\ref{s3:thmT16s} with $t:=\tJS_l$ and $\rho:=\rho_l$. Its
hypotheses hold:
\begin{itemize}
\item $s\ge2^{135}tL^{28}\rho^{-5}$ by row~5;
\item $\rho N\ge L^2$, since by row~2 at $\lambda=\lambda_r$
  (Lemma~\ref{s3:lemCOLJV}(ii)), $\bar M^4\le\lambda^{31.7}$, so
  $\rho_lN\ge\bar M^{-4}\lambda^{103}/2\ge\lambda^{71}/2\ge4\lambda^2\ge L^2$.
\end{itemize}
Each class fails with probability at most
$2^{86}\tJS_lL^{19}\rho_l^{-3}N^{-3}$. By row~6, the sum over $\IJS(Y)$
is at most $N^{-2}/4$. The final sentence of (b) is
Lemma~\ref{s3:lemMonotone}(iii).

\emph{(c)} Condition as in (b). The sets $V_{l,c,\uslot}$ are independent
of the colouring. For a fixed index, apply Theorem~\ref{s3:thmT16s} to
the fixed class $\LU_{Y,l,c,\uslot}$, an $N$-vertex
$(2^{-6},s_r/(16k))$-expander, with $t:=t_Y=\lceil\lambda^{1.6}\rceil$
(so $1\le t\le2\lambda^{1.6}$) and
$\rho:=\rho_{l,c,\uslot}\ge1/(12L^5)$. Its hypotheses hold:
\begin{itemize}
\item $s_r/(16k)\ge2^{135}tL^{28}\rho^{-5}$ by row~4 (with row~9), since
  $\rho^{-5}\le(12L^5)^5$;
\item $\rho N\ge N/(12L^5)\ge\lambda^{103}/(24(2\lambda)^5)\ge L^2$.
\end{itemize}
Each index fails with probability at most
$2^{86}tL^{19}(12L^5)^3N^{-3}$. The theorem is applied to one index at
a time, so dependence across indices is irrelevant. By row~7, the union
over $\IU(Y)$ is at most $N^{-2}/4$. Together with the failure
probability of (a), (c) fails with probability at most $N^{-2}/2$.

\emph{(e)} The inequalities are row~8 of Table~\ref{s3:tabCOLJV}, with
$s_Y=s_r$ for standalone $Y$. On (a), $\Own_Y$ is a
$(2^{-5},s_r/4)$-expander if $Y$ is standalone, and the own classes are
$(2^{-6},s_r/(16\kown))$-expanders if $Y$ is light. Also $H_Y=X_Y$ is a
$(2^{-6},s_r/2)$-expander for light $Y$, by
Proposition~\ref{s2:propStructure}(i).

\emph{(g)} This holds by construction. If $r\le R-2$, each edge of
$\Lend_Y$ receives exactly one index; if $r\ge R-1$, $\Lend_Y$ has no
classes and is junk for $Y$'s own device. Each edge of $\Own_Y$ receives
exactly one own label.

\emph{Probabilities.} By the bounds above, (a) or (b) fails with
probability at most $N^{-2}/4+N^{-2}/4=N^{-2}/2$, and (e) holds on (a).
Adding the U-union of (c) gives at most $3N^{-2}/4\le N^{-2}$.
\end{proof}

\section{Vortex tools}\label{s4:sec}

This section proves three absorption statements for a single vertex set carrying a spanning robust expander:
Lemma~TPV (the transparent P-vortex, used on standalone pre-parts), Lemma~PV (the four-phase P-vortex, used on light
parts) and Theorem~VX$^+$ (vortex absorption with arbitrary extra edges, used on demoted light parts). All three are
proved from Theorem~16$^*$ (Theorem~\ref{s3:thmT16s}), Lemma~HB
(Lemma~\ref{s3:lemHB}), \BMnames's Corollary~22 (Cited result~\ref{s1:citCor22}) and the elementary long-cycle
bound (Fact~\ref{s1:factEG0}); Lemma~TPV and Theorem~VX$^+$ also use Lemma~15$^+$ (Lemma~\ref{s3:lemL15p}), while
Lemma~PV takes its own classes as given expanders and does not. None of them refers to the hierarchy: each is a statement about one vertex set $Z$, one
expander on $Z$ and edge sets inside $Z$. In each of them the random labels and the good event are fixed
\emph{before} the edge set to be decomposed is revealed, and the conclusion then holds for \emph{every} admissible edge
set; this is the form in which the later sections use them.

Throughout the section, graphs are simple (Convention~\ref{s1:convGraphs}); this is used, for instance, in
Corollary~22, in the distinctness of the far ends $u_1\ne u_2$ of a cherry (two distinct edges at $w$ have distinct
far ends), in the fact that a closed trail in a simple graph has length at least~$3$, in the per-vertex bound
$\deg_{H_0}(x)\le|Z|-1$ of Lemma~PV(c), and in the finishes of the three runs, through Fact~\ref{s1:factEG0}, whose
proof uses that a vertex of degree at least $d$ has at least $d$ neighbours and that a graph on $h$
vertices has at most $\binom h2$ edges.

\subsection{Conventions and elementary observations}\label{s4:ssecConv}

\begin{convention}[vortex-local notation]\label{s4:convVortex}
Inside one run of Lemma~TPV, Lemma~PV or Theorem~VX$^+$ below (a \emph{vortex run}),
$Z$ denotes the vertex set of the run, $N:=|Z|$ and $L:=\log_2|Z|$. All run-local objects are written in sans-serif:
the number of steps $\vJ$; the reservoir colour classes $\vR_{j,c}$ and the reserve class $\vM$; the level $\vlev(v)\in\{0,1,\dots,\vJ\}$
and the label $\vkap(v)$ of a vertex; the sets $\vU_j$ (vertices still active before step $j$), $\vW_j$ (vertices
retiring at step $j$, for $0\le j<\vJ$; the vertices of level $\vJ$ stay in $\vU_{\vJ}$ and are treated by the
finish), $\vZ_j$ (reserve far ends of step $j$) and $\vV_{j,c}$ (connector interiors of step $j$ and
phase $c$); the graphs $\vH_j$ (edges not yet covered before step $j$) and edge sets $\vF_j$ (edges of $\vH_j$ meeting
$\vW_j$); the partition $Z=\vP\sqcup\vQ$ of Lemma~TPV; the reserve sets $\vC^j_w$ (or $\vC_w$); the parameters
$\vm$, $\vb$ (of Lemma~HB) and $\vt$ (the multiplicity used with Theorem~16$^*$); the path families $\vPaths_{j,c}$
(produced by Corollary~22) and $\vPathsQ_{j,c}$ (after moving ends and splitting). The letter $j$ is the step index
and $c$ the phase. The number of colours of a run is written $k$, and the good events of a run are called
(G1)--(G5) (numbered alike in the three runs; a run need not use all five); these letters, the invariants
(Inv1)--(Inv3) and the path properties (P1)--(P4) (and (P5) in Lemma~PV) are local to the run
(the good events (G$k$) are unrelated to the galactic conditions \Gc{k} of Condition~\ref{s1:condGamma}). None of
these symbols has a meaning outside the run in which it is defined. Auxiliary symbols that are local to a
single estimate inside one of the three proofs (such as $\chi$, $\mathcal K_{j,c}$, $V'$ and $I_u$ below) are
written in italic or calligraphic type, not in sans-serif.
\deps{none.}
\end{convention}


The three proofs below share five elementary observations. They are recorded here once, with tags, and proved in
full; they are parts of the proofs of Lemma~TPV, Lemma~PV and Theorem~VX$^+$ below, not
separate statements.

\smallskip\noindent\textbf{Trail splitting}\namedlabel{s4:eqSplit}{\textup{(S)}}~\textup{(S)}.
\emph{Let $T=x_0x_1\cdots x_q$ ($q\ge0$) be a trail (a walk with pairwise distinct edges) in a simple graph and put
$\mathrm{rep}(T):=(q+1)-|\{x_0,\dots,x_q\}|$. Then $E(T)$ is the disjoint union of the edge sets of at most $\mathrm{rep}(T)$ cycles,
each of length at least $3$, and of at most one path; the path is present if and only if $x_0\ne x_q$, and then its
ends are $x_0$ and $x_q$. Every vertex of every piece is a vertex of $T$. In particular, if the inner vertices
$x_1,\dots,x_{q-1}$ are pairwise distinct, then $\mathrm{rep}(T)\le2$ and $T$ splits into at most two cycles and at most one
path, so into at most three pieces.}

\emph{Proof.} Induction on $q$. If $\mathrm{rep}(T)=0$, all $x_i$ are distinct, so $T$ is empty ($q=0$) or a path with ends
$x_0\ne x_q$. Otherwise choose $i<i'$ with $x_i=x_{i'}$ and $i'-i$ minimal. Then $x_i,x_{i+1},\dots,x_{i'-1}$ are pairwise
distinct, and $i'-i\ge3$: $i'-i=1$ would be a loop, and $i'-i=2$ would give $x_ix_{i+1}=x_{i+1}x_{i+2}$, the same
edge twice, which a trail does not contain. So $C:=x_ix_{i+1}\cdots x_{i'}$ is a cycle of length $i'-i\ge3$. Delete it:
$T':=x_0\cdots x_ix_{i'+1}\cdots x_q$ is again a trail (its step from $x_i$ to $x_{i'+1}$ is the edge $x_{i'}x_{i'+1}$ of
$T$), $E(T)=E(C)\sqcup E(T')$, and $T'$ begins at $x_0$ and ends at $x_q$ (if $i'=q$ it ends at $x_i=x_q$). The
vertices $x_{i+1},\dots,x_{i'-1}$ are pairwise distinct and differ from $x_i$, so at most $i'-i-1$ distinct vertices
disappear while $i'-i$ entries are removed; hence $\mathrm{rep}(T')\le\mathrm{rep}(T)-1$, and the induction hypothesis applied to
$T'$ finishes the proof. If $x_1,\dots,x_{q-1}$ are pairwise distinct, then $|\{x_0,\dots,x_q\}|\ge q-1$, so
$\mathrm{rep}(T)\le2$. \hfill$\lozenge$

\smallskip\noindent\textbf{Path ends}\namedlabel{s4:eqEnds}{\textup{(E)}}~\textup{(E)}.
\emph{Let $\mathcal P$ be a decomposition of an edge set $F$ into (non-trivial) paths. For every vertex $v$ the number
of paths of $\mathcal P$ having $v$ as an end is congruent to $\deg_F(v)$ modulo $2$. If every vertex is an end of at
most two paths of $\mathcal P$ and all edges of $F$ have both ends in a set $W$, then $|\mathcal P|\le|W|$.}

\emph{Proof.} A path contributes $2$ to the degree of each of its inner vertices and $1$ to the degree of each of its
two (distinct) ends. For the second part count path ends: $2|\mathcal P|\le2|W|$. \hfill$\lozenge$

\noindent By Cited result~\ref{s1:citCor22}, every (simple) graph, in particular every edge set $F$, has a
decomposition $\mathcal P$ as in (E) in which every vertex is an end of at most two paths; we call it a
\emph{Corollary-22 decomposition} of $F$.

\smallskip\noindent\textbf{Monotone comparison}\namedlabel{s4:eqMono}{\textup{(MC)}}~\textup{(MC)}.
\emph{Let $S$ be a finite set and let $\mathcal K$ be a family of subsets of $S$ that is closed upwards: $T\in\mathcal K$
and $T\subseteq T'\subseteq S$ imply $T'\in\mathcal K$. For $r=(r_x)_{x\in S}\in[0,1]^S$ put
\[
\Psi_{\mathcal K}(r):=\sum_{T\in\mathcal K}\ \prod_{x\in T}r_x\prod_{x\in S\setminus T}(1-r_x).
\]
If $p,q\in[0,1]^S$ satisfy $p_x\le q_x$ for every $x\in S$, then $\Psi_{\mathcal K}(p)\le\Psi_{\mathcal K}(q)$. If $V$
is a random subset of $S$ such that the events $\{x\in V\}$ \textup{(}$x\in S$\textup{)} are mutually independent, then
$\Prob(V\in\mathcal K)=\Psi_{\mathcal K}(q^V)$, where $q^V\in[0,1]^S$ is given by $q^V_x:=\Prob(x\in V)$. In
particular, let $\rho\in[0,1]$, let $V$ be a random subset of $S$ such that the events $\{x\in V\}$
\textup{(}$x\in S$\textup{)} are mutually independent and $\Prob(x\in V)\ge\rho$ for every $x\in S$, and let $V'$ be a
$\rho$-random subset of $S$, that is, a random subset of $S$ that contains each element of $S$ independently with
probability $\rho$. Then $\Prob(V\notin\mathcal K)\le\Prob(V'\notin\mathcal K)$ \textup{(}$V$ and $V'$ need not be
defined on the same probability space; only the two probabilities are compared\textup{)}.}

\emph{Proof.} For the comparison write $S=\{x_1,\dots,x_{|S|}\}$, and for $0\le i\le|S|$ let $q^{(i)}\in[0,1]^S$ agree
with $q$ at $x_1,\dots,x_i$ and with $p$ at $x_{i+1},\dots,x_{|S|}$. Then $q^{(0)}=p$ and $q^{(|S|)}=q$, and for $1\le i\le|S|$
the vectors $q^{(i-1)}$ and $q^{(i)}$ agree except possibly at $x_i$, where
$q^{(i-1)}_{x_i}=p_{x_i}\le q_{x_i}=q^{(i)}_{x_i}$. So it suffices to show that $\Psi_{\mathcal K}$ is non-decreasing
in each coordinate when the other coordinates are fixed. Fix $y\in S$ and values $r_x\in[0,1]$
($x\in S\setminus\{y\}$), and for $z\in[0,1]$ let $r(z)\in[0,1]^S$ have entry $r_x$ at $x\ne y$ and entry $z$ at $y$.
Every subset of $S$ is, for exactly one $T\subseteq S\setminus\{y\}$, equal to $T$ or to $T\cup\{y\}$. With
$\omega(T):=\prod_{x\in T}r_x\prod_{x\in S\setminus(T\cup\{y\})}(1-r_x)\ge0$ this gives
\begin{align*}
\Psi_{\mathcal K}(r(z))&=\sum_{T\subseteq S\setminus\{y\}}\omega(T)\Bigl((1-z)\,\ind{T\in\mathcal K}+z\,\ind{T\cup\{y\}\in\mathcal K}\Bigr)\\
&=\sum_{T\subseteq S\setminus\{y\}}\omega(T)\Bigl(\ind{T\in\mathcal K}+z\bigl(\ind{T\cup\{y\}\in\mathcal K}-\ind{T\in\mathcal K}\bigr)\Bigr).
\end{align*}
As $\mathcal K$ is closed upwards, $\ind{T\cup\{y\}\in\mathcal K}\ge\ind{T\in\mathcal K}$ for every $T$, so every
summand, and hence $\Psi_{\mathcal K}(r(z))$, is a non-decreasing function of $z$. For the formula, let $V$ be as
stated. For $T\subseteq S$ the event $\{V=T\}$ is the intersection of the events $\{x\in V\}$ ($x\in T$) and
$\{x\notin V\}$ ($x\in S\setminus T$); these events are mutually independent, because replacing some of finitely many
mutually independent events by their complements preserves mutual independence. So
$\Prob(V=T)=\prod_{x\in T}q^V_x\prod_{x\in S\setminus T}(1-q^V_x)$, and summing over the pairwise disjoint events
$\{V=T\}$, $T\in\mathcal K$, gives $\Prob(V\in\mathcal K)=\Psi_{\mathcal K}(q^V)$. For the last claim, the events
$\{x\in V'\}$ \textup{(}$x\in S$\textup{)} are mutually independent and $q^{V'}_x=\rho\le q^V_x$ for every $x\in S$; so,
by the formula and the comparison,
$\Prob(V'\in\mathcal K)=\Psi_{\mathcal K}(q^{V'})\le\Psi_{\mathcal K}(q^V)=\Prob(V\in\mathcal K)$. Taking complements
gives the claim. \hfill$\lozenge$

\smallskip\noindent\textbf{Binomial domination}\namedlabel{s4:eqDom}{\textup{(B)}}~\textup{(B)}.
\emph{Let $X=\sum_{i=1}^{m'}I_i$ be a sum of independent indicator variables with $\Prob(I_i=1)\ge p$ for every $i$,
and let $a>0$ with $m'p\ge2a$. Then $\Prob(X<a)\le\exp(-m'p/8)$.}

\emph{Proof.} As $m'p\ge2a>0$, we have $m'\ge1$ and $0<p\le\Prob(I_1=1)\le1$. The set $V:=\{i\in[m']:I_i=1\}$ is a
random subset of $[m']$ whose events $\{i\in V\}=\{I_i=1\}$ \textup{(}$i\in[m']$\textup{)} are mutually independent,
each of probability at least $p$, and $X=|V|$. The family $\mathcal K:=\{T\subseteq[m']:|T|\ge a\}$ is closed upwards.
Let $V'$ be a $p$-random subset of $[m']$; then $|V'|\sim\Bin(m',p)$, with mean $\mu=m'p\ge2a$. By observation
\textup{(MC)} (with $S:=[m']$ and $\rho:=p$), $\Prob(X<a)=\Prob(V\notin\mathcal K)\le\Prob(V'\notin\mathcal
K)=\Prob(|V'|<a)$, and by the lower tail of Cited result~\ref{s1:citChernoff} with $\delta=1/2$,
$\Prob(|V'|<a)\le\Prob(|V'|<\mu/2)\le\exp(-\mu/8)$. \hfill$\lozenge$

\smallskip\noindent\textbf{Closing}\namedlabel{s4:eqClose}{\textup{(C)}}~\textup{(C)}.
\emph{Let $Q$ be a path of length at least $2$ with ends $x\ne y$, and let $Q'$ be an $x$--$y$ path whose inner
vertices avoid $V(Q)$ and which shares no edge with $Q$. Then $Q\cup Q'$ is a cycle of length at least $3$.}

\emph{Proof.} The only common vertices of $Q$ and $Q'$ are $x$ and $y$, so $Q\cup Q'$ is a cycle; its length is
$|E(Q)|+|E(Q')|\ge2+1$. \hfill$\lozenge$

\smallskip\noindent\textbf{Finite probability spaces.} Every random object of this section takes finitely many
values. The labels of a run are finitely many independent random variables, each with finitely many values: in
Lemma~TPV and Theorem~VX$^+$ a uniform colouring of $E(O)$ with $k$ colours, and in all three statements the levels,
with values in $\{0,1,\dots,\vJ\}$, and the labels $\vkap$. The $\rho$-random subsets of a finite set $S$ used with
\textup{(MC)} and \textup{(B)} may be taken to be the identity map on the finite probability space of all subsets
$T\subseteq S$, where $T$ has probability $\prod_{x\in T}\rho\prod_{x\in S\setminus T}(1-\rho)$ (an empty product is
$1$); identifying each $T$ with its indicator vector, this space is the product $\{0,1\}^S$ of $|S|$ copies of the
two-point space $\{0,1\}$ in which $1$ has probability $\rho$. So every probability and expectation in this section is
a finite sum over a finite product of finite probability spaces. Where a proof below fixes the colouring of a run,
the probability over the remaining labels, which are independent of the colouring, is a finite sum as well, and the
unconditional probability is the average of these sums over the finitely many (equally likely) colourings.

\smallskip
The absolute constant $\Nzero$ (Definition~\ref{s1:defConstants}) is chosen so large that the finitely many
explicit inequalities listed as \emph{size conditions} in the three proofs of this section hold for every
$N\ge\Nzero$; each of them is an inequality between explicit functions of $N$ that is true for all sufficiently
large $N$. In every application in this paper the vertex set of a vortex run
is a pre-part or a light part, which has at least $\Nzero$ vertices by condition \ref{s1:condG3} of Condition~\ref{s1:condGamma}.

\subsection{Lemma TPV: the transparent P-vortex}\label{s4:ssecTPV}

\begin{lemma}[Lemma TPV (transparent P-vortex)]\label{s4:lemTPV}
Let $Z$ be a set of $N\ge\Nzero$ vertices and $L:=\log_2N$. Let $O$ be a spanning $(\eps_O,s)$-expander on $Z$ with
$\eps_O\in[2^{-7},2^{-5}]$ (only $\eps_O\ge2^{-7}$ is used in the proof) and $s\ge2^{150}L^{42}$, and let $Z=\vP\sqcup\vQ$ be a partition. Put
\[
\vJ:=\lfloor\log_2(L/6)\rfloor,\qquad k:=3\vJ+1,\qquad \vm:=\lceil L^6\rceil,\qquad
\vb:=\lceil2^8L^2\vm\rceil,\qquad \vt:=2^{10}L^8,
\]
($\vb$ is the bound of Lemma~HB at the worst case $\eps_O=2^{-7}$),
and fix, as a function of $O$ only, sets $\Aw(w)\subseteq N_O(w)$ ($w\in Z$) with $|\Aw(w)|=\vm$ such that every vertex
lies in at most $\vb$ of them (they exist by Lemma~\ref{s3:lemHB}; see the proof). The \emph{labels} of the run are
the following independent random variables:
\begin{itemize}
\item[(a)] a uniform colouring of $E(O)$ with the $k$ colours $\vR_{j,c}$ ($0\le j<\vJ$, $c\in[3]$) and $\vM$; the
  colour classes are regarded as spanning subgraphs of $O$ on $Z$ and denoted by the same letters;
\item[(b)] for every $v\in\vP$ a level $\vlev(v)\in\{0,1,\dots,\vJ\}$ with $\Prob(\vlev(v)\ge j)=2^{-j}$ for
  $0\le j\le\vJ$ (that is, $\Prob(\vlev(v)=j)=2^{-(j+1)}$ for $0\le j<\vJ$ and $\Prob(\vlev(v)=\vJ)=2^{-\vJ}$);
\item[(c)] for every $v\in Z$ a label $\vkap(v)\in\{0,1,2,3\}$ with $\Prob(\vkap(v)=0)=1/2$ and
  $\Prob(\vkap(v)=c)=1/6$ for $c\in[3]$.
\end{itemize}
Put $\vU_j:=\vQ\cup\{v\in\vP:\vlev(v)\ge j\}$, $\vW_j:=\{v\in\vP:\vlev(v)=j\}$,
$\vZ_j:=\{u\in\vU_{j+1}:\vkap(u)=0\}$ and $\vV_{j,c}:=\{u\in\vU_{j+1}:\vkap(u)=c\}$. There is an event
$\mathcal G_{\mathrm{TPV}}$, determined by $(Z,O,\vP)$ and the labels only, with
\begin{equation}\label{s4:eqTPVprob}
\Prob(\mathcal G_{\mathrm{TPV}})\ \ge\ \frac12-\eta_{\mathrm{TPV}}(N),\qquad
\eta_{\mathrm{TPV}}(N):=2LN^{-5}+2^{96}L^{31}N^{-3}+NL\,e^{-3L^4/8}\ \le\ \frac1{100},
\end{equation}
such that on $\mathcal G_{\mathrm{TPV}}$ the following holds for \emph{every} edge set $E$ with $E(O)\subseteq E$ all
of whose edges have both ends in $Z$: there is a partition $E=\EV\sqcup\EQ$ with the properties
\begin{itemize}
\item[\textup{(T1)}]\namedlabel{s4:T1}{\textup{(T1)}} every edge of $E$ with both ends in $\vP$ lies in $\EV$;
\item[\textup{(T2)}]\namedlabel{s4:T2}{\textup{(T2)}} every edge of $\EV$ with an end in $\vQ$ is an edge of $O$;
\item[\textup{(T3)}]\namedlabel{s4:T3}{\textup{(T3)}} $\EV$ decomposes into at most $169|\vP|$ objects, and
  $\EV=\emptyset$ if $\vP=\emptyset$;
\item[\textup{(T4)}]\namedlabel{s4:T4}{\textup{(T4)}} (no arcs) all objects of (T3) consist of edges of $E$; every
  path formed during the run is closed into a cycle within the step in which it is formed, so the run leaves no path
  (``arc'') to be closed by any other device.
\end{itemize}
The partition and the decomposition are obtained from $E$ and the labels by a deterministic procedure.
\deps{\ref{s3:lemL15p}, \ref{s3:thmT16s}, \ref{s3:lemHB}, \ref{s3:lemMonotone}, \ref{s1:citCor22},
\ref{s1:factEG0}, \ref{s1:citChernoff}, \ref{s1:citMarkov}, \ref{s1:condG3}, \ref{s4:convVortex}, \ref{s1:convGraphs}, \ref{s1:defConstants}.}
\end{lemma}

\begin{proof}
If $\vP=\emptyset$, let $\mathcal G_{\mathrm{TPV}}$ be the sure event and put $\EV:=\emptyset$, $\EQ:=E$; then
(T1) is vacuous and (T2)--(T4) hold trivially. From now on $\vP\ne\emptyset$.

\emph{Size conditions.} We use: (i) $L\ge2^{10}$; (ii) $N\ge L^3$; (iii) $4\cdot48\le L$ (a consequence of (i)); (iv)
$\eta_{\mathrm{TPV}}(N)\le1/100$. They hold for $N\ge\Nzero$ by the choice of $\Nzero$.

\emph{Parameters.} By the definition of $\vJ$, $6\cdot2^{\vJ}\le L<12\cdot2^{\vJ}$, so
\begin{equation}\label{s4:eqTPVJ}
2^{-\vJ}\ge 6/L,\qquad 2^{-\vJ}<12/L,\qquad 1\le\vJ\le\log_2L,\qquad k=3\vJ+1\le L
\end{equation}
(the last two by (i)). Since $s\ge2^{150}L^{42}\ge2\lceil L^6\rceil=2\vm$, Lemma~\ref{s3:lemHB} applied to $X:=O$,
$\eps':=\eps_O\ge2^{-7}$ and $m:=\vm$ gives sets $\Aw(w)\subseteq N_O(w)$, $|\Aw(w)|=\vm$, such that every vertex lies in
at most $\lceil2\vm L^2/\eps_O\rceil\le\lceil2^8L^2\vm\rceil=\vb$ of them; we fix the lexicographically first such
family (for a fixed ordering of $Z$), so that it is a function of $O$. Moreover
\begin{equation}\label{s4:eqTPVt}
2+\vb\ \le\ 2+2^8L^2(L^6+1)+1\ \le\ 2^{10}L^8=\vt .
\end{equation}
Directly from the definitions: $\vU_0=Z$; $\vU_{j+1}\subseteq\vU_j$ and $\vU_j=\vU_{j+1}\sqcup\vW_j$;
$\vW_j\subseteq\vP$ and $\vW_j\cap\vU_{j+1}=\emptyset$; and $\vZ_j,\vV_{j,1},\vV_{j,2},\vV_{j,3}$ partition
$\vU_{j+1}$.

\emph{Good events.} Let $\mathcal G_{\mathrm{TPV}}$ be the intersection of the following four events.
\begin{itemize}
\item[(G1)] Every colour class ($\vR_{j,c}$ and $\vM$) is a spanning $(\eps_O,s/(2k))$-expander on $Z$. Since
  $s\ge2^{150}L^{42}\ge40kL$, Lemma~\ref{s3:lemL15p} (with $X:=O$ and $k$ colours) shows that each class fails with
  probability at most $2N^{-5}$; so (G1) fails with probability at most $2kN^{-5}\le2LN^{-5}$. Note that
  $s/(2k)\ge s/(2L)\ge2^{149}L^{41}$.
\item[(G2)] For all $0\le j<\vJ$ and $c\in[3]$, the class $\vR_{j,c}$ is $(2^{12}L^4,\vt)$-path connected through
  $\vV_{j,c}$ (for every multiset of pairs of distinct vertices of $Z$ in which each vertex lies in at most $\vt$
  pairs there are edge-disjoint paths of $\vR_{j,c}$ of length at most $2^{12}L^4$ joining the pairs, with all inner
  vertices in $\vV_{j,c}$). Fix a colouring $\chi$ in (G1) and a pair $(j,c)$, and let $\mathcal K_{j,c}$ be the
  family of sets $T\subseteq Z$ such that $\vR_{j,c}$ is $(2^{12}L^4,\vt)$-path connected through $T$. It is
  determined by $\chi$. It is closed upwards by Lemma~\ref{s3:lemMonotone}(ii), applied with $G=G':=\vR_{j,c}$,
  $\ell=\ell':=2^{12}L^4$, $t=t':=\vt$ and with $T\subseteq T'\subseteq Z=V(\vR_{j,c})$ as the two sets through
  which path connectivity is taken (the colour classes are spanning). For $v\in Z$ the event $v\in\vV_{j,c}$
  depends only on $\vkap(v)$ and, if $v\in\vP$, on $\vlev(v)$, so these events are mutually independent over $v$
  and independent of the colouring, and
  \[
  \Prob(v\in\vV_{j,c})=\begin{cases}1/6\ \ge\ 1/L,& v\in\vQ,\\ 2^{-(j+1)}/6\ \ge\ 2^{-\vJ}/6\ \ge\ 1/L,& v\in\vP,\end{cases}
  \]
  by \eqref{s4:eqTPVJ} and size condition (i). Let $V'$ be a $(1/L)$-random subset of $Z$ (each vertex independently
  with probability exactly $1/L$), and apply Theorem~\ref{s3:thmT16s} to $X:=\vR_{j,c}$ (an $N$-vertex
  $(\eps_O,s/(2k))$-expander with $\eps_O\in[2^{-7},1]$), $V:=V'$, $\rho:=1/L$ and $t:=\vt$: its hypotheses
  $\rho N=N/L\ge L^2$ (by size condition (ii)) and
  $s/(2k)\ge2^{149}L^{41}\ge2^{135}\cdot2^{10}L^8\cdot L^{28}\cdot L^{5}=2^{135}\vt L^{28}\rho^{-5}$ hold. So
  $\Prob(V'\notin\mathcal K_{j,c})\le2^{86}\vt L^{19}\rho^{-3}N^{-3}=2^{96}L^{30}N^{-3}$. By observation
  \textup{(MC)} (with $S:=Z$, $\mathcal K:=\mathcal K_{j,c}$, the random set $\vV_{j,c}$ in place of $V$, and
  $\rho:=1/L$), also $\Prob(\vV_{j,c}\notin\mathcal K_{j,c})\le2^{96}L^{30}N^{-3}$, the probability being over the
  labels (b) and (c): for the fixed colouring $\chi$, $\vR_{j,c}$ fails to be $(2^{12}L^4,\vt)$-path connected
  through $\vV_{j,c}$ with probability at most $2^{96}L^{30}N^{-3}$. As the labels (b), (c) are independent of the
  colouring, a union bound over the $3\vJ$ pairs $(j,c)$ shows that the probability that (G1) holds and (G2) fails
  is at most $\sum_{\chi}\Prob(\text{the colouring is }\chi)\cdot3\vJ\cdot2^{96}L^{30}N^{-3}\le2^{96}L^{31}N^{-3}$,
  the sum running over the colourings $\chi$ in (G1), and using $3\vJ\le L$. By
  Lemma~\ref{s3:lemMonotone}(iii) the event (G2) depends only on the classes $\vR_{j,c}$ and the sets $\vV_{j,c}$,
  and on it the pair multisets may be chosen after all labels are revealed, as they will be below.
\item[(G4)] For every $0\le j<\vJ$ and every $w\in\vP$ the set
  $\vC^j_w:=\{wu:\ u\in \Aw(w),\ wu\in\vM,\ u\in\vZ_j\}$ has at least $6$ elements. For $u\in \Aw(w)$ let
  $I_u:=\ind{wu\in\vM\text{ and }u\in\vZ_j}$. Since $u\ne w$, the variables $I_u$ ($u\in \Aw(w)$) depend on pairwise
  distinct edges $wu$ and distinct vertices $u$, so they are independent; and
  $\Prob(I_u=1)=\frac1k\cdot\Prob(u\in\vU_{j+1})\cdot\frac12\ge\frac1L\cdot2^{-\vJ}\cdot\frac12\ge\frac3{L^2}$ by
  \eqref{s4:eqTPVJ} (for $u\in\vQ$, $\Prob(u\in\vU_{j+1})=1$). As $\vm\cdot3/L^2\ge3L^4\ge12$, observation
  \textup{(B)} gives $\Prob(|\vC^j_w|<6)\le e^{-3L^4/8}$. A union bound over the at most $N\vJ\le NL$ pairs $(w,j)$
  bounds the failure probability of (G4) by $NLe^{-3L^4/8}$.
\item[(G5)] $\sum_{j<\vJ}|\vP\cap\vU_j|\le8|\vP|$ and $|\vP\cap\vU_{\vJ}|\le48|\vP|/L$. Here
  $\Ex|\vP\cap\vU_j|=2^{-j}|\vP|$, so $\Ex\sum_{j<\vJ}|\vP\cap\vU_j|<2|\vP|$ and, by \eqref{s4:eqTPVJ},
  $\Ex|\vP\cap\vU_{\vJ}|<12|\vP|/L$. By Markov's inequality (Cited result~\ref{s1:citMarkov}) each of the two
  inequalities fails with probability less than $1/4$.
\end{itemize}
All four events are determined by $(Z,O,\vP)$ (through $O$, the sets $\Aw(w)$ and the partition) and the labels; none
involves $E$. Summing the failure probabilities gives \eqref{s4:eqTPVprob}. From now on we fix labels in
$\mathcal G_{\mathrm{TPV}}$ and an arbitrary edge set $E\supseteq E(O)$ inside $Z$.

\emph{The process.} All choices below (reserved edges, classes, Corollary-22 decompositions, connectors, the final
decomposition) are made by fixed rules, so the procedure is deterministic given $E$ and the labels. An edge is
\emph{used} once it has been placed into an output object. Put $\vH_0:=\{e\in E:\ e\text{ has both ends in }\vP\}$
and, for $j\ge0$, $\vH_{j+1}:=\vH_j\setminus\{\text{edges used at step }j\}$; thus $\vH_j$ is the set of unused edges
of $\vH_0$. Before step $j$ ($0\le j\le\vJ$) we maintain:
\begin{itemize}
\item[(Inv1)] every edge of $\vH_j$ has both ends in $\vP\cap\vU_j$;
\item[(Inv2)] no edge of $\vR_{j',c}$ with $j'\ge j$, $c\in[3]$ and both ends in $\vU_{j'+1}$ has been used;
\item[(Inv3)] no $\vM$-edge $wu$ with $w\in\vW_{j'}$ for some $j'\ge j$ and $u\in\vU_{j'+1}$ has been used.
\end{itemize}
For $j=0$ nothing has been used and $\vU_0=Z$, so (Inv1)--(Inv3) hold.

\emph{Step $j$} ($0\le j<\vJ$). Let $\vF_j$ be the set of edges of $\vH_j$ meeting $\vW_j$. By (Inv1), every edge of
$\vF_j$ has both ends in $\vP\cap\vU_j$.
\begin{itemize}
\item[(i)] \emph{Reserve.} For each $w\in\vW_j$ choose six distinct edges $e^{c,i}_w\in\vC^j_w$ ($c\in[3]$,
  $i\in[2]$); this is possible by (G4), since $w\in\vP$. Each $e^{c,i}_w=wu$ is an $\vM$-edge with $w\in\vW_j$ and
  $u\in\vZ_j\subseteq\vU_{j+1}$, so it is unused by (Inv3), and $w$ is its only end in $\vW_j$; hence reserved
  edges of distinct $w$ are distinct. If $u\in\vP$ then $e^{c,i}_w\in\vH_0$ (as $E(O)\subseteq E$), so
  $e^{c,i}_w\in\vH_j$ and $e^{c,i}_w\in\vF_j$; if $u\in\vQ$ it is an edge of $O$ outside $\vH_0$. The $\vZ_j$-end $u$
  of a reserved edge is called its \emph{far end}.
\item[(ii)] \emph{Classes.} For every non-reserved edge of $\vF_j$ choose an end $w\in\vW_j$ (either one if both ends
  lie in $\vW_j$), call the other end $x$, and give the edge a class $\gamma\in[3]\setminus\{\vkap(x)\}$ (possible as
  $\vkap(x)$ is a single value). Let $\vF_{j,c}$ be the set of non-reserved edges of class $c$. Then
  \begin{equation}\label{s4:eqTPVclass}
  V(\vF_{j,c})\cap\vV_{j,c}=\emptyset\qquad(c\in[3]):
  \end{equation}
  an end in $\vW_j$ is not in $\vU_{j+1}\supseteq\vV_{j,c}$, and the other end $x$ of a class-$c$ edge has
  $\vkap(x)\ne c$.
\item[(iii)] \emph{Paths.} Let $\vPaths_{j,c}$ be a Corollary-22 decomposition of $\vF_{j,c}$. As
  $V(\vF_{j,c})\subseteq\vP\cap\vU_j$, observation \textup{(E)} gives $|\vPaths_{j,c}|\le|\vP\cap\vU_j|$.
\item[(iv)] \emph{Ends in $\vW_j$.} For $w\in\vW_j$ and $c\in[3]$ let $d\in\{0,1,2\}$ be the number of paths of
  $\vPaths_{j,c}$ having $w$ as an end.
  \begin{itemize}
  \item If $d=2$, append $e^{c,1}_w$ at $w$ to one of these paths and $e^{c,2}_w$ at $w$ to the other.
  \item If $d=1$, append $e^{c,1}_w$ at $w$ to that path and output $e^{c,2}_w$ as a single edge.
  \item If $d=0$, form the \emph{cherry} $u_1wu_2$ from $e^{c,1}_w=wu_1$ and $e^{c,2}_w=wu_2$ ($u_1\ne u_2$, since the
    two edges are distinct).
  \end{itemize}
  This produces, for each $c$, a family of \emph{trails}: every path of $\vPaths_{j,c}$ with exactly one reserved
  edge appended at each of its ends lying in $\vW_j$ (and nothing appended at its other ends), and the cherries.
  If a path has both ends in $\vW_j$, the two appended edges belong to distinct vertices of $\vW_j$ and are distinct.
  In every trail $x_0x_1\cdots x_q$ the inner vertices $x_1,\dots,x_{q-1}$ are vertices of one path of
  $\vPaths_{j,c}$ (or the centre $w$ of a cherry), hence pairwise distinct, and the edges are pairwise distinct.
\item[(v)] \emph{Splitting.} By observation \textup{(S)} every trail splits into at most two cycles (which are output)
  and at most one path, with the same two ends as the trail; let $\vPathsQ_{j,c}$ be the set of these paths. They
  satisfy:
  \begin{itemize}
  \item[(P1)] no vertex of a path in $\vPathsQ_{j,c}$ lies in $\vV_{j,c}$: its vertices are vertices of $\vF_{j,c}$
    (by \eqref{s4:eqTPVclass}), cherry centres (in $\vW_j$) or far ends (in $\vZ_j$, where $\vkap=0$);
  \item[(P2)] the two ends of a path in $\vPathsQ_{j,c}$ are distinct and lie in $\vU_{j+1}$: an end of a trail is
    either an end, not in $\vW_j$, of a path of $\vPaths_{j,c}$, hence in $(\vP\cap\vU_j)\setminus\vW_j\subseteq
    \vU_{j+1}$, or a far end, in $\vZ_j\subseteq\vU_{j+1}$; distinctness is part of \textup{(S)};
  \item[(P3)] every path in $\vPathsQ_{j,c}$ has length at least $2$: each of its edges meets $\vW_j$ (edges of
    $\vF_j$ by definition, reserved edges at $w$), while by (P2) neither of its ends lies in $\vW_j$, so no single
    edge can join its two ends;
  \item[(P4)] every vertex $x$ is an end of at most $2+\vb\le\vt$ paths of $\vPathsQ_{j,c}$: each trail yields at
    most one path, whose ends are ends of the trail; $x$ is an end of at most two paths of $\vPaths_{j,c}$; and $x$ is
    the far end of a reserved edge $e^{c,i}_w$ only if $x\in \Aw(w)$, which holds for at most $\vb$ vertices $w$, while
    for each such $w$ at most one of the two distinct edges $e^{c,1}_w,e^{c,2}_w$ at $w$ has far end $x$. Then use
    \eqref{s4:eqTPVt}.
  \end{itemize}
\item[(vi)] \emph{Closing.} For each $c\in[3]$ the ends of the paths of $\vPathsQ_{j,c}$ form a multiset of pairs of
  distinct vertices in which every vertex lies in at most $\vt$ pairs (P2, P4). By (G2) there are pairwise
  edge-disjoint paths $Q'\subseteq\vR_{j,c}$, one for each $Q\in\vPathsQ_{j,c}$, joining the two ends of $Q$, of length
  at most $2^{12}L^4$, with all inner vertices in $\vV_{j,c}$. All vertices of $Q'$ lie in $\vU_{j+1}$ (ends by (P2),
  inner vertices in $\vV_{j,c}\subseteq\vU_{j+1}$), so every edge of $Q'$ is an $\vR_{j,c}$-edge with both ends in
  $\vU_{j+1}$; it is unused by (Inv2), it is not in $\vF_j$ (which meets $\vW_j$), and it is not reserved (reserved
  edges are $\vM$-edges). The inner vertices of $Q'$ avoid $V(Q)$ by (P1), and $Q$ has length at least $2$ by (P3);
  so by observation \textup{(C)}, $Q\cup Q'$ is a cycle, which is output.
\end{itemize}
The edges used at step $j$ are: all of $\vF_j$ (every non-reserved edge lies on a path of some $\vPaths_{j,c}$,
every reserved edge is appended, in a cherry or a single edge, and every trail edge ends in a split cycle or in a
closed cycle), the reserved edges with far end in $\vQ$, and the connector edges. Within the step the output objects
are pairwise edge-disjoint: trails of one class are edge-disjoint and classes are disjoint; connectors of one class
are edge-disjoint by (G2), connectors of distinct classes lie in distinct colour classes, and connectors meet neither
$\vF_j$ nor the reserved edges. No edge used at step $j$ was used before: $\vF_j\subseteq\vH_j$, reserved edges by
(Inv3), connectors by (Inv2).

\emph{The invariants persist.} (Inv1) for $j+1$: an edge of $\vH_{j+1}$ is an edge of $\vH_j$ not in $\vF_j$, so it
has both ends in $(\vP\cap\vU_j)\setminus\vW_j=\vP\cap\vU_{j+1}$. (Inv2), (Inv3) for $j+1$: step $j$ used only edges
meeting $\vW_j$ and $\vR_{j,\cdot}$-edges. An edge protected by (Inv2) or (Inv3) for an index $j'\ge j+1$ has both
ends in $\vU_{j'}\subseteq\vU_{j+1}$ (in (Inv3), $w\in\vW_{j'}\subseteq\vU_{j'}$), so it does not meet $\vW_j$, and
it is an $\vR_{j',c}$-edge with $j'\ne j$ or an $\vM$-edge; so it was not used at step $j$.

\emph{Cost of step $j$.} There are at most $3|\vW_j|$ single edges (at most one per pair $(w,c)$). For each $c$ there
are at most $|\vPaths_{j,c}|+|\vW_j|\le|\vP\cap\vU_j|+|\vW_j|$ trails, each giving at most three objects (at most two
split cycles and at most one closed cycle). So step $j$ outputs at most
\begin{equation}\label{s4:eqTPVstep}
3|\vW_j|+3\cdot3\bigl(|\vP\cap\vU_j|+|\vW_j|\bigr)\ =\ 12\,|\vW_j|+9\,|\vP\cap\vU_j|
\end{equation}
objects. The sets $\vW_j$ ($j<\vJ$) are pairwise disjoint subsets of $\vP$, so $\sum_{j<\vJ}|\vW_j|\le|\vP|$, and
$\sum_{j<\vJ}|\vP\cap\vU_j|\le8|\vP|$ by (G5). Hence, by \eqref{s4:eqTPVstep}, the steps $0\le j<\vJ$ together
output at most $12|\vP|+9\cdot8|\vP|=84|\vP|$ objects.

\emph{Finish.} By (Inv1) for $j=\vJ$, every edge of $\vH_{\vJ}$ has both ends in $\vP\cap\vU_{\vJ}$, and
$|\vP\cap\vU_{\vJ}|\le48|\vP|/L$ by (G5). Apply Fact~\ref{s1:factEG0}(b) with $F:=\vH_{\vJ}$, $W:=\vP\cap\vU_{\vJ}$,
$\beta:=48$ and $\alpha:=|\vP|\le N$; its hypothesis $N>1$ holds as $L\ge2^{10}$ by size condition (i), and its hypothesis $4\beta\le L$ is size
condition (iii). So $\vH_{\vJ}$ decomposes into at most
$48|\vP|$ objects, which are output. The total number of objects is at most $84|\vP|+48|\vP|=132|\vP|\le169|\vP|$.

\emph{Conclusion.} Let $\EV$ be the set of used edges and $\EQ:=E\setminus\EV$. Every used edge lies in $\vH_0\subseteq
E$ or is a reserved or connector edge, which lies in $E(O)\subseteq E$; so $\EV\subseteq E$ and the output is a
decomposition of $\EV$ into at most $169|\vP|$ objects, which is (T3).
(T1): every edge of $E$ with both ends in $\vP$ lies in $\vH_0$, and every edge of $\vH_0$ is used (edges leave
$\vH_j$ only by being used, and $\vH_{\vJ}$ is used in the finish).
(T2): an edge of $\EV$ with an end in $\vQ$ is not in $\vH_0$, so it is a reserved edge or a connector edge, hence an
edge of $O$.
(T4): all vertices of the Corollary-22 paths lie in $\vP$ (their edges lie in $\vH_0$), every path of
$\vPathsQ_{j,c}$ is closed in step $j$, and the finish outputs only cycles and single edges; no path is left open.

\emph{Remark.} Property (P4) bounds the pair multiplicity by $2+\vb$ (the cruder bound $2+2\vb$ is also below
$\vt$, so nothing changes quantitatively). Property (P3) is what makes $Q\cup Q'$ a cycle of length at least $3$; the
case of a single edge closed by a single connector edge cannot occur.
\end{proof}


\subsection{Lemma PV: the four-phase P-vortex}\label{s4:ssecPV}

In Lemma~TPV every path is closed inside the run. In the P-vortex a set $\Rt$ of \emph{rooted} vertices never
retires, and paths whose two ends are rooted are not closed inside the run: they are returned as \emph{arcs}, each
carrying one of four phases, to be closed later by a different device through sets of vertices that are marked by an
external phase map $\extph$. The only property of $\extph$ that the lemma uses is the hypothesis (E1$'$) below.

\begin{lemma}[Lemma PV (the four-phase P-vortex)]\label{s4:lemPV}
Let $Z$ be a set of $N\ge\Nzero$ vertices and $L:=\log_2N$. Put
\[
\vJ:=\lfloor\log_2(L/8)\rfloor,\qquad \vm:=\lceil L^6\rceil,\qquad \vb:=\lceil2^7L^2\vm\rceil,\qquad
\vt:=2^9L^8 .
\]
The following data are fixed (not random):
\begin{itemize}
\item a spanning $(2^{-6},s_O)$-expander $O$ on $Z$ with $2\vm\le s_O$, and sets $\Aw(w)\subseteq N_O(w)$ ($w\in Z$)
  with $|\Aw(w)|=\vm$ such that every vertex lies in at most $\vb$ of them (such sets exist by
  Lemma~\ref{s3:lemHB} with $\eps'=2^{-6}$ and $m=\vm$);
\item pairwise edge-disjoint graphs $\vR_{j,c}$ ($0\le j<\vJ$, $c\in[4]$) and $\vM$ on $Z$, the \emph{own classes},
  each a spanning $(2^{-6},s')$-expander on $Z$ with $s'\ge2^{145}L^{41}$;
\item a partition $Z=\Rt\sqcup\Pl$ into \emph{rooted} and \emph{parentless} vertices;
\item an external phase map $\extph:Z\to[4]\cup\{*\}$.
\end{itemize}
For $w\in Z$ let $\vC_w:=\{wu:\ u\in\Aw(w),\ wu\in E(\vM),\ \extph(u)=*\}$, and assume
\begin{itemize}
\item[(E1$'$)] $|\vC_w|\ge L^5/8$ for every $w\in\Pl$.
\end{itemize}
The \emph{labels} of the run are the following independent random variables: for every $v\in\Pl$ a level
$\vlev(v)\in\{0,1,\dots,\vJ\}$ with $\Prob(\vlev(v)\ge j)=2^{-j}$ for $0\le j\le\vJ$ (that is,
$\Prob(\vlev(v)=j)=2^{-(j+1)}$ for $0\le j<\vJ$ and $\Prob(\vlev(v)=\vJ)=2^{-\vJ}$); for every $v\in Z$ a label
$\vkap(v)\in\{0,1,2,3,4\}$ with $\Prob(\vkap(v)=0)=1/2$ and $\Prob(\vkap(v)=c)=1/8$ for $c\in[4]$. Put
\begin{gather*}
\vU_j:=\Rt\cup\{v\in\Pl:\vlev(v)\ge j\},\qquad \vW_j:=\{v\in\Pl:\vlev(v)=j\},\\
\vZ_j:=\{u\in\vU_{j+1}:\vkap(u)=0\},\qquad \vV_{j,c}:=\{u\in\vU_{j+1}:\vkap(u)=c\}.
\end{gather*}
Let $\mathcal G_{\mathrm{PV}}$ be the event that the following three conditions hold:
\begin{itemize}
\item[(G2)] for all $0\le j<\vJ$ and $c\in[4]$, $\vR_{j,c}$ is $(2^{12}L^4,\vt)$-path connected through $\vV_{j,c}$;
\item[(G4)] for all $0\le j<\vJ$ and $w\in\Pl$, the set $\vC^j_w:=\{wu\in\vC_w:\ u\in\vZ_j\}$ has at least $8$
  elements;
\item[(G5)] $\sum_{j<\vJ}|\Pl\cap\vU_j|\le8|\Pl|$ and $|\Pl\cap\vU_{\vJ}|\le64|\Pl|/L$.
\end{itemize}
The event $\mathcal G_{\mathrm{PV}}$ is determined by the own classes, the sets $\Aw(w)$, the partition, $\extph$ and
the labels, and
\begin{equation}\label{s4:eqPVprob}
\Prob(\mathcal G_{\mathrm{PV}})\ \ge\ \frac12-\eta_{\mathrm{PV}}(N),\qquad
\eta_{\mathrm{PV}}(N):=2^{95}L^{31}N^{-3}+NL\,e^{-L^4/16}\ \le\ \frac1{100}.
\end{equation}
On $\mathcal G_{\mathrm{PV}}$, for \emph{every} edge set $H_0$ all of whose edges have both ends in $Z$ and which
contains $E(\vM)$ and every $E(\vR_{j,c})$, there is a partition $H_0=H^{\mathrm{obj}}\sqcup H^{\mathrm{arc}}$ and a
family $\mathfrak A$ of paths (\emph{arcs}) such that:
\begin{itemize}
\item[(a)] $H^{\mathrm{obj}}$ decomposes into at most $369|\Pl|$ objects;
\item[(b)] $H^{\mathrm{arc}}$ is the disjoint union of the edge sets of the arcs; the arcs are pairwise edge-disjoint
  paths of length at least $1$, each has two distinct ends, both in $\Rt$, and each carries a phase $c\in[4]$ such
  that every vertex $x$ of a phase-$c$ arc, ends included, satisfies $\extph(x)\ne c$;
\item[(c)] $|\mathfrak A|\le4(\vJ+1)N$, and every vertex $x$ is an end of at most $\deg_{H_0}(x)\le|Z|-1$ arcs;
\item[(d)] if $\Rt=\emptyset$ then $\mathfrak A=\emptyset$; if $\Pl=\emptyset$ then $H^{\mathrm{obj}}=\emptyset$.
\end{itemize}
The partition, the decomposition and the arcs are obtained from $H_0$ and the labels by a deterministic procedure.
\deps{\ref{s3:thmT16s}, \ref{s3:lemHB}, \ref{s3:lemMonotone}, \ref{s1:citCor22}, \ref{s1:factEG0},
\ref{s1:citChernoff}, \ref{s1:citMarkov}, \ref{s1:condG3}, \ref{s4:convVortex}, \ref{s4:lemTPV}, \ref{s1:convGraphs}, \ref{s1:defConstants}.}
\end{lemma}

\begin{proof}
The proof follows that of Lemma~\ref{s4:lemTPV}, with three changes: rooted vertices never retire; the classing uses
four phases and also avoids the external phases of both ends; and paths with two rooted ends that do not come from
cherries are returned as arcs instead of being closed.

\emph{Size conditions.} We use: (i) $L\ge2^{10}$; (ii) $N\ge L^3$; (iii) $4\cdot64\le L$ (a consequence of (i)); (iv)
$\eta_{\mathrm{PV}}(N)\le1/100$. They hold for $N\ge\Nzero$ by the choice of $\Nzero$.

\emph{Parameters.} By the definition of $\vJ$, $8\cdot2^{\vJ}\le L<16\cdot2^{\vJ}$, so
\begin{equation}\label{s4:eqPVJ}
2^{-\vJ}\ge8/L,\qquad 2^{-\vJ}<16/L,\qquad 1\le\vJ\le\log_2L,\qquad 4\vJ\le L
\end{equation}
(the last two by (i)). Lemma~\ref{s3:lemHB} with $\eps'=2^{-6}$ and $m=\vm\le s_O/2$ gives sets as in the statement
with multiplicity at most $\lceil2\vm L^2/2^{-6}\rceil=\vb$. Since $\vb\le2^7L^2(L^6+1)+1$,
\begin{equation}\label{s4:eqPVt}
2+\vb\ \le\ 2^7L^8+2^7L^2+3\ \le\ 2^9L^8=\vt .
\end{equation}
As in Lemma~\ref{s4:lemTPV}: $\vU_0=Z$, $\vU_{j+1}\subseteq\vU_j$, $\vU_j=\vU_{j+1}\sqcup\vW_j$,
$\vW_j\subseteq\Pl$, $\vW_j\cap\vU_{j+1}=\emptyset$, $\Rt\subseteq\vU_j$ for every $j$, and
$\vZ_j,\vV_{j,1},\dots,\vV_{j,4}$ partition $\vU_{j+1}$.

\emph{Probability of $\mathcal G_{\mathrm{PV}}$.} If $\Pl=\emptyset$, (G4) and (G5) hold surely, and the bound for (G2)
below covers this case as well; in the items (G4) and (G5) below we assume $\Pl\ne\emptyset$.
\begin{itemize}
\item[(G2)] Fix $(j,c)$, and let $\mathcal K_{j,c}$ be the family of sets $T\subseteq Z$ such that $\vR_{j,c}$ is
  $(2^{12}L^4,\vt)$-path connected through $T$. It is fixed, since the own classes are fixed data. It is closed
  upwards by Lemma~\ref{s3:lemMonotone}(ii), applied with $G=G':=\vR_{j,c}$, $\ell=\ell':=2^{12}L^4$, $t=t':=\vt$ and
  with $T\subseteq T'\subseteq Z=V(\vR_{j,c})$ as the two sets through which path connectivity is taken (the own
  classes are spanning). For $v\in Z$ the event $v\in\vV_{j,c}$ depends only on $\vkap(v)$ and, if $v\in\Pl$, on
  $\vlev(v)$, so these events are mutually independent over $v$; and $\Prob(v\in\vV_{j,c})=1/8\ge1/L$ if $v\in\Rt$
  and $\Prob(v\in\vV_{j,c})=2^{-(j+1)}/8\ge2^{-\vJ}/8\ge1/L$ if $v\in\Pl$, by \eqref{s4:eqPVJ} and size condition
  (i). Let $V'$ be a $(1/L)$-random subset of $Z$.
  Theorem~\ref{s3:thmT16s} applies to $X:=\vR_{j,c}$ (a $(2^{-6},s')$-expander), $V:=V'$, $\rho:=1/L$, $t:=\vt$,
  because $\rho N\ge L^2$ by size condition (ii) and $s'\ge2^{145}L^{41}\ge2^{135}\cdot2^9L^8\cdot L^{28}\cdot
  L^5=2^{135}\vt L^{28}\rho^{-5}$; so
  $\Prob(V'\notin\mathcal K_{j,c})\le2^{86}\cdot2^9L^8\cdot L^{19}\cdot L^3\cdot N^{-3}=2^{95}L^{30}N^{-3}$. By
  observation \textup{(MC)} (with $S:=Z$, $\mathcal K:=\mathcal K_{j,c}$, the random set $\vV_{j,c}$ in place of $V$,
  and $\rho:=1/L$), also $\Prob(\vV_{j,c}\notin\mathcal K_{j,c})\le2^{95}L^{30}N^{-3}$: $\vR_{j,c}$ fails to be
  $(2^{12}L^4,\vt)$-path connected through $\vV_{j,c}$ with probability at most $2^{95}L^{30}N^{-3}$. There are
  $4\vJ\le L$ pairs $(j,c)$. By Lemma~\ref{s3:lemMonotone}(iii) the pair multisets may be chosen after the labels
  are revealed.
\item[(G4)] Fix $j<\vJ$ and $w\in\Pl$. The set $\vC_w$ is fixed and has at least $L^5/8$ elements by (E1$'$). For
  $wu\in\vC_w$ the events $u\in\vZ_j$ depend on the labels of distinct vertices $u\ne w$, so they are independent,
  and $\Prob(u\in\vZ_j)=\frac12\Prob(u\in\vU_{j+1})\ge\frac12\min(1,2^{-\vJ})\ge4/L$ by \eqref{s4:eqPVJ}. As
  $(L^5/8)(4/L)=L^4/2\ge16$, observation \textup{(B)} gives $\Prob(|\vC^j_w|<8)\le e^{-L^4/16}$. There are at most
  $N\vJ\le NL$ pairs $(w,j)$.
\item[(G5)] $\Ex\sum_{j<\vJ}|\Pl\cap\vU_j|<2|\Pl|$ and $\Ex|\Pl\cap\vU_{\vJ}|=2^{-\vJ}|\Pl|<16|\Pl|/L$ by
  \eqref{s4:eqPVJ}; by Markov's inequality (Cited result~\ref{s1:citMarkov}) each inequality of (G5) fails with
  probability less than $1/4$.
\end{itemize}
Summing gives \eqref{s4:eqPVprob}. The three events involve only the own classes, $\vC_w$ (hence $\Aw(w)$, $\vM$
and $\extph$), the partition and the labels; none involves $H_0$. Fix labels in $\mathcal G_{\mathrm{PV}}$ and an edge
set $H_0$ as in the statement.

\emph{The process.} All choices are made by fixed rules. An edge is \emph{used} once it has been placed into an
output object or into an arc. Put $\vH_{j+1}:=\vH_j\setminus\{\text{edges used at step }j\}$, where $\vH_0:=H_0$.
Before step $j$ ($0\le j\le\vJ$) we maintain:
\begin{itemize}
\item[(Inv1)] every edge of $\vH_j$ has both ends in $\vU_j$;
\item[(Inv2)] every edge of $\vR_{j',c}$ with $j'\ge j$, $c\in[4]$ and both ends in $\vU_{j'+1}$ lies in $\vH_j$;
\item[(Inv3)] every $\vM$-edge $wu$ with $w\in\vW_{j'}$ for some $j'\ge j$ and $u\in\vU_{j'+1}$ lies in $\vH_j$.
\end{itemize}
They hold for $j=0$ because $\vU_0=Z$ and $H_0$ contains all own classes.

\emph{Step $j$} ($0\le j<\vJ$). Let $\vF_j$ be the set of edges of $\vH_j$ meeting $\vW_j$; by (Inv1) all their ends lie
in $\vU_j$.
\begin{itemize}
\item[(i)] \emph{Reserve.} For each $w\in\vW_j$ choose eight distinct edges $e^{c,i}_w\in\vC^j_w$ ($c\in[4]$,
  $i\in[2]$), possible by (G4). Each is an $\vM$-edge $wu$ with $w\in\vW_j$, $u\in\vZ_j\subseteq\vU_{j+1}$, so it lies
  in $\vH_j$ by (Inv3), hence in $\vF_j$; $w$ is its only end in $\vW_j$, so reserved edges of distinct $w$ are
  distinct. Its far end $u$ satisfies $\vkap(u)=0$ and $\extph(u)=*$.
\item[(ii)] \emph{Phases.} For every non-reserved edge of $\vF_j$ choose an end $w\in\vW_j$, call the other end $x$,
  and give the edge a phase $\gamma\in[4]\setminus\{\vkap(x),\extph(x),\extph(w)\}$; at most three of the four values
  are excluded, so a phase exists. Let $\vF_{j,c}$ be the set of non-reserved edges of phase $c$. Then, for every
  $c\in[4]$,
  \begin{equation}\label{s4:eqPVclass}
  V(\vF_{j,c})\cap\vV_{j,c}=\emptyset\qquad\text{and}\qquad \extph(y)\ne c\ \text{ for every } y\in V(\vF_{j,c}),
  \end{equation}
  the first as in \eqref{s4:eqTPVclass}, the second because both ends of a phase-$c$ edge were constrained.
\item[(iii)] \emph{Paths.} Let $\vPaths_{j,c}$ be a Corollary-22 decomposition of $\vF_{j,c}$; by \textup{(E)},
  $|\vPaths_{j,c}|\le|\vU_j|\le N$.
\item[(iv)] \emph{Ends in $\vW_j$.} Exactly as in step (iv) of the proof of Lemma~\ref{s4:lemTPV}, for every
  $w\in\vW_j$ and $c\in[4]$: if $w$ is an end of two paths of $\vPaths_{j,c}$, append $e^{c,1}_w$ to one and
  $e^{c,2}_w$ to the other at $w$; if of one, append $e^{c,1}_w$ and output $e^{c,2}_w$ as a single edge; if of none,
  form the cherry $u_1wu_2$ from $e^{c,1}_w=wu_1$ and $e^{c,2}_w=wu_2$. The resulting trails have pairwise distinct
  inner vertices and pairwise distinct edges.
\item[(v)] \emph{Splitting.} By \textup{(S)} each trail splits into at most two cycles (output) and at most one path
  with the same ends as the trail; let $\vPathsQ_{j,c}$ be the set of these paths. Exactly as in the proof of
  Lemma~\ref{s4:lemTPV} (with \eqref{s4:eqPVclass} and \eqref{s4:eqPVt} in place of \eqref{s4:eqTPVclass} and
  \eqref{s4:eqTPVt}): (P1) no vertex of a path in $\vPathsQ_{j,c}$ lies in $\vV_{j,c}$; (P2) its two ends are
  distinct and lie in $\vU_{j+1}$ (an end is an end of a path of $\vPaths_{j,c}$ not in $\vW_j$, hence in
  $\vU_j\setminus\vW_j=\vU_{j+1}$, or a far end); (P3) it has length at least $2$; (P4) every vertex is an end of at
  most $2+\vb\le\vt$ paths of $\vPathsQ_{j,c}$. Moreover:
  \begin{itemize}
  \item[(P5)] if $Q\in\vPathsQ_{j,c}$ does not come from a cherry, then $\extph(y)\ne c$ for every vertex $y$ of $Q$:
    the vertices of its trail are vertices of $\vF_{j,c}$ (use \eqref{s4:eqPVclass}) and far ends (where
    $\extph=*$).
  \end{itemize}
\item[(vi)] \emph{Arcs and closing.} A path $Q\in\vPathsQ_{j,c}$ that does not come from a cherry and has both ends in
  $\Rt$ is declared an \emph{arc of phase $c$}; its edges are put into $H^{\mathrm{arc}}$. Every other path of
  $\vPathsQ_{j,c}$ (one with an end in $\Pl$, or one coming from a cherry) is closed: the multiset of their end pairs
  has every vertex in at most $\vt$ pairs by (P4), so by (G2) there are pairwise edge-disjoint connectors
  $Q'\subseteq\vR_{j,c}$ joining the ends of these paths, of length at most $2^{12}L^4$, with inner vertices in
  $\vV_{j,c}$. As in step (vi) of the proof of Lemma~\ref{s4:lemTPV}, every connector edge is an
  $\vR_{j,c}$-edge with both ends in $\vU_{j+1}$, so it lies in $\vH_j$ by (Inv2), is not in $\vF_j$ and is not
  reserved; by (P1), (P3) and \textup{(C)}, $Q\cup Q'$ is a cycle, which is output.
\end{itemize}
Closing cherries inside the run (rather than declaring a cherry with two rooted ends an arc) guarantees (b) for the
centre $w$ of the cherry, for which $\extph(w)=c$ is not excluded.

The edges used at step $j$ are exactly the edges of $\vF_j$ and the connector edges, all of them edges of $\vH_j$;
as in the proof of Lemma~\ref{s4:lemTPV}, the objects and arcs of the step are pairwise edge-disjoint. The invariants
persist by the argument given there: an edge of $\vH_{j+1}$ lies in $\vH_j\setminus\vF_j$, so has both ends in
$\vU_j\setminus\vW_j=\vU_{j+1}$; step $j$ used only edges meeting $\vW_j$ and $\vR_{j,\cdot}$-edges, while edges
protected by (Inv2), (Inv3) for indices $j'\ge j+1$ have both ends in $\vU_{j+1}$ and are not $\vR_{j,\cdot}$-edges.

\emph{Cost of step $j$.} There are at most $4|\vW_j|$ single edges. For a phase $c$, a trail produces objects only if
(1) it is a cherry or has an appended edge, or (2) it is a path of $\vPaths_{j,c}$ without appended edges having an
end in $\Pl$. A trail of type (1) contains at least one of the at most $2|\vW_j|$ reserved edges of phase $c$ used in
(iv), and each of these edges lies in exactly one trail (a trail may contain two of them: a cherry, or a path of
$\vPaths_{j,c}$ with both ends in $\vW_j$); so there are at most $2|\vW_j|$ such trails. An end in $\Pl$ of a trail of type (2) is not in $\vW_j$ (ends in $\vW_j$ receive appended edges), so it
lies in $\Pl\cap\vU_{j+1}$, and each vertex is an end of at most two paths of $\vPaths_{j,c}$; so there are at most
$2|\Pl\cap\vU_{j+1}|$ such trails. Every remaining trail is a path of $\vPaths_{j,c}$ with both ends in $\Rt$ and
without appended edges; it is not split ($\mathrm{rep}=0$), and it becomes an arc. Each trail of type (1) or (2) gives
at most three objects (at most two split cycles and at most one closed cycle). Hence, using
$\vW_j\sqcup(\Pl\cap\vU_{j+1})=\Pl\cap\vU_j$, step $j$ outputs at most
\begin{equation}\label{s4:eqPVstep}
4|\vW_j|+4\cdot3\bigl(2|\Pl\cap\vU_{j+1}|+2|\vW_j|\bigr)\ \le\ 28\,|\Pl\cap\vU_j|
\end{equation}
objects. (No separate allowance is needed for split cycles of trails with two rooted ends: by the case analysis
just given such trails are either of type (1), and already counted, or not split at all. The bound
$28|\Pl\cap\vU_j|$ is the one used in (a) below.)

\emph{Finish.} By (Inv1) for $j=\vJ$, every edge of $\vH_{\vJ}$ has both ends in $\vU_{\vJ}=\Rt\sqcup \Pl_{\vJ}$, where
$\Pl_{\vJ}:=\Pl\cap\vU_{\vJ}$. Split $\vH_{\vJ}=E_1\sqcup E_2$, where $E_1$ is the set of edges with both ends in
$\Pl_{\vJ}$; every edge of $E_2$ has an end in $\Rt$.
\begin{itemize}
\item By (G5), $|\Pl_{\vJ}|\le64|\Pl|/L$. Apply Fact~\ref{s1:factEG0}(b) with $F:=E_1$, $W:=\Pl_{\vJ}$, $\beta:=64$ and
  $\alpha:=|\Pl|\le N$; its hypothesis $N>1$ holds as $L\ge2^{10}$ by size condition (i), and its hypothesis $4\beta\le L$
  is size condition (iii). So $E_1$ decomposes into at most $64|\Pl|$
  objects, which are output.
\item Give every edge $xy\in E_2$ a phase $c\in[4]\setminus\{\extph(x),\extph(y)\}$ (at most two values are excluded)
  and let $E_{2,c}$ be the set of edges of phase $c$. Take a Corollary-22 decomposition of each $E_{2,c}$. Let $T$ be
  one of its paths. If an end $p$ of $T$ lies in $\Pl_{\vJ}$, the edge of $T$ at $p$ lies in $E_2$, so its other end
  lies in $\Rt$; output this edge as a single edge (\emph{stripping}). After stripping at the ends of $T$ lying in
  $\Pl_{\vJ}$, what remains of $T$ has no edges or is a path of length at least $1$ with two distinct ends in $\Rt$ (every end
  of $T$ lies in $\Rt\sqcup \Pl_{\vJ}$, and a stripped end is replaced by a vertex of $\Rt$); in the latter case it is
  declared an arc of phase $c$. Every vertex $y$ of such an arc is an end of an edge of phase $c$, so
  $\extph(y)\ne c$. Since every vertex of $\Pl_{\vJ}$ is an end of at most two paths of each decomposition, at most
  $2|\Pl_{\vJ}|$ edges are stripped per phase, i.e.\ at most $8|\Pl_{\vJ}|\le16|\Pl|$ single
  edges in total (also $8|\Pl_{\vJ}|\le8|\Pl|$; the weaker factor $16$ suffices, and nothing
  depends on it).
\end{itemize}

\emph{Conclusion.} Every edge of $H_0$ is used exactly once: edges leave $\vH_j$ only by being used, $\vH_{\vJ}$ is
used entirely in the finish, and every used edge lies in $H_0$ (reserved edges lie in $\vF_j\subseteq\vH_j$ and
connector edges in $\vH_j$). Let $H^{\mathrm{arc}}$ be the union of the edge sets of the arcs and
$H^{\mathrm{obj}}:=H_0\setminus H^{\mathrm{arc}}$, the union of the output objects.

(a) By \eqref{s4:eqPVstep} and (G5), the steps $0\le j<\vJ$ together output at most
$28\sum_{j<\vJ}|\Pl\cap\vU_j|\le28\cdot8|\Pl|=224|\Pl|$ objects, and the finish outputs at most $64|\Pl|+16|\Pl|$
objects. So the number of objects is at most $224|\Pl|+64|\Pl|+16|\Pl|=304|\Pl|\le369|\Pl|$.

(b) Arcs are pairwise edge-disjoint (each edge is used once), they are paths of length at least $1$ with two
distinct ends in $\Rt$ (by (P2) and the rule in (vi), and by the finish), and every vertex of a phase-$c$ arc has
$\extph\ne c$ (by (P5), since arcs never come from cherries, and by the finish).

(c) At step $j$, each arc of phase $c$ comes from a distinct path of $\vPaths_{j,c}$, so there are at most $N$ of
them; in the finish there are at most $|\vU_{\vJ}|\le N$ arcs per phase by \textup{(E)}. Hence
$|\mathfrak A|\le4\vJ N+4N=4(\vJ+1)N$. For the per-vertex bound, let $x\in Z$. By (b) the arcs are pairwise
edge-disjoint paths with two distinct ends, and their edges lie in $H^{\mathrm{arc}}\subseteq H_0$. An arc with end
$x$ contains exactly one edge at $x$ (a path meets each of its two ends in exactly one edge), and distinct arcs are
edge-disjoint; so $x$ is an end of at most $\deg_{H_0}(x)$ arcs. Every edge of $H_0$ at $x$ joins $x$ to a vertex of
$Z\setminus\{x\}$, and distinct edges at $x$ have distinct other ends (edge sets are sets of edges of a simple graph,
Convention~\ref{s1:convGraphs}(a)); hence $\deg_{H_0}(x)\le|Z|-1$. (Counting only ends of Corollary-22 paths and
far ends of reserved edges, which would bound the arc ends at $x$ by $8\vJ+8+\vb$, does not suffice. That count
misses the arcs left by stripping in the finish: such an arc can end at a vertex $x\in\Rt$ that is an interior vertex of its
Corollary-22 path, and Corollary~22 does not limit the number of such paths at $x$. The bound through degrees covers
all arcs.)

(d) Arcs have their ends in $\Rt$, so there are none if $\Rt=\emptyset$. If $\Pl=\emptyset$ then every $\vW_j$ is
empty, the steps use nothing, $\Pl_{\vJ}=\emptyset$, and $E_1=\emptyset$; no edge is stripped, and all of
$H_0=\vH_{\vJ}=E_2$ is partitioned into arcs, so $H^{\mathrm{obj}}=\emptyset$.
\end{proof}

\subsection{Theorem \texorpdfstring{VX$^+$}{VX+}: vortex absorption with arbitrary extra edges}\label{s4:ssecVX}

\begin{theorem}[Theorem VX$^+$ (sharpened vortex absorption with arbitrary extra edges)]\label{s4:thmVXp}
Let $Z$ be a set of $N\ge\Nzero$ vertices, $L:=\log_2N$, and let $O$ be a spanning $(\eps_O,s)$-expander on $Z$, where
either
\[
\eps_O=2^{-5}\ \text{ and }\ s\ge2^{146}L^{38}\log_2L,\qquad\text{or}\qquad \eps_O=2^{-6}\ \text{ and }\
s\ge2^{151}L^{38}\log_2L .
\]
Put $\vJ:=\lfloor\log_2(L/6)\rfloor$, $k:=3\vJ+1$, $\vm:=\lceil L^3\rceil$, $\vb:=\lceil2\vm L^2/\eps_O\rceil$ (that
is, $\lceil64L^2\vm\rceil$, resp.\ $\lceil128L^2\vm\rceil$) and $\vt:=2^8L^5$, resp.\ $\vt:=2^9L^5$. The labels of the
run are independent: a uniform colouring of $E(O)$ with the $k$ colours $\vR_{j,c}$ ($0\le j<\vJ$, $c\in[3]$) and
$\vM$; for every $v\in Z$ a level $\vlev(v)\in\{0,1,\dots,\vJ\}$ with $\Prob(\vlev(v)\ge j)=2^{-j}$ for $0\le j\le\vJ$ (that is,
$\Prob(\vlev(v)=j)=2^{-(j+1)}$ for $0\le j<\vJ$ and $\Prob(\vlev(v)=\vJ)=2^{-\vJ}$); for every $v\in Z$ a
label $\vkap(v)\in\{0,1,2,3\}$ with $\Prob(\vkap(v)=0)=1/2$ and $\Prob(\vkap(v)=c)=1/6$ ($c\in[3]$). There is an
event $\mathcal G_{\mathrm{VX}}$, determined by $O$ and the labels only, with
\begin{equation}\label{s4:eqVXprob}
\begin{gathered}
\Prob(\mathcal G_{\mathrm{VX}})\ \ge\ 1-\eta_{\mathrm{VX}}(N)\ \ge\ \frac12,\qquad\text{where}\\
\eta_{\mathrm{VX}}(N):=2LN^{-5}+2^{95}L^{28}N^{-3}+NL\,e^{-3L^2/(32\log_2L)}+L\,e^{-N/(5000L)},
\end{gathered}
\end{equation}
such that on $\mathcal G_{\mathrm{VX}}$ the following holds simultaneously for \emph{every} graph $G$ with $V(G)=Z$ and
$E(O)\subseteq E(G)$ (the edges of $E(G)\setminus E(O)$ are arbitrary): $E(G)$ decomposes into at most
\[
38.4\,N+13\,N\ \le\ 52N\ \le\ 80N
\]
objects, all of which are cycles except at most $N+13N/L=N+o(N)$ single edges. In particular
$\f(G)\le80N$.
\deps{\ref{s3:lemL15p}, \ref{s3:thmT16s}, \ref{s3:lemHB}, \ref{s3:lemMonotone}, \ref{s1:citCor22},
\ref{s1:factEG0}, \ref{s1:citChernoff}, \ref{s1:condG3}, \ref{s4:convVortex}, \ref{s1:convGraphs}, \ref{s1:defConstants}.}
\end{theorem}

\begin{proof}
\emph{Size conditions.} We use: (i) $L\ge2^{10}$; (ii) $N\ge L^3$; (iii) $4\cdot13\le L$ (a consequence of (i)); (iv)
$\eta_{\mathrm{VX}}(N)\le1/100$. They hold for $N\ge\Nzero$ by the choice of $\Nzero$.

\emph{Parameters.} By the definition of $\vJ$, $6\cdot2^{\vJ}\le L<12\cdot2^{\vJ}$, so $2^{-\vJ}\ge6/L$,
$2^{-\vJ}<12/L$, $1\le\vJ$ and $\vJ+1\le L$; moreover $k=3\vJ+1\le3\log_2L+1\le4\log_2L$, so $2k\le8\log_2L$ and
$k\le L$. Since $s\ge2\lceil L^3\rceil=2\vm$, Lemma~\ref{s3:lemHB} (with $\eps'=\eps_O\ge2^{-7}$, $m=\vm$) gives sets
$\Aw(w)\subseteq N_O(w)$, $|\Aw(w)|=\vm$, with every vertex in at most $\vb$ of them; we fix the lexicographically
first such family. As $\vm\le L^3+1$, we have $\vb\le2^6L^5+2^6L^2+1$ (resp.\ $\vb\le2^7L^5+2^7L^2+1$), hence
\begin{equation}\label{s4:eqVXt}
2+\vb\ \le\ 2L^2\vm/\eps_O+3\ \le\ \vt .
\end{equation}
Put $\vU_j:=\{v:\vlev(v)\ge j\}$, $\vW_j:=\vU_j\setminus\vU_{j+1}$, $\vZ_j:=\{u\in\vU_{j+1}:\vkap(u)=0\}$ and
$\vV_{j,c}:=\{u\in\vU_{j+1}:\vkap(u)=c\}$; thus $\vU_0=Z$, $\vW_j\cap\vU_{j+1}=\emptyset$ and
$\vZ_j,\vV_{j,1},\vV_{j,2},\vV_{j,3}$ partition $\vU_{j+1}$.

\emph{Good events.} Let $\mathcal G_{\mathrm{VX}}$ be the intersection of:
\begin{itemize}
\item[(G1)] every colour class is a spanning $(\eps_O,s/(2k))$-expander on $Z$. As $s\ge40kL$,
  Lemma~\ref{s3:lemL15p} bounds the failure probability by $2kN^{-5}\le2LN^{-5}$. Note
  $s/(2k)\ge s/(8\log_2L)\ge2^{143}L^{38}$ (resp.\ $\ge2^{148}L^{38}$).
\item[(G2)] every $\vR_{j,c}$ is $(2^{12}L^4,\vt)$-path connected through $\vV_{j,c}$. Here the event
  $v\in\vV_{j,c}$ depends only on $(\vlev(v),\vkap(v))$, so these events are mutually independent over $v$ and
  independent of the colouring, and $\Prob(v\in\vV_{j,c})=2^{-(j+1)}/6\ge2^{-\vJ}/6\ge1/L$ for every $v$.
  Fix a colouring $\chi$ in (G1) and a pair $(j,c)$, and let $\mathcal K_{j,c}$ be the family of sets
  $T\subseteq Z$ such that $\vR_{j,c}$ is $(2^{12}L^4,\vt)$-path connected through $T$. It is determined by $\chi$.
  It is closed upwards by Lemma~\ref{s3:lemMonotone}(ii), applied with $G=G':=\vR_{j,c}$, $\ell=\ell':=2^{12}L^4$,
  $t=t':=\vt$ and with $T\subseteq T'\subseteq Z=V(\vR_{j,c})$ as the two sets through which path connectivity is
  taken (the colour classes are spanning). Let $V'$ be a $(1/L)$-random subset of $Z$ and apply Theorem~\ref{s3:thmT16s}
  with $X:=\vR_{j,c}$, $V:=V'$, $\rho:=1/L$, $t:=\vt$: its hypotheses are $N/L\ge L^2$ (by size condition (ii)) and
  $s/(2k)\ge2^{135}\vt L^{28}L^5$, i.e.\ $2^{143}L^{38}$ for $\vt=2^8L^5$ and $2^{144}L^{38}$ for $\vt=2^9L^5$, both
  covered by (G1). So $\Prob(V'\notin\mathcal K_{j,c})\le2^{86}\vt L^{19}L^3N^{-3}\le2^{95}L^{27}N^{-3}$. By
  observation \textup{(MC)} (with $S:=Z$, $\mathcal K:=\mathcal K_{j,c}$, the random set $\vV_{j,c}$ in place of $V$,
  and $\rho:=1/L$), also $\Prob(\vV_{j,c}\notin\mathcal K_{j,c})\le2^{95}L^{27}N^{-3}$, the probability being over the
  levels and the labels $\vkap$. As these are independent of the colouring, a union bound over the $3\vJ\le L$ pairs
  $(j,c)$ shows that the probability that (G1) holds and (G2) fails is at most
  $\sum_{\chi}\Prob(\text{the colouring is }\chi)\cdot3\vJ\cdot2^{95}L^{27}N^{-3}\le2^{95}L^{28}N^{-3}$, the sum
  running over the colourings $\chi$ in (G1). By
  Lemma~\ref{s3:lemMonotone}(iii), pair multisets may be chosen after the labels are revealed.
\item[(G3)] $|\vU_j|\le1.01\cdot2^{-j}N$ for $0\le j\le\vJ$. $|\vU_j|$ is binomial with mean
  $2^{-j}N\ge2^{-\vJ}N\ge6N/L$; by Cited result~\ref{s1:citChernoff} with $\delta=1/100$ it exceeds $1.01$ times its
  mean with probability at most $e^{-10^{-4}\cdot6N/(3L)}\le e^{-N/(5000L)}$; there are $\vJ+1\le L$ values of $j$.
\item[(G4)] for every $0\le j<\vJ$ and every $w\in Z$ the set
  $\vC^j_w:=\{wu:\ u\in\Aw(w),\ wu\in\vM,\ u\in\vZ_j\}$ has at least $9$ elements. For $u\in\Aw(w)$ the events
  $\{wu\in\vM,\ u\in\vZ_j\}$ are independent (distinct edges $wu$ and distinct vertices $u\ne w$), each of probability
  $\frac1k\cdot2^{-(j+1)}\cdot\frac12\ge\frac1{4\log_2L}\cdot\frac3L$. As
  $\vm\cdot\frac{3}{4L\log_2L}\ge\frac{3L^2}{4\log_2L}\ge18$, observation \textup{(B)} bounds the failure
  probability for one pair $(w,j)$ by $e^{-3L^2/(32\log_2L)}$; there are at most $NL$ pairs.
\end{itemize}
These events are determined by $O$ (through the colouring and $\Aw(\cdot)$) and the labels; summing the failure
probabilities gives \eqref{s4:eqVXprob}, and $\eta_{\mathrm{VX}}(N)\le1/100\le1/2$ by (iv). Fix labels in
$\mathcal G_{\mathrm{VX}}$ and a graph $G$ with $V(G)=Z$ and $E(O)\subseteq E(G)$.

\emph{The process.} All choices are made by fixed rules. Put $\vH_0:=E(G)$ and let $\vH_{j+1}$ be $\vH_j$ minus the
edges used at step $j$. Before step $j$ ($0\le j\le\vJ$) we maintain: (Inv1) every edge of $\vH_j$ has both ends in
$\vU_j$; (Inv2) every $\vR_{j',c}$-edge with $j'\ge j$ and both ends in $\vU_{j'+1}$ lies in $\vH_j$; (Inv3) every
$\vM$-edge $wu$ with $w\in\vW_{j'}$ for some $j'\ge j$ and $u\in\vU_{j'+1}$ lies in $\vH_j$. They hold for $j=0$
since $\vH_0\supseteq E(O)$ and $\vU_0=Z$.

\emph{Step $j$} ($0\le j<\vJ$). Let $\vF_j$ be the set of edges of $\vH_j$ meeting $\vW_j$. By (Inv3),
$\vC^j_w\subseteq\vF_j$ for every $w\in\vW_j$; every edge of $\vC^j_w$ has exactly one end in $\vW_j$ (its other end
lies in $\vZ_j\subseteq\vU_{j+1}$), so the sets $\vC^j_w$ ($w\in\vW_j$) are pairwise disjoint.
\begin{itemize}
\item[(a)] \emph{Parity.} For every $w\in\vW_j$ with $\deg_{\vF_j}(w)$ odd, output one edge of $\vC^j_w$ as a single
  edge (the \emph{parity edge} of $w$). Let $\vF'_j$ be $\vF_j$ minus the parity edges. As each parity edge has
  exactly one end in $\vW_j$, $\deg_{\vF'_j}(w)$ is even for every $w\in\vW_j$.
\item[(b)] \emph{Reserved and flexible edges.} For every $w\in\vW_j$ choose eight further distinct edges of
  $\vC^j_w$ (possible by (G4), since $|\vC^j_w|\ge9$): \emph{reserved} edges $e^{c,1}_w,e^{c,2}_w$ ($c\in[3]$) and
  \emph{flexible} edges $f_w,f'_w$. These \emph{chosen} edges of distinct $w$ are distinct; the end of a chosen edge
  in $\vZ_j$ is its \emph{far end}, and $\vkap=0$ there.
\item[(c)] \emph{Classes.} Every edge of $\vF'_j$ that is not chosen receives, once and for all, a class: choose an
  end $w\in\vW_j$, call the other end $x$, and pick $\gamma\in[3]\setminus\{\vkap(x)\}$. For $w\in\vW_j$ let
  $d_c(w)$ be the number of such classed edges of class $c$ at $w$ (an edge with both ends in $\vW_j$ counts at both
  ends). Then $d_1(w)+d_2(w)+d_3(w)=\deg_{\vF'_j}(w)-8$ is even, so the number of odd $d_c(w)$ is $0$ or $2$. If
  $d_c(w)$ and $d_{c'}(w)$ are odd ($c\ne c'$), give $f_w$ class $c$ and $f'_w$ class $c'$; otherwise give both
  class $1$. A flexible edge has only one end in $\vW_j$, so these choices at different $w$ do not interact. Let
  $\vF_{j,c}$ be the set of class-$c$ edges (classed and flexible; reserved edges excluded). Then for all
  $w\in\vW_j$ and $c\in[3]$, $\deg_{\vF_{j,c}}(w)$ is even, and $V(\vF_{j,c})\cap\vV_{j,c}=\emptyset$ (ends in
  $\vW_j$ are outside $\vU_{j+1}$, other ends $x$ of classed edges have $\vkap(x)\ne c$, far ends have $\vkap=0$).
\item[(d)] \emph{Paths.} Let $\vPaths_{j,c}$ be a Corollary-22 decomposition of $\vF_{j,c}$. By \textup{(E)} and
  (Inv1), $|\vPaths_{j,c}|\le|\vU_j|$, and every $w\in\vW_j$ is an end of an even number, hence of $0$ or $2$, paths
  of $\vPaths_{j,c}$.
\item[(e)] \emph{Moving ends out of $\vW_j$.} For $w\in\vW_j$ and $c\in[3]$: if $w$ is an end of two paths of
  $\vPaths_{j,c}$, append $e^{c,1}_w$ at $w$ to one and $e^{c,2}_w$ at $w$ to the other; if of none, form the cherry
  $u_1wu_2$ from $e^{c,1}_w=wu_1$, $e^{c,2}_w=wu_2$ ($u_1\ne u_2$). Every resulting trail $x_0\cdots x_q$ has pairwise
  distinct edges and pairwise distinct inner vertices (those of one path of $\vPaths_{j,c}$, or the centre of a
  cherry). By \textup{(S)} it splits into at most two cycles, which are output, and at most one path with the ends
  $x_0\ne x_q$; let $\vPathsQ_{j,c}$ be the set of these paths. Every $Q\in\vPathsQ_{j,c}$ satisfies:
  (P1) $V(Q)\cap\vV_{j,c}=\emptyset$ (its vertices are vertices of $\vF_{j,c}$, cherry centres in $\vW_j$, or far
  ends); (P2) its ends are distinct and lie in $\vU_{j+1}$ (ends of paths of $\vPaths_{j,c}$ outside $\vW_j$, which
  lie in $\vU_j\setminus\vW_j=\vU_{j+1}$, or far ends in $\vZ_j$); (P3) its length is at least $2$ (all its edges meet
  $\vW_j$, its ends do not); (P4) every vertex $x$ is an end of at most $2+\vb$ paths of $\vPathsQ_{j,c}$ (at most
  two ends of paths of $\vPaths_{j,c}$, plus at most one reserved edge $e^{c,i}_w$ with far end $x$ for each of the at
  most $\vb$ vertices $w$ with $x\in\Aw(w)$; each trail yields at most one path, with the trail's ends).
\item[(f)] \emph{Closing.} For each $c$, the end pairs of $\vPathsQ_{j,c}$ form a multiset of pairs of distinct
  vertices with every vertex in at most $2+\vb\le\vt$ pairs, by (P2), (P4) and \eqref{s4:eqVXt}. By (G2) there are
  pairwise edge-disjoint connectors $Q'\subseteq\vR_{j,c}$ joining them, of length at most $2^{12}L^4$, with inner
  vertices in $\vV_{j,c}\subseteq\vU_{j+1}$. All vertices of a connector lie in $\vU_{j+1}$, so its edges are
  $\vR_{j,c}$-edges with both ends in $\vU_{j+1}$: they lie in $\vH_j$ by (Inv2), are not in $\vF_j$ and are not
  chosen edges (which are $\vM$-edges). By (P1), (P3) and \textup{(C)}, $Q\cup Q'$ is a cycle, which is output.
\end{itemize}
Every edge of $\vF_j$ is used exactly once at step $j$ (parity edges as single edges; reserved edges appended or in
cherries; classed and flexible edges on paths of $\vPaths_{j,c}$; every trail edge ends in a split cycle or a closed
cycle), connectors of one class are edge-disjoint and connectors of distinct classes lie in distinct colour classes,
and connectors avoid $\vF_j$. The invariants persist: an edge of $\vH_{j+1}$ lies in
$\vH_j\setminus\vF_j$, so both its ends lie in $\vU_{j+1}$; step $j$ used only edges meeting $\vW_j$ and
$\vR_{j,\cdot}$-edges, whereas the edges protected by (Inv2), (Inv3) for indices $j'\ge j+1$ have both ends in
$\vU_{j+1}$ and are $\vR_{j',\cdot}$- or $\vM$-edges.

\emph{Cost.} Step $j$ outputs at most $|\vW_j|$ parity edges and, for each $c$, at most
$|\vPaths_{j,c}|+|\vW_j|\le|\vU_j|+|\vW_j|$ trails, each giving at most three objects. As $\vW_j\subseteq\vU_j$, this
is at most $|\vW_j|+9(|\vU_j|+|\vW_j|)\le19|\vU_j|$ objects. By (Inv1), after the last step every edge of
$\vH_{\vJ}$ has both ends in $\vU_{\vJ}$, and $|\vU_{\vJ}|\le1.01\cdot2^{-\vJ}N<1.01\cdot12N/L\le13N/L$ by (G3) and
$2^{-\vJ}<12/L$. By Fact~\ref{s1:factEG0}(b) with $F:=\vH_{\vJ}$, $W:=\vU_{\vJ}$, $\beta:=13$ and $\alpha:=N$ (its hypothesis $N>1$ holds as $L\ge2^{10}$ by
size condition (i), and its hypothesis $4\beta\le L$ is size condition (iii)), $\vH_{\vJ}$ decomposes into at most $13N$ objects, at most $|\vU_{\vJ}|\le13N/L$ of which are
single edges. Every edge of $E(G)=\vH_0$ is used exactly once. By (G3) the total number of objects is at most
\[
\sum_{j<\vJ}19\cdot1.01\cdot2^{-j}N+13N\ \le\ 38.4\,N+13N\ \le\ 52N .
\]
The only single edges are the parity edges, at most
$\sum_j|\vW_j|\le N$ in total since the sets $\vW_j$ are disjoint, and the at most $13N/L$ single edges of the
finish; all other objects are cycles. So there are $N+o(N)$ single edges; the finish may add single edges
beyond the at most $N$ parity edges.
\end{proof}

\begin{remark}[Lemma TPV as a vortex-absorption statement; not used by the main argument]\label{s4:remTPVVX}
Let $O$ be a spanning $(\eps_O,s)$-expander on a set $Z$ of $N\ge\Nzero$ vertices with $\eps_O\in[2^{-7},2^{-5}]$ and
$s\ge2^{150}L^{42}$. Apply Lemma~\ref{s4:lemTPV} with $\vP:=Z$ and $\vQ:=\emptyset$. On
$\mathcal G_{\mathrm{TPV}}$, for every edge set $E\supseteq E(O)$ inside $Z$, property (T1) gives $\EV=E$, and (T3)
gives a decomposition of $E$ into at most $169N$ objects. So Lemma~\ref{s4:lemTPV} contains an absorption statement of
the same kind as Theorem~\ref{s4:thmVXp}, with $169$ in place of $80$ and the threshold $2^{150}L^{42}$ in place of
$2^{146}L^{38}\log_2L$ (resp.\ $2^{151}L^{38}\log_2L$); for $\eps_O=2^{-6}$ both thresholds apply to the same
expanders once $s\ge2^{150}L^{42}$. Consequently, wherever Theorem~\ref{s4:thmVXp} is applied in this paper to an
expander that also meets the threshold of Lemma~\ref{s4:lemTPV}, either result may be used; the choice changes only
the coefficient ($80$ or $169$) of the vertex mass to which the result is applied. In this paper
Theorem~\ref{s4:thmVXp} is applied only to parts whose total mass enters the final bound through a term $\varepsilon n$ with
$\varepsilon$ depending only on $\Dstar$ and tending to $0$ as $\Dstar\to\infty$, so the choice does not affect any constant of the cost line. This remark is not used in the proofs; the
paper applies Theorem~\ref{s4:thmVXp}.
\deps{\ref{s4:lemTPV}, \ref{s4:thmVXp}.}
\end{remark}

%% ======================================================================
%% Section s5: Light parts: Theorem U machinery and K-RED
%% (part of the proof of the main theorem)
%% ======================================================================
\section{Light parts: Theorem U machinery and K-RED}\label{s5:sec}

This section treats the light parts of the hierarchy $\HBtp$. Each light part runs the four-phase P-vortex of Lemma~\ref{s4:lemPV} on its own edges. The paths that the vortex cannot close inside the part (the \emph{arcs}, whose two ends are rooted in light parts two or more rounds back) are chained over one quotient multigraph whose nodes are all admissible parents of all earlier rounds. The connectors run through lent colour classes of the parents, inside random vertex sets (\emph{zones}) that are pairwise disjoint over all parents of all rounds. Lemma K-RED, the last lemma of this section, collects the result: all edges outside the standalone pre-parts and outside the externally lent edges decompose into $(\Dstar/2+\cKRED+\epsK(\Dstar))n+\cVX\,\mathrm{dem}+\cPV\,\mathrm{lp}$ objects, where $\mathrm{dem}$ and $\mathrm{lp}$ are stage-1 quantities whose weighted sum has expectation at most $\epsU(\Dstar)n$.

\paragraph{Setting of this section.} Throughout, $G$ is a graph on $n\ge\Nzero$ vertices with $d_1\ge\Dstar$; $\Dstar$ satisfies \Gc{1}--\Gc{4} (Condition~\ref{s1:condGamma}); a valid $\HBtp$ run on $G$ is fixed (Definition~\ref{s2:defHBtp}); and the stage-1 lending data of Definition~\ref{s3:defCOL} are drawn. Every lemma below is stated for every outcome of the stage-1 data that it does not name; probabilities are taken only over the random labels named in each statement. For a light part $Y$ we write $r=r(Y)$ for its round, $L_Y=\log_2|Y|$, $H_Y=X_Y$, and we use $J_Y=\lfloor\log_2(L_Y/8)\rfloor$ and $\kown=4J_Y+1\le L_Y$ from Definition~\ref{s3:defCOL}(iii). We shall use repeatedly the following consequences of Section~\ref{s2:sec}: for a light part $Y$ of round $r$,
\begin{equation}\label{s5:eqLY}
|Y|\ge P_r/2\ge\lambda_r^{103}/2,\qquad L_Y\ge 103\log_2\lambda_r-1\ge 102\log_2\lambda_r,\qquad R-r\le 2\logs d_r+2\le\log_2\lambda_r ,
\end{equation}
by Proposition~\ref{s2:propStructure}(iv), $P_r=\lceil\lambda_r^{\Cp}\rceil$ with $\Cp=103$, and Lemma~\ref{s2:lemTower}(a),(d). In particular $R-r\le L_Y/102$. Moreover, for every vertex $v$,
\begin{equation}\label{s5:eqZp}
\sum_{Y\ni v}L_Y^{-2}\ <\ \tfrac12 ,
\end{equation}
the sum over all light parts $Y$ containing $v$. Indeed, by Proposition~\ref{s2:propStructure}(iv), $v$ lies in at most one light part per round; let $Y_r$ denote the light part of round $r$ containing $v$, if there is one. By \eqref{s5:eqLY}, $L_{Y_r}^{-2}\le(102\log_2\lambda_r)^{-2}$. For $r<R$, Lemma~\ref{s2:lemTower}(a) gives $\lambda_r\ge d_{r+1}^{1/\Aexp}=2^{\lambda_{r+1}/\Aexp}$, hence
\[
\log_2\lambda_r\ \ge\ \lambda_{r+1}/\Aexp\ =\ 2^{\log_2\lambda_{r+1}}/\Aexp\ \ge\ 2\log_2\lambda_{r+1},
\]
because $x\defeq\log_2\lambda_{r+1}\ge\log_2\log_2\Dstar$, so that \Gc{1}(a),(b) (Condition~\ref{s1:condG1}) at $x$ give $x\ge2^8$ and $2^{x}\ge2^{14}\Aexp x^3\ge 2\Aexp x$. Hence $\log_2\lambda_r\ge 2^{R-r}\log_2\lambda_R\ge 2^{R-r}\cdot 2^8$, by \Gc{1}(a) at $\log_2\lambda_R$, and
\[
\sum_{Y\ni v}L_Y^{-2}\ \le\ \sum_{r\le R}(102\log_2\lambda_r)^{-2}\ \le\ (102\cdot 2^8)^{-2}\sum_{k\ge 0}4^{-k}\ \le\ \tfrac43\,(102\cdot 2^8)^{-2}\ <\ \tfrac12 .
\]

\subsection{Zones}

\begin{definition}[vertex choice, sublabels and zones; stage 1b]\label{s5:defZones}
For a light part $Y$ of round $r$ (Definition~\ref{s2:defAncestors}) put $\zp_Y\defeq L_Y^{-2}$ and $\Tslot_Y\defeq\lceil L_Y^2\rceil$ (as in Definition~\ref{s3:defCOL}). Independently for every vertex $v$ of $G$, and independently of all other stage-1 data (the colourings and the JS labels of Definition~\ref{s3:defCOL}, and any further stage-1 labels), draw
\[
\mathrm{choice}(v)\in\{\mathrm{none}\}\cup\{Y:\ Y\text{ a light part with }v\in Y\text{ and }r(Y)\le R-2\},
\]
with $\Prob(\mathrm{choice}(v)=Y)=\zp_Y$ for each such $Y$ and $\mathrm{choice}(v)=\mathrm{none}$ with the remaining probability (which is at least $1/2$, since $\sum_Y\zp_Y\le\sum_{Y\ni v}L_Y^{-2}<1/2$ by \eqref{s5:eqZp}, the first sum over the light parts $Y$ available to $\mathrm{choice}(v)$). If $\mathrm{choice}(v)=Y$, the vertex $v$ also draws, independently and uniformly, a \emph{sublabel}
\[
(l,c,\uslot)\qquad\text{with}\qquad r(Y)+2\le l\le R,\quad c\in[4],\quad \uslot\in[\Tslot_Y].
\]
The sublabels available to $Y$ are exactly the indices of the family $\IU$ of $Y$ in Definition~\ref{s3:defCOL}(ii). For such an index put
\[
\Zone_{Y,l,c,\uslot}\defeq\{v\in Y:\ \mathrm{choice}(v)=Y\text{ and the sublabel of $v$ is }(l,c,\uslot)\}.
\]
For a round $l$ and a vertex $v$ we say that $v$ \emph{carries a round-$l$ zone} if $\mathrm{choice}(v)\ne\mathrm{none}$ and the sublabel of $v$ has first coordinate $l$; the \emph{phase} of that zone is then the second coordinate $c$. Equivalently, $v$ carries a round-$l$ zone of phase $c$ iff $v\in\Zone_{Y,l,c,\uslot}$ for some light part $Y$ and some $\uslot$. A \emph{round-$l$ phase-$c$ zone} is any set $\Zone_{Y,l,c,\uslot}$.
\deps{\ref{s2:defAncestors}, \ref{s3:defCOL}, \ref{s2:propStructure}, \ref{s2:lemTower}, \ref{s1:condG1}}
\end{definition}

\begin{lemma}[zones]\label{s5:lemZones}
\begin{enumerate}
\item[(i)] For every vertex $v$, $\sum_{Y}\zp_Y\le 1/2$, the sum over all light parts $Y\ni v$.
\item[(ii)] For every light part $Y$ of round $r\le R-2$ and every sublabel $(l,c,\uslot)$ of $Y$, the set $\Zone_{Y,l,c,\uslot}$ is exactly a $\rho_Y$-random subset of $Y$ (each vertex of $Y$ independently, product measure), where
\[
\rho_Y\defeq\frac{\zp_Y}{(R-r-1)\cdot 4\cdot\Tslot_Y}\ \ge\ \frac{1}{12L_Y^5}.
\]
\item[(iii)] All zones $\Zone_{Y,l,c,\uslot}$, over all light parts $Y$ and all sublabels $(l,c,\uslot)$ of $Y$, are pairwise disjoint.
\item[(iv)] The family of all zones is independent of the colourings (i)--(iii) and of the JS labels (iv) of Definition~\ref{s3:defCOL}.
\end{enumerate}
\deps{\ref{s5:defZones}, \ref{s2:propStructure}, \ref{s2:lemTower}, \ref{s1:condG1}, \ref{s3:defCOL}}
\end{lemma}

\begin{proof}
(i) Since $\zp_Y=L_Y^{-2}$, this is \eqref{s5:eqZp}, which is proved in the setting paragraph of this section. In particular the distribution of $\mathrm{choice}(v)$ in Definition~\ref{s5:defZones} is well defined, and $\Prob(\mathrm{choice}(v)=\mathrm{none})\ge 1/2$.

(ii) Let $Y$ have round $r\le R-2$ and fix a sublabel $(l,c,\uslot)$ of $Y$. The number of sublabels of $Y$ is $(R-r-1)\cdot4\cdot\Tslot_Y$, since $l$ ranges over the $R-r-1\ge 1$ values $r+2,\dots,R$. For $v\in Y$,
\begin{align*}
\Prob\bigl(v\in\Zone_{Y,l,c,\uslot}\bigr)&=\Prob(\mathrm{choice}(v)=Y)\cdot\Prob\bigl(\text{sublabel }(l,c,\uslot)\,\big|\,\mathrm{choice}(v)=Y\bigr)\\
&=\frac{\zp_Y}{(R-r-1)\cdot4\cdot\Tslot_Y}=\rho_Y .
\end{align*}
The event $\{v\in\Zone_{Y,l,c,\uslot}\}$ is a function of the pair $(\mathrm{choice}(v),\text{sublabel of }v)$, and these pairs are independent over $v$. So $\Zone_{Y,l,c,\uslot}$ contains each vertex of $Y$ independently with probability exactly $\rho_Y$. For the lower bound, \eqref{s5:eqLY} gives $R-r-1\le L_Y/102$, and $\Tslot_Y=\lceil L_Y^2\rceil\le 2L_Y^2$. Hence
\[
\rho_Y\ \ge\ \frac{L_Y^{-2}}{(L_Y/102)\cdot4\cdot 2L_Y^2}\ =\ \frac{102}{8L_Y^5}\ \ge\ \frac{1}{12L_Y^5}.
\]

(iii) Each vertex $v$ has a single value $\mathrm{choice}(v)$ and, if it is not $\mathrm{none}$, a single sublabel. So $v$ lies in at most one set $\Zone_{Y,l,c,\uslot}$, namely the one with $Y=\mathrm{choice}(v)$ and $(l,c,\uslot)$ its sublabel. Distinct zones are therefore disjoint.

(iv) The zones are functions of the variables $\mathrm{choice}(v)$ and of the sublabels, which by Definition~\ref{s5:defZones} are drawn independently of the colourings and of the JS labels.
\end{proof}

\subsection{Stage-2 statuses of light parts and their probabilities}

\begin{definition}[stage-2 statuses of light parts]\label{s5:defStages}
Fix an outcome of the stage-1 data (the colourings and labels of Definition~\ref{s3:defCOL} and the choices and sublabels of Definition~\ref{s5:defZones}). Let $Y$ be a light part of round $r$. Then $H_Y=X_Y$ is a $(2^{-6},s_r/2)$-expander on $Y$ (Definition~\ref{s2:defAncestors} and Proposition~\ref{s2:propStructure}(i)). Put $\vm_Y\defeq\lceil L_Y^6\rceil$ and $\vb_Y\defeq\lceil 2^7L_Y^2\vm_Y\rceil$. Since $s_r\ge\lambda_r^{100}$ and $L_Y\le 2\lambda_r$ (Lemma~\ref{s2:lemTower}(b)), we have $2\vm_Y\le 2^8\lambda_r^6+2\le s_r/2$. So Lemma~\ref{s3:lemHB}, applied to $X_Y$ with $\eps'=2^{-6}$ and $m=\vm_Y$ (then $b=\lceil 2^7L_Y^2\vm_Y\rceil=\vb_Y$), provides sets $\Aw_Y(w)\subseteq N_{X_Y}(w)$, $w\in Y$, with $|\Aw_Y(w)|=\vm_Y$ such that every vertex lies in at most $\vb_Y$ of them. We fix one such family by a fixed rule; it depends only on the run. Let $\vM_Y$ denote the own colour class $\vM$ of $Y$ (Definition~\ref{s3:defCOL}(iii)).
\begin{itemize}
\item \textup{(E1)} holds for $Y$ if for every round $l$ with $r\le l\le R$ and every $w\in Y$:
\[
\#\{u\in\Aw_Y(w):\ wu\in\vM_Y\ \text{and $u$ carries no round-$l$ zone}\}\ \ge\ L_Y^5/8 .
\]
Only the instance $l=r$ is used later in this section; the instances $l>r$ (the child rounds of $Y$) are included so that one event serves every round, and they are paid for in the failure bound below.
\item $Y$ is \emph{demoted} if (E1) fails for $Y$ or a COL(a) event fails for $Y$, that is, if one of $\Own_Y$, $\Lend_Y$ is not a $(2^{-6},s_r/8)$-expander on $Y$, or some lent class of $Y$ is not a $(2^{-6},s_{\mathrm{lent}})$-expander on~$Y$ with $s_{\mathrm{lent}}\defeq s_r/(16\klend(Y))$, or some own class of $Y$ is not a $(2^{-6},s_{\mathrm{own}})$-expander on $Y$ with $s_{\mathrm{own}}\defeq s_r/(16\kown)$ (Lemma~\ref{s3:lemCOL}(a) with $\eps_Y=2^{-6}$, $s_Y=s_r/2$).
\item $Y$ is \emph{parent-bad} if for some $(l,c,\uslot)$ with $r+2\le l\le R$, $c\in[4]$, $\uslot\in[\Tslot_Y]$, the class $\LU_{Y,l,c,\uslot}$ is not $(2^{12}L_Y^4,\,t_Y)$-path connected through $\Zone_{Y,l,c,\uslot}$ (in the multiset sense of Cited result~\ref{s1:citDef7}), where $t_Y=\lceil\lambda_r^{1.6}\rceil$ is the U-multiplicity parameter of Definition~\ref{s3:defCOL}. If $r\ge R-1$, the family $\IU$ of $Y$ is empty (Definition~\ref{s3:defCOL}(ii)), so $Y$ has no sublabels and no zones and is never parent-bad; Lemma~\ref{s5:lemZones}(ii) is then vacuous, and Lemma~\ref{s5:lemE1}(c) for $Y$ reduces to (b).
\item A \emph{good parent} is a light part that is neither demoted nor parent-bad. Put
\[
\Bad\defeq\{\text{demoted light parts}\}\cup\{\text{parent-bad light parts}\}.
\]
\item For a non-demoted light part $Z$ of round $l$ put
\[
\Rt(Z)\defeq Z\cap\bigcup\{Y:\ Y\text{ a good parent of round }\le l-2\},\qquad \Pl(Z)\defeq Z\setminus\Rt(Z),
\]
and give every $v\in\Rt(Z)$ the fixed parent $\parU(v)$, the good parent of the smallest round $\le l-2$ that contains $v$. (It is unique: $v$ lies in at most one light part per round, Proposition~\ref{s2:propStructure}(iv).) Since $v$ lies in at most one light part of round $l$, $\parU(v)$ depends only on $v$ and on the round $l$ of the part in which it is used.
\item $\mathrm{dem}\defeq\sum_{Y\text{ demoted}}|Y|$ and $\mathrm{lp}\defeq\sum_{Y\in\Bad}|Y|\,(R-r(Y))$.
\end{itemize}
A pair $(v,Z)$ with $Z$ a light part of round $l$ and $v\in\Pl(Z)$ is exactly a pair that is parentless with respect to the set $\Bad$ in the sense of Proposition~\ref{s2:propParentless}(iii): no light part outside $\Bad$ of a round $\le l-2$ contains $v$. Hence, by Proposition~\ref{s2:propParentless}(iii),
\begin{equation}\label{s5:eqPl}
\sum_{Z}|\Pl(Z)|\ \le\ 2n+\mathrm{lp},
\end{equation}
the sum over all non-demoted light parts $Z$ of all rounds. The statuses (E1), demoted, parent-bad, good parent, the set $\Bad$, the sets $\Rt(Z)$, $\Pl(Z)$, the map $\parU$, and the numbers $\mathrm{dem}$, $\mathrm{lp}$ are deterministic functions of the stage-1 outcome: (E1) depends only on the colouring of $Y$ and on the choices and sublabels, COL(a) only on the colouring of $Y$, and parent-bad only on the colouring of $Y$ and the zones of $Y$. None of them depends on any edge set formed later (remainders, junk, lent edges used by consumers) or on any later random choice.
\deps{\ref{s5:defZones}, \ref{s3:defCOL}, \ref{s3:lemHB}, \ref{s2:propParentless}, \ref{s2:defAncestors}, \ref{s2:propStructure}, \ref{s2:lemTower}, \ref{s1:citDef7}}
\end{definition}

In \eqref{s5:eqPl}, note that a vertex $v$ can be parentless only in its light parts at rounds $j_1(v)$ and $j_1(v)+1$, and at those later rounds at which its light part of round $j_1(v)$ lies in $\Bad$, i.e.\ is demoted \emph{or} parent-bad. Both kinds of bad parts are counted in $\mathrm{lp}$.

\begin{lemma}[bad probabilities]\label{s5:lemE1}
Let $Y$ be a light part of round $r$. Over the stage-1 randomness:
\begin{enumerate}
\item[(a)] $\Prob(\text{(E1) fails for }Y)\le 3L_Y|Y|\exp(-L_Y^5/32)$;
\item[(b)] $\Prob(Y\text{ is demoted})\le|Y|^{-2}/2$;
\item[(c)] $\Prob(Y\in\Bad)\le|Y|^{-2}$.
\end{enumerate}
\deps{\ref{s3:lemCOL}, \ref{s3:lemHB}, \ref{s5:lemZones}, \ref{s5:defStages}, \ref{s1:citChernoffGen}, \ref{s3:defCOL}, \ref{s3:lemL15p}, \ref{s3:lemCOLJV}, \ref{s2:propStructure}, \ref{s2:lemTower}, \ref{s1:condG1}}
\end{lemma}

\begin{proof}
(a) Fix a round $l$ with $r\le l\le R$ and a vertex $w\in Y$. For $u\in\Aw_Y(w)$ let $\chi_u$ be the indicator of the event that $wu\in\vM_Y$ and $u$ carries no round-$l$ zone. The sets $\Aw_Y(w)$ are fixed (Definition~\ref{s5:defStages}), so the only randomness in $\chi_u$ is the colour of the edge $wu$ and the pair (choice, sublabel) of the vertex $u$. The edges $wu$, $u\in\Aw_Y(w)$, are distinct, colours of distinct edges are independent (Definition~\ref{s3:defCOL}), the pairs (choice, sublabel) of distinct vertices are independent, and the two families are independent of each other (Lemma~\ref{s5:lemZones}(iv)). Hence the $\chi_u$, $u\in\Aw_Y(w)$, are independent. By the uniform split and the uniform own colouring of Definition~\ref{s3:defCOL}(i),(iii), $\Prob(wu\in\vM_Y)=\frac12\cdot\frac1{\kown}\ge\frac{1}{2L_Y}$, because $\kown=4J_Y+1\le L_Y$. By Lemma~\ref{s5:lemZones}(i), $\Prob(u\text{ carries a round-$l$ zone})\le\Prob(\mathrm{choice}(u)\ne\mathrm{none})\le 1/2$. Hence $\Prob(\chi_u=1)\ge 1/(4L_Y)$, and the count $\mathrm{cnt}_{w,l}\defeq\sum_{u\in\Aw_Y(w)}\chi_u$ is a sum of $\vm_Y$ independent indicators with mean $\mu\ge\vm_Y/(4L_Y)\ge L_Y^5/4$. By the Chernoff lower tail for sums of independent indicators (Cited result~\ref{s1:citChernoffGen}) with $\delta=1/2$,
\[
\Prob\bigl(\mathrm{cnt}_{w,l}<L_Y^5/8\bigr)\ \le\ \Prob\bigl(\mathrm{cnt}_{w,l}<\mu/2\bigr)\ \le\ \exp(-\mu/8)\ \le\ \exp(-L_Y^5/32).
\]
The number of rounds $l$ with $r\le l\le R$ is $R-r+1\le 2\logs d_r+3\le L_Y/102+1\le 3L_Y$ by \eqref{s5:eqLY}. A union bound over these rounds and over the $|Y|$ vertices $w$ gives~(a).

(b) COL(a) is the conclusion of three applications of Lemma~15$^+$ (Lemma~\ref{s3:lemL15p}) in the proof of Lemma~\ref{s3:lemCOL}: with $k=2$ to $X_Y$ (giving $\Own_Y$, $\Lend_Y$); with $k=\klend(Y)$ to $\Lend_Y$ (the lent classes); with $k=\kown$ to $\Own_Y$ (the own classes). The hypotheses $s\ge 40kL$ of these applications are row~3 of Table~\ref{s3:tabCOLJV} (Lemma~\ref{s3:lemCOLJV}). If $r\ge R-1$, then $\klend(Y)=0$ and $\Lend_Y$ has no classes (Definition~\ref{s3:defCOL}(ii)), so only the applications with $k=2$ and $k=\kown$ occur. The second and third colourings are uniform given the split (Definition~\ref{s3:defCOL}(ii),(iii)), so Lemma~15$^+$ applies to them conditionally on the split, on the event that $\Lend_Y$ (respectively $\Own_Y$) is a $(2^{-6},s_r/8)$-expander; if that event fails, COL(a) has already failed in the first application. In each application each colour class fails with probability at most $2|Y|^{-5}$. Summing over the $2+\klend(Y)+\kown$ classes,
\[
\Prob(\text{COL(a) fails for }Y)\ \le\ 2\bigl(2+\klend(Y)+\kown\bigr)|Y|^{-5}\ \le\ 6|Y|^{-4},
\]
since $\klend(Y)\le\lambda_r^{3.3}\le|Y|$ if $r\le R-2$ (Lemma~\ref{s3:lemCOLJV}(i) and \eqref{s5:eqLY}), $\klend(Y)=0$ if $r\ge R-1$, and $\kown\le L_Y\le|Y|$. By (a) and $|Y|=2^{L_Y}$, $\Prob(\text{(E1) fails})\le 3L_Y2^{L_Y}e^{-L_Y^5/32}\le 2^{-4L_Y}=|Y|^{-4}$, because $\ln(3L_Y)+5L_Y\ln2\le L_Y^5/32$ for $L_Y\ge 10$ (and $L_Y\ge 102\cdot 2^8$ by \eqref{s5:eqLY} and \Gc{1}(a)). Hence $\Prob(Y\text{ demoted})\le 7|Y|^{-4}\le|Y|^{-2}/2$.

(c) By Lemma~\ref{s5:lemZones}(ii),(iv), for every sublabel $(l,c,\uslot)$ of $Y$ the set $V_{l,c,\uslot}\defeq\Zone_{Y,l,c,\uslot}\subseteq V(Y)$ contains each vertex of $Y$ independently with probability $\rho_Y\ge 1/(12L_Y^5)$, independently of the colouring of $Y$. (The zones of $Y$ are dependent across indices, since they are disjoint; Lemma~\ref{s3:lemCOL}(c) allows this.) So Lemma~\ref{s3:lemCOL}(c), applied with this family, gives $\Prob(Y\text{ parent-bad})\le|Y|^{-2}/2$. With (b), $\Prob(Y\in\Bad)\le\Prob(Y\text{ demoted})+\Prob(Y\text{ parent-bad})\le|Y|^{-2}$.
\end{proof}

\begin{lemma}[expected demotion and lost-parent mass]\label{s5:lemExpect}
Over the stage-1 randomness,
\[
\Ex\bigl[\cVX\,\mathrm{dem}+\cPV\,\mathrm{lp}\bigr]\ \le\ \epsU(\Dstar)\,n,\qquad\text{where}\qquad \epsU(\Dstar)\defeq 2462\,(\log_2\Dstar)^{-205}.
\]
Moreover $\epsU(\Dstar)\to0$ as $\Dstar\to\infty$.
\deps{\ref{s5:lemE1}, \ref{s5:defStages}, \ref{s2:propOV}, \ref{s2:lemTower}, \ref{s2:lemLacunary}, \ref{s1:condG1}, \ref{s2:propStructure}}
\end{lemma}

\begin{proof}
By linearity and Definition~\ref{s5:defStages},
\[
\Ex\bigl[\cVX\,\mathrm{dem}+\cPV\,\mathrm{lp}\bigr]=\sum_{Y}|Y|\Bigl(80\,\Prob(Y\text{ demoted})+369\,(R-r(Y))\,\Prob(Y\in\Bad)\Bigr),
\]
the sum over all light parts $Y$. By Lemma~\ref{s5:lemE1}(b),(c), both probabilities are at most $|Y|^{-2}$; by \eqref{s5:eqLY}, $R-r\le 2\logs d_r+2\le\log_2\lambda_r$. Hence the term of $Y$ is at most
\[
|Y|^{-1}\bigl(80+369\,(2\logs d_r+2)\bigr)\ \le\ |Y|^{-1}\,449\log_2\lambda_r\ \le\ 898\,\lambda_r^{-103}\log_2\lambda_r ,
\]
using $|Y|\ge\lambda_r^{103}/2$. Light parts of round $r$ are ancestors of round $r$, and there are at most $1.37n/P_r\le 1.37n\lambda_r^{-103}$ of them (Proposition~\ref{s2:propOV}, (K1), with the $\HBtp$ constant). Therefore round $r$ contributes at most $1.37\cdot898\,n\,\lambda_r^{-206}\log_2\lambda_r\le 1231\,n\,\lambda_r^{-205}$. Let $F(x)\defeq x^{-205}$ on $[\log_2\Dstar,\infty)$. It is non-increasing, and $F(2^{x/\Aexp})=2^{-205x/\Aexp}\le x^{-205}/2=F(x)/2$ for $x\ge\log_2\Dstar$. Indeed, put $\mu\defeq\log_2x\ge\log_2\log_2\Dstar$; then \Gc{1}(a),(b) (Condition~\ref{s1:condG1}) at $\mu$ give $\mu\ge2^8$ and $x=2^\mu\ge2^{14}\Aexp\mu^3$, so $x/\Aexp\ge2^{14}\mu^3\ge\mu+1=\log_2x+1$ and $2^{-205x/\Aexp}\le2^{-205}x^{-205}\le x^{-205}/2$. (This is also the case $a=205$ of Lemma~\ref{s2:lemLacunary}(iv).) By Lemma~\ref{s2:lemLacunary}(ii), $\sum_{r\le R}F(\lambda_r)\le 2F(\log_2\Dstar)$. So the total is at most $2\cdot1231\,n\,(\log_2\Dstar)^{-205}=\epsU(\Dstar)n$. Finally, $\epsU(\Dstar)=2462\,(\log_2\Dstar)^{-205}\to0$ as $\Dstar\to\infty$.
\end{proof}

\paragraph{How Lemma~\ref{s5:lemExpect} is used.} It is an expectation bound over stage~1 only. It is used through a single Markov bound, with multiplier $3$, on one combined weighted stage-1 functional that contains $\cVX\,\mathrm{dem}+\cPV\,\mathrm{lp}$ as a summand (not through separate Markov bounds for $\mathrm{dem}$ and $\mathrm{lp}$: with multiplier $2$ each, the two failure probabilities could add up to $1$, and no outcome would be guaranteed). All statements below hold for every stage-1 outcome; the bounds on $\mathrm{dem}$ and $\mathrm{lp}$ enter only through that functional.

\subsection{The child side}

\begin{lemma}[child side]\label{s5:lemChild}
Fix a stage-1 outcome and a non-demoted light part $Z$ of round $l$. Apply Lemma~\ref{s4:lemPV} to the vertex set $Z$ with the following data:
\begin{itemize}
\item $O\defeq H_Z=X_Z$, and the sets $\Aw(w)\defeq\Aw_Z(w)$ of Definition~\ref{s5:defStages} (so $\vm=\lceil L_Z^6\rceil$, $\vb=\lceil2^7L_Z^2\vm\rceil$);
\item as own classes, the classes $\vR_{j,c}$ ($0\le j<J_Z$, $c\in[4]$) and $\vM=\vM_Z$ of the colouring of $\Own_Z$ in Definition~\ref{s3:defCOL}(iii);
\item $\Rt\defeq\Rt(Z)$ and $\Pl\defeq\Pl(Z)$ (Definition~\ref{s5:defStages});
\item $\extph(v)\defeq c$ if $v$ carries a round-$l$ zone and $c$ is its phase, and $\extph(v)\defeq *$ if $v$ carries no round-$l$ zone.
\end{itemize}
Then all hypotheses of Lemma~\ref{s4:lemPV} hold, including \textup{(E1$'$)}. Let $\mathcal{G}_Z$ be its good event $G_{\mathrm{PV}}$ for these data. It is an event over the stage-3 labels of $Z$ (the levels $\vlev$ on $\Pl(Z)$ and the labels $\vkap$ on $Z$); it is determined by the stage-1 outcome and these labels; and its conditional probability given the stage-1 outcome is at least $1/2-o(1)\ge1/3$. On $\mathcal{G}_Z$ the following holds for \emph{every} edge set $H_0(Z)$ with $\Own_Z\subseteq H_0(Z)\subseteq E_l(Z)$: there is a partition $H_0(Z)=H^{\mathrm{obj}}(Z)\sqcup H^{\mathrm{arc}}(Z)$ such that
\begin{enumerate}
\item[(a)] $H^{\mathrm{obj}}(Z)$ decomposes into at most $\cPV|\Pl(Z)|$ objects;
\item[(b)] $H^{\mathrm{arc}}(Z)$ is the edge set of a family of pairwise edge-disjoint paths (the \emph{arcs} of $Z$), each with at least one edge and both ends in $\Rt(Z)$, and each carrying a phase $c\in[4]$, such that every vertex of a phase-$c$ arc, ends included, lies in no round-$l$ phase-$c$ zone;
\item[(c)] $Z$ has at most $14(J_Z+1)|Z|$ arcs, and every vertex $x$ is an end of at most $\deg_{H_0(Z)}(x)\le|Z|-1$ arcs of $Z$;
\item[(d)] if $\Rt(Z)=\emptyset$, in particular if $l\le2$, then $Z$ has no arcs.
\end{enumerate}
\deps{\ref{s4:lemPV}, \ref{s5:defStages}, \ref{s5:lemZones}, \ref{s3:lemCOL}, \ref{s2:propStructure}, \ref{s3:defCOL}, \ref{s1:condG3}}
\end{lemma}

\begin{proof}
\emph{Hypotheses.} Put $N\defeq|Z|$ and $L\defeq L_Z=\log_2N$.
\begin{itemize}
\item $N\ge P_l/2\ge\Nzero$, by Proposition~\ref{s2:propStructure}(iv) and \Gc{3} (Condition~\ref{s1:condG3}).
\item $O=X_Z$ is a spanning $(2^{-6},s_l/2)$-expander on $Z$ (Proposition~\ref{s2:propStructure}(i)), and $2\vm\le s_l/2$ (Definition~\ref{s5:defStages}). The sets $\Aw_Z(w)$ are Lemma-HB sets of $O$ with $\vm=\lceil L^6\rceil$ and $\vb=\lceil 2^7L^2\vm\rceil$, as Lemma~\ref{s4:lemPV} requires.
\item $Z$ is not demoted, so COL(a) holds for $Z$: every own class is a spanning $(2^{-6},s')$-expander on $Z$ with $s'\defeq s_l/(16\kown)$ (Lemma~\ref{s3:lemCOL}(a)). By Lemma~\ref{s3:lemCOL}(e), $s'\ge 2^{145}L^{41}$ and $2\lceil L^6\rceil\le s'$. The index range of the own classes is $0\le j<J_Z=\lfloor\log_2(L/8)\rfloor$, $c\in[4]$, which is the range $0\le j<\vJ$ of Lemma~\ref{s4:lemPV}. The own classes are fixed in stage~1, before the stage-3 labels of $Z$ are drawn.
\item $\Rt(Z)\sqcup\Pl(Z)=Z$ by definition.
\item $\extph$ is a well-defined map $Z\to[4]\cup\{*\}$, because every vertex carries at most one zone (Lemma~\ref{s5:lemZones}(iii)).
\item \textup{(E1$'$)}. For $u\in Z$ we have $\extph(u)=*$ iff $u$ carries no round-$l$ zone. Since $Z$ is not demoted, (E1) holds for $Z$; its instance at the round $l=r(Z)$ of $Z$ states that for every $w\in Z$, in particular for every $w\in\Pl(Z)$,
\[
\#\{u\in\Aw_Z(w):\ wu\in\vM_Z,\ \extph(u)=*\}\ \ge\ L^5/8 ,
\]
which is \textup{(E1$'$)}.
\end{itemize}

\emph{The good event.} By Lemma~\ref{s4:lemPV}, $G_{\mathrm{PV}}$ depends only on the own classes, on $\Rt$, on $\extph$ and on the stage-3 labels of $Z$. The own classes are stage-1 data; $\Rt(Z)$ is a deterministic function of stage~1 (Definition~\ref{s5:defStages}); $\extph$ is a function of the choices and sublabels (stage~1). So $\mathcal G_Z$ is determined by the stage-1 outcome and the stage-3 labels of $Z$. Its probability is at least $1/2-o(1)$ by Lemma~\ref{s4:lemPV}, and the $o(1)$ term is at most $1/100$ by \Gc{3}, so the probability is at least $1/3$. In particular $\mathcal G_Z$ does not depend on $H_0(Z)$.

\emph{Conclusion.} Let $H_0(Z)$ satisfy $\Own_Z\subseteq H_0(Z)\subseteq E_l(Z)$. Every edge of $E_l(Z)$ has both ends in $Z$ (Proposition~\ref{s2:propStructure}(iii)), and $\Own_Z$ is the union of the own classes. So $H_0\defeq H_0(Z)$ is an edge set inside $Z$ containing all own classes, and on $G_{\mathrm{PV}}$ Lemma~\ref{s4:lemPV} gives $H_0(Z)=H^{\mathrm{obj}}\sqcup H^{\mathrm{arc}}$ with its properties (a)--(d). By Lemma~\ref{s4:lemPV}(b), every arc is a path with at least one edge and has two distinct ends. Now:
\begin{itemize}
\item (a) is Lemma~\ref{s4:lemPV}(a), since $369=\cPV$.
\item (b): by Lemma~\ref{s4:lemPV}(b) the arcs are pairwise edge-disjoint paths with both ends in $\Rt$, each with a phase $c\in[4]$, and every vertex $x$ of a phase-$c$ arc has $\extph(x)\ne c$. If $x$ lay in a round-$l$ phase-$c$ zone $\Zone_{Y,l,c,\uslot}$, then $x$ would carry a round-$l$ zone of phase $c$, i.e.\ $\extph(x)=c$. So $x$ lies in no round-$l$ phase-$c$ zone.
\item (c): Lemma~\ref{s4:lemPV}(c) bounds the number of arcs by $4(\vJ+1)N\le14(\vJ+1)N=14(J_Z+1)|Z|$ (the weaker constant $14$ is kept so that the constants of this section are unchanged) and the number of arcs ending at a vertex $x$ by $\deg_{H_0}(x)=\deg_{H_0(Z)}(x)\le N-1=|Z|-1$.
\item (d): Lemma~\ref{s4:lemPV}(d) gives no arcs when $\Rt=\emptyset$. If $l\le2$ there is no round $\le l-2$, so $\Rt(Z)=\emptyset$.
\end{itemize}
Since $\mathcal G_Z$ does not depend on $H_0(Z)$, on $\mathcal G_Z$ the conclusion holds simultaneously for all admissible $H_0(Z)$; in particular it holds for whatever set $H_0(Z)$ the later processing order of this section produces. The stage-3 labels of distinct parts are independent, so outcomes in the events $\mathcal G_Z$ can be fixed part by part; no union bound over parts is used.
\end{proof}

\subsection{The parent side: chaining over U-bundles}

For a round $l\ge3$ put $\bar J_l\defeq\max\{J_Z:\ Z\text{ a light part of round }l\}$ (and $\bar J_l\defeq0$ if there is none), and let $\nu_l$ be the number of ancestors of rounds $\le l-2$, as in Proposition~\ref{s2:propOV}.

\begin{lemma}[parent side: multi-parent chaining over U-bundles]\label{s5:lemParent}
Fix a stage-1 outcome and, for every non-demoted light part $Z$, a stage-3 outcome in the event $\mathcal G_Z$ of Lemma~\ref{s5:lemChild}. Fix a round $l$ with $3\le l\le R$ and, for every non-demoted light part $Z$ of round $l$, an edge set $H_0(Z)$ with $\Own_Z\subseteq H_0(Z)\subseteq E_l(Z)$; the arcs of $Z$ are those given by Lemma~\ref{s5:lemChild} for this $H_0(Z)$. For $c\in[4]$ the \emph{U-bundle} $\Bdl_{l,c}$ is the set of all phase-$c$ arcs of all non-demoted light parts of round $l$; $E(\Bdl_{l,c})$ denotes the union of their edge sets. Then for every $c\in[4]$ there is a set $\LentU_{l,c}$ of edges, disjoint from $E(\Bdl_{l,c})$, such that $E(\Bdl_{l,c})\cup\LentU_{l,c}$ decomposes into cycles and single edges, where:
\begin{enumerate}
\item[(i)] $\LentU_{l,c}\subseteq\bigcup_Y\bigcup_{\uslot\in[\Tslot_Y]}\LU_{Y,l,c,\uslot}$, the union over the good parents $Y$ of rounds $\le l-2$; $\LentU_{l,c}$ is the union of the edge sets of pairwise edge-disjoint \emph{connectors}, one per transition: a connector for a transition at the good parent $Y$ with slot $\uslot$ is a path in $\LU_{Y,l,c,\uslot}$ of length at most $2^{12}L_Y^4$ that joins the two vertices of the transition and has all its interior vertices in $\Zone_{Y,l,c,\uslot}$;
\item[(ii)] the number of objects, summed over $c\in[4]$, is at most
\[
4\cdot14\,(\bar J_l+1)\,M_l(M_l+1)\,\nu_l\;+\;\mathrm{cap}_l,\qquad \mathrm{cap}_l\le\frac{14\,(\bar J_l+1)\,n}{(102\log_2\lambda_{l-2})^2},
\]
where $\mathrm{cap}_l$ is the number of \emph{visit-capping pieces}; and
\[
\sum_{l=3}^{R}\Bigl(4\cdot14\,(\bar J_l+1)M_l(M_l+1)\nu_l+\mathrm{cap}_l\Bigr)\ \le\ \epsChain(\Dstar)\,n,
\]
where
\[
\epsChain(\Dstar)\defeq\frac{614}{\Dstar}+\frac{\log_2\bigl(2\Aexp\log_2(\Aexp\log_2\log_2\Dstar)\bigr)}{371\,(\log_2\log_2\Dstar)^2};
\]
\item[(iii)] no other edge is used: every object consists of edges of $E(\Bdl_{l,c})\cup\LentU_{l,c}$, and every such edge lies in exactly one object.
\end{enumerate}
The sets $\LentU_{l,c}$ for distinct pairs $(l,c)$ are pairwise disjoint, and the objects and $\LentU_{l,c}$ depend only on the stage-1 outcome and on the bundle $\Bdl_{l,c}$. Moreover $\epsChain(\Dstar)\to0$ as $\Dstar\to\infty$.
\deps{\ref{s5:lemChild}, \ref{s5:defStages}, \ref{s5:lemZones}, \ref{s3:defCOL}, \ref{s3:lemCOLJV}, \ref{s3:lemMonotone}, \ref{s1:citEuler}, \ref{s2:propOV}, \ref{s2:lemTower}, \ref{s2:propStructure}, \ref{s1:condG1}, \ref{s1:citDef7}}
\end{lemma}

\begin{proof}
Fix $l$ and $c$; the letter $Z$ ranges over the non-demoted light parts of round $l$ and $Y$ over the good parents of rounds $\le l-2$.

\emph{Step 0: facts used.}
\begin{itemize}
\item[(F-a)] $|Z|\le M_l$, since $Z\subseteq Z^0$ and $|Z^0|\le M_l$ (Proposition~\ref{s2:propStructure}(ii)). Hence $L_Z\le\log_2M_l$, and $J_Z+1\le\log_2L_Z-2\le\log_2\log_2M_l$.
\item[(F-b)] \emph{(Not used.)} If $Y$ has round $r\le l-2$ then $L_Y\ge L_Z$. Indeed $|Y|\ge P_r/2$ (Proposition~\ref{s2:propStructure}(iv)), $P_r\ge P_{l-2}$ (as $P_{r'}\ge2P_{r'+1}$) and $P_{l-2}/2\ge M_l\log^4M_l$ (both by Lemma~\ref{s2:lemTower}(b)), so $|Y|\ge M_l\ge|Z|$. ((F-b) is not used in this proof: the multiplicity bound of Step~5 is $M_l-1$, which does not involve $L_Z$. It is recorded for completeness.)
\item[(F-c)] Light parts of round $l$ are pairwise vertex-disjoint (Proposition~\ref{s2:propStructure}(iv)), and the edge sets $E_l(Z)$, $E_r(Y)$ of distinct light parts are pairwise disjoint (Proposition~\ref{s2:propStructure}(iii)).
\item[(F-d)] By Lemma~\ref{s5:lemChild}(b),(c), every arc of $Z$ is a path with at least one edge, hence with two distinct ends; both ends lie in $\Rt(Z)$; the arcs of $Z$ are pairwise edge-disjoint; $Z$ has at most $14(J_Z+1)|Z|\le \mathrm{ncl}_l\defeq14(\bar J_l+1)M_l$ arcs; every vertex $v$ is an end of at most $\deg_{H_0(Z)}(v)\le|Z|-1\le M_l-1$ arcs of $Z$ (the last step by (F-a)); every vertex of an arc of $Z$ lies in $Z$ (its edges lie in $H_0(Z)\subseteq E_l(Z)$, Proposition~\ref{s2:propStructure}(iii)); and every vertex of a phase-$c$ arc lies in no round-$l$ phase-$c$ zone.
\end{itemize}
For an end $v$ of an arc of $Z$ we have $v\in\Rt(Z)$, so $\parU(v)$ is defined; it is a good parent of round $\le l-2$ containing $v$ (Definition~\ref{s5:defStages}).

\emph{Step 1: classes.} For each $Z$ list its phase-$c$ arcs in a fixed order as $a_{Z,1},\dots,a_{Z,\mathrm{na}(Z)}$, with $\mathrm{na}(Z)\le \mathrm{ncl}_l$. For $1\le i\le \mathrm{ncl}_l$ the \emph{class} $i$ is $\mathcal A_i\defeq\{a_{Z,i}:\ \mathrm{na}(Z)\ge i\}$. Each class contains at most one arc of each $Z$, so by (F-c) the arcs of one class are pairwise vertex-disjoint. Every arc of $\Bdl_{l,c}$ lies in exactly one class.

\emph{Step 2: the quotient multigraphs.} For each class $i$ let $\UQ_{l,c,i}$ be the multigraph (loops allowed) whose vertices are the good parents of rounds $\le l-2$, with one edge $e_a$ joining $\parU(v)$ and $\parU(v')$ for each arc $a\in\mathcal A_i$ with ends $v,v'$ (a loop if $\parU(v)=\parU(v')$). Good parents are ancestors of rounds $\le l-2$, so $\UQ_{l,c,i}$ has at most $\nu_l$ vertices, and $\nu_l\le2.74n/P_{l-2}$ (Lemma~\ref{s2:lemTower}(b), which completes (K1) of Proposition~\ref{s2:propOV}).

\emph{Step 3: $T$-join and Euler trails.} Let $\mathrm{Odd}_i$ be the set of vertices of odd degree in $\UQ_{l,c,i}$ (a loop adds $2$ to the degree). Every component contains an even number of them. Fix a spanning forest $\mathcal F_i$ of $\UQ_{l,c,i}$ (it contains no loop). By Cited result~\ref{s1:citEuler}, applied to a spanning tree of each component, there is a set $\mathcal T_i\subseteq E(\mathcal F_i)$ in which every vertex has odd degree iff it lies in $\mathrm{Odd}_i$. Then all degrees of $\UQ_{l,c,i}-\mathcal T_i$ are even, and $|\mathcal T_i|\le|E(\mathcal F_i)|\le\nu_l-1$. For every arc $a$ with $e_a\in\mathcal T_i$, output each edge of $a$ as a single edge; $a$ is a path in $Z$ and so has at most $|Z|-1\le M_l-1$ edges. Every component of $\UQ_{l,c,i}-\mathcal T_i$ that has an edge is connected with all degrees even, so it has a closed Euler trail (Cited result~\ref{s1:citEuler}; the statement is valid for multigraphs with loops). This gives at most $\nu_l$ closed trails for class $i$, which together use every edge of $\UQ_{l,c,i}-\mathcal T_i$ exactly once.

We record a closed trail $\mathcal W$ as a cyclic sequence $(a_1,\dots,a_k)$, $k\ge1$ (in this proof $j$ indexes positions along a trail), of distinct arcs of one class, each with an orientation: $a_j$ is traversed from its end $v_j^-$ to its end $v_j^+$, and $\parU(v_j^+)=\parU(v_{j+1}^-)$ for all $j$ (indices modulo $k$). The \emph{transition} at position $j$ is the pair $(v_j^+,v_{j+1}^-)$, and its \emph{node} is $\parU(v_j^+)$.

\begin{claim}[transitions]
In every closed trail of this kind, every transition consists of two distinct vertices of its node $Y$.
\end{claim}
\begin{proof}
Both vertices have parent $Y$, so both lie in $Y$. If $k\ge2$, then $a_j\ne a_{j+1}$ are distinct arcs of one class, which are vertex-disjoint (Step~1), so $v_j^+\ne v_{j+1}^-$. If $k=1$, the transition is $(v_1^+,v_1^-)$, the two ends of one arc, which are distinct by (F-d).
\end{proof}

\emph{Step 4: visit capping.} For a closed trail $\mathcal W$ and a node $Y$ let $\mathrm{vis}_Y(\mathcal W)$ be the number of transitions of $\mathcal W$ at $Y$. If $\mathrm{vis}_Y(\mathcal W)>\Tslot_Y$, \emph{split} $\mathcal W$ at $Y$ as follows. Write $\mathrm{vis}\defeq\mathrm{vis}_Y(\mathcal W)$. Rotate the cyclic sequence so that its last transition (position $k$) is at $Y$, and let $j_1<\dots<j_{\mathrm{vis}}=k$ be the positions of the transitions at $Y$. They cut $\mathcal W$ into $\mathrm{vis}$ consecutive segments $(a_1,\dots,a_{j_1})$, $(a_{j_1+1},\dots,a_{j_2})$, \dots, $(a_{j_{\mathrm{vis}-1}+1},\dots,a_{j_{\mathrm{vis}}})$; each segment begins with an arc whose first end has parent $Y$ and ends with an arc whose last end has parent $Y$. Group consecutive segments into $\lceil\mathrm{vis}/\Tslot_Y\rceil$ blocks of at most $\Tslot_Y$ segments. A block $(a_p,\dots,a_q)$ is again a closed trail of the same kind: its transitions are the transitions of $\mathcal W$ at positions $p,\dots,q-1$ together with the new transition $(v_q^+,v_p^-)$ at $Y$. By the argument of the transition claim, $v_q^+\ne v_p^-$ (two distinct arcs of one class, or the two ends of one arc). The block has as many transitions at $Y$ as it has segments, hence at most $\Tslot_Y$; and each transition of $\mathcal W$ at a node $Y'\ne Y$ lies in exactly one block. So a split at $Y$ replaces $\mathcal W$ by $\lceil\mathrm{vis}/\Tslot_Y\rceil\le1+\mathrm{vis}/\Tslot_Y$ closed trails, partitions the arcs of $\mathcal W$ among them, leaves every transition at every other node unchanged, and does not increase $\mathrm{vis}_{Y'}$ of any trail for $Y'\ne Y$.

Process the nodes $Y$ one after another in a fixed order and split every current trail at the current node. Since the number of $Y$-transitions of a trail never increases through splits at other nodes, at the end every trail $\mathcal W$ satisfies $\mathrm{vis}_Y(\mathcal W)\le\Tslot_Y$ for every $Y$. When $Y$ is processed, the number of new trails is at most $\sum_{\mathcal W}\mathrm{vis}_Y(\mathcal W)/\Tslot_Y$, summed over the trails present at that time; and $\sum_{\mathcal W}\mathrm{vis}_Y(\mathcal W)$, the total number of transitions at $Y$ in class $i$ of phase $c$, is not changed by any split. Call it $\mathrm{tr}_{c,i,Y}$ (for the fixed $l$).

A closed trail with $k$ arcs has exactly $k$ transitions, and splits preserve the total number of arcs. Hence, summing over all nodes $Y$, phases $c\in[4]$ and classes $i$, the total number of transitions $\sum_{Y,c,i}\mathrm{tr}_{c,i,Y}$ is at most the number of arcs of all non-demoted light parts of round $l$, which by (F-d), (F-a) and (F-c) is at most $\sum_Z14(J_Z+1)|Z|\le14(\bar J_l+1)\,n$ (the parts $Z$ are pairwise disjoint, so $\sum_Z|Z|\le n$). Every node $Y$ has round $r\le l-2$, so $\Tslot_Y\ge L_Y^2\ge(102\log_2\lambda_r)^2\ge(102\log_2\lambda_{l-2})^2$, by \eqref{s5:eqLY} and $\lambda_r\ge\lambda_{l-2}$ (Lemma~\ref{s2:lemTower}(a)). Therefore the number of visit-capping pieces over all $c$ and $i$ is
\[
\mathrm{cap}_l\ \le\ \sum_{Y,c,i}\frac{\mathrm{tr}_{c,i,Y}}{\Tslot_Y}\ \le\ \frac{14(\bar J_l+1)\,n}{(102\log_2\lambda_{l-2})^2}.
\]

\emph{Step 5: slots and pair multisets.} For every final trail $\mathcal W$ and every node $Y$, give the transitions of $\mathcal W$ at $Y$ pairwise distinct slots $\uslot\in[\Tslot_Y]$, by a fixed rule; this is possible because $\mathrm{vis}_Y(\mathcal W)\le\Tslot_Y$. For each good parent $Y$ of round $\le l-2$ and each $\uslot\in[\Tslot_Y]$ let $\mathfrak P_{Y,\uslot}$ be the multiset of the transitions at $Y$ with slot $\uslot$, over all final trails of all classes $i$ (for the fixed $l,c$), each regarded as an unordered pair. By the transition claim its elements are pairs of distinct vertices of $Y=V(\LU_{Y,l,c,\uslot})$; the same pair may occur several times, so $\mathfrak P_{Y,\uslot}$ is a multiset as in Cited result~\ref{s1:citDef7}. A vertex $v$ occurs in a pair at $Y$ only if $\parU(v)=Y$. We bound the number of its occurrences.
\begin{itemize}
\item \emph{Each arc end lies in at most one transition.} In a closed trail $(a_1,\dots,a_k)$ of the kind recorded in Step~3, the transitions are the $k$ pairs $(v_j^+,v_{j+1}^-)$ (indices modulo $k$); so the end $v_j^+$ of $a_j$ lies in the transition at position $j$ only, and the end $v_j^-$ in the transition at position $j-1$ only (for $k=1$ both ends lie in the single transition $(v_1^+,v_1^-)$, once each). Thus every end of every arc of a trail lies in exactly one transition of that trail. This applies to every final trail, including the blocks created in Step~4: a block $(a_p,\dots,a_q)$ is a closed trail of the same kind, and its transitions are exactly the pairs $(v_j^+,v_{j+1}^-)$, $p\le j<q$, together with $(v_q^+,v_p^-)$; the new transition carries the two ends $v_q^+$ and $v_p^-$, whose transitions at positions $q$ and $p-1$ of $\mathcal W$ are not transitions of the block. The final trails partition the arcs that are not output in Step~3, and the arcs output in Step~3 lie on no trail. Hence every arc end lies in at most one transition over all final trails of all classes.
\item \emph{Counting.} By the transition claim, $v$ is at most one of the two vertices of any pair, and each occurrence of $v$ in a pair is an arc end at $v$ (the end $v_j^+$ of $a_j$ or the end $v_{j+1}^-$ of $a_{j+1}$). By the first bullet, distinct occurrences are distinct arc ends, and an arc has at most one end equal to $v$, since its two ends are distinct (F-d). So the number of pairs of $\bigcup_{\uslot}\mathfrak P_{Y,\uslot}$ containing $v$, counted with multiplicity, is at most the number of arcs of $\Bdl_{l,c}$ having $v$ as an end. By (F-d) and (F-c), all of them are arcs of the unique round-$l$ light part $Z(v)$ containing $v$, and there are at most $|Z(v)|-1\le M_l-1$ of them.
\item \emph{Comparison with $t_Y$.} Let $r=r(Y)\le l-2$. By Lemma~\ref{s2:lemTower}(b) (row~12 of Table~\ref{s3:tabCOLJV}) and Lemma~\ref{s2:lemTower}(a), $M_l\le\lambda_{l-2}^{1.6}\le\lambda_r^{1.6}\le\lceil\lambda_r^{1.6}\rceil=t_Y$ (Definition~\ref{s3:defCOL}).
\end{itemize}
So $v$ lies in at most $M_l-1<t_Y$ pairs of $\bigcup_{\uslot}\mathfrak P_{Y,\uslot}$, and in particular of each $\mathfrak P_{Y,\uslot}$.

\emph{Step 6: routing.} Let $Y$ be a good parent of round $r\le l-2$ and $\uslot\in[\Tslot_Y]$. Then $(l,c,\uslot)$ is an index of $\IU$ for $Y$, since $r+2\le l\le R$. As $Y$ is not parent-bad, $\LU_{Y,l,c,\uslot}$ is $(2^{12}L_Y^4,t_Y)$-path connected through $\Zone_{Y,l,c,\uslot}$ (Definition~\ref{s5:defStages}). This property depends only on the pair $(\LU_{Y,l,c,\uslot},\Zone_{Y,l,c,\uslot})$, which is fixed by the stage-1 outcome, so it may be applied to the multiset $\mathfrak P_{Y,\uslot}$, which is formed afterwards (Lemma~\ref{s3:lemMonotone}(iii)). By Step~5 every vertex lies in at most $M_l-1\le t_Y$ of its pairs. Hence there are pairwise edge-disjoint paths $\mathcal C_\pi\subseteq\LU_{Y,l,c,\uslot}$, one for each element $\pi$ of $\mathfrak P_{Y,\uslot}$ (counted with multiplicity), such that $\mathcal C_\pi$ joins the two vertices of $\pi$, has length at most $2^{12}L_Y^4$, and has all its interior vertices in $\Zone_{Y,l,c,\uslot}$. These paths are the connectors. Define $\LentU_{l,c}$ as the union of the edge sets of all connectors, over all $Y$ and $\uslot$. This proves the inclusion in (i) and the property of the connectors in (i).

\emph{Step 7: the cycles.} Let $\mathcal W=(a_1,\dots,a_k)$ be a final trail, and let $\mathcal C_j$ be the connector of its transition $(v_j^+,v_{j+1}^-)$, which lies at the node $Y_j$ and has slot $\uslot_j$. Let $C(\mathcal W)$ be the closed walk that traverses $a_1$ from $v_1^-$ to $v_1^+$, then $\mathcal C_1$ from $v_1^+$ to $v_2^-$, then $a_2$, then $\mathcal C_2$, and so on, and returns to $v_1^-$ along $\mathcal C_k$.

\begin{claim}[cycles]
$C(\mathcal W)$ is a cycle of $G$ of length at least $3$.
\end{claim}
\begin{proof}
(1) The arcs $a_1,\dots,a_k$ are paths with pairwise disjoint vertex sets, since they are distinct arcs of one class (Step~1). (2) The interior vertices of $\mathcal C_j$ lie in $\Zone_{Y_j,l,c,\uslot_j}$, which is a round-$l$ phase-$c$ zone; every vertex of every arc of $\Bdl_{l,c}$ lies in no such zone (F-d). So no interior vertex of a connector lies on an arc of $\mathcal W$. (3) For $j\ne j'$ the interiors of $\mathcal C_j$ and $\mathcal C_{j'}$ are disjoint: if $Y_j\ne Y_{j'}$, zones of distinct light parts are disjoint; if $Y_j=Y_{j'}$, then $\uslot_j\ne\uslot_{j'}$ (Step~5), and the two zones are again disjoint (Lemma~\ref{s5:lemZones}(iii) in both cases). (4) Each $\mathcal C_j$ is a path from $v_j^+$ to $v_{j+1}^-$, and these two vertices are distinct (transition claim), so $\mathcal C_j$ has at least one edge. By (1)--(4), $C(\mathcal W)$ visits each of its vertices once before closing, so it is a cycle with at least two edges. Suppose it had exactly two. Then $k=1$, $a_1$ is the single edge $v_1^-v_1^+$ and $\mathcal C_1$ is the single edge $v_1^+v_1^-$. As $G$ is simple, these are the same edge. But $a_1\subseteq E_l(Z)$ for a round-$l$ light part $Z$, while $\mathcal C_1\subseteq E(H_{Y_1})\subseteq E_{r}(Y_1)$ with $r=r(Y_1)\le l-2$, and these edge sets are disjoint (F-c). This is a contradiction, so the length is at least $3$.
\end{proof}

\emph{Step 8: exact cover.} The objects for $(l,c)$ are the single edges of Step~3 and the cycles $C(\mathcal W)$ of Step~7.
\begin{itemize}
\item \emph{Arc edges.} Arcs of one $Z$ are pairwise edge-disjoint (F-d), and arcs of distinct $Z$ lie in the disjoint sets $E_l(Z)$ (F-c). Each arc lies in exactly one class; there it is either an edge $e_a\in\mathcal T_i$, and then its edges are output as single edges, or it lies on exactly one Euler trail and, after the splits (which partition the arcs of a trail), on exactly one final trail $\mathcal W$, and then its edges lie on $C(\mathcal W)$ only.
\item \emph{Connector edges.} For fixed $(Y,\uslot)$ the connectors are pairwise edge-disjoint (Step~6). Connectors of distinct $(Y,\uslot)$ lie in distinct classes $\LU_{Y,l,c,\uslot}$: for a fixed $Y$ these are distinct colour classes of one colouring, and for distinct $Y$ they lie in the disjoint sets $E(H_Y)\subseteq E_{r(Y)}(Y)$ (F-c). Each connector belongs to exactly one transition and hence to exactly one cycle $C(\mathcal W)$.
\item \emph{Arc edges versus connector edges.} Arc edges lie in sets $E_l(Z)$ and connector edges in sets $E_r(Y)$ with $r\le l-2$; these are disjoint (F-c). So $\LentU_{l,c}\cap E(\Bdl_{l,c})=\emptyset$.
\end{itemize}
Hence the objects partition $E(\Bdl_{l,c})\cup\LentU_{l,c}$, and no other edge is used; this is (iii). For distinct pairs $(l,c)\ne(l',c')$ the sets $\LentU_{l,c}$ and $\LentU_{l',c'}$ are disjoint: within one good parent $Y$ they lie in distinct colour classes of one colouring (the index $(l,c,\uslot)$ differs), and for distinct parents $Y\ne Y'$ they lie in the disjoint sets $E(H_Y)\subseteq E_{r(Y)}(Y)$ and $E(H_{Y'})\subseteq E_{r(Y')}(Y')$ (F-c). Every choice made above is a fixed rule applied to the stage-1 outcome and to $\Bdl_{l,c}$.

\emph{Step 9: counting.} For a fixed class $i$, Step~3 outputs at most $(\nu_l-1)(M_l-1)$ single edges and at most $\nu_l$ closed trails, together at most $\nu_lM_l$ objects before capping. There are at most $\mathrm{ncl}_l=14(\bar J_l+1)M_l$ classes and $4$ phases. With the capping pieces of Step~4, the number of objects summed over $c$ is at most
\[
4\cdot14(\bar J_l+1)M_l\cdot\nu_lM_l+\mathrm{cap}_l\ \le\ 4\cdot14(\bar J_l+1)M_l(M_l+1)\nu_l+\mathrm{cap}_l ,
\]
which is (ii) for the round $l$. For the sum over $l$, first let $3\le l\le R$. By (F-a), $\bar J_l+1\le\log_2\log_2M_l\le M_l$; moreover $M_l+1\le2M_l$, $\nu_l\le2.74n/P_{l-2}$ (Step~2), and $P_{l-2}\ge M_l^{13}$ (Lemma~\ref{s2:lemTower}(b)). So
\[
4\cdot14(\bar J_l+1)M_l(M_l+1)\nu_l\ \le\ 56\cdot2\cdot2.74\,n\,M_l^{3}/M_l^{13}\ \le\ 307\,n/M_l ,
\]
and $\sum_l307n/M_l\le614n/\Dstar$ because $\sum_l1/M_l\le2/\Dstar$ (Lemma~\ref{s2:lemTower}(b)).

Next, the capping pieces. By Step~4 and (F-a), $\bar J_l+1\le\log_2\log_2M_l$. Moreover $M_l\le(\Aexp\log\lambda_{l-2})^{2\Aexp}$ (Lemma~\ref{s2:lemTower}(b)), so that $\log_2M_l\le2\Aexp\log_2(\Aexp\log_2\lambda_{l-2})$. With
\[
\eta(x)\defeq\frac{\log_2\bigl(2\Aexp\log_2(\Aexp\log_2x)\bigr)}{(\log_2x)^2}\qquad(x\ge\log_2\Dstar),
\]
this gives $\mathrm{cap}_l\le\frac{14}{102^2}\,n\,\eta(\lambda_{l-2})\le\frac{n}{743}\,\eta(\lambda_{l-2})$. We claim that $\eta$ is non-increasing on $[\log_2\Dstar,\infty)$ and that $\eta(2^{x/\Aexp})\le\eta(x)/2$ there. Write $x_1\defeq\log_2x\ge\log_2\log_2\Dstar$, so that $x_1\ge2^8$ by \Gc{1}(a) at $x_1$, $x_2\defeq\log_2(\Aexp x_1)\ge14$ and $x_3\defeq\log_2(2\Aexp x_2)\ge11$, so that $\eta=x_3/x_1^2$. Then
\[
\frac{d}{dx_1}\ln\eta\ =\ \frac{1}{x_3\,x_2\,x_1\,(\ln2)^2}-\frac2{x_1}\ <\ 0 .
\]
Further, $\log_2(\Aexp\log_2 2^{x/\Aexp})=\log_2x=x_1$, so $\eta(2^{x/\Aexp})=\Aexp^2\log_2(2\Aexp x_1)/x^2$, while $\eta(x)/2=x_3/(2x_1^2)$ with $x_3\ge1$. Since $x=2^{x_1}$ and $2\Aexp^2x_1^2\log_2(2\Aexp x_1)\le4\Aexp^3x_1^3\le(2^{14}\Aexp x_1^3)^2\le2^{2x_1}$ (by $\log_2z\le z$ and \Gc{1}(b) at $x_1$), we get $\eta(2^{x/\Aexp})\le\eta(x)/2$. As $\lambda_{r'}\ge2^{\lambda_{r'+1}/\Aexp}$ (Lemma~\ref{s2:lemTower}(a)), $\eta(\lambda_{r'})\le\eta(2^{\lambda_{r'+1}/\Aexp})\le\eta(\lambda_{r'+1})/2$ for $r'<R$. Hence $\sum_{l=3}^R\eta(\lambda_{l-2})=\sum_{r'=1}^{R-2}\eta(\lambda_{r'})\le2\eta(\lambda_{R-2})\le2\eta(\log_2\Dstar)$, because $\lambda_{R-2}\ge\log_2\Dstar$. Altogether
\[
\sum_{l=3}^R\mathrm{cap}_l\ \le\ \frac{2}{743}\,\eta(\log_2\Dstar)\,n\ \le\ \frac{\log_2\bigl(2\Aexp\log_2(\Aexp\log_2\log_2\Dstar)\bigr)}{371\,(\log_2\log_2\Dstar)^2}\,n .
\]
Adding the two sums gives $\epsChain(\Dstar)n$, which is (ii). Finally $\eta(\log_2\Dstar)\to0$ and $614/\Dstar\to0$ as $\Dstar\to\infty$, so $\epsChain(\Dstar)\to0$.
\end{proof}

\subsection{Demoted parts and Lemma K-RED}

\begin{lemma}[fallback for demoted light parts]\label{s5:lemDemoted}
Fix a stage-1 outcome and let $Y$ be a demoted light part of round $r$. Then $Y$ is not a good parent; in particular no class $\LU_{Y,\cdot}$ of $Y$ is used in Lemma~\ref{s5:lemParent}. Let $\mathcal G^{\mathrm{VX}}_Y$ be the good event of Theorem~\ref{s4:thmVXp} for $O\defeq X_Y$, an event over the VX$^+$ labels of $Y$ (stage~3) that depends only on $X_Y$ and these labels and has probability at least $1-\eta_{\mathrm{VX}}(|Y|)\ge1/2$. On $\mathcal G^{\mathrm{VX}}_Y$, for every edge set $E$ with $E(X_Y)\subseteq E\subseteq E_r(Y)$, the set $E$ decomposes into at most $\cVX|Y|$ objects.
\deps{\ref{s4:thmVXp}, \ref{s2:propStructure}, \ref{s3:lemCOLJV}, \ref{s2:lemTower}, \ref{s1:condG1}, \ref{s1:condG3}, \ref{s5:defStages}, \ref{s5:lemParent}}
\end{lemma}

\begin{proof}
A demoted part lies in $\Bad$, so it is not a good parent (Definition~\ref{s5:defStages}), and Lemma~\ref{s5:lemParent}(i) uses only classes of good parents. Put $N\defeq|Y|$ and $L\defeq L_Y$. By Proposition~\ref{s2:propStructure}(i),(iv), $X_Y$ is a spanning $(2^{-6},s_r/2)$-expander on $Y$ and $N\ge P_r/2\ge\Nzero$ (\Gc{3}). The threshold of Theorem~\ref{s4:thmVXp} at $\eps=2^{-6}$ is $s_r/2\ge2^{151}L^{38}\log_2L$; this is row~8(c) of Table~\ref{s3:tabCOLJV} (Lemma~\ref{s3:lemCOLJV}), and it holds in every round $r\le R$, including $r\ge R-1$. Explicitly, $s_r\ge\lambda_r^{100}$ (Proposition~\ref{s2:propStructure}(i)) and $L\le2\lambda_r$ (Lemma~\ref{s2:lemTower}(b), valid for every ancestor of every round $r\le R$) give $2^{152}L^{38}\log_2L\le2^{190}\lambda_r^{38}\cdot2\lambda_r=2^{191}\lambda_r^{39}\le\lambda_r^{100}\le s_r$, since $\log_2\lambda_r\ge2^8$ (\Gc{1}(a)). Theorem~\ref{s4:thmVXp} provides random labels and a good event depending only on $O=X_Y$, of probability at least $1-\eta_{\mathrm{VX}}(N)\ge1/2$, since $N\ge\Nzero$ by \Gc{3}. On this event its conclusion holds simultaneously for every graph on the vertex set $Y$ that contains $O$. Every edge of $E_r(Y)$ has both ends in $Y$ (Proposition~\ref{s2:propStructure}(iii)), so for $E(X_Y)\subseteq E\subseteq E_r(Y)$ the graph $(Y,E)$ contains $O$, and Theorem~\ref{s4:thmVXp} gives $\f(E)\le80N=\cVX|Y|$.
\end{proof}

Put $\Kstd\defeq\bigsqcup_{l}\bigsqcup_{Z\in\Std_l}E_l(Z)$. Here $\Std_l$ is the deterministic family of Definition~\ref{s2:defHBtp} (it contains the GC-parts and the pre-parts failing (L1), and it does \emph{not} contain the random demotions of Definition~\ref{s5:defStages}); so $\Kstd$ is a function of the run alone, and every demoted light part is still a light part.

\begin{lemma}[Lemma K-RED]\label{s5:lemKRED}
Fix the valid run, a stage-1 outcome, and stage-3 outcomes in the good events of all light parts: the event $\mathcal G_Z$ of Lemma~\ref{s5:lemChild} for every non-demoted light part $Z$, and the event $\mathcal G^{\mathrm{VX}}_Z$ of Lemma~\ref{s5:lemDemoted} for every demoted light part $Z$. Let $(\Lentext(Y))_Y$ be \emph{any} family, indexed by the light parts $Y$, such that for every light part $Y$
\[
\Lentext(Y)\ \subseteq\ \Bigl(\bigcup_{l,j}\LJS_{Y,l,j}\Bigr)\cup\Bigl(\bigcup_l\LJV_{Y,l}\Bigr),\qquad\text{and}\qquad \Lentext(Y)=\emptyset\ \text{ if $Y$ is demoted.}
\]
Process the rounds $l=R,R-1,\dots,1$ (children before parents). At round $l$:
\begin{enumerate}
\item[(1)] for every non-demoted light part $Z$ of round $l$, let $\LentU(Z)\defeq E(H_Z)\cap\bigcup_{l'\ge l+2,\ c\in[4]}\LentU_{l',c}$ be the set of U-lent edges of $Z$ used by Lemma~\ref{s5:lemParent} at the rounds already processed, put
\[
H_0(Z)\defeq E_l(Z)\setminus\bigl(\LentU(Z)\cup\Lentext(Z)\bigr),
\]
and decompose $H_0(Z)$ by Lemma~\ref{s5:lemChild} into $H^{\mathrm{obj}}(Z)$ (as objects) and arcs;
\item[(2)] for every demoted light part $Z$ of round $l$, decompose $E_l(Z)$ by Lemma~\ref{s5:lemDemoted};
\item[(3)] if $l\ge3$, apply Lemma~\ref{s5:lemParent} to the bundles $\Bdl_{l,c}$, $c\in[4]$, formed from the arcs of step~(1).
\end{enumerate}
Output the long cycles $\bigcup_l\Cyc_l$ and the edges of $E_0$ (as single edges) directly. Then the objects produced form a decomposition of
\[
E(G)\setminus\Kstd\setminus\bigcup_Y\Lentext(Y)
\]
into at most
\[
\bigl(\Dstar/2+\cKRED+\epsK(\Dstar)\bigr)\,n\;+\;\cVX\,\mathrm{dem}\;+\;\cPV\,\mathrm{lp}
\]
objects, where $\epsK(\Dstar)\defeq\epsChain(\Dstar)$. Itemized: long cycles and $E_0$, at most $n+\Dstar n/2$; child vortices, at most $\cPV\sum_Z|\Pl(Z)|\le\cPV(2n+\mathrm{lp})$; U-chaining, at most $\epsChain(\Dstar)n$; demoted parts, at most $\cVX\,\mathrm{dem}$. The exact sum of these items is at most $(\Dstar/2+\cKREDexact+\epsK(\Dstar))n+\cVX\,\mathrm{dem}+\cPV\,\mathrm{lp}$; the constant $\cKRED$ is kept for the cost line. No edge of $\Kstd$ and no edge of any $\Lentext(Y)$ is used. Moreover, the objects produced at round $l$ depend only on the stage-1 and stage-3 outcomes, on the sets $\Lentext(Z)$ of the light parts $Z$ of round $l$, and on what was produced at the rounds $>l$.
\deps{\ref{s5:lemChild}, \ref{s5:lemParent}, \ref{s5:lemDemoted}, \ref{s2:propStructure}, \ref{s2:propParentless}, \ref{s2:lemEL}, \ref{s3:defCOL}, \ref{s5:defStages}}
\end{lemma}

\begin{proof}
\emph{The edge partition.} By Proposition~\ref{s2:propStructure}(iii),
\[
E(G)\ =\ \bigsqcup_l\Cyc_l\ \sqcup\ E_0\ \sqcup\ \bigsqcup_{Y\text{ light}}E_{r(Y)}(Y)\ \sqcup\ \Kstd ,
\]
where $\Cyc_l$ is read as an edge set. For every light $Y$, $\Lentext(Y)\subseteq\Lend_Y\subseteq E(H_Y)\subseteq E_{r(Y)}(Y)$ (Definition~\ref{s3:defCOL} and Proposition~\ref{s2:propStructure}(iii)). Hence the set to be decomposed is
\[
\mathcal D\ \defeq\ \bigsqcup_l\Cyc_l\ \sqcup\ E_0\ \sqcup\ \bigsqcup_{Y\text{ light}}\bigl(E_{r(Y)}(Y)\setminus\Lentext(Y)\bigr).
\]
The edges of a light part $Y$ of round $r$ have both ends in $Y$ (Proposition~\ref{s2:propStructure}(iii)). The two kinds of edges combined by Lemma~\ref{s5:lemParent} are separated also by Lemma~\ref{s2:lemEL}: an arc edge of a round-$l$ part lies in $E(G_l)$, so it does not have both ends in any parent $Y$ of round $\le l-2$, whereas every connector edge of $Y$ has both ends in $Y$. The construction below applies Lemma~\ref{s5:lemDemoted} and Lemma~\ref{s5:lemChild} only to edge sets of light parts, and Lemma~\ref{s5:lemParent} only to arc edges and U-lent edges of light parts; it never applies anything to an edge set $E_l(Z)$ with $Z\in\Std_l$.

\emph{The construction is well defined.} At round $l$, the set $\LentU(Z)$ of a round-$l$ light part $Z$ consists of edges used by Lemma~\ref{s5:lemParent} at rounds $l'\ge l+2$, all of which were processed before round $l$; Lemma~\ref{s5:lemParent} at a round $l'$ uses classes only of good parents of rounds $\le l'-2$, so no later round adds edges of $Z$ to any $\LentU_{l',c}$. Thus $\LentU(Z)$ is final when $Z$ is processed. The hypotheses of Lemma~\ref{s5:lemChild} hold for $H_0(Z)$: $H_0(Z)\subseteq E_l(Z)$, and $\LentU(Z)\cup\Lentext(Z)\subseteq\Lend_Z$ (the U-lent classes and the JS- and JV-lent classes are colour classes of $\Lend_Z$, Definition~\ref{s3:defCOL}(ii)), which is disjoint from $\Own_Z$; so $\Own_Z\subseteq H_0(Z)$. The stage-3 outcome of $Z$ lies in $\mathcal G_Z$, which does not depend on $H_0(Z)$. The hypotheses of Lemma~\ref{s5:lemParent} at round $l$ are those of Lemma~\ref{s5:lemChild} for the round-$l$ parts. For a demoted $Z$ of round $l$, $\Lentext(Z)=\emptyset$ by assumption and no edge of $Z$ lies in any $\LentU_{l',c}$ (Lemma~\ref{s5:lemDemoted}), so all of $E_l(Z)$ is available, and $E(X_Z)\subseteq E_l(Z)$ (Proposition~\ref{s2:propStructure}(iii)); Lemma~\ref{s5:lemDemoted} applies with $E\defeq E_l(Z)$. The objects produced at round $l$ are functions of the stage-1 and stage-3 outcomes, of $\Lentext(Z)$ and $\LentU(Z)$ for the round-$l$ parts $Z$, and of the bundles $\Bdl_{l,c}$, which are formed from these; and $\LentU(Z)$ is determined by the rounds $>l$. This proves the last assertion.

\emph{Every edge of $\mathcal D$ lies in exactly one object, and no other edge is used.}
\begin{itemize}
\item Long cycles and $E_0$ are output directly, once each.
\item Let $Y$ be a demoted light part of round $r$. All of $E_r(Y)$ is covered by the objects of Lemma~\ref{s5:lemDemoted}, each edge once. No other object uses an edge of $E_r(Y)$: the objects of Lemmas~\ref{s5:lemChild} and~\ref{s5:lemDemoted} for another part $Z'$ use only edges of $E_{r(Z')}(Z')$; the objects of Lemma~\ref{s5:lemParent} use only arc edges of non-demoted parts and U-lent edges of good parents (Lemma~\ref{s5:lemParent}(i),(iii)).
\item Let $Y$ be a non-demoted light part of round $r$. The U-lent classes are disjoint from the JS- and JV-lent classes, so $E_r(Y)\setminus\Lentext(Y)=\LentU(Y)\sqcup H_0(Y)$. An edge of $\LentU(Y)$ lies in $\LentU_{l,c}$ for exactly one pair $(l,c)$ (these sets are pairwise disjoint), with $l\ge r+2$, and is covered by exactly one object of Lemma~\ref{s5:lemParent} at $(l,c)$. An edge $e\in H_0(Y)$ lies either in $H^{\mathrm{obj}}(Y)$, and is covered by exactly one of its objects, or on exactly one arc of $Y$, of some phase $c$; that arc belongs to $\Bdl_{r,c}$ (and $r\ge3$, by Lemma~\ref{s5:lemChild}(d)), and $e$ is covered by exactly one object of Lemma~\ref{s5:lemParent} at $(r,c)$. No other object uses $e$: it is not in any $\LentU_{l',c'}$ (those edges were removed as $\LentU(Y)$), it is an arc edge only of $Y$'s own arcs, and objects of other parts use their own edges.
\item Edges of $\Kstd$ and of the sets $\Lentext(Y)$ are used by no object: Lemma~\ref{s5:lemDemoted} is applied only to demoted light parts, which have $\Lentext=\emptyset$; Lemma~\ref{s5:lemChild} is applied to $H_0(Z)$, which excludes $\Lentext(Z)$; Lemma~\ref{s5:lemParent} uses arc edges (contained in the sets $H_0(Z)$) and U-lent edges (disjoint from JS- and JV-lent classes); and all these edges lie in the edge sets of light parts, which are disjoint from $\Kstd$. Zones are vertex sets and carry no edges.
\end{itemize}
Every object is a cycle of length at least $3$ or a single edge, by Lemmas~\ref{s5:lemChild}, \ref{s5:lemParent} and~\ref{s5:lemDemoted} and the definition of $\Cyc_l$.

\emph{Cost.}
\begin{itemize}
\item Long cycles: at most $n$; single edges of $E_0$: $|E_0|<\Dstar n/2$ (Proposition~\ref{s2:propStructure}(iii)).
\item Child vortices: by Lemma~\ref{s5:lemChild}(a), at most $\cPV\sum_Z|\Pl(Z)|$ over the non-demoted light parts $Z$. By \eqref{s5:eqPl}, which is Proposition~\ref{s2:propParentless}(iii) applied with the set $\Bad$ of Definition~\ref{s5:defStages} (demoted \emph{or} parent-bad parts), this is at most $\cPV(2n+\mathrm{lp})=738n+\cPV\,\mathrm{lp}$.
\item U-chaining: by Lemma~\ref{s5:lemParent}(ii), at most $\epsChain(\Dstar)n$ over all rounds $l\ge3$ and phases $c$.
\item Demoted parts: by Lemma~\ref{s5:lemDemoted}, at most $\cVX\sum_{Y\text{ demoted}}|Y|=\cVX\,\mathrm{dem}$.
\end{itemize}
The total is at most $(\Dstar/2+1+738+\epsChain(\Dstar))n+\cVX\,\mathrm{dem}+\cPV\,\mathrm{lp}=(\Dstar/2+\cKREDexact+\epsK(\Dstar))n+\cVX\,\mathrm{dem}+\cPV\,\mathrm{lp}$, which is at most the stated bound with $\cKRED$ in place of $\cKREDexact$.
\end{proof}

\paragraph{The for-all form and the lent edges.} Lemma~\ref{s5:lemKRED} holds for \emph{every} family $(\Lentext(Y))$ of the stated kind, and the objects of round $l$ depend on that family only through the sets $\Lentext(Z)$ of the round-$l$ parts. The lemma is applied later with $\Lentext(Y)$ equal to the set of JS-lent edges of $Y$ used on output cycles of junction chaining together with the JV-lent junction edges of $Y$ used by lifted cycles, at rounds $\ge r(Y)+2$; such a family is of the stated kind, and it is fixed before round $r(Y)$ is processed, which is all that the last assertion of the lemma requires. With the used U-lent edges this gives the full set of used lent edges of $Y$, all of which are excluded from $H_0(Y)$; in particular, used JV junction edges are removed from $H_0$ and are not covered twice. Unused lent edges of $Y$ stay in $H_0(Y)$ as junk, which Lemma~\ref{s5:lemChild} accepts because it holds for every $H_0(Y)$ between $\Own_Y$ and $E_{r(Y)}(Y)$.

\begin{remark}[the constants $369$ and $745$]\label{s5:remConstants}
(a) The constant $\cPV=369$ of Lemma~\ref{s4:lemPV} is an upper bound for $28\cdot8+64+16=304$. The proof of Lemma~\ref{s4:lemPV} shows that a vortex step $j$ costs at most $4|\vW_j|+12(2|\Pl\cap\vU_{j+1}|+2|\vW_j|)\le28|\Pl\cap\vU_j|$ objects (equation~\eqref{s4:eqPVstep}, using $\Pl\cap\vU_j=\vW_j\sqcup(\Pl\cap\vU_{j+1})$). The good event (G5) gives $\sum_j|\Pl\cap\vU_j|\le8|\Pl|$, so all steps together cost at most $28\cdot8|\Pl|=224|\Pl|$. The finish costs at most $64|\Pl|$ objects for the edges with both ends in $\Pl\cap\vU_{\vJ}$ (Fact~\ref{s1:factEG0}(b), applied to at most $64|\Pl|/L$ vertices) and at most $8|\Pl\cap\vU_{\vJ}|\le16|\Pl|$ stripped single edges (proof of Lemma~\ref{s4:lemPV}; the bound $8|\Pl\cap\vU_{\vJ}|\le8|\Pl|$ also holds, and the weaker factor $16$ suffices; nothing depends on it). So at most $304|\Pl|$ objects are output; the constant $369$ is not optimized, and it is the value used in the cost line.

(b) The constant $\cKRED=745$ exceeds the exact accounting $\cKREDexact=739=1+2\cdot369$ of Lemma~\ref{s5:lemKRED} by $6$, so $745\ge1+2\cdot369+\epsK(\Dstar)$ whenever $\epsK(\Dstar)\le6$. This holds under \Gc{4} (Condition~\ref{s1:condG4}), whose left-hand side contains $\epsK$ as a summand. The bound of Lemma~\ref{s5:lemKRED} carries $\epsK$ separately in any case, so this slack is not needed; the value $745$ is the one used in the cost line.

(c) The constants $369$ and $745$ depend on the hierarchy only through three facts: each vertex lies in at most one light part per round (Proposition~\ref{s2:propStructure}(iv); used in Lemma~\ref{s5:lemZones}(i), in the vertex-disjointness of the arcs of one class, and in \eqref{s5:eqPl}); the parentless count $\sum_Z|\Pl(Z)|\le2n+\mathrm{lp}$ (Proposition~\ref{s2:propParentless}(iii); used in \eqref{s5:eqPl}; the size comparison $|Y|\ge|Z|L_Z^4$ of Proposition~\ref{s2:propParentless}(ii) is not used here, cf.\ (F-b) in the proof of Lemma~\ref{s5:lemParent}); and the overlap counts, namely at most $1.37n/P_r$ ancestors of round $r$ (Proposition~\ref{s2:propOV}, (K1)) and $\nu_l\le2.74n/P_{l-2}$ (Lemma~\ref{s2:lemTower}(b)), which enter only the terms $\epsU(\Dstar)n$ and $\epsChain(\Dstar)n$. The value $369$ is combinatorics of the four-phase P-vortex and does not depend on colour probabilities; $80$ is the constant of Theorem~\ref{s4:thmVXp}; and neither constant depends on the Markov multiplier used to fix the stage-1 outcome.
\deps{\ref{s4:lemPV}, \ref{s5:lemKRED}, \ref{s2:propStructure}, \ref{s2:propParentless}, \ref{s2:propOV}, \ref{s2:lemTower}, \ref{s4:thmVXp}, \ref{s1:condG4}, \ref{s1:factEG0}}
\end{remark}

\section{Chaining in standalone parts: HCC-P and JS-LC; hub concentration; the MIX-C assembly}\label{s6:sec}

This section contains the part of the proof that treats the standalone pre-parts of a valid $\HBtp$ run. It has four parts.
\begin{itemize}
\item A purely combinatorial \emph{routing engine}: Lemma GATE, clusters, Lemma MED, Lemma EQ-LPT, Lemma PAR and Lemma HCC-P. HCC-P closes layered systems of clusters into few cycles through junction paths.
\item The \emph{hub-concentration} bounds CONC and CONC-L. For every valid $\HBtp$ run and every designation they give $\sum_{(Y,l)}m_{Y,l}\le\epsCONC(\Dstar)\,n$.
\item The stage-2 \emph{lending statuses}, and Lemma JS-LC, which chains the beads of all standalone parts of one round through lent classes of ancestors two or more rounds back. The single edges left over form the \emph{J-set} $J_l$, whose interface is Lemma J$^+$.
\item The \emph{MIX-C assembly}. For every J-consumer (a rule that disposes of the J-set, defined in the last subsection) it turns all of this into a decomposition of $E(G)$ with an explicit cost.
\end{itemize}
There are no VX-parts ($\mathcal V=\emptyset$): every standalone pre-part is chained.

Throughout, $G$ is the input graph (Convention~\ref{s1:convGraphs}). Orientations, directed paths and directed cycles are those of digraphs obtained by orienting edges of $G$; each edge receives exactly one direction. A digraph is \emph{acyclic} if it has no directed cycle.

\subsection{The routing engine: GATE, clusters, MED, EQ-LPT, PAR, HCC-P}

\begin{lemma}[Lemma GATE, forward direction]\label{s6:lemGATE}
Let $F\subseteq E(G)$ and let $\vec F$ be an orientation of $F$ with $d^+_{\vec F}(v)=d^-_{\vec F}(v)$ at every vertex $v$. Let $\mathcal W$ be a set of arcs of $\vec F$ such that $\vec F-\mathcal W$ is acyclic. Then $F$ decomposes into at most $|\mathcal W|$ cycles of $G$. More precisely, $\vec F$ is the arc-disjoint union of directed cycles, each of length at least $3$, each the orientation of a cycle of $G$, and each containing an arc of $\mathcal W$; there are at most $|\mathcal W|$ of them.
\deps{\ref{s1:defObject}}
\end{lemma}

\begin{proof}
We first show that $\vec F$ is an arc-disjoint union of directed cycles of length at least $3$, by induction on the number of arcs. If $\vec F$ has no arc there is nothing to prove.

Otherwise let $v_0\to v_1\to\dots\to v_k$ ($k\ge1$) be a directed path in $\vec F$ with pairwise distinct vertices and with $k$ maximal. The vertex $v_k$ has in-degree at least $1$, so by balance it has an out-arc $v_k\to x$. By maximality of the path, $x=v_i$ for some $i\le k$. Since $G$ has no loops, $i\neq k$. So $C:=v_i\to v_{i+1}\to\dots\to v_k\to v_i$ is a directed cycle with distinct vertices, of length $k-i+1\ge2$.

Suppose the length were $2$. Then $i=k-1$, and $\vec F$ would contain the arcs $v_{k-1}\to v_k$ and $v_k\to v_{k-1}$. These arcs come from two distinct edges of $F$ with the same ends, which is impossible since $G$ is simple. So $C$ has length at least $3$, and its underlying edge set is a cycle of $G$.

Deleting the arcs of $C$ lowers in- and out-degree by $1$ at every vertex of $C$ and changes nothing elsewhere. So the remaining digraph is still balanced, and the induction hypothesis applies to it.

Let $C_1,\dots,C_p$ be the directed cycles obtained. Their underlying cycles partition $F$, so they form a decomposition of $F$ into cycles in the sense of Definition~\ref{s1:defObject}. Each $C_i$ contains an arc of $\mathcal W$, since otherwise $C_i$ would be a directed cycle of $\vec F-\mathcal W$. The $C_i$ are arc-disjoint, so $p\le|\mathcal W|$.
\end{proof}

\begin{definition}[clusters]\label{s6:defCluster}
A \emph{cluster} is a triple $\clu=(A_\clu,U_\clu,B_\clu)$ with the following parts.
\begin{itemize}
\item $A_\clu$ (the \emph{hubs} of $\clu$) and $U_\clu$ (the \emph{ports} of $\clu$) are disjoint vertex sets. Put $V(\clu):=A_\clu\cup U_\clu$.
\item $B_\clu\subseteq E(G[A_\clu\cup U_\clu])$ is the set of \emph{beads}. No bead has both ends in $A_\clu$.
\end{itemize}
The \emph{bead count} is $b(\clu):=\sum_{h\in A_\clu}\deg_{B_\clu}(h)/2+e(B_\clu[U_\clu])$. The cluster $\clu$ is \emph{parity-clean} if every hub has even $B_\clu$-degree.

An \emph{admissible orientation} of $\clu$ is an acyclic orientation of $B_\clu$ in which every hub $h$ has $d^+(h)=d^-(h)$. Given an admissible orientation, put for every port $u$
\[
\exc(u):=d^+(u)-d^-(u),\qquad \exc(u)^+:=\max(\exc(u),0),\qquad \exc(u)^-:=\max(-\exc(u),0),
\]
and define the \emph{load} $\Phi(\clu):=\sum_{u\in U_\clu}\exc(u)^+$.

Let $\clu$ be parity-clean and let $\prec$ be a linear order on $U_\clu$. The \emph{median orientation} $\MED(\prec)$ orients $B_\clu$ as follows.
\begin{itemize}
\item At a hub $h$ with $B_\clu$-neighbours $u_1\prec u_2\prec\dots\prec u_{2d}$, it orients $u_i\to h$ for $i\le d$ and $h\to u_i$ for $i>d$. These neighbours are ports, since no bead joins two hubs.
\item A port--port bead $uv$ with $u\prec v$ is oriented $u\to v$ ($\prec$-upwards).
\end{itemize}
Every bead is oriented exactly once: a bead at a hub has its other end at a port, so it is oriented by the rule at that hub only.
\deps{none}
\end{definition}

\begin{lemma}[Lemma MED]\label{s6:lemMED}
Let $\clu$ be a parity-clean cluster and $\prec$ a linear order on $U_\clu$.
\begin{enumerate}
\item[(a)] $\MED(\prec)$ is an admissible orientation. There is a function $\pot:A_\clu\cup U_\clu\to(0,1)$ that strictly increases along every arc of $\MED(\prec)$.
\item[(b)] For every admissible orientation of $\clu$, in particular for $\MED(\prec)$, and every port $u$: $|\exc(u)|\le\deg_{B_\clu}(u)$ and $\exc(u)\equiv\deg_{B_\clu}(u)\pmod 2$. Moreover $\sum_{u\in U_\clu}\exc(u)=0$ and $\Phi(\clu)\le b(\clu)$.
\end{enumerate}
\deps{\ref{s6:defCluster}}
\end{lemma}

\begin{proof}
(a) Let $\mathrm{rk}:U_\clu\to\{1,\dots,|U_\clu|\}$ be the bijection realising $\prec$, and put $\pot(u):=\mathrm{rk}(u)/(|U_\clu|+1)$ for ports.

For a hub $h$ with $2d\ge2$ neighbours $u_1\prec\dots\prec u_{2d}$, put $\pot(h):=(\mathrm{rk}(u_d)+\tfrac12)/(|U_\clu|+1)$. Since $\mathrm{rk}(u_d)\le|U_\clu|-1$, this lies in $(0,1)$. A hub without beads gets $\pot(h):=\tfrac12$.

Check the arcs.
\begin{itemize}
\item An arc $u_i\to h$ ($i\le d$) has $\mathrm{rk}(u_i)\le\mathrm{rk}(u_d)<\mathrm{rk}(u_d)+\tfrac12$.
\item An arc $h\to u_i$ ($i>d$) has $\mathrm{rk}(u_i)\ge\mathrm{rk}(u_{d+1})\ge\mathrm{rk}(u_d)+1$, because ranks are distinct integers.
\item Port--port arcs go up in rank.
\end{itemize}
So $\pot$ strictly increases along every arc, and $\MED(\prec)$ has no directed cycle. Every hub of degree $2d$ has exactly $d$ in-arcs and $d$ out-arcs. Hence $\MED(\prec)$ is admissible.

(b) Fix an admissible orientation. $\exc(u)$ is a sum of $\pm1$ over the $\deg_{B_\clu}(u)$ beads at $u$, which gives the bound $|\exc(u)|\le\deg_{B_\clu}(u)$ and the parity claim.

Every arc contributes $+1$ to $d^+-d^-$ at its tail and $-1$ at its head, so $\sum_{v\in V(\clu)}(d^+(v)-d^-(v))=0$. Hubs contribute $0$, hence $\sum_{u\in U_\clu}\exc(u)=0$.

Finally,
\begin{align*}
\Phi(\clu)=\sum_u\exc(u)^+&\le\sum_{u\in U_\clu}d^+(u)=\#\{\text{arcs port}\to\text{hub}\}+e(B_\clu[U_\clu])\\
&=\sum_{h\in A_\clu}d^-(h)+e(B_\clu[U_\clu])=b(\clu).
\end{align*}
Here we used that every arc with tail at a port ends at a hub or at a port, and that $d^-(h)=\deg_{B_\clu}(h)/2$ at every hub.
\end{proof}

\begin{lemma}[Lemma EQ-LPT]\label{s6:lemEQLPT}
Let $\Phi_1\ge\Phi_2\ge\dots\ge\Phi_m\ge0$ be loads of $m$ items, and let $1\le k\le m$. Place the items in the order $1,2,\dots,m$ into $k$ initially empty layers. Each item goes to a layer of currently minimum load, and to an empty layer whenever one exists (``greedy longest-processing-time placement''). Then all $k$ layers are non-empty, and the final layer loads satisfy $\max_j\Phi^{\mathrm{lay}}_j-\min_j\Phi^{\mathrm{lay}}_j\le\Phi_1$.
\deps{none}
\end{lemma}

\begin{proof}
An empty layer has load $0$, which is minimum. So the rule is consistent, and the first $k$ items go to the $k$ distinct empty layers. Since $k\le m$, every layer is non-empty.

Let $j^*$ be a layer of maximum final load and $x$ the last item placed in it. When $x$ was placed, the load of $j^*$ was at most the load of every layer at that moment. That is at most the final load of every layer, because loads only grow. Hence
\[
\max_j\Phi^{\mathrm{lay}}_j=(\text{load of }j^*\text{ before }x)+\Phi_x\le\min_j\Phi^{\mathrm{lay}}_j+\Phi_1.
\qedhere
\]
\end{proof}

\begin{lemma}[Lemma PAR, port--port parity]\label{s6:lemPAR}
Let $E_{ab}$ be a bipartite graph with sides $V_a,V_b$, and let $V(E_{ab})$ be the set of vertices incident with an edge of $E_{ab}$. Call a connected component of $E_{ab}$ \emph{odd} if it has an odd number of edges. In each odd component choose one edge that is either a non-bridge or a pendant edge (an edge with an end of degree $1$) of that component; such an edge exists. Let $J'$ be the set of chosen edges. Then, for \emph{every} such choice,
\[
|J'|\le\#\{\text{odd components}\}\le|V(E_{ab})|/2,
\]
and there is a partition $E_{ab}\setminus J'=S_a\sqcup S_b$ such that every vertex of $V_b$ has even $S_a$-degree and every vertex of $V_a$ has even $S_b$-degree.
\deps{\ref{s1:citEuler}}
\end{lemma}

\begin{proof}
\emph{Existence of the edge.} Let $C$ be a component with at least one edge. If $C$ contains a cycle, every edge of that cycle is a non-bridge. Otherwise $C$ is a tree with at least one edge; it has a leaf, and the edge at the leaf is pendant.

\emph{The count.} Every odd component has at least one edge, hence at least two vertices of $V(E_{ab})$. Components are vertex-disjoint, so there are at most $|V(E_{ab})|/2$ of them.

\emph{Even components after deletion.} Let $C$ be odd with chosen edge $e$.
\begin{itemize}
\item If $e$ is a non-bridge, $C-e$ is connected and has $|E(C)|-1$ edges, an even number.
\item If $e$ is pendant, $C-e$ consists of an isolated vertex and a connected graph with $|E(C)|-1$ edges.
\end{itemize}
Hence every component $C'$ of $E_{ab}\setminus J'$ that has an edge has an even number of edges.

\emph{The partition.} Fix such a component $C'$. Prescribe parities $p(v):=0$ for $v\in V_b\cap V(C')$ and $p(u):=\deg_{C'}(u)\bmod2$ for $u\in V_a\cap V(C')$. Every edge of $C'$ has exactly one end in $V_a$, so
\[
\sum_{w\in V(C')}p(w)\equiv\sum_{u\in V_a\cap V(C')}\deg_{C'}(u)=|E(C')|\equiv0\pmod2.
\]
So $T:=\{w:p(w)=1\}$ has even size. By Cited result~\ref{s1:citEuler}, a spanning tree of the connected graph $C'$ contains a $T$-join $S(C')$: an edge set with $\deg_{S(C')}(w)$ odd exactly for $w\in T$.

Put $S_a:=\bigcup_{C'}S(C')$ and $S_b:=(E_{ab}\setminus J')\setminus S_a$. Then:
\begin{itemize}
\item every $v\in V_b$ has even $S_a$-degree;
\item every $u\in V_a$ has $\deg_{S_b}(u)=\deg_{E_{ab}\setminus J'}(u)-\deg_{S_a}(u)\equiv\deg_{C'}(u)-p(u)\equiv0\pmod 2$, where $C'$ is the component of $u$.
\end{itemize}
Isolated vertices have degree $0$ in both.
\end{proof}

\begin{lemma}[Lemma HCC-P, layered hub-cluster chaining with path junctions]\label{s6:lemHCCP}
Let $k\ge1$, and suppose the following data are given.
\begin{itemize}
\item \emph{Layers} $\mathcal K_0,\dots,\mathcal K_{k-1}$: finite families of clusters, each cluster with a fixed admissible orientation. Put $\LayP_j:=\bigcup_{\clu\in\mathcal K_j}U_\clu$.
\item \emph{Junction sets} $T_0,\dots,T_{k-1}$: pairwise disjoint vertex sets.
\item A \emph{padding} $\pad:\bigcup_j\LayP_j\to\mathbb Z_{\ge0}$, with $\pad\equiv0$ if $k=1$.
\end{itemize}
Define the demands $\dem^-(u):=\exc(u)^-+\pad(u)$ and $\dem^+(u):=\exc(u)^++\pad(u)$, and the padded loads $\Phi_j:=\sum_{u\in\LayP_j}\dem^+(u)$. For $k=1$ put $X^{\mathrm{out}}:=\{u\in\LayP_0:\exc(u)<0\}$ and $X^{\mathrm{in}}:=\{u\in\LayP_0:\exc(u)>0\}$. Indices of layers and junction sets are taken modulo $k$. Assume:
\begin{itemize}
\item[(D)] the sets $A_\clu$ and $U_\clu$ (over all clusters of all layers) and $T_0,\dots,T_{k-1}$ are pairwise disjoint;
\item[(P)] every cluster is parity-clean;
\item[(JC-P)] for each $j$ there is a family $\mathcal P_j$ of pairwise edge-disjoint paths of $G$ with the following properties.
  \begin{itemize}
  \item Each path runs from a vertex of $\LayP_j$ to a vertex of $\LayP_{j+1}$; for $k=1$, from $X^{\mathrm{out}}$ to $X^{\mathrm{in}}$.
  \item All interior vertices of each path lie in $T_j$.
  \item Every $u\in\LayP_j$ is the first vertex of exactly $\dem^-(u)$ paths of $\mathcal P_j$, and every $v\in\LayP_{j+1}$ is the last vertex of exactly $\dem^+(v)$ paths of $\mathcal P_j$. For $k=1$ this reads: $u\in X^{\mathrm{out}}$ starts exactly $\exc(u)^-$ paths and $v\in X^{\mathrm{in}}$ ends exactly $\exc(v)^+$ paths.
  \end{itemize}
  The edge sets $E(\mathcal P_0),\dots,E(\mathcal P_{k-1})$ and all bead sets $B_\clu$ are pairwise disjoint.
\end{itemize}
Then:
\begin{enumerate}
\item[(i)] all $\Phi_j$ are equal to a common value $\Phi$, and $\Phi=|\mathcal P_j|$ for every $j$;
\item[(ii)] there are sets $F_j\subseteq E(\mathcal P_j)$ with $E(\mathcal P_j)\setminus F_j\subseteq E(G[T_j])$ such that $\bigcup_\clu B_\clu\cup\bigcup_jF_j$ decomposes into at most $\Phi$ cycles of $G$, each of length at least $\max(3,k)$.
\end{enumerate}
The edges of $E(\mathcal P_j)\setminus F_j$ are called \emph{returned}; they are not used.
\deps{\ref{s6:lemGATE}, \ref{s6:lemMED}, \ref{s6:defCluster}}
\end{lemma}

\begin{proof}
We treat $k\ge2$ and $k=1$ together. For $k=1$, read ``$\LayP_j$ as first vertices'' as $X^{\mathrm{out}}$ and ``$\LayP_{j+1}$ as last vertices'' as $X^{\mathrm{in}}$.

\emph{Step 0: the loads.} Each cluster has an admissible orientation, so $\sum_{u\in U_\clu}\exc(u)=0$ by Lemma~\ref{s6:lemMED}(b). Hence $\sum_{u\in\LayP_j}\exc(u)^-=\sum_{u\in\LayP_j}\exc(u)^+$ for every $j$. Counting first vertices and last vertices of paths gives
\[
|\mathcal P_j|=\sum_{u\in\LayP_j}\dem^-(u)=\sum_{u\in\LayP_j}(\exc(u)^++\pad(u))=\Phi_j,
\qquad
|\mathcal P_j|=\sum_{v\in\LayP_{j+1}}\dem^+(v)=\Phi_{j+1}.
\]
Hence $\Phi_j=\Phi_{j+1}$ for all $j$ modulo $k$, and all padded loads equal $\Phi:=|\mathcal P_j|$. For $k=1$ the same count gives $|\mathcal P_0|=\sum\exc^-=\sum\exc^+=\Phi_0$. This proves (i).

\emph{Step 1: orient and cancel.} Orient every path of $\mathcal P_j$ from its first to its last vertex, and let $D_j$ be the resulting set of arcs. The paths are edge-disjoint, so every edge of $E(\mathcal P_j)$ receives exactly one direction. We record three facts.
\begin{itemize}
\item No vertex of $\LayP_j$ is an interior vertex of a path of $\mathcal P_j$ (interiors lie in $T_j$, which avoids all ports by (D)). No vertex of $\LayP_j$ is a last vertex: last vertices lie in $\LayP_{j+1}$, which is disjoint from $\LayP_j$ by (D) when $k\ge2$ (the clusters of different layers are different clusters, and for $k=2$ the layers $j+1$ and $j-1$ coincide but differ from $j$); for $k=1$, $X^{\mathrm{out}}\cap X^{\mathrm{in}}=\emptyset$. Symmetrically, no vertex of $\LayP_{j+1}$ is a first or interior vertex.
\item Hence in $D_j$ every $u\in\LayP_j$ has out-degree $\dem^-(u)$ and in-degree $0$, and every $v\in\LayP_{j+1}$ has in-degree $\dem^+(v)$ and out-degree $0$.
\item Every $y\in T_j$ is never a first or last vertex (these are ports), and each path through $y$ contributes one in-arc and one out-arc at $y$. So $d^+_{D_j}(y)=d^-_{D_j}(y)$. All vertices of $D_j$ lie in $\LayP_j\cup T_j\cup\LayP_{j+1}$.
\end{itemize}
A directed cycle of $D_j$ cannot pass through a port, since ports are sources or sinks of $D_j$. So every directed cycle of $D_j$ has all its vertices in $T_j$.

Delete directed cycles from $D_j$ one at a time until no directed cycle remains. Call the result $D'_j$ and let $F_j$ be its underlying edge set. Each deletion lowers in- and out-degree by one at the vertices of a cycle inside $T_j$ and touches no arc at a port. So the three facts above hold for $D'_j$, $D'_j$ is acyclic, and $E(\mathcal P_j)\setminus F_j\subseteq E(G[T_j])$.

\emph{Step 2: the Eulerian digraph.} Let $\vec F$ consist of the arcs of the fixed admissible orientations of all clusters together with $D'_0,\dots,D'_{k-1}$. Its underlying edge set is $F:=\bigcup_\clu B_\clu\cup\bigcup_jF_j$.
\begin{itemize}
\item These are distinct edges of $G$. The clusters are vertex-disjoint by (D), so their bead sets are disjoint; the remaining disjointness is part of (JC-P).
\item Every hub is balanced. A hub is not a port and lies in no $T_j$, so it meets only arcs of its own cluster, which are balanced at hubs by admissibility.
\item Every $y\in T_j$ is balanced. By (D), $y$ lies in no cluster and in no $T_i$ with $i\neq j$. The arcs of $D'_i$ have both ends in $\LayP_i\cup T_i\cup\LayP_{i+1}$, so $y$ meets only arcs of $D'_j$, which are balanced at $y$.
\item Every port is balanced. Let $u$ be a port of a cluster of layer $j$, with $k\ge2$. By (D) and the fact just recalled, $u$ meets arcs of its own cluster, out-arcs of $D'_j$ (as a vertex of $\LayP_j$) and in-arcs of $D'_{j-1}$ (as a vertex of $\LayP_{(j-1)+1}$), and no others; $j-1\neq j$ modulo $k$. Hence
\[
d^+_{\vec F}(u)-d^-_{\vec F}(u)=\exc(u)+\dem^-(u)-\dem^+(u)=\exc(u)+\exc(u)^--\exc(u)^+=0.
\]
For $k=1$: a port with $\exc(u)<0$ has $\exc(u)^-=-\exc(u)$ out-arcs in $D'_0$; a port with $\exc(u)>0$ has $\exc(u)$ in-arcs; a port with $\exc(u)=0$ has none. In each case the port is balanced.
\end{itemize}

\emph{Step 3: a potential.} For each cluster $\clu$ fix a function $\pot_\clu:V(\clu)\to(0,1)$ that strictly increases along the arcs of its admissible orientation. One exists because the orientation is acyclic: take the position in a topological order divided by $|V(\clu)|+1$ (for $\MED$ one may use Lemma~\ref{s6:lemMED}(a)). For each $j$ fix a numbering $\mathrm{tp}_j:T_j\to\{1,\dots,|T_j|\}$ that increases along the arcs of the acyclic digraph $D'_j[T_j]$. Define
\[
\Psi(v):=3j+\pot_\clu(v)\ \ (v\in V(\clu),\ \clu\in\mathcal K_j),\qquad
\Psi(y):=3j+1+\frac{\mathrm{tp}_j(y)}{|T_j|+1}\ \ (y\in T_j).
\]
This is well defined by (D). Moreover $\Psi(v)\in(3j,3j+1)$ on layer $j$, and $\Psi(y)\in(3j+1,3j+2)$ on $T_j$. Now check the arcs.
\begin{itemize}
\item Bead arcs increase $\Psi$.
\item An arc of $D'_j$ from $u\in\LayP_j$ to $y\in T_j$ has $\Psi(u)<3j+1<\Psi(y)$. An arc inside $T_j$ increases $\mathrm{tp}_j$.
\item For $j\le k-2$, an arc of $D'_j$ with head $v\in\LayP_{j+1}$ has tail in $T_j\cup\LayP_j$, and $\Psi(\text{tail})<3j+2<3j+3<\Psi(v)$.
\end{itemize}
Let $\mathcal W$ be the set of arcs of $D'_{k-1}$ whose head lies in $\LayP_0$ (for $k=1$: the arcs of $D'_0$ whose head is a port). Every arc outside $\mathcal W$ increases $\Psi$, so $\vec F-\mathcal W$ is acyclic.

An arc of $D'_{k-1}$ whose head is a port is the last arc of its path, because interior vertices lie in $T_{k-1}$. Each path of $\mathcal P_{k-1}$ has exactly one last arc, which ends at a port and is never deleted by the cancellation of Step 1. Hence $|\mathcal W|=|\mathcal P_{k-1}|=\Phi$.

\emph{Step 4: count and length.} By Lemma~\ref{s6:lemGATE}, $F$ decomposes into at most $|\mathcal W|=\Phi$ cycles of $G$. Each of them is the orientation of a directed cycle $C$ of $\vec F$ of length at least $3$ containing an arc of $\mathcal W$.

For the bound $k$ on the length, call $[3j,3j+3)$ the $j$-th block. Every arc of $\vec F-\mathcal W$ either has both ends in the same block, or is an arc of some $D'_j$ ($j\le k-2$) from block $j$ to block $j+1$ with head in $\LayP_{j+1}$. Take a maximal subpath of $C$ that contains no arc of $\mathcal W$. It starts at the head of a $\mathcal W$-arc, a vertex of $\LayP_0$ (block $0$), and ends at the tail of a $\mathcal W$-arc, a vertex of $T_{k-1}\cup\LayP_{k-1}$ (block $k-1$). Along it $\Psi$ increases, so it contains at least $k-1$ block-crossing arcs. Together with the $\mathcal W$-arc that follows it, $C$ has at least $k$ arcs. This proves (ii).
\end{proof}

\begin{lemma}[union of systems]\label{s6:lemHCCglob}
Let $E_1,\dots,E_q\subseteq E(G)$ be pairwise disjoint, and suppose each $E_i$ decomposes into $a_i$ objects. Then $E_1\cup\dots\cup E_q$ decomposes into $\sum_ia_i$ objects.

In particular, suppose several systems satisfy the hypotheses of Lemma~\ref{s6:lemHCCP} separately, each with its own layers, junction sets and path families, and the beads and the path edges of distinct systems are pairwise disjoint. Then, with the sets $F_j$ given by Lemma~\ref{s6:lemHCCP} for each system, the union over all systems of the beads together with the $F_j$ decomposes into at most the sum of their values $\Phi$ of cycles of $G$. No vertex-disjointness across systems is needed: a vertex may be a hub, a port or a junction vertex in several systems. Hypothesis (D) is used only inside one system, through the potential $\Psi$ of that system.
\deps{\ref{s1:factAdd}, \ref{s6:lemHCCP}}
\end{lemma}

\begin{proof}
The union of the given decompositions is a family of objects whose edge sets partition $\bigcup_iE_i$, because the $E_i$ are pairwise disjoint. This is the argument of the proof of Fact~\ref{s1:factAdd}(b) (concatenation of decompositions of disjoint edge sets); the statement of Fact~\ref{s1:factAdd}(b) itself is the corresponding bound on $\f$.

For the second statement, apply Lemma~\ref{s6:lemHCCP} to each system separately. Its proof uses (D) only for the vertices of that system and of its own junction sets.
\end{proof}

\subsection{Designations and hub concentration: CONC and CONC-L}

From here on a valid $\HBtp$ run on $G$ is fixed (Definition~\ref{s2:defHBtp}), with the notation of Definition~\ref{s2:defAncestors}. For an ancestor $Y$ we write $r=r(Y)$ for its round. All logarithms are to base $2$.

\begin{definition}[designation and class data; $\mathcal V=\emptyset$]\label{s6:defDesign}
A \emph{designation} $\delta$ assigns, for every round $l\ge3$, to every standalone pre-part $Z\in\Std_l$ and every classed port $u\in Q_Z$ an ancestor $Y(u)\in\anc_l(u)$, the \emph{class} of $u$. Here $\delta$ is any deterministic function of the run. There are no VX-parts: every $Z\in\Std_l$ is treated as a MIX-part. Every classed port $u$ of round $l$ lies in $Q_Z$ for exactly one $Z\in\Std_l$ (port sets of one round are disjoint), and we write $Z_u$ for that part. Since $Y(u)\in\anc_l(u)$, we have $u\in V(Y(u))$ and $r(Y(u))\le l-2$.

For $l\ge3$, an ancestor $Y$ of round $r$ and any vertex $h$, define:
\begin{itemize}
\item the \emph{class-$Y$ degree} $d_{Y,l}(h):=\#\{hu\in E_l(Z):\ Z\in\Std_l,\ u\in Q_Z,\ Y(u)=Y\}$, and $m_{Y,l}:=\max_hd_{Y,l}(h)$, the maximum over all vertices $h$ of $G$;
\item $d^*_{Y,l}:=\#\{u:\ u\text{ a classed port of round }l,\ Y(u)=Y,\ u\in\Dups_{r}\}$;
\item for light $Y$, the \emph{guest core degree}
\[
\alpha_{Y,l}:=\max_{x\in S_Y}\#\{u\in Y\setminus\Dups_r:\ u\text{ a classed port of round }l,\ Y(u)=Y,\ xu\in E_l(Z_u)\},
\]
with $\alpha_{Y,l}:=0$ if $S_Y=\emptyset$;
\item the \emph{deterministic bead graph}
\[
\Bead_{Y,l}:=\{hu\in E_l(Z):\ Z\in\Std_l,\ u\in Q_Z,\ Y(u)=Y,\ \text{and }h\notin Q_Z\text{ or }Y(h)\neq Y\};
\]
\item $\gamma_l:=\lfloor P_{l-2}/M_l\rfloor$; the pair $(Y,l)$ is \emph{giant} iff some connected component of $\Bead_{Y,l}$ has more than $2\gamma_l$ edges;
\item the \emph{class counts} $\cagg_{h,l}:=\#\{Y:\ d_{Y,l}(h)\ge1\}$; for $Z\in\Std_l$ and a fresh port $x\in F_Z$, $c_x(Z):=\#\{Y(u):\ u\in Q_Z,\ xu\in E_l(Z)\}$; and for $u\in Q_Z$, $\cpp(u):=\#\{Y(v):\ v\in Q_Z,\ uv\in E_l(Z)\}$.
\end{itemize}
For $l\le2$ there are no classed ports ($\anc_l(x)=\emptyset$), and all these quantities vanish. All of them are deterministic functions of the run and $\delta$.
\deps{\ref{s2:defAncestors}, \ref{s2:defHBtp}}
\end{definition}

The following elementary facts about $\logs$ are used in the proofs of the next two theorems. Put $\mathsf T_0:=1$ and $\mathsf T_{i+1}:=2^{\mathsf T_i}$. For $x\ge1$ and an integer $k\ge0$, $\logs x\le k$ iff $x\le\mathsf T_k$. For $x>1$, $\logs x=1+\logs(\log x)$. And $\logs x\le1+\log x$ for all $x\ge1$: this holds for $x\le4$ by inspection, and for $x>4$ by induction, since $2+\log\log x\le1+\log x$ when $x\ge4$.

Put $k_*:=\logs\Dstar$. By \ref{s1:condG1}(a) at $\mu=\log\log\Dstar$, $\log\log\Dstar\ge2^8>16$, so $\Dstar>2^{65536}=\mathsf T_5$ and $k_*\ge6$.

For $d\ge\Dstar$ put
\[
\mathsf a(d):=\frac{2\logs d+2}{\log d},\qquad \mathsf b(d):=\frac{2\logs d+1}{\log\log d},\qquad \mathsf c(d):=\frac{2\logs d+1}{(\log d)^{1/2}}.
\]
We use two facts. First, for every valid run and every $r<R$,
\begin{equation}\label{s6:eqTowerHalf}
\mathsf a(d_r)\le\tfrac12\mathsf a(d_{r+1}),\qquad \mathsf b(d_r)\le\tfrac12\mathsf b(d_{r+1}),\qquad \mathsf c(d_r)\le\tfrac12\mathsf c(d_{r+1}).
\end{equation}
Second, for every $d\ge\Dstar$,
\begin{equation}\label{s6:eqTowerEnd}
\mathsf a(d)\le\frac{2k_*+4}{\log\Dstar},\qquad \mathsf b(d)\le\frac{2k_*+3}{\log\log\Dstar},\qquad \mathsf c(d)\le\frac{2k_*+3}{(\log\Dstar)^{1/2}}.
\end{equation}

\begin{proof}[Proof of \eqref{s6:eqTowerHalf} and \eqref{s6:eqTowerEnd}]
For \eqref{s6:eqTowerHalf}, put $x:=\lambda_r=\log d_r$ and $y:=\lambda_{r+1}$. Then $y\ge\log\Dstar$, so $v\defeq\log y\ge\log\log\Dstar$, and \ref{s1:condG1}(a),(b) at $v$ give $v\ge2^8$ and $y=2^v\ge2^{14}\Aexp v^3$. Moreover $x\ge d_{r+1}^{1/\Aexp}=2^{y/\Aexp}$ by Lemma~\ref{s2:lemTower}(a). Also $\logs d_r=1+\logs x\le2+\log x$, and $\logs d_r=2+\logs(\log x)\le3+\log\log x$.

The functions $t\mapsto(2\log t+6)/t$, $t\mapsto(7+2\log t)/t$ and $t\mapsto(2\log t+5)/t^{1/2}$ are decreasing for $t\ge4$. Hence:
\begin{itemize}
\item $\mathsf a(d_r)\le(2\log x+6)/x\le(2y/\Aexp+6)/2^{y/\Aexp}\le1/y\le\mathsf a(d_{r+1})/2$;
\item with $\mu:=\log x\ge y/\Aexp$: $\mathsf b(d_r)\le(7+2\log\mu)/\mu\le\Aexp(7+2\log y)/y\le1/(2\log y)\le\mathsf b(d_{r+1})/2$;
\item $\mathsf c(d_r)\le(2\log x+5)/x^{1/2}\le(2y/\Aexp+5)/2^{y/(2\Aexp)}\le1/(2y^{1/2})\le\mathsf c(d_{r+1})/2$.
\end{itemize}
The middle inequalities hold because $y/(2\Aexp)\ge2^{13}v^3\ge4+2v$ and $y\ge2^{14}\Aexp v^3\ge2\Aexp v(7+2v)$. The first gives $2^{y/(2\Aexp)}\ge2^{4+2v}=16y^2$, hence $y(2y/\Aexp+6)\le8y^2\le2^{y/\Aexp}$ and $2y^{1/2}(2y/\Aexp+5)\le14y^{3/2}\le2^{y/(2\Aexp)}$; the second is $\Aexp(7+2\log y)/y\le1/(2\log y)$.

For \eqref{s6:eqTowerEnd}, note $\Dstar\le\mathsf T_{k_*}$. There are two cases.
\begin{itemize}
\item If $\log d\le\mathsf T_{k_*}$, then $d\le\mathsf T_{k_*+1}$, so $\logs d\le k_*+1$, while $\log d\ge\log\Dstar$ and $\log\log d\ge\log\log\Dstar$. The three bounds follow.
\item If $\log d>\mathsf T_{k_*}\ge\Dstar$, put $x:=\log d>\Dstar$. Then $\logs d\le2+\log x$ and $\logs d\le3+\log\log x$. Hence $\mathsf a(d)\le(2\log x+6)/x\le(2\log\Dstar+6)/\Dstar\le1/\log\Dstar$; similarly $\mathsf b(d)\le(7+2\log\log x)/\log x\le(7+2\log\log\Dstar)/\log\Dstar\le1/\log\log\Dstar$, and $\mathsf c(d)\le(2\log\Dstar+5)/\Dstar^{1/2}\le(\log\Dstar)^{-1/2}$. The last inequality of each of the three chains holds because \ref{s1:condG1}(a),(b), applied at $\mu=\log\Dstar$ and at $\mu=\log\log\Dstar$ (both at least $\log\log\Dstar$), give $t\ge2^8$ and $2^t\ge2^{14}\Aexp t^3$ for $t\in\{\log\Dstar,\log\log\Dstar\}$; hence, as $t\ge1$, $t(2t+6)\le t(2t+7)\le9t^2\le2^t$ for both values of $t$, and $t(2t+5)^2\le49t^3\le2^t$ for $t=\log\Dstar$.\qedhere
\end{itemize}
\end{proof}

\begin{theorem}[Theorem CONC, standalone ancestors]\label{s6:thmCONC}
For every valid $\HBtp$ run and every designation $\delta$:
\begin{enumerate}
\item[(i)] for every $Y\in\Std_r$, every $l\ge r+1$ and every vertex $h$: $\#\{u\in Y^0\setminus\Dups_r:\ hu\in E(G_l)\}\le\tau_r-1$;
\item[(ii)] $m_{Y,l}\le\tau_r-1+d^*_{Y,l}$ for every $Y\in\Std_r$ and every $l$;
\item[(iii)] $\displaystyle\sum_{(Y,l):\ Y\text{ standalone}}m_{Y,l}\le2^{\sigma+14}\,n\,\frac{\logs\Dstar}{\log\Dstar}+100\,\eps\,n\,\frac{\logs\Dstar}{\Cp\log\log\Dstar}$;
\item[(iv)] (i)--(iii) hold for every designation, including designations that mix light and standalone classes: (i) bounds, for every vertex $h$, the number of $G_l$-edges from $h$ into $Y^0\setminus\Dups_r$, whatever the classes of the ports involved.
\end{enumerate}
\deps{\ref{s2:propOrigin}, \ref{s2:lemTower}, \ref{s2:propOV}, \ref{s2:lemLacunary}, \ref{s6:defDesign}, \ref{s2:defHBtp}, \ref{s1:condG1}}
\end{theorem}

\begin{proof}
(i) By \ref{s2:R1} and \ref{s2:R5} of Definition~\ref{s2:defHBtp}, $G_{l}\subseteq G_{r+1}$ for $l\ge r+1$. By Proposition~\ref{s2:propOrigin}(a), a standalone $Y$ has no edges of type $(\alpha)$. So every $G_{r+1}$-edge $hu$ with $u\in Y^0\setminus\Dups_r$ is of type $(\beta)$; in particular $h\notin Y^0$. By Proposition~\ref{s2:propOrigin}(b) there are at most $\tau_r-1$ of them in $G_l$.

(ii) Fix $h$. Every edge counted in $d_{Y,l}(h)$ is an edge $hu\in E_l(Z)\subseteq E(G_l)$ with $Y(u)=Y$, hence $u\in V(Y)=Y^0$. Those with $u\notin\Dups_r$ number at most $\tau_r-1$ by (i). Those with $u\in\Dups_r$ go to distinct ports $u$ ($G$ is simple) counted by $d^*_{Y,l}$. Taking the maximum over $h$ gives (ii).

(iii) \emph{The $\tau$-term.} $m_{Y,l}=0$ unless $3\le l$ and $r+2\le l\le R$, so at most $R-r-1$ values of $l$ occur for each $Y$ of round $r$. We use the looser bound $R-r-1<R-r\le2\logs d_r+2$ (Lemma~\ref{s2:lemTower}(a)), which matches the numerator of $\mathsf a$; the sharper count is not needed. By Proposition~\ref{s2:propOV}(K1), $|\Std_r|\le1.37n/P_r$. By Lemma~\ref{s2:lemTower}(c), $\tau_r/P_r\le2^{\sigma+11}/\lambda_r$. Hence
\[
\sum_{(Y,l),\ Y\in\Std}(\tau_r-1)\le\sum_{r\le R}1.37\,\frac{n}{P_r}(2\logs d_r+2)\tau_r\le1.37\cdot2^{\sigma+11}\,n\sum_{r\le R}\mathsf a(d_r).
\]
By \eqref{s6:eqTowerHalf} and Lemma~\ref{s2:lemLacunary}(i), $\sum_r\mathsf a(d_r)\le2\mathsf a(d_R)$. By \eqref{s6:eqTowerEnd} and $k_*\ge6$, $\mathsf a(d_R)\le(2k_*+4)/\log\Dstar\le\tfrac83k_*/\log\Dstar$. So the $\tau$-term is at most $1.37\cdot\tfrac{16}{3}\cdot2^{\sigma+11}k_*n/\log\Dstar\le2^{\sigma+14}n\logs\Dstar/\log\Dstar$, since $1.37\cdot16/3<8$.

\emph{The $d^*$-term.} Fix $l$. A classed port $u$ of round $l$ is counted in $d^*_{Y,l}$ only for $Y=Y(u)$, and then $u\in V(Y)\cap\Dups_{r(Y)}\subseteq Y^0\cap\Dups_{r(Y)}$. The map $u\mapsto(Y(u),u)$ is injective. Hence
\[
\sum_Yd^*_{Y,l}\le\sum_{r\le l-2}\ \sum_{Y^0\text{ a pre-part of round }r}|Y^0\cap\Dups_r|\le\sum_{r\le l-2}\frac{16\eps n}{\log P_r},
\]
by Proposition~\ref{s2:propOV}(K3). Here distinct ancestors of round $r$ have distinct pre-parts, and a standalone ancestor is its own pre-part.

Summing over $l$, each $r$ occurs for at most $R-r-1\le2\logs d_r+1$ values of $l$. Moreover $\log P_r\ge\Cp\log\lambda_r=\Cp\log\log d_r$. Hence
\[
\sum_{(Y,l)}d^*_{Y,l}\le\frac{16\eps n}{\Cp}\sum_{r\le R}\mathsf b(d_r)\le\frac{32\eps n}{\Cp}\mathsf b(d_R)\le\frac{32\eps n(2k_*+3)}{\Cp\log\log\Dstar}\le\frac{80\,\eps\,n\,\logs\Dstar}{\Cp\log\log\Dstar},
\]
by \eqref{s6:eqTowerHalf}, \eqref{s6:eqTowerEnd} and Lemma~\ref{s2:lemLacunary}(i), using $2k_*+3\le2.5k_*$ for $k_*\ge6$. Adding the two terms, via (ii), proves (iii); the constant $100$ in (iii) is deliberately looser than the $80$ proved here.

(iv) Nothing in (i)--(iii) used how $\delta$ chooses among the ancestors in $\anc_l(u)$. The bound (i) is a statement about $G_l$ and $Y^0$ alone. The bound (ii) holds for every $Y\in\Std_r$ whatever other ports have other classes. The sum in (iii) runs over the pairs $(Y,l)$ with $Y$ standalone, however the light classes are chosen.
\end{proof}

\begin{theorem}[Theorem CONC-L]\label{s6:thmCONCL}
Let $Y$ be a light part of round $r$ of a valid $\HBtp$ run, let $l\ge r+1$, and let $h$ be a vertex.
\begin{enumerate}
\item[(i)] If $h\notin S_Y$: $\#\{u\in Y\setminus\Dups_r:\ hu\in E(G_l)\}\le\tau_r-1$, and there are no such $u$ if $h\in Y$.
\item[(ii)] If $h\in S_Y$: every such $hu$ is an edge of $X^0_Y$, and their number is at most $e_{X^0_Y}(h,Y)$.
\item[(iii)] For every designation, $m_{Y,l}\le\max(\tau_r-1,\alpha_{Y,l})+d^*_{Y,l}$.
\item[(iv)] $\alpha_{Y,l}<\thGC_r(Y^0)$. For every valid run and every designation, summing over \emph{all} ancestors (light and standalone, GC-parts included),
\[
\sum_{(Y,l)}m_{Y,l}\le\epsCONC(\Dstar)\,n,
\]
where
\[
\epsCONC(\Dstar):=2^{\sigma+15}\frac{\logs\Dstar}{\log\Dstar}+200\,\eps\,\frac{\logs\Dstar}{\Cp\log\log\Dstar}+\frac{4(2\logs\Dstar+2)}{(\log\Dstar)^{1/2}}.
\]
$\epsCONC$ depends only on $\Dstar$ (and on $\sigma,\Cp,\eps$), and $\epsCONC(\Dstar)\to0$ as $\Dstar\to\infty$.
\end{enumerate}
\deps{\ref{s2:propOrigin}, \ref{s2:lemSEP}, \ref{s2:lemThinCut}, \ref{s2:lemGC}, \ref{s2:lemEL}, \ref{s2:lemTower}, \ref{s2:lemLacunary}, \ref{s6:thmCONC}, \ref{s2:propOV}, \ref{s2:defHBtp}, \ref{s6:defDesign}, \ref{s1:condG1}}
\end{theorem}

\begin{proof}
Recall $Y^0=Y\sqcup S_Y$ and $V(Y)=Y$. Let $u\in Y\setminus\Dups_r$ and let $hu\in E(G_l)\subseteq E(G_{r+1})$. By Proposition~\ref{s2:propOrigin}(a) (which rests on Lemmas~\ref{s2:lemSEP} and~\ref{s2:lemThinCut}), $hu$ is of exactly one of two types: $(\beta)$, with $h\notin Y^0$; or $(\alpha)$, with $h\in S_Y$ and $hu\in E(X^0_Y)$. The three cases $h\notin Y^0$, $h\in Y$ and $h\in S_Y$ are exhaustive.

(i) If $h\notin Y^0$, all such edges are of type $(\beta)$, and Proposition~\ref{s2:propOrigin}(b) bounds their number by $\tau_r-1$. If $h\in Y$, the edge $hu$ would be an edge of $G_l$ ($l>r$) with both ends in $V(Y)$, which Lemma~\ref{s2:lemEL} forbids.

(ii) If $h\in S_Y\subseteq Y^0$, the edge cannot be of type $(\beta)$, so it is of type $(\alpha)$ and lies in $E(X^0_Y)$. Distinct $u$ give distinct $X^0_Y$-edges from $h$ into $Y$, so their number is at most $e_{X^0_Y}(h,Y)$.

(iii) Fix $h$. A port $u$ with $Y(u)=Y$ lies in $V(Y)=Y$. The ports in $\Dups_r$ contribute at most $d^*_{Y,l}$, as in the proof of Theorem~\ref{s6:thmCONC}(ii). The others contribute at most $\tau_r-1$ if $h\notin Y^0$, nothing if $h\in Y$, and at most $\alpha_{Y,l}$ if $h\in S_Y$ (by definition of $\alpha_{Y,l}$, since the counted edges $hu$ lie in $E_l(Z_u)$). Exactly one case applies to $h$.

(iv) The edges counted in $\alpha_{Y,l}$ at $x\in S_Y$ are distinct $X^0_Y$-edges from $x$ into $Y=Y^0\setminus S_Y$, by (ii). By rule (GC) and Lemma~\ref{s2:lemGC}(ii), $e_{X^0_Y}(x,Y)<\thGC_r(Y^0)$. Hence $\alpha_{Y,l}<\thGC_r(Y^0)$, and by (iii),
\[
m_{Y,l}\le(\tau_r-1)+(\thGC_r(Y^0)-1)+d^*_{Y,l}\quad(Y\text{ light}),\qquad m_{Y,l}\le\tau_r-1+d^*_{Y,l}\quad(Y\text{ standalone}),
\]
the latter by Theorem~\ref{s6:thmCONC}(ii); GC-parts are standalone (Lemma~\ref{s2:lemGC}(iii)). We sum the three kinds of terms over all pairs $(Y,l)$ with $r(Y)+2\le l\le R$ (for other $l$, $m_{Y,l}=0$).

\emph{$\tau$-terms.} The number of ancestors of round $r$ (light and standalone) is at most $1.37n/P_r$ (Proposition~\ref{s2:propOV}(K1)). The computation in the proof of Theorem~\ref{s6:thmCONC}(iii) applies verbatim and gives at most $2^{\sigma+14}n\logs\Dstar/\log\Dstar$.

\emph{$d^*$-terms.} The injectivity argument in the proof of Theorem~\ref{s6:thmCONC}(iii) used only $u\in V(Y(u))\subseteq Y(u)^0$ and Proposition~\ref{s2:propOV}(K3), which sums over all pre-parts of round $r$. So it applies to all ancestors together and gives at most $80\eps n\logs\Dstar/(\Cp\log\log\Dstar)$.

\emph{$\thGC$-terms.} For light $Y$ of round $r$, $\thGC_r(Y^0)\le|Y^0|\lambda_r^{-1/2}+1$. The pre-parts $Y^0$ of distinct light parts of round $r$ are distinct, so by Proposition~\ref{s2:propOV}(K1)
\[
\sum_{Y\text{ light, round }r}\thGC_r(Y^0)\le1.37\,n\,\lambda_r^{-1/2}+1.37\,\frac n{P_r}\le1.38\,n\,\lambda_r^{-1/2},
\]
using $P_r\ge\lambda_r^{\Cp}$. Each $Y$ occurs for at most $R-r-1\le2\logs d_r+1$ values of $l$ (Lemma~\ref{s2:lemTower}(a)). By \eqref{s6:eqTowerHalf}, \eqref{s6:eqTowerEnd} and Lemma~\ref{s2:lemLacunary}(i),
\[
\sum_{(Y,l),\ Y\text{ light}}\thGC_r(Y^0)\le1.38\,n\sum_{r\le R}\mathsf c(d_r)\le2.76\,n\,\mathsf c(d_R)\le\frac{2.76\,n(2k_*+3)}{(\log\Dstar)^{1/2}}\le\frac{4n(2\logs\Dstar+2)}{(\log\Dstar)^{1/2}}.
\]
The last step holds because $2.76(2k+3)\le4(2k+2)$ for all $k\ge1$.

Adding the three sums gives $\sum_{(Y,l)}m_{Y,l}\le\epsCONC(\Dstar)n$. In fact the first two terms of $\epsCONC$ carry a spare factor $2$.

Each term of $\epsCONC$ tends to $0$ as $\Dstar\to\infty$, because $\logs x$ grows more slowly than every iterate of $\log$. The bound uses nothing about $\delta$ beyond $u\in V(Y(u))$, so it holds for every designation.
\end{proof}

\subsection{Lending statuses and lost ports}

\begin{definition}[stage-2 lending statuses, retirement sets and the functional $\XU$]\label{s6:defLending}
Fix a valid run, a designation $\delta$ and a stage-1 outcome: the colourings and JS labels of Definition~\ref{s3:defCOL}, and the zones of Definition~\ref{s5:defZones}.
\begin{itemize}
\item An ancestor $Y$ is \emph{lend-bad} iff $Y$ is a demoted light part (Definition~\ref{s5:defStages}), or an event of Lemma~\ref{s3:lemCOL}(a) fails for $Y$, or an event of Lemma~\ref{s3:lemCOL}(b) fails for $Y$. Otherwise $Y$ is \emph{lend-good}. Thus for a lend-good $Y$ the classes $\Own_Y$, $\Lend_Y$ and all lent classes are the expanders of Lemma~\ref{s3:lemCOL}(a), and every $\LJS_{Y,l,j}$ is $(2^{12}L_Y^4,\tJS_l)$-path connected through $T_j(Y,l)$.
\item For $l\ge3$ and $Z\in\Std_l$, the set of \emph{lost} ports is
\[
\Lost_Z:=\{u\in Q_Z:\ Y(u)\text{ is lend-bad, or }\lab_{Y(u),l}(u)\neq*\}.
\]
Put $\Ret_Z:=A_Z\cup F_Z\cup\Lost_Z$ (the \emph{retirement set} of $Z$ for Lemma~\ref{s4:lemTPV}) and $\Qs_Z:=Q_Z\setminus\Lost_Z$. For $l\le2$ put $\Lost_Z:=\emptyset$, $\Ret_Z:=A_Z\cup F_Z=Z^0$ and $\Qs_Z:=\emptyset$.

Since $A_Z=Z^0\cap D_l$ and $U_Z=Z^0\setminus D_l=F_Z\sqcup Q_Z$, we have $Z^0=\Ret_Z\sqcup\Qs_Z$, and $A_Z,F_Z,\Lost_Z$ are pairwise disjoint.
\item $O_Z:=\Own_Z$ if the ancestor $Z$ is lend-good, and $O_Z:=X^0_Z$ otherwise.
\item $\Lost_l:=\bigcup_{Z\in\Std_l}\Lost_Z$.
\item $\XU:=80\,\mathrm{dem}+369\,\mathrm{lp}+\sum_l(169+M_l)|\Lost_l|$, with $\mathrm{dem}$ and $\mathrm{lp}$ as in Definition~\ref{s5:defStages}.
\end{itemize}
All of these are deterministic functions of the run, $\delta$ and the stage-1 outcome. For every classed port $u\in\Qs_Z$, the class $Y(u)$ is lend-good and $u\notin\bigcup_jT_j(Y(u),l)$.
\deps{\ref{s3:defCOL}, \ref{s3:lemCOL}, \ref{s5:defStages}, \ref{s6:defDesign}, \ref{s5:defZones}}
\end{definition}

\begin{lemma}[lost ports; the constant (K6)]\label{s6:lemLost}
Fix a valid run and a designation $\delta$. Probabilities and expectations are over the stage-1 outcome.
\begin{enumerate}
\item[(i)] for every classed port $u$ of round $l\ge3$,
\[
\Prob(u\in\Lost)\le\KJS_l\rho_l+\Prob(Y(u)\text{ lend-bad})\le M_l^{-2}+|V(Y(u))|^{-2}\le2M_l^{-2};
\]
\item[\textup{(K6)}]\namedlabel{s6:K6}{\textup{(K6)}} $\Ex|\Lost_l|\le2n/M_l^2$ for every round $l$;
\item[(iii)] $\Ex\XU\le\epsU(\Dstar)\,n+\sum_l2n(169+M_l)/M_l^2\le\epsU(\Dstar)\,n+5.5\,n/\Dstar$.
\end{enumerate}
\deps{\ref{s6:defLending}, \ref{s3:lemCOL}, \ref{s5:lemE1}, \ref{s5:lemExpect}, \ref{s2:propStructure}, \ref{s2:lemTower}, \ref{s3:defCOL}}
\end{lemma}

\begin{proof}
(i) Let $Y:=Y(u)$, of round $r\le l-2$. The designation is deterministic, so $Y$ is fixed.

\emph{The label.} By Definition~\ref{s3:defCOL}(iv), $\lab_{Y,l}(u)\neq*$ has probability $\KJS_l\rho_l=M_l^2\cdot M_l^{-4}=M_l^{-2}$.

\emph{Lend-bad.} By the union bound, $\Prob(Y\text{ lend-bad})$ is at most the probability that an event of Lemma~\ref{s3:lemCOL}(a) or~(b) fails, plus, for light $Y$, the probability that $Y$ is demoted. The first is at most $|V(Y)|^{-2}/2$ by Lemma~\ref{s3:lemCOL}. The second is at most $|Y|^{-2}/2=|V(Y)|^{-2}/2$ by Lemma~\ref{s5:lemE1}(b), since $V(Y)=Y$ for a light part (Definition~\ref{s2:defAncestors}). Hence $\Prob(Y\text{ lend-bad})\le|V(Y)|^{-2}$.

\emph{Size of $V(Y)$.} By Proposition~\ref{s2:propStructure}(iv), $|V(Y)|\ge P_r/2$. By Lemma~\ref{s2:lemTower}(b), $P_r\ge P_{l-2}$ (as $P_r\ge2P_{r+1}$) and $P_{l-2}\ge M_l^{13}$. Hence $|V(Y)|^{-2}\le4M_l^{-26}\le M_l^{-2}$.

(K6) The port sets $U_Z$ of distinct $Z\in\Std_l$ are pairwise disjoint (Proposition~\ref{s2:propStructure}(iv)), so round $l$ has at most $n$ classed ports. By (i) and linearity of expectation, $\Ex|\Lost_l|\le2n/M_l^2$.

(iii) By Lemma~\ref{s5:lemExpect}, $\Ex[80\,\mathrm{dem}+369\,\mathrm{lp}]\le\epsU(\Dstar)n$. By (K6),
\[
\Ex\sum_l(169+M_l)|\Lost_l|\le\sum_l\Big(\frac{338n}{M_l^2}+\frac{2n}{M_l}\Big)\le\Big(2+\frac{338}{2^{40}}\Big)n\sum_l\frac1{M_l}\le\Big(2+\frac{338}{2^{40}}\Big)\frac{2n}{\Dstar}\le\frac{5.5\,n}{\Dstar},
\]
using $M_l\ge2^{40}$ and $\sum_l1/M_l\le2/\Dstar$ (Lemma~\ref{s2:lemTower}(b)).
\end{proof}

\subsection{Junction synchronization: Lemma JS-LC and the J-interface}

\begin{lemma}[Lemma JS-LC; one round $l\ge3$; cherry split for giant pairs]\label{s6:lemJSLC}
Fix a valid run, a designation $\delta$, a stage-1 outcome with its stage-2 data (Definition~\ref{s6:defLending}), and a round $3\le l\le R$. For every $Z\in\Std_l$ let $B_Z\subseteq E_l(Z)$ be a set of edges, each of which has an end in $\Qs_Z$. (In the assembly $B_Z:=\EQ(Z)$, which has this property by \ref{s4:T1} of Lemma~\ref{s4:lemTPV}.)

Hypothesis: no edge of any class $\LJS_{Y,l,j}$ ($0\le j<\KJS_l$, $Y$ a lend-good ancestor of a round $\le l-2$) has been used.

Then there are a set $J_l\subseteq\bigcup_ZB_Z$ and a set $\LentJS_l\subseteq\bigcup_{Y,j}\LJS_{Y,l,j}$ such that $\bigcup_ZB_Z\cup\LentJS_l$ is partitioned into the single edges of $J_l$ and at most
\[
126\,\frac n{M_l}+1.5\sum_{(Y,l)\text{ giant}}m_{Y,l}\quad\text{cycles}.
\]
Here the sum is over the ancestors $Y$ for which the pair $(Y,l)$ is giant (Definition~\ref{s6:defDesign}), and $m_{Y,l}$ is the maximal class-$Y$ degree of Definition~\ref{s6:defDesign}. The proof gives the constant $15$ in place of $126$; the looser $126$ is kept because it is the value used downstream (see the end of Step~8). Each cycle consists of edges of $\bigcup_ZB_Z$ and edges of $\bigcup_j\LJS_{Y,l,j}$ for a single ancestor $Y$. Moreover
\begin{equation}\label{s6:eqJbound}
|J_l|\le\sum_{h\in D_l}\cagg_{h,l}+\sum_{Z\in\Std_l}\Big[\sum_{x\in F_Z}c_x(Z)+(M_l-1)|\Lost_Z|+\frac12\sum_{u\in Q_Z}\cpp(u)\Big].
\end{equation}
The right side of \eqref{s6:eqJbound} is computed from the sets $E_l(Z)$, not from the $B_Z$. The bound \eqref{s6:eqJbound} is recorded for completeness and is not used later: Section~\ref{s7:sec} bounds $J_l$ through \Jp{1} and \Jp{2} of Lemma~\ref{s6:lemJplus}. The edges of $\bigcup_{Y,j}\LJS_{Y,l,j}$ outside $\LentJS_l$, among them the returned edges inside the graphs $G[T_j(Y,l)]$, are not used; they are junk for their owner $Y$.
\deps{\ref{s6:lemPAR}, \ref{s6:lemMED}, \ref{s6:lemEQLPT}, \ref{s6:lemHCCP}, \ref{s6:lemHCCglob}, \ref{s6:defDesign}, \ref{s6:defLending}, \ref{s3:lemCOL}, \ref{s3:lemMonotone}, \ref{s2:lemEL}, \ref{s2:lemCap}, \ref{s2:propOV}, \ref{s2:lemTower}, \ref{s2:defHBtp}, \ref{s2:propStructure}, \ref{s3:defCOL}, \ref{s1:citDef7}}
\end{lemma}

The hypothesis refers only to the classes indexed by the current round $l$. Classes $\LJS_{Y,l',j}$ with $l'>l$ of the same ancestors have already been used at earlier steps of the order $R,R-1,\dots$, and that is allowed.

\begin{proof}
Write $\KJS:=\KJS_l=M_l^2$ and $t:=\tJS_l=2M_l+2$. Recall two facts. For $u\in\Qs_Z$, the class $Y(u)$ is lend-good and $\lab_{Y(u),l}(u)=*$ (Definition~\ref{s6:defLending}). Every vertex has degree at most $M_l-1$ in every $E_l(Z)$ (Lemma~\ref{s2:lemCap}(ii)).

\emph{Step 1: types.} Let $e\in B_Z$. By hypothesis $e$ has an end $u\in\Qs_Z$. Both ends lie in $Z^0$, since $E_l(Z)\subseteq E(G[Z^0])$ by \ref{s2:R5}. Since $Z^0=\Ret_Z\sqcup\Qs_Z$, there are two cases for the other end $v$:
\begin{itemize}
\item $v\in\Ret_Z=A_Z\sqcup F_Z\sqcup\Lost_Z$. Then $e$ is a \emph{centre--port edge} with centre $v$.
\item $v\in\Qs_Z$. Then $e$ is a \emph{port--port edge}, and $Y(v)\neq Y(u)$: otherwise both ends of $e\in E(G_l)$ would lie in $V(Y(u))$ with $l>r(Y(u))$, contrary to Lemma~\ref{s2:lemEL}.
\end{itemize}
By the same lemma, for every edge $hu\in E_l(Z)$ with $u\in Q_Z$ and $Y(u)=Y$, the vertex $h$ lies outside $V(Y)$. In particular $h\notin T_j(Y,l)$, since $T_j(Y,l)\subseteq V(Y)$.

\emph{Step 2: parity; the set $J_l$.}

(2a) For each $Z\in\Std_l$ and each unordered pair $\{a,b\}$ of distinct classes, let $E_{ab}(Z)$ be the set of port--port edges of $B_Z$ joining a class-$a$ port of $\Qs_Z$ to a class-$b$ port of $\Qs_Z$. It is bipartite with sides $V_a$ (class-$a$ ports) and $V_b$ (class-$b$ ports). Apply Lemma~\ref{s6:lemPAR} with any admissible choice of the deleted edges. This gives $J'_{ab}(Z)$ and $E_{ab}(Z)\setminus J'_{ab}(Z)=S_a\sqcup S_b$. The edges of $S_a$ are \emph{assigned to class $a$}: in them the class-$b$ end is a \emph{pseudo-hub}, of even $S_a$-degree. Symmetrically, $S_b$ is assigned to class $b$.

(2b) For each class $Y$ let $R^0_Y$ consist of all centre--port edges $hu\in B_Z$ ($Z\in\Std_l$) with $Y(u)=Y$, together with all port--port edges assigned to $Y$ in (2a). Every edge of $R^0_Y$ has exactly one end that is a class-$Y$ port of some $\Qs_Z$, its \emph{port}; its other end is its \emph{centre}. A centre lies in $\Ret_Z$ or is a port of $\Qs_Z$ of a class $\neq Y$. By Step~1, centres of class-$Y$ edges lie outside $V(Y)$, while ports lie in $V(Y)$; so no vertex is both a port and a centre of class $Y$.

For each $Y$ and each vertex $h$ that is the centre of an odd number of edges of $R^0_Y$, move one such edge (any one) into $J_l$. Let $R_Y$ be the rest: the \emph{realized class-$Y$ beads}. Finally, $J_l$ consists of all edges $J'_{ab}(Z)$ and all moved edges.

\emph{Only centres in $\Ret_Z$ are ever odd.} A pseudo-hub $v\in\Qs_Z$ of class $b$ lies in exactly one part (it is not in $D_l$). It is a centre of class-$Y$ edges only through the edges of $E_{Yb}(Z)$ assigned to $Y$, and it has even degree in those by Lemma~\ref{s6:lemPAR}. A moved edge $hu$ changes the class-$Y$ centre degree of $h$ only: its other end $u$ is a port, which carries no parity constraint. So after the moves every centre has even degree in every $R_Y$.

\emph{Counting $J_l$.}
\begin{itemize}
\item PAR deletions: $\sum_{\{a,b\}}|J'_{ab}(Z)|\le\frac12\sum_{\{a,b\}}|V(E_{ab}(Z))|$. A port $u\in\Qs_Z$ lies in $V(E_{ab}(Z))$ only for $\{a,b\}=\{Y(u),Y(v)\}$ with $uv\in B_Z\subseteq E_l(Z)$ and $v\in Q_Z$. This happens for at most $\cpp(u)$ pairs, so the PAR deletions number at most $\frac12\sum_{u\in Q_Z}\cpp(u)$.
\item Moves at a hub $h\in D_l$: at most one per class $Y$ for which $h$ is the centre of an edge $hu\in E_l(Z)$ with $u\in Q_Z$ and $Y(u)=Y$, i.e.\ with $d_{Y,l}(h)\ge1$. So at most $\cagg_{h,l}$ moves, summed over all parts containing $h$.
\item Moves at a fresh centre $x\in F_Z$: $x\notin D_l$ lies in one part only, so at most $c_x(Z)$ moves.
\item Moves at a lost centre $v\in\Lost_Z$: at most the number of edges of $E_l(Z)$ at $v$, which is at most $M_l-1$.
\end{itemize}
This proves \eqref{s6:eqJbound}. Every quantity on its right is defined from the sets $E_l(Z)$.

\emph{Step 3: realized components.} Fix an ancestor $Y$ with $R_Y\neq\emptyset$. Then $Y$ is the class of a port in some $\Qs_Z$, so $Y$ is lend-good and $r(Y)\le l-2$.

Every edge $hu\in R_Y$ lies in $\Bead_{Y,l}$: $hu\in E_l(Z)$, $u\in\Qs_Z\subseteq Q_Z$, $Y(u)=Y$, and $h$ is not a class-$Y$ port of $Q_Z$ (a class-$Y$ vertex of $Q_Z$ lies in $V(Y)$, and $h\notin V(Y)$). Hence every connected component of the graph $R_Y$ lies inside a component of $\Bead_{Y,l}$ and has at most as many edges. Call a component of $R_Y$ \emph{giant} if it has more than $2\gamma_l$ edges. Giant components exist only if $(Y,l)$ is giant; whether $(Y,l)$ is giant is decided by the deterministic graph $\Bead_{Y,l}$.

Every edge of $R_Y$ joins a port (a class-$Y$ port $u$ of some $\Qs_Z$, with $u\in V(Y)\setminus\bigcup_jT_j(Y,l)$) to a centre (a vertex outside $V(Y)$). The degree of a centre $h$ in $R_Y$ is at most $d_{Y,l}(h)\le m_{Y,l}$. The degree of a port $u$ in $R_Y$ is at most $\deg_{E_l(Z_u)}(u)\le M_l-1$. All edges of $R_Y$ at a port lie in the component of that port.

\emph{Step 4: non-giant components; the system $\mathcal S_0(Y,l)$.} For every non-giant component $C$ of $R_Y$, let $\clu_C:=(\text{centres of }C,\ \text{ports of }C,\ E(C))$.
\begin{itemize}
\item This is a cluster: centres and ports are disjoint, and every bead joins a port to a centre. It is parity-clean by Step~2.
\item Fix any linear order of its ports and orient $\clu_C$ by $\MED$. By Lemma~\ref{s6:lemMED}, this is admissible and $\Phi(\clu_C)\le b(\clu_C)=|E(C)|/2\le\gamma_l$; here $b(\clu_C)=|E(C)|/2$ because $\clu_C$ has no port--port beads. Also $|\exc(u)|\le M_l-1$.
\item $\Phi(\clu_C)\ge1$: an acyclic orientation with an arc has a source, which must be a port of positive excess.
\end{itemize}
If $R_Y$ has no non-giant component, $\mathcal S_0(Y,l)$ is empty and contributes no cycles, with $\Phi(\mathcal S_0(Y,l))\defeq0$ in this case; the rest of Step~4 assumes that there is at least one cluster.
Let $\Sigma\Phi$ and $\Phi_{\max}$ be the sum and the maximum of these loads, and put
\[
k_0:=\min\bigl(\KJS,\max(1,\lfloor\Sigma\Phi/(2\Phi_{\max})\rfloor)\bigr).
\]
Then $k_0$ is at most the number of clusters. List the clusters in non-increasing order of $\Phi(\clu_C)$, as Lemma~\ref{s6:lemEQLPT} requires, and layer them by that lemma, with loads $\Phi(\clu_C)$, into $k_0$ non-empty layers with raw loads $\Phi^{\mathrm{raw}}_j$, where $\max_j\Phi^{\mathrm{raw}}_j-\min_j\Phi^{\mathrm{raw}}_j\le\Phi_{\max}$.
\begin{itemize}
\item If $k_0=1$: no padding, and the system has depth $1$ with load $\Sigma\Phi<4\Phi_{\max}$.
\item If $k_0\ge2$: put $\Phi^*:=\max_j\Phi^{\mathrm{raw}}_j$ and $\Pi_j:=\Phi^*-\Phi^{\mathrm{raw}}_j\le\Phi_{\max}$. Since $k_0\le\Sigma\Phi/(2\Phi_{\max})$, every raw load satisfies $\Phi^{\mathrm{raw}}_j\ge\Sigma\Phi/k_0-\Phi_{\max}\ge\Phi_{\max}\ge\Pi_j$. By Lemma~\ref{s6:lemMED}(b), $\Phi^{\mathrm{raw}}_j=\frac12\sum_{u\in\LayP_j}|\exc(u)|\le\frac12(M_l-1)|\LayP_j|$. So $\Pi_j$ units of padding can be spread over the ports of layer $j$ with $\pad(u)\le\lceil(M_l-1)/2\rceil\le\lceil M_l/2\rceil$. All padded loads then equal $\Phi^*\le\Sigma\Phi/k_0+\Phi_{\max}$.
\end{itemize}
In all cases the common padded load of $\mathcal S_0(Y,l)$ satisfies
\begin{equation}\label{s6:eqSzeroCost}
\Phi(\mathcal S_0(Y,l))\le1.5\,\Sigma\Phi/\KJS+5\gamma_l.
\end{equation}
Indeed: if $k_0=\KJS$ then $\Phi_{\max}\le\Sigma\Phi/(2\KJS)$; if $2\le k_0<\KJS$ then $k_0\ge\Sigma\Phi/(2\Phi_{\max})-1\ge\Sigma\Phi/(4\Phi_{\max})$, so $\Phi^*\le5\Phi_{\max}$; if $k_0=1$ the load is $<4\Phi_{\max}$; and $\Phi_{\max}\le\gamma_l$.

A port $u$ of layer $j$ has demand $\dem^-(u)$ at junction $j$ and $\dem^+(u)$ at junction $j-1$. Both are at most $|\exc(u)|+\pad(u)\le(M_l-1)+\lceil M_l/2\rceil$.

\emph{Step 5: giant components; the cherry split.} At every centre $h$ of a giant component of $R_Y$, pair the (evenly many) edges of $R_Y$ at $h$ arbitrarily into \emph{cherries} $\{hu,hu'\}$, written $u$--$h$--$u'$. Since $G$ is simple, $u\neq u'$. Each cherry is a parity-clean cluster $(\{h\},\{u,u'\},\{hu,hu'\})$. Its orientation $u\to h\to u'$ is admissible, with $\exc(u)=+1$, $\exc(u')=-1$ and load $1$. Every edge of a giant component lies in exactly one cherry.

Two cherries \emph{conflict} if they share a vertex. A centre is never an end of a cherry, because centres lie outside $V(Y)$ and ends inside. A cherry at $h$ conflicts with at most $\deg_{R_Y}(h)/2-1\le m_{Y,l}/2-1$ other cherries at $h$. Each of its two ends $u$ has at most $M_l-1$ edges in $R_Y$, each in at most one cherry, so it is an end of at most $M_l-2$ other cherries. So the conflict degree is at most $m_{Y,l}/2+2M_l-5$. Greedy colouring uses
\[
g_{Y,l}\le m_{Y,l}/2+2M_l-4
\]
colours, and each colour class $\mathcal S_i(Y,l)$ ($1\le i\le g_{Y,l}$) is a set of pairwise vertex-disjoint cherries. Let $\eta_i$ be the number of cherries in $\mathcal S_i(Y,l)$.

Each class $\mathcal S_i(Y,l)$ is made into a layered system of depth
\[
k_i:=\min\bigl(\KJS,\max(1,\lfloor\eta_i/2\rfloor)\bigr).
\]
The $\max(1,\cdot)$ is needed: $\eta_i\le3$ gives $k_i=1$. Since $k_i\le\eta_i$, Lemma~\ref{s6:lemEQLPT} with unit loads (trivially in non-increasing order) gives $k_i$ non-empty layers whose loads differ by at most $1$, i.e.\ equal $\lfloor\eta_i/k_i\rfloor$ or $\lceil\eta_i/k_i\rceil$. There are three cases.
\begin{itemize}
\item $k_i=1$ ($\eta_i\le3$): no padding, $X^{\mathrm{out}}$ the ports of excess $-1$, $X^{\mathrm{in}}$ those of excess $+1$, and load $\Phi(\mathcal S_i)=\eta_i\le3$.
\item $2\le k_i=\lfloor\eta_i/2\rfloor<\KJS$ ($\eta_i\ge4$): every layer of load $\lfloor\eta_i/k_i\rfloor<\lceil\eta_i/k_i\rceil$ gets padding $1$ on one of its ports. The load is $\Phi(\mathcal S_i)=\lceil\eta_i/\lfloor\eta_i/2\rfloor\rceil\le3$, since $\eta_i/\lfloor\eta_i/2\rfloor\le2\eta_i/(\eta_i-1)\le8/3$.
\item $k_i=\KJS$: padding as in the previous case, and $\Phi(\mathcal S_i)=\lceil\eta_i/\KJS\rceil\le\eta_i/\KJS+1$.
\end{itemize}
So $\pad\le1$, $\Phi(\mathcal S_i)\le\eta_i/\KJS+3$, and every port of a cherry system has demand at most $2$ at each junction. Put $\sco_{Y,l}:=(3g_{Y,l}-6M_l)^+$ (a symbol local to this proof), and $g_{Y,l}:=\sco_{Y,l}:=0$ if $R_Y$ has no giant component. Summing over the classes of $(Y,l)$,
\begin{equation}\label{s6:eqCherryCost}
\sum_i\Phi(\mathcal S_i(Y,l))\le\frac{\#\{\text{cherries of }(Y,l)\}}{\KJS}+3g_{Y,l}\le\frac{\#\{\text{cherries of }(Y,l)\}}{\KJS}+6M_l+\sco_{Y,l},
\end{equation}
and $\sco_{Y,l}=(3g_{Y,l}-6M_l)^+\le(1.5m_{Y,l}-12)^+\le1.5m_{Y,l}$. If $R_Y$ has a giant component, then the component of $\Bead_{Y,l}$ containing it has more than $2\gamma_l$ edges (Step~3), so $(Y,l)$ is giant; hence $\sco_{Y,l}=0$ unless $(Y,l)$ is giant.

\emph{Step 6: joint routing.} The \emph{systems of $(Y,l)$} are $\mathcal S_0(Y,l)$ (if non-empty) and the cherry classes $\mathcal S_i(Y,l)$. For a system $\mathcal S$ of depth $k_{\mathcal S}$ and a junction $j<k_{\mathcal S}$, its \emph{junction-$j$ pairs} are the pairs of a fixed bijection between two multisets:
\begin{itemize}
\item if $k_{\mathcal S}\ge2$: the out-units of layer $j$ ($\dem^-(u)$ copies of each $u\in\LayP_j$) and the in-units of layer $j+1$ modulo $k_{\mathcal S}$ ($\dem^+(v)$ copies of each $v\in\LayP_{j+1}$);
\item if $k_{\mathcal S}=1$ and $j=0$: $\exc(u)^-$ copies of each $u\in X^{\mathrm{out}}$ and $\exc(v)^+$ copies of each $v\in X^{\mathrm{in}}$.
\end{itemize}
Let $\mathfrak P_j(Y,l)$ be the multiset union of the junction-$j$ pairs of all systems of $(Y,l)$ of depth $>j$.

\begin{claim}[the joint-routing facts]
For every lend-good $Y$ with $R_Y\neq\emptyset$ and every $j<\KJS$:
\begin{enumerate}
\item[(a)] The bijections exist. Every pair consists of two distinct vertices of $V(Y)\setminus\bigcup_{j'}T_{j'}(Y,l)$.
\item[(b)] Every edge of $R_Y$ is a bead of exactly one system of $(Y,l)$. Every port of $R_Y$ is a port of $\mathcal S_0(Y,l)$ or of cherry systems, never of both. A port is an end of at most $M_l-1$ cherries, and so lies in at most $M_l-1$ cherry systems.
\item[(c)] Every vertex lies in at most $\max\bigl(2M_l-2,\ (M_l-1)+\lceil M_l/2\rceil\bigr)=2M_l-2\le t$ pairs of $\mathfrak P_j(Y,l)$ (the equality uses $M_l\ge2$, which holds since $M_l\ge2^{40}$ by~\ref{s2:R2}). If $M_l\ge4$, the value $2M_l-2$ can occur only at a port that has demand $2$ at junction $j$ in each of $M_l-1$ cherry systems (for $M_l\in\{2,3\}$ the bound $(M_l-1)+\lceil M_l/2\rceil$ for ports of $\mathcal S_0(Y,l)$ also equals $2M_l-2$).
\item[(d)] There are pairwise edge-disjoint paths in $\LJS_{Y,l,j}$, one for each pair of $\mathfrak P_j(Y,l)$ and joining its two vertices. Each has length at most $2^{12}L_Y^4$ and all its interior vertices in $T_j(Y,l)$.
\item[(e)] No port or centre of a system of $(Y,l)$ lies in any $T_{j'}(Y,l)$.
\end{enumerate}
\end{claim}
Vertices of systems of \emph{other} pairs $(Y',l)$ may lie in $T_j(Y,l)$ and may be interior vertices of the paths in (d). This is harmless, because hypothesis (D) of Lemma~\ref{s6:lemHCCP} is needed only inside one system (Lemma~\ref{s6:lemHCCglob}).

\begin{proof}[Proof of the claim]
(a) In a system of depth $k\ge2$, the padded loads of all layers are equal (Steps~4 and~5). By Lemma~\ref{s6:lemMED}(b), $\sum_{\LayP_j}\exc^-=\sum_{\LayP_j}\exc^+$. So the number of out-units of layer $j$ equals its padded load, which equals the padded load of layer $j+1$, which is the number of its in-units. For depth $1$, $\sum\exc^-=\sum\exc^+$.

The two vertices of a pair are distinct: for $k\ge2$, the layers $j$ and $j+1$ consist of different, pairwise vertex-disjoint clusters (for $k=2$, layer $j+1$ is layer $j-1\neq j$); for $k=1$, $X^{\mathrm{out}}\cap X^{\mathrm{in}}=\emptyset$. The vertices of pairs are class-$Y$ ports $u\in\Qs_Z$. They lie in $V(Y)$ because $Y=Y(u)\in\anc_l(u)$, and outside every $T_{j'}(Y,l)$ because $\lab_{Y,l}(u)=*$.

(b) Every edge of $R_Y$ lies in exactly one component: in one cluster of $\mathcal S_0(Y,l)$ if the component is non-giant, and otherwise in exactly one cherry, which lies in exactly one class. All edges of $R_Y$ at a port $u$ lie in the component of $u$, so $u$ is a port of $\mathcal S_0(Y,l)$ or of cherries, not both. Distinct cherries with end $u$ use distinct edges of $R_Y$ at $u$, so there are at most $M_l-1$ of them. A class contains at most one of them, since its cherries are vertex-disjoint.

(c) In a system of depth $k\ge2$, a port of layer $j_u$ occurs in junction-$j_u$ pairs ($\dem^-$ times) and in junction-$(j_u-1)$ pairs ($\dem^+$ times), and $j_u\neq j_u-1$ modulo $k$. In a system of depth $1$ it occurs only as a member of $X^{\mathrm{out}}$ or of $X^{\mathrm{in}}$. So at a fixed junction a port occurs in at most $\max(\dem^-,\dem^+)$ pairs of each system containing it. Centres occur in no pair.

By Steps~4 and~5 and part~(b), a port occurs in at most $(M_l-1)+\lceil M_l/2\rceil$ pairs (if it belongs to $\mathcal S_0$) or at most $2(M_l-1)$ pairs (if it belongs to cherry systems). For $M_l\ge2$, $\lceil M_l/2\rceil\le M_l-1$, so both are at most $2M_l-2<t=2M_l+2$.

(d) $Y$ is lend-good. By Definition~\ref{s6:defLending} and Lemma~\ref{s3:lemCOL}(b), $\LJS_{Y,l,j}$ is $(2^{12}L_Y^4,t)$-path connected through $T_j(Y,l)$, as a graph on $V(Y)$. By Cited result~\ref{s1:citDef7} (see Remark~\ref{s3:remMultiset}) this is a statement about \emph{every multiset} of pairs of distinct vertices of $V(Y)$ in which every vertex lies in at most $t$ pairs, with endpoints unrestricted. By Lemma~\ref{s3:lemMonotone}(iii), the property depends only on $(\LJS_{Y,l,j},T_j(Y,l))$, so the multiset may be chosen after the stage-1 outcome and after Steps~1--5. By the hypothesis of the lemma, no edge of $\LJS_{Y,l,j}$ has been used, so the whole class is available. Apply this to $\mathfrak P_j(Y,l)$, which is admissible by (a) and (c).

(e) Ports: see (a). Centres lie outside $V(Y)\supseteq T_{j'}(Y,l)$ (Step~1).
\end{proof}

\emph{Step 7: HCC-P per system, and the union.} Fix a system $\mathcal S$ of $(Y,l)$ of depth $k=k_{\mathcal S}$. For $j<k$ let $\mathcal P_j(\mathcal S)$ be the paths of the joint-routing claim, part~(d), for the junction-$j$ pairs of $\mathcal S$, each oriented from its out-unit to its in-unit. We check the hypotheses of Lemma~\ref{s6:lemHCCP} for $\mathcal S$, with its layers, junction sets $T_j:=T_j(Y,l)$ ($j<k$) and its padding.
\begin{itemize}
\item[(D)] The clusters of $\mathcal S$ are pairwise vertex-disjoint: distinct components for $\mathcal S_0$, vertex-disjoint cherries for a class. Centres lie outside $V(Y)$ and ports inside, so no vertex is both. Neither meets $\bigcup_jT_j(Y,l)$, by part~(e) of the joint-routing claim. The sets $T_j(Y,l)$ are pairwise disjoint (Definition~\ref{s3:defCOL}(iv)).
\item[(P)] Step~2 for $\mathcal S_0$; centres of cherries have degree $2$.
\item[(JC-P)] By parts (a) and (d) of the joint-routing claim, $\mathcal P_j(\mathcal S)$ is a family of edge-disjoint paths from $\LayP_j$ to $\LayP_{j+1}$ (from $X^{\mathrm{out}}$ to $X^{\mathrm{in}}$ if $k=1$), with interiors in $T_j(Y,l)$. By the choice of the bijection, every $u$ is the first vertex of exactly $\dem^-(u)$ of them and the last vertex of exactly $\dem^+(u)$. The path families for distinct $j$ lie in distinct classes $\LJS_{Y,l,j}$, which are disjoint. All of them lie in $\Lend_Y\subseteq E(H_Y)\subseteq E_{r}(Y)$, while beads lie in $\bigcup_ZE_l(Z)$. These sets are disjoint by the partition of Proposition~\ref{s2:propStructure}(iii).
\end{itemize}
Lemma~\ref{s6:lemHCCP} gives sets $F_j(\mathcal S)\subseteq E(\mathcal P_j(\mathcal S))$ such that the beads of $\mathcal S$ together with $\bigcup_jF_j(\mathcal S)$ decompose into at most $\Phi(\mathcal S)$ cycles. The returned edges $E(\mathcal P_j(\mathcal S))\setminus F_j(\mathcal S)$ lie in $G[T_j(Y,l)]$.

Now consider all systems of all $(Y,l)$ together. Their bead sets are pairwise disjoint (part~(b) of the joint-routing claim, and each edge of $\bigcup_ZB_Z\setminus J_l$ lies in exactly one $R_Y$). Their path edges are pairwise disjoint:
\begin{itemize}
\item within one $(Y,l)$ and one $j$, by the joint routing, part~(d) of the claim;
\item for different $j$, because different classes are used;
\item for different $Y\neq Y'$, because $E(H_Y)\cap E(H_{Y'})=\emptyset$ ($E(H_Y)\subseteq E_{r(Y)}(Y)$, and these sets are disjoint by Proposition~\ref{s2:propStructure}(iii)).
\end{itemize}
Path edges are disjoint from beads, as shown above. By Lemma~\ref{s6:lemHCCglob}, the union decomposes into at most $\sum_{\mathcal S}\Phi(\mathcal S)$ cycles.

Put $\LentJS_l:=\bigcup_{\mathcal S}\bigcup_jF_j(\mathcal S)\subseteq\bigcup_{Y,j}\LJS_{Y,l,j}$. Since $\bigcup_ZB_Z=J_l\sqcup\bigsqcup_YR_Y$ and $\bigsqcup_YR_Y$ is the set of all beads of all systems, $\bigcup_ZB_Z\cup\LentJS_l$ is partitioned into the single edges of $J_l$ and these cycles. Each cycle lies inside one system, so it uses beads and lent edges of one ancestor $Y$ only.

\emph{Step 8: the count.} Every bead has exactly one end that is a port of its class (the other end is a centre, possibly a pseudo-hub). Count beads at that end. The classed ports of round $l$ are pairwise distinct vertices, because port sets of one round are disjoint (Proposition~\ref{s2:propStructure}(iv)). Each port has at most $M_l-1$ beads. So there are at most $n(M_l-1)$ beads in total. Hence $\sum_Y\Sigma\Phi(\mathcal S_0(Y,l))\le\frac12n(M_l-1)$ (as $\Phi(\clu_C)\le|E(C)|/2$), and the total number of cherries is at most $\frac12n(M_l-1)$.

The number of ancestors $Y$ of rounds $\le l-2$ is $\nu_l\le2.74n/P_{l-2}$ (Lemma~\ref{s2:lemTower}(b), which completes (K1) of Proposition~\ref{s2:propOV}). Moreover $\gamma_l\le P_{l-2}/M_l$ and $P_{l-2}\ge M_l^{13}$ (Lemma~\ref{s2:lemTower}(b)). By \eqref{s6:eqSzeroCost} and \eqref{s6:eqCherryCost}, the number of cycles is at most
\[
\frac{1.5\cdot\frac12n(M_l-1)}{M_l^2}+5\,\frac{P_{l-2}}{M_l}\cdot\frac{2.74n}{P_{l-2}}+\frac{\frac12n(M_l-1)}{M_l^2}+6M_l\cdot\frac{2.74n}{P_{l-2}}+\sum_{(Y,l)\text{ giant}}\sco_{Y,l}
\le15\,\frac n{M_l}+\sum_{(Y,l)\text{ giant}}\sco_{Y,l}.
\]
The four terms are at most $0.75n/M_l$, $13.7n/M_l$, $0.5n/M_l$ and $16.44\,n/M_l^{12}$, and $14.95+16.44\,M_l^{-11}\le15$ because $M_l\ge2^{40}$. By Step~5, $\sco_{Y,l}\le1.5m_{Y,l}$ for every giant $(Y,l)$. Since $15\le126$, the number of cycles is at most $126n/M_l+1.5\sum_{(Y,l)\text{ giant}}m_{Y,l}$, which proves the lemma as stated.

\emph{On the constant.} The proof gives the constant $15$. The lemma states the looser constant $126$ on purpose: $126$ is the value carried into Theorem~\ref{s6:thmMIXC}(c) (the term $252\,n/\Dstar$ in $\varepsilon_{\mathrm M}$) and from there into Section~7. Any absolute constant would do there, because the term is $O(n/\Dstar)$. Replacing $126$ by $15$ would only lower $252/\Dstar$ to $30/\Dstar$ in those places.

\emph{The sharper bound.} Step~8 proves, as a by-product, the sharper bound of at most $15n/M_l+\sum_{(Y,l)\text{ giant}}\sco_{Y,l}$ cycles, where $\sco_{Y,l}=(3g_{Y,l}-6M_l)^+\le(1.5m_{Y,l}-12)^+$ (Step~5). The stated bound follows from it, and it is the form used by the only consumer, Theorem~\ref{s6:thmMIXC}(c); nothing downstream uses the sharper bound.

No cost term is paid per layer or per component. The cost of a system is its common padded load $\Phi(\mathcal S)$. Padding adds paths, not cycles, and returned edges are junk. The only per-(centre, class) costs are the single edges moved into $J_l$ in Step~2, and these are counted in \eqref{s6:eqJbound}.
\end{proof}

\begin{remark}[the star split is not used]\label{s6:remStar}
For a giant pair $(Y,l)$ one could instead use a ``star split'' with a concentration threshold $\tau'_l$ and a hub-concentration functional $\mathrm{HC}$. That option is not used in this paper, and neither $\tau'_l$ nor $\mathrm{HC}$ appears. Its cost is bounded only by $\frac12\mathrm{HC}_{Y,l}$, which is not controlled here; the cherry split costs at most $\#\{\text{cherries}\}/\KJS_l+6M_l+1.5m_{Y,l}$ by \eqref{s6:eqCherryCost}. Neither split emits single edges: all single edges of JS-LC come from Step~2, before any split. So $J_l$ does not depend on this choice, and neither does any property of $J_l$ proved below.
\deps{\ref{s6:lemJSLC}}
\end{remark}

\begin{lemma}[Lemma J$^+$: the J-interface]\label{s6:lemJplus}
In Lemma~\ref{s6:lemJSLC}, the set $J_l$ consists exactly of the deletions of Step~2 of its proof:
\begin{itemize}
\item the PAR deletions, per part and per pair of classes;
\item then one bead per (centre, class) with odd realized class degree, where the degree at a hub is aggregated over all standalone parts of round $l$.
\end{itemize}
This holds for \emph{every} choice of which edges are deleted. Steps~3--8 delete nothing.

Moreover $J_l=\Jlost_l\sqcup\Jhub_l\sqcup\Jfr_l\sqcup\Jpar_l$, where, for $Z\in\Std_l$:
\begin{itemize}
\item a $\Jpar$-edge has both ends in $\Qs_Z$, with different classes;
\item a $\Jhub$-edge is an edge $hu$ with $h\in A_Z\subseteq D_l$ and $u\in\Qs_Z$;
\item a $\Jfr$-edge is an edge $xu$ with $x\in F_Z$ and $u\in\Qs_Z$;
\item a $\Jlost$-edge is an edge $vu$ with $v\in\Lost_Z$ (a lost centre) and $u\in\Qs_Z$.
\end{itemize}
The \emph{class} of a $\Jhub$-, $\Jfr$- or $\Jlost$-edge $hu$ is $Y(u)$. The following properties hold.
\begin{itemize}
\item[\Jp{1}]\namedlabel{s6:J1}{\Jp{1}} For each hub $h\in D_l$ and each class $Y$, $J_l$ contains at most one $\Jhub$-edge of class $Y$ at $h$, aggregated over all parts of round $l$. For each fresh centre $x$ and each class $Y$, $J_l$ contains at most one $\Jfr$-edge of class $Y$ at $x$. (The same holds at lost centres.)
\item[\Jp{2}]\namedlabel{s6:J2}{\Jp{2}} Every J-edge lying in $E_l(Z)$ has an end in $\Qs_Z$. For every $Z$ and every vertex, at most $M_l-1$ J-edges at that vertex lie in $E_l(Z)$. In particular a port or a fresh centre (a vertex outside $D_l$) carries at most $M_l-1$ J-edges of round $l$. Also $|J_l|\le n(M_l-1)$.
\item[\Jp{3}]\namedlabel{s6:J3}{\Jp{3}} The sets $D_l$ (hubs), $\bigcup_ZF_Z$ (fresh centres) and $\bigcup_Z\Qs_Z$ (classed non-lost ports) are pairwise disjoint. Each vertex outside $D_l$ that lies in a round-$l$ pre-part lies in exactly one such pre-part, so it has one role and, if classed, one class.
\end{itemize}
In addition, the following interface facts hold.
\begin{enumerate}
\item[(i)] The types above are exhaustive and exclusive.
\item[(ii)] (Size cap at fresh centres.) A fresh centre carries at most $M_l-1$ $\Jfr$-edges.
\item[(iii)] (Disjointness from \ref{s2:R5}.) $E(H_Y)\subseteq E_{r}(Y)$ for every ancestor $Y$ of round $r$, and the sets $E_r(Y)$ ($Y$ light) and $E_l(Z)$ ($Z\in\Std_l$) are pairwise disjoint.
\item[(iv)] $u\in V(Y(u))$ for every classed port $u$.
\item[(v)] In the colourings of Definition~\ref{s3:defCOL}, colours of distinct edges are independent.
\end{enumerate}
\deps{\ref{s6:lemJSLC}, \ref{s6:lemPAR}, \ref{s6:defDesign}, \ref{s2:propStructure}, \ref{s2:lemEL}, \ref{s2:lemCap}, \ref{s3:defCOL}, \ref{s2:defHBtp}, \ref{s6:lemHCCP}, \ref{s6:defLending}}
\end{lemma}

\begin{proof}
\emph{The composition of $J_l$.} In the proof of Lemma~\ref{s6:lemJSLC}, single edges are created only in Step~2. Steps~3--8 form clusters from the sets $R_Y$, route paths and apply Lemma~\ref{s6:lemHCCP}, whose output consists of cycles only. Every bead lies in some system and is covered, and padding adds paths, not edges. Step~2 allowed any admissible choice in Lemma~\ref{s6:lemPAR} and any choice of the moved edge, and nothing below depends on these choices.

\emph{Types.} A PAR deletion lies in some $E_{ab}(Z)$, so it joins two ports of $\Qs_Z$ of different classes: a $\Jpar$-edge. A moved edge is a class-$Y$ edge whose centre has odd degree in $R^0_Y$. By Step~2 of that proof this centre lies in $\Ret_Z=A_Z\sqcup F_Z\sqcup\Lost_Z$, and the port lies in $\Qs_Z$. The edge is a $\Jhub$-, $\Jfr$- or $\Jlost$-edge according as the centre lies in $A_Z=Z^0\cap D_l$, in $F_Z$ or in $\Lost_Z$. These cases are exclusive because $A_Z$, $F_Z$, $\Lost_Z$ and $\Qs_Z$ are pairwise disjoint (Definition~\ref{s6:defLending}). This proves (i).

\Jp{1}: The moves are made once per pair (centre, class), with the class-$Y$ degree at a hub counted over all parts of round $l$. PAR deletions are port--port edges and never $\Jhub$- or $\Jfr$-edges. A fresh centre lies in one part only.

\Jp{2}: $J_l\subseteq\bigcup_ZB_Z$, and every edge of $B_Z$ has an end in $\Qs_Z$. By Lemma~\ref{s2:lemCap}(ii), every vertex has at most $M_l-1$ edges of $E_l(Z)$ for each $Z$. A vertex outside $D_l$ lies in at most one round-$l$ pre-part, which gives the ``in particular''. For the total, charge every J-edge to an end in $\Qs_Z$. The classed ports of round $l$ are distinct vertices (Proposition~\ref{s2:propStructure}(iv)), and each receives at most $M_l-1$ charges. So $|J_l|\le n(M_l-1)$.

\Jp{3}: $A_Z\subseteq D_l$, while $F_Z$ and $\Qs_Z$ are disjoint subsets of $U_Z=Z^0\setminus D_l$. The sets $U_Z$ of distinct $Z\in\Std_l$ are pairwise disjoint, and every vertex outside $D_l$ lies in at most one round-$l$ pre-part (Proposition~\ref{s2:propStructure}(iv)).

(ii) is the degree bound of \Jp{2} at a fresh centre. (iii) is Proposition~\ref{s2:propStructure}(iii). (iv) holds because $Y(u)\in\anc_l(u)$ (Definition~\ref{s6:defDesign}), and $\anc_l(u)$ consists of ancestors whose vertex set contains $u$. (v) holds because every colouring in Definition~\ref{s3:defCOL} is uniform and independent over edges. By Lemma~\ref{s2:lemEL}, the two ends of a $\Jpar$-edge indeed have different classes.
\end{proof}

\begin{definition}[round-$l$ J-consumer]\label{s6:defJconsumer}
A \emph{J-consumer} is a rule that acts at each round $3\le l\le R$. Given everything constructed at rounds $>l$, the stage-1 outcome and the set $J_l$ (and possibly fresh randomness), it outputs a family $\Obj_l$ of objects and a set $\LentJV_l$ of edges such that:
\begin{itemize}
\item[\JC{1}]\namedlabel{s6:JC1}{\JC{1}} $\LentJV_l\subseteq\bigcup\{\LJV_{Y,l}:\ Y\text{ a lend-good ancestor of a round}\le l-2\}$;
\item[\JC{2}]\namedlabel{s6:JC2}{\JC{2}} the objects of $\Obj_l$ are pairwise edge-disjoint cycles and single edges of $G$, and their union is exactly $J_l\cup\LentJV_l$;
\item[\JC{3}]\namedlabel{s6:JC3}{\JC{3}} $\Obj_l$ uses no other edge.
\end{itemize}
By the exclusive-consumer rule COL(g) of Definition~\ref{s3:defCOL}, the classes $\LJV_{Y,l}$ are used by nothing except the round-$l$ J-consumer. For example, the rule ``$\Obj_l:=$ the edges of $J_l$ as single edges, $\LentJV_l:=\emptyset$'' is a J-consumer. The J-consumer used for the final count is the junction-vertex round step constructed later in the paper. Here $J_l$ has the properties of Lemma~\ref{s6:lemJplus}.
\deps{\ref{s3:defCOL}, \ref{s6:lemJplus}}
\end{definition}

\subsection{The MIX-C assembly}

\begin{construction}[the processing order $R,R-1,\dots,1$, and the Lent sets]\label{s6:consOrder}
\emph{Input.}
\begin{itemize}
\item A valid $\HBtp$ run on $G$ with $n\ge\Nzero$ and $d_1\ge\Dstar$, and a designation $\delta$.
\item A stage-1 outcome with its stage-2 data (Definitions~\ref{s5:defStages} and~\ref{s6:defLending}).
\item Stage-3 outcomes: labels for Lemma~\ref{s4:lemTPV} for every $Z\in\Std$; child-vortex labels (Lemma~\ref{s5:lemChild}) for every non-demoted light part; labels for Theorem~\ref{s4:thmVXp} (Lemma~\ref{s5:lemDemoted}) for every demoted light part. Each is fixed in its good event.
\item A J-consumer (Definition~\ref{s6:defJconsumer}).
\end{itemize}

\emph{The Lent sets} are defined recursively along the construction. For an ancestor $Y$ of round $r$:
\begin{itemize}
\item $\LentU(Y)$: the U-lent edges of $Y$ used by Lemma~\ref{s5:lemParent} at rounds $l\ge r+2$ (light $Y$ only; $\LentU(Y):=\emptyset$ for standalone $Y$);
\item $\LentJS(Y):=\bigcup_{l\ge r+2}\bigl(\LentJS_l\cap\bigcup_j\LJS_{Y,l,j}\bigr)$, the JS-lent edges of $Y$ on output cycles of JS-LC (returned edges are not included);
\item $\LentJV(Y):=\bigcup_{l\ge r+2}\bigl(\LentJV_l\cap E(H_Y)\bigr)$;
\item $\Lent(Y):=\LentU(Y)\cup\LentJS(Y)\cup\LentJV(Y)$.
\end{itemize}
For standalone $Y$ put $\Erem(Y):=E_r(Y)\setminus\Lent(Y)$. For light non-demoted $Y$ put $H_0(Y):=E_r(Y)\setminus\Lent(Y)$.

\emph{The construction.} For $l=R,R-1,\dots,1$, in this order, do:
\begin{enumerate}
\item[(1)] For every $Z\in\Std_l$, the set $\Lent(Z)$ is already determined. Its edges lie in lent classes of $Z$ indexed by rounds $\ge l+2$, and by COL(g) of Definition~\ref{s3:defCOL} the only consumers of these classes act at steps of rounds $\ge l+2$, which are complete. Apply Lemma~\ref{s4:lemTPV} on its fixed good event, with vertex set $Z^0$, expander $O:=O_Z$, edge set $E:=\Erem(Z)$, $\vP:=\Ret_Z$ and $\vQ:=\Qs_Z$. This gives $\Erem(Z)=\EV(Z)\sqcup\EQ(Z)$ and a decomposition of $\EV(Z)$ into at most $169|\Ret_Z|$ objects, which are output.
\item[(2)] Every light part $Y$ of round $l$ that is not demoted runs its child vortex (Lemma~\ref{s5:lemChild}) on $H_0(Y)$. Every demoted light part $Y$ of round $l$ runs Theorem~\ref{s4:thmVXp} on $E_l(Y)$ (Lemma~\ref{s5:lemDemoted}). If $l\ge3$, then for every phase $c\in[4]$ the U-bundle $\Bdl_{l,c}$ is chained through the U-classes of good parents of rounds $\le l-2$ (Lemma~\ref{s5:lemParent}), and the used U-lent edges are recorded. The objects are output.
\item[(3)] If $l\ge3$: after all round-$l$ applications of Lemma~\ref{s4:lemTPV}, apply Lemma~\ref{s6:lemJSLC} with $B_Z:=\EQ(Z)$ for all $Z\in\Std_l$. Its cycles are output, and $\LentJS_l$ is recorded. This yields $J_l$.
\item[(4)] If $l\ge3$: apply the J-consumer to $J_l$. The objects $\Obj_l$ are output, and $\LentJV_l$ is recorded.
\end{enumerate}
Finally, the cycles of every $\Cyc_l$ are output, and every edge of $E_0$ is output as a single edge.

Rounds $1$ and $2$ have no classed ports ($\anc_l(x)=\emptyset$ for $l\le2$), so $J_2=J_1=\emptyset$ and steps (3)--(4) are void there.
\deps{\ref{s4:lemTPV}, \ref{s5:lemChild}, \ref{s5:lemParent}, \ref{s5:lemDemoted}, \ref{s6:lemJSLC}, \ref{s6:defJconsumer}, \ref{s6:defLending}, \ref{s5:defStages}, \ref{s4:thmVXp}, \ref{s3:defCOL}}
\end{construction}

The construction is well defined because each set $\Lent(Y)$ is determined by rounds $\ge r(Y)+2$, and these precede round $r(Y)$ in the order $R,R-1,\dots,1$. This is part~(i) of the following lemma.

\begin{lemma}[Lent sets]\label{s6:lemLent}
In Construction~\ref{s6:consOrder}, for every ancestor $Y$ of round $r$:
\begin{enumerate}
\item[(i)] $\Lent(Y)\subseteq\Lend_Y$, and $\Lent(Y)$ depends only on the steps at rounds $\ge r+2$. So it is final when $Y$ is processed at round $r$.
\item[(ii)] If $Y$ is standalone and lend-bad, then $\Lent(Y)=\emptyset$: every port of class $Y$ is lost, so $Y$ has no JS system, and \JC{1} excludes its JV classes. If $Y$ is light and demoted, then $\Lent(Y)=\emptyset$.
\item[(iii)] $\Erem(Z)\supseteq E(O_Z)$ for every $Z\in\Std$, and $H_0(Y)\supseteq\Own_Y$ for every non-demoted light $Y$.
\item[(iv)] Put $\Lentext(Y):=\LentJS(Y)\cup\LentJV(Y)$ for light $Y$. The family $(\Lentext(Y))_Y$ satisfies the hypothesis of Lemma~\ref{s5:lemKRED}, and $H_0(Y)=E_r(Y)\setminus(\LentU(Y)\cup\Lentext(Y))$ is the set $H_0(Y)$ of that lemma. The sets $\LentU(Y)$, $\LentJS(Y)$ and $\LentJV(Y)$ are pairwise disjoint.
\end{enumerate}
\deps{\ref{s6:consOrder}, \ref{s6:defLending}, \ref{s6:defJconsumer}, \ref{s6:lemJSLC}, \ref{s5:lemParent}, \ref{s5:lemKRED}, \ref{s3:defCOL}, \ref{s6:lemJplus}, \ref{s5:defStages}, \ref{s4:lemTPV}, \ref{s5:lemChild}, \ref{s5:lemDemoted}, \ref{s4:thmVXp}}
\end{lemma}

\begin{proof}
(i) By Lemma~\ref{s5:lemParent}(i), the used U-lent edges lie in U-classes $\LU_{Y',l,c,\uslot}$ of good parents $Y'$, so $\LentU(Y)\subseteq\bigcup\LU_{Y,l,c,\uslot}\subseteq\Lend_Y$. By definition $\LentJS(Y)\subseteq\bigcup\LJS_{Y,l,j}\subseteq\Lend_Y$.

For $\LentJV(Y)$: by \JC{1}, $\LentJV_l\subseteq\bigcup_{Y'}\LJV_{Y',l}$, and $\LJV_{Y',l}\subseteq E(H_{Y'})$. For $Y'\neq Y$ the sets $E(H_{Y'})$ and $E(H_Y)$ are disjoint (Lemma~\ref{s6:lemJplus}(iii)). Hence $\LentJV_l\cap E(H_Y)\subseteq\LJV_{Y,l}\subseteq\Lend_Y$.

The classes $\LU_{Y,l,\cdot,\cdot}$, $\LJS_{Y,l,\cdot}$ and $\LJV_{Y,l}$ exist only for $r+2\le l\le R$ (Definition~\ref{s3:defCOL}(ii)). By COL(g) of Definition~\ref{s3:defCOL}, each is used only by its round-$l$ consumer: the U-bundle $(l,c)$ in step~(2), the JS-LC system of $(Y,l)$ in step~(3), or the round-$l$ J-consumer in step~(4).

No other step uses an edge of $E(H_Y)$ before round $r$. At a round $l'>r$:
\begin{itemize}
\item step~(1) uses only edges of $\Erem(Z')\subseteq E_{l'}(Z')$;
\item the child vortices and VX$^+$ use only edges of $E_{l'}(Y')$;
\item JS-LC uses beads in $E_{l'}(\cdot)$ and lent edges of classes $\LJS_{\cdot,l',\cdot}$;
\item the J-consumer uses only $J_{l'}\cup\LentJV_{l'}$ (\JC{3}).
\end{itemize}
The sets $E_{l'}(\cdot)$ are disjoint from $E_r(Y)\supseteq E(H_Y)$ (Lemma~\ref{s6:lemJplus}(iii)). So $\Lent(Y)$ is determined by steps of rounds $\ge r+2$, all of which come before round $r$.

(ii) Let $Y$ be lend-bad, standalone or light.
\begin{itemize}
\item A JS-LC system of $(Y,l)$ exists only if $R_Y\neq\emptyset$ at round $l$, i.e.\ only if $Y$ is the class of a port of some $\Qs_Z$. Since $Y$ is lend-bad, Definition~\ref{s6:defLending} puts every port $u$ with $Y(u)=Y$ into $\Lost_Z$, so $Y$ is not the class of any port of any $\Qs_Z$. Hence $R_Y=\emptyset$ at every round $l$, no pair is routed in any $\LJS_{Y,l,j}$, and $\LentJS(Y)=\emptyset$.
\item By \JC{1} and the disjointness used in (i), $\LentJV_l\cap E(H_Y)\subseteq\LJV_{Y,l}$ can be non-empty only if $Y$ is lend-good. So $\LentJV(Y)=\emptyset$.
\item $\LentU(Y)=\emptyset$ for standalone $Y$ by definition. For a demoted light $Y$ it is empty because Lemma~\ref{s5:lemParent} uses only U-classes of good parents, and a demoted part is not a good parent (Definition~\ref{s5:defStages}). A demoted light part is lend-bad by Definition~\ref{s6:defLending}.
\end{itemize}

(iii) Let $Z\in\Std_l$.
\begin{itemize}
\item If $Z$ is lend-good, then $O_Z=\Own_Z$. Moreover $E(\Own_Z)\subseteq E(X^0_Z)=E(H_Z)$, and $E(H_Z)\subseteq E_l(Z)$ by Lemma~\ref{s6:lemJplus}(iii). By (i), $\Lent(Z)\subseteq\Lend_Z$, and $\Lend_Z\cap\Own_Z=\emptyset$ (Definition~\ref{s3:defCOL}(i)). So $E(O_Z)\subseteq E_l(Z)\setminus\Lent(Z)=\Erem(Z)$.
\item If $Z$ is lend-bad, then $\Lent(Z)=\emptyset$ by (ii), so $\Erem(Z)=E_l(Z)\supseteq E(X^0_Z)=E(O_Z)$.
\end{itemize}
For a non-demoted light $Y$: $\Own_Y\subseteq E(X_Y)\subseteq E_r(Y)$ and $\Own_Y\cap\Lend_Y=\emptyset$, so $\Own_Y\subseteq H_0(Y)$ by (i).

(iv) By the argument of (i), $\Lentext(Y)\subseteq\bigcup_{l,j}\LJS_{Y,l,j}\cup\bigcup_l\LJV_{Y,l}$. By (ii), $\Lentext(Y)=\emptyset$ if $Y$ is demoted. $\LentU(Y)$ is by definition the set of U-lent edges of $Y$ used by Lemma~\ref{s5:lemParent}, exactly as in Lemma~\ref{s5:lemKRED}. So $H_0(Y)=E_r(Y)\setminus\Lent(Y)=E_r(Y)\setminus(\LentU(Y)\cup\Lentext(Y))$. The three sets lie in distinct colour classes of the colouring of $\Lend_Y$ (Definition~\ref{s3:defCOL}(ii)), so they are pairwise disjoint.
\end{proof}

\begin{theorem}[Theorem MIX-C: the assembly with $\mathcal V=\emptyset$, deterministic form]\label{s6:thmMIXC}
Fix the following data:
\begin{itemize}
\item a valid $\HBtp$ run on $G$ with $n\ge\Nzero$ and $d_1\ge\Dstar$, and a designation $\delta$;
\item \emph{any} stage-1 outcome;
\item stage-3 outcomes in the good events of Lemma~\ref{s4:lemTPV} (every $Z\in\Std$), of Lemma~\ref{s4:lemPV} through Lemma~\ref{s5:lemChild} (every non-demoted light part), and of Theorem~\ref{s4:thmVXp} through Lemma~\ref{s5:lemDemoted} (every demoted light part);
\item \emph{any} J-consumer.
\end{itemize}
Run Construction~\ref{s6:consOrder}. Then:
\begin{enumerate}
\item[(a)] (Order.) At step~(1) of round $l$, $\Lent(Z)$ is final for every $Z\in\Std_l$. JS-LC at round $l$ uses only classes $\LJS_{Y,l,\cdot}$, and the J-consumer only classes $\LJV_{Y,l}$, of ancestors $Y$ of rounds $\le l-2$, which are processed later. The hypotheses of Lemma~\ref{s6:lemJSLC} hold at every round $l\ge3$, and every device is applied within its hypotheses.
\item[(b)] (Coverage.) The output is a decomposition of $E(G)$: every edge of $G$ lies in exactly one output object. In detail:
  \begin{itemize}
  \item for a light $Y$ of round $r$, $E_r(Y)=\LentU(Y)\sqcup\LentJS(Y)\sqcup\LentJV(Y)\sqcup H_0(Y)$ if $Y$ is not demoted, and $\Lent(Y)=\emptyset$ if $Y$ is demoted;
  \item for $Z\in\Std_l$, $E_l(Z)=\Lent(Z)\sqcup\EV(Z)\sqcup(\EQ(Z)\setminus J_l)\sqcup(J_l\cap E_l(Z))$, where the third set consists of the beads of $Z$;
  \item every lent edge is used by at most one object;
  \item lent edges that are unused or returned are covered by the device of their owner, and every device tolerates such junk;
  \item the bound \eqref{s6:eqJbound} on $|J_l|$ is computed on the sets $E_l(Z)$ with lent edges included, so lending never increases it (that bound is not used later; see the remark after it).
  \end{itemize}
\item[(c)] (Cost.) The number of output objects is at most
\[
\Bigl(\frac{\Dstar}2+\cKRED+\cFresh\Bigr)n+\varepsilon_{\mathrm M}(\Dstar)\,n+\Bigl[80\,\mathrm{dem}+369\,\mathrm{lp}+169\sum_l|\Lost_l|\Bigr]+1.5\sum_{(Y,l)}m_{Y,l}+\sum_l|\Obj_l|,
\]
where $\varepsilon_{\mathrm M}(\Dstar):=\epsK(\Dstar)+169\,\epsA+252/\Dstar$, with $\epsA=31\eps/(\Cp\log\log\Dstar)$ as defined in Lemma~\ref{s2:lemTower}(e). Thus $\varepsilon_{\mathrm M}$ is a function of $\Dstar$ alone; it depends neither on the run nor on $n$. Moreover $1.5\sum_{(Y,l)}m_{Y,l}\le1.5\,\epsCONC(\Dstar)\,n$ by Theorem~\ref{s6:thmCONCL}(iv).
\end{enumerate}
The itemization of (c) is:
\begin{itemize}
\item K-RED: $(\Dstar/2+\cKRED+\epsK)n+80\,\mathrm{dem}+369\,\mathrm{lp}$, including $\Cyc$ and $E_0$;
\item TPV: $169\sum(|A_Z|+|F_Z|+|\Lost_Z|)\le169(2+\epsA)n+169\sum_l|\Lost_l|$;
\item JS-LC cycles: $\sum_l126n/M_l+1.5\sum m\le252n/\Dstar+1.5\sum m$;
\item J-consumer: $\sum_l|\Obj_l|$. The set $J_l$ is handed entirely to the J-consumer.
\end{itemize}
\deps{\ref{s6:consOrder}, \ref{s4:lemTPV}, \ref{s5:lemKRED}, \ref{s5:lemDemoted}, \ref{s6:lemJSLC}, \ref{s6:lemJplus}, \ref{s6:defJconsumer}, \ref{s6:lemLent}, \ref{s6:thmCONCL}, \ref{s2:propStructure}, \ref{s2:propParentless}, \ref{s2:propOV}, \ref{s2:lemTower}, \ref{s3:defCOL}, \ref{s3:lemCOL}, \ref{s4:lemPV}, \ref{s5:lemChild}, \ref{s4:thmVXp}, \ref{s2:defHBtp}, \ref{s1:condG3}, \ref{s6:defDesign}, \ref{s6:defLending}}
\end{theorem}

\begin{proof}
\emph{(a) Order.}
\begin{itemize}
\item \emph{Finality.} By Lemma~\ref{s6:lemLent}(i), $\Lent(Z)$ for $Z\in\Std_l$ depends only on steps of rounds $\ge l+2$. So it is final at step~(1) of round $l$. The same holds for $H_0(Y)$ of every light part $Y$ of round $l$ at step~(2).
\item \emph{Classes used.} At round $l$, JS-LC routes only in classes $\LJS_{Y,l,j}$ with $Y$ the class of a non-lost port of round $l$, so $Y\in\anc_l(u)$ has round $\le l-2$. The J-consumer uses only classes $\LJV_{Y,l}$ with $r(Y)\le l-2$, by \JC{1}. Such $Y$ are processed at round $r(Y)<l$, i.e.\ later.
\item \emph{Hypotheses of Lemma~\ref{s6:lemJSLC}.} Take $B_Z:=\EQ(Z)\subseteq\Erem(Z)\subseteq E_l(Z)$. By \ref{s2:R5} every edge of $E_l(Z)$ lies inside $Z^0=\Ret_Z\sqcup\Qs_Z$. By \ref{s4:T1} of Lemma~\ref{s4:lemTPV}, every edge of $\Erem(Z)$ with both ends in $\Ret_Z$ lies in $\EV(Z)$. So every edge of $B_Z$ has an end in $\Qs_Z$. By COL(g) of Definition~\ref{s3:defCOL}, the only consumer of $\LJS_{Y,l,j}$ is the JS-LC system of $(Y,l)$ at step~(3) of round $l$. By the argument in the proof of Lemma~\ref{s6:lemLent}(i), no other step up to that moment uses an edge of $E(H_Y)$. So no edge of these classes has been used.
\item \emph{Lemma~\ref{s4:lemTPV} at step (1).}
  \begin{itemize}
  \item $|Z^0|\ge P_l\ge\Nzero$ by \ref{s1:condG3}.
  \item If $Z$ is lend-good, then $O_Z=\Own_Z$. By Lemma~\ref{s3:lemCOL}(a) it is a spanning $(2^{-5},s_l/4)$-expander on $Z^0$, and by Lemma~\ref{s3:lemCOL}(e) we have $s_l/4\ge2^{150}L_Z^{42}$.
  \item If $Z$ is lend-bad, $O_Z=X^0_Z$ is a spanning $(2^{-5},s_l)$-expander on $Z^0$ (Proposition~\ref{s2:propStructure}(i)), and $s_l\ge s_l/4\ge2^{150}L_Z^{42}$: the inequality of Lemma~\ref{s3:lemCOL}(e) involves only $s_l$ and $L_Z$, not the colouring.
  \item $E=\Erem(Z)\supseteq E(O_Z)$ (Lemma~\ref{s6:lemLent}(iii)), and all its edges lie inside $Z^0$.
  \item The good event of Lemma~\ref{s4:lemTPV} is determined by $(Z^0,O_Z,\Ret_Z)$ and the stage-3 labels, all fixed before step~(1). On it, the conclusion holds for every such $E$. So TPV tolerates the junk (unused and returned lent edges) left in $\Erem(Z)$.
  \end{itemize}
\item \emph{Light parts at step (2).} Let $Y$ be a light part of round $l$. If $Y$ is not demoted, then $\Own_Y\subseteq H_0(Y)\subseteq E_l(Y)$ by Lemma~\ref{s6:lemLent}(iii). This is what Lemma~\ref{s5:lemChild} needs, and the child vortex covers every such $H_0(Y)$. If $Y$ is demoted, then $\Lent(Y)=\emptyset$ by Lemma~\ref{s6:lemLent}(ii). So Lemma~\ref{s5:lemDemoted} applies to $E:=E_l(Y)\supseteq E(X_Y)$, which may contain arbitrary extra edges.
\item \emph{Step (4).} The J-consumer receives $J_l$, which has the properties of Lemma~\ref{s6:lemJplus}.
\end{itemize}

\emph{(b) Coverage.} Split the output objects into four families.
\begin{itemize}
\item[$\mathcal O_{\mathrm K}$:] the cycles of all $\Cyc_l$, the edges of $E_0$, and the objects of step~(2) at all rounds;
\item[$\mathcal O_{\mathrm T}$:] the objects of step~(1);
\item[$\mathcal O_{\mathrm S}$:] the cycles of step~(3);
\item[$\mathcal O_{\mathrm J}$:] $\bigsqcup_l\Obj_l$.
\end{itemize}

\emph{$\mathcal O_{\mathrm K}$ is the output of Lemma~\ref{s5:lemKRED}.} By Lemma~\ref{s6:lemLent}(iv), the final family $(\Lentext(Y))_Y$ satisfies the hypothesis of Lemma~\ref{s5:lemKRED}, and its $H_0(Y)$ are those of the construction. That lemma processes rounds $R,\dots,1$. Its step at round $l$ depends only on $H_0(Y)$ for light parts $Y$ of round $l$ and on its own earlier steps. At round $l$ the sets $H_0(Y)$ of the construction are already final (part (a)). So the objects of step~(2) over all rounds, together with $\Cyc$ and $E_0$, are exactly the output of Lemma~\ref{s5:lemKRED} for this family. Hence $\mathcal O_{\mathrm K}$ decomposes
\[
E_{\mathrm K}:=E(G)\setminus\Kstd\setminus\textstyle\bigsqcup_{Y\text{ light}}\Lentext(Y)
=\bigsqcup_l\Cyc_l\sqcup E_0\sqcup\bigsqcup_{Y\text{ light}}\bigl(E_r(Y)\setminus\Lentext(Y)\bigr),
\]
where the equality is Proposition~\ref{s2:propStructure}(iii).

\emph{The other three families.}
\begin{itemize}
\item By \ref{s4:T3} and \ref{s4:T4}, $\mathcal O_{\mathrm T}$ decomposes $\bigsqcup_Z\EV(Z)$.
\item By Lemma~\ref{s6:lemJSLC}, at each round $l\ge3$ the cycles of step~(3) decompose $\bigl(\bigsqcup_{Z\in\Std_l}\EQ(Z)\setminus J_l\bigr)\sqcup\LentJS_l$.
\item By \JC{2}, $\Obj_l$ decomposes $J_l\sqcup\LentJV_l$. These two sets are disjoint, because $J_l\subseteq\bigcup_ZE_l(Z)$ and $\LentJV_l\subseteq\bigcup_YE(H_Y)$ with $r(Y)\le l-2$ (Lemma~\ref{s6:lemJplus}(iii)).
\end{itemize}
Every edge of $\LentJS_l$ lies in exactly one class $\LJS_{Y,l,j}$. By \JC{1} and Lemma~\ref{s6:lemJplus}(iii), every edge of $\LentJV_l$ lies in exactly one $E(H_Y)$. Distinct $l$ use distinct classes. Hence
\[
\bigsqcup_l\LentJS_l=\bigsqcup_Y\LentJS(Y),\qquad\bigsqcup_l\LentJV_l=\bigsqcup_Y\LentJV(Y),
\]
with all unions disjoint.

\emph{Putting the pieces together.} Since $\Erem(Z)=\EV(Z)\sqcup\EQ(Z)$ and $J_l\cap E_l(Z)\subseteq\EQ(Z)$,
\[
\EV(Z)\sqcup(\EQ(Z)\setminus J_l)\sqcup(J_l\cap E_l(Z))=\Erem(Z)=E_l(Z)\setminus\Lent(Z),
\]
and $\Lent(Z)=\LentJS(Z)\sqcup\LentJV(Z)$ for standalone $Z$. For light non-demoted $Y$,
\[
E_r(Y)\setminus\Lentext(Y)=\LentU(Y)\sqcup H_0(Y)\qquad\text{(Lemma~\ref{s6:lemLent}(iv))},
\]
and for demoted $Y$, $\Lentext(Y)=\emptyset$. Therefore the edge sets of the four families are pairwise disjoint, and their union is
\[
E_{\mathrm K}\sqcup\bigsqcup_{Z\in\Std}E_l(Z)\sqcup\bigsqcup_{Y\text{ light}}\Lentext(Y)=E(G),
\]
again by Proposition~\ref{s2:propStructure}(iii).

Each family decomposes its own edge set, so together they decompose $E(G)$. In particular every lent edge lies in at most one object. A lent edge that is not used (not in $\Lent(Y)$), including every returned edge of JS-LC, stays in $\Erem(Y)$ or $H_0(Y)$ and is covered by the owner's device, which tolerates it by part~(a). The last item of (b) is the remark after \eqref{s6:eqJbound}. Its right side involves only $E_l(Z)$, $\delta$ and $\Lost_Z$, not $\Lent$, the TPV split or earlier rounds.

\emph{(c) Cost.}
\begin{itemize}
\item \emph{$\mathcal O_{\mathrm K}$.} By Lemma~\ref{s5:lemKRED}, $|\mathcal O_{\mathrm K}|\le(\Dstar/2+\cKRED+\epsK(\Dstar))n+80\,\mathrm{dem}+369\,\mathrm{lp}$.
\item \emph{$\mathcal O_{\mathrm T}$.} By \ref{s4:T3}, $|\mathcal O_{\mathrm T}|\le\sum_l\sum_{Z\in\Std_l}169|\Ret_Z|=169\sum_{l,Z}(|A_Z|+|F_Z|+|\Lost_Z|)$. Here $\sum_{l,Z}|A_Z|\le\sum_l15.2\eps n/\log P_l\le\epsA n$ by Proposition~\ref{s2:propOV}(K2) and Lemma~\ref{s2:lemTower}(e). Also $\sum_{l,Z}|F_Z|\le2n$ by Proposition~\ref{s2:propParentless}(i). Finally $\sum_{Z\in\Std_l}|\Lost_Z|=|\Lost_l|$, because the $\Lost_Z$ of one round lie in the pairwise disjoint port sets $U_Z$ (Proposition~\ref{s2:propStructure}(iv)). So $|\mathcal O_{\mathrm T}|\le(\cFresh+169\epsA)n+169\sum_l|\Lost_l|$, with $\cFresh=169\cdot2$.
\item \emph{$\mathcal O_{\mathrm S}$.} By Lemma~\ref{s6:lemJSLC} and $\sum_l1/M_l\le2/\Dstar$ (Lemma~\ref{s2:lemTower}(b)),
\[
|\mathcal O_{\mathrm S}|\le\sum_{l\ge3}\Big(126\frac n{M_l}+1.5\sum_{(Y,l)\text{ giant}}m_{Y,l}\Big)\le\frac{252\,n}{\Dstar}+1.5\sum_{(Y,l)}m_{Y,l}.
\]
\item \emph{$\mathcal O_{\mathrm J}$.} $|\mathcal O_{\mathrm J}|=\sum_l|\Obj_l|$.
\end{itemize}
Adding the four bounds gives (c), since $\epsK+169\epsA+252/\Dstar=\varepsilon_{\mathrm M}(\Dstar)$. The final inequality is Theorem~\ref{s6:thmCONCL}(iv), which holds for every valid run and every designation.
\end{proof}

\section{The junction-vertex quotient, ultra hubs, and the induction}\label{s7:sec}

This section supplies the J-consumer of Definition~\ref{s6:defJconsumer}, bounds
the size of the quotient graphs it produces, and closes the argument by the
layered quotient induction.

\paragraph{Setting of this section.}
Throughout, $\Dstar$ satisfies \Gc{1}--\Gc{4}, $G$ is a graph with $n\ge\Nzero$
vertices and $d_1\ge\Dstar$, a valid $\HBtp$ run on $G$ is fixed
(Proposition~\ref{s2:propExists}), a designation $\delta$ is fixed
(Definition~\ref{s6:defDesign}), and there are no VX-parts
($\mathcal{V}=\emptyset$), so that every standalone pre-part is a MIX-part. The
rounds of the run are $1,\dots,R$; the round step below acts at the rounds
$3\le l\le R$ (rounds $1$ and $2$ have no classed ports). For a classed port $u$
of round $l$ (that is, $u\in Q_Z$ for some $Z\in\Std_l$) we write $Y(u)$ for its
class and $r(u)$ for the round of $Y(u)$; since $Y(u)\in\anc_l(u)$ we have
$r(u)\le l-2$ and $u\in V(Y(u))$. For an ancestor $Y$ of round $r$, $H_Y$ is its
expander ($X_Y$ if $Y$ is light, $X^0_Y$ if $Y$ is standalone), $s_Y$ is $s_r/2$
or $s_r$ respectively, and $p_Y=1/(2\klend(Y))$ (Definition~\ref{s3:defCOL}).
The functions $\epsOne$ and $\epsTwo$ defined below contain the function
$\epsA=31\eps/(\Cp\log\log\Dstar)$ of Lemma~\ref{s2:lemTower}(e), which depends
on $\Dstar$ only and satisfies $\sum_{l\le R}15.2\eps/\log P_l\le\epsA$; the
function $\epsX$ does not contain $\epsA$.

%% ----------------------------------------------------------------------
\subsection{Round-split pools and candidate junctions}

\begin{definition}[Round-split pool labels]\label{s7:defPool}
Every vertex $v$ of $G$ independently draws a \emph{pool label}
\[
\plab(v)\in\{\bot\}\cup\{(l,r):3\le l\le R,\ 1\le r\le l-2\},\qquad
\Prob\bigl(\plab(v)=(l,r)\bigr)=q_l\,\pi_{l,r},
\]
where $q_l\defeq M_l^{-2}$ and $\pi_{l,r}\defeq 2^{-(l-1-r)}$, and
$\Prob(\plab(v)=\bot)\defeq 1-\sum_{l,r}q_l\pi_{l,r}$. This is a probability
distribution: for each $l$ we have $\sum_{r=1}^{l-2}\pi_{l,r}=1-2^{-(l-2)}<1$,
and $\sum_l q_l\le\sum_l M_l^{-1}\le 2/\Dstar<1$ by
Lemma~\ref{s2:lemTower}(b). Put
\[
\Pool_{l,r}\defeq\plab^{-1}(l,r),\qquad
\Pool_l\defeq\bigcup_{r=1}^{l-2}\Pool_{l,r}\ \ \text{(a disjoint union)},
\]
so that $\Prob(v\in\Pool_l)=q_l(1-2^{-(l-2)})\le q_l$ for every $v$. The sets
$\Pool_l$ ($3\le l\le R$) are pairwise disjoint, because every vertex has one
label. For $w\in\Pool_l$ let $r(w)$ denote the unique $r$ with
$w\in\Pool_{l,r}$. The pool labels are stage (1d) of the randomness schedule.
They are independent of all other stage-1 data, in particular of the colourings
and labels of Definition~\ref{s3:defCOL}.
\deps{\ref{s2:lemTower}, \ref{s3:defCOL}}
\end{definition}

\begin{definition}[Candidates and JV-bad ports]\label{s7:defCand}
Let $u$ be a classed port of round $l\ge3$, with class $Y\defeq Y(u)$ of round
$r\defeq r(u)\le l-2$. Its \emph{candidate set} is
\[
\Cand_l(u)\defeq\{\,w\in\Pool_{l,r(u)}:\ uw\in\LJV_{Y,l}\,\}.
\]
Since $\LJV_{Y,l}\subseteq\Lend_Y\subseteq E(H_Y)$ and $H_Y$ is a graph on
$V(Y)$, automatically $\Cand_l(u)\subseteq N_{H_Y}(u)$ and
$N_{H_Y}(u)\subseteq V(Y)$. The
candidate sets are always taken in the class graph $H_{Y(u)}$ itself; this is
used in Lemma~MULT below. Put
\[
\Hcd_l\defeq\frac{\lambda_{l-2}^{95}}{8M_l^{2}}.
\]
The port $u$ is \emph{JV-bad} if
\[
|\Cand_l(u)|<\tfrac12\,\Ex|\Cand_l(u)|,
\]
and \emph{JV-good} otherwise. Here $\Ex$ is the unconditional expectation over
stage~1, which is a number determined by the run and $\delta$. The definition
applies to every classed port, whether or not it lies in $\Lost_l$. JV-badness
is a deterministic function of stage~1; it is one of the stage-2 statuses. The
number $\Hcd_l$ is not a threshold of the definition: it is a lower bound for
the threshold, since $\tfrac12\Ex|\Cand_l(u)|\ge\Hcd_l$ for every classed port
$u$ of round $l$ (Lemma~\ref{s7:lemCand}(ii) below). Consequently every JV-good
port has $|\Cand_l(u)|\ge\Hcd_l$ (Lemma~\ref{s7:lemCand}(iv)).
\deps{\ref{s7:defPool}, \ref{s3:defCOL}, \ref{s6:defDesign}}
\end{definition}

\begin{lemma}[Candidate counts]\label{s7:lemCand}
Let $u$ be a classed port of round $l\ge3$ with class $Y$ of round $r$. Then:
\begin{enumerate}
\item[(i)] $|\Cand_l(u)|$ is a sum of independent Bernoulli variables, one for
each $w\in N_{H_Y}(u)$, each with success probability $q_l\pi_{l,r}p_Y$;
\item[(ii)] its mean satisfies
\[
\Ex|\Cand_l(u)|=q_l\,\pi_{l,r}\,p_Y\,\deg_{H_Y}(u)
\;\ge\;\frac{\lambda_r^{96}\,2^{-(l-r)}}{M_l^{2}}
\;\ge\;2\Hcd_l ;
\]
\item[(iii)] $\Prob(u\text{ is JV-bad})\le\exp(-\Hcd_l/4)$;
\item[(iv)] if $u$ is JV-good then $|\Cand_l(u)|\ge\Hcd_l\ge2^{10}M_l^{10}$.
\end{enumerate}
\deps{\ref{s7:defCand}, \ref{s7:defPool}, \ref{s3:defCOL}, \ref{s3:lemCOLJV}, \ref{s2:propStructure}, \ref{s2:lemTower}, \ref{s1:citChernoffGen}}
\end{lemma}

\begin{proof}
(i) By definition,
\[
|\Cand_l(u)|=\sum_{w\in N_{H_Y}(u)}\ind{\plab(w)=(l,r)}\cdot\ind{uw\in\LJV_{Y,l}} .
\]
The pool labels are independent across vertices. The colours of the edges of
$H_Y$ are independent across edges (Definition~\ref{s3:defCOL}), and the edges
$uw$ for distinct $w$ are distinct. The two families are independent of each
other (Definition~\ref{s7:defPool}). So the summands for distinct $w$ are
functions of disjoint sets of independent variables, hence independent. For a
fixed $w$ the two indicators are independent, with
$\Prob(\plab(w)=(l,r))=q_l\pi_{l,r}$. Moreover $\Prob(uw\in\LJV_{Y,l})=p_Y$: the
edge $uw\in E(H_Y)$ lies in $\Lend_Y$ with probability $1/2$, and it then
receives the colour $l\in\IJV$ with probability $1/\klend(Y)$; note that
$r+2\le l\le R$, so $l$ is an index of $\IJV$.

(ii) The first equality follows from (i). Every $H_Y$ has minimum degree larger
than $s_Y\ge s_r/2$, and $s_r\ge\lambda_r^{100}$
(Proposition~\ref{s2:propStructure}(i)). Also $p_Y\ge\lambda_r^{-4}$
(Lemma~\ref{s3:lemCOLJV}). Hence
\[
\Ex|\Cand_l(u)|\ \ge\ M_l^{-2}\,2^{-(l-1-r)}\,\lambda_r^{-4}\,\frac{\lambda_r^{100}}{2}
\;=\;\frac{\lambda_r^{96}\,2^{-(l-r)}}{M_l^{2}} .
\]
Since $2\Hcd_l=\lambda_{l-2}^{95}/(4M_l^2)$, the second inequality is equivalent
to
\begin{equation}\label{s7:eqCandExp}
\lambda_r^{96}\ \ge\ 2^{\,l-r-2}\,\lambda_{l-2}^{95}.
\end{equation}
By Lemma~\ref{s2:lemTower}(a), $\lambda_r\ge\lambda_{l-2}$ because $r\le l-2$.
By Lemma~\ref{s2:lemTower}(d),
$\lambda_r\ge2^{2\logs d_r+2}\ge2^{R-r}\ge2^{l-r}\ge2^{l-r-2}$. Multiplying
$\lambda_r^{95}\ge\lambda_{l-2}^{95}$ by $\lambda_r\ge2^{l-r-2}$ gives
\eqref{s7:eqCandExp}. (The chain of exponents is
$2\logs d_r+2\ge R-r\ge l-r$; only the exponent $l-r-2$ is needed.)

(iii) Put $\mu\defeq\Ex|\Cand_l(u)|$. By definition, $u$ is JV-bad exactly
when $|\Cand_l(u)|<\mu/2$. By (i) and the lower-tail Chernoff
bound for sums of independent indicators, with deviation parameter $1/2$
(Cited result~\ref{s1:citChernoffGen}),
$\Prob(|\Cand_l(u)|<\mu/2)\le\exp(-\mu/8)\le\exp(-\Hcd_l/4)$, where the last
step uses $\mu/2\ge\Hcd_l$ from (ii).

(iv) If $u$ is JV-good then $|\Cand_l(u)|\ge\tfrac12\Ex|\Cand_l(u)|\ge\Hcd_l$,
by the definition of JV-good and (ii); this is the first inequality. The second
is the line
``$\Hcd_l\ge2^{10}M_l^{10}$'' of the table of Lemma~\ref{s3:lemCOLJV}. It also
follows directly from $M_l\le\lambda_{l-2}^{1.6}$
(Lemma~\ref{s2:lemTower}(b)): then
$8M_l^2\cdot2^{10}M_l^{10}=2^{13}M_l^{12}\le2^{13}\lambda_{l-2}^{19.2}\le\lambda_{l-2}^{95}$,
because $\lambda_{l-2}\ge M_l^{1/1.6}\ge2^{25}$.
\end{proof}

%% ----------------------------------------------------------------------
\subsection{The randomness schedule and the round step}

\begin{definition}[The randomness schedule]\label{s7:defSchedule}
All random choices of the construction are made in the following order,
summarized in Table~\ref{s7:tabSchedule}.
\begin{itemize}
\item \emph{Stage 1} consists of four mutually independent families:
(1a) the COL-JV colourings of Definition~\ref{s3:defCOL}(i)--(iii);
(1b) the vertex choices and sublabels of Definition~\ref{s5:defZones};
(1c) the JS labels of Definition~\ref{s3:defCOL}(iv);
(1d) the pool labels of Definition~\ref{s7:defPool}.
\item \emph{Selection} (1e), which is not a random family: a stage-1 outcome is
then fixed in the event $\{\Xp\le3\,\Ex\Xp\}$. Here
$\Xp\ge0$ is the stage-1 functional introduced later in this section, a
deterministic function of the run, $\delta$ and stage~1. After (1e), stage~1 is
a fixed outcome: no later step of the construction uses the law of stage~1 or
conditions on the event of (1e); the numbers $\Ex\Xp$ and $\Ex|\Cand_l(u)|$
that appear later are unconditional expectations, fixed in advance.
\item \emph{Stage 2} is deterministic given stage~1: demotion, parent-bad and
$\Bad$ (Definition~\ref{s5:defStages}); lend-bad, $\Lost$, $\Ret$ and $\Qs$
(Definition~\ref{s6:defLending}); JV-bad (Definition~\ref{s7:defCand}).
\item \emph{Stage 3} is drawn part by part: the TPV labels of every standalone
pre-part, the P-vortex labels of every non-demoted light part, and the VX$^+$
labels of every demoted light part. Given stage~1, the label families of distinct
parts are independent. Each part's labels are fixed inside that part's good event.
\item \emph{Rounds} $l=R,R-1,\dots,3$. The round randomness $\xi_l$ consists of
the following mutually independent variables:
a uniformly random bijection $\eta_h:[4M_l]\to[4M_l]$ for every vertex $h$;
a uniformly random $3$-subset $\zeta_{h,u}\subseteq[4M_l]$ for every ordered pair
$(h,u)$ of distinct vertices;
and a uniformly random linear order $\prec_u$ of $V(G)$ for every vertex $u$.
The dependence of $\eta_h,\zeta_{h,u},\prec_u$ on $l$ is suppressed. The draw
$\xi_l$ is fresh: it is independent of stage~1, of stage~3 and of all $\xi_{l'}$
with $l'>l$.
\end{itemize}
The \emph{past} $\Past_l$ at round $l$ is the fixed outcome of stage~1, of
stage~3 and of the draws $\xi_{l'}$ ($l'>l$), together with everything these
determine: in particular the decompositions fixed at the rounds $l'>l$ and the
set $J_l$. It is a fixed outcome, not a random $\sigma$-algebra. We write
$\Ex[\,\cdot\mid\Past_l]$ and $\Prob(\,\cdot\mid\Past_l)$ for expectation and
probability over $\xi_l$ with $\Past_l$ fixed. The fixed rules of
Construction~\ref{s7:consRound} read $\xi_l$ only as stated there. At each round, $\xi_l$ is chosen in
the event of step (f) of the round step. All draws in stage~1, stage~3 and the
rounds are mutually independent. The outcomes actually used are chosen inside
events (stage~1 in $\{\Xp\le3\,\Ex\Xp\}$, stage~3 in the good events, each
$\xi_l$ in the event of step (f)); after each choice only deterministic
properties of the chosen outcome are used, and every later probability is over
fresh variables with the past fixed.
\deps{\ref{s3:defCOL}, \ref{s5:defZones}, \ref{s7:defPool}, \ref{s6:defLending}, \ref{s5:defStages}, \ref{s7:defCand}, \ref{s4:lemTPV}, \ref{s4:lemPV}, \ref{s4:thmVXp}, \ref{s5:lemChild}, \ref{s5:lemDemoted}}
\end{definition}

\begin{table}[ht]
\centering
\small
\begin{tabular}{|p{1.5cm}|p{5.9cm}|p{3.6cm}|p{3.8cm}|}
\hline
\textbf{Stage} & \textbf{Drawn or decided} & \textbf{Depends on} & \textbf{How it is fixed}\\
\hline
1a & COL-JV colourings of every ancestor & independent per ancestor and per edge & part of the stage-1 outcome\\
1b & vertex choices and sublabels (zones) & independent per vertex & part of the stage-1 outcome\\
1c & JS labels $\lab_{Y,l}(y)$ & independent & part of the stage-1 outcome\\
1d & pool labels $\plab(v)$ & independent per vertex & part of the stage-1 outcome\\
1e & choice of the stage-1 outcome & the functional $\Xp$ & $\Xp\le3\,\Ex\Xp$ (probability $\ge2/3$)\\
\hline
2 & demotion, parent-bad, lend-bad, $\Lost$, $\Ret$, $\Qs$, JV-bad & deterministic in stage~1 & determined\\
\hline
3 & TPV, P-vortex and VX$^+$ labels & independent per part & inside each part's good event (probability $\ge1/2-o(1)$)\\
\hline
$l=R,\dots,3$ & $\xi_l=(\eta_h,\zeta_{h,u},\prec_u)$: HUB lists and injections & fresh draw, $\Past_l$ fixed & event of step (f) (probability $\ge1/2$)\\
\hline
\end{tabular}
\caption{The randomness schedule (Definition~\ref{s7:defSchedule}).}
\label{s7:tabSchedule}
\end{table}

\begin{construction}[The JV$^{+*}$ round step at round $l\ge3$; the J-consumer]\label{s7:consRound}
Fix $3\le l\le R$ and the past $\Past_l$. The input is the set
$J_l=\Jlost_l\cup\Jhub_l\cup\Jfr_l\cup\Jpar_l$ produced by JS-LC at round $l$,
with the properties \Jp{1}--\Jp{3} and (i)--(v) of Lemma~\ref{s6:lemJplus}.

\emph{Fixed rules.} The rules and orders called fixed in steps (b)--(e), (g)
and (h) below are deterministic functions, chosen once and for all before any
randomness is drawn, with the following arguments.
\begin{itemize}
\item \emph{Functions of $\Past_l$ alone:} the pairing of the live $\Jfr$ items
at a fresh centre and the choice of its unpaired fresh leg, and the order in
which the PAR objects are coloured, in (b); the split of the live items of a
non-ultra hub into groups, together with the numbering of the groups, in (c).
In particular these rules do not read $\xi_l$.
\item \emph{Functions of $\Past_l$ and of the lists}, that is, of the variables
$\eta_h$ and $\zeta_{h,u}$ of $\xi_l$: the choice of the system of distinct
representatives in (d); the processing order and the order of $V(G)$ in (e1);
and the order $e_1,\dots,e_q$ of $E'(u)$ in (e2).
\item None of the rules of steps (b)--(e) reads the orders $\prec_u$ of
$\xi_l$; within steps (a)--(f), the orders $\prec_u$ of $\xi_l$ enter only
through the assignment $e_i\mapsto w_i$ of (e2) (and through what is computed
from it). Through $\Past_l$ the rules may depend on the draws
$\xi_{l'}$ of the rounds $l'>l$, orders included, and in general they do,
since $J_l$ does.
\item The rank order in (g) and the rule choosing $\Dec_l$ in (h) are not
restricted: they are arbitrary deterministic functions of $\Past_l$ and
$\xi_l$. They necessarily depend on $\xi_l$, since the layers and $Q_l$ do.
\end{itemize}
Rules with these arguments exist, for instance lexicographically first choices
for fixed orders of the finite sets involved (the greedy colouring of (b)
succeeds for every order, by part (i) of the next lemma, and a split into groups
as in (c) always exists). Every statement of this section about the
round step holds for every choice of the fixed rules that obeys these
restrictions. The restrictions are used in the claims of steps (c) and (e2), in
Lemma~\ref{s7:lemWellDef}(iii) (the grouping does not read $\eta_h$), and in
Lemmas~\ref{s7:lemCC}(i), \ref{s7:lemUltra}(i) and \ref{s7:lemPay}(c) (nothing
before the assignment $e_i\mapsto w_i$ reads the orders $\prec_u$), as well as
in the opening paragraph of the subsection ``Quotient size and payments'' and
in the proof of Lemma~\ref{s7:lemUHsplit}(ii) (the copies of non-ultra hubs are
deterministic given the lists).

\emph{Items.} The \emph{items} are the edges of $J_l\setminus\Jlost_l$. The
\emph{vertices} of an item are its port end or ends, which lie in $\Qs_Z$, and
its centre. The centre is a hub $h\in D_l$ for a $\Jhub$ item and a fresh centre
$x\in F_Z$ for a $\Jfr$ item. A $\Jpar$ item $u$--$v$ has two port ends and no
centre. A $\Jhub$ item is written $(h,u)$, with $h$ its hub and $u$ its port.

\smallskip\noindent
(a) \emph{Payments bounded by stage-1 weights.} The following are \emph{paid},
that is, output as single edges:
(a1) every edge of $\Jlost_l$;
(a2) every item with a vertex in $\Pool_l$;
(a3) every item with a JV-bad port end.
The remaining items are \emph{live}. A vertex is \emph{live} if it is a vertex of
a live item, and $\Live$ denotes the set of live vertices. By (a2),
$\Live\cap\Pool_l=\emptyset$. The sets paid in (a1)--(a3) depend on $J_l$, hence
on the past; only their sizes are bounded by stage-1 functionals
(this is proved in the payment lemma below).

\smallskip\noindent
(b) \emph{PAR objects.} Every live $\Jpar$ item $u$--$v$ is a \emph{PAR object}
with \emph{ends} at $u$ and at $v$; here $u,v\in\Qs_Z$ and $Y(u)\ne Y(v)$. At each
fresh centre $x$, the live $\Jfr$ items at $x$ are paired, by a fixed rule, into
\emph{cherries} $u$--$x$--$u'$. Here $u\ne u'$ because $G$ is simple, so
distinct items at $x$ have distinct port ends. A cherry is a PAR object with ends
at $u$ and $u'$ and \emph{middle} $x$. If the number of live $\Jfr$ items at $x$
is odd, one of them, chosen by a fixed rule, is paid: it is the \emph{unpaired
fresh leg} of $x$. The PAR objects are then coloured greedily, in a fixed order,
with colours from $[3M_l]$, so that two PAR objects sharing a port, or sharing a
middle, receive different colours.

\smallskip\noindent
(c) \emph{HUB lists.} For a hub $h\in D_l$ let $\clive_h$ be its number of live
$\Jhub$ items. Put
\[
\thult_l\defeq\bigl\lfloor M_l\Hcd_l/7\bigr\rfloor,\qquad \KHUB_l\defeq4M_l .
\]
The hub $h$ is \emph{ultra} if $\clive_h>\thult_l$, and \emph{non-ultra}
otherwise. For a non-ultra $h$ with $\clive_h\ge1$, its live items are split, by
a fixed rule, into
\[
k_h\defeq\bigl\lceil 8\clive_h/\Hcd_l\bigr\rceil
\]
\emph{groups} of at most $\lceil\Hcd_l/8\rceil$ items each. This is possible
because $k_h\lceil\Hcd_l/8\rceil\ge\clive_h$. The $i$-th group receives the
\emph{list} $\{\eta_h(3i-2),\eta_h(3i-1),\eta_h(3i)\}\subseteq[4M_l]$; this needs
$3k_h\le4M_l$, which is part (ii) of the next lemma. The lists of the
groups of $h$ are pairwise disjoint, and, since the groups and their numbering
are functions of $\Past_l$ while $\eta_h$ is independent of $\Past_l$, they form
a uniformly random sequence of
$k_h$ pairwise disjoint $3$-subsets of $[4M_l]$. Lists of distinct hubs are
independent. For an ultra $h$, every live item $(h,u)$ is its own group, with
list $\zeta_{h,u}$.

\smallskip\noindent
(d) \emph{SDR.} Every port $u$ carrying at least one live hub item chooses, by a
fixed rule, pairwise distinct colours for its live hub items, each item's colour
taken from the list of its group (a system of distinct representatives). If no
such choice exists, all live hub items of $u$ are paid (an \emph{SDR failure} at
$u$). The items that receive a colour are the \emph{coloured hub items}. In
particular, every port carries at most one coloured hub item of each HUB colour.

\smallskip\noindent
(e) \emph{Junctions.} Initially $\Used(u)\defeq\emptyset$ for every port $u$ and
$\Used_\kappa(h)\defeq\emptyset$ for every hub $h$ and every $\kappa\in[4M_l]$.
\begin{itemize}
\item[(e1)] The coloured hub items of non-ultra hubs are processed in a fixed
order. The item $(h,u)$ of colour $\kappa$ receives as its \emph{junction} the
first vertex $w$, in a fixed order of $V(G)$, of
$\Cand_l(u)\setminus\Used(u)\setminus\Used_\kappa(h)$. Then $w$ is added to
$\Used(u)$ and to $\Used_\kappa(h)$.
\item[(e2)] For every port $u$ let $E'(u)$ be the set of the ends at $u$ of the
PAR objects and of the coloured hub items $(h,u)$ with $h$ ultra, listed in a
fixed order $e_1,\dots,e_q$. Let $w_1\prec_u w_2\prec_u\cdots$ be the elements of
$\Cand_l(u)\setminus\Used(u)$ in the order $\prec_u$; the end $e_i$ receives the
junction $w_i$. This is well defined by part (v) of the next lemma. With
$\Past_l$ and the lists fixed, the map $e_i\mapsto w_i$ is a uniformly random
injection of $E'(u)$ into $\Cand_l(u)\setminus\Used(u)$, and these injections are
independent over ports: by the restrictions on the fixed rules, $\Used(u)$, the
set $E'(u)$ and its order $e_1,\dots,e_q$ are determined by $\Past_l$ and the
lists, and the orders $\prec_u$ are independent of both. No payment is made for
coincidences of junctions at ultra hubs.
\item[(e3)] \emph{Loops.} A PAR object whose two ends received the same junction
is paid. A looped $\Jpar$ item costs one single edge; a looped cherry costs its
two edges.
\end{itemize}
Every end of a coloured hub item or of an unpaid PAR object now has a junction
$w$. It is a candidate of the end's port $u$, so $w\in\Cand_l(u)\subseteq\Pool_l$
and $uw\in\LJV_{Y(u),l}$. The edge $uw$ is the \emph{junction edge} of that end.

\smallskip\noindent
(f) \emph{Markov choice.} Let $\payrd_l$ be the number of edges paid in (d) and
(e3), and let $\copies_l\defeq|V(Q_l)|$, with $Q_l$ as in (g). Both are
functions of $\Past_l$ and $\xi_l$. The round randomness $\xi_l$ is chosen in the
event
\[
\bigl\{\copies_l\le4\,\Ex[\copies_l\mid\Past_l]\bigr\}\cap
\bigl\{\payrd_l\le4\,\Ex[\payrd_l\mid\Past_l]\bigr\}.
\]
By Markov's inequality in the form~(b) of Cited result~\ref{s1:citMarkov},
taken conditionally as in its item~(c) ($\xi_l$ is independent of the variables
whose outcome $\Past_l$ fixes, Definition~\ref{s7:defSchedule}), this event has
probability at least $1/2$ given $\Past_l$ (each of the two events fails with
probability at most $1/4$).

\smallskip\noindent
(g) \emph{Layers, sub-layers and the quotient.} For $\kappa\in[3M_l]$ the
\emph{PAR layer} $\BP_\kappa$ is the multigraph on the vertex set
$\{[w]:w\in\Pool_l\}$ of formal copies. It has one edge $[w_1][w_2]$ for every
unpaid PAR object of colour $\kappa$ whose ends have the junctions $w_1$ and
$w_2$; here $w_1\ne w_2$ by (e3). For $\kappa\in[4M_l]$ the \emph{HUB layer}
$\BH_\kappa$ is the bipartite multigraph with sides $D_l$ and
$\{[w]:w\in\Pool_l\}$. It has one edge $h[w]$ for every coloured hub item $(h,u)$
of colour $\kappa$ whose junction is $w$. \emph{Rank split:} in every layer, the
edges joining the same two vertices are listed in a fixed order, and the $i$-th
edge of such a list has \emph{rank} $i$. For a layer $B$ and $\iota\ge1$, the
\emph{sub-layer} $B^{(\iota)}$ is the spanning subgraph of $B$ formed by its edges
of rank $\iota$. The \emph{quotient} $Q_l$ is the vertex-disjoint union of all
sub-layers $B^{\mathrm{P},(\iota)}_\kappa$ ($\kappa\in[3M_l]$, $\iota\ge1$) and
$B^{\mathrm{H},(\iota)}_\kappa$ ($\kappa\in[4M_l]$, $\iota\ge1$), with the
isolated vertices of each sub-layer deleted. Vertices of distinct sub-layers are
distinct vertices of $Q_l$; hence every cycle of $Q_l$ lies in one sub-layer.

\smallskip\noindent
(h) \emph{Decompose and lift.} Let $\Dec_l$ be a decomposition of $E(Q_l)$ into
$\f(Q_l)$ objects, chosen by a fixed deterministic rule (for instance the
lexicographically first one for a fixed order of $E(Q_l)$). Each member of
$\Dec_l$ is \emph{lifted} to objects of $G$ as follows.
\begin{itemize}
\item A single edge of a PAR sub-layer is a PAR object. It is output as its one
($\Jpar$) or two (cherry) J-edges, each as a single edge. A single edge $h[w]$ of
a HUB sub-layer is a coloured hub item $(h,u)$, and it is output as the single
edge $hu$.
\item A cycle $[w_0][w_1]\cdots[w_{m-1}]$ of a PAR sub-layer, with $m\ge3$:
indices are taken mod $m$, and the edge $[w_i][w_{i+1}]$ is a PAR object $o_i$.
Let $p_i$ be the end of $o_i$ with junction $w_i$ and $q_i$ its end with junction
$w_{i+1}$; they are well defined because $w_i\ne w_{i+1}$. The lift is the closed
walk
\[
w_0\,p_0\,(x_0)\,q_0\,w_1\,p_1\,(x_1)\,q_1\,w_2\ \cdots\ q_{m-1}\,w_0 ,
\]
where $x_i$ is the middle of $o_i$ if $o_i$ is a cherry, and is omitted if $o_i$
is a $\Jpar$ item.
\item A cycle $h_0[w_0]h_1[w_1]\cdots h_{m-1}[w_{m-1}]h_0$ of a HUB sub-layer,
with $m\ge2$: the edge $h_i[w_i]$ is a coloured hub item $(h_i,u_i)$ with
junction $w_i$, and the edge $[w_i]h_{i+1}$ is a coloured hub item
$(h_{i+1},u'_i)$ with junction $w_i$ (indices mod $m$). The lift is the closed
walk
\[
h_0\,u_0\,w_0\,u'_0\,h_1\,u_1\,w_1\,u'_1\,h_2\ \cdots\ u'_{m-1}\,h_0 .
\]
\end{itemize}
Put $\Obj_l\defeq{}$(the single edges paid in (a), (b), (d) and (e3)) together
with (the lifted objects of all members of $\Dec_l$), and let $\LentJV_l$ be the
set of junction edges that lie on lifted cycles. A junction edge assigned in (e)
that does not lie on a lifted cycle is not used by the round step; it stays with
its owner as junk.
\deps{\ref{s6:lemJplus}, \ref{s6:defJconsumer}, \ref{s7:defCand}, \ref{s7:defPool}, \ref{s7:defSchedule}, \ref{s1:citMarkov}}
\end{construction}

\begin{lemma}[The round step is well defined]\label{s7:lemWellDef}
For every past $\Past_l$ ($3\le l\le R$) the following hold.
\begin{enumerate}
\item[(i)] \emph{PAR palette.} Every PAR object shares a port or a middle with
fewer than $3M_l$ other PAR objects, so the greedy colouring of (b) succeeds.
Consequently every port carries at most one PAR object of each colour, and
every fresh centre is the middle of at most one cherry of each colour.
\item[(ii)] \emph{Disjoint lists.} Every non-ultra hub $h$ has
$k_h\le\lceil8M_l/7\rceil$ and $3k_h\le4M_l$.
\item[(iii)] \emph{Cuckoo SDR.} For every port $u$,
$\Prob(\text{SDR failure at }u\mid\Past_l)\le(\KHUB_l)^{-4}=(4M_l)^{-4}$.
\item[(iv)] \emph{Greedy step.} Step (e1) always finds a junction. Two coloured
items of the same non-ultra hub $h$ with the same colour $\kappa$ receive
distinct junctions.
\item[(v)] \emph{Injections.} At every port $u$ carrying a live item,
$|\Cand_l(u)\setminus\Used(u)|\ge\Hcd_l-(M_l-1)\ge\Hcd_l/2\ge M_l>|E'(u)|$, so
step (e2) is well defined.
\end{enumerate}
\deps{\ref{s7:consRound}, \ref{s6:lemJplus}, \ref{s7:lemCand}, \ref{s3:lemCOLJV}, \ref{s1:citHall}, \ref{s2:defHBtp}}
\end{lemma}

\begin{proof}
Throughout, $M_l\ge2^{40}$ by the definition of $M_l$, and
$\Hcd_l\ge2^{10}M_l^{10}$ by Lemma~\ref{s7:lemCand}(iv) (equivalently, by the
table of Lemma~\ref{s3:lemCOLJV}). By \Jp{2} of Lemma~\ref{s6:lemJplus}, every
port carries at most $M_l-1$ J-edges, hence at most $M_l-1$ live items. By
Lemma~\ref{s6:lemJplus}(ii), every fresh centre carries at most $M_l-1$ $\Jfr$
edges, hence is the middle of at most $(M_l-1)/2$ cherries.

(i) A PAR object has at most two ports and at most one middle. Each of its ports
carries at most $M_l-2$ further live items, and each further PAR object at that
port uses one of them. Its middle, if any, is the middle of at most
$(M_l-1)/2-1$ further cherries. So the object conflicts with at most
$2(M_l-2)+(M_l-1)/2<3M_l$ others. When it is coloured, fewer than $3M_l$ colours
are therefore blocked, and a colour of $[3M_l]$ remains. The consequences are
the defining property of the colouring.

(ii) If $h$ is non-ultra then $\clive_h\le\thult_l\le M_l\Hcd_l/7$, so
$k_h=\lceil8\clive_h/\Hcd_l\rceil\le\lceil8M_l/7\rceil$. Hence
$3k_h\le3(8M_l/7+1)=24M_l/7+3\le4M_l$, because $M_l\ge21/4$.

(iii) Fix $u$ and let $(h_1,u),\dots,(h_m,u)$ be its live hub items, so that
$m\le M_l-1\le K/4$ with $K\defeq\KHUB_l=4M_l$. These items are distinct edges at
$u$ and $G$ is simple, so the hubs $h_1,\dots,h_m$ are distinct. With $\Past_l$
fixed, it is determined which hubs are ultra, how the items are grouped, and
which group each item belongs to (the grouping rule of step (c) reads $\Past_l$
only, by the restrictions on the fixed rules in
Construction~\ref{s7:consRound}). The list of $(h_i,u)$ is a function of
$\eta_{h_i}$ if $h_i$ is non-ultra, and equals $\zeta_{h_i,u}$ if $h_i$ is ultra.
Since the $h_i$ are distinct, the $m$ lists are independent. Each list is a
uniformly random $3$-subset of $[K]$; for a non-ultra hub this holds because the
sequence of its lists is a uniformly random sequence of pairwise disjoint
$3$-subsets, so each single list is uniform. By Hall's theorem (Cited
result~\ref{s1:citHall}, applied with the $m$ items as the index set and, for
each item $(h_i,u)$, its list as the set of that item), an SDR fails to exist
only if some set of $s$ items has at most $s-1$ colours in the union of their
lists. Since every list has $3$ elements, this forces $s\ge4$. The union is
then contained in some $(s-1)$-subset of $[K]$. The union bound over $s$, over
the $s$-sets of items and over the $(s-1)$-subsets of $[K]$ gives
\[
\Prob(\text{SDR failure at }u\mid\Past_l)\ \le\ \sum_{s=4}^{m}T_s,\qquad
T_s\defeq\binom{m}{s}\binom{K}{s-1}\Bigl(\binom{s-1}{3}\Big/\binom{K}{3}\Bigr)^{s}.
\]
We use $\binom{m}{s}\le(em/s)^s\le(eK/(4s))^s$,
$\binom{K}{s-1}\le(eK/(s-1))^{s-1}$, and
\[
\binom{s-1}{3}\Big/\binom{K}{3}=\frac{(s-1)(s-2)(s-3)}{K(K-1)(K-2)}\le\Bigl(\frac{s-1}{K}\Bigr)^3 ,
\]
which holds because $(s-1-j)/(K-j)\le(s-1)/K$ for $j=1,2$ and $s\le K+1$.
Multiplying, and using $(s-1)^{2s+1}\le s^{2s+1}$,
\[
T_s\ \le\ e^{2s-1}\,4^{-s}\,s^{-s}\,(s-1)^{2s+1}\,K^{-s-1}\ \le\
a_s\defeq\frac1e\cdot\frac sK\cdot\Bigl(\frac{e^2s}{4K}\Bigr)^{s}.
\]
First, $a_4=4e^7K^{-5}<4400\,K^{-5}$. Second, for $s\ge4$,
\[
\frac{a_{s+1}}{a_s}=\frac{s+1}{s}\cdot\frac{e^2(s+1)}{4K}\cdot\Bigl(\frac{s+1}{s}\Bigr)^{s}
\le\frac54\cdot\frac{e^2(s+1)}{4K}\cdot e=\frac{5e^3}{16}\cdot\frac{s+1}{K},
\]
which is at most $\frac{5e^3}{16}(\frac1{13}+\frac1K)<0.49$ when $4\le s<K/13$ and
$K\ge10^4$. Hence $a_s\le2^{4-s}a_4$ for $4\le s\le\lceil K/13\rceil$, and the sum
of these terms is at most $2a_4<8800\,K^{-5}$. Third, for $K/13<s\le m\le K/4$ we
have $s/K\le1/4$, so $a_s\le\frac1{4e}(e^2/16)^s\le0.47^{s}$, and the sum of
these terms is at most $0.47^{K/13}/0.53\le2\cdot0.47^{K/13}$. Altogether
\[
\Prob(\text{SDR failure at }u\mid\Past_l)\ \le\ 8800\,K^{-5}+2\cdot0.47^{K/13}\ \le\ K^{-4}
\qquad(K\ge10^4),
\]
since $8800K^{-5}\le0.88K^{-4}$ and $2\cdot0.47^{K/13}\le0.1K^{-4}$ for
$K\ge10^4$. Here $K=4M_l\ge2^{42}$.

(iv) Items reach (e1) only if they are live, so their ports are JV-good and
$|\Cand_l(u)|\ge\Hcd_l$ (Lemma~\ref{s7:lemCand}(iv)). When the item $(h,u)$ of
colour $\kappa$ is processed, $\Used(u)$ contains one junction for each earlier
item at $u$, so $|\Used(u)|\le M_l-2$. The lists of the groups of $h$ are pairwise
disjoint, so all coloured items of $h$ with colour $\kappa$ lie in the one group
whose list contains $\kappa$. That group has at most $\lceil\Hcd_l/8\rceil$
items, so $|\Used_\kappa(h)|\le\lceil\Hcd_l/8\rceil-1<\Hcd_l/8$. Hence at least
$\Hcd_l-(M_l-2)-\Hcd_l/8>0$ choices remain. The second statement holds because
each junction chosen for $h$ in colour $\kappa$ is added to $\Used_\kappa(h)$
and is excluded afterwards.

(v) After (e1), $\Used(u)$ consists of one junction for each coloured item of a
non-ultra hub at $u$. The elements of $E'(u)$ correspond to further live items at
$u$: a PAR end at $u$ corresponds to the $\Jpar$ item or the leg of the cherry at
$u$, and an ultra hub end is itself an item. So
$|\Used(u)|+|E'(u)|\le M_l-1$. Every port carrying a live item is JV-good, so
$|\Cand_l(u)\setminus\Used(u)|\ge\Hcd_l-(M_l-1)\ge\Hcd_l/2\ge M_l>|E'(u)|$,
using $\Hcd_l\ge2^{10}M_l^{10}\ge2M_l$.

\emph{Remark on numerics (not used).} An exact evaluation of the union bound
$K^5\sum_sT_s$ with $m=M_l-1$ has its maximum, about $0.0502$, near $M_l=19$;
it tends to $9/256$ as $M_l\to\infty$, which is the limit of $K^5T_4$, where
$T_4=\binom{m}{4}/\binom{K}{3}^3\sim9m^4/K^9$. None of these numerical values is used: only the analytic bound (iii) is
used, and it holds since $M_l\ge2^{40}$.
\end{proof}

%% ----------------------------------------------------------------------
\subsection{The quotient is simple, and every decomposition lifts}

For a HUB colour $\kappa$, a hub $h$ and $w\in\Pool_l$, let $m_\kappa(h,w)$ be
the number of edges $h[w]$ of $\BH_\kappa$, that is, the number of coloured hub
items $(h,u)$ of colour $\kappa$ with junction $w$.

\begin{lemma}[Lemma MULT]\label{s7:lemMULT}
Let $3\le l\le R$, $\kappa\in[4M_l]$, let $h$ be a hub and let
$w\in\Pool_{l,r}$. Then
\[
m_\kappa(h,w)\ \le\ \mult_r(w)\ \le\ \mu_r(w).
\]
Moreover, the number of sub-layers $B^{\mathrm{H},(\iota)}_\kappa$ ($\iota\ge1$)
in which $[w]$ is not isolated equals $\max_{h}m_\kappa(h,w)$, and the number of
those in which $h$ is not isolated equals $\max_{w}m_\kappa(h,w)$.
\deps{\ref{s7:consRound}, \ref{s6:lemJplus}, \ref{s7:defCand}, \ref{s7:defPool}, \ref{s2:propOV}}
\end{lemma}

\begin{proof}
Let $(h,u_1),\dots,(h,u_m)$ be the coloured hub items of colour $\kappa$ at $h$
with junction $w$, so $m=m_\kappa(h,w)$. By step (e) of
Construction~\ref{s7:consRound}, the junction of $(h,u_i)$ lies in
$\Cand_l(u_i)\subseteq\Pool_{l,r(u_i)}$ (Definition~\ref{s7:defCand}). Every
vertex has a single pool label (Definition~\ref{s7:defPool}), and
$w\in\Pool_{l,r}$, so $r(u_i)=r$ for all $i$. Moreover
$w\in\Cand_l(u_i)\subseteq N_{H_{Y(u_i)}}(u_i)\subseteq V(Y(u_i))$; this is where
it matters that candidates are taken in the class graph $H_{Y(u_i)}$ itself. So
each $Y(u_i)$ is an ancestor of round $r$ whose vertex set contains $w$. The
items $(h,u_i)$ are distinct $\Jhub$ edges at $h$ of round $l$. By \Jp{1} of
Lemma~\ref{s6:lemJplus}, for each class $Y$ there is at most one $\Jhub$ edge of
class $Y$ at $h$ in round $l$, aggregated over all MIX-parts. So the classes
$Y(u_1),\dots,Y(u_m)$ are pairwise distinct. Hence $m\le\mult_r(w)$, the number
of ancestors of round $r$ containing $w$, and $\mult_r(w)\le\mu_r(w)$ by
Proposition~\ref{s2:propOV}.

For the second statement: $[w]$ is not isolated in $B^{\mathrm{H},(\iota)}_\kappa$
if and only if some edge of rank $\iota$ is incident to $[w]$. By the definition
of ranks, this holds if and only if some pair $\{h,[w]\}$ carries at least
$\iota$ edges, that is, $\max_h m_\kappa(h,w)\ge\iota$. The statement for $h$ is
proved in the same way.
\end{proof}

\begin{lemma}[$Q_l$ is simple]\label{s7:lemSimple}
Every sub-layer of Construction~\ref{s7:consRound}(g) is a simple graph: it has
no parallel edges, PAR sub-layers have no loops, and HUB sub-layers are
bipartite between hubs and junction copies $[w]$, where the hub $h$ and the
junction $w$ of any edge $h[w]$ are distinct vertices of $G$. Consequently $Q_l$
is a simple graph without isolated vertices.
\deps{\ref{s7:consRound}}
\end{lemma}

\begin{proof}
By the definition of ranks, a sub-layer contains at most one edge of rank
$\iota$ between any two vertices, so it has no parallel edges. An edge of a PAR
layer joins $[w_1]$ and $[w_2]$ with $w_1\ne w_2$, because looped PAR objects
were paid in (e3); so PAR sub-layers have no loops. An edge $h[w]$ of a HUB
layer joins the side $D_l$ to the side $\{[w]\}$. Moreover $h$ is the hub of a
coloured, hence live, item, so $h$ is a live vertex and $h\notin\Pool_l$ by (a2),
while $w\in\Pool_l$. So $h\ne w$ as vertices of $G$, and HUB sub-layers have no
loops even when $[w]$ is identified with $w$. So every sub-layer is simple.
$Q_l$ is a vertex-disjoint union of the sub-layers with their isolated vertices
deleted, so it is simple and has no isolated vertices.
\end{proof}

\begin{lemma}[Lemma JV-L: the lift]\label{s7:lemLift}
Let $3\le l\le R$, fix any past, any $\xi_l$, and let $\Live$ be the set of live
vertices of round $l$.
\begin{enumerate}
\item[(i)] $\Live\cap\Pool_l=\emptyset$. Every port carries at most one PAR
object of each PAR colour and at most one coloured hub item of each HUB colour.
Hence, although ports are not vertices of $Q_l$, at most one edge of each
sub-layer is an object or item with an end at a given port. Every fresh centre is the middle of at
most one cherry of each PAR colour. Hubs, fresh centres and ports are pairwise
disjoint.
\item[(ii)] For \emph{every} decomposition $\Dec$ of $E(Q_l)$ into cycles and
single edges, the lift of each cycle of $\Dec$ (Construction~\ref{s7:consRound}(h))
is a cycle of $G$, and the lift of each single edge of $\Dec$ consists of one or
two single edges of $G$. The lifts of distinct members of $\Dec$ are pairwise
edge-disjoint. Hence the lift of $\Dec$ consists of at most $2|\Dec|$ pairwise
edge-disjoint objects of $G$.
\item[(iii)] Every item is either paid or lies in exactly one edge of $Q_l$.
Every junction edge $uw$ is the junction edge of at most one item-end, and that
end is at $u$, never at $w$. The edge $uw$ lies in exactly one lent class, namely
$\LJV_{Y(u),l}$, where $Y(u)$ is a lend-good ancestor of round at most $l-2$; the
only consumer of this class is the round-$l$ step. Junction edges lie in
$E_{r}(Y)$ for ancestors $Y$ of rounds $r\le l-2$, items lie in $E_l(Z)$ for
$Z\in\Std_l$, and these edge sets are disjoint.
\item[(iv)] Consequently the round step of Construction~\ref{s7:consRound},
applied at every round $3\le l\le R$, is a J-consumer:
$\Obj_l$ and $\LentJV_l$ satisfy \JC{1}--\JC{3} of
Definition~\ref{s6:defJconsumer}. The objects of $\Obj_l$ other than paid single
edges number at most $2|\Dec_l|=2\f(Q_l)$.
\end{enumerate}
\deps{\ref{s7:consRound}, \ref{s7:lemSimple}, \ref{s7:lemWellDef}, \ref{s6:lemJplus}, \ref{s6:defJconsumer}, \ref{s3:defCOL}, \ref{s2:propStructure}, \ref{s6:defLending}}
\end{lemma}

\begin{proof}
(i) The first claim is (a2) of Construction~\ref{s7:consRound}. The claims about
ports and fresh centres are Lemma~\ref{s7:lemWellDef}(i) together with the SDR
of step (d). Since every edge of a sub-layer of colour $\kappa$ is a PAR object or
a coloured hub item of colour $\kappa$, at most one edge of each sub-layer is an
object or item with an end at a given port. Disjointness of the three roles is \Jp{3} of
Lemma~\ref{s6:lemJplus}.

(ii) \emph{PAR cycles.} Let $[w_0]\cdots[w_{m-1}]$ be a cycle of a PAR sub-layer
of colour $\kappa$. By Lemma~\ref{s7:lemSimple} the sub-layer is simple, so
$m\ge3$. Distinct vertices of the sub-layer are copies of distinct pooled
vertices, so $w_0,\dots,w_{m-1}$ are distinct. The objects $o_0,\dots,o_{m-1}$
are distinct, because they are distinct edges of the cycle, and all have colour
$\kappa$. By (i) a port carries at most one object of colour $\kappa$, so ports
of distinct objects are distinct. The two ends of one PAR object lie at distinct
ports ($u\ne v$ for a $\Jpar$ item, $u\ne u'$ for a cherry), so $p_i\ne q_i$.
Hence the $2m$ ports $p_i,q_i$ are pairwise distinct. The middles $x_i$ of the
cherries among the $o_i$ are distinct, because a fresh centre is the middle of at
most one cherry of colour $\kappa$. Ports lie in $\Qs_Z$ and middles in $F_Z$, so
by (i) no port is a middle. Ports and middles are live, junctions are pooled,
and $\Live\cap\Pool_l=\emptyset$. So all vertices of the closed walk
$w_0p_0(x_0)q_0w_1\cdots q_{m-1}w_0$ are pairwise distinct. Consecutive vertices
are adjacent in $G$: $w_ip_i$ and $q_iw_{i+1}$ are junction edges (edges of
$H_{Y(p_i)}$ and $H_{Y(q_i)}$), and $p_iq_i$ or $p_ix_i,x_iq_i$ are the J-edges
of $o_i$. The walk has length at least $3m\ge9$. A closed walk of length at
least $3$ with pairwise distinct vertices is a cycle of $G$.

\emph{HUB cycles.} Let $h_0[w_0]h_1[w_1]\cdots h_{m-1}[w_{m-1}]h_0$ be a cycle of
a HUB sub-layer of colour $\kappa$. The sub-layer is simple and bipartite
(Lemma~\ref{s7:lemSimple}), so the cycle has even length at least $4$, and
$m\ge2$. The hubs $h_i$ are pairwise distinct, and so are the junctions $w_i$.
The $2m$ items $(h_i,u_i)$ and $(h_{i+1},u'_i)$ are distinct, because they are
distinct edges of the cycle, and all have colour $\kappa$. A port carries at most
one coloured hub item of colour $\kappa$, so the $2m$ ports $u_i,u'_i$ are
pairwise distinct; in particular $u_i\ne u'_i$. Hubs and ports are disjoint by
(i), and hubs and ports are live while junctions are pooled. So all vertices of
the closed walk $h_0u_0w_0u'_0h_1\cdots u'_{m-1}h_0$ are distinct. Consecutive
vertices are adjacent: $h_iu_i$ and $u'_ih_{i+1}$ are items, and $u_iw_i$,
$w_iu'_i$ are junction edges. The length is $4m\ge8$, so the lift is a cycle of
$G$.

\emph{Single edges.} A single edge of a PAR sub-layer is lifted to the one or two
J-edges of its PAR object, and a single edge of a HUB sub-layer to one item. So
every member of $\Dec$ lifts to at most two objects.

\emph{Edge-disjointness.} Every live item that is not paid is a coloured hub
item or belongs to exactly one unpaid PAR object, so it lies in exactly one edge
of $Q_l$. Each edge of $Q_l$ lies in exactly one member of $\Dec$. So the item
edges used by the lifts of distinct members are distinct. A junction edge $uw$
on a lifted cycle is the junction edge of an end at $u$ with junction $w$. The
ends at $u$ have pairwise distinct junctions: in (e1) every chosen junction is
added to $\Used(u)$ and is excluded later, and in (e2) the map is injective into
$\Cand_l(u)\setminus\Used(u)$. So $uw$ determines the end at $u$, hence the
item, the edge of $Q_l$ and the member of $\Dec$. It cannot also be the junction
edge of an end at $w$: $w$ is pooled, so every item with vertex $w$ was paid in
(a2), and no end lies at $w$. Finally, junction edges are not items: a junction
edge lies in $E(H_Y)\subseteq E_r(Y)$ for an ancestor $Y$ of round $r\le l-2$,
an item lies in $E_l(Z)$ with $Z\in\Std_l$, and these edge sets are disjoint
(Proposition~\ref{s2:propStructure}(iii); Lemma~\ref{s6:lemJplus}(iii)).

(iii) The first two claims were shown in the proof of (ii); every item that is
not in an edge of $Q_l$ was paid in (a2), (a3), (b), (d) or (e3). Let $uw$ be a
junction edge of an end at $u$. Then $w\in\Cand_l(u)$, so
$uw\in\LJV_{Y(u),l}\subseteq E(H_{Y(u)})$. The graphs $H_Y$ of distinct ancestors
are edge-disjoint, since $E(H_Y)\subseteq E_r(Y)$ and these sets are pairwise
disjoint (Proposition~\ref{s2:propStructure}(iii)). Inside $\Lend_{Y(u)}$ the
edge has exactly one colour. So $uw$ lies in exactly one lent class, namely
$\LJV_{Y(u),l}$. By COL(g) of Definition~\ref{s3:defCOL}, the only consumer of
this class is the round-$l$ J-consumer. The port $u$ carries an end, so it is
live, and hence $u\in\Qs_Z=Q_Z\setminus\Lost_Z$. Since $\Lost_Z$ contains every
classed port whose class is lend-bad, $Y(u)$ is lend-good. Finally
$r(u)\le l-2$. The disjointness claim was shown in (ii).

(iv) \JC{1}: $\LentJV_l$ consists of junction edges on lifted cycles, and by
(iii) each lies in $\LJV_{Y,l}$ for a lend-good ancestor $Y$ of round at most
$l-2$. \JC{2}: the paid single edges are J-edges that lie in no edge of $Q_l$;
they are distinct from each other and from all edges used by the lifts, which
consist of items in edges of $Q_l$ and of junction edges. Together with (ii), the
objects of $\Obj_l$ are pairwise edge-disjoint cycles and single edges. Their
union contains every edge of $\Jlost_l$ (paid in (a1)) and every item: each item
is paid or lies in an edge of $Q_l$, whose member of $\Dec_l$ lifts to objects
containing all items of that edge. The only other edges used are the junction
edges on lifted cycles, which form $\LentJV_l$ by definition. So the union is
exactly $J_l\cup\LentJV_l$. \JC{3} holds because no other edge is used. The
bound $2|\Dec_l|=2\f(Q_l)$ is (ii) for $\Dec=\Dec_l$. Junction edges that were
assigned but do not lie on a lifted cycle are not in $\LentJV_l$; they stay in
$E_r(Y)\setminus\Lent(Y)$ as junk for their owner.
\end{proof}

%% ----------------------------------------------------------------------
\subsection{Quotient size and payments}

In this subsection $3\le l\le R$ and the past $\Past_l$ are fixed. We say that
\emph{the lists are fixed} when, in addition, the variables $\eta_h$ and
$\zeta_{h,u}$ of $\xi_l$ are fixed. Then steps (a)--(d) and (e1) of
Construction~\ref{s7:consRound}, the sets $E'(u)$ and their orders
$e_1,\dots,e_q$ are determined (by the restrictions on the fixed rules), and
only the orders $\prec_u$ are random. By step (e2), the junctions of the ends at distinct ports are then
independent, and the junction of an end in $E'(u)$ is uniform on
$\Cand_l(u)\setminus\Used(u)$, a set of size at least $\Hcd_l/2$
(Lemma~\ref{s7:lemWellDef}(v)). By \Jp{2} of Lemma~\ref{s6:lemJplus},
$|J_l|\le n(M_l-1)$, hence $|J_l|\le1.37n(M_l-1)$; in this subsection we
deliberately use the looser form (it only enlarges constants). Hence the number of PAR
objects and the number of hub items are each at most $1.37nM_l$. For a PAR colour $\kappa$ and $w\ne w'$ in
$\Pool_l$, let $m_\kappa(w,w')$ be the number of edges $[w][w']$ of $\BP_\kappa$.
As in Lemma~\ref{s7:lemMULT}, $[w]$ is not isolated in exactly
$\max_{w'}m_\kappa(w,w')$ of the sub-layers $B^{\mathrm{P},(\iota)}_\kappa$.

\begin{lemma}[Lemma CC: PAR copies]\label{s7:lemCC}
Let the past and the lists be fixed. Fix a PAR colour $\kappa$ and $w\in\Pool_l$.
Condition further on the indicators
\[
I_x\defeq\ind{\text{the end at $x$ of the $\kappa$-object at $x$ receives the junction }w}
\]
for every port $x$ carrying a PAR object of colour $\kappa$. Let $S_w$ be the
set of PAR objects of colour $\kappa$ with exactly one end whose junction is $w$.
\begin{enumerate}
\item[(i)] The partner ports $v_o$ of the objects $o\in S_w$ (the ports of the
ends whose junction is not $w$) are pairwise distinct. Under the conditioning,
the junctions of the partner ends are independent, and the junction of the
partner end of $o$ is uniform on $(\Cand_l(v_o)\setminus\Used(v_o))\setminus\{w\}$.
In particular it equals a given $w'$ with probability at most $3/\Hcd_l$.
\item[(ii)] For every integer $t\ge1$,
\[
\Ex\Bigl[\max_{w'}m_\kappa(w,w')\ \Big|\ \Past_l,\text{lists},(I_x)_x\Bigr]
\ \le\ t\cdot\ind{S_w\ne\emptyset}+\frac{6e|S_w|}{\Hcd_l}+4e\,2^{-t}|S_w| .
\]
\item[(iii)] With $t\defeq\tCC_l\defeq\lceil2\log_2M_l\rceil$, the number of
vertices of $Q_l$ lying in PAR sub-layers satisfies
\[
\Ex\bigl[\#\{\text{PAR copies}\}\ \big|\ \Past_l,\text{lists}\bigr]\ \le\
3M_l(\tCC_l+1)|\Pool_l|+44.7\,\frac{nM_l}{\Hcd_l}+29.8\,\frac{n}{M_l}.
\]
\end{enumerate}
These bounds hold for every past, in particular for every realized $J_l$, and for
every outcome of the lists, the SDR and the greedy step (e1): only the orders
$\prec_u$ are random.
\deps{\ref{s7:consRound}, \ref{s7:lemWellDef}, \ref{s1:citChernoffGen}, \ref{s6:lemJplus}}
\end{lemma}

\begin{proof}
The edges at $[w]$ in $\BP_\kappa$ are exactly the objects in $S_w$. Indeed, an
object with both ends at junction $w$ is looped and was paid in (e3), and an
object with no end at junction $w$ is not incident to $[w]$. So
$m_\kappa(w,w')$ is the number of $o\in S_w$ whose partner end has junction
$w'$.

(i) A PAR object has its two ends at distinct ports, and every port carries at
most one PAR object of colour $\kappa$ (Lemma~\ref{s7:lemWellDef}(i)). So every
such port $x$ carries exactly one end of colour $\kappa$, the indicator $I_x$ is
well defined, and the partner ports of distinct objects of $S_w$ are distinct.
The set $S_w$ and the partner ports are determined by the indicators. With the
past and the lists fixed, $\Cand_l(x)\setminus\Used(x)$, $E'(x)$ and its order
are fixed for every port $x$ (restrictions on the fixed rules in
Construction~\ref{s7:consRound}), so the junction of the $\kappa$-end at a
port $x$ is a function of $\prec_x$ alone, and the orders of distinct ports are independent.
So conditioning on $(I_x)_x$ keeps the partner ends independent, and it
conditions the partner end at $v_o$ only on the event $I_{v_o}=0$. The junction
of that end is uniform on $C(v_o)\defeq\Cand_l(v_o)\setminus\Used(v_o)$, since
it is the image of a fixed element under a uniform injection. Given $I_{v_o}=0$
it is therefore uniform on $C(v_o)\setminus\{w\}$. By
Lemma~\ref{s7:lemWellDef}(v), $|C(v_o)\setminus\{w\}|\ge\Hcd_l/2-1\ge\Hcd_l/3$.

(ii) If $S_w=\emptyset$ the left side is $0$. Otherwise, by (i), each
$m_\kappa(w,w')$ is, under the conditioning, a sum of independent indicators
with mean $\mu_{w'}\le3|S_w|/\Hcd_l$, and $\sum_{w'}\mu_{w'}=|S_w|$. Put
$T\defeq\max(t,\lceil6e|S_w|/\Hcd_l\rceil)$. For every integer $j\ge T$ we have
$j\ge2e\mu_{w'}$. By Cited result~\ref{s1:citChernoffGen},
$\Prob(m_\kappa(w,w')\ge j)\le(e\mu_{w'}/j)^j\le(e\mu_{w'}/j)\,2^{1-j}$. Hence,
writing $m$ for $m_\kappa(w,w')$,
\begin{align*}
\Ex\bigl[m\,\ind{m\ge T}\bigr]
&=T\,\Prob(m\ge T)+\sum_{j>T}\Prob(m\ge j)
\le e\mu_{w'}2^{1-T}+\sum_{j>T}\frac{e\mu_{w'}}{j}2^{1-j}\\
&\le e\mu_{w'}\,2^{2-T}\le e\mu_{w'}\,2^{2-t}.
\end{align*}
Moreover
\[
\max_{w'}m_\kappa(w,w')\le(T-1)+\sum_{w'}m_\kappa(w,w')\,\ind{m_\kappa(w,w')\ge T}:
\]
if the maximum is below $T$ it is at most $T-1$, and otherwise it is one of the
summands. Taking conditional expectations,
\[
\Ex\Bigl[\max_{w'}m_\kappa(w,w')\ \Big|\ \cdot\Bigr]\le T-1+4e\,2^{-t}|S_w|
\le t+\frac{6e|S_w|}{\Hcd_l}+4e\,2^{-t}|S_w|,
\]
because $T-1\le\max(t-1,6e|S_w|/\Hcd_l)$.

(iii) The number of PAR copies is
$\sum_{\kappa\in[3M_l]}\sum_{w\in\Pool_l}\max_{w'}m_\kappa(w,w')$. Averaging (ii)
over the indicators and summing over the at most $3M_l|\Pool_l|$ pairs
$(\kappa,w)$ gives at most
$3M_l\,t\,|\Pool_l|+\bigl(6e/\Hcd_l+4e\,2^{-t}\bigr)\,\Ex\sum_{\kappa,w}|S_w|$.
Every PAR object lies in at most two of the sets $S_w$ (one for each end), so
$\sum_{\kappa,w}|S_w|\le2\cdot1.37nM_l=2.74\,nM_l$. With $t=\tCC_l$ we have
$2^{-t}\le M_l^{-2}$, and $6e\cdot2.74<44.7$ and $4e\cdot2.74<29.8$. Finally
$t\le t+1$.
\end{proof}

\begin{lemma}[Ultra-hub copies]\label{s7:lemUltra}
Let the past and the lists be fixed.
\begin{enumerate}
\item[(i)] Let $h$ be an ultra hub, $\kappa\in[4M_l]$, and let
$\omega_1,\dots,\omega_N$ be the junctions of the coloured hub items of $h$ of
colour $\kappa$. These items lie at $N$ distinct ports, because distinct items
at $h$ are distinct edges $hu$ and $G$ is simple. The $\omega_i$ are independent,
and each is uniform on a set of size at least $\Hcd_l/2$.
\item[(ii)] If $N\ge1$ then, with $\mu^*\defeq2N/\Hcd_l$,
$\Ex[\max_w m_\kappa(h,w)\mid\Past_l,\text{lists}]\le\log_2N+2e\mu^*+8$.
\item[(iii)] The number of vertices of HUB sub-layers that are copies of $h$
satisfies
$\Ex[\,\cdot\mid\Past_l,\text{lists}]\le4M_l(\log_2\clive_h+8)+4e\,\clive_h/\Hcd_l$.
\item[(iv)] Summed over all ultra hubs, the expected number of hub copies of
ultra hubs, given $\Past_l$ and the lists, is at most
\[
320\,\frac{nM_l^{3}\,(\log_2\thult_l+8)}{\lambda_{l-2}^{95}} .
\]
\end{enumerate}
The bound depends on the past only through deterministic counts, so it holds for
every past and every realized $J_l$.
\deps{\ref{s7:consRound}, \ref{s7:lemWellDef}, \ref{s6:lemJplus}, \ref{s1:citChernoffGen}, \ref{s7:lemMULT}, \ref{s7:lemCand}}
\end{lemma}

\begin{proof}
(i) Items of ultra hubs receive their junctions in (e2). The junction of the
item $(h,u_i)$ is the image of its end under the uniformly random injection at
$u_i$, so it is uniform on $\Cand_l(u_i)\setminus\Used(u_i)$, a set of size at
least $\Hcd_l/2$ by Lemma~\ref{s7:lemWellDef}(v). The ports $u_i$ are distinct,
and injections at distinct ports are independent, so the $\omega_i$ are
independent. (Distinctness of the ports comes from the simplicity of $G$; it
does not use \Jp{1}.)

(ii) Let $m_w\defeq\#\{i:\omega_i=w\}$. It is a sum of independent indicators
with mean $\mu_w\le2N/\Hcd_l=\mu^*$, and $\sum_w\mu_w=N$. Put
$T\defeq\max(\lceil2e\mu^*\rceil,\lceil\log_2N\rceil+1)$. For $j\ge T$ we have
$e\mu_w/j\le1/2$, so by Cited result~\ref{s1:citChernoffGen}
$\Prob(m_w\ge j)\le(e\mu_w/j)^j\le(e\mu_w/j)2^{1-j}$, and summing over $w$ gives
$\sum_w\Prob(m_w\ge j)\le(eN/j)2^{1-j}$. Hence
\[
\Ex\bigl[\max_wm_w\bigr]=\sum_{j\ge1}\Prob\bigl(\max_wm_w\ge j\bigr)
\le(T-1)+\sum_{j\ge T}\frac{eN}{j}2^{1-j}\le(T-1)+\frac{eN}{T}2^{2-T}
\le T-1+\frac{2e}{T},
\]
using $2^{-T}\le2^{-\log_2N-1}=1/(2N)$. Finally
$T-1\le\max(2e\mu^*,\log_2N+1)\le\log_2N+2e\mu^*+1$ and $2e/T\le2e<7$.

(iii) By Lemma~\ref{s7:lemMULT}, the number of copies of $h$ in the sub-layers of
colour $\kappa$ is $\max_wm_\kappa(h,w)$. Let $N_\kappa$ be the number of
coloured items of $h$ of colour $\kappa$. Then $N_\kappa\le\clive_h$ and
$\sum_\kappa N_\kappa\le\clive_h$, and at most $4M_l$ colours have
$N_\kappa\ge1$. By (ii) the expected number of copies of $h$ is at most
$\sum_{\kappa:N_\kappa\ge1}\bigl(\log_2\clive_h+8\bigr)+2e\sum_\kappa2N_\kappa/\Hcd_l
\le4M_l(\log_2\clive_h+8)+4e\,\clive_h/\Hcd_l$.

(iv) The function $g(x)\defeq(\log_2x+8)/x$ is decreasing on $[1,\infty)$, since
$x^2g'(x)=1/\ln2-8-\log_2x<0$ there. For an ultra hub, $\clive_h>\thult_l\ge1$,
so $\log_2\clive_h+8=\clive_h\,g(\clive_h)\le\clive_h(\log_2\thult_l+8)/\thult_l$.
By \Jp{2} of Lemma~\ref{s6:lemJplus},
$\sum_h\clive_h\le|J_l|\le1.37nM_l$. We have
$M_l\Hcd_l/7=\lambda_{l-2}^{95}/(56M_l)$, and $M_l\Hcd_l\ge2^{10}M_l^{11}$
(Lemma~\ref{s7:lemCand}(iv)), so $\lambda_{l-2}^{95}/M_l\ge2^{453}$. Hence
$\thult_l\ge\lambda_{l-2}^{95}/(56M_l)-1\ge\lambda_{l-2}^{95}/(57M_l)$, and also
$\thult_l\ge2^{8}$. Summing (iii) over ultra hubs,
\[
\sum_{h\text{ ultra}}\Bigl(4M_l(\log_2\clive_h+8)+\frac{4e\,\clive_h}{\Hcd_l}\Bigr)
\le4M_l\cdot1.37nM_l\cdot\frac{57M_l}{\lambda_{l-2}^{95}}(\log_2\thult_l+8)
+4e\cdot1.37nM_l\cdot\frac{8M_l^2}{\lambda_{l-2}^{95}} .
\]
The first term is $312.36\,nM_l^3(\log_2\thult_l+8)/\lambda_{l-2}^{95}$ and the
second is at most $119.2\,nM_l^3/\lambda_{l-2}^{95}$. Since
$\log_2\thult_l+8\ge16$, the second term is at most
$7.5\,nM_l^3(\log_2\thult_l+8)/\lambda_{l-2}^{95}$, and $312.36+7.5<320$.
\end{proof}

\begin{lemma}[The stage-1 weight $\XV$]\label{s7:lemVstar}
For $3\le l\le R$ put
\[
X_{\mathrm{V},l}\defeq3M_l(\tCC_l+1)|\Pool_l|+4M_l\sum_{w\in\Pool_l}\mult_{r(w)}(w),
\qquad \XV\defeq\sum_{l=3}^{R}X_{\mathrm{V},l}.
\]
Each $X_{\mathrm{V},l}$ is a deterministic function of the run and the pool
labels, and
\[
\Ex X_{\mathrm{V},l}\le\frac{(6\log_2M_l+12)\,n}{M_l},\qquad
\Ex\XV\le2.1\,\frac{(6\log_2\Dstar+12)\,n}{\Dstar}.
\]
\deps{\ref{s7:defPool}, \ref{s2:propOV}, \ref{s2:lemTower}}
\end{lemma}

\begin{proof}
The quantities $M_l$, $\tCC_l$ and $\mult_r(w)$ are determined by the run, and
$\Pool_l$ and $r(w)$ by the pool labels. Since $\tCC_l+1\le2\log_2M_l+2$ and
$\Ex|\Pool_l|\le q_ln=n/M_l^2$, the first summand has expectation at most
$(6\log_2M_l+6)n/M_l$. For the second,
\[
\Ex\sum_{w\in\Pool_l}\mult_{r(w)}(w)=\sum_{w}\sum_{r=1}^{l-2}q_l\pi_{l,r}\,\mult_r(w)
=q_l\sum_{r=1}^{l-2}\pi_{l,r}\sum_w\mult_r(w)\le q_l\cdot1.37n ,
\]
since $\sum_w\mult_r(w)\le1.37n$ for every $r$ (Proposition~\ref{s2:propOV}) and
$\sum_r\pi_{l,r}<1$. So the second summand has expectation at most
$4M_l\cdot1.37n/M_l^2=5.48\,n/M_l$. Adding, $\Ex X_{\mathrm{V},l}\le(6\log_2M_l+11.48)n/M_l$.
Summing over $l$ and using $\sum_l\psi(M_l)\le2.1\,\psi(\Dstar)$ for
$\psi(x)=(6\log_2x+12)/x$ (Lemma~\ref{s2:lemTower}(b)) gives the bound for
$\Ex\XV$.
\end{proof}

\begin{lemma}[Payments]\label{s7:lemPay}
For every past and every $3\le l\le R$ the following hold.
\begin{enumerate}
\item[(a1)] $|\Jlost_l|\le(M_l-1)|\Lost_l|$.
\item[(a2)] The number of items with a vertex in $\Pool_l$ is at most
$\sum_{v\in\Pool_l}\omega_l(v)$, where $\omega_l(v)\defeq\cagg_{v,l}$ if
$v\in D_l$ and $\omega_l(v)\defeq M_l$ otherwise.
\item[(a3)] The number of items with a JV-bad port end is at most $M_l-1$ times
the number of JV-bad classed ports of round $l$.
\item[(b)] Every fresh centre has at most one unpaired fresh leg per round, so
the number of unpaired fresh legs over all rounds is at most
$\sum_{l,Z}|F_Z|\le2n$.
\item[(c)] $\Ex[\text{SDR-failure payments}\mid\Past_l]\le n(M_l-1)(4M_l)^{-4}\le n/(256M_l^3)$,
and $\Ex[\text{loop-paid edges}\mid\Past_l]\le2\cdot(2/\Hcd_l)\cdot\#\{\text{PAR objects}\}\le5.5\,nM_l/\Hcd_l$.
\item[(d)] If $\xi_l$ lies in the event of Construction~\ref{s7:consRound}(f),
then $\payrd_l\le4\bigl(n/(256M_l^3)+5.5\,nM_l/\Hcd_l\bigr)\le2n/M_l$; if this
holds at every round, then $\sum_l\payrd_l\le9n/\Dstar$.
\item[(e)] No other item is paid. In particular, no item is paid because of a
coincidence of junctions at an ultra hub, and no item is paid for a concentrated
pair.
\end{enumerate}
\deps{\ref{s7:consRound}, \ref{s7:lemWellDef}, \ref{s6:lemJplus}, \ref{s2:propParentless}, \ref{s2:lemTower}}
\end{lemma}

\begin{proof}
We use the properties of Lemma~\ref{s6:lemJplus}.

(a1) Every $\Jlost$ edge joins a vertex of $\Lost_l$ to a port, and every vertex
carries at most $M_l-1$ J-edges by \Jp{2}.

(a2) Let $v\in\Pool_l$. If $v\in D_l$, then $v$ is a hub and not a port or a
fresh centre, by \Jp{3}. The items with vertex $v$ are therefore $\Jhub$ items
$(v,u)$: $\Jpar$ items have both ends in $\Qs_Z$, and $\Jfr$ items have their
centre in $F_Z$. By \Jp{1} there is at most one such item per class $Y$. Each
such class has $d_{Y,l}(v)\ge1$, because the item is an edge $vu\in E_l(Z)$ with
$u\in Q_Z$ and $Y(u)=Y$. So there are at most $\cagg_{v,l}$ items with vertex
$v$. If $v\notin D_l$, then $v$ carries at most $M_l-1\le M_l$ J-edges by \Jp{2}.

(a3) A port carries at most $M_l-1$ items by \Jp{2}.

(b) The pairing at a fresh centre $x$ leaves at most one live $\Jfr$ item
unpaired. Every fresh centre of round $l$ lies in $F_Z$ for some $Z\in\Std_l$,
and $\sum_{l,Z}|F_Z|\le2n$ by Proposition~\ref{s2:propParentless}(i).

(c) By \Jp{3}, every non-hub vertex has one role in one part per round, so there
are at most $n$ ports. A port $u$ pays at most $M_l-1$ items for an SDR failure,
and it fails with probability at most $(4M_l)^{-4}$ given $\Past_l$
(Lemma~\ref{s7:lemWellDef}(iii)). This gives the first bound. For loops, fix the
lists. A PAR object has its two ends at distinct ports $p\ne q$. Their junctions
are independent, and the junction at $q$ is uniform on a set of size at least
$\Hcd_l/2$ (Lemma~\ref{s7:lemWellDef}(v)). So the probability that the two
junctions coincide is
$\sum_w\Prob(p\to w)\Prob(q\to w)\le\max_w\Prob(q\to w)\le2/\Hcd_l$. A looped
object costs at most two edges (a cherry costs two, a $\Jpar$ item one), and
there are at most $1.37nM_l$ PAR objects (\Jp{2}). So the expected number of
loop-paid edges is at most $4\cdot1.37\,nM_l/\Hcd_l\le5.5\,nM_l/\Hcd_l$. This
bound does not depend on the lists, so it also holds given $\Past_l$ alone.

(d) The first inequality is the definition of the event in (f), together with
(c). Next, $4\cdot5.5\,nM_l/\Hcd_l=176\,nM_l^3/\lambda_{l-2}^{95}$. By
Lemma~\ref{s2:lemTower}(b), $M_l\le\lambda_{l-2}^{1.6}$, so
$\lambda_{l-2}\ge M_l^{1/1.6}\ge2^{25}$ and
$176\,M_l^3/\lambda_{l-2}^{95}\le176\,\lambda_{l-2}^{-90.2}\le\lambda_{l-2}^{-1.6}\le1/M_l$.
Also $4n/(256M_l^3)\le n/M_l$. Hence $\payrd_l\le2n/M_l$, and
$\sum_l\payrd_l\le2n\sum_l1/M_l\le4n/\Dstar\le9n/\Dstar$ by
Lemma~\ref{s2:lemTower}(b).

(e) In Construction~\ref{s7:consRound}, single edges are paid only in (a1),
(a2), (a3), (b), (d) and (e3).
\end{proof}

\begin{definition}[The stage-1 functional $\Xp$]\label{s7:defXprime}
With $\omega_l$ as in Lemma~\ref{s7:lemPay}(a2), put
\[
\Xpool\defeq\sum_{l=3}^{R}\Bigl[\sum_{v\in\Pool_l}\omega_l(v)+(M_l-1)\cdot\#\{\text{JV-bad classed ports of round }l\}\Bigr],
\]
let $\XV$ be as in Lemma~\ref{s7:lemVstar}, and let $\XU$ be as in
Definition~\ref{s6:defLending}. Then
\[
\Xp\defeq\XU+\Xpool+\XV .
\]
Each of the three terms is non-negative and is a deterministic function of the
run, $\delta$ and the stage-1 outcome. Indeed, $\omega_l(v)$ is determined by the
run and $\delta$ (through $\cagg_{v,l}$), $\Pool_l$ is determined by the pool
labels, and JV-badness is determined by stage~1. $\Xp$ is the functional used in
stage (1e) of the schedule.
\deps{\ref{s6:defLending}, \ref{s7:lemVstar}, \ref{s7:lemPay}}
\end{definition}

\begin{lemma}[Expectation of $\Xp$]\label{s7:lemEXprime}
$\Ex\Xp\le\epsX(\Dstar)\,n$, where
\[
\epsX(\Dstar)\defeq\epsU(\Dstar)+\frac{10.3}{\Dstar}+2.1\,\frac{6\log_2\Dstar+12}{\Dstar}
\]
and $\epsU$ is the function of Lemma~\ref{s5:lemExpect} (through
Lemma~\ref{s6:lemLost}). More precisely, $\Ex\XU\le\epsU(\Dstar)n+5.5n/\Dstar$,
$\Ex\Xpool\le4.8n/\Dstar$ and $\Ex\XV\le2.1(6\log_2\Dstar+12)n/\Dstar$.
\deps{\ref{s7:defXprime}, \ref{s6:lemLost}, \ref{s7:lemCand}, \ref{s7:lemVstar}, \ref{s2:propOV}, \ref{s2:propStructure}, \ref{s2:lemTower}, \ref{s5:lemExpect}, \ref{s2:lemCap}}
\end{lemma}

\begin{proof}
The bound on $\Ex\XU$ is Lemma~\ref{s6:lemLost}, and the bound on $\Ex\XV$ is
Lemma~\ref{s7:lemVstar}. For $\Xpool$, fix a round $l$.

\emph{Pooled weights.} The weights $\omega_l(v)$ are determined by the run and
$\delta$, so they are independent of the pool labels. Hence
$\Ex\sum_{v\in\Pool_l}\omega_l(v)\le q_l\sum_v\omega_l(v)\le q_l\bigl(\sum_{h\in D_l}\cagg_{h,l}+nM_l\bigr)$.
Now $\cagg_{h,l}\le\sum_Yd_{Y,l}(h)$, and $\sum_Yd_{Y,l}(h)$ counts the edges
$hu\in E_l(Z)$ with $Z\in\Std_l$ and $u\in Q_Z$. Summing over $h$ counts each such
edge at most once for each of its classed-port ends. So
$\sum_h\cagg_{h,l}\le\sum_{Z\in\Std_l}\sum_{u\in Q_Z}\deg_{E_l(Z)}(u)$. Every
vertex has $E_l(Z)$-degree at most $M_l-1$ (Lemma~\ref{s2:lemCap}(ii)), and
$\sum_{Z\in\Std_l}|Z^0|\le1.37n$ by Proposition~\ref{s2:propOV}. Hence
$\sum_{h\in D_l}\cagg_{h,l}\le1.37nM_l$, and the expected pooled weight is at
most $2.37\,n/M_l$.

\emph{JV-bad ports.} There are at most $n$ classed ports of round $l$, because
$Q_Z\subseteq U_Z$ and the port sets $U_Z$ of distinct $Z\in\Std_l$ are pairwise
disjoint (Proposition~\ref{s2:propStructure}(iv)); this is the count used in
the proofs of \ref{s6:K6} of Lemma~\ref{s6:lemLost} and of \Jp{2} of
Lemma~\ref{s6:lemJplus}. By Lemma~\ref{s7:lemCand}(iii) each is JV-bad with
probability at most $\exp(-\Hcd_l/4)$, and $\exp(-\Hcd_l/4)\le M_l^{-3}$ because
$\Hcd_l\ge2^{10}M_l^{10}$ (Lemma~\ref{s7:lemCand}(iv)). So the expected JV-bad
weight is at most
\[
(M_l-1)\cdot n\cdot M_l^{-3}\ \le\ n/M_l^2\ \le\ 0.03\,n/M_l .
\]
The pooled-weight bound above counts the same classed ports by the
deliberately looser $\sum_{Z\in\Std_l}|Z^0|\le1.37n$; this only enlarges its
constant.

Hence $\Ex\Xpool\le\sum_l2.4\,n/M_l\le4.8\,n/\Dstar$, using
$\sum_l1/M_l\le2/\Dstar$ (Lemma~\ref{s2:lemTower}(b)). Adding the three bounds
gives $\Ex\Xp\le\epsU(\Dstar)n+(5.5+4.8)n/\Dstar+2.1(6\log_2\Dstar+12)n/\Dstar=\epsX(\Dstar)n$.
\end{proof}

\begin{lemma}[Lemma UH$^*$-split: quotient size]\label{s7:lemUHsplit}
For $3\le l\le R$ put
\[
\dett_l\defeq3|D_l|+78\,\frac{nM_l}{\Hcd_l}+30\,\frac{n}{M_l}
+320\,\frac{nM_l^{3}(\log_2\thult_l+8)}{\lambda_{l-2}^{95}},
\]
a deterministic function of the run.
\begin{enumerate}
\item[(i)] $X_{\mathrm{V},l}$ is the stage-1 weight of Lemma~\ref{s7:lemVstar}.
\item[(ii)] For every past, $\Ex[\copies_l\mid\Past_l]\le X_{\mathrm{V},l}+\dett_l$.
\item[(iii)] If the stage-1 outcome satisfies $\Xp\le3\,\Ex\Xp$ and at every round
$\xi_l$ lies in the event of Construction~\ref{s7:consRound}(f), then
\[
\sum_{l=3}^{R}|V(Q_l)|\ \le\ 12\,\Ex\Xp+4\sum_{l=3}^{R}\dett_l .
\]
\item[(iv)] Put $F(x)\defeq(3184+30400\log_2x)\,x^{-90.2}$ for $x\ge1$.
Then $\sum_l\dett_l\le3\epsA n+60n/\Dstar+2F(\log_2\Dstar)\,n$, and
$2F(\log_2\Dstar)<2^{-1000}$. Hence, in the situation of (iii),
\[
\sum_{l=3}^{R}|V(Q_l)|\le\thetaQ\,n,\qquad
\thetaQ\defeq\epsTwo(\Dstar)\defeq12\,\epsX(\Dstar)+4\Bigl(3\epsA+\frac{60}{\Dstar}+2F(\log_2\Dstar)\Bigr),
\]
with $\epsA=31\eps/(\Cp\log\log\Dstar)$ as in Lemma~\ref{s2:lemTower}(e). So
$\thetaQ$ depends on $\Dstar$ only, and $\thetaQ\le1/4$ by
condition~\ref{s1:condG4}.
\end{enumerate}
\deps{\ref{s7:lemCC}, \ref{s7:lemMULT}, \ref{s7:lemUltra}, \ref{s7:lemVstar}, \ref{s7:lemEXprime}, \ref{s7:defSchedule}, \ref{s2:propOV}, \ref{s2:lemTower}, \ref{s2:lemLacunary}, \ref{s1:condG4}, \ref{s1:condG1}, \ref{s7:consRound}, \ref{s7:lemWellDef}, \ref{s2:defHBtp}, \ref{s6:lemJplus}}
\end{lemma}

\begin{proof}
(ii) Fix the past and the lists. The vertices of $Q_l$ are of four kinds.

\emph{PAR copies.} Their conditional expectation is at most
$3M_l(\tCC_l+1)|\Pool_l|+44.7\,nM_l/\Hcd_l+29.8\,n/M_l$ by
Lemma~\ref{s7:lemCC}(iii).

\emph{Copies of non-ultra hubs} (deterministic given the lists). Let $h$ be
non-ultra. By Lemma~\ref{s7:lemWellDef}(iv) (see step (e1) of
Construction~\ref{s7:consRound}), $m_\kappa(h,w)\le1$ for all $\kappa$ and $w$.
So by Lemma~\ref{s7:lemMULT}, $h$ has exactly one copy in each colour that one of
its items uses, and none in the other colours. The colours of its items lie in
the union of its $k_h$ lists, which has $3k_h$ elements. So the number of copies
of non-ultra hubs is at most
$\sum_{h:\clive_h\ge1}3k_h\le\sum_{h\in D_l}3(1+8\clive_h/\Hcd_l)\le3|D_l|+24\cdot1.37\,nM_l/\Hcd_l$,
since $\sum_h\clive_h\le1.37nM_l$. Here $24\cdot1.37=32.88<32.9$.

\emph{Junction copies in HUB sub-layers} (the count depends on the junctions,
but the bound below is deterministic). By
Lemma~\ref{s7:lemMULT}, in colour $\kappa$ the vertex $[w]$ has
$\max_hm_\kappa(h,w)\le\mult_{r(w)}(w)$ copies. Summing over the $4M_l$ HUB
colours and $w\in\Pool_l$ gives at most $4M_l\sum_{w\in\Pool_l}\mult_{r(w)}(w)$.

\emph{Copies of ultra hubs.} Their conditional expectation is at most
$320\,nM_l^3(\log_2\thult_l+8)/\lambda_{l-2}^{95}$ by Lemma~\ref{s7:lemUltra}(iv).

The first summands of the PAR bound and of the junction bound add up to
$X_{\mathrm{V},l}$. The remaining terms add up to at most $\dett_l$, because
$44.7+32.9\le78$ and $29.8\le30$. Averaging over the lists gives (ii).

(iii) By the choice in (f) and by (ii),
$|V(Q_l)|=\copies_l\le4\,\Ex[\copies_l\mid\Past_l]\le4(X_{\mathrm{V},l}+\dett_l)$
at every round. Summing over $l$,
$\sum_l|V(Q_l)|\le4\XV+4\sum_l\dett_l\le4\Xp+4\sum_l\dett_l\le12\,\Ex\Xp+4\sum_l\dett_l$,
because $\XU,\Xpool\ge0$ and $\Xp\le3\Ex\Xp$ (stage (1e) of
Definition~\ref{s7:defSchedule}).

(iv) By Proposition~\ref{s2:propOV}(K2), $|D_l|\le7.6\eps n/\log P_l$. So
$\sum_l3|D_l|\le3n\sum_l7.6\eps/\log P_l\le1.5\epsA n\le3\epsA n$ by
Lemma~\ref{s2:lemTower}(e) (the factor~$2$ of slack is deliberate; it only
enlarges $\epsTwo$). Next,
$\sum_l30n/M_l\le60n/\Dstar$ by Lemma~\ref{s2:lemTower}(b). For the remaining
terms, $78\,nM_l/\Hcd_l=624\,nM_l^3/\lambda_{l-2}^{95}$, and
$\log_2\thult_l\le\log_2(M_l\Hcd_l/7)\le95\log_2\lambda_{l-2}$. With
$M_l\le\lambda_{l-2}^{1.6}$ (Lemma~\ref{s2:lemTower}(b)), these terms are at most
$nF(\lambda_{l-2})$, with $F$ as in (iv).
Here $\lambda_{l-2}\ge M_l^{1/1.6}\ge2^{25}$. On $[2^{25},\infty)$, $F$ is
decreasing, and $F(2x)\le F(x)/2$, since the first factor at most doubles while
$x^{-90.2}$ is multiplied by $2^{-90.2}$. By Lemma~\ref{s2:lemLacunary}(ii),
$\lambda_{r}\ge2^{\lambda_{r+1}/A}\ge2\lambda_{r+1}$ for consecutive rounds,
because $\lambda_{r+1}\ge2^{25}$. So $X_l\defeq F(\lambda_{l-2})$
($3\le l\le R$) satisfies $X_l\le X_{l+1}/2$. By
Lemma~\ref{s2:lemLacunary}(i), $\sum_lX_l\le2X_R=2F(\lambda_{R-2})$ (if $R\le2$
the sum is empty). Moreover $\lambda_{R-2}\ge\lambda_R=\log_2d_R\ge\log_2\Dstar$,
by Lemma~\ref{s2:lemTower}(a) and because round $R$ has $d_R\ge\Dstar$ by
rule~\ref{s2:R0} of Definition~\ref{s2:defHBtp}; and $\log_2\Dstar\ge2^{256}\ge2^{25}$ by \Gc{1}(a) at $\mu=\log_2\log_2\Dstar$. Since $F$
is decreasing on $[2^{25},\infty)$, we get
\[
\sum_lX_l\ \le\ 2F(\log_2\Dstar)\ \le\ 2F(2^{25})\ \le\ 2\cdot2^{20}\cdot2^{-2255}\ <\ 2^{-1000}.
\]
Together, $\sum_l\dett_l\le3\epsA n+60n/\Dstar+2F(\log_2\Dstar)\,n$. The bound
on $\sum_l|V(Q_l)|$ follows from (iii) and $\Ex\Xp\le\epsX(\Dstar)n$
(Lemma~\ref{s7:lemEXprime}). The functions $\epsX$, $\epsA$ and
$F(\log_2\Dstar)$ depend on $\Dstar$ only, and $\thetaQ\le1/4$ is part of
condition~\ref{s1:condG4}.
\end{proof}

%% ----------------------------------------------------------------------
\subsection{One joint outcome, the cost, and Theorem \texorpdfstring{JV$^{+*}$}{JV+*}}

\begin{lemma}[One joint outcome exists]\label{s7:lemOneOutcome}
There is a joint outcome, consisting of a stage-1 outcome, stage-3 outcomes for
all parts, and draws $\xi_R,\dots,\xi_3$, such that:
\begin{enumerate}
\item[(1)] $\Xp\le3\,\Ex\Xp$;
\item[(2)] the stage-3 labels of every part lie in that part's good event: the
event of Lemma~\ref{s4:lemTPV} for every standalone pre-part (with the data of
step (1) of Construction~\ref{s6:consOrder}); the event of Lemma~\ref{s4:lemPV}
as used in Lemma~\ref{s5:lemChild} for every non-demoted light part; and the event
of Theorem~\ref{s4:thmVXp} as used in Lemma~\ref{s5:lemDemoted} for every
demoted light part;
\item[(3)] at every round $3\le l\le R$, $\xi_l$ lies in the event of
Construction~\ref{s7:consRound}(f). Consequently
$\copies_l\le4(X_{\mathrm{V},l}+\dett_l)$ and
$\payrd_l\le4\bigl(n/(256M_l^3)+5.5\,nM_l/\Hcd_l\bigr)$.
\end{enumerate}
Such an outcome is obtained as follows.
(i) Choose the stage-1 outcome in $\{\Xp\le3\Ex\Xp\}$; stage~2 is then
determined.
(ii) For every part, choose its stage-3 labels in its good event.
(iii) For $l=R,R-1,\dots,3$ in turn: run steps (1)--(3) of
Construction~\ref{s6:consOrder} at round $l$, which are deterministic given the
past; draw $\xi_l$ and choose it in the event of
Construction~\ref{s7:consRound}(f); build $Q_l$, fix $\Dec_l$, lift it, and record
$\LentJV_l$ (step (4) of Construction~\ref{s6:consOrder}).
(iv) At rounds $2$ and $1$, run steps (1)--(2) of Construction~\ref{s6:consOrder};
here $J_2=J_1=\emptyset$.
\deps{\ref{s7:defSchedule}, \ref{s7:consRound}, \ref{s6:consOrder}, \ref{s4:lemTPV}, \ref{s4:lemPV}, \ref{s4:thmVXp}, \ref{s5:lemChild}, \ref{s5:lemDemoted}, \ref{s7:lemEXprime}, \ref{s7:lemPay}, \ref{s7:lemUHsplit}, \ref{s1:citMarkov}, \ref{s1:condG3}}
\end{lemma}

\begin{proof}
\emph{Step (i).} $\Xp$ is a non-negative random variable on a finite probability
space. By Markov's inequality (Cited result~\ref{s1:citMarkov}),
$\Prob(\Xp\le3\Ex\Xp)\ge2/3>0$, so the event is non-empty. Fix an outcome in it;
this gives (1). Stage~2 is a deterministic function of stage~1
(Definition~\ref{s7:defSchedule}).

\emph{Step (ii).} With stage~1 fixed, the good event of each part is determined
by stage-1 and stage-2 data together with that part's own stage-3 labels.
\begin{itemize}
\item For a standalone pre-part $Z$, Lemma~\ref{s4:lemTPV} is applied with $O_Z$
and $\vP=\Ret_Z$, both fixed at stage~2, and its good event is determined by
$(Z^0,O_Z,\Ret_Z)$ and the labels.
\item For a non-demoted light part, the good event of Lemma~\ref{s4:lemPV}
depends only on the own classes, $\Rt$, $\extph$ and the labels
(Lemma~\ref{s5:lemChild}).
\item For a demoted light part, the good event of Theorem~\ref{s4:thmVXp}
depends only on $O=X_Y$ and the labels (Lemma~\ref{s5:lemDemoted}).
\end{itemize}
Each of these events has probability at least $1/2-o(1)$, and by
condition~\ref{s1:condG3} every such $o(1)$ is at most $1/100$. So each event
is non-empty, whatever the stage-1 outcome; only this non-emptiness is used. The events of distinct parts involve disjoint families of labels,
so choosing each part's labels inside its own event gives one stage-3 outcome in
all of the events at once; no union bound is needed. None of these events
depends on the edge sets that the devices later receive: the conclusions of
Lemma~\ref{s4:lemTPV}, Lemma~\ref{s4:lemPV} and Theorem~\ref{s4:thmVXp} hold on
their good events for \emph{every} admissible edge set (every $E\supseteq E(O)$,
every $H_0$ containing the own classes, every graph containing $O$). So the
events do not depend on junk, on $\Erem$, on $H_0$ or on later rounds. This
gives (2).

\emph{Step (iii).} At round $l$ the past $\Past_l$ is fixed, and steps (1)--(3)
of Construction~\ref{s6:consOrder} are deterministic given it; in particular
$J_l$ is determined. The draw $\xi_l$ is independent of the past, and
$\copies_l$ and $\payrd_l$ are non-negative functions of $\xi_l$ with the past
fixed. By Markov's inequality (Cited result~\ref{s1:citMarkov}), the event of
Construction~\ref{s7:consRound}(f) has probability at least $1/2$. This holds for
every past and needs nothing beyond Markov's inequality. So a choice of $\xi_l$
in the event exists, whatever happened before. On this event, the explicit
bounds of (3) follow from Lemma~\ref{s7:lemUHsplit}(ii) and
Lemma~\ref{s7:lemPay}(c), which hold for \emph{every} past; this is the only
place where the uniformity in the past is used.

\emph{Step (iv).} Rounds $1$ and $2$ have no classed ports, since every port
there is fresh. So JS-LC produces no J-set, and there is no round step.
\end{proof}

\begin{proposition}[Cost on the chosen outcome]\label{s7:propCost}
On an outcome as in Lemma~\ref{s7:lemOneOutcome}, with the round step of
Construction~\ref{s7:consRound} as the J-consumer, the decomposition of $E(G)$
given by Theorem~\ref{s6:thmMIXC} has at most
\[
\Bigl(\frac{\Dstar}{2}+1085+\epsOne(\Dstar)\Bigr)n+2\sum_{l=3}^{R}\f(Q_l)
\]
objects, where
\[
\epsOne(\Dstar)\defeq\epsK(\Dstar)+169\,\epsA+\frac{252}{\Dstar}+1.5\,\epsCONC(\Dstar)+3\,\epsX(\Dstar)+\frac{9}{\Dstar},
\]
with $\epsA=31\eps/(\Cp\log\log\Dstar)$ as in Lemma~\ref{s2:lemTower}(e). The
terms are itemized in Table~\ref{s7:tabCost}. The constant is $745+338+2=1085$, and
$\epsOne$ is counted once. Consequently
$\Czero=\Dstar/2+1085+\epsOne(\Dstar)\le\Dstar/2+1091$ under \Gc{4}. No term
involves $\cEG$. The only Markov multipliers are the fixed $3$ (stage 1) and $4$
(per round). There is no term for VX-parts ($\mathcal{V}=\emptyset$).
\deps{\ref{s6:thmMIXC}, \ref{s6:thmCONCL}, \ref{s7:lemLift}, \ref{s7:lemPay}, \ref{s7:lemEXprime}, \ref{s7:lemOneOutcome}, \ref{s5:lemKRED}, \ref{s2:propParentless}, \ref{s2:propOV}, \ref{s7:consRound}, \ref{s2:lemTower}, \ref{s1:condG4}, \ref{s6:defLending}, \ref{s7:defXprime}}
\end{proposition}

\begin{table}[ht]
\centering
\small
\begin{tabular}{|p{5.0cm}|p{6.2cm}|p{3.6cm}|}
\hline
\textbf{Item} & \textbf{Bound on the number of objects} & \textbf{Source}\\
\hline
Long cycles, $E_0$ and light parts (K-RED) & $(\Dstar/2+745+\epsK)n+[80\,\mathrm{dem}+369\,\mathrm{lp}]$ & Lemma~\ref{s5:lemKRED}\\
TPV in standalone pre-parts & $169(2+\epsA)n+[169\sum_l|\Lost_l|]$ & Thm.~\ref{s6:thmMIXC}(c); Props.~\ref{s2:propParentless}, \ref{s2:propOV}\\
JS-LC cycles & $252\,n/\Dstar+1.5\,\epsCONC(\Dstar)\,n$ & Thms.~\ref{s6:thmMIXC}(c), \ref{s6:thmCONCL}\\
$\Jlost$ items & $[\sum_l(M_l-1)|\Lost_l|]$ & Lemma~\ref{s7:lemPay}(a1)\\
Pooled and JV-bad items & $[\Xpool]$ & Lemma~\ref{s7:lemPay}(a2),(a3)\\
Unpaired fresh legs & $2n$ & Lemma~\ref{s7:lemPay}(b)\\
SDR failures and loops & $9n/\Dstar$ & Lemma~\ref{s7:lemPay}(d)\\
Lift of the quotients & $2\sum_l|\Dec_l|=2\sum_l\f(Q_l)$ & Lemma~\ref{s7:lemLift}\\
\hline
All bracketed terms together & $\le\XU+\Xpool\le\Xp\le3\,\Ex\Xp\le3\,\epsX(\Dstar)\,n$ & Lemmas~\ref{s7:lemEXprime}, \ref{s7:lemOneOutcome}\\
\hline
\end{tabular}
\caption{The cost line on the chosen outcome (Proposition~\ref{s7:propCost}).
Summing the rows gives $(\Dstar/2+745+338+2)n+\epsOne(\Dstar)n+2\sum_l\f(Q_l)$.}
\label{s7:tabCost}
\end{table}

\begin{proof}
By Lemma~\ref{s7:lemLift}(iv), the round step is a J-consumer. The outcome of
Lemma~\ref{s7:lemOneOutcome} has stage-3 labels in all good events. So
Theorem~\ref{s6:thmMIXC} applies to the run, $\delta$, the chosen stage-1 and
stage-3 outcomes and this J-consumer, and it yields a decomposition of $E(G)$. By
Theorem~\ref{s6:thmMIXC}(c), its number of objects is at most
\[
\Bigl(\frac{\Dstar}2+745+338\Bigr)n+\eps_M(\Dstar)\,n
+\Bigl[80\,\mathrm{dem}+369\,\mathrm{lp}+169\sum_l|\Lost_l|\Bigr]
+1.5\sum_{(Y,l)}m_{Y,l}+\sum_l|\Obj_l| ,
\]
with $\eps_M(\Dstar)=\epsK(\Dstar)+169\epsA+252/\Dstar$. This itemization
combines the K-RED bound of Lemma~\ref{s5:lemKRED}; the TPV bound
\[
169\sum_{l,Z}\bigl(|A_Z|+|F_Z|+|\Lost_Z|\bigr),
\]
with $\sum|F_Z|\le2n$ (Proposition~\ref{s2:propParentless}(i)) and
$\sum|A_Z|\le\epsA n$ (Proposition~\ref{s2:propOV}(K2) and
Lemma~\ref{s2:lemTower}(e)); and the JS-LC bound. By
Theorem~\ref{s6:thmCONCL}(iv),
$1.5\sum_{(Y,l)}m_{Y,l}\le1.5\,\epsCONC(\Dstar)\,n$.

The objects of $\Obj_l$ are the paid single edges and the lifted objects. By
Lemma~\ref{s7:lemPay}, the paid single edges number at most:
$(M_l-1)|\Lost_l|$ in (a1); the round-$l$ summand of $\Xpool$ in (a2) and (a3);
at most $2n$ unpaired fresh legs over all rounds; and $\payrd_l$ SDR-failure and
loop edges, with $\sum_l\payrd_l\le9n/\Dstar$ by Lemma~\ref{s7:lemPay}(d), since
the outcome satisfies Lemma~\ref{s7:lemOneOutcome}(3). By
Lemma~\ref{s7:lemLift}(iv), the lifted objects number at most $2\f(Q_l)$. Hence
\[
\sum_l|\Obj_l|\le\sum_l(M_l-1)|\Lost_l|+\Xpool+2n+\frac{9n}{\Dstar}+2\sum_l\f(Q_l).
\]
The bracketed terms and the terms $\sum_l(M_l-1)|\Lost_l|$ and $\Xpool$ together
equal
$80\,\mathrm{dem}+369\,\mathrm{lp}+\sum_l(169+M_l-1)|\Lost_l|+\Xpool\le\XU+\Xpool\le\Xp$.
By Lemma~\ref{s7:lemOneOutcome}(1) and Lemma~\ref{s7:lemEXprime},
$\Xp\le3\Ex\Xp\le3\epsX(\Dstar)n$. Adding everything gives at most
\[
\Bigl(\frac{\Dstar}2+745+338+2\Bigr)n
+\Bigl(\epsK+169\epsA+\frac{252}{\Dstar}+1.5\epsCONC+3\epsX+\frac9{\Dstar}\Bigr)n
+2\sum_l\f(Q_l),
\]
which is the claim, since $745+338+2=1085$. Every term of $\epsOne$ is a function
of $\Dstar$ alone, and $\epsOne(\Dstar)\le6$ is part
of \Gc{4}, which gives $\Czero\le\Dstar/2+1091$.
\end{proof}

\begin{theorem}[Theorem JV$^{+*}$]\label{s7:thmJVps}
Assume that $\Dstar$ satisfies \Gc{1}--\Gc{4}. Let $G$ be a graph with
$n\ge\Nzero$ vertices and $d_1\ge\Dstar$. Take any valid $\HBtp$ run on $G$
(Proposition~\ref{s2:propExists}), any designation $\delta$, and no VX-parts
($\mathcal{V}=\emptyset$). Then there are simple graphs $Q_3,\dots,Q_R$ without
isolated vertices such that
\begin{enumerate}
\item[(1)] $\f(G)\le\Czero n+2\sum_{l=3}^{R}\f(Q_l)$, where
$\Czero\defeq\Dstar/2+1085+\epsOne(\Dstar)$;
\item[(2)] $\sum_{l=3}^{R}|V(Q_l)|\le\thetaQ\,n$, where
$\thetaQ=\epsTwo(\Dstar)\le1/4$.
\end{enumerate}
Here $\epsOne$ and $\epsTwo$ are the explicit functions of $\Dstar$ defined in
Proposition~\ref{s7:propCost} and Lemma~\ref{s7:lemUHsplit}. Items at ultra hubs
cost no collision payments. They pay only the standard costs (pooled vertices,
JV-bad ports, SDR failures), and they add quotient vertices only through
Lemma~\ref{s7:lemUHsplit}.
\deps{\ref{s7:propCost}, \ref{s7:lemUHsplit}, \ref{s7:lemSimple}, \ref{s7:lemLift}, \ref{s7:lemOneOutcome}, \ref{s6:thmMIXC}, \ref{s2:propExists}, \ref{s7:lemPay}, \ref{s7:lemUltra}, \ref{s1:condGamma}}
\end{theorem}

\begin{proof}
Take the outcome of Lemma~\ref{s7:lemOneOutcome} and let $Q_3,\dots,Q_R$ be the
quotients built on it. By Lemma~\ref{s7:lemSimple} each $Q_l$ is simple and has
no isolated vertices. By Lemma~\ref{s7:lemLift}(iv), the round step is a
J-consumer, so Theorem~\ref{s6:thmMIXC} gives a decomposition of $E(G)$. By
Proposition~\ref{s7:propCost} it has at most $\Czero n+2\sum_l\f(Q_l)$ objects,
which proves (1), since $\f(G)$ is at most the number of objects of any
decomposition. The outcome satisfies the hypotheses of
Lemma~\ref{s7:lemUHsplit}(iii): $\Xp\le3\Ex\Xp$, and the event of (f) holds at
every round. So Lemma~\ref{s7:lemUHsplit}(iv) gives (2). The statement about
ultra hubs is Lemma~\ref{s7:lemPay}(e) together with Lemma~\ref{s7:lemUltra}.
\end{proof}

%% ----------------------------------------------------------------------
\subsection{The layered quotient induction and the proof of the Main Theorem}

\begin{theorem}[Theorem HI$''$ (layered quotient induction)]\label{s7:thmHI}
Let $C\ge\Dstar/2$ and $\vartheta\in[0,1/2)$ be constants with the
following property. Every graph $G$ without isolated vertices, with $n\ge\Nzero$
vertices and $d_1=2|E(G)|/n\ge\Dstar$, admits finitely many simple graphs
$Q_1,\dots,Q_k$ ($k\ge0$) such that
\[
\f(G)\le C\,n+2\sum_{i=1}^{k}\f(Q_i)\qquad\text{and}\qquad\sum_{i=1}^{k}|V(Q_i)|\le\vartheta\,n .
\]
Then $\f(G)\le c\,|V(G)|$ for every graph $G$, where
$c\defeq\max\bigl(C/(1-2\vartheta),\,\Nzero/2\bigr)$.
\deps{\ref{s1:factAdd}, \ref{s1:defObject}, \ref{s1:defConstants}}
\end{theorem}

\begin{proof}
Suppose the conclusion fails, and let $G$ be a counterexample with
$n\defeq|V(G)|$ minimal, so that $\f(G)>cn$. The graph without vertices has
$\f=0$, so $n\ge1$.

(1) \emph{$G$ has no isolated vertex.} If $v$ were isolated, then
$\f(G-v)=\f(G)$ by Fact~\ref{s1:factAdd}(d), so
$\f(G-v)>cn\ge c(n-1)=c|V(G-v)|$, and $G-v$ would be a smaller counterexample.

(2) \emph{$n\ge\Nzero$.} Otherwise $n-1<\Nzero$, and by Fact~\ref{s1:factAdd}(a),
$\f(G)\le|E(G)|\le n(n-1)/2<(\Nzero/2)\,n\le cn$.

(3) \emph{$d_1\ge\Dstar$.} Otherwise, by Fact~\ref{s1:factAdd}(a),
$\f(G)\le|E(G)|=d_1n/2<(\Dstar/2)\,n\le Cn\le cn$, because
$c\ge C/(1-2\vartheta)\ge C$. This is the trivial bound $\f(G)\le|E(G)|$; no
other estimate is used for graphs of small average degree.

(4) By (1)--(3), the hypothesis applies to $G$ and gives simple graphs
$Q_1,\dots,Q_k$ with $\f(G)\le Cn+2\sum_i\f(Q_i)$ and
$\sum_i|V(Q_i)|\le\vartheta n<n$. Each $Q_i$ has fewer vertices than $G$, so by
the minimality of $n$, $\f(Q_i)\le c|V(Q_i)|$.

(5) Hence $\f(G)\le Cn+2c\sum_i|V(Q_i)|\le Cn+2c\vartheta n
\le c(1-2\vartheta)n+2c\vartheta n=cn$, where we used $C\le c(1-2\vartheta)$.
This contradicts $\f(G)>cn$.

The constants $C$ and $\vartheta$ do not depend on $c$, and minimality is applied
only to the graphs $Q_i$, which have strictly fewer vertices than $G$.
\end{proof}

\begin{lemma}[The galactic conditions are satisfiable]\label{s7:lemGammaSat}
Let $\Nzero$ be as in Definition~\ref{s1:defConstants}(ii). Then:
\begin{enumerate}
\item[(i)] each item (a)--(f) of \Gc{1} (Condition~\ref{s1:condG1}) holds for
all sufficiently large real $\mu$; hence there is $\mu_1\ge1$ such that all of them
hold for every $\mu\ge\mu_1$, and \Gc{1} holds for every $\Dstar>2$ with
$\log_2\log_2\Dstar\ge\mu_1$;
\item[(ii)] $\epsOne(\Dstar)\to0$ and $\epsTwo(\Dstar)\to0$ as $\Dstar\to\infty$;
\item[(iii)] there is $D_0$ such that every $\Dstar\ge D_0$ satisfies
\Gc{1}--\Gc{4} simultaneously.
\end{enumerate}
In particular a constant $\Dstar$ as in Definition~\ref{s1:defConstants}(iii)
exists. Only this existence is used: neither $\mu_1$ nor $D_0$ is computed.
\deps{\ref{s1:condGamma}, \ref{s7:propCost}, \ref{s7:lemUHsplit}, \ref{s5:lemExpect}, \ref{s5:lemKRED}, \ref{s2:lemTower}, \ref{s7:lemEXprime}, \ref{s6:thmCONCL}, \ref{s1:defConstants}, \ref{s3:lemCOL}, \ref{s5:lemE1}, \ref{s3:lemCOLJVev}}
\end{lemma}

\begin{proof}
(i) We use inequalities of type (E) as defined in Lemma~\ref{s3:lemCOLJVev},
which has no hypothesis: by part~(i) of that lemma, every inequality of type (E)
holds for all sufficiently large $\mu$. For $\mu\ge1$ we have
$0\le\log_2\mu\le\log_2(\Aexp\mu)$, and:
item~(a) holds for every $\mu\ge2^8$;
item~(b), that is $\mu-14-\log_2\Aexp-3\log_2\mu\ge0$, follows from the
inequality $\mu-(14+\log_2\Aexp)-3\log_2(\Aexp\mu)\ge0$ of type (E);
item~(c) is the inequality $1.6\mu-2\Aexp\log_2(\Aexp\mu)\ge0$ of type (E);
item~(d), that is $\mu-8-2\log_2\mu\ge0$, follows from the inequality
$\mu-8-2\log_2(\Aexp\mu)\ge0$ of type (E);
item~(e), which is \Gstar, is (after taking $\log_2$ of both sides) the
inequality $36\mu-240-46\Aexp\log_2(\Aexp\mu)\ge0$ of type (E);
and every inequality of item~(f) holds for all sufficiently large $\mu$ by
Lemma~\ref{s3:lemCOLJVev}(iii).
So each of the finitely many inequalities of \Gc{1} holds for all $\mu$ at least
some threshold; let $\mu_1\ge1$ be at least all of these thresholds. If
$\Dstar>2$ and $\log_2\log_2\Dstar\ge\mu_1$, every $\mu\ge\log_2\log_2\Dstar$
satisfies $\mu\ge\mu_1$, so \Gc{1} holds.

(ii) By Proposition~\ref{s7:propCost} and Lemmas~\ref{s7:lemEXprime}
and~\ref{s7:lemUHsplit}, $\epsOne$ and $\epsTwo$ are finite sums of non-negative
explicit functions of $\Dstar$ with fixed coefficients:
$1/\Dstar$; $(6\log_2\Dstar+12)/\Dstar$;
$\epsA=31\eps/(\Cp\log\log\Dstar)$ (Lemma~\ref{s2:lemTower}(e));
$\epsCONC(\Dstar)$; $\epsU(\Dstar)$; $\epsK(\Dstar)$; and, in $\epsTwo$ only,
$F(\log_2\Dstar)$ with $F(x)=(3184+30400\log_2x)\,x^{-90.2}$ as in
Lemma~\ref{s7:lemUHsplit}(iv), which tends to $0$ because $F(x)\to0$ as
$x\to\infty$. The function $\epsCONC$ is the explicit expression
of Theorem~\ref{s6:thmCONCL}(iv). It is a combination of
$\logs\Dstar/\log\Dstar$, $\logs\Dstar/\log\log\Dstar$ and
$(2\logs\Dstar+2)/(\log\Dstar)^{1/2}$, each of which tends to $0$. The functions
$\epsU$ (Lemma~\ref{s5:lemExpect}) and $\epsK$ (Lemma~\ref{s5:lemKRED}) are, by
their defining statements, explicit functions of $\Dstar$ that tend to $0$ as
$\Dstar\to\infty$. Every other listed function tends to $0$. Hence
$\epsOne(\Dstar)\to0$ and $\epsTwo(\Dstar)\to0$.

(iii) Put $D_1\defeq2^{2^{\mu_1}}$; then $D_1\ge4>2$, as $\mu_1\ge1$. By (i),
\Gc{1} holds for every $\Dstar\ge D_1$, since then $\log_2\log_2\Dstar\ge\mu_1$.
\Gc{2}(a) holds for every $\Dstar\ge2^{117}$, and \Gc{2}(b),(c) are implied by
\Gc{1} (Condition~\ref{s1:condGamma}). \Gc{3} holds for every
$\Dstar\ge2^{(2\Nzero)^{1/\Cp}}$; here $\Nzero$ is fixed before $\Dstar$. By (ii)
there is $D''$ with $\epsOne(\Dstar)\le6$ and $\epsTwo(\Dstar)\le1/4$ for all
$\Dstar\ge D''$, which is \Gc{4}. Take
$D_0\defeq\max\bigl(D_1,2^{117},2^{(2\Nzero)^{1/\Cp}},D''\bigr)$.
\end{proof}

\begin{corollary}[Corollary JV$^{+*}$-EG; proof of the Main Theorem]\label{s7:thmMainProof}
Theorem~\ref{s1:thmMain} holds.
\deps{\ref{s7:thmHI}, \ref{s7:thmJVps}, \ref{s2:propExists}, \ref{s7:lemGammaSat}, \ref{s6:defDesign}, \ref{s1:condGamma}, \ref{s1:defConstants}}
\end{corollary}

\begin{proof}
Let $\Dstar$ be as in Theorem~\ref{s1:thmMain}, that is, a constant satisfying
\Gc{1}--\Gc{4}. Such constants exist by Lemma~\ref{s7:lemGammaSat}, and the
argument below applies to each of them. Let
$\Czero=\Dstar/2+1085+\epsOne(\Dstar)$, $\thetaQ=\epsTwo(\Dstar)$ and
$\cEG=\max\bigl(\Czero/(1-2\thetaQ),\Nzero/2\bigr)$, as in
Theorem~\ref{s1:thmMain}. We apply Theorem~\ref{s7:thmHI} with $C\defeq\Czero$
and $\vartheta\defeq\thetaQ$. Its assumptions on the constants hold:
$\Czero\ge\Dstar/2$, since $\epsOne\ge0$; and $0\le\thetaQ\le1/4<1/2$.

Let $G$ be a graph without isolated vertices, with $n\ge\Nzero$ vertices and
$d_1\ge\Dstar$. By Proposition~\ref{s2:propExists}, $G$ has a valid $\HBtp$
run. A designation exists: for every standalone pre-part $Z$ of a round $l\ge3$
and every $u\in Q_Z$, the set $\anc_l(u)$ is non-empty by the definition of
classed ports, and $Y(u)$ can be chosen in it by a fixed rule
(Definition~\ref{s6:defDesign}). With no VX-parts, Theorem~\ref{s7:thmJVps} gives
simple graphs $Q_3,\dots,Q_R$ with $\f(G)\le\Czero n+2\sum_l\f(Q_l)$ and
$\sum_l|V(Q_l)|\le\thetaQ n$. This is the hypothesis of
Theorem~\ref{s7:thmHI}, which therefore gives $\f(G)\le\cEG|V(G)|$ for every
graph $G$. In particular $\f(n)\le\cEG n$ for every $n$, so $\f(n)=O(n)$, where
$\cEG$ depends only on $\sigma,\Cp,\Aexp,\eps,\Nzero$ and $\Dstar$.

This completes the proof of Theorem~\ref{s1:thmMain}.
\end{proof}

\begin{remark}[Why the argument is not circular]\label{s7:remNonCirc}
(1) $\Czero$ and $\thetaQ$ are fixed functions of $\Dstar$
(Proposition~\ref{s7:propCost}, Lemma~\ref{s7:lemUHsplit}), and $\Dstar$ is fixed
before any graph is considered. The constant $\cEG$ enters only as the
constant $c$ of Theorem~\ref{s7:thmHI} (throughout its statement and
proof); it plays no other role in any construction or bound of
Sections~\ref{s2:sec}--\ref{s7:sec}.
(2) Minimality is applied only to the quotients $Q_l$. They are simple and have
at most $\thetaQ n\le n/4<n$ vertices in total, and the induction uses about
them only the inequality $\f(Q_l)\le\cEG|V(Q_l)|$.
(3) The quotient $Q_l$ depends on the decompositions $\Dec_{l'}$ ($l'>l$), through
$J_l$ and through junk. This is harmless. Each $\Dec_{l'}$ is fixed at round $l'$
by a deterministic rule. The Markov choice at round $l$ succeeds with
probability at least $1/2$ whatever the past. Every bound of
Lemmas~\ref{s7:lemCC}, \ref{s7:lemUltra} and~\ref{s7:lemPay} holds for every
past.
(4) No other link calls itself. Lemma~TPV, Lemma~PV and Theorem~VX$^+$
finish their runs with the elementary long-cycle bound of
Fact~\ref{s1:factEG0}, which is proved directly and does not use the statement
being proved; all three are unconditional.
\deps{\ref{s7:thmHI}, \ref{s7:thmJVps}, \ref{s7:propCost}, \ref{s7:lemUHsplit}, \ref{s7:lemCC}, \ref{s7:lemUltra}, \ref{s7:lemPay}, \ref{s7:lemSimple}, \ref{s7:consRound}, \ref{s4:lemTPV}, \ref{s4:lemPV}, \ref{s4:thmVXp}, \ref{s1:condG4}, \ref{s1:factEG0}}
\end{remark}

%% ======================================================================
%% Section 8: Formal verification. Input by ms.tex after \section{The junction-vertex quotient, ultra hubs, and the induction}\label{s7:sec}

This section supplies the J-consumer of Definition~\ref{s6:defJconsumer}, bounds
the size of the quotient graphs it produces, and closes the argument by the
layered quotient induction.

\paragraph{Setting of this section.}
Throughout, $\Dstar$ satisfies \Gc{1}--\Gc{4}, $G$ is a graph with $n\ge\Nzero$
vertices and $d_1\ge\Dstar$, a valid $\HBtp$ run on $G$ is fixed
(Proposition~\ref{s2:propExists}), a designation $\delta$ is fixed
(Definition~\ref{s6:defDesign}), and there are no VX-parts
($\mathcal{V}=\emptyset$), so that every standalone pre-part is a MIX-part. The
rounds of the run are $1,\dots,R$; the round step below acts at the rounds
$3\le l\le R$ (rounds $1$ and $2$ have no classed ports). For a classed port $u$
of round $l$ (that is, $u\in Q_Z$ for some $Z\in\Std_l$) we write $Y(u)$ for its
class and $r(u)$ for the round of $Y(u)$; since $Y(u)\in\anc_l(u)$ we have
$r(u)\le l-2$ and $u\in V(Y(u))$. For an ancestor $Y$ of round $r$, $H_Y$ is its
expander ($X_Y$ if $Y$ is light, $X^0_Y$ if $Y$ is standalone), $s_Y$ is $s_r/2$
or $s_r$ respectively, and $p_Y=1/(2\klend(Y))$ (Definition~\ref{s3:defCOL}).
The functions $\epsOne$ and $\epsTwo$ defined below contain the function
$\epsA=31\eps/(\Cp\log\log\Dstar)$ of Lemma~\ref{s2:lemTower}(e), which depends
on $\Dstar$ only and satisfies $\sum_{l\le R}15.2\eps/\log P_l\le\epsA$; the
function $\epsX$ does not contain $\epsA$.

%% ----------------------------------------------------------------------
\subsection{Round-split pools and candidate junctions}

\begin{definition}[Round-split pool labels]\label{s7:defPool}
Every vertex $v$ of $G$ independently draws a \emph{pool label}
\[
\plab(v)\in\{\bot\}\cup\{(l,r):3\le l\le R,\ 1\le r\le l-2\},\qquad
\Prob\bigl(\plab(v)=(l,r)\bigr)=q_l\,\pi_{l,r},
\]
where $q_l\defeq M_l^{-2}$ and $\pi_{l,r}\defeq 2^{-(l-1-r)}$, and
$\Prob(\plab(v)=\bot)\defeq 1-\sum_{l,r}q_l\pi_{l,r}$. This is a probability
distribution: for each $l$ we have $\sum_{r=1}^{l-2}\pi_{l,r}=1-2^{-(l-2)}<1$,
and $\sum_l q_l\le\sum_l M_l^{-1}\le 2/\Dstar<1$ by
Lemma~\ref{s2:lemTower}(b). Put
\[
\Pool_{l,r}\defeq\plab^{-1}(l,r),\qquad
\Pool_l\defeq\bigcup_{r=1}^{l-2}\Pool_{l,r}\ \ \text{(a disjoint union)},
\]
so that $\Prob(v\in\Pool_l)=q_l(1-2^{-(l-2)})\le q_l$ for every $v$. The sets
$\Pool_l$ ($3\le l\le R$) are pairwise disjoint, because every vertex has one
label. For $w\in\Pool_l$ let $r(w)$ denote the unique $r$ with
$w\in\Pool_{l,r}$. The pool labels are stage (1d) of the randomness schedule.
They are independent of all other stage-1 data, in particular of the colourings
and labels of Definition~\ref{s3:defCOL}.
\deps{\ref{s2:lemTower}, \ref{s3:defCOL}}
\end{definition}

\begin{definition}[Candidates and JV-bad ports]\label{s7:defCand}
Let $u$ be a classed port of round $l\ge3$, with class $Y\defeq Y(u)$ of round
$r\defeq r(u)\le l-2$. Its \emph{candidate set} is
\[
\Cand_l(u)\defeq\{\,w\in\Pool_{l,r(u)}:\ uw\in\LJV_{Y,l}\,\}.
\]
Since $\LJV_{Y,l}\subseteq\Lend_Y\subseteq E(H_Y)$ and $H_Y$ is a graph on
$V(Y)$, automatically $\Cand_l(u)\subseteq N_{H_Y}(u)$ and
$N_{H_Y}(u)\subseteq V(Y)$. The
candidate sets are always taken in the class graph $H_{Y(u)}$ itself; this is
used in Lemma~MULT below. Put
\[
\Hcd_l\defeq\frac{\lambda_{l-2}^{95}}{8M_l^{2}}.
\]
The port $u$ is \emph{JV-bad} if
\[
|\Cand_l(u)|<\tfrac12\,\Ex|\Cand_l(u)|,
\]
and \emph{JV-good} otherwise. Here $\Ex$ is the unconditional expectation over
stage~1, which is a number determined by the run and $\delta$. The definition
applies to every classed port, whether or not it lies in $\Lost_l$. JV-badness
is a deterministic function of stage~1; it is one of the stage-2 statuses. The
number $\Hcd_l$ is not a threshold of the definition: it is a lower bound for
the threshold, since $\tfrac12\Ex|\Cand_l(u)|\ge\Hcd_l$ for every classed port
$u$ of round $l$ (Lemma~\ref{s7:lemCand}(ii) below). Consequently every JV-good
port has $|\Cand_l(u)|\ge\Hcd_l$ (Lemma~\ref{s7:lemCand}(iv)).
\deps{\ref{s7:defPool}, \ref{s3:defCOL}, \ref{s6:defDesign}}
\end{definition}

\begin{lemma}[Candidate counts]\label{s7:lemCand}
Let $u$ be a classed port of round $l\ge3$ with class $Y$ of round $r$. Then:
\begin{enumerate}
\item[(i)] $|\Cand_l(u)|$ is a sum of independent Bernoulli variables, one for
each $w\in N_{H_Y}(u)$, each with success probability $q_l\pi_{l,r}p_Y$;
\item[(ii)] its mean satisfies
\[
\Ex|\Cand_l(u)|=q_l\,\pi_{l,r}\,p_Y\,\deg_{H_Y}(u)
\;\ge\;\frac{\lambda_r^{96}\,2^{-(l-r)}}{M_l^{2}}
\;\ge\;2\Hcd_l ;
\]
\item[(iii)] $\Prob(u\text{ is JV-bad})\le\exp(-\Hcd_l/4)$;
\item[(iv)] if $u$ is JV-good then $|\Cand_l(u)|\ge\Hcd_l\ge2^{10}M_l^{10}$.
\end{enumerate}
\deps{\ref{s7:defCand}, \ref{s7:defPool}, \ref{s3:defCOL}, \ref{s3:lemCOLJV}, \ref{s2:propStructure}, \ref{s2:lemTower}, \ref{s1:citChernoffGen}}
\end{lemma}

\begin{proof}
(i) By definition,
\[
|\Cand_l(u)|=\sum_{w\in N_{H_Y}(u)}\ind{\plab(w)=(l,r)}\cdot\ind{uw\in\LJV_{Y,l}} .
\]
The pool labels are independent across vertices. The colours of the edges of
$H_Y$ are independent across edges (Definition~\ref{s3:defCOL}), and the edges
$uw$ for distinct $w$ are distinct. The two families are independent of each
other (Definition~\ref{s7:defPool}). So the summands for distinct $w$ are
functions of disjoint sets of independent variables, hence independent. For a
fixed $w$ the two indicators are independent, with
$\Prob(\plab(w)=(l,r))=q_l\pi_{l,r}$. Moreover $\Prob(uw\in\LJV_{Y,l})=p_Y$: the
edge $uw\in E(H_Y)$ lies in $\Lend_Y$ with probability $1/2$, and it then
receives the colour $l\in\IJV$ with probability $1/\klend(Y)$; note that
$r+2\le l\le R$, so $l$ is an index of $\IJV$.

(ii) The first equality follows from (i). Every $H_Y$ has minimum degree larger
than $s_Y\ge s_r/2$, and $s_r\ge\lambda_r^{100}$
(Proposition~\ref{s2:propStructure}(i)). Also $p_Y\ge\lambda_r^{-4}$
(Lemma~\ref{s3:lemCOLJV}). Hence
\[
\Ex|\Cand_l(u)|\ \ge\ M_l^{-2}\,2^{-(l-1-r)}\,\lambda_r^{-4}\,\frac{\lambda_r^{100}}{2}
\;=\;\frac{\lambda_r^{96}\,2^{-(l-r)}}{M_l^{2}} .
\]
Since $2\Hcd_l=\lambda_{l-2}^{95}/(4M_l^2)$, the second inequality is equivalent
to
\begin{equation}\label{s7:eqCandExp}
\lambda_r^{96}\ \ge\ 2^{\,l-r-2}\,\lambda_{l-2}^{95}.
\end{equation}
By Lemma~\ref{s2:lemTower}(a), $\lambda_r\ge\lambda_{l-2}$ because $r\le l-2$.
By Lemma~\ref{s2:lemTower}(d),
$\lambda_r\ge2^{2\logs d_r+2}\ge2^{R-r}\ge2^{l-r}\ge2^{l-r-2}$. Multiplying
$\lambda_r^{95}\ge\lambda_{l-2}^{95}$ by $\lambda_r\ge2^{l-r-2}$ gives
\eqref{s7:eqCandExp}. (The chain of exponents is
$2\logs d_r+2\ge R-r\ge l-r$; only the exponent $l-r-2$ is needed.)

(iii) Put $\mu\defeq\Ex|\Cand_l(u)|$. By definition, $u$ is JV-bad exactly
when $|\Cand_l(u)|<\mu/2$. By (i) and the lower-tail Chernoff
bound for sums of independent indicators, with deviation parameter $1/2$
(Cited result~\ref{s1:citChernoffGen}),
$\Prob(|\Cand_l(u)|<\mu/2)\le\exp(-\mu/8)\le\exp(-\Hcd_l/4)$, where the last
step uses $\mu/2\ge\Hcd_l$ from (ii).

(iv) If $u$ is JV-good then $|\Cand_l(u)|\ge\tfrac12\Ex|\Cand_l(u)|\ge\Hcd_l$,
by the definition of JV-good and (ii); this is the first inequality. The second
is the line
``$\Hcd_l\ge2^{10}M_l^{10}$'' of the table of Lemma~\ref{s3:lemCOLJV}. It also
follows directly from $M_l\le\lambda_{l-2}^{1.6}$
(Lemma~\ref{s2:lemTower}(b)): then
$8M_l^2\cdot2^{10}M_l^{10}=2^{13}M_l^{12}\le2^{13}\lambda_{l-2}^{19.2}\le\lambda_{l-2}^{95}$,
because $\lambda_{l-2}\ge M_l^{1/1.6}\ge2^{25}$.
\end{proof}

%% ----------------------------------------------------------------------
\subsection{The randomness schedule and the round step}

\begin{definition}[The randomness schedule]\label{s7:defSchedule}
All random choices of the construction are made in the following order,
summarized in Table~\ref{s7:tabSchedule}.
\begin{itemize}
\item \emph{Stage 1} consists of four mutually independent families:
(1a) the COL-JV colourings of Definition~\ref{s3:defCOL}(i)--(iii);
(1b) the vertex choices and sublabels of Definition~\ref{s5:defZones};
(1c) the JS labels of Definition~\ref{s3:defCOL}(iv);
(1d) the pool labels of Definition~\ref{s7:defPool}.
\item \emph{Selection} (1e), which is not a random family: a stage-1 outcome is
then fixed in the event $\{\Xp\le3\,\Ex\Xp\}$. Here
$\Xp\ge0$ is the stage-1 functional introduced later in this section, a
deterministic function of the run, $\delta$ and stage~1. After (1e), stage~1 is
a fixed outcome: no later step of the construction uses the law of stage~1 or
conditions on the event of (1e); the numbers $\Ex\Xp$ and $\Ex|\Cand_l(u)|$
that appear later are unconditional expectations, fixed in advance.
\item \emph{Stage 2} is deterministic given stage~1: demotion, parent-bad and
$\Bad$ (Definition~\ref{s5:defStages}); lend-bad, $\Lost$, $\Ret$ and $\Qs$
(Definition~\ref{s6:defLending}); JV-bad (Definition~\ref{s7:defCand}).
\item \emph{Stage 3} is drawn part by part: the TPV labels of every standalone
pre-part, the P-vortex labels of every non-demoted light part, and the VX$^+$
labels of every demoted light part. Given stage~1, the label families of distinct
parts are independent. Each part's labels are fixed inside that part's good event.
\item \emph{Rounds} $l=R,R-1,\dots,3$. The round randomness $\xi_l$ consists of
the following mutually independent variables:
a uniformly random bijection $\eta_h:[4M_l]\to[4M_l]$ for every vertex $h$;
a uniformly random $3$-subset $\zeta_{h,u}\subseteq[4M_l]$ for every ordered pair
$(h,u)$ of distinct vertices;
and a uniformly random linear order $\prec_u$ of $V(G)$ for every vertex $u$.
The dependence of $\eta_h,\zeta_{h,u},\prec_u$ on $l$ is suppressed. The draw
$\xi_l$ is fresh: it is independent of stage~1, of stage~3 and of all $\xi_{l'}$
with $l'>l$.
\end{itemize}
The \emph{past} $\Past_l$ at round $l$ is the fixed outcome of stage~1, of
stage~3 and of the draws $\xi_{l'}$ ($l'>l$), together with everything these
determine: in particular the decompositions fixed at the rounds $l'>l$ and the
set $J_l$. It is a fixed outcome, not a random $\sigma$-algebra. We write
$\Ex[\,\cdot\mid\Past_l]$ and $\Prob(\,\cdot\mid\Past_l)$ for expectation and
probability over $\xi_l$ with $\Past_l$ fixed. The fixed rules of
Construction~\ref{s7:consRound} read $\xi_l$ only as stated there. At each round, $\xi_l$ is chosen in
the event of step (f) of the round step. All draws in stage~1, stage~3 and the
rounds are mutually independent. The outcomes actually used are chosen inside
events (stage~1 in $\{\Xp\le3\,\Ex\Xp\}$, stage~3 in the good events, each
$\xi_l$ in the event of step (f)); after each choice only deterministic
properties of the chosen outcome are used, and every later probability is over
fresh variables with the past fixed.
\deps{\ref{s3:defCOL}, \ref{s5:defZones}, \ref{s7:defPool}, \ref{s6:defLending}, \ref{s5:defStages}, \ref{s7:defCand}, \ref{s4:lemTPV}, \ref{s4:lemPV}, \ref{s4:thmVXp}, \ref{s5:lemChild}, \ref{s5:lemDemoted}}
\end{definition}

\begin{table}[ht]
\centering
\small
\begin{tabular}{|p{1.5cm}|p{5.9cm}|p{3.6cm}|p{3.8cm}|}
\hline
\textbf{Stage} & \textbf{Drawn or decided} & \textbf{Depends on} & \textbf{How it is fixed}\\
\hline
1a & COL-JV colourings of every ancestor & independent per ancestor and per edge & part of the stage-1 outcome\\
1b & vertex choices and sublabels (zones) & independent per vertex & part of the stage-1 outcome\\
1c & JS labels $\lab_{Y,l}(y)$ & independent & part of the stage-1 outcome\\
1d & pool labels $\plab(v)$ & independent per vertex & part of the stage-1 outcome\\
1e & choice of the stage-1 outcome & the functional $\Xp$ & $\Xp\le3\,\Ex\Xp$ (probability $\ge2/3$)\\
\hline
2 & demotion, parent-bad, lend-bad, $\Lost$, $\Ret$, $\Qs$, JV-bad & deterministic in stage~1 & determined\\
\hline
3 & TPV, P-vortex and VX$^+$ labels & independent per part & inside each part's good event (probability $\ge1/2-o(1)$)\\
\hline
$l=R,\dots,3$ & $\xi_l=(\eta_h,\zeta_{h,u},\prec_u)$: HUB lists and injections & fresh draw, $\Past_l$ fixed & event of step (f) (probability $\ge1/2$)\\
\hline
\end{tabular}
\caption{The randomness schedule (Definition~\ref{s7:defSchedule}).}
\label{s7:tabSchedule}
\end{table}

\begin{construction}[The JV$^{+*}$ round step at round $l\ge3$; the J-consumer]\label{s7:consRound}
Fix $3\le l\le R$ and the past $\Past_l$. The input is the set
$J_l=\Jlost_l\cup\Jhub_l\cup\Jfr_l\cup\Jpar_l$ produced by JS-LC at round $l$,
with the properties \Jp{1}--\Jp{3} and (i)--(v) of Lemma~\ref{s6:lemJplus}.

\emph{Fixed rules.} The rules and orders called fixed in steps (b)--(e), (g)
and (h) below are deterministic functions, chosen once and for all before any
randomness is drawn, with the following arguments.
\begin{itemize}
\item \emph{Functions of $\Past_l$ alone:} the pairing of the live $\Jfr$ items
at a fresh centre and the choice of its unpaired fresh leg, and the order in
which the PAR objects are coloured, in (b); the split of the live items of a
non-ultra hub into groups, together with the numbering of the groups, in (c).
In particular these rules do not read $\xi_l$.
\item \emph{Functions of $\Past_l$ and of the lists}, that is, of the variables
$\eta_h$ and $\zeta_{h,u}$ of $\xi_l$: the choice of the system of distinct
representatives in (d); the processing order and the order of $V(G)$ in (e1);
and the order $e_1,\dots,e_q$ of $E'(u)$ in (e2).
\item None of the rules of steps (b)--(e) reads the orders $\prec_u$ of
$\xi_l$; within steps (a)--(f), the orders $\prec_u$ of $\xi_l$ enter only
through the assignment $e_i\mapsto w_i$ of (e2) (and through what is computed
from it). Through $\Past_l$ the rules may depend on the draws
$\xi_{l'}$ of the rounds $l'>l$, orders included, and in general they do,
since $J_l$ does.
\item The rank order in (g) and the rule choosing $\Dec_l$ in (h) are not
restricted: they are arbitrary deterministic functions of $\Past_l$ and
$\xi_l$. They necessarily depend on $\xi_l$, since the layers and $Q_l$ do.
\end{itemize}
Rules with these arguments exist, for instance lexicographically first choices
for fixed orders of the finite sets involved (the greedy colouring of (b)
succeeds for every order, by part (i) of the next lemma, and a split into groups
as in (c) always exists). Every statement of this section about the
round step holds for every choice of the fixed rules that obeys these
restrictions. The restrictions are used in the claims of steps (c) and (e2), in
Lemma~\ref{s7:lemWellDef}(iii) (the grouping does not read $\eta_h$), and in
Lemmas~\ref{s7:lemCC}(i), \ref{s7:lemUltra}(i) and \ref{s7:lemPay}(c) (nothing
before the assignment $e_i\mapsto w_i$ reads the orders $\prec_u$), as well as
in the opening paragraph of the subsection ``Quotient size and payments'' and
in the proof of Lemma~\ref{s7:lemUHsplit}(ii) (the copies of non-ultra hubs are
deterministic given the lists).

\emph{Items.} The \emph{items} are the edges of $J_l\setminus\Jlost_l$. The
\emph{vertices} of an item are its port end or ends, which lie in $\Qs_Z$, and
its centre. The centre is a hub $h\in D_l$ for a $\Jhub$ item and a fresh centre
$x\in F_Z$ for a $\Jfr$ item. A $\Jpar$ item $u$--$v$ has two port ends and no
centre. A $\Jhub$ item is written $(h,u)$, with $h$ its hub and $u$ its port.

\smallskip\noindent
(a) \emph{Payments bounded by stage-1 weights.} The following are \emph{paid},
that is, output as single edges:
(a1) every edge of $\Jlost_l$;
(a2) every item with a vertex in $\Pool_l$;
(a3) every item with a JV-bad port end.
The remaining items are \emph{live}. A vertex is \emph{live} if it is a vertex of
a live item, and $\Live$ denotes the set of live vertices. By (a2),
$\Live\cap\Pool_l=\emptyset$. The sets paid in (a1)--(a3) depend on $J_l$, hence
on the past; only their sizes are bounded by stage-1 functionals
(this is proved in the payment lemma below).

\smallskip\noindent
(b) \emph{PAR objects.} Every live $\Jpar$ item $u$--$v$ is a \emph{PAR object}
with \emph{ends} at $u$ and at $v$; here $u,v\in\Qs_Z$ and $Y(u)\ne Y(v)$. At each
fresh centre $x$, the live $\Jfr$ items at $x$ are paired, by a fixed rule, into
\emph{cherries} $u$--$x$--$u'$. Here $u\ne u'$ because $G$ is simple, so
distinct items at $x$ have distinct port ends. A cherry is a PAR object with ends
at $u$ and $u'$ and \emph{middle} $x$. If the number of live $\Jfr$ items at $x$
is odd, one of them, chosen by a fixed rule, is paid: it is the \emph{unpaired
fresh leg} of $x$. The PAR objects are then coloured greedily, in a fixed order,
with colours from $[3M_l]$, so that two PAR objects sharing a port, or sharing a
middle, receive different colours.

\smallskip\noindent
(c) \emph{HUB lists.} For a hub $h\in D_l$ let $\clive_h$ be its number of live
$\Jhub$ items. Put
\[
\thult_l\defeq\bigl\lfloor M_l\Hcd_l/7\bigr\rfloor,\qquad \KHUB_l\defeq4M_l .
\]
The hub $h$ is \emph{ultra} if $\clive_h>\thult_l$, and \emph{non-ultra}
otherwise. For a non-ultra $h$ with $\clive_h\ge1$, its live items are split, by
a fixed rule, into
\[
k_h\defeq\bigl\lceil 8\clive_h/\Hcd_l\bigr\rceil
\]
\emph{groups} of at most $\lceil\Hcd_l/8\rceil$ items each. This is possible
because $k_h\lceil\Hcd_l/8\rceil\ge\clive_h$. The $i$-th group receives the
\emph{list} $\{\eta_h(3i-2),\eta_h(3i-1),\eta_h(3i)\}\subseteq[4M_l]$; this needs
$3k_h\le4M_l$, which is part (ii) of the next lemma. The lists of the
groups of $h$ are pairwise disjoint, and, since the groups and their numbering
are functions of $\Past_l$ while $\eta_h$ is independent of $\Past_l$, they form
a uniformly random sequence of
$k_h$ pairwise disjoint $3$-subsets of $[4M_l]$. Lists of distinct hubs are
independent. For an ultra $h$, every live item $(h,u)$ is its own group, with
list $\zeta_{h,u}$.

\smallskip\noindent
(d) \emph{SDR.} Every port $u$ carrying at least one live hub item chooses, by a
fixed rule, pairwise distinct colours for its live hub items, each item's colour
taken from the list of its group (a system of distinct representatives). If no
such choice exists, all live hub items of $u$ are paid (an \emph{SDR failure} at
$u$). The items that receive a colour are the \emph{coloured hub items}. In
particular, every port carries at most one coloured hub item of each HUB colour.

\smallskip\noindent
(e) \emph{Junctions.} Initially $\Used(u)\defeq\emptyset$ for every port $u$ and
$\Used_\kappa(h)\defeq\emptyset$ for every hub $h$ and every $\kappa\in[4M_l]$.
\begin{itemize}
\item[(e1)] The coloured hub items of non-ultra hubs are processed in a fixed
order. The item $(h,u)$ of colour $\kappa$ receives as its \emph{junction} the
first vertex $w$, in a fixed order of $V(G)$, of
$\Cand_l(u)\setminus\Used(u)\setminus\Used_\kappa(h)$. Then $w$ is added to
$\Used(u)$ and to $\Used_\kappa(h)$.
\item[(e2)] For every port $u$ let $E'(u)$ be the set of the ends at $u$ of the
PAR objects and of the coloured hub items $(h,u)$ with $h$ ultra, listed in a
fixed order $e_1,\dots,e_q$. Let $w_1\prec_u w_2\prec_u\cdots$ be the elements of
$\Cand_l(u)\setminus\Used(u)$ in the order $\prec_u$; the end $e_i$ receives the
junction $w_i$. This is well defined by part (v) of the next lemma. With
$\Past_l$ and the lists fixed, the map $e_i\mapsto w_i$ is a uniformly random
injection of $E'(u)$ into $\Cand_l(u)\setminus\Used(u)$, and these injections are
independent over ports: by the restrictions on the fixed rules, $\Used(u)$, the
set $E'(u)$ and its order $e_1,\dots,e_q$ are determined by $\Past_l$ and the
lists, and the orders $\prec_u$ are independent of both. No payment is made for
coincidences of junctions at ultra hubs.
\item[(e3)] \emph{Loops.} A PAR object whose two ends received the same junction
is paid. A looped $\Jpar$ item costs one single edge; a looped cherry costs its
two edges.
\end{itemize}
Every end of a coloured hub item or of an unpaid PAR object now has a junction
$w$. It is a candidate of the end's port $u$, so $w\in\Cand_l(u)\subseteq\Pool_l$
and $uw\in\LJV_{Y(u),l}$. The edge $uw$ is the \emph{junction edge} of that end.

\smallskip\noindent
(f) \emph{Markov choice.} Let $\payrd_l$ be the number of edges paid in (d) and
(e3), and let $\copies_l\defeq|V(Q_l)|$, with $Q_l$ as in (g). Both are
functions of $\Past_l$ and $\xi_l$. The round randomness $\xi_l$ is chosen in the
event
\[
\bigl\{\copies_l\le4\,\Ex[\copies_l\mid\Past_l]\bigr\}\cap
\bigl\{\payrd_l\le4\,\Ex[\payrd_l\mid\Past_l]\bigr\}.
\]
By Markov's inequality in the form~(b) of Cited result~\ref{s1:citMarkov},
taken conditionally as in its item~(c) ($\xi_l$ is independent of the variables
whose outcome $\Past_l$ fixes, Definition~\ref{s7:defSchedule}), this event has
probability at least $1/2$ given $\Past_l$ (each of the two events fails with
probability at most $1/4$).

\smallskip\noindent
(g) \emph{Layers, sub-layers and the quotient.} For $\kappa\in[3M_l]$ the
\emph{PAR layer} $\BP_\kappa$ is the multigraph on the vertex set
$\{[w]:w\in\Pool_l\}$ of formal copies. It has one edge $[w_1][w_2]$ for every
unpaid PAR object of colour $\kappa$ whose ends have the junctions $w_1$ and
$w_2$; here $w_1\ne w_2$ by (e3). For $\kappa\in[4M_l]$ the \emph{HUB layer}
$\BH_\kappa$ is the bipartite multigraph with sides $D_l$ and
$\{[w]:w\in\Pool_l\}$. It has one edge $h[w]$ for every coloured hub item $(h,u)$
of colour $\kappa$ whose junction is $w$. \emph{Rank split:} in every layer, the
edges joining the same two vertices are listed in a fixed order, and the $i$-th
edge of such a list has \emph{rank} $i$. For a layer $B$ and $\iota\ge1$, the
\emph{sub-layer} $B^{(\iota)}$ is the spanning subgraph of $B$ formed by its edges
of rank $\iota$. The \emph{quotient} $Q_l$ is the vertex-disjoint union of all
sub-layers $B^{\mathrm{P},(\iota)}_\kappa$ ($\kappa\in[3M_l]$, $\iota\ge1$) and
$B^{\mathrm{H},(\iota)}_\kappa$ ($\kappa\in[4M_l]$, $\iota\ge1$), with the
isolated vertices of each sub-layer deleted. Vertices of distinct sub-layers are
distinct vertices of $Q_l$; hence every cycle of $Q_l$ lies in one sub-layer.

\smallskip\noindent
(h) \emph{Decompose and lift.} Let $\Dec_l$ be a decomposition of $E(Q_l)$ into
$\f(Q_l)$ objects, chosen by a fixed deterministic rule (for instance the
lexicographically first one for a fixed order of $E(Q_l)$). Each member of
$\Dec_l$ is \emph{lifted} to objects of $G$ as follows.
\begin{itemize}
\item A single edge of a PAR sub-layer is a PAR object. It is output as its one
($\Jpar$) or two (cherry) J-edges, each as a single edge. A single edge $h[w]$ of
a HUB sub-layer is a coloured hub item $(h,u)$, and it is output as the single
edge $hu$.
\item A cycle $[w_0][w_1]\cdots[w_{m-1}]$ of a PAR sub-layer, with $m\ge3$:
indices are taken mod $m$, and the edge $[w_i][w_{i+1}]$ is a PAR object $o_i$.
Let $p_i$ be the end of $o_i$ with junction $w_i$ and $q_i$ its end with junction
$w_{i+1}$; they are well defined because $w_i\ne w_{i+1}$. The lift is the closed
walk
\[
w_0\,p_0\,(x_0)\,q_0\,w_1\,p_1\,(x_1)\,q_1\,w_2\ \cdots\ q_{m-1}\,w_0 ,
\]
where $x_i$ is the middle of $o_i$ if $o_i$ is a cherry, and is omitted if $o_i$
is a $\Jpar$ item.
\item A cycle $h_0[w_0]h_1[w_1]\cdots h_{m-1}[w_{m-1}]h_0$ of a HUB sub-layer,
with $m\ge2$: the edge $h_i[w_i]$ is a coloured hub item $(h_i,u_i)$ with
junction $w_i$, and the edge $[w_i]h_{i+1}$ is a coloured hub item
$(h_{i+1},u'_i)$ with junction $w_i$ (indices mod $m$). The lift is the closed
walk
\[
h_0\,u_0\,w_0\,u'_0\,h_1\,u_1\,w_1\,u'_1\,h_2\ \cdots\ u'_{m-1}\,h_0 .
\]
\end{itemize}
Put $\Obj_l\defeq{}$(the single edges paid in (a), (b), (d) and (e3)) together
with (the lifted objects of all members of $\Dec_l$), and let $\LentJV_l$ be the
set of junction edges that lie on lifted cycles. A junction edge assigned in (e)
that does not lie on a lifted cycle is not used by the round step; it stays with
its owner as junk.
\deps{\ref{s6:lemJplus}, \ref{s6:defJconsumer}, \ref{s7:defCand}, \ref{s7:defPool}, \ref{s7:defSchedule}, \ref{s1:citMarkov}}
\end{construction}

\begin{lemma}[The round step is well defined]\label{s7:lemWellDef}
For every past $\Past_l$ ($3\le l\le R$) the following hold.
\begin{enumerate}
\item[(i)] \emph{PAR palette.} Every PAR object shares a port or a middle with
fewer than $3M_l$ other PAR objects, so the greedy colouring of (b) succeeds.
Consequently every port carries at most one PAR object of each colour, and
every fresh centre is the middle of at most one cherry of each colour.
\item[(ii)] \emph{Disjoint lists.} Every non-ultra hub $h$ has
$k_h\le\lceil8M_l/7\rceil$ and $3k_h\le4M_l$.
\item[(iii)] \emph{Cuckoo SDR.} For every port $u$,
$\Prob(\text{SDR failure at }u\mid\Past_l)\le(\KHUB_l)^{-4}=(4M_l)^{-4}$.
\item[(iv)] \emph{Greedy step.} Step (e1) always finds a junction. Two coloured
items of the same non-ultra hub $h$ with the same colour $\kappa$ receive
distinct junctions.
\item[(v)] \emph{Injections.} At every port $u$ carrying a live item,
$|\Cand_l(u)\setminus\Used(u)|\ge\Hcd_l-(M_l-1)\ge\Hcd_l/2\ge M_l>|E'(u)|$, so
step (e2) is well defined.
\end{enumerate}
\deps{\ref{s7:consRound}, \ref{s6:lemJplus}, \ref{s7:lemCand}, \ref{s3:lemCOLJV}, \ref{s1:citHall}, \ref{s2:defHBtp}}
\end{lemma}

\begin{proof}
Throughout, $M_l\ge2^{40}$ by the definition of $M_l$, and
$\Hcd_l\ge2^{10}M_l^{10}$ by Lemma~\ref{s7:lemCand}(iv) (equivalently, by the
table of Lemma~\ref{s3:lemCOLJV}). By \Jp{2} of Lemma~\ref{s6:lemJplus}, every
port carries at most $M_l-1$ J-edges, hence at most $M_l-1$ live items. By
Lemma~\ref{s6:lemJplus}(ii), every fresh centre carries at most $M_l-1$ $\Jfr$
edges, hence is the middle of at most $(M_l-1)/2$ cherries.

(i) A PAR object has at most two ports and at most one middle. Each of its ports
carries at most $M_l-2$ further live items, and each further PAR object at that
port uses one of them. Its middle, if any, is the middle of at most
$(M_l-1)/2-1$ further cherries. So the object conflicts with at most
$2(M_l-2)+(M_l-1)/2<3M_l$ others. When it is coloured, fewer than $3M_l$ colours
are therefore blocked, and a colour of $[3M_l]$ remains. The consequences are
the defining property of the colouring.

(ii) If $h$ is non-ultra then $\clive_h\le\thult_l\le M_l\Hcd_l/7$, so
$k_h=\lceil8\clive_h/\Hcd_l\rceil\le\lceil8M_l/7\rceil$. Hence
$3k_h\le3(8M_l/7+1)=24M_l/7+3\le4M_l$, because $M_l\ge21/4$.

(iii) Fix $u$ and let $(h_1,u),\dots,(h_m,u)$ be its live hub items, so that
$m\le M_l-1\le K/4$ with $K\defeq\KHUB_l=4M_l$. These items are distinct edges at
$u$ and $G$ is simple, so the hubs $h_1,\dots,h_m$ are distinct. With $\Past_l$
fixed, it is determined which hubs are ultra, how the items are grouped, and
which group each item belongs to (the grouping rule of step (c) reads $\Past_l$
only, by the restrictions on the fixed rules in
Construction~\ref{s7:consRound}). The list of $(h_i,u)$ is a function of
$\eta_{h_i}$ if $h_i$ is non-ultra, and equals $\zeta_{h_i,u}$ if $h_i$ is ultra.
Since the $h_i$ are distinct, the $m$ lists are independent. Each list is a
uniformly random $3$-subset of $[K]$; for a non-ultra hub this holds because the
sequence of its lists is a uniformly random sequence of pairwise disjoint
$3$-subsets, so each single list is uniform. By Hall's theorem (Cited
result~\ref{s1:citHall}, applied with the $m$ items as the index set and, for
each item $(h_i,u)$, its list as the set of that item), an SDR fails to exist
only if some set of $s$ items has at most $s-1$ colours in the union of their
lists. Since every list has $3$ elements, this forces $s\ge4$. The union is
then contained in some $(s-1)$-subset of $[K]$. The union bound over $s$, over
the $s$-sets of items and over the $(s-1)$-subsets of $[K]$ gives
\[
\Prob(\text{SDR failure at }u\mid\Past_l)\ \le\ \sum_{s=4}^{m}T_s,\qquad
T_s\defeq\binom{m}{s}\binom{K}{s-1}\Bigl(\binom{s-1}{3}\Big/\binom{K}{3}\Bigr)^{s}.
\]
We use $\binom{m}{s}\le(em/s)^s\le(eK/(4s))^s$,
$\binom{K}{s-1}\le(eK/(s-1))^{s-1}$, and
\[
\binom{s-1}{3}\Big/\binom{K}{3}=\frac{(s-1)(s-2)(s-3)}{K(K-1)(K-2)}\le\Bigl(\frac{s-1}{K}\Bigr)^3 ,
\]
which holds because $(s-1-j)/(K-j)\le(s-1)/K$ for $j=1,2$ and $s\le K+1$.
Multiplying, and using $(s-1)^{2s+1}\le s^{2s+1}$,
\[
T_s\ \le\ e^{2s-1}\,4^{-s}\,s^{-s}\,(s-1)^{2s+1}\,K^{-s-1}\ \le\
a_s\defeq\frac1e\cdot\frac sK\cdot\Bigl(\frac{e^2s}{4K}\Bigr)^{s}.
\]
First, $a_4=4e^7K^{-5}<4400\,K^{-5}$. Second, for $s\ge4$,
\[
\frac{a_{s+1}}{a_s}=\frac{s+1}{s}\cdot\frac{e^2(s+1)}{4K}\cdot\Bigl(\frac{s+1}{s}\Bigr)^{s}
\le\frac54\cdot\frac{e^2(s+1)}{4K}\cdot e=\frac{5e^3}{16}\cdot\frac{s+1}{K},
\]
which is at most $\frac{5e^3}{16}(\frac1{13}+\frac1K)<0.49$ when $4\le s<K/13$ and
$K\ge10^4$. Hence $a_s\le2^{4-s}a_4$ for $4\le s\le\lceil K/13\rceil$, and the sum
of these terms is at most $2a_4<8800\,K^{-5}$. Third, for $K/13<s\le m\le K/4$ we
have $s/K\le1/4$, so $a_s\le\frac1{4e}(e^2/16)^s\le0.47^{s}$, and the sum of
these terms is at most $0.47^{K/13}/0.53\le2\cdot0.47^{K/13}$. Altogether
\[
\Prob(\text{SDR failure at }u\mid\Past_l)\ \le\ 8800\,K^{-5}+2\cdot0.47^{K/13}\ \le\ K^{-4}
\qquad(K\ge10^4),
\]
since $8800K^{-5}\le0.88K^{-4}$ and $2\cdot0.47^{K/13}\le0.1K^{-4}$ for
$K\ge10^4$. Here $K=4M_l\ge2^{42}$.

(iv) Items reach (e1) only if they are live, so their ports are JV-good and
$|\Cand_l(u)|\ge\Hcd_l$ (Lemma~\ref{s7:lemCand}(iv)). When the item $(h,u)$ of
colour $\kappa$ is processed, $\Used(u)$ contains one junction for each earlier
item at $u$, so $|\Used(u)|\le M_l-2$. The lists of the groups of $h$ are pairwise
disjoint, so all coloured items of $h$ with colour $\kappa$ lie in the one group
whose list contains $\kappa$. That group has at most $\lceil\Hcd_l/8\rceil$
items, so $|\Used_\kappa(h)|\le\lceil\Hcd_l/8\rceil-1<\Hcd_l/8$. Hence at least
$\Hcd_l-(M_l-2)-\Hcd_l/8>0$ choices remain. The second statement holds because
each junction chosen for $h$ in colour $\kappa$ is added to $\Used_\kappa(h)$
and is excluded afterwards.

(v) After (e1), $\Used(u)$ consists of one junction for each coloured item of a
non-ultra hub at $u$. The elements of $E'(u)$ correspond to further live items at
$u$: a PAR end at $u$ corresponds to the $\Jpar$ item or the leg of the cherry at
$u$, and an ultra hub end is itself an item. So
$|\Used(u)|+|E'(u)|\le M_l-1$. Every port carrying a live item is JV-good, so
$|\Cand_l(u)\setminus\Used(u)|\ge\Hcd_l-(M_l-1)\ge\Hcd_l/2\ge M_l>|E'(u)|$,
using $\Hcd_l\ge2^{10}M_l^{10}\ge2M_l$.

\emph{Remark on numerics (not used).} An exact evaluation of the union bound
$K^5\sum_sT_s$ with $m=M_l-1$ has its maximum, about $0.0502$, near $M_l=19$;
it tends to $9/256$ as $M_l\to\infty$, which is the limit of $K^5T_4$, where
$T_4=\binom{m}{4}/\binom{K}{3}^3\sim9m^4/K^9$. None of these numerical values is used: only the analytic bound (iii) is
used, and it holds since $M_l\ge2^{40}$.
\end{proof}

%% ----------------------------------------------------------------------
\subsection{The quotient is simple, and every decomposition lifts}

For a HUB colour $\kappa$, a hub $h$ and $w\in\Pool_l$, let $m_\kappa(h,w)$ be
the number of edges $h[w]$ of $\BH_\kappa$, that is, the number of coloured hub
items $(h,u)$ of colour $\kappa$ with junction $w$.

\begin{lemma}[Lemma MULT]\label{s7:lemMULT}
Let $3\le l\le R$, $\kappa\in[4M_l]$, let $h$ be a hub and let
$w\in\Pool_{l,r}$. Then
\[
m_\kappa(h,w)\ \le\ \mult_r(w)\ \le\ \mu_r(w).
\]
Moreover, the number of sub-layers $B^{\mathrm{H},(\iota)}_\kappa$ ($\iota\ge1$)
in which $[w]$ is not isolated equals $\max_{h}m_\kappa(h,w)$, and the number of
those in which $h$ is not isolated equals $\max_{w}m_\kappa(h,w)$.
\deps{\ref{s7:consRound}, \ref{s6:lemJplus}, \ref{s7:defCand}, \ref{s7:defPool}, \ref{s2:propOV}}
\end{lemma}

\begin{proof}
Let $(h,u_1),\dots,(h,u_m)$ be the coloured hub items of colour $\kappa$ at $h$
with junction $w$, so $m=m_\kappa(h,w)$. By step (e) of
Construction~\ref{s7:consRound}, the junction of $(h,u_i)$ lies in
$\Cand_l(u_i)\subseteq\Pool_{l,r(u_i)}$ (Definition~\ref{s7:defCand}). Every
vertex has a single pool label (Definition~\ref{s7:defPool}), and
$w\in\Pool_{l,r}$, so $r(u_i)=r$ for all $i$. Moreover
$w\in\Cand_l(u_i)\subseteq N_{H_{Y(u_i)}}(u_i)\subseteq V(Y(u_i))$; this is where
it matters that candidates are taken in the class graph $H_{Y(u_i)}$ itself. So
each $Y(u_i)$ is an ancestor of round $r$ whose vertex set contains $w$. The
items $(h,u_i)$ are distinct $\Jhub$ edges at $h$ of round $l$. By \Jp{1} of
Lemma~\ref{s6:lemJplus}, for each class $Y$ there is at most one $\Jhub$ edge of
class $Y$ at $h$ in round $l$, aggregated over all MIX-parts. So the classes
$Y(u_1),\dots,Y(u_m)$ are pairwise distinct. Hence $m\le\mult_r(w)$, the number
of ancestors of round $r$ containing $w$, and $\mult_r(w)\le\mu_r(w)$ by
Proposition~\ref{s2:propOV}.

For the second statement: $[w]$ is not isolated in $B^{\mathrm{H},(\iota)}_\kappa$
if and only if some edge of rank $\iota$ is incident to $[w]$. By the definition
of ranks, this holds if and only if some pair $\{h,[w]\}$ carries at least
$\iota$ edges, that is, $\max_h m_\kappa(h,w)\ge\iota$. The statement for $h$ is
proved in the same way.
\end{proof}

\begin{lemma}[$Q_l$ is simple]\label{s7:lemSimple}
Every sub-layer of Construction~\ref{s7:consRound}(g) is a simple graph: it has
no parallel edges, PAR sub-layers have no loops, and HUB sub-layers are
bipartite between hubs and junction copies $[w]$, where the hub $h$ and the
junction $w$ of any edge $h[w]$ are distinct vertices of $G$. Consequently $Q_l$
is a simple graph without isolated vertices.
\deps{\ref{s7:consRound}}
\end{lemma}

\begin{proof}
By the definition of ranks, a sub-layer contains at most one edge of rank
$\iota$ between any two vertices, so it has no parallel edges. An edge of a PAR
layer joins $[w_1]$ and $[w_2]$ with $w_1\ne w_2$, because looped PAR objects
were paid in (e3); so PAR sub-layers have no loops. An edge $h[w]$ of a HUB
layer joins the side $D_l$ to the side $\{[w]\}$. Moreover $h$ is the hub of a
coloured, hence live, item, so $h$ is a live vertex and $h\notin\Pool_l$ by (a2),
while $w\in\Pool_l$. So $h\ne w$ as vertices of $G$, and HUB sub-layers have no
loops even when $[w]$ is identified with $w$. So every sub-layer is simple.
$Q_l$ is a vertex-disjoint union of the sub-layers with their isolated vertices
deleted, so it is simple and has no isolated vertices.
\end{proof}

\begin{lemma}[Lemma JV-L: the lift]\label{s7:lemLift}
Let $3\le l\le R$, fix any past, any $\xi_l$, and let $\Live$ be the set of live
vertices of round $l$.
\begin{enumerate}
\item[(i)] $\Live\cap\Pool_l=\emptyset$. Every port carries at most one PAR
object of each PAR colour and at most one coloured hub item of each HUB colour.
Hence, although ports are not vertices of $Q_l$, at most one edge of each
sub-layer is an object or item with an end at a given port. Every fresh centre is the middle of at
most one cherry of each PAR colour. Hubs, fresh centres and ports are pairwise
disjoint.
\item[(ii)] For \emph{every} decomposition $\Dec$ of $E(Q_l)$ into cycles and
single edges, the lift of each cycle of $\Dec$ (Construction~\ref{s7:consRound}(h))
is a cycle of $G$, and the lift of each single edge of $\Dec$ consists of one or
two single edges of $G$. The lifts of distinct members of $\Dec$ are pairwise
edge-disjoint. Hence the lift of $\Dec$ consists of at most $2|\Dec|$ pairwise
edge-disjoint objects of $G$.
\item[(iii)] Every item is either paid or lies in exactly one edge of $Q_l$.
Every junction edge $uw$ is the junction edge of at most one item-end, and that
end is at $u$, never at $w$. The edge $uw$ lies in exactly one lent class, namely
$\LJV_{Y(u),l}$, where $Y(u)$ is a lend-good ancestor of round at most $l-2$; the
only consumer of this class is the round-$l$ step. Junction edges lie in
$E_{r}(Y)$ for ancestors $Y$ of rounds $r\le l-2$, items lie in $E_l(Z)$ for
$Z\in\Std_l$, and these edge sets are disjoint.
\item[(iv)] Consequently the round step of Construction~\ref{s7:consRound},
applied at every round $3\le l\le R$, is a J-consumer:
$\Obj_l$ and $\LentJV_l$ satisfy \JC{1}--\JC{3} of
Definition~\ref{s6:defJconsumer}. The objects of $\Obj_l$ other than paid single
edges number at most $2|\Dec_l|=2\f(Q_l)$.
\end{enumerate}
\deps{\ref{s7:consRound}, \ref{s7:lemSimple}, \ref{s7:lemWellDef}, \ref{s6:lemJplus}, \ref{s6:defJconsumer}, \ref{s3:defCOL}, \ref{s2:propStructure}, \ref{s6:defLending}}
\end{lemma}

\begin{proof}
(i) The first claim is (a2) of Construction~\ref{s7:consRound}. The claims about
ports and fresh centres are Lemma~\ref{s7:lemWellDef}(i) together with the SDR
of step (d). Since every edge of a sub-layer of colour $\kappa$ is a PAR object or
a coloured hub item of colour $\kappa$, at most one edge of each sub-layer is an
object or item with an end at a given port. Disjointness of the three roles is \Jp{3} of
Lemma~\ref{s6:lemJplus}.

(ii) \emph{PAR cycles.} Let $[w_0]\cdots[w_{m-1}]$ be a cycle of a PAR sub-layer
of colour $\kappa$. By Lemma~\ref{s7:lemSimple} the sub-layer is simple, so
$m\ge3$. Distinct vertices of the sub-layer are copies of distinct pooled
vertices, so $w_0,\dots,w_{m-1}$ are distinct. The objects $o_0,\dots,o_{m-1}$
are distinct, because they are distinct edges of the cycle, and all have colour
$\kappa$. By (i) a port carries at most one object of colour $\kappa$, so ports
of distinct objects are distinct. The two ends of one PAR object lie at distinct
ports ($u\ne v$ for a $\Jpar$ item, $u\ne u'$ for a cherry), so $p_i\ne q_i$.
Hence the $2m$ ports $p_i,q_i$ are pairwise distinct. The middles $x_i$ of the
cherries among the $o_i$ are distinct, because a fresh centre is the middle of at
most one cherry of colour $\kappa$. Ports lie in $\Qs_Z$ and middles in $F_Z$, so
by (i) no port is a middle. Ports and middles are live, junctions are pooled,
and $\Live\cap\Pool_l=\emptyset$. So all vertices of the closed walk
$w_0p_0(x_0)q_0w_1\cdots q_{m-1}w_0$ are pairwise distinct. Consecutive vertices
are adjacent in $G$: $w_ip_i$ and $q_iw_{i+1}$ are junction edges (edges of
$H_{Y(p_i)}$ and $H_{Y(q_i)}$), and $p_iq_i$ or $p_ix_i,x_iq_i$ are the J-edges
of $o_i$. The walk has length at least $3m\ge9$. A closed walk of length at
least $3$ with pairwise distinct vertices is a cycle of $G$.

\emph{HUB cycles.} Let $h_0[w_0]h_1[w_1]\cdots h_{m-1}[w_{m-1}]h_0$ be a cycle of
a HUB sub-layer of colour $\kappa$. The sub-layer is simple and bipartite
(Lemma~\ref{s7:lemSimple}), so the cycle has even length at least $4$, and
$m\ge2$. The hubs $h_i$ are pairwise distinct, and so are the junctions $w_i$.
The $2m$ items $(h_i,u_i)$ and $(h_{i+1},u'_i)$ are distinct, because they are
distinct edges of the cycle, and all have colour $\kappa$. A port carries at most
one coloured hub item of colour $\kappa$, so the $2m$ ports $u_i,u'_i$ are
pairwise distinct; in particular $u_i\ne u'_i$. Hubs and ports are disjoint by
(i), and hubs and ports are live while junctions are pooled. So all vertices of
the closed walk $h_0u_0w_0u'_0h_1\cdots u'_{m-1}h_0$ are distinct. Consecutive
vertices are adjacent: $h_iu_i$ and $u'_ih_{i+1}$ are items, and $u_iw_i$,
$w_iu'_i$ are junction edges. The length is $4m\ge8$, so the lift is a cycle of
$G$.

\emph{Single edges.} A single edge of a PAR sub-layer is lifted to the one or two
J-edges of its PAR object, and a single edge of a HUB sub-layer to one item. So
every member of $\Dec$ lifts to at most two objects.

\emph{Edge-disjointness.} Every live item that is not paid is a coloured hub
item or belongs to exactly one unpaid PAR object, so it lies in exactly one edge
of $Q_l$. Each edge of $Q_l$ lies in exactly one member of $\Dec$. So the item
edges used by the lifts of distinct members are distinct. A junction edge $uw$
on a lifted cycle is the junction edge of an end at $u$ with junction $w$. The
ends at $u$ have pairwise distinct junctions: in (e1) every chosen junction is
added to $\Used(u)$ and is excluded later, and in (e2) the map is injective into
$\Cand_l(u)\setminus\Used(u)$. So $uw$ determines the end at $u$, hence the
item, the edge of $Q_l$ and the member of $\Dec$. It cannot also be the junction
edge of an end at $w$: $w$ is pooled, so every item with vertex $w$ was paid in
(a2), and no end lies at $w$. Finally, junction edges are not items: a junction
edge lies in $E(H_Y)\subseteq E_r(Y)$ for an ancestor $Y$ of round $r\le l-2$,
an item lies in $E_l(Z)$ with $Z\in\Std_l$, and these edge sets are disjoint
(Proposition~\ref{s2:propStructure}(iii); Lemma~\ref{s6:lemJplus}(iii)).

(iii) The first two claims were shown in the proof of (ii); every item that is
not in an edge of $Q_l$ was paid in (a2), (a3), (b), (d) or (e3). Let $uw$ be a
junction edge of an end at $u$. Then $w\in\Cand_l(u)$, so
$uw\in\LJV_{Y(u),l}\subseteq E(H_{Y(u)})$. The graphs $H_Y$ of distinct ancestors
are edge-disjoint, since $E(H_Y)\subseteq E_r(Y)$ and these sets are pairwise
disjoint (Proposition~\ref{s2:propStructure}(iii)). Inside $\Lend_{Y(u)}$ the
edge has exactly one colour. So $uw$ lies in exactly one lent class, namely
$\LJV_{Y(u),l}$. By COL(g) of Definition~\ref{s3:defCOL}, the only consumer of
this class is the round-$l$ J-consumer. The port $u$ carries an end, so it is
live, and hence $u\in\Qs_Z=Q_Z\setminus\Lost_Z$. Since $\Lost_Z$ contains every
classed port whose class is lend-bad, $Y(u)$ is lend-good. Finally
$r(u)\le l-2$. The disjointness claim was shown in (ii).

(iv) \JC{1}: $\LentJV_l$ consists of junction edges on lifted cycles, and by
(iii) each lies in $\LJV_{Y,l}$ for a lend-good ancestor $Y$ of round at most
$l-2$. \JC{2}: the paid single edges are J-edges that lie in no edge of $Q_l$;
they are distinct from each other and from all edges used by the lifts, which
consist of items in edges of $Q_l$ and of junction edges. Together with (ii), the
objects of $\Obj_l$ are pairwise edge-disjoint cycles and single edges. Their
union contains every edge of $\Jlost_l$ (paid in (a1)) and every item: each item
is paid or lies in an edge of $Q_l$, whose member of $\Dec_l$ lifts to objects
containing all items of that edge. The only other edges used are the junction
edges on lifted cycles, which form $\LentJV_l$ by definition. So the union is
exactly $J_l\cup\LentJV_l$. \JC{3} holds because no other edge is used. The
bound $2|\Dec_l|=2\f(Q_l)$ is (ii) for $\Dec=\Dec_l$. Junction edges that were
assigned but do not lie on a lifted cycle are not in $\LentJV_l$; they stay in
$E_r(Y)\setminus\Lent(Y)$ as junk for their owner.
\end{proof}

%% ----------------------------------------------------------------------
\subsection{Quotient size and payments}

In this subsection $3\le l\le R$ and the past $\Past_l$ are fixed. We say that
\emph{the lists are fixed} when, in addition, the variables $\eta_h$ and
$\zeta_{h,u}$ of $\xi_l$ are fixed. Then steps (a)--(d) and (e1) of
Construction~\ref{s7:consRound}, the sets $E'(u)$ and their orders
$e_1,\dots,e_q$ are determined (by the restrictions on the fixed rules), and
only the orders $\prec_u$ are random. By step (e2), the junctions of the ends at distinct ports are then
independent, and the junction of an end in $E'(u)$ is uniform on
$\Cand_l(u)\setminus\Used(u)$, a set of size at least $\Hcd_l/2$
(Lemma~\ref{s7:lemWellDef}(v)). By \Jp{2} of Lemma~\ref{s6:lemJplus},
$|J_l|\le n(M_l-1)$, hence $|J_l|\le1.37n(M_l-1)$; in this subsection we
deliberately use the looser form (it only enlarges constants). Hence the number of PAR
objects and the number of hub items are each at most $1.37nM_l$. For a PAR colour $\kappa$ and $w\ne w'$ in
$\Pool_l$, let $m_\kappa(w,w')$ be the number of edges $[w][w']$ of $\BP_\kappa$.
As in Lemma~\ref{s7:lemMULT}, $[w]$ is not isolated in exactly
$\max_{w'}m_\kappa(w,w')$ of the sub-layers $B^{\mathrm{P},(\iota)}_\kappa$.

\begin{lemma}[Lemma CC: PAR copies]\label{s7:lemCC}
Let the past and the lists be fixed. Fix a PAR colour $\kappa$ and $w\in\Pool_l$.
Condition further on the indicators
\[
I_x\defeq\ind{\text{the end at $x$ of the $\kappa$-object at $x$ receives the junction }w}
\]
for every port $x$ carrying a PAR object of colour $\kappa$. Let $S_w$ be the
set of PAR objects of colour $\kappa$ with exactly one end whose junction is $w$.
\begin{enumerate}
\item[(i)] The partner ports $v_o$ of the objects $o\in S_w$ (the ports of the
ends whose junction is not $w$) are pairwise distinct. Under the conditioning,
the junctions of the partner ends are independent, and the junction of the
partner end of $o$ is uniform on $(\Cand_l(v_o)\setminus\Used(v_o))\setminus\{w\}$.
In particular it equals a given $w'$ with probability at most $3/\Hcd_l$.
\item[(ii)] For every integer $t\ge1$,
\[
\Ex\Bigl[\max_{w'}m_\kappa(w,w')\ \Big|\ \Past_l,\text{lists},(I_x)_x\Bigr]
\ \le\ t\cdot\ind{S_w\ne\emptyset}+\frac{6e|S_w|}{\Hcd_l}+4e\,2^{-t}|S_w| .
\]
\item[(iii)] With $t\defeq\tCC_l\defeq\lceil2\log_2M_l\rceil$, the number of
vertices of $Q_l$ lying in PAR sub-layers satisfies
\[
\Ex\bigl[\#\{\text{PAR copies}\}\ \big|\ \Past_l,\text{lists}\bigr]\ \le\
3M_l(\tCC_l+1)|\Pool_l|+44.7\,\frac{nM_l}{\Hcd_l}+29.8\,\frac{n}{M_l}.
\]
\end{enumerate}
These bounds hold for every past, in particular for every realized $J_l$, and for
every outcome of the lists, the SDR and the greedy step (e1): only the orders
$\prec_u$ are random.
\deps{\ref{s7:consRound}, \ref{s7:lemWellDef}, \ref{s1:citChernoffGen}, \ref{s6:lemJplus}}
\end{lemma}

\begin{proof}
The edges at $[w]$ in $\BP_\kappa$ are exactly the objects in $S_w$. Indeed, an
object with both ends at junction $w$ is looped and was paid in (e3), and an
object with no end at junction $w$ is not incident to $[w]$. So
$m_\kappa(w,w')$ is the number of $o\in S_w$ whose partner end has junction
$w'$.

(i) A PAR object has its two ends at distinct ports, and every port carries at
most one PAR object of colour $\kappa$ (Lemma~\ref{s7:lemWellDef}(i)). So every
such port $x$ carries exactly one end of colour $\kappa$, the indicator $I_x$ is
well defined, and the partner ports of distinct objects of $S_w$ are distinct.
The set $S_w$ and the partner ports are determined by the indicators. With the
past and the lists fixed, $\Cand_l(x)\setminus\Used(x)$, $E'(x)$ and its order
are fixed for every port $x$ (restrictions on the fixed rules in
Construction~\ref{s7:consRound}), so the junction of the $\kappa$-end at a
port $x$ is a function of $\prec_x$ alone, and the orders of distinct ports are independent.
So conditioning on $(I_x)_x$ keeps the partner ends independent, and it
conditions the partner end at $v_o$ only on the event $I_{v_o}=0$. The junction
of that end is uniform on $C(v_o)\defeq\Cand_l(v_o)\setminus\Used(v_o)$, since
it is the image of a fixed element under a uniform injection. Given $I_{v_o}=0$
it is therefore uniform on $C(v_o)\setminus\{w\}$. By
Lemma~\ref{s7:lemWellDef}(v), $|C(v_o)\setminus\{w\}|\ge\Hcd_l/2-1\ge\Hcd_l/3$.

(ii) If $S_w=\emptyset$ the left side is $0$. Otherwise, by (i), each
$m_\kappa(w,w')$ is, under the conditioning, a sum of independent indicators
with mean $\mu_{w'}\le3|S_w|/\Hcd_l$, and $\sum_{w'}\mu_{w'}=|S_w|$. Put
$T\defeq\max(t,\lceil6e|S_w|/\Hcd_l\rceil)$. For every integer $j\ge T$ we have
$j\ge2e\mu_{w'}$. By Cited result~\ref{s1:citChernoffGen},
$\Prob(m_\kappa(w,w')\ge j)\le(e\mu_{w'}/j)^j\le(e\mu_{w'}/j)\,2^{1-j}$. Hence,
writing $m$ for $m_\kappa(w,w')$,
\begin{align*}
\Ex\bigl[m\,\ind{m\ge T}\bigr]
&=T\,\Prob(m\ge T)+\sum_{j>T}\Prob(m\ge j)
\le e\mu_{w'}2^{1-T}+\sum_{j>T}\frac{e\mu_{w'}}{j}2^{1-j}\\
&\le e\mu_{w'}\,2^{2-T}\le e\mu_{w'}\,2^{2-t}.
\end{align*}
Moreover
\[
\max_{w'}m_\kappa(w,w')\le(T-1)+\sum_{w'}m_\kappa(w,w')\,\ind{m_\kappa(w,w')\ge T}:
\]
if the maximum is below $T$ it is at most $T-1$, and otherwise it is one of the
summands. Taking conditional expectations,
\[
\Ex\Bigl[\max_{w'}m_\kappa(w,w')\ \Big|\ \cdot\Bigr]\le T-1+4e\,2^{-t}|S_w|
\le t+\frac{6e|S_w|}{\Hcd_l}+4e\,2^{-t}|S_w|,
\]
because $T-1\le\max(t-1,6e|S_w|/\Hcd_l)$.

(iii) The number of PAR copies is
$\sum_{\kappa\in[3M_l]}\sum_{w\in\Pool_l}\max_{w'}m_\kappa(w,w')$. Averaging (ii)
over the indicators and summing over the at most $3M_l|\Pool_l|$ pairs
$(\kappa,w)$ gives at most
$3M_l\,t\,|\Pool_l|+\bigl(6e/\Hcd_l+4e\,2^{-t}\bigr)\,\Ex\sum_{\kappa,w}|S_w|$.
Every PAR object lies in at most two of the sets $S_w$ (one for each end), so
$\sum_{\kappa,w}|S_w|\le2\cdot1.37nM_l=2.74\,nM_l$. With $t=\tCC_l$ we have
$2^{-t}\le M_l^{-2}$, and $6e\cdot2.74<44.7$ and $4e\cdot2.74<29.8$. Finally
$t\le t+1$.
\end{proof}

\begin{lemma}[Ultra-hub copies]\label{s7:lemUltra}
Let the past and the lists be fixed.
\begin{enumerate}
\item[(i)] Let $h$ be an ultra hub, $\kappa\in[4M_l]$, and let
$\omega_1,\dots,\omega_N$ be the junctions of the coloured hub items of $h$ of
colour $\kappa$. These items lie at $N$ distinct ports, because distinct items
at $h$ are distinct edges $hu$ and $G$ is simple. The $\omega_i$ are independent,
and each is uniform on a set of size at least $\Hcd_l/2$.
\item[(ii)] If $N\ge1$ then, with $\mu^*\defeq2N/\Hcd_l$,
$\Ex[\max_w m_\kappa(h,w)\mid\Past_l,\text{lists}]\le\log_2N+2e\mu^*+8$.
\item[(iii)] The number of vertices of HUB sub-layers that are copies of $h$
satisfies
$\Ex[\,\cdot\mid\Past_l,\text{lists}]\le4M_l(\log_2\clive_h+8)+4e\,\clive_h/\Hcd_l$.
\item[(iv)] Summed over all ultra hubs, the expected number of hub copies of
ultra hubs, given $\Past_l$ and the lists, is at most
\[
320\,\frac{nM_l^{3}\,(\log_2\thult_l+8)}{\lambda_{l-2}^{95}} .
\]
\end{enumerate}
The bound depends on the past only through deterministic counts, so it holds for
every past and every realized $J_l$.
\deps{\ref{s7:consRound}, \ref{s7:lemWellDef}, \ref{s6:lemJplus}, \ref{s1:citChernoffGen}, \ref{s7:lemMULT}, \ref{s7:lemCand}}
\end{lemma}

\begin{proof}
(i) Items of ultra hubs receive their junctions in (e2). The junction of the
item $(h,u_i)$ is the image of its end under the uniformly random injection at
$u_i$, so it is uniform on $\Cand_l(u_i)\setminus\Used(u_i)$, a set of size at
least $\Hcd_l/2$ by Lemma~\ref{s7:lemWellDef}(v). The ports $u_i$ are distinct,
and injections at distinct ports are independent, so the $\omega_i$ are
independent. (Distinctness of the ports comes from the simplicity of $G$; it
does not use \Jp{1}.)

(ii) Let $m_w\defeq\#\{i:\omega_i=w\}$. It is a sum of independent indicators
with mean $\mu_w\le2N/\Hcd_l=\mu^*$, and $\sum_w\mu_w=N$. Put
$T\defeq\max(\lceil2e\mu^*\rceil,\lceil\log_2N\rceil+1)$. For $j\ge T$ we have
$e\mu_w/j\le1/2$, so by Cited result~\ref{s1:citChernoffGen}
$\Prob(m_w\ge j)\le(e\mu_w/j)^j\le(e\mu_w/j)2^{1-j}$, and summing over $w$ gives
$\sum_w\Prob(m_w\ge j)\le(eN/j)2^{1-j}$. Hence
\[
\Ex\bigl[\max_wm_w\bigr]=\sum_{j\ge1}\Prob\bigl(\max_wm_w\ge j\bigr)
\le(T-1)+\sum_{j\ge T}\frac{eN}{j}2^{1-j}\le(T-1)+\frac{eN}{T}2^{2-T}
\le T-1+\frac{2e}{T},
\]
using $2^{-T}\le2^{-\log_2N-1}=1/(2N)$. Finally
$T-1\le\max(2e\mu^*,\log_2N+1)\le\log_2N+2e\mu^*+1$ and $2e/T\le2e<7$.

(iii) By Lemma~\ref{s7:lemMULT}, the number of copies of $h$ in the sub-layers of
colour $\kappa$ is $\max_wm_\kappa(h,w)$. Let $N_\kappa$ be the number of
coloured items of $h$ of colour $\kappa$. Then $N_\kappa\le\clive_h$ and
$\sum_\kappa N_\kappa\le\clive_h$, and at most $4M_l$ colours have
$N_\kappa\ge1$. By (ii) the expected number of copies of $h$ is at most
$\sum_{\kappa:N_\kappa\ge1}\bigl(\log_2\clive_h+8\bigr)+2e\sum_\kappa2N_\kappa/\Hcd_l
\le4M_l(\log_2\clive_h+8)+4e\,\clive_h/\Hcd_l$.

(iv) The function $g(x)\defeq(\log_2x+8)/x$ is decreasing on $[1,\infty)$, since
$x^2g'(x)=1/\ln2-8-\log_2x<0$ there. For an ultra hub, $\clive_h>\thult_l\ge1$,
so $\log_2\clive_h+8=\clive_h\,g(\clive_h)\le\clive_h(\log_2\thult_l+8)/\thult_l$.
By \Jp{2} of Lemma~\ref{s6:lemJplus},
$\sum_h\clive_h\le|J_l|\le1.37nM_l$. We have
$M_l\Hcd_l/7=\lambda_{l-2}^{95}/(56M_l)$, and $M_l\Hcd_l\ge2^{10}M_l^{11}$
(Lemma~\ref{s7:lemCand}(iv)), so $\lambda_{l-2}^{95}/M_l\ge2^{453}$. Hence
$\thult_l\ge\lambda_{l-2}^{95}/(56M_l)-1\ge\lambda_{l-2}^{95}/(57M_l)$, and also
$\thult_l\ge2^{8}$. Summing (iii) over ultra hubs,
\[
\sum_{h\text{ ultra}}\Bigl(4M_l(\log_2\clive_h+8)+\frac{4e\,\clive_h}{\Hcd_l}\Bigr)
\le4M_l\cdot1.37nM_l\cdot\frac{57M_l}{\lambda_{l-2}^{95}}(\log_2\thult_l+8)
+4e\cdot1.37nM_l\cdot\frac{8M_l^2}{\lambda_{l-2}^{95}} .
\]
The first term is $312.36\,nM_l^3(\log_2\thult_l+8)/\lambda_{l-2}^{95}$ and the
second is at most $119.2\,nM_l^3/\lambda_{l-2}^{95}$. Since
$\log_2\thult_l+8\ge16$, the second term is at most
$7.5\,nM_l^3(\log_2\thult_l+8)/\lambda_{l-2}^{95}$, and $312.36+7.5<320$.
\end{proof}

\begin{lemma}[The stage-1 weight $\XV$]\label{s7:lemVstar}
For $3\le l\le R$ put
\[
X_{\mathrm{V},l}\defeq3M_l(\tCC_l+1)|\Pool_l|+4M_l\sum_{w\in\Pool_l}\mult_{r(w)}(w),
\qquad \XV\defeq\sum_{l=3}^{R}X_{\mathrm{V},l}.
\]
Each $X_{\mathrm{V},l}$ is a deterministic function of the run and the pool
labels, and
\[
\Ex X_{\mathrm{V},l}\le\frac{(6\log_2M_l+12)\,n}{M_l},\qquad
\Ex\XV\le2.1\,\frac{(6\log_2\Dstar+12)\,n}{\Dstar}.
\]
\deps{\ref{s7:defPool}, \ref{s2:propOV}, \ref{s2:lemTower}}
\end{lemma}

\begin{proof}
The quantities $M_l$, $\tCC_l$ and $\mult_r(w)$ are determined by the run, and
$\Pool_l$ and $r(w)$ by the pool labels. Since $\tCC_l+1\le2\log_2M_l+2$ and
$\Ex|\Pool_l|\le q_ln=n/M_l^2$, the first summand has expectation at most
$(6\log_2M_l+6)n/M_l$. For the second,
\[
\Ex\sum_{w\in\Pool_l}\mult_{r(w)}(w)=\sum_{w}\sum_{r=1}^{l-2}q_l\pi_{l,r}\,\mult_r(w)
=q_l\sum_{r=1}^{l-2}\pi_{l,r}\sum_w\mult_r(w)\le q_l\cdot1.37n ,
\]
since $\sum_w\mult_r(w)\le1.37n$ for every $r$ (Proposition~\ref{s2:propOV}) and
$\sum_r\pi_{l,r}<1$. So the second summand has expectation at most
$4M_l\cdot1.37n/M_l^2=5.48\,n/M_l$. Adding, $\Ex X_{\mathrm{V},l}\le(6\log_2M_l+11.48)n/M_l$.
Summing over $l$ and using $\sum_l\psi(M_l)\le2.1\,\psi(\Dstar)$ for
$\psi(x)=(6\log_2x+12)/x$ (Lemma~\ref{s2:lemTower}(b)) gives the bound for
$\Ex\XV$.
\end{proof}

\begin{lemma}[Payments]\label{s7:lemPay}
For every past and every $3\le l\le R$ the following hold.
\begin{enumerate}
\item[(a1)] $|\Jlost_l|\le(M_l-1)|\Lost_l|$.
\item[(a2)] The number of items with a vertex in $\Pool_l$ is at most
$\sum_{v\in\Pool_l}\omega_l(v)$, where $\omega_l(v)\defeq\cagg_{v,l}$ if
$v\in D_l$ and $\omega_l(v)\defeq M_l$ otherwise.
\item[(a3)] The number of items with a JV-bad port end is at most $M_l-1$ times
the number of JV-bad classed ports of round $l$.
\item[(b)] Every fresh centre has at most one unpaired fresh leg per round, so
the number of unpaired fresh legs over all rounds is at most
$\sum_{l,Z}|F_Z|\le2n$.
\item[(c)] $\Ex[\text{SDR-failure payments}\mid\Past_l]\le n(M_l-1)(4M_l)^{-4}\le n/(256M_l^3)$,
and $\Ex[\text{loop-paid edges}\mid\Past_l]\le2\cdot(2/\Hcd_l)\cdot\#\{\text{PAR objects}\}\le5.5\,nM_l/\Hcd_l$.
\item[(d)] If $\xi_l$ lies in the event of Construction~\ref{s7:consRound}(f),
then $\payrd_l\le4\bigl(n/(256M_l^3)+5.5\,nM_l/\Hcd_l\bigr)\le2n/M_l$; if this
holds at every round, then $\sum_l\payrd_l\le9n/\Dstar$.
\item[(e)] No other item is paid. In particular, no item is paid because of a
coincidence of junctions at an ultra hub, and no item is paid for a concentrated
pair.
\end{enumerate}
\deps{\ref{s7:consRound}, \ref{s7:lemWellDef}, \ref{s6:lemJplus}, \ref{s2:propParentless}, \ref{s2:lemTower}}
\end{lemma}

\begin{proof}
We use the properties of Lemma~\ref{s6:lemJplus}.

(a1) Every $\Jlost$ edge joins a vertex of $\Lost_l$ to a port, and every vertex
carries at most $M_l-1$ J-edges by \Jp{2}.

(a2) Let $v\in\Pool_l$. If $v\in D_l$, then $v$ is a hub and not a port or a
fresh centre, by \Jp{3}. The items with vertex $v$ are therefore $\Jhub$ items
$(v,u)$: $\Jpar$ items have both ends in $\Qs_Z$, and $\Jfr$ items have their
centre in $F_Z$. By \Jp{1} there is at most one such item per class $Y$. Each
such class has $d_{Y,l}(v)\ge1$, because the item is an edge $vu\in E_l(Z)$ with
$u\in Q_Z$ and $Y(u)=Y$. So there are at most $\cagg_{v,l}$ items with vertex
$v$. If $v\notin D_l$, then $v$ carries at most $M_l-1\le M_l$ J-edges by \Jp{2}.

(a3) A port carries at most $M_l-1$ items by \Jp{2}.

(b) The pairing at a fresh centre $x$ leaves at most one live $\Jfr$ item
unpaired. Every fresh centre of round $l$ lies in $F_Z$ for some $Z\in\Std_l$,
and $\sum_{l,Z}|F_Z|\le2n$ by Proposition~\ref{s2:propParentless}(i).

(c) By \Jp{3}, every non-hub vertex has one role in one part per round, so there
are at most $n$ ports. A port $u$ pays at most $M_l-1$ items for an SDR failure,
and it fails with probability at most $(4M_l)^{-4}$ given $\Past_l$
(Lemma~\ref{s7:lemWellDef}(iii)). This gives the first bound. For loops, fix the
lists. A PAR object has its two ends at distinct ports $p\ne q$. Their junctions
are independent, and the junction at $q$ is uniform on a set of size at least
$\Hcd_l/2$ (Lemma~\ref{s7:lemWellDef}(v)). So the probability that the two
junctions coincide is
$\sum_w\Prob(p\to w)\Prob(q\to w)\le\max_w\Prob(q\to w)\le2/\Hcd_l$. A looped
object costs at most two edges (a cherry costs two, a $\Jpar$ item one), and
there are at most $1.37nM_l$ PAR objects (\Jp{2}). So the expected number of
loop-paid edges is at most $4\cdot1.37\,nM_l/\Hcd_l\le5.5\,nM_l/\Hcd_l$. This
bound does not depend on the lists, so it also holds given $\Past_l$ alone.

(d) The first inequality is the definition of the event in (f), together with
(c). Next, $4\cdot5.5\,nM_l/\Hcd_l=176\,nM_l^3/\lambda_{l-2}^{95}$. By
Lemma~\ref{s2:lemTower}(b), $M_l\le\lambda_{l-2}^{1.6}$, so
$\lambda_{l-2}\ge M_l^{1/1.6}\ge2^{25}$ and
$176\,M_l^3/\lambda_{l-2}^{95}\le176\,\lambda_{l-2}^{-90.2}\le\lambda_{l-2}^{-1.6}\le1/M_l$.
Also $4n/(256M_l^3)\le n/M_l$. Hence $\payrd_l\le2n/M_l$, and
$\sum_l\payrd_l\le2n\sum_l1/M_l\le4n/\Dstar\le9n/\Dstar$ by
Lemma~\ref{s2:lemTower}(b).

(e) In Construction~\ref{s7:consRound}, single edges are paid only in (a1),
(a2), (a3), (b), (d) and (e3).
\end{proof}

\begin{definition}[The stage-1 functional $\Xp$]\label{s7:defXprime}
With $\omega_l$ as in Lemma~\ref{s7:lemPay}(a2), put
\[
\Xpool\defeq\sum_{l=3}^{R}\Bigl[\sum_{v\in\Pool_l}\omega_l(v)+(M_l-1)\cdot\#\{\text{JV-bad classed ports of round }l\}\Bigr],
\]
let $\XV$ be as in Lemma~\ref{s7:lemVstar}, and let $\XU$ be as in
Definition~\ref{s6:defLending}. Then
\[
\Xp\defeq\XU+\Xpool+\XV .
\]
Each of the three terms is non-negative and is a deterministic function of the
run, $\delta$ and the stage-1 outcome. Indeed, $\omega_l(v)$ is determined by the
run and $\delta$ (through $\cagg_{v,l}$), $\Pool_l$ is determined by the pool
labels, and JV-badness is determined by stage~1. $\Xp$ is the functional used in
stage (1e) of the schedule.
\deps{\ref{s6:defLending}, \ref{s7:lemVstar}, \ref{s7:lemPay}}
\end{definition}

\begin{lemma}[Expectation of $\Xp$]\label{s7:lemEXprime}
$\Ex\Xp\le\epsX(\Dstar)\,n$, where
\[
\epsX(\Dstar)\defeq\epsU(\Dstar)+\frac{10.3}{\Dstar}+2.1\,\frac{6\log_2\Dstar+12}{\Dstar}
\]
and $\epsU$ is the function of Lemma~\ref{s5:lemExpect} (through
Lemma~\ref{s6:lemLost}). More precisely, $\Ex\XU\le\epsU(\Dstar)n+5.5n/\Dstar$,
$\Ex\Xpool\le4.8n/\Dstar$ and $\Ex\XV\le2.1(6\log_2\Dstar+12)n/\Dstar$.
\deps{\ref{s7:defXprime}, \ref{s6:lemLost}, \ref{s7:lemCand}, \ref{s7:lemVstar}, \ref{s2:propOV}, \ref{s2:propStructure}, \ref{s2:lemTower}, \ref{s5:lemExpect}, \ref{s2:lemCap}}
\end{lemma}

\begin{proof}
The bound on $\Ex\XU$ is Lemma~\ref{s6:lemLost}, and the bound on $\Ex\XV$ is
Lemma~\ref{s7:lemVstar}. For $\Xpool$, fix a round $l$.

\emph{Pooled weights.} The weights $\omega_l(v)$ are determined by the run and
$\delta$, so they are independent of the pool labels. Hence
$\Ex\sum_{v\in\Pool_l}\omega_l(v)\le q_l\sum_v\omega_l(v)\le q_l\bigl(\sum_{h\in D_l}\cagg_{h,l}+nM_l\bigr)$.
Now $\cagg_{h,l}\le\sum_Yd_{Y,l}(h)$, and $\sum_Yd_{Y,l}(h)$ counts the edges
$hu\in E_l(Z)$ with $Z\in\Std_l$ and $u\in Q_Z$. Summing over $h$ counts each such
edge at most once for each of its classed-port ends. So
$\sum_h\cagg_{h,l}\le\sum_{Z\in\Std_l}\sum_{u\in Q_Z}\deg_{E_l(Z)}(u)$. Every
vertex has $E_l(Z)$-degree at most $M_l-1$ (Lemma~\ref{s2:lemCap}(ii)), and
$\sum_{Z\in\Std_l}|Z^0|\le1.37n$ by Proposition~\ref{s2:propOV}. Hence
$\sum_{h\in D_l}\cagg_{h,l}\le1.37nM_l$, and the expected pooled weight is at
most $2.37\,n/M_l$.

\emph{JV-bad ports.} There are at most $n$ classed ports of round $l$, because
$Q_Z\subseteq U_Z$ and the port sets $U_Z$ of distinct $Z\in\Std_l$ are pairwise
disjoint (Proposition~\ref{s2:propStructure}(iv)); this is the count used in
the proofs of \ref{s6:K6} of Lemma~\ref{s6:lemLost} and of \Jp{2} of
Lemma~\ref{s6:lemJplus}. By Lemma~\ref{s7:lemCand}(iii) each is JV-bad with
probability at most $\exp(-\Hcd_l/4)$, and $\exp(-\Hcd_l/4)\le M_l^{-3}$ because
$\Hcd_l\ge2^{10}M_l^{10}$ (Lemma~\ref{s7:lemCand}(iv)). So the expected JV-bad
weight is at most
\[
(M_l-1)\cdot n\cdot M_l^{-3}\ \le\ n/M_l^2\ \le\ 0.03\,n/M_l .
\]
The pooled-weight bound above counts the same classed ports by the
deliberately looser $\sum_{Z\in\Std_l}|Z^0|\le1.37n$; this only enlarges its
constant.

Hence $\Ex\Xpool\le\sum_l2.4\,n/M_l\le4.8\,n/\Dstar$, using
$\sum_l1/M_l\le2/\Dstar$ (Lemma~\ref{s2:lemTower}(b)). Adding the three bounds
gives $\Ex\Xp\le\epsU(\Dstar)n+(5.5+4.8)n/\Dstar+2.1(6\log_2\Dstar+12)n/\Dstar=\epsX(\Dstar)n$.
\end{proof}

\begin{lemma}[Lemma UH$^*$-split: quotient size]\label{s7:lemUHsplit}
For $3\le l\le R$ put
\[
\dett_l\defeq3|D_l|+78\,\frac{nM_l}{\Hcd_l}+30\,\frac{n}{M_l}
+320\,\frac{nM_l^{3}(\log_2\thult_l+8)}{\lambda_{l-2}^{95}},
\]
a deterministic function of the run.
\begin{enumerate}
\item[(i)] $X_{\mathrm{V},l}$ is the stage-1 weight of Lemma~\ref{s7:lemVstar}.
\item[(ii)] For every past, $\Ex[\copies_l\mid\Past_l]\le X_{\mathrm{V},l}+\dett_l$.
\item[(iii)] If the stage-1 outcome satisfies $\Xp\le3\,\Ex\Xp$ and at every round
$\xi_l$ lies in the event of Construction~\ref{s7:consRound}(f), then
\[
\sum_{l=3}^{R}|V(Q_l)|\ \le\ 12\,\Ex\Xp+4\sum_{l=3}^{R}\dett_l .
\]
\item[(iv)] Put $F(x)\defeq(3184+30400\log_2x)\,x^{-90.2}$ for $x\ge1$.
Then $\sum_l\dett_l\le3\epsA n+60n/\Dstar+2F(\log_2\Dstar)\,n$, and
$2F(\log_2\Dstar)<2^{-1000}$. Hence, in the situation of (iii),
\[
\sum_{l=3}^{R}|V(Q_l)|\le\thetaQ\,n,\qquad
\thetaQ\defeq\epsTwo(\Dstar)\defeq12\,\epsX(\Dstar)+4\Bigl(3\epsA+\frac{60}{\Dstar}+2F(\log_2\Dstar)\Bigr),
\]
with $\epsA=31\eps/(\Cp\log\log\Dstar)$ as in Lemma~\ref{s2:lemTower}(e). So
$\thetaQ$ depends on $\Dstar$ only, and $\thetaQ\le1/4$ by
condition~\ref{s1:condG4}.
\end{enumerate}
\deps{\ref{s7:lemCC}, \ref{s7:lemMULT}, \ref{s7:lemUltra}, \ref{s7:lemVstar}, \ref{s7:lemEXprime}, \ref{s7:defSchedule}, \ref{s2:propOV}, \ref{s2:lemTower}, \ref{s2:lemLacunary}, \ref{s1:condG4}, \ref{s1:condG1}, \ref{s7:consRound}, \ref{s7:lemWellDef}, \ref{s2:defHBtp}, \ref{s6:lemJplus}}
\end{lemma}

\begin{proof}
(ii) Fix the past and the lists. The vertices of $Q_l$ are of four kinds.

\emph{PAR copies.} Their conditional expectation is at most
$3M_l(\tCC_l+1)|\Pool_l|+44.7\,nM_l/\Hcd_l+29.8\,n/M_l$ by
Lemma~\ref{s7:lemCC}(iii).

\emph{Copies of non-ultra hubs} (deterministic given the lists). Let $h$ be
non-ultra. By Lemma~\ref{s7:lemWellDef}(iv) (see step (e1) of
Construction~\ref{s7:consRound}), $m_\kappa(h,w)\le1$ for all $\kappa$ and $w$.
So by Lemma~\ref{s7:lemMULT}, $h$ has exactly one copy in each colour that one of
its items uses, and none in the other colours. The colours of its items lie in
the union of its $k_h$ lists, which has $3k_h$ elements. So the number of copies
of non-ultra hubs is at most
$\sum_{h:\clive_h\ge1}3k_h\le\sum_{h\in D_l}3(1+8\clive_h/\Hcd_l)\le3|D_l|+24\cdot1.37\,nM_l/\Hcd_l$,
since $\sum_h\clive_h\le1.37nM_l$. Here $24\cdot1.37=32.88<32.9$.

\emph{Junction copies in HUB sub-layers} (the count depends on the junctions,
but the bound below is deterministic). By
Lemma~\ref{s7:lemMULT}, in colour $\kappa$ the vertex $[w]$ has
$\max_hm_\kappa(h,w)\le\mult_{r(w)}(w)$ copies. Summing over the $4M_l$ HUB
colours and $w\in\Pool_l$ gives at most $4M_l\sum_{w\in\Pool_l}\mult_{r(w)}(w)$.

\emph{Copies of ultra hubs.} Their conditional expectation is at most
$320\,nM_l^3(\log_2\thult_l+8)/\lambda_{l-2}^{95}$ by Lemma~\ref{s7:lemUltra}(iv).

The first summands of the PAR bound and of the junction bound add up to
$X_{\mathrm{V},l}$. The remaining terms add up to at most $\dett_l$, because
$44.7+32.9\le78$ and $29.8\le30$. Averaging over the lists gives (ii).

(iii) By the choice in (f) and by (ii),
$|V(Q_l)|=\copies_l\le4\,\Ex[\copies_l\mid\Past_l]\le4(X_{\mathrm{V},l}+\dett_l)$
at every round. Summing over $l$,
$\sum_l|V(Q_l)|\le4\XV+4\sum_l\dett_l\le4\Xp+4\sum_l\dett_l\le12\,\Ex\Xp+4\sum_l\dett_l$,
because $\XU,\Xpool\ge0$ and $\Xp\le3\Ex\Xp$ (stage (1e) of
Definition~\ref{s7:defSchedule}).

(iv) By Proposition~\ref{s2:propOV}(K2), $|D_l|\le7.6\eps n/\log P_l$. So
$\sum_l3|D_l|\le3n\sum_l7.6\eps/\log P_l\le1.5\epsA n\le3\epsA n$ by
Lemma~\ref{s2:lemTower}(e) (the factor~$2$ of slack is deliberate; it only
enlarges $\epsTwo$). Next,
$\sum_l30n/M_l\le60n/\Dstar$ by Lemma~\ref{s2:lemTower}(b). For the remaining
terms, $78\,nM_l/\Hcd_l=624\,nM_l^3/\lambda_{l-2}^{95}$, and
$\log_2\thult_l\le\log_2(M_l\Hcd_l/7)\le95\log_2\lambda_{l-2}$. With
$M_l\le\lambda_{l-2}^{1.6}$ (Lemma~\ref{s2:lemTower}(b)), these terms are at most
$nF(\lambda_{l-2})$, with $F$ as in (iv).
Here $\lambda_{l-2}\ge M_l^{1/1.6}\ge2^{25}$. On $[2^{25},\infty)$, $F$ is
decreasing, and $F(2x)\le F(x)/2$, since the first factor at most doubles while
$x^{-90.2}$ is multiplied by $2^{-90.2}$. By Lemma~\ref{s2:lemLacunary}(ii),
$\lambda_{r}\ge2^{\lambda_{r+1}/A}\ge2\lambda_{r+1}$ for consecutive rounds,
because $\lambda_{r+1}\ge2^{25}$. So $X_l\defeq F(\lambda_{l-2})$
($3\le l\le R$) satisfies $X_l\le X_{l+1}/2$. By
Lemma~\ref{s2:lemLacunary}(i), $\sum_lX_l\le2X_R=2F(\lambda_{R-2})$ (if $R\le2$
the sum is empty). Moreover $\lambda_{R-2}\ge\lambda_R=\log_2d_R\ge\log_2\Dstar$,
by Lemma~\ref{s2:lemTower}(a) and because round $R$ has $d_R\ge\Dstar$ by
rule~\ref{s2:R0} of Definition~\ref{s2:defHBtp}; and $\log_2\Dstar\ge2^{256}\ge2^{25}$ by \Gc{1}(a) at $\mu=\log_2\log_2\Dstar$. Since $F$
is decreasing on $[2^{25},\infty)$, we get
\[
\sum_lX_l\ \le\ 2F(\log_2\Dstar)\ \le\ 2F(2^{25})\ \le\ 2\cdot2^{20}\cdot2^{-2255}\ <\ 2^{-1000}.
\]
Together, $\sum_l\dett_l\le3\epsA n+60n/\Dstar+2F(\log_2\Dstar)\,n$. The bound
on $\sum_l|V(Q_l)|$ follows from (iii) and $\Ex\Xp\le\epsX(\Dstar)n$
(Lemma~\ref{s7:lemEXprime}). The functions $\epsX$, $\epsA$ and
$F(\log_2\Dstar)$ depend on $\Dstar$ only, and $\thetaQ\le1/4$ is part of
condition~\ref{s1:condG4}.
\end{proof}

%% ----------------------------------------------------------------------
\subsection{One joint outcome, the cost, and Theorem \texorpdfstring{JV$^{+*}$}{JV+*}}

\begin{lemma}[One joint outcome exists]\label{s7:lemOneOutcome}
There is a joint outcome, consisting of a stage-1 outcome, stage-3 outcomes for
all parts, and draws $\xi_R,\dots,\xi_3$, such that:
\begin{enumerate}
\item[(1)] $\Xp\le3\,\Ex\Xp$;
\item[(2)] the stage-3 labels of every part lie in that part's good event: the
event of Lemma~\ref{s4:lemTPV} for every standalone pre-part (with the data of
step (1) of Construction~\ref{s6:consOrder}); the event of Lemma~\ref{s4:lemPV}
as used in Lemma~\ref{s5:lemChild} for every non-demoted light part; and the event
of Theorem~\ref{s4:thmVXp} as used in Lemma~\ref{s5:lemDemoted} for every
demoted light part;
\item[(3)] at every round $3\le l\le R$, $\xi_l$ lies in the event of
Construction~\ref{s7:consRound}(f). Consequently
$\copies_l\le4(X_{\mathrm{V},l}+\dett_l)$ and
$\payrd_l\le4\bigl(n/(256M_l^3)+5.5\,nM_l/\Hcd_l\bigr)$.
\end{enumerate}
Such an outcome is obtained as follows.
(i) Choose the stage-1 outcome in $\{\Xp\le3\Ex\Xp\}$; stage~2 is then
determined.
(ii) For every part, choose its stage-3 labels in its good event.
(iii) For $l=R,R-1,\dots,3$ in turn: run steps (1)--(3) of
Construction~\ref{s6:consOrder} at round $l$, which are deterministic given the
past; draw $\xi_l$ and choose it in the event of
Construction~\ref{s7:consRound}(f); build $Q_l$, fix $\Dec_l$, lift it, and record
$\LentJV_l$ (step (4) of Construction~\ref{s6:consOrder}).
(iv) At rounds $2$ and $1$, run steps (1)--(2) of Construction~\ref{s6:consOrder};
here $J_2=J_1=\emptyset$.
\deps{\ref{s7:defSchedule}, \ref{s7:consRound}, \ref{s6:consOrder}, \ref{s4:lemTPV}, \ref{s4:lemPV}, \ref{s4:thmVXp}, \ref{s5:lemChild}, \ref{s5:lemDemoted}, \ref{s7:lemEXprime}, \ref{s7:lemPay}, \ref{s7:lemUHsplit}, \ref{s1:citMarkov}, \ref{s1:condG3}}
\end{lemma}

\begin{proof}
\emph{Step (i).} $\Xp$ is a non-negative random variable on a finite probability
space. By Markov's inequality (Cited result~\ref{s1:citMarkov}),
$\Prob(\Xp\le3\Ex\Xp)\ge2/3>0$, so the event is non-empty. Fix an outcome in it;
this gives (1). Stage~2 is a deterministic function of stage~1
(Definition~\ref{s7:defSchedule}).

\emph{Step (ii).} With stage~1 fixed, the good event of each part is determined
by stage-1 and stage-2 data together with that part's own stage-3 labels.
\begin{itemize}
\item For a standalone pre-part $Z$, Lemma~\ref{s4:lemTPV} is applied with $O_Z$
and $\vP=\Ret_Z$, both fixed at stage~2, and its good event is determined by
$(Z^0,O_Z,\Ret_Z)$ and the labels.
\item For a non-demoted light part, the good event of Lemma~\ref{s4:lemPV}
depends only on the own classes, $\Rt$, $\extph$ and the labels
(Lemma~\ref{s5:lemChild}).
\item For a demoted light part, the good event of Theorem~\ref{s4:thmVXp}
depends only on $O=X_Y$ and the labels (Lemma~\ref{s5:lemDemoted}).
\end{itemize}
Each of these events has probability at least $1/2-o(1)$, and by
condition~\ref{s1:condG3} every such $o(1)$ is at most $1/100$. So each event
is non-empty, whatever the stage-1 outcome; only this non-emptiness is used. The events of distinct parts involve disjoint families of labels,
so choosing each part's labels inside its own event gives one stage-3 outcome in
all of the events at once; no union bound is needed. None of these events
depends on the edge sets that the devices later receive: the conclusions of
Lemma~\ref{s4:lemTPV}, Lemma~\ref{s4:lemPV} and Theorem~\ref{s4:thmVXp} hold on
their good events for \emph{every} admissible edge set (every $E\supseteq E(O)$,
every $H_0$ containing the own classes, every graph containing $O$). So the
events do not depend on junk, on $\Erem$, on $H_0$ or on later rounds. This
gives (2).

\emph{Step (iii).} At round $l$ the past $\Past_l$ is fixed, and steps (1)--(3)
of Construction~\ref{s6:consOrder} are deterministic given it; in particular
$J_l$ is determined. The draw $\xi_l$ is independent of the past, and
$\copies_l$ and $\payrd_l$ are non-negative functions of $\xi_l$ with the past
fixed. By Markov's inequality (Cited result~\ref{s1:citMarkov}), the event of
Construction~\ref{s7:consRound}(f) has probability at least $1/2$. This holds for
every past and needs nothing beyond Markov's inequality. So a choice of $\xi_l$
in the event exists, whatever happened before. On this event, the explicit
bounds of (3) follow from Lemma~\ref{s7:lemUHsplit}(ii) and
Lemma~\ref{s7:lemPay}(c), which hold for \emph{every} past; this is the only
place where the uniformity in the past is used.

\emph{Step (iv).} Rounds $1$ and $2$ have no classed ports, since every port
there is fresh. So JS-LC produces no J-set, and there is no round step.
\end{proof}

\begin{proposition}[Cost on the chosen outcome]\label{s7:propCost}
On an outcome as in Lemma~\ref{s7:lemOneOutcome}, with the round step of
Construction~\ref{s7:consRound} as the J-consumer, the decomposition of $E(G)$
given by Theorem~\ref{s6:thmMIXC} has at most
\[
\Bigl(\frac{\Dstar}{2}+1085+\epsOne(\Dstar)\Bigr)n+2\sum_{l=3}^{R}\f(Q_l)
\]
objects, where
\[
\epsOne(\Dstar)\defeq\epsK(\Dstar)+169\,\epsA+\frac{252}{\Dstar}+1.5\,\epsCONC(\Dstar)+3\,\epsX(\Dstar)+\frac{9}{\Dstar},
\]
with $\epsA=31\eps/(\Cp\log\log\Dstar)$ as in Lemma~\ref{s2:lemTower}(e). The
terms are itemized in Table~\ref{s7:tabCost}. The constant is $745+338+2=1085$, and
$\epsOne$ is counted once. Consequently
$\Czero=\Dstar/2+1085+\epsOne(\Dstar)\le\Dstar/2+1091$ under \Gc{4}. No term
involves $\cEG$. The only Markov multipliers are the fixed $3$ (stage 1) and $4$
(per round). There is no term for VX-parts ($\mathcal{V}=\emptyset$).
\deps{\ref{s6:thmMIXC}, \ref{s6:thmCONCL}, \ref{s7:lemLift}, \ref{s7:lemPay}, \ref{s7:lemEXprime}, \ref{s7:lemOneOutcome}, \ref{s5:lemKRED}, \ref{s2:propParentless}, \ref{s2:propOV}, \ref{s7:consRound}, \ref{s2:lemTower}, \ref{s1:condG4}, \ref{s6:defLending}, \ref{s7:defXprime}}
\end{proposition}

\begin{table}[ht]
\centering
\small
\begin{tabular}{|p{5.0cm}|p{6.2cm}|p{3.6cm}|}
\hline
\textbf{Item} & \textbf{Bound on the number of objects} & \textbf{Source}\\
\hline
Long cycles, $E_0$ and light parts (K-RED) & $(\Dstar/2+745+\epsK)n+[80\,\mathrm{dem}+369\,\mathrm{lp}]$ & Lemma~\ref{s5:lemKRED}\\
TPV in standalone pre-parts & $169(2+\epsA)n+[169\sum_l|\Lost_l|]$ & Thm.~\ref{s6:thmMIXC}(c); Props.~\ref{s2:propParentless}, \ref{s2:propOV}\\
JS-LC cycles & $252\,n/\Dstar+1.5\,\epsCONC(\Dstar)\,n$ & Thms.~\ref{s6:thmMIXC}(c), \ref{s6:thmCONCL}\\
$\Jlost$ items & $[\sum_l(M_l-1)|\Lost_l|]$ & Lemma~\ref{s7:lemPay}(a1)\\
Pooled and JV-bad items & $[\Xpool]$ & Lemma~\ref{s7:lemPay}(a2),(a3)\\
Unpaired fresh legs & $2n$ & Lemma~\ref{s7:lemPay}(b)\\
SDR failures and loops & $9n/\Dstar$ & Lemma~\ref{s7:lemPay}(d)\\
Lift of the quotients & $2\sum_l|\Dec_l|=2\sum_l\f(Q_l)$ & Lemma~\ref{s7:lemLift}\\
\hline
All bracketed terms together & $\le\XU+\Xpool\le\Xp\le3\,\Ex\Xp\le3\,\epsX(\Dstar)\,n$ & Lemmas~\ref{s7:lemEXprime}, \ref{s7:lemOneOutcome}\\
\hline
\end{tabular}
\caption{The cost line on the chosen outcome (Proposition~\ref{s7:propCost}).
Summing the rows gives $(\Dstar/2+745+338+2)n+\epsOne(\Dstar)n+2\sum_l\f(Q_l)$.}
\label{s7:tabCost}
\end{table}

\begin{proof}
By Lemma~\ref{s7:lemLift}(iv), the round step is a J-consumer. The outcome of
Lemma~\ref{s7:lemOneOutcome} has stage-3 labels in all good events. So
Theorem~\ref{s6:thmMIXC} applies to the run, $\delta$, the chosen stage-1 and
stage-3 outcomes and this J-consumer, and it yields a decomposition of $E(G)$. By
Theorem~\ref{s6:thmMIXC}(c), its number of objects is at most
\[
\Bigl(\frac{\Dstar}2+745+338\Bigr)n+\eps_M(\Dstar)\,n
+\Bigl[80\,\mathrm{dem}+369\,\mathrm{lp}+169\sum_l|\Lost_l|\Bigr]
+1.5\sum_{(Y,l)}m_{Y,l}+\sum_l|\Obj_l| ,
\]
with $\eps_M(\Dstar)=\epsK(\Dstar)+169\epsA+252/\Dstar$. This itemization
combines the K-RED bound of Lemma~\ref{s5:lemKRED}; the TPV bound
\[
169\sum_{l,Z}\bigl(|A_Z|+|F_Z|+|\Lost_Z|\bigr),
\]
with $\sum|F_Z|\le2n$ (Proposition~\ref{s2:propParentless}(i)) and
$\sum|A_Z|\le\epsA n$ (Proposition~\ref{s2:propOV}(K2) and
Lemma~\ref{s2:lemTower}(e)); and the JS-LC bound. By
Theorem~\ref{s6:thmCONCL}(iv),
$1.5\sum_{(Y,l)}m_{Y,l}\le1.5\,\epsCONC(\Dstar)\,n$.

The objects of $\Obj_l$ are the paid single edges and the lifted objects. By
Lemma~\ref{s7:lemPay}, the paid single edges number at most:
$(M_l-1)|\Lost_l|$ in (a1); the round-$l$ summand of $\Xpool$ in (a2) and (a3);
at most $2n$ unpaired fresh legs over all rounds; and $\payrd_l$ SDR-failure and
loop edges, with $\sum_l\payrd_l\le9n/\Dstar$ by Lemma~\ref{s7:lemPay}(d), since
the outcome satisfies Lemma~\ref{s7:lemOneOutcome}(3). By
Lemma~\ref{s7:lemLift}(iv), the lifted objects number at most $2\f(Q_l)$. Hence
\[
\sum_l|\Obj_l|\le\sum_l(M_l-1)|\Lost_l|+\Xpool+2n+\frac{9n}{\Dstar}+2\sum_l\f(Q_l).
\]
The bracketed terms and the terms $\sum_l(M_l-1)|\Lost_l|$ and $\Xpool$ together
equal
$80\,\mathrm{dem}+369\,\mathrm{lp}+\sum_l(169+M_l-1)|\Lost_l|+\Xpool\le\XU+\Xpool\le\Xp$.
By Lemma~\ref{s7:lemOneOutcome}(1) and Lemma~\ref{s7:lemEXprime},
$\Xp\le3\Ex\Xp\le3\epsX(\Dstar)n$. Adding everything gives at most
\[
\Bigl(\frac{\Dstar}2+745+338+2\Bigr)n
+\Bigl(\epsK+169\epsA+\frac{252}{\Dstar}+1.5\epsCONC+3\epsX+\frac9{\Dstar}\Bigr)n
+2\sum_l\f(Q_l),
\]
which is the claim, since $745+338+2=1085$. Every term of $\epsOne$ is a function
of $\Dstar$ alone, and $\epsOne(\Dstar)\le6$ is part
of \Gc{4}, which gives $\Czero\le\Dstar/2+1091$.
\end{proof}

\begin{theorem}[Theorem JV$^{+*}$]\label{s7:thmJVps}
Assume that $\Dstar$ satisfies \Gc{1}--\Gc{4}. Let $G$ be a graph with
$n\ge\Nzero$ vertices and $d_1\ge\Dstar$. Take any valid $\HBtp$ run on $G$
(Proposition~\ref{s2:propExists}), any designation $\delta$, and no VX-parts
($\mathcal{V}=\emptyset$). Then there are simple graphs $Q_3,\dots,Q_R$ without
isolated vertices such that
\begin{enumerate}
\item[(1)] $\f(G)\le\Czero n+2\sum_{l=3}^{R}\f(Q_l)$, where
$\Czero\defeq\Dstar/2+1085+\epsOne(\Dstar)$;
\item[(2)] $\sum_{l=3}^{R}|V(Q_l)|\le\thetaQ\,n$, where
$\thetaQ=\epsTwo(\Dstar)\le1/4$.
\end{enumerate}
Here $\epsOne$ and $\epsTwo$ are the explicit functions of $\Dstar$ defined in
Proposition~\ref{s7:propCost} and Lemma~\ref{s7:lemUHsplit}. Items at ultra hubs
cost no collision payments. They pay only the standard costs (pooled vertices,
JV-bad ports, SDR failures), and they add quotient vertices only through
Lemma~\ref{s7:lemUHsplit}.
\deps{\ref{s7:propCost}, \ref{s7:lemUHsplit}, \ref{s7:lemSimple}, \ref{s7:lemLift}, \ref{s7:lemOneOutcome}, \ref{s6:thmMIXC}, \ref{s2:propExists}, \ref{s7:lemPay}, \ref{s7:lemUltra}, \ref{s1:condGamma}}
\end{theorem}

\begin{proof}
Take the outcome of Lemma~\ref{s7:lemOneOutcome} and let $Q_3,\dots,Q_R$ be the
quotients built on it. By Lemma~\ref{s7:lemSimple} each $Q_l$ is simple and has
no isolated vertices. By Lemma~\ref{s7:lemLift}(iv), the round step is a
J-consumer, so Theorem~\ref{s6:thmMIXC} gives a decomposition of $E(G)$. By
Proposition~\ref{s7:propCost} it has at most $\Czero n+2\sum_l\f(Q_l)$ objects,
which proves (1), since $\f(G)$ is at most the number of objects of any
decomposition. The outcome satisfies the hypotheses of
Lemma~\ref{s7:lemUHsplit}(iii): $\Xp\le3\Ex\Xp$, and the event of (f) holds at
every round. So Lemma~\ref{s7:lemUHsplit}(iv) gives (2). The statement about
ultra hubs is Lemma~\ref{s7:lemPay}(e) together with Lemma~\ref{s7:lemUltra}.
\end{proof}

%% ----------------------------------------------------------------------
\subsection{The layered quotient induction and the proof of the Main Theorem}

\begin{theorem}[Theorem HI$''$ (layered quotient induction)]\label{s7:thmHI}
Let $C\ge\Dstar/2$ and $\vartheta\in[0,1/2)$ be constants with the
following property. Every graph $G$ without isolated vertices, with $n\ge\Nzero$
vertices and $d_1=2|E(G)|/n\ge\Dstar$, admits finitely many simple graphs
$Q_1,\dots,Q_k$ ($k\ge0$) such that
\[
\f(G)\le C\,n+2\sum_{i=1}^{k}\f(Q_i)\qquad\text{and}\qquad\sum_{i=1}^{k}|V(Q_i)|\le\vartheta\,n .
\]
Then $\f(G)\le c\,|V(G)|$ for every graph $G$, where
$c\defeq\max\bigl(C/(1-2\vartheta),\,\Nzero/2\bigr)$.
\deps{\ref{s1:factAdd}, \ref{s1:defObject}, \ref{s1:defConstants}}
\end{theorem}

\begin{proof}
Suppose the conclusion fails, and let $G$ be a counterexample with
$n\defeq|V(G)|$ minimal, so that $\f(G)>cn$. The graph without vertices has
$\f=0$, so $n\ge1$.

(1) \emph{$G$ has no isolated vertex.} If $v$ were isolated, then
$\f(G-v)=\f(G)$ by Fact~\ref{s1:factAdd}(d), so
$\f(G-v)>cn\ge c(n-1)=c|V(G-v)|$, and $G-v$ would be a smaller counterexample.

(2) \emph{$n\ge\Nzero$.} Otherwise $n-1<\Nzero$, and by Fact~\ref{s1:factAdd}(a),
$\f(G)\le|E(G)|\le n(n-1)/2<(\Nzero/2)\,n\le cn$.

(3) \emph{$d_1\ge\Dstar$.} Otherwise, by Fact~\ref{s1:factAdd}(a),
$\f(G)\le|E(G)|=d_1n/2<(\Dstar/2)\,n\le Cn\le cn$, because
$c\ge C/(1-2\vartheta)\ge C$. This is the trivial bound $\f(G)\le|E(G)|$; no
other estimate is used for graphs of small average degree.

(4) By (1)--(3), the hypothesis applies to $G$ and gives simple graphs
$Q_1,\dots,Q_k$ with $\f(G)\le Cn+2\sum_i\f(Q_i)$ and
$\sum_i|V(Q_i)|\le\vartheta n<n$. Each $Q_i$ has fewer vertices than $G$, so by
the minimality of $n$, $\f(Q_i)\le c|V(Q_i)|$.

(5) Hence $\f(G)\le Cn+2c\sum_i|V(Q_i)|\le Cn+2c\vartheta n
\le c(1-2\vartheta)n+2c\vartheta n=cn$, where we used $C\le c(1-2\vartheta)$.
This contradicts $\f(G)>cn$.

The constants $C$ and $\vartheta$ do not depend on $c$, and minimality is applied
only to the graphs $Q_i$, which have strictly fewer vertices than $G$.
\end{proof}

\begin{lemma}[The galactic conditions are satisfiable]\label{s7:lemGammaSat}
Let $\Nzero$ be as in Definition~\ref{s1:defConstants}(ii). Then:
\begin{enumerate}
\item[(i)] each item (a)--(f) of \Gc{1} (Condition~\ref{s1:condG1}) holds for
all sufficiently large real $\mu$; hence there is $\mu_1\ge1$ such that all of them
hold for every $\mu\ge\mu_1$, and \Gc{1} holds for every $\Dstar>2$ with
$\log_2\log_2\Dstar\ge\mu_1$;
\item[(ii)] $\epsOne(\Dstar)\to0$ and $\epsTwo(\Dstar)\to0$ as $\Dstar\to\infty$;
\item[(iii)] there is $D_0$ such that every $\Dstar\ge D_0$ satisfies
\Gc{1}--\Gc{4} simultaneously.
\end{enumerate}
In particular a constant $\Dstar$ as in Definition~\ref{s1:defConstants}(iii)
exists. Only this existence is used: neither $\mu_1$ nor $D_0$ is computed.
\deps{\ref{s1:condGamma}, \ref{s7:propCost}, \ref{s7:lemUHsplit}, \ref{s5:lemExpect}, \ref{s5:lemKRED}, \ref{s2:lemTower}, \ref{s7:lemEXprime}, \ref{s6:thmCONCL}, \ref{s1:defConstants}, \ref{s3:lemCOL}, \ref{s5:lemE1}, \ref{s3:lemCOLJVev}}
\end{lemma}

\begin{proof}
(i) We use inequalities of type (E) as defined in Lemma~\ref{s3:lemCOLJVev},
which has no hypothesis: by part~(i) of that lemma, every inequality of type (E)
holds for all sufficiently large $\mu$. For $\mu\ge1$ we have
$0\le\log_2\mu\le\log_2(\Aexp\mu)$, and:
item~(a) holds for every $\mu\ge2^8$;
item~(b), that is $\mu-14-\log_2\Aexp-3\log_2\mu\ge0$, follows from the
inequality $\mu-(14+\log_2\Aexp)-3\log_2(\Aexp\mu)\ge0$ of type (E);
item~(c) is the inequality $1.6\mu-2\Aexp\log_2(\Aexp\mu)\ge0$ of type (E);
item~(d), that is $\mu-8-2\log_2\mu\ge0$, follows from the inequality
$\mu-8-2\log_2(\Aexp\mu)\ge0$ of type (E);
item~(e), which is \Gstar, is (after taking $\log_2$ of both sides) the
inequality $36\mu-240-46\Aexp\log_2(\Aexp\mu)\ge0$ of type (E);
and every inequality of item~(f) holds for all sufficiently large $\mu$ by
Lemma~\ref{s3:lemCOLJVev}(iii).
So each of the finitely many inequalities of \Gc{1} holds for all $\mu$ at least
some threshold; let $\mu_1\ge1$ be at least all of these thresholds. If
$\Dstar>2$ and $\log_2\log_2\Dstar\ge\mu_1$, every $\mu\ge\log_2\log_2\Dstar$
satisfies $\mu\ge\mu_1$, so \Gc{1} holds.

(ii) By Proposition~\ref{s7:propCost} and Lemmas~\ref{s7:lemEXprime}
and~\ref{s7:lemUHsplit}, $\epsOne$ and $\epsTwo$ are finite sums of non-negative
explicit functions of $\Dstar$ with fixed coefficients:
$1/\Dstar$; $(6\log_2\Dstar+12)/\Dstar$;
$\epsA=31\eps/(\Cp\log\log\Dstar)$ (Lemma~\ref{s2:lemTower}(e));
$\epsCONC(\Dstar)$; $\epsU(\Dstar)$; $\epsK(\Dstar)$; and, in $\epsTwo$ only,
$F(\log_2\Dstar)$ with $F(x)=(3184+30400\log_2x)\,x^{-90.2}$ as in
Lemma~\ref{s7:lemUHsplit}(iv), which tends to $0$ because $F(x)\to0$ as
$x\to\infty$. The function $\epsCONC$ is the explicit expression
of Theorem~\ref{s6:thmCONCL}(iv). It is a combination of
$\logs\Dstar/\log\Dstar$, $\logs\Dstar/\log\log\Dstar$ and
$(2\logs\Dstar+2)/(\log\Dstar)^{1/2}$, each of which tends to $0$. The functions
$\epsU$ (Lemma~\ref{s5:lemExpect}) and $\epsK$ (Lemma~\ref{s5:lemKRED}) are, by
their defining statements, explicit functions of $\Dstar$ that tend to $0$ as
$\Dstar\to\infty$. Every other listed function tends to $0$. Hence
$\epsOne(\Dstar)\to0$ and $\epsTwo(\Dstar)\to0$.

(iii) Put $D_1\defeq2^{2^{\mu_1}}$; then $D_1\ge4>2$, as $\mu_1\ge1$. By (i),
\Gc{1} holds for every $\Dstar\ge D_1$, since then $\log_2\log_2\Dstar\ge\mu_1$.
\Gc{2}(a) holds for every $\Dstar\ge2^{117}$, and \Gc{2}(b),(c) are implied by
\Gc{1} (Condition~\ref{s1:condGamma}). \Gc{3} holds for every
$\Dstar\ge2^{(2\Nzero)^{1/\Cp}}$; here $\Nzero$ is fixed before $\Dstar$. By (ii)
there is $D''$ with $\epsOne(\Dstar)\le6$ and $\epsTwo(\Dstar)\le1/4$ for all
$\Dstar\ge D''$, which is \Gc{4}. Take
$D_0\defeq\max\bigl(D_1,2^{117},2^{(2\Nzero)^{1/\Cp}},D''\bigr)$.
\end{proof}

\begin{corollary}[Corollary JV$^{+*}$-EG; proof of the Main Theorem]\label{s7:thmMainProof}
Theorem~\ref{s1:thmMain} holds.
\deps{\ref{s7:thmHI}, \ref{s7:thmJVps}, \ref{s2:propExists}, \ref{s7:lemGammaSat}, \ref{s6:defDesign}, \ref{s1:condGamma}, \ref{s1:defConstants}}
\end{corollary}

\begin{proof}
Let $\Dstar$ be as in Theorem~\ref{s1:thmMain}, that is, a constant satisfying
\Gc{1}--\Gc{4}. Such constants exist by Lemma~\ref{s7:lemGammaSat}, and the
argument below applies to each of them. Let
$\Czero=\Dstar/2+1085+\epsOne(\Dstar)$, $\thetaQ=\epsTwo(\Dstar)$ and
$\cEG=\max\bigl(\Czero/(1-2\thetaQ),\Nzero/2\bigr)$, as in
Theorem~\ref{s1:thmMain}. We apply Theorem~\ref{s7:thmHI} with $C\defeq\Czero$
and $\vartheta\defeq\thetaQ$. Its assumptions on the constants hold:
$\Czero\ge\Dstar/2$, since $\epsOne\ge0$; and $0\le\thetaQ\le1/4<1/2$.

Let $G$ be a graph without isolated vertices, with $n\ge\Nzero$ vertices and
$d_1\ge\Dstar$. By Proposition~\ref{s2:propExists}, $G$ has a valid $\HBtp$
run. A designation exists: for every standalone pre-part $Z$ of a round $l\ge3$
and every $u\in Q_Z$, the set $\anc_l(u)$ is non-empty by the definition of
classed ports, and $Y(u)$ can be chosen in it by a fixed rule
(Definition~\ref{s6:defDesign}). With no VX-parts, Theorem~\ref{s7:thmJVps} gives
simple graphs $Q_3,\dots,Q_R$ with $\f(G)\le\Czero n+2\sum_l\f(Q_l)$ and
$\sum_l|V(Q_l)|\le\thetaQ n$. This is the hypothesis of
Theorem~\ref{s7:thmHI}, which therefore gives $\f(G)\le\cEG|V(G)|$ for every
graph $G$. In particular $\f(n)\le\cEG n$ for every $n$, so $\f(n)=O(n)$, where
$\cEG$ depends only on $\sigma,\Cp,\Aexp,\eps,\Nzero$ and $\Dstar$.

This completes the proof of Theorem~\ref{s1:thmMain}.
\end{proof}

\begin{remark}[Why the argument is not circular]\label{s7:remNonCirc}
(1) $\Czero$ and $\thetaQ$ are fixed functions of $\Dstar$
(Proposition~\ref{s7:propCost}, Lemma~\ref{s7:lemUHsplit}), and $\Dstar$ is fixed
before any graph is considered. The constant $\cEG$ enters only as the
constant $c$ of Theorem~\ref{s7:thmHI} (throughout its statement and
proof); it plays no other role in any construction or bound of
Sections~\ref{s2:sec}--\ref{s7:sec}.
(2) Minimality is applied only to the quotients $Q_l$. They are simple and have
at most $\thetaQ n\le n/4<n$ vertices in total, and the induction uses about
them only the inequality $\f(Q_l)\le\cEG|V(Q_l)|$.
(3) The quotient $Q_l$ depends on the decompositions $\Dec_{l'}$ ($l'>l$), through
$J_l$ and through junk. This is harmless. Each $\Dec_{l'}$ is fixed at round $l'$
by a deterministic rule. The Markov choice at round $l$ succeeds with
probability at least $1/2$ whatever the past. Every bound of
Lemmas~\ref{s7:lemCC}, \ref{s7:lemUltra} and~\ref{s7:lemPay} holds for every
past.
(4) No other link calls itself. Lemma~TPV, Lemma~PV and Theorem~VX$^+$
finish their runs with the elementary long-cycle bound of
Fact~\ref{s1:factEG0}, which is proved directly and does not use the statement
being proved; all three are unconditional.
\deps{\ref{s7:thmHI}, \ref{s7:thmJVps}, \ref{s7:propCost}, \ref{s7:lemUHsplit}, \ref{s7:lemCC}, \ref{s7:lemUltra}, \ref{s7:lemPay}, \ref{s7:lemSimple}, \ref{s7:consRound}, \ref{s4:lemTPV}, \ref{s4:lemPV}, \ref{s4:thmVXp}, \ref{s1:condG4}, \ref{s1:factEG0}}
\end{remark}
.
%% Uses only the preamble of ms.tex; defines no macro, environment or
%% counter. Labels use the prefix s8:.
%% ======================================================================
\section{Formal verification}\label{s8:sec}
\begingroup\sloppy

Theorem~\ref{s1:thmMain}, in the form $\f(n)=O(n)$, has been formalized and
proved in the Lean~4 proof assistant \cite{Lean4} with its mathematical library
Mathlib \cite{Mathlib}. The formal target is not a statement of our own: it is
the theorem \texttt{Erdos184.\allowbreak erdos\_184} of the formal-conjectures
repository \cite{FC} (file
\nolinkurl{FormalConjectures/ErdosProblems/184.lean}, commit
\texttt{2424bb48}). Up to typesetting it reads:
\begin{quote}\small\ttfamily
open scoped Classical in\\
theorem erdos\_184 :\\
\hspace*{1em}$\exists$ f : $\mathbb N\to\mathbb R$,\\
\hspace*{2em}(f =O[atTop] fun n : $\mathbb N$ $\mapsto$ (n : $\mathbb R$)) $\wedge$\\
\hspace*{2em}$\forall$ \symbol{123}V : Type*\symbol{125} [Fintype V] [DecidableEq V] (G : SimpleGraph V),\\
\hspace*{2em}$\exists$ (D : Finset G.Subgraph),\\
\hspace*{3em}($\forall$ H $\in$ D, IsCycleOrEdge H.coe) $\wedge$\\
\hspace*{3em}IsDecomposition G D $\wedge$\\
\hspace*{3em}(D.card : $\mathbb R$) $\le$ f (Fintype.card V)
\end{quote}
(in namespace \texttt{Erdos184}, under \texttt{open Filter SimpleGraph}).
Here \texttt{IsCycleOrEdge H} means that the finite simple graph \texttt{H} is
connected and $2$-regular, or has exactly one edge; \texttt{IsDecomposition G D}
means that the edge sets of the members of \texttt{D} are pairwise disjoint with
union $E(G)$. A connected $2$-regular finite simple graph is a cycle of length at
least~$3$, so the statement says that some $f(n)=O(n)$ bounds, for every finite
simple graph on $n$ vertices, the size of a partition of its edge set into cycles
and single edges: this is the ``in particular'' part of
Theorem~\ref{s1:thmMain} (take $f(n)=\cEG n$).

The Lean proof of \texttt{Erdos184.\allowbreak erdos\_184} depends only on the
standard axioms \texttt{propext}, \texttt{Classical.choice} and
\texttt{Quot.sound}. No published result is assumed: the cited inputs of
Section~\ref{s1:sec} that the proof uses are proved in Lean, except Hall's
theorem (Cited result~\ref{s1:citHall}), which is taken from Mathlib. The proof
has been checked by the Lean kernel and by the independent checker
\emph{comparator} \cite{Comparator} in release mode, which verifies that the
proved theorem has exactly the statement of formal-conjectures, that no axioms
other than the three above are used, and that the proof is accepted on a kernel
replay. The code is available at
\url{https://github.com/steelwheel01/erdos-gallai-lean} \cite{EGRepo} (release
tag \texttt{release-2026-10-01}, commit
\texttt{a11bb45847101e0d1595089bc390c59f82ad1135}; verification run
\texttt{36795612102} of its GitHub Actions release workflow). The formalization
follows this paper section by section; a table matching every numbered
statement of Sections~\ref{s1:sec}--\ref{s7:sec} with its Lean counterpart, where
there is one, is given in the file \nolinkurl{CORRESPONDENCE.md} of that repository.
\par\endgroup


\section*{Acknowledgements}
The proof in this paper and its Lean formalization were developed with
substantial assistance from the AI system Claude (Anthropic), under the
author's direction. The author takes full responsibility for the content of
this paper.

\bibliographystyle{amsplain}
\bibliography{refs}
\end{document}
